% =====================================================================
%  R. Nielsen, LOGIK OG PSYCHOLOGIE
%  Cursus for Universitetsaaret 1867--68.
%  Kjøbenhavn: Forlagt af den Gyldendalske Boghandel (F. Hegel).
%  Trykt hos J. H. Schultz. 1867.   [VI] + 482 pp.
%
%  Diplomatic transcription from the Google Books scan
%  (`bibliotek/Nielsen, Rasmus/Logik-1867.pdf`, 502 PDF pages).
%
%  This header is the edition's apparatus. The transcription harness
%  (pagemap.py, ocr.sh, check.py, splice.py, joints.py) is not tracked in
%  git and is unrunnable without the scan, so everything an outside reader
%  needs in order to check this text against the scan is recorded here.
%
% ---------------------------------------------------------------------
%  1. PAGE MAP  (verified 12 Sep 2026; see pagemap.py)
% ---------------------------------------------------------------------
%
%      printed 1--482   =   PDF page + 16      (PDF 17--498)
%
%  Uniform: no offset change anywhere in the book. The folio in the head of
%  every one of PDF 17--498 was read mechanically and compared against
%  printed = PDF - 16; 479 of 482 agree exactly and the three that do not
%  are accounted for:
%
%    printed 1   (PDF 17)  carries no folio -- opening page of the treatise.
%                          Fixed from below by the "2" on PDF 18.
%    printed 8   (PDF 24)  folio dropped by the ABBYY layer only; tesseract
%                          -l Fraktur reads "8" cleanly.
%    printed 238 (PDF 254) THE FOLIO IS MISPRINTED "233". See errata below.
%
%  Front matter (PDF pages; unnumbered in the print except the synopsis,
%  which carries its own 3--6):
%      PDF  1   Google digitisation notice -- not part of the book
%      PDF  9   title page
%      PDF 11   title page again (the scan images the leaf twice)
%      PDF 13--15  synopsis "Logik og Psychologie", printed pp. 3--5
%      PDF 16   Indhold, printed p. 6
%
% ---------------------------------------------------------------------
%  2. THE BOOK'S TWO CONFLICTING WITNESSES TO ITS OWN STRUCTURE
% ---------------------------------------------------------------------
%
%  The front matter and the body text do not agree, and this edition
%  follows the BODY, recording the divergences rather than normalising
%  them away. Three separate disagreements, all confirmed on the scan:
%
%  (a) The body opens (p. 1) with a three-deep heading that the book never
%      continues:
%              Første Deel.  /  Videns Idee.
%              Første Afsnit.  /  Den subjective Viden.
%              Første Kapitel.
%              Viden og den Vidende: Subjectivitetens Logik.
%      There is no "Anden Deel", no "Andet Afsnit" and no "Andet Kapitel"
%      anywhere in the 482 pages. The divisions that DO continue are the
%      I / II / III of the synopsis. So p. 1's heading apparatus is the
%      opening of a projected larger system of which this Cursus prints
%      only the first instalment. It is transcribed as printed.
%
%  (b) The Indhold (p. 6) calls the third part
%              "Psychologie og Biologie"
%      where the head over p. 277 reads
%              "Biologie og Psychologie".
%      Both readings confirmed. Neither is normalised to the other; the
%      disagreement is a datum.
%
%  (c) The Indhold's coverage STOPS AT p. 477. The closing section
%              III.  Philosophien og de akademiske Studier   (pp. 477--482)
%      is not listed in it at all, and is printed under a second "III."
%      (the first being the "III." over p. 277).
%
%  Division heads as PRINTED IN THE BODY, with the pages verified from the
%  scan (the Indhold's own page figures are given in brackets where they
%  differ, and differ only by rounding a division back to the first page of
%  its part):
%
%      Første Deel. Videns Idee.                                  p.   1
%        Første Afsnit. Den subjective Viden.                     p.   1
%          Første Kapitel.
%          Viden og den Vidende: Subjectivitetens Logik.          p.   1
%            [introductory, no head]                         pp.   1-- 11
%            A. Den umiddelbare Subjectivitet                pp.  12-- 53   [Indhold: 1--53]
%            B. Den reflecterende Subjectivitet               pp.  54--132
%            C. Den begribende Subjectivitet                  pp. 133--263
%      II. Psychologie.                                           p. 264
%            Indledende Sætninger                             pp. 264--265
%            Plantelivet                                      pp. 266--267
%            Dyrelivet                                        pp. 268--272
%            Menneskelivet: Psychologien i Grundtræk          pp. 273--276
%      III. Biologie og Psychologie.                              p. 277
%            a) Det organiske Formaal                         pp. 277--285
%            b) Livets Trin                                   pp. 286--422
%            c) Sjæl og Legeme                                pp. 423--476
%      III. Philosophien og de akademiske Studier.            pp. 477--482
%
%  Finer divisions are run-in in the print -- numbered paragraphs whose
%  number and rubric begin a paragraph in the text ("77. Mulige Fremskridt.
%  At angive Videns Grændser ...") and Greek-lettered subdivisions
%  (αα, ββ). They are transcribed run-in, as printed, and are NOT promoted
%  to sectioning commands.
%
% ---------------------------------------------------------------------
%  3. ERRATA -- printer's defects, transcribed as printed
% ---------------------------------------------------------------------
%
%  Per house practice the compositor is not silently corrected: each defect
%  stands in the text with a % comment at the site, and is listed here.
%
%    p. 238  The FOLIO is misprinted "233". The text is continuous with
%            p. 237 before it and p. 239 after it, so only the numeral is
%            wrong -- a 3 set for an 8. Both witnesses read 233.
%            The page is printed p. 238 and is marked as such.
%
%    p.  60  The first sub-head of division B is printed
%                    d) Forstandsreflexion.
%            where the series requires a). The glyph is unambiguous: it has
%            the ascender and curved top-stroke of Fraktur d, checked at
%            600 dpi against the plainly two-bowled b of "b) Væsensreflexion"
%            (p. 76) and the c of "c) Virkelighedsreflexion" (p. 105). So this
%            is not a doubtful reading but a compositor's error, and a well
%            evidenced one: division B's other two limbs are lettered b) and
%            c), which require an a) before them, and the synopsis on p. 3
%            gives B the limbs a) Forstandsreflexion, b) Væsens-,
%            c) Virkelighedsreflexion. The likeliest mechanism is that the
%            compositor carried the count on from division A, which had ended
%            at c) on p. 35, and then corrected himself for the rest of B.
%            TRANSCRIBED AS PRINTED (d), per house practice, with a comment
%            at the site.
%
%    Further wrong sorts and dropped points logged in situ by the batches, all
%    transcribed as printed: p. 33 "din Vider" for "Viden."; p. 34 missing
%    point after "Principet"; p. 61 "guddommmelig" (three m); p. 66 "of" for
%    "af"; p. 67 "a" for "af" and "om" for "som"; p. 68 "og" doubled across a
%    line break; p. 69 a footnote with no printed call mark; p. 73 „Ting an
%    sich“ against p. 74's „Ding an sich“; p. 78 "Erscheinigung"; p. 83
%    "Phænomerne" for "Phænomenerne"; p. 103 "Arogonit" for "Aragonit";
%    p. 108 "stode" for "støde"; p. 109 "Tunkle" for "Dunkle".
%
%    THE QUOTE-MARK LEDGER. check.py counts „ against “ over the whole file,
%    and the running balance IS NOT ZERO, because the compositor set four
%    unmatched marks. Each was verified on the image; none is to be "fixed" by
%    adding or deleting a mark the compositor did not set. The ledger:
%
%      p. 139   +1   A quotation opens „det Problem ... and is never closed
%                    separately: the single “ at p. 140 ("endnu ikke denne.“")
%                    serves both nestings.
%      p. 234   -1   The opening „ before "Alvidenheden" was dropped; the
%                    closing mark is set.
%      p. 287   -1   A closing “ after "Sind wir im Innern.“" with no opening
%                    partner anywhere.
%      p. 304   -1   Likewise in the German note carried over from p. 303: the
%                    closing high mark is set, the opening low mark is not.
%                                                        ---------
%                    running balance through p. 327:        -2
%
%            So a NONZERO balance from check.py is expected. What would be a
%            real signal is a CHANGE in it that no new entry here explains.
%
%  [further defects appended as the batches find them]
%
% ---------------------------------------------------------------------
%  4. TRANSCRIPTION CONVENTIONS
% ---------------------------------------------------------------------
%
%  * 1867 orthography exactly as printed: Subjectivitet, Videnskaben,
%    Psychologie, høiere, aa, Christendommen, capitalised nouns. Not
%    modernised, and not made self-consistent where the book is not.
%  * Danish quotation marks „…“ low-high, as printed. Em-dash ---.
%  * THE COMPOSITOR SETS PLAIN o FOR ø, INCONSISTENTLY, AND IT IS NOT AN OCR
%    ARTEFACT. Every instance below was checked on the image against a plainly
%    slashed ø elsewhere on the same page, and on pp. 270 and 273 the plain and
%    the slashed form stand in the same passage. All are transcribed as printed
%    and none is normalised:
%        hoist (pp. 1--11, 178, 180, 184 "hoieste"), Oiekast, Hoiere, selvlos,
%        gron (150), Omboining(er) (151, 155), høist WITH the slash (226),
%        nodvendig (270, beside nødvendig on the same page), Udodelighed (268,
%        273), allerforste/forste (269), Grundsporgsmaalet, Doken/Doden,
%        godtgjore (273), religios (276), villielos (127), Strommen (132),
%        stode (108).
%    Note that the book also prints hoist in the opening pages and høist on
%    p. 226, so the variation is not a drift but genuinely unsystematic.
%    check.py deliberately carries NO "dropped ø" test for this reason.
%  * Letterspaced emphasis (Sperrsatz) -> \emph{}. Latin, Greek and German
%    set in antiqua against the surrounding Fraktur -> \textit{}. Emphasis
%    is decided per occurrence against the image, never by word identity:
%    a name may be antiqua on one page and Fraktur on the next.
%  * Fraktur has ONE sort for I and J. It is transcribed I throughout
%    (Idee, Indhold, Individet), per house convention.
%  * The id-est mark is ɔ: (U+0254 followed by a colon), not "o:".
%  * Footnotes are marked *) in the print; see \thefootnote below.
%  * Herbart's formulae in the closing chapters are set as displayed or
%    in-line mathematics reproducing the printed arrangement.
%  * SECTIONING LADDER, fixed after the first five batches. The print's own
%    labels carry the hierarchy, so the LaTeX level does not have to:
%        \part*        I / II / III  and  Første Deel
%        \chapter*     Første Afsnit; a) b) c) of part III
%        \subsection*  A / B / C
%        \subsubsection*  every lettered and Greek-lettered display head
%                          inside A/B/C -- a) b) c) d) ... and αα) ββ) alike.
%    THE SUB-HEAD SERIES. Each of A, B, C is divided a) b) c), and each of
%    those into α) β) γ) -- a regular three-by-three, confirmed on the scan
%    and agreeing with the synopsis (p. 3):
%        A  a) Sensualistisk Subjectivitet            p.  24
%           b) Idealistisk Subjectivitet              p.  29
%           c) Subjectivitetsproblemet                p.  35
%        B  a) Forstandsreflexion                     p.  60  [PRINTED "d)"]
%             α) Sættende og forudsættende Reflexion  p.  62
%             β) Modsættende og eenhedsættende Refl.  p.  66
%             γ) Iboende og oversvævende Reflexion    p.  70
%           b) Væsensreflexion                        p.  76
%             α) Den phænomenale Reflexion            p.  77
%             β) Den grundende Reflexion              p.  85
%             γ) Den begrundende Reflexion            p.  92
%           c) Virkelighedsreflexion                  p. 105
%             α) Den substantielle Reflexion          p. 106
%             β) Den causale Reflexion                p. 114
%             γ) Den totaliserende Reflexion          p. 123
%        C  a) Begriben og Begreb: logisk Ideebevægelse   p. 143
%             α) Negationsbestemmelser: Begrebet i
%                logisk Form                          p. 146
%             β) Positionsbestemmelser: Begrebet i
%                antilogisk Form                      p. 155
%             γ) Begrebets Realitet: den fuldkomne
%                Begriben                             p. 166
%           b) Begreb og Idee: Begribningens Phaser   p. 177
%                [Indhold, p. 6: "Idealsystemet og Realsystemet"]
%             α) Videns Oprindelighed:
%                universalia ante rem                 p. 178
%             β) Magtens Oprindelighed:
%                universalia post rem                 p. 193
%             γ) Magtens og Videns oprindelige Eenhed:
%                universalia in re                    p. 207
%           c) Det logiske Ideal: Begriben og Viden   p. 230
%                [the synopsis, p. 5, titles this limb only by its matter]
%    Beneath the Greek letters the print uses RUN-IN numbered rubrics, set
%    letterspaced at the head of a paragraph and left there, not promoted:
%    e.g. "2) Realsystemet." (p. 220), "3) Ideal- og Realsystemets
%    teleologiske Eenhed." (p. 224), "7) Vidensidealet" (p. 253).
%    Greek-lettered heads are a level below the lettered ones; both are set
%    \subsubsection* and are told apart by their printed label.
%  * A THIRD EMPHASIS DEVICE, found on p. 59: fat-face (halbfette) Fraktur in
%    running text, distinct from letterspacing and from antiqua ->
%    \textbf{}. Calibrate against plain occurrences of the same word on the
%    same page before using it.
%  * Page markers: % --- p. N --- on its own line where printed page N
%    begins. check.py depends on them; every page carries one.
%  * A word the compositor broke across a line is rejoined with \- and a
%    %  to eat the newline, never left as "Cau- salitetsforholdet".
%
% ---------------------------------------------------------------------
%  5. THE OCR WITNESSES
% ---------------------------------------------------------------------
%
%  Three, because on this scan none of them dominates (see ocr.sh):
%    1. tesseract -l Fraktur at 300 ppi  (a 2x downsample of the scan)
%    2. tesseract -l Fraktur at 600 ppi  (native; the scan is 600 ppi
%       bitonal jbig2, one image per page, already trimmed to the type
%       block -- it needs no cropping)
%    3. the PDF's own ABBYY layer, which read much of the book as antiqua
%       and so fails in an unrelated way (k->t/f, sk->st, æ->a, ſ->f).
%  Where all three agree the reading is safe; where two differ, the image
%  decides. No reading is settled by a witness alone against the image.
% =====================================================================

\documentclass[11pt]{book}
\usepackage[utf8]{inputenc}
\usepackage[T1]{fontenc}
\usepackage{libertinus}
\usepackage{libertinust1math}
\usepackage{textalpha}
\usepackage{graphicx}
% The id-est mark ɔ: (U+0254 + colon). It first appears on p. 296 -- the book
% writes "d. e." and "d. v. s." everywhere before that -- and U+0254 has no T1
% slot, so without this mapping pdflatex halts on it.
\DeclareUnicodeCharacter{0254}{\reflectbox{c}}
\usepackage{amsmath}
\usepackage[danish]{babel}
\usepackage{fancyhdr}
\setlength{\headheight}{28pt}
\addtolength{\topmargin}{-13.5pt}
\pagestyle{fancy}
\fancyhf{}
\fancyhead[LE,RO]{\thepage}
\fancyhead[RE]{\textit{Logik og Psychologie}}
\fancyhead[LO]{\textit{\leftmark}}

% The print marks footnotes *), and a page that carries two marks the second **)
% (first seen on p. 118). NOT \fnsymbol (math mode, gives U+2217) and NOT
% footmisc[perpage], which resets on the TYPESET page, not the printed one. The
% counter is instead reset by hand on every printed page: a
% \setcounter{footnote}{0}% line follows each  % --- p. N ---  marker. The
% trailing % on that line is load-bearing -- a word the compositor broke across
% the page break is joined with \-% before the marker, and without the trailing
% % the \setcounter line's newline would reappear as a space inside the word.
\renewcommand{\thefootnote}{\ifcase\value{footnote}\or *)\or **)\or ***)\else
  \arabic{footnote})\fi}

\usepackage{setspace}
\setstretch{1.3}

\title{Logik og Psychologie\\[0.5em]
  \large Cursus for Universitetsaaret 1867--68}
\author{R.\ Nielsen}
\date{Kjøbenhavn: Forlagt af den Gyldendalske Boghandel (F.\ Hegel).\\
  Trykt hos J.\ H.\ Schultz. 1867}

\begin{document}
\maketitle
\tableofcontents

%% ============================================================
%% FRONT MATTER: the course synopsis, printed pp. 3--6
%% Its own numbering, and its headings disagree with the body's
%% (see §2 of the header). Transcribed from the image, not from the
%% contents-page OCR layer.
%% ============================================================

\frontmatter\pagestyle{empty}
\chapter*{[Synopsis]}

% --- syn. p. 3 ---
% FRONT-MATTER SYNOPSIS, printed pp. 3--6 (PDF 13--16). Its own pagination,
% independent of the body's 1--482, so outside pagemap.py's range; the four
% pages were rendered and read directly (300 and 600 dpi) and OCR'd by hand
% with the three witnesses.
% Printed p. 3 carries NO folio (opening page of the synopsis); pp. 4 and 5
% carry centred antiqua folios; p. 6 (Indhold) carries none.
% No footnote, no Greek and no ɔ: mark anywhere on pp. 3--6; the id-est
% formula used here is the spelled-out „d. e.“ (pp. 4, 5).
\section*{Logik og Psychologie.}
% Head as printed: centred fat-face Fraktur with a full stop, NOT
% letterspaced (checked at 600 dpi).

Som alle Videnskabers fælles Grundlag afgiver Logiken særegne
Støttepunkter for hver Videnskab især og da navnlig ogsaa for
Psychologien. Saaledes er Læren om Subjectivitetens Oprindelighed af
særegen Betydning, naar det gjælder om Indsigt i Livets Væsen. Læren om
Reflexionen tjener Psychologien deels som „Forstandsreflexion“
(S. 60--70) til Opklaring af Forstandslivet, deels som Væsens- og
Virkelighedsreflexion (S. 77--123) til nærmere Belysning af Forholdet
mellem Reflecteren og Virken i de organiske Processer; endelig afgiver
Læren om den begribende Subjectivitet uundværlige Bidrag til grundig
Opfattelse af Ideernes Psychologie og til dybere Forstaaelse af det
selvbevidste Livs absolute Gyldighed.
% „Forstandsreflexion“ is quoted, not letterspaced — the only emphasis
% decision on p. 3 that the image overturned against a first reading of the
% OCR. „Væsens- og Virkelighedsreflexion“ is a SUSPENDED compound (printed
% with the Fraktur double stroke „Væſens=“ at the line end) and keeps its
% hyphen and its space; it is not joined.
% The internal references print Fraktur „S.“ / „Side“ with antiqua figures.

\subsection*{I.\\ Subjectiv Logik i Omrids.}
% „I.“ centred antiqua; „Subjectiv Logik i Omrids.“ centred fat-face
% Fraktur, not letterspaced.

\subsubsection*{A.\quad Subjectivitetens Oprindelighed.}
% The letter „A.“ is antiqua, the rubric letterspaced Fraktur. This is
% STRUCTURAL letterspacing carried by position, as with the body's division
% heads, so it is not marked \emph{}. Note the full stop: the Indhold
% (p. 6) prints a colon in the same place.

a) At Subjectiviteten ikke kan afledes af nogetsomhelst udvortes eller
indvortes Andet, lære vi af de Modsigelser, hvori Sensualismen har gjort
sig skyldig (Side 24--29);\par\noindent
% Structure as printed on p. 3: only a) is indented; the b) and c) clauses
% each begin a FRESH LINE FLUSH LEFT, the preceding line being justified
% out to the full measure. On pp. 4--5 the lettered clauses are ordinary
% indented paragraphs. The lettered enumerators a) b) c) are antiqua
% throughout the synopsis.
b) at det menneskelige Selv ikke kan være sit eget Ophav, viser Jegets
Selvmodsigelse i den subjective Idealisme (S. 29--35);\par\noindent
c) Subjectivitetsproblemets Løsning leder til den Indsigt, at den
menneskelige Subjectivitet har sin Oprindelse fra den guddommelige.

% --- syn. p. 4 ---
\subsubsection*{B.\quad Den reflecterende Subjectivitet.}

Subjectiviteten bestaaer i en ideel Virken, hvorunder Selvheden
uophørlig speiler Andet i sig og sig i Andet, altsaa i en uendelig
Reflecteren: Subjectiviteten er væsentlig reflecterende.

a) Forsaavidt den reflecterende Subjectivitet sætter og forudsætter en
umiddelbar Adskillelse mellem Andethed og Selvhed, mellem de Gjenstande,
hvorover den reflecterer, og sine egne Reflexioner, charakteriserer
Reflexionen sig som \emph{Forstandsreflexion}.

b) Forsaavidt den forstandigt reflecterende Subjectivitet indseer, at de
Egenskaber og Relationer, den maa tillægge Tingene, ere Tingenes egne
Reflexionsbestemmelser, Bestemmelser, hvori Tingene reflectere hinanden
indbyrdes, er Reflexionen objectiv, og den objective Reflexion at
bestemme som \emph{Væsensreflexion}.
% Printed „Væſens=reflexion“ broken across a line end, BOTH halves
% letterspaced, so the emphasis is not split; the line-break hyphen is an
% artefact of the setting and the word is joined.

c) For Væsensreflexionen er der endnu en Rest af Umiddelbarhed, nemlig
den haarde Umiddelbarhed, der charakteriserer det objectivt Virkelige,
tilbage at overvinde. Ved at optage, gjennemtrænge og assimilere sig
denne Umiddelbarhed naaer den reflecterende Subjectivitet omsider det
Vendepunkt, da den udtrykkelig veed, at Reflecteren er en Virken og al
Virken væsentlig en Reflecteren: den fuldt udviklede Reflexion er
\emph{Virkelighedsreflexion}.

\subsubsection*{C.\quad Den begribende Subjectivitet.}

Naar Virkelighedsreflexionen har tilegnet sig de for alt udvortes
Virkeligt afgjørende Umiddelbarhedselemenier, har den med det Samme
% printer's defect, p. 4: „Umiddelbarhedselemenier“ stands as printed,
% where the sense requires „Umiddelbarhedselementer“ — an i sort set for a
% t. Settled at 600 dpi on two POSITIVE features of the glyph: a separated
% i-dot, and no ascender; calibrated against the certain t of „udvortes“
% and „Virkeligt“ on the same line, which show an ascender and no dot.
% Both tesseract witnesses read i; the ABBYY layer silently normalises to
% „elementer“. Not corrected here.
optaget en Modsætning, der tvinger den til at gaae ud over sig selv og
forvandle sin væsentlige Reflecteren til virkelig Begriben.

a) Fra Væsenssiden er Reflexionen \emph{logisk}; fra Virkelighedssiden
faaer den en Modstand i det Virkeliges punktuelle Umiddelbarhed (Dette
mod Dette) d. e. i \emph{det Antilogiske}. For at magte det Logiske og
det Antilogiske
% --- syn. p. 5 ---
i Begrebets Eenhed, maa Subjectiviteten virke som Selvhedseenhed af
Viden og Magt, d. e. som ontologisk, guddommelig Subjectivitet.
% The sentence runs across the p. 4/5 break, so no blank line above the
% marker. „Selvheds=eenhed“ is broken at the line end on p. 5 and joined.
% The emphasis on p. 4 begins at „det“, one word before the line end
% („i d e t  A n t i l o g i s k e.“), and the immediately following
% „det Logiske og det Antilogiske“ is NOT letterspaced — the print
% distinguishes the two occurrences and both stand as printed.

b) Forsaavidt den ontologiske Subjectivitet begriber Tilværelsens Indhold
i logiske, ved antilogiske Reflexer articulerede Former, have vi
\emph{Idealsystemet}, den Verden af virksomme Ideer, hvori al Virkelighed
oprindelig præformeres og anlægges. I denne Begriben er Magten
\emph{ideel}. Forsaavidt den altformaaende Subjectivitet begriber
Tilværelsens Indhold i udvortes Virkelighed altsaa i antilogiske, ved
logiske Former articulerede Kjendsgjerninger, have vi
\emph{Realsystemet}. I denne Begriben er Magten \emph{reel}.
% „præ=formeres“ broken at a line end and joined. I/J: „Idealsystemet“,
% „Ideer“, „Indhold“, „I denne“ take I; „Jegets“ (p. 3) takes J, per the
% house convention recorded in the header.

c) Ved den begrebne Modsætning af ideel og reel Magt indgaae de to
Verdener, den ideale og den reale, i en saa gjennemgaaende
Vexelvirkning, at den ene bliver Middel, naar den anden bliver Maal: den
fuldendte Begriben er \emph{teleologisk}.

\begin{center}\rule{0.10\textwidth}{0.4pt}\end{center}
% A short centred rule closes the synopsis below the last line of p. 5;
% the rest of p. 5 is blank.

% --- syn. p. 6 ---
\clearpage
\section*{Indhold.}
% Head as printed: centred fat-face LETTERSPACED Fraktur with a full stop,
% over a short centred rule (reproduced below). Transcribed from the image,
% not from the contents-page OCR layer: both tesseract witnesses read the
% head „Iudhold.“, the column head „Site./Sitr.“ for „Side.“, and „Sjal“
% for „Sjæl“ in the last entry.
\begin{center}\rule{0.10\textwidth}{0.4pt}\end{center}

{\setlength{\parindent}{0pt}

\begin{center}I.\\ {\bfseries Subjectiv Logik i Omrids.}\end{center}

A.\quad Subjectivitetens Oprindelighed:\hfill{\small Side.}\par\smallskip
% As printed: a COLON here, where the synopsis head on p. 3 has a full
% stop. The A. line carries no leaders and no figures; the column head
% „Side.“ stands at the right margin of this line only.
\hangindent=2.2em\hangafter=1\hspace*{1.1em}Den umiddelbare Subjectivitet\dotfill 1--53.\par\smallskip
B.\quad Den reflecterende Subjectivitet\dotfill 54--132.\par\smallskip
C.\quad Den begribende Subjectivitet\dotfill 133--263.\par\medskip
% The A./B./C. rubrics and the „Den umiddelbare Subjectivitet“ entry are
% letterspaced — structural letterspacing carried by position, as with the
% body's division heads, so not marked \emph{}. The part II entries below
% are NOT letterspaced; the part III entries are. Recorded once here
% rather than per entry.
% On the figures: the Indhold rounds A. back to p. 1, the first page of
% its part, where the body's head „A. Den umiddelbare Subjectivitet“
% stands on p. 12 (pp. 1--11 are the introductory matter of the Kapitel).
% This is the rounding already noted in the header, not a new divergence.

\begin{center}II.\\ {\bfseries Psychologie.}\end{center}

Indledende Sætninger\dotfill 264--266.\par\smallskip
Plantelivet \dotfill 266--268.\par\smallskip
Dyrelivet \dotfill 268--273.\par\smallskip
Menneskelivet: Psychologien i Grundtræk\dotfill 273--277.\par\medskip

\begin{center}III.\\ {\bfseries Psychologie og Biologie.}\end{center}
% DIVERGENCE 1 (deliberate, not normalised). The Indhold heads the third
% part
%       „Psychologie og Biologie.“
% where the head printed over body p. 277 reads
%       „Biologie og Psychologie.“
% Both readings confirmed at 600 dpi. Each stands as printed in its own
% place; the disagreement is a datum.

a) Det organiske Formaal\dotfill 277--286.\par\smallskip
b) Livets Trin\dotfill 286--423.\par\smallskip
% printer, p. 6: the closing parenthesis of „b)“ is imperfectly inked and
% reads like a brace — both tesseract witnesses read it as a closing
% brace. It is a closing-parenthesis sort, defectively printed, and is
% transcribed „b)“.
c) Sjæl og Legeme \dotfill 423--477.\par\medskip
% „Sjæl“, not „Sjal“: the æ ligature confirmed at 600 dpi against the
% certain „aa“ of „Formaal“ two lines above (the ligature is distinctly
% wider, with the open angled left bowl). Both tesseract witnesses read
% „Sjal“ — exactly the contents-page OCR trap the playbook warns of.
%
% DIVERGENCE 2 (deliberate, not normalised). The Indhold's coverage STOPS
% AT p. 477. The closing section
%       „Philosophien og de akademiske Studier“   (pp. 477--482),
% printed in the body under a SECOND „III.“, is not listed in the Indhold
% at all. Transcribed exactly as printed: the missing section is NOT
% supplied here.
}

\medskip
\begin{center}\rule{0.14\textwidth}{1.2pt}\end{center}
% A short, thick centred rule closes the Indhold; the rest of p. 6 is
% blank. Not transcribed on these pages: the „Digitized by Google“
% footer of the scan.
\mainmatter\pagestyle{fancy}

%% ============================================================
%% FØRSTE DEEL: VIDENS IDEE  (pp. 1--263)
%% ============================================================

\part*{Første Deel.\\ Videns Idee.}
\addcontentsline{toc}{part}{Første Deel. Videns Idee}
\markboth{Videns Idee}{Videns Idee}

\chapter*{Første Afsnit.\\ Den subjective Viden.}
\addcontentsline{toc}{chapter}{Første Afsnit. Den subjective Viden}

\section*{Første Kapitel.\\ Viden og den Vidende: Subjectivitetens Logik.}
\addcontentsline{toc}{section}{Første Kapitel. Viden og den Vidende: Subjectivitetens Logik}

% --- p. 1 ---
\setcounter{footnote}{0}%
% The three display heads (Første Deel / Videns Idee; Første Afsnit / Den
% subjective Viden; Første Kapitel / Viden og den Vidende: Subjectivitetens
% Logik) stand at the head of this page and are already set by the skeleton.
% The signature mark at the foot of p. 1 (`Grundideernes Logik.' + `1') is a
% printer's mark and is not transcribed.
Hvormed skal Philosophien begynde? Dette Grundspørgsmaal er i den
nyere Tid blevet opfattet og besvaret paa hoist forskjellig Maade.
% `hoist' printed with a plain o (no slash) -- verified at 600 dpi against the
% clearly slashed ø of `spørgsmaal', `nødvendigviis' and `indrømme' on this
% same page. The book is inconsistent in this; transcribed as printed.
Forsaavidt Talen er om en systematisk Begyndelse, maa man vistnok
indrømme, at Philosophien skal begynde baade med og uden
Forudsætning, at den paa een Gang skal begynde med Alt og med Intet;
men det kommer nu rigtignok væsentlig an paa, hvorledes denne Fordring
nærmere forstaaes. At begynde med Alt og dog med Intet er, overfladisk
seet, unegtelig en Modsigelse; men dersom man paa Grund af Modsigelsen
vilde afvise Fordringen, vilde man med det Samme afvise al egentlig
Begyndelse, thi Modsigelsen ligger i Begyndelsen selv. Philosophien
maa, dersom dens Begyndelse skal have methodisk Gyldighed,
nødvendigviis begynde med en Forudsætning, der paa en skjult Maade
indeholder Alt, nemlig Alt, hvad Videnskaben formaaer at udvikle i
videnskabelig Form; thi i streng
% --- p. 2 ---
\setcounter{footnote}{0}%
videnskabelig Form kan, under den fremadstræbende Udvikling, Intet
optages udenfra. Men med Alt til Forudsætning har Udviklingens
Begyndelse dog Intet, den tør fastsætte som et umiddelbart Givet. Hvad
skulde et saadant Givet vel egentlig være? En Antagelse, hvormed
Udviklingen begyndte, men som ikke selv blev udviklet! Altsaa skulde
Udviklingen begynde med at forbigaae en Udvikling, d. v. s. begynde med
at springe sit eget første Led, den egentlige Begyndelse, over; men
saaledes at begynde med at forbigaae den egentlige Begyndelse, er jo
ogsaa en Modsigelse. Det i Begyndelsen Modsigende fremkommer derved, at
Begyndelsen selv opfattes som et Værende, som et i sig afsluttet
Begreb, medens dette Begreb dog er i uendelig Vorden uendelig relativt.
Der begyndes med at forudsætte Alt, idet der begyndes med at ophæve al
Forudsætning; thi der begyndes med den Forudsætning, at Alt skal
udvikles, og at Intet endnu er udviklet.

Begyndelsesbegrebet, det Begreb, hvormed Udviklingen begynder, kan
altsaa ikke være Begrebet af den tomme Begyndelse, et Begreb, der
skulde fremkomme derved, at der begyndtes med en Begyndelse til at
bestemme, hvad Begyndelse er; nei, Begyndelsesbegrebet maa, ifølge
Sagens Natur, være det første adæqvate Udviklingsbegreb, det Begreb,
som, for Philosophiens Vedkommende, paa den meest afgjørende Maade
tilfredsstiller Begyndelsens Fordring. Hvilket er nu dette Begreb?
Hegel svarer: „\emph{det er den rene Væren}“; vi svare: „\emph{det er
den rene Viden}“. Begyndelsen skeer altsaa med, at
Begyndelsesbegrebet selv bliver prøvet.

Om den rene Væren gjælder det unegtelig, at den ikke blot er „den
yderste Abstraction af alt Værende“, men tillige den almindeligste og
meest omfangsrige Bestemmelse, hvorunder alt Værende kan henføres; den
rene Væren kan da med Hensyn paa det Værende siges at være det
indholdsrigeste og dog fattigste af alle Begreber. Væren omfatter Alt,
og indeholder dog, for sig betragtet, Intet; forsaavidt
% --- p. 3 ---
\setcounter{footnote}{0}%
% The foot of p. 3 carries the signature mark `1*' (half sheet); not transcribed.
kan den rene Væren jo vistnok opfattes som Begyndelsesbegreb; men
Begyndelsen selv er fra dette Synspunkt eensidig opfattet. Som den
yderste Abstraction af alt Værende maa jo den rene Væren forudsætte
Abstractionen, nemlig den Abstraction, hvorved den selv fremkommer; men
at forudsætte Abstractionen er at forudsætte Tænkningen: den rene Væren
forudsætter den rene Tænken. Saa lidt som Begyndelsen kan skee med en
Tænken uden Væren, ligesaa lidt kan den skee med en Væren uden Tænken;
men, naar Begyndelsen selv forudsætter Eenhed af Tænken og Væren, saa
begyndes her jo ikke med den rene Væren, men med den rene Viden.

Ved at begynde med Væren vil Hegel udrense alt det Endelige, det
Tilfældige, det blot Psychologiske og Subjective, der ellers klæber ved
den menneskelige Tænkning. Men idet vor Tanke paa denne Maade begynder
med at tænke den almene objective Væren, bliver den selv abstract
objectiv, gjør sig uden videre til Eet med Væren, og glemmer i den Grad
sit eget subjective Væsen, at alle Tankebestemmelser tage sig ud som
rene Værensbestemmelser, alle Tankebevægelser som Værensbevægelser.
Dette er en logisk Mystification.

Man har fra den reale Værens Standpunkt, og det ikke uden Grund,
bebreidet Hegel, at den Væren, hvormed hans Philosophie begynder, ikke
er en objectiv, væsentlig og virkelig Væren, men en abstract subjectiv,
en reen Tankeværen, en almindelig Tanke, saa at Tanken, under
Foregivende af at udvikle lutter Værensbestemmelser, fra Først til
Sidst bevæger sig i Kredsen af sine egne Former og saaledes, istedenfor
at begribe det Værende, spinder sig ind i sit eget Væv.
Mystificationen beroer da forsaavidt paa en Miskjendelse af det
Objective.

Fra den ideale Videns Standpunkt kan man imidlertid med ligesaa megen
Ret bebreide Hegel, at han fra Grunden af miskjender Tænkningens, den
objective Tænknings subjective Væsen; thi han behandler jo Sagen, som
om den tænkende
% --- p. 4 ---
\setcounter{footnote}{0}%
Subjectivitet lod sig ved „den immanente Methode“ nedsætte til en
ørkesløs Tilskuer, der i en objectiv Logik kunde blive Vidne til,
hvorledes den objective Væren formaaer at tænke sig selv, det objective
Væsen at reflectere sig selv, det objective Begreb at begribe sig selv.
Mystificationen beroer da i denne Henseende paa en Miskjendelse af det
Subjective.

Skal Mystificationen hæves, maa Philosophien blive sig tydeligt
bevidst, hvad det vil sige, at den udtrykkelig begynder, ikke med den
rene Væren, men med den rene Viden. Ved at sætte Begyndelsen i Viden,
er Tænkningens subjective og Værens objective Væsen ligelig hævdet.
Idet Begyndelsen udgaaer fra Viden, udgaaer den fra en tænkende og
vidende Subjectivitet, og har i Subjectiviteten sin ideelle
Selvstændighed ligeover for den objective Væren, paa hvis Bestemmelse
den væsentlig gaaer ud. Vistnok kan den objective, reale Væren, ifølge
denne Begyndelse, ikke faae Betydning uden derved, at den bliver
Gjenstand for Viden; men herved er det objective Element paa ingen
Maade blevet tilbagetrængt. Tvertimod! den reale Væren er væsentlig
hævdet; thi medens det ingenlunde ved sig selv er indlysende, at Væren
overalt maa forudsætte Viden, er det derimod ved sig selv indlysende,
at Viden overalt maa forudsætte Væren. Af det Cartesianske:
% the three Latin tags are set in antiqua against the surrounding Fraktur;
% the Danish words between them (`følger et:', `men ikke omvendt et:') are
% Fraktur, as is `Cartesianske'. Verified at 600 dpi.
\textit{cogito, ergo sum} følger et: \textit{cogitat, ergo est}, men
ikke omvendt et: \textit{est, ergo cogitat}.

At Viden som Begyndelsesbegreb forudsætter Alt og indeholder Intet, er
da ogsaa let at paavise. I Viden maa nemlig Alt, hvad der paa nogen
Maade skal blive Gjenstand for philosophisk Behandling, nødvendigviis
opgaae. Gaves der nu en Art Værende, der ganske faldt udenfor Viden,
da maatte dette Værende jo ogsaa falde udenfor Philosophien. Et
saadant Værende maatte da --- thi her spørges ikke om min og din
ufuldkomne Viden, men om Videns evige Væsen --- opfattes af den
Vidende, som det, hvorom Intet, slet Intet kunde vides. Men at vide om
Noget, at man Intet veed
% --- p. 5 ---
\setcounter{footnote}{0}%
om det, er en reen Modsigelse: at vide om en Ting, at Intet vides om
den, er jo dog at vide Noget om den. Vor rene Viden er i sin
Begyndelse desuagtet ligesaa tom, som Hegels rene Væren. Dette følger
ligefrem deraf, at her er abstraheret fra alt Indhold. Her vides vel,
at her er en Viden; her vides i al Almindelighed, hvad Viden er,
nemlig: en Forskjelseenhed af Tanken og Væren; men Eenheden selv er
svævende og ubestemt, fordi Tænkningen i denne Eenhed endnu er ligesaa
ubestemt og svævende som Væren. Er den rene Væren et Intet, fordi den
forsvinder i en ubestemt Tanken, og den ubestemte Tanken igjen et
Intet, fordi den forsvinder i en ubestemt Væren: da maa Viden som
Tænkens og Værens endnu ubestemte Eenhed i denne sin Begyndelse jo
absolut være en Viden af negativt Indhold, en Viden af egen
Ubestemthed, af egen Tomhed, en Viden af Intet.

Som Begyndelsesbegreb maa Viden tillige være Bevægelsesbegreb; denne
Fordring opfyldes da derved, at det endnu Ubestemte bestemmes, at
Forskjelseenheden af Tænken og Væren udvikles. Væren for sig kan ikke
være Bevægelsesbegreb; den rene Væren kan, naar der abstraheres fra den
abstraherende Tænkning, umulig bevæge sig selv. Bevægelsen er, ifølge
Hegel, en Overgang fra Væren til Intet, og fra Intet til Væren; en
saadan Overgang er imidlertid kun en Værensbevægelse, forsaavidt det er
en Tankebevægelse; men en ved Tanken bestemt Værensbevægelse er jo
udtrykkelig en Vidensbevægelse. Kunde Væren uden Tænken gjøre sig selv
til Intet, og det saaledes fremkomne Intet igjen vende sig om og
forvandle sig til Væren: da vilde den første logiske Bevægelse
unegtelig være, ikke blot objectiv, i objectiv Eenhed med Væren, men
saa objectiv, at den faldt udenfor Viden; for at udvikle Ideen af „den
absolute Viden“, maatte Logiken altsaa begynde med en Bevægelse, der
faldt udenfor Viden, og saaledes fra Ikke-Viden gjøre en Overgang til
Viden; men at gjøre en
% --- p. 6 ---
\setcounter{footnote}{0}%
Overgang fra en absolut Ikke-Viden til en absolut Viden er en absolut
Umulighed.

Det er denne Umulighed, der paa en energisk, om end noget sophistisk
Maade indskærpes af de Gamle, naar det hedder: „Den Vidende kan,
forsaavidt han er vidende, Intet faae at vide; thi som Vidende har han
jo allerede Viden, og hvad man allerede har, kan man ikke faae“.
Indvender man, „at En jo meget vel kan vide Noget uden at vide Alt, og
just derved, at han allerede veed Noget, være forberedt paa at faae
Noget at vide, han ikke veed“, saa gives herpaa følgende Svar: „Ingen
kan kaldes en Vidende uden i Forhold til det, han veed; i Forhold til
det, han ikke veed, er han en Uvidende. Det blev altsaa ikke den
Vidende, men den Uvidende, der fik Noget at vide. Men heller ikke den
Uvidende kan, forsaavidt han er uvidende, faae Noget at vide. For at
faae Noget at vide, maatte den absolut Uvidende jo absolut begynde med
en Viden om sin Uvidenhed; for at begynde med en Viden om sin
Uvidenhed, maatte han begynde med en Viden om en Ikke-Viden; men at
begynde med en Viden om en Ikke-Viden er at begynde, ikke som en
Uvidende, men som en Vidende.“

Det Sande heri er, at Viden i det Uendelige forudsætter sig selv,
bevæger sig fra sig selv gjennem alt Andet til sig selv, saa at dens
hele Udvikling væsentlig er en Selvudvikling. At Viden --- ikke min
eller din ufuldkomne Viden, men Viden i sin Oprindelighed ---
forudsætter sig selv og begynder med sig selv, vil da med andre Ord
sige, at den har sit Princip i sig selv, at dens Væsen er Selvhed.
Dersom Viden ikke selv var Princip, kunde den umulig være den egentlige
og første Begyndelse. Hvor nær \emph{Princip} er beslægtet med
% `Princip' and `Begyndelse' letterspaced here; `Begyndelse' is spaced in both
% halves of its printed line break (`Be-' / `gyndelse'), verified at 600 dpi.
% The bare `Princip' two lines above is NOT spaced.
\emph{Begyndelse}, skjønnes jo allerede deraf, at Ord som
\textit{ἀρχή} og \textit{principium} i den frie Sprogbrug ogsaa betyde
Begyndelse. Kan Viden ikke være sin egen Begyndelse, dersom den ikke
er
% --- p. 7 ---
\setcounter{footnote}{0}%
sit eget Princip, saa kan den heller ikke være sit eget Princip, dersom
den ikke er sin egen Begyndelse.

Men endskjøndt her er Eenhed af Princip og Begyndelse, er dog Principet
selv ikke uden videre det Samme som Begyndelsen: Eenheden er en
Forskjelseenhed. Fremhæves Forskjellen, da forholder Principet sig til
Begyndelsen, som det første Bevægende til Bevægelsen. Begyndelsen er
ikke det første Bevægende, men den første Bevægelse; idet Bevægelsen
begynder, falder Begyndelsen ikke udenfor, men sammen med Bevægelsen.
Det er dette Punkt, Eleaterne oversaae; og kun, idet de oversaae dette
Punkt, kom de conseqvent til at negte Bevægelsen. For Achilles kan
% `For Achilles' printed with a plain capital O (no slash), while the lowercase
% `før' six words later plainly carries its slash: verified at 600 dpi.
% Transcribed as printed; cf. `hoist' (p. 1), `Oiekast' and `Hoiere' (p. 10).
naae Skildpadden, maa han allerede have tilbagelagt den halve Vei; men
før han kan have tilbagelagt det Halve af Veien, maa han have
tilbagelagt det Halve af det Halve af det Halve i det Uendelige. „For
at en Bevægelse kan fuldendes, maa den naturligviis være begyndt:
sættes nu den begyndende Bevægelse identisk med den fuldendte, saa have
vi som Modsigelsernes Cardinalpunkt den vistnok urimelige Sætning: „for
at en Bevægelse kan fuldendes, maa den allerede forud være fuldendt“.
Man fordrer, med andre Ord, en absolut Begyndelse; men at fordre en
absolut Begyndelse er en Selvmodsigelse, thi ingen Begyndelse er eller
kan være absolut.“%
% printed footnote mark: *)
\footnote{Forelæsn. ov. „Phil. Propæd.“ 1860--61. S. 70.}

Det i Begyndelsen Absolute er Principet, og det i Principet selv
Relative er Begyndelsen. Da nu Begyndelsen er Bevægelse, saa falder
Principet paa een Gang udenfor og indenfor Bevægelsen. At det Alt
bevægende Princip falder udenfor Bevægelsen, har Aristoteles, skjøndt
rigtignok paa en aldeles eensidig Maade, udtrykkelig fremhævet, idet
han har opfattet det første Bevægende som det i sig selv Ubevægede
(τὸ πρῶτον κινοῦν ἀκίνητον ὄν). At opfatte Forholdet
% --- p. 8 ---
\setcounter{footnote}{0}%
af det Bevægede og det Ubevægede saaledes, at de umiddelbart udelukke
hinanden, er at opfatte det ideeløst; Bevægelse forudsætter Bevægelse i
det Uendelige; kun fra det i sig Bevægede kan en Bevægelse udgaae: det
Alt bevægende Princip maa selv være i Bevægelse. „I den ufuldkomne
Energie falde Hvile og Bevægelse ud fra hinanden; det Hvilende bevæger
sig ikke; det sig Bevægende hviler ikke; det Hvilende maa sættes i
Bevægelse, det Bevægede gaae til Hvile. I den fuldkomne Energie falder
Bevægelsen derimod sammen med Hvilen: her er Hvile i Bevægelsen,
Bevægelse i Hvilen.“%
% printed footnote mark: *)
\footnote{Phil. Propæd. S. 149.}

Skal Principet bestemmes som det i sin Bevægelse Ubevægede, da maa
Principbevægelsen nødvendig opfattes som Selvbevægelse, og Principet
selv som oprindelig Selvhed, som vidende Eenhed. Kun som vidende
Eenhed, som den sig i Alt vidende Viden er Viden sit eget Princip; men
den sig i Alt vidende Viden er i reen Almindelighed den vidende
Subjectivitet. Saa lidt som Viden kan bestemmes og udvikles uden
derved, at den ved sit Andet bestemmer sig selv som en Viden af Noget,
ligesaalidt kan den bestemmes og udvikles som en Viden af sit Andet,
uden derved, at den bestemmes som en oprindelig Viden af sig selv.
Tages Andet, det Reale, der skal udgjøre Videns Indhold, bort, saa
fremkommer den tomme Viden, en Viden af Intet, der ikke tillige er en
Viden af Noget; men en Viden af Intet uden en Viden af Noget er en
absolut Ikke-Viden, en Tilintetgjørelse af al Viden. Det reale Andet,
der i Virkeligheden skal udgjøre Videns Indhold, maa altsaa ligge i
Videns eget Princip; men Videns Princip er den vidende Subjectivitet;
det objective Indhold hører oprindelig med i Subjectiviteten.

At Principet for alt Subjectivt ogsaa er Principet for alt Objectivt,
at Objectiviteten med sit uendelige Værens-%
% A genuine compound hyphen, not a line-break hyphen: the continuation on
% p. 9 opens with the capital I/J sort (`Indhold'), verified at 600 dpi.
% Hence `Værens-Indhold', hyphen retained.
% --- p. 9 ---
\setcounter{footnote}{0}%
Indhold forsaavidt maa være ligesaa oprindelig som Subjectiviteten med
sin uendelige Viden, er uimodsigeligt; men derfor er det ligefuldt en
Miisforstaaelse, og det en grov Misforstaaelse, naar man med Schelling
% The two spellings differ in the print -- `Miisforstaaelse' then
% `Misforstaaelse' -- within one sentence; both verified at 600 dpi and
% transcribed as printed.
vil søge et bagved begge liggende Princip%
% printed footnote mark: *) ; this page carries ONE footnote, in two paragraphs
\footnote{„Fordyber man sig i den Tanke, at Subject og Object gjensidig
betinge og begrændse hinanden (jvfr. Schelling: Vom Ich als Princip der
Philosophie, oder über das Unbedingte im menschlichen Wissen), saa
forsvinde de begge som Momenter i en ubetinget Eenhed, en absolut
Identitet, et Subject-Object, i hvilket Forskjellen oprindelig er
begrundet. Men ved at føre den ideal-reale Udvikling tilbage til et
Princip, der er ligesaa abstract dualistisk som monistisk, ligesaa
umiddelbart subjectivt som objectivt, nødes man til at personificere
den dunkle Realgrund, insinuere Forstandens Lysglimt i den selvløse
drømmende Fornuft, anticipere Friheden uden at vedkjende sig
Anticipationen, og saaledes overføre Subjectivitetens Prædicater paa
den Forudsætning, der endnu ingen Subjectivitet kan have.“ Den
propædeutiske Logik. Kjøbhvn 1845. S. 206.\par
Paa andre Steder, navnlig i „System des transcendentalen Idealismus“
har Schelling unegtelig lagt Eftertryk paa Jeget som Videns Princip;
man seer heri Eftervirkningen af Fichte. „Der Begriff des Ich, d. h.
der Act, wodurch das Denken überhaupt sich zum Objecte wird, und das
Ich selbst (das Object, sind absolut Eins; außer diesem Act ist Ich
% printer's defect: the parenthesis opened at `(das Object,' is never closed
% (Fichte's text has `(das Object)'). Transcribed as printed.
nichts.“ \ldots{} „Ich als reiner Act, als reines Thun, ist im Wissen
selbst nicht objectiv, deswegen, weil es Princip alles Wissens ist.
Soll es Object des Wissens werden, muß dieß durch eine vom gemeinen
Wissen ganz verschiedene Art zu wissen geschehen.“ \ldots{} „Ich ist
nichts Anderes als ein sich selbst zum Object werdendes Produciren.
Die Wissenschaft kann von nichts Objectivem ausgehen; sie muß vom
„Nicht-Objectivem, das sich selbst zum Objecte wird“, ausgehen.“ ---
Men til nogen klar og methodisk Udvikling af, hvorledes Subjectiviteten
oprindelig forholder sig til sit objective Indhold, og hvad
Vidensprincipet egentlig betyder, har den schellingske Philosophie,
trods alle sine Tilløb, ikke bragt det.}, enten man saa iøvrigt vil
kalde dette
% --- p. 10 ---
\setcounter{footnote}{0}%
Princip „Indifferensen“ eller „det Absolute“ eller „ein unvordenkliches
Sein“. At gjøre Indifferensen til Princip er at forudsætte en
oprindelig Eenhed, hvori Modsætningen imellem det Subjective og det
Objective endnu ikke er blevet til, en Grundeenhed, hvoraf denne
Modsætning selv skulde lade sig aflede. Nu tager det sig ved første
Oiekast jo vistnok ud, som om man ved en saadan Fremgangsmaade
% `Oiekast' and the three `Hoiere' below are printed with a plain capital O
% (no slash), verified at 600 dpi against the slashed ø of `indrømme' and
% `Oprindelighed' on this page. Transcribed as printed.
upartisk vilde hævde det Subjective lige Ret med det Objective, idet
man nemlig deducerer begge af et tredie Hoiere. Men Upartiskheden er
kun tilsyneladende; ved at begynde med Indifferensen som det tredie
Hoiere, har man i Virkeligheden begyndt med at tage Parti for det
Objective. Det tredie Hoiere eller rettere: det første Dybere, der
skal ligge til Grund baade for det Objective og det Subjective, er jo
netop, naar vi see nærmere til, intet Andet end det første, endnu
ubestemte Objective selv, det Objective, hvoraf Alt skal deduceres, i
Modsætning til det afledede, med det afledede Subjective coordinerede,
Objective. At deducere det Subjective er i Principet selv ligesaa
umuligt, som at deducere det Objective. Istedenfor ved abstracte
Skyggerids af en indifferent Baggrund at give det Objective et
eensidigt Forspring for det Subjective, maae vi tvertimod, med
Anerkjendelse af begges Oprindelighed, indrømme det Subjective den
ideale Prioritet, forsaavidt Subjectiviteten som vidende Selvhed dog
altid er og bliver det Overgribende.%
% printed footnote mark: *) ; this single note runs over onto p. 11
\footnote{„Subjectivitetens udelukkende Kriterium er den begribende
Selvhed, der bestaaer i Reflexion; Objectivitetens derimod en selvlos
% `selvlos' printed with a plain o here, though the same author's `selvløse'
% on p. 9 carries its slash. Verified at 600 dpi; transcribed as printed.
Gjenstændighed, der mangler Reflexion. Uden Objectivitet vilde
Subjectets Selvbegriben forsvinde i Tomhed, thi alle Forestillinger og
Anskuelser referere sig til Objectiviteten. Uden Subjectivitet vilde
Objectivitetens Begreb være utænkeligt, thi Objectet er, egentlig talt,
hverken for sig selv, ei heller for et andet Object, da al Reflexion
mangler. Sige vi: det ene Object refererer sig til det andet, saa
anvende vi en Reflexionsbestemmelse, og beholde det Ene
\textit{in mente}, naar det Andet forestilles. Vel er der i den
objective Verden en væsentlig Sammenhæng; men da den enkelte Gjenstand
hverken har sig selv i sin Magt, ei heller de andre \textit{in mente}.
% printer's defect: a full stop where the construction requires a comma
% (`da \ldots\ saa'). The mark is a round dot on the baseline with no descending
% tail, unlike the plainly tailed comma of `Magt,' on the same line; verified
% at 600 dpi. Transcribed as printed.
saa kunne Objecterne ikke selv objectivere sig for hverandre, ikke
magte deres egen Totalitet, men maae derimod magtes og objectiveres af
en begribende Subjectivitet.“ Propæd. Logik. S. 204--205.}

% --- p. 11 ---
\setcounter{footnote}{0}%
Da nu alle Begyndelsens Forudsætninger, de objective med de subjective,
ligge skjulte i Videns Princip, den vidende Subjectivitet, saa maa
Udviklingen af det oprindelige Forhold mellem Viden og den Vidende
nødvendig falde sammen med disse Forudsætningers systematiske
Udvikling. Eftersom da Forudsætningerne optræde i abstract
Forskjellighed, i indbyrdes Gjennemsigtighed og i totaliseret Eenhed,
vil Principet udvikle sig for os som \emph{umiddelbar}, som
\emph{reflecterende} og som \emph{begribende Subjectivitet}.
Særkjendet for den \emph{umiddelbare Subjectivitet} er i det Hele taget
% Emphasis extents settled at 600 dpi: `psychologiske Standpunkt' here and
% `Standpunkt' after `begribende Subjectivitets' below are NOT letterspaced,
% though the preceding words are. Transcribed as printed, not regularised.
dens psychologiske Standpunkt. Fra dette Standpunkt af arbeider den
sig efterhaanden dybere og dybere ind i almene Væsensbestemmelser og
ontologiske Opgaver. Særkjendet for den \emph{reflecterende
Subjectivitet} er dens bevidste Opfattelse af sin egen ontologiske
Væsenhed: den reflecterende Subjectivitet er væsentlig ontologisk. Med
den gjennemførte Reflexion føres Subjectiviteten vistnok tilbage til
sit første psychologiske Udgangspunkt; men saaledes, at den nu
erkjender saavel sin evige Grund som sin timelige Oprindelse. At
Subjectiviteten, ontologisk seet, er evig og guddommelig, men,
psychologisk seet, timelig og menneskelig, vækker Bevidstheden om en
Dualisme i dens eget Væsen. Denne Dualisme maa da finde sin Forklaring
paa den \emph{begribende Subjectivitets} Standpunkt. Kun Begrebet af
en fuldkommen Begriben kan lede den tænkende Subjectivitet til at see
sit dobbelte Ansigt i Ideens Speil, og derved til at indsee, hvad det
væsentlig er, hvorpaa al Indsigt, Begriben og Viden beroer.
\subsection*{A.\\ Den umiddelbare Subjectivitet.}
\addcontentsline{toc}{subsection}{A. Den umiddelbare Subjectivitet}

% --- p. 12 ---
\setcounter{footnote}{0}%
Viden er Forskjelseenhed af Tænken og Væren. Med Hensyn paa den
rene Væren er Viden umiddelbar Anskuen; med Hensyn paa den rene
Tænken er den Negation af al umiddelbar Anskuen. Den ved Væren
bestemte Anskuen er Anskuelse; den ved Tænkningen ophævede Anskuen
er Negation. Anskuen og Tænken forholde sig som Væren og Intet; den
umiddelbare Væren kan ikke tænkes; det rene Intet kan ikke anskues.
Intet at anskue er det Samme som ikke at anskue; men at tænke Intet
er ingenlunde det Samme som ikke at tænke. Formedelst Anskuelsen
bliver Tænkningen bestemt ved Væren; formedelst Negationen bliver
Væren bestemt ved Tænkningen: alle Anskuelsesbestemmelser ere
positive, alle Tankebestemmelser negative. At Væren objectivt
omsættes i Intet og Intet i Væren, beroer da oprindelig derpaa, at
Anskuen subjectivt omsættes i Tænken og Tænken i Anskuen. Allerede
heraf sees, hvad det vil sige, at Subjectiviteten er det Overgribende.
Dersom den objective Væren ikke paa nogen Maade lod sig bestemme,
kunde den umulig være; men hvorved skal nu Væren bestemmes? Ved sin
Negation, ved Intet, altsaa ved Tanken! Men Tanken er kun Tanke
derved, at den tænkes, og den tænkes kun derved, at der er en
tænkende Subjectivitet; forsaavidt Væren bestemmes ved sin Negation,
bliver den altsaa bestemt objectivt derved, at den bestemmes
subjectivt. Eller skal Væren maaskee bestemmes, ikke ved Negation,
men ved Position saaledes, at Væren selv objectivt adskiller sig fra
Væren, som det ene Værende fra det andet, som Noget fra Andet? Men
at faae Noget adskilt fra Andet uden Negation er jo aldeles umuligt!
Adskillelsen af Noget og Andet beroer netop derpaa, at Noget er, hvad
Andet ikke er, og omvendt. Skulde nu Væren objectivt bestemme sig ved
sig selv, da maatte den ogsaa objectivt kunne negere sig
% --- p. 13 ---
\setcounter{footnote}{0}%
selv; men den objective Negation er en subjectiv Bestemmen, en Act af
den tænkende Subjectivitet. Eller er det maaskee et logisk Feilgreb,
at lade Væren, den rene almindelige Væren, bestemme sig til et
forskjelligt Værende, og derved til Adskillelse af Noget og Andet:
saa lader os engang vende Forholdet om og antage, at der oprindelig
er en Vrimmel af Detteværende, absolut Tilværende, hvoraf alle
Værens- og Vidensbestemmelser blive at aflede. Istedenfor at begynde
med en subjectiv Eenhed, begynde vi altsaa med en objectiv
Mangfoldighed --- ikke af abstracte og forsaavidt uselvstændige
Værensbestemmelser, men af selvstændigt Værende. I denne
Mangfoldighed af selvstændigt Værende har da hvert Led en Bestemthed,
det hverken har givet sig selv eller erholdt fra andet Værende, en
oprindelig Bestemthed, der ligesaa lidt lader sig henføre til en
Bestemmen, som det paagjældende Lede Tilværen til en Tilbliven; vi
have da i denne abstract objective Forudsætning et Mangfoldighedens
Grundlag, der paa een Gang minder om Demokrits Atomer og Herbarts
Realer.

Men skal en Mangfoldighed af Bestemmelser, der ikke ere fremkomne
og ikke kunne fremkomme ved nogensomhelst objectivt bestemmende
Indflydelse, desuagtet lade sig bestemme, da kan det alene skee
derved, at Bestemmelserne selv vise sig bestemmende for en
oversvævende Opfatning, for en ved dem eiendommelig bestemt Viden,
d. v. s. for en Anskuelse. At de absolute, i sig selv
ubestemmelige, Bestemtheder kun ere bestemmende for en tilsvarende
Anskuelse, vil ved nærmere Betragtning let indlyse. Det absolut
bestemte Atom, det umiddelbart tilværende Reale, kan jo i sin
absolute Bestemthed hverken være bestemt ved sig selv eller ved de
andre ligesaa absolut bestemte Atomer. At Atomet ikke kan være
bestemt ved nogen Selvbestemmen, følger ligefrem deraf, at dets
Bestemthed falder umiddelbart sammen med dets Tilværen; at det da
heller ikke kan være bestemt ved andre bestemmende
% --- p. 14 ---
\setcounter{footnote}{0}%
Atomer, er en Selvfølge, naar Bestemtheden skal gaae i Eet med den
bestemte Tilværen, og denne Tilværen være absolut. Det er saa langt
fra, at de absolut bestemte Atomer skulde kunne udøve nogensomhelst
Indflydelse paa hinanden og derved i nogen Maade bestemme hinanden,
at de tvertimod unddrage sig ethvert muligt objectivt Forhold. Skulde
de virkelig komme i Forhold til hinanden, maatte de naturligviis
blive afhængige af hinanden; men ved at blive afhængige af hinanden,
maatte de jo blive relative; men hvorledes skulde de blive relative,
naar de oprindelig ere absolute?

Fixideen med absolute Atomer, umiddelbart tilværende Realer
o. desl., hidrører ikke fra en virkelig objectiv Væren, men fra en
eensidig subjectiv Tænken, en Tænken, der er bedaaret af Anskuelsen.

At Anskuelsen, uden at vide deraf, kan bedaare Tænkningen, skeer
imidlertid derved, at Tænkningen, uden at tænke derpaa, forvansker
Anskuelsen. Tanken, den abstracte Tanke, kommer aldrig i umiddelbar
eller ligefrem Berøring med det Værende selv; enhver Forhandling
mellem Tænken og Væren mægles ved Anskuelsen. I sin rene
Oprindelighed er Anskuelsen naiv; den bestemmes af det Værende som
det er, opfatter det Givne saaledes som det nu engang viser sig, og
seer ikke Andet. At ingen Anskuelse kan bestemmes af det Værende uden
igjen at bestemme det, derom veed Anskuelsen selv i sin Uskyldighed
Intet; det er en Hemmelighed, der først opgaaer for Tanken. Men
Tanken er negativ: det Negative bestemmes kun ved det Positive,
forsaavidt det bestemmer det Positive og derved sig selv. Tanken
bestemmes af Anskuelsen, idet den negerer Anskuelsen for at bestemmes
af det Værende; men da det udtrykkelig er igjennem en Negation,
altsaa igjennem dens egen negative Selvbestemmen, at den kan
bestemmes af det Værende, saa har den bestemmende Tanke jo allerede
herved faaet den umiddelbare Værensbestemmelse i sin Magt, forsaavidt
den har faaet Anskuelsen i sin Magt.
% --- p. 15 ---
\setcounter{footnote}{0}%
Er Tanken nu med denne sin formelle, paa Negationen beroende Magt
ikke desto mindre selv uklar over Negationen og det Negative, saa
skeer det let, at den i negativ Naivetet umiddelbart sætter et
Positivt, og fixerer det i reen Umiddelbarhed; den fixerer da sig
selv med det Samme, eftersom det saaledes fixerede Positive jo dog
vitterlig er af negativ Oprindelse.

Det er ikke den umiddelbare, i sin Umiddelbarhed naive Anskuelse,
der har opdaget Atomerne og fastsat Realerne; tvertimod! de absolute
Atomer ere tænkte Anskuelser, og man behøver blot at ophæve den
sidste Rest af et umiddelbart Anskueligt for at forvandle de absolute
Atomer til absolute Positioner, altsaa til rene Realer. Det er
Tanken, som i den Indbildning, at den har faaet umiddelbart fat paa
det Værende selv, benytter Anskuelsen og tvinger Forestillingerne, de
individuelle Anskuelser, ind i sit eget Hjernespind.

Det er den i Forandringen skjulte Modsigelse, hvis Opløsning her er
et Tankens Problem. Paa den naive Videns Standpunkt har det ingen
Vanskelighed med Forandringen; ifølge den daglige Erfaring, at
Tingene idelig forandre sig, ere Anskuelsen og Tanken enige med
hinanden om, at Noget naturligviis maa blive til Andet, eftersom det
jo aldrig staaer ene, aldrig beroer paa sig selv alene, men overalt
er i Forbindelse med Andet, støtter sig til Andet, og har som dette,
i sin eiendommelige Fremtræden bestemte, Noget altid et Andet skjult
i sig; Anskuelsen og Tænkningen ere paa en naturlig Maade enige om,
at der i Tilværelsen vel er en Modsætning mellem Noget og Andet, men
at Modsætningen er relativ, og at det netop er denne Modsætningens
Relativitet, der udtaler sig i Forandringen.

Anderledes, naar Tanken med sin abstracte Negation, sit Ikke og
Ikke-Ikke begynder at adskille sig fra Anskuelsen, medens den dog
endnu er naivt sammenvoxet med den. Paa dette Standpunkt af negativ
Naivetet bliver For\-%
% --- p. 16 ---
\setcounter{footnote}{0}%
andringsproblemet, skjøndt under Foregivende af, at det er et reent
Værensproblem, til et eensidigt Forstandsproblem, en fixerende Opgave
for den sig selv fixerende Tænkning. --- Noget, Detteværende, er da i
sin umiddelbare Modsætning til Andet, der jo ogsaa er et
Detteværende, givet for Anskuelsen og derved bestemt som Gjenstand
for Tænkningen; hvad gjør nu Tænkningen ved ethvert saadant
Detteværende? Den tager det positivt, bestemmer det negativt, og
udgiver saa, uden at ændse det i Bestemmelsen Negative, det negativt
Bestemte for et i sig Positivt, for en saa absolut Position, at deri
end ikke skal være Skygge af Negation. Dette er naturligviis en uhyre
Modsigelse; og dog er det netop denne Modsigelse, den abstract
kritiske Tænkning lægger til Grund, det Selvmodsigelsens Grundlag,
hvorpaa den oprindelig stiller sig, for at afvise alle Tilværelsens
og da ogsaa Forandringens Modsigelser.

Naar $A$ som et Noget, $B$ som et Andet, hvert for sig, anskues som
et Detteværende, hvad kan den rene Tænkning da lære os om ethvert af
disse Led? At $A$ er $A$ og ikke $B$, at $B$ er $B$ og ikke $A$! Man
seer øieblikkelig, at Tanken er aldeles negativ, at den af sig selv
ikke veed det Allerringeste om det Positive. Vistnok siger Tanken, at
$A$ er $A$; men naar der saa spørges: hvad er da $A$? saa gjentager
den kun hvad den har fra Anskuelsen. Hvad har den da fra Anskuelsen?
At $A$, Detteværende, er et Positivt; men ingenlunde, at det er et
absolut Positivt. Det absolut Positive, et positivt Noget, hvori der
end ikke er Anlæg til Negation eller Andethed, er en Opfindelse af
den abstracte Tænkning, men en Opfindelse, som Tænkningen desuagtet
henfører til Anskuelsen, for at faae den til at gjælde for det
Værende. Hvorved skeer det nemlig, at det ved Anskuelsen givne
Positive bliver, endskjøndt det kun er relativt, til et absolut
% printed note mark: *) ; the note begins at the foot of p. 16 and runs on,
% below a rule, through the whole of p. 17 (whose body text is only two lines).
Positivt?\footnote{Det vil i en senere Sammenhæng blive oplyst, at
dette for Tanken, som det synes, absolut Positive er i sin objective
Relativitet et Magtens Element, hvorved Viden bestemmes, en
Realnegation af det blot Logiske, en antilogisk Factor, der er
afgjørende for Forholdet mellem Begreb og Existens. Den kantiske
„Ding an sich“ fremkommer just derved, at den antilogiske Factor
særlig opfattes og fastholdes for sig. Den herbartske Manie med de
absolute Realer og Bestræbelsen for at vride og dreie det Logiske
saaledes, at Tankebestemmelserne faae Udseende af at være rene
Existensbestemmelser, hidrører øiensynlig fra en Reflexion, som er
blendet af den antilogiske Reflex, Existensen kaster, og som bilder
sig ind, at dette Antilogiske just er det rette Logiske, og som saa,
for at naae det rette Logiske, arbeider paa at drive Tankerne udenfor
% letterspacing breaks off at the line end: only "Foran-" is spaced, "dringen" is not
Tænkningen og reflectere sig ud over al Reflecteren. \emph{Foran}dringen
være et Exempel. Som Overgang fra Noget til Andet er Forandringen en
kategorisk Bestemmelse ved alt Tilværende. $A$, et Tilværende, maa
altsaa, ifølge sit Begreb, kunne gjennemløbe Tilstandene: $A_0$,
$A_1$, $A_2$ o. s. v. Den Modsigelse, at $A$ under Forandringen er
det Samme og dog ikke det Samme, finder sin logiske Opløsning i
Forandringens eget Begreb. Men mærkes nu Naturen af det Antilogiske,
hvilken Forandring maa der da ikke foregaae med Opfattelsen af
Begrebet: Forandring! I existentiel Umiddelbarhed er $A = A$, hvilket
udtrykkelig strider imod, at $A$ skulde kunne blive enten $A_0$ eller
$A_1$ eller $A_2$; fremdeles er $A_0 = A_0$, hvilket strider imod, at
dets Væren skulde gaae over i dets Ikke-Væren, saa at det kunde blive
til $A_1$ o. s. v. Paa denne Maade kan Tanken, der endnu ikke har
gjennemskuet sit Forhold til Existensen, martre sig med at ville
forstaae sig selv som en Existens, altsaa som Ikke-Tanke, ved at
ville forstaae sig selv som Tanke, altsaa som Ikke-Existens.
Tænkningens hele atomistiske Unatur kan iøvrigt sees af „Realerne“.
Realet er ikke et materielt, men et logisk Atom. Det materielle Atom
er som Gjenstand for Tanken, som Tankens Realnegation, det, hvorover
der tænkes, men som ikke selv lader sig tænke, et antilogisk Element,
og har som saadant sin ydre Tilværelse i Rummet. I Tankens Opfattelse
af Atomet er det Logiske altsaa forenet med det Antilogiske. Realet
derimod, det logiske Atom, er behæftet med den ufordragelige
Modsigelse, at dets logiske Væsen absolut skal være antilogisk.} Det
skeer,
% --- p. 17 ---
\setcounter{footnote}{0}%
mærkværdigt nok, ved lutter Negationer! Da $A$, Detteværende, er
ifølge Anskuelsen $A$, saa indseer Tanken, at naar
% the foot of p. 17 carries the signature mark "Grundideernes Logik." + "2" (not transcribed)
% --- p. 18 ---
\setcounter{footnote}{0}%
$A$ er, hvad $A$ er, saa er $A$ ikke, hvad $A$ ikke er. Med samme
Klarhed indseer Tanken, at naar $B$, Detteværende, ifølge Anskuelsen
er $B$, saa kan dette $B$ kun være, hvad det er, ved ikke at være,
hvad det ikke er. Anskuelsens umiddelbare og i al sin Umiddelbarhed
dog relative Forskjel mellem $A$ og $B$ er da paa denne Maade
forvandlet til en absolut Modsætning. At $A$ ikke er, hvad $A$ ikke
er, betyder, at $A$ ikke er $B$ og i al Evighed aldrig kan endog blot
nærme sig til at blive $B$; at $B$ ikke er, hvad $B$ ikke er, beviser
med samme Eftertryk, at $B$ i al Evighed aldrig kan endog blot nærme
sig til at blive $A$. De absolut Positive: $A$ og $B$, dette Noget og
dette Andet, ere da ved de gjentagne Negationer, ret som ved logiske
Besværgelser, blevne adskilte ved et saa svælgende Dyb, at enhver
Overgang fra Noget til Andet, altsaa enhver mulig Forandring, er i
Sagen selv bleven aldeles umulig.

Alle Vanskeligheder med saadanne rene Begreber som: Vorden,
Forandring, Bevægelse o. desl., hidrøre væsentlig fra en forvirret
Opfatning af Forholdet mellem Anskuen og Tænken. At den rene Tænken
skulde, som Hegel antager, kunne bevæge sig selv uden Hjælp fra
Anskuelse og Forestilling, er ligesaa usandt, som at den rene Væren
skulde være istand til at sætte enten sig selv eller Tanken i
Bevægelse. Naar Hegel desuagtet faaer den rene Tanke, som det synes,
til at bevæge sig uden al Anskuen, da skeer det aabenbart kun derved,
at den af Tanken ophævede Anskuen desuagtet bevares i Tanken, saa at
den fra Bevægelsen udelukkede Anskuelse dog, som en logisk
% Latin in antiqua; broken "expel-/las" over the printed line, joined here
Stadfæstelse af det Gamle: \textit{naturam expellas furca}, er med i
Bevægelsen. Hvor utænkelig derimod Bevægelsen er, naar Tanken skarpt
vil skille sig fra Anskuelsen, kan fra mange Sider oplyses, saavel
ved Herbarts som ved Eleaternes Fremgangsmaade. I enhver nok saa reen
Tankebevægelse er der i Virkeligheden altid en, om end uendelig
% "Bianskuelse": initial read as Fraktur B at 600 dpi; cf. "Biforestilling" p. 19
flygtig, Bianskuelse af Rum og Tid. Skal der
% --- p. 19 ---
\setcounter{footnote}{0}%
f. Ex. gjøres Overgang fra Væren til Intet, da maa man naturligviis
begynde med at holde de to Begreber Væren og Intet ud fra hinanden,
eller, som det ogsaa udtrykkes, udenfor hinanden; men saadanne
% letterspacing breaks off at the line end: "ud fra, uden-" is spaced, "for," is not
Bestemmelser som: \emph{ud fra}, \emph{uden}for, beroe aabenbart paa
en Biforestilling om Rum; da nu desuden Overgangen selv: Væren,
Intet, Væren, betegner, at Væren vexler med Intet, ligesom Intet
igjen med Væren, og følger af Intet, ligesom Intet igjen af Væren,
maa denne Vexlen og Følgen nødvendig forudsætte, ikke blot en
Biforestilling om Rum, men ogsaa, og det ligesaa oprindelig, en
Biforestilling om Tid: al virkelig Tænken forudsætter Rummet og
foregaaer i Tiden.

Biforestillinger af Tid og Rum træde endnu stærkere frem ved
Bestemmelsen af Noget og Andet, efterdi Noget og Andet ere concretere
Udtryk end Væren og Intet. Skulle Noget og Andet som to „Realer“
holdes ud fra hinanden og ved Siden af hinanden, saa paatrænger
Rumsanskuelsen sig af sig selv. Da nu de to Realer, hvor bestemt de
end stilles imod hinanden, ikke desto mindre skulle fastholdes, ikke
som udvortes Ting, men som absolut positive Tankegjenstande, ja som
absolut positive Tanker, maae de nødvendigviis holdes ud fra
hinanden, og det ikke i et almindeligt Rum med naturlige Mellemrum,
men i et reent Tankerum, et Rum, hvori Tanken frit kan nedlægge sit
„Imellem“, uden at her derfor egentlig maa tænkes paa noget
„Imellem“. Paa denne Maade har Herbart for Realernes Skyld beriget
% printed note mark: *) ; the note begins at the foot of p. 19 and runs on
% through the foot of p. 20
Metaphysiken med et intelligibelt Rum\footnote{„Det ene, for sig
værende Reale ligger aldeles udenfor, og ikke indenfor eller tildeels
indenfor det andet. De mangfoldige Realer, eller, hvad der allerede
er tilstrækkeligt, tvende Realer kunne være sammen, men ogsaa ikke
sammen. Have vi saaledes $A$ og $B$, kunne vi med $A$ forbinde
Billedet af $B$, eller omvendt. Herved have vi faaet fire Begreber,
nemlig de to virkelige $A$ og $B$ og deres Billeder. Efterat denne
Operation er foretaget, kan man atter forbinde det virkelige $A$ med
et allerede forbundet Par, nemlig $B$ og Billedet af $A$, saa at
disse tvende nu ere sammen, medens Billedet af $B$ ikke er sammen med
disse. Vi ville intet Mellemrum mellem $A$ og $B$; de ere ikke
sammen, men der er alligevel Intet imellem dem. Klagerne over, at man
ikke kan forestille sig det, hjælpe her Intet. Begrebet skal forblive
reent; og vi begjære ingen Billeder, som om man kunde anskue dem; men
vi fordre Begreber og deres Sammenknytning.“ Imod dem, der ville
indføre almindelige Rumsforestillinger, udtaler Herbart sig paa
følgende Maade: „Idet vi fordrede, at man skulde tænke $A$ og $B$ som
\emph{ikke samværende}, have Hine efter deres vante Viis rykket $A$
og $B$ i en vilkaarlig Afstand ud fra hinanden, ret som om der
allerede var Rum nok, hvoraf man kunde indskyde en vilkaarlig
Størrelse mellem $A$ og $B$. Denne vistnok uundgaaelige Indskyden
forbyde vi ikke desto mindre. Alt hvad der paa nogen Maade kunde være
\emph{mellem} $A$ og $B$, naar det ikke er hine tomme Billeder, som
vi selv frembragte, skal forsvinde, og maa i Tanken atter
tilintetgjøres.“ Noget mere Vilkaarligt, Søgt og Forskruet end hvad
den iøvrigt skarpsindige og i andre Henseender af Videnskaben
fortjente Herbart her (Metaphys. II. 119, flgd.) har sammenbragt, vil
Philosophiens Historie ikke let frembyde. Jvfr. P. S. V. Heegaard,
den herbartske Philosophie S. 25--30; S. 111--121.}.
% the foot of p. 19 carries the signature mark "2*" (not transcribed)

% --- p. 20 ---
\setcounter{footnote}{0}%
Men hvad har nu et saadant Rum egentlig at betyde? I skarp
Adskillelse fra Sandserummet er det intelligible Rum naturligviis en
tankeløs Tanke, Anskuelsen af en tilintetgjort Anskuen. At
Sandserummet selv har en intelligibel Side, maa indrømmes; men naar
det intelligible Rum opfattes som Sandserummets egen tænkelige Side,
da forsoner det netop Anskuen med Tanken og afgiver et væsentligt
Medium for det Anskueliges Overgang i det Tænkelige, det Positives
Overgang i det Negative, Værens Overgang i sit Intet, Nogets Overgang
i sit Andet.

Medens Hegel lægger Bevægelsen ind i den rene Tænken og Herbart
stræber at forlige Tænkningen med en forestillet
% --- p. 21 ---
\setcounter{footnote}{0}%
Forandring ved at kunstle med Realerne og det intelligible Rum, have
Eleaterne --- hvad der da ogsaa, naar der endelig skal sættes Splid
mellem Tænken og Anskuen, er det eneste Consequente --- paa det
Strengeste udelukket Kategorierne Bevægelse, Vorden og Forandring fra
% "skudt": sk-ligature, faded at the left edge; "sludt" is not a word
den rene Tænken og skudt alle slige Sandselighedsbestemmelser ind
under den blotte Menen (\textit{δόξα}). Hvad der kan anskues, kan
altsaa, ifølge denne Synsmaade, ikke tænkes. Her anskues f. Ex. en
vis Afstand, en Veilængde; hvad er nu, i skarp Modsætning til det
saaledes Anskuede, Tankens Opgave? At anbringe sin Negation, at sætte
Grændse! Paa den angivne Veilængde sættes da utallige Grændser, og
dog kan Længden selv ikke sammensættes af Grændser; i denne Splid
mellem Anskuen og Tanken møde vi igjen vor „Achilles“. Den blot
grændsesættende Tanke er som uden Bevægelse; kun ved Anskuelsens
Hjælp kan Tanken bevæge sig, kun den sig bevægende Tanke kan tænke
Bevægelsen.

„I Grændsen reflecteres det Begrændsede. At tænke Grændsen uden at
underforstaae et derved Begrændset, er ligesaa umuligt, som at fatte
Betydningen af en Negtelse uden at tænke paa det, som derved bliver
negtet. Den umiddelbare Nulgrændse bliver ogsaa ligefrem sat under
den Forudsætning, at der i Anskuelsen er et Givet, hvis Grændse den
er. At saadanne Grændser som Punkter, Linier og Flader, uagtet de i
% letterspacing breaks off at the line end: only "anskues i" is spaced
og for sig ere rene Negationer, dog \emph{anskues i} positiv Form:
Punktet som Liniens mindste Deel, Linien som Fladens Rand, Fladen som
en Side af Legemet, er et slaaende Beviis for, at Grændsen gaaer i
Eet med det Begrændsede. Men naar Grændsen nu ikke desto mindre
adskiller sig fra det Begrændsede som Negationen fra det Positive,
naar denne Adskillelse er i lige Grad gyldig for Tanken og
Anskuelsen: saa indeholdes heri den Opgave at følge det, som skal
begrændses, til Grændsen. Hvorledes naaer man nu fra Grændse til
Grændse? Ved en Bevægelse! Her er
% --- p. 22 ---
\setcounter{footnote}{0}%
oprindelig ingen anden Udvei, intet andet Middel. Forsøger man i
Tanken at afsætte t. Ex. paa en Linie to Grændsepunkter, $A$ og $B$,
i en vis Afstand, da maa man, om end ubevidst, foretage en Bevægelse,
idet man, for at anskue den vilkaarlig valgte Afstand, lader Blikket
gjennemløbe Veien fra $A$ til $B$. Blikket maa fæste sig, hvor det
vil, det møder paa sin Vei kun Grændse, thi hvert Punkt i Linien er i
og for sig et Grændsepunkt \dots{} At sætte Grændse umiddelbart ved
Grændse, er umuligt, thi idet der abstraheres fra det Begrændsede,
falde de to Punkter sammen \dots{} Imellem to paa hinanden følgende
Grændser er her altsaa et Begrændset. Dersom nu dette Begrændsede
faldt umiddelbart sammen med Grændserne, saa at her intet, slet intet
Mellemværende var, maatte Grændserne selv falde sammen og kunde ikke
følge paa hinanden; faldt det Begrændsede som et umiddelbart
Tilværende mellem de derfra forskjellige Grændser, saa var her en
endelig Afstand, altsaa utallige Grændser mellem to paa hinanden
umiddelbart følgende Grændser, hvilket er en Modsigelse: i den
dialektiske Eenhed af Grændsen og det Begrændsede er al Bevægelse
% printed note mark: *)
begrundet.“\footnote{Forelæsninger over „Phil. Propæd.“ 1860--61.
S. 56--58.}

Vistnok er al Tænken i sig selv en Bevægelse, men den blot
grændsesættende Tanke bevæger sig just paa en saadan Maade, at den
ikke selv kan mærke sin egen Bevægelse. For at Bevægelsen skal kunne
mærkes, maa der i Grændsen være et Mere end Grændsen. Til Indsigt i
Tankebevægelsen er just derfor Læren om Differentialet, det
mathematiske Element, af afgjørende logisk Betydning; thi
Differentialet, det intellectuelle Anskuelsesmoment, er Grændse og
dog mere end Grændse. Differentialet er hverken Nul eller endelig
Størrelse, hverken et abstract Negativt eller et abstract Positivt,
hverken et Produkt af Tænken for sig uden Anskuen eller af Anskuen
for sig uden Tænken: det mathematiske Element er
% --- p. 23 ---
\setcounter{footnote}{0}%
et Bevægelsens Punkt i punktuel Bevægelse, Anskuelsens og
Tænkningens, det Endeliges og Uendeliges, det Positives og Negatives
% printed note mark: *)
urolige Eenhed\footnote{„Vistnok gjør det en væsentlig Forskjel, om
Overgangen anskues i Rummet eller blot forestilles som arithmetisk
Overgang fra Endeligt til 0, fra Endeligt til $\infty$; men dette
staaer fast: uden Reflexer af Rum og Tid ingen Bevægelse, uden
Bevægelse ingen Qvantiteren, uden Qvantiteren intet Qvantum, uden
Qvantum ingen Mathematik. At udelukke Bevægelseskategorien fra den
rene Mathematik og endda beholde Kategorierne: „Grændse,“
„Overgang“, „Tilnærmelse“, er en Modsigelse. Hvorledes Kategorierne i
Mathematiken vælges, betones og anvendes, forsaavidt det technisk
gjælder om den meest hensigtssvarende Udførelse af exacte
Operationer, maa naturligviis afgjøres i Mathematiken selv; men
Operationsbegrebernes logiske Sammenhæng og dialektiske Overgange
undersøges i Philosophien. Hænder det saa --- thi det kan jo hændes
--- at en Fagmand, en Routineret, hvis Begrebsmechanik forlængst er
gaaet i Eet med hans Operationsmechanik, i en saadan Anledning bliver
høirøstet og ved høirøstet Anvendelse af det saraceniske Dilemma
udtrykkelig fordrer, at Philosophien skal opgive Dialektiken og
ubetinget hylde Begrebsmechaniken: da kan en saadan Hændelse dog ikke
videre afbryde, men snarere fremskynde den dialektiske Udvikling.“
Forelæsninger over „Phil. Propæd.“ 1860--61 S. 59 flgd. Jvfr. Phil.
og Math. S. 23--24; Begrebet Nøiagtighed S. 45--47. Hr. Prof. Steens
„Afregning“, gjennemseet af R. N., S. 7, Bilag S. 29.}.

Staaer det nu fast, at det sande Forhold saavel imellem Tænken og
Væren, som imellem Tænken og Vorden, Tænken og Forandring, Tænken og
Bevægelse, kun kan mægles ved Anskuelsen, saa indsees heraf, at
Subjectivitetens Forhold til den virkelige Verden og derigjennem til
sig selv da ogsaa maa, fra Videns Standpunkt, være at bestemme ved
det principielle Forhold imellem Anskuelse og Tænkning. Men
Anskuelsen er i sin første concrete Umiddelbarhed en Sandsning,
nemlig en Sandsning af det umiddelbart Objective, en første
Objectivering. Ved at give sig hen i Sands\-%
% --- p. 24 ---
\setcounter{footnote}{0}%
ning og Sandseindtryk vil den vidende Subjectivitet da nedsætte
Tænkningen til en tom Afskygning af det Sandselige, og miskjende sin
egen Oprindelighed: denne Eensidighed er fuldstændig udtalt i den
gjennemførte \emph{Sensualisme}. Ved ubetinget Fastholden af sin egen
uendelige Selvtænken vil den vidende Subjectivitet ikke alene
nedsætte Sandsningen og det Sandselige til forsvindende
Bibestemmelser i den rene Tænken, men tillige lægge an paa at
forvandle al objectiv Virkelighed til en blot negativ Betingelse for
sin egen Selvudvikling: denne Eensidighed har fundet sit consequente
Udtryk i den \emph{subjective Idealisme}. Idet den umiddelbare
Subjectivitet da paa denne Maade sees at kaste sig fra det ene Extrem
over i det andet, vil det under Bevægelsen indlyse, hvad der da
nødvendig maa antages at ligge i dens oprindelige Natur og udgjøre
dens ontologiske Charakteristik. Altsaa:

% display head, centred; "a)" in antiqua, the rest letterspaced Fraktur
\subsubsection*{a) Sensualistisk Subjectivitet.}

Ved sensualistisk Subjectivitet kan her i denne Sammenhæng ikke
tænkes paa en i Sandseanskuelser fortabt, i det Sandselige hvilende
Bevidsthed; en saadan med sandseligt Indhold mættet Subjectivitet har
kun ringe Betydning for en væsentlig Bestemmelse af Viden og den
Vidende. Sandselighedsphilosophien, den Sensualisme, hvorom her
handles, har altid til Forudsætning en bestemt endelig Adskillelse af
Sandsning og Tænkning, en Adskillelse, der allerede minder om, at den
fra det Sandselige abstraherende Tænkning er bleven tom og mat og
ufrugtbar. Sensualismens principielle Opgave er: \emph{at aflede
Tænkningen af Sandsningen og Sandsningen af det Sandselige; med andre
Ord: at deducere Subjectiviteten af det sandselig Objective.}

Hvor umuligt det er at løse en saadan Opgave, sees tydeligt nok af
% the print has one sort for I and J; "J." resolved as the initial of John Locke
de derpaa rettede Forsøg. Naar J. Locke, for at forklare „Ideernes“
Oprindelse, begynder med den tomme Sjæl som med „et reent Papir“, en
„\textit{tabula rasa}“, saa
% --- p. 25 ---
\setcounter{footnote}{0}%
begynder han alligevel med apriorisk-subjective Forudsætninger,
forsaavidt han nemlig gaaer ud fra en given Adskillelse mellem
udvortes og indvortes Sandsning, mellem „Sensation“ og „Reflexion“.

Istedenfor at begynde med en tom Sjæl, det vil sige: med en Sjæl
uden Anlæg, begynder Locke tvertimod, skjøndt han ikke ret vil vide
deraf, med en Sjæl fuld af Anlæg, ja med et heelt System af sjælelige
Kjendsgjerninger. Thi han forefinder jo Sandsningen som en
Kjendsgjerning og lader den gjælde som Kjendsgjerning; han forefinder
Reflexionen som Kjendsgjerning, og lader den gjælde som
Kjendsgjerning. Og hvormange psychologiske Kjendsgjerninger
indeholdes her saa ikke, naar vi see nærmere til, i disse to angivne!
Det er saa langt fra, at Locke har formaaet at deducere det
Subjective af det Ikke-Subjective, at han meget mere begynder, ikke
blot med en Subjectivitet, men med en Subjectivitet, der er i
oprindelig Selvvirksomhed.

Lockes nærmeste Forgængere, Baco, Hobbes, Gassendi, gik i det Hele
ud fra den Forudsætning, at vore subjective Begreber kunde berigtiges
og beriges ved Erfaringens Kjendsgjerninger, uden at undersøge,
hvorledes det egentlig forholdt sig med Begrebernes psychologiske
Oprindelse. Paa dette Punkt er Locke bleven opmærksom, og ved denne
% English title in antiqua; broken "under-/standing" over the printed line, joined
Sagens Vending kan hans \textit{Essay concerning human understanding}
vel siges at betegne et Vendepunkt i den nyere Philosophies Historie.

Men saa eensidig er Locke i sit Forsøg beskjæftiget med
Gjendrivelsen af den cartesianske Lære om „medfødte Ideer“, at han af
Iver for at vise, hvorledes Begreberne efterhaanden maae være dannede
under Indtryk, der skyldes Erfaringen og følgelig ere \textit{a
posteriori}, ganske forglemmer at lægge Vægt paa det Aprioriske og,
om man saa vil, „Medfødte“ i den Subjectivitetsvirksomhed, uden
hvilken Begrebernes Dannelse, trods alle Sandseindtryk og empiriske
Kjendsgjerninger, vilde være
% --- p. 26 ---
\setcounter{footnote}{0}%
umulig. Det er da saa langt fra, at Locke har formaaet at aflede
Sandsningen af det Sandselige, at han, for at bestemme, hvad det
Sandselige er, udtrykkelig har maattet forudsætte Sandsningen; han er
saa langt fra at have afledet Tænkningen af Sandsningen, at han
tvertimod har maattet bruge Tænkningen for at forklare Sandsningen;
saa langt fra at have bragt noget Activt ud af det blot Passive, at
han, uden at vide deraf, just har givet os Beviser for, at den
passive, til \textit{tabula rasa} reducerede Subjectivitet oprindelig
er og maa være activ.

Den umiddelbare Selvstændighed, som Locke havde indrømmet
„Reflexionen“ ved Siden af „Sensationen“, var imidlertid saa
iøinefaldende og stridende mod Sensualismens Væsen, at sindrige Forsøg
paa at faae den bortforklaret ikke kunde udeblive. Det er i denne
Henseende et virkeligt Fremskridt, naar Condillac fører Reflexionen
tilbage til Sensationen og saaledes gjør Alvor af at deducere
Tænkningen af Sandsningen. Sandsningen, hedder det, er i sin første
Umiddelbarhed og ubestemte Almindelighed Fornemmelse. Paatrænger nu en
enkelt Fornemmelse sig med eiendommelig Styrke, saa fremkommer
Opmærksomheden; men den saaledes vakte Opmærksomhed er ikke en fra
Fornemmelsen forskjellig Sjæleevne; tvertimod! det er Fornemmelsen
selv, der har forvandlet sig til Opmærksomhed. Fornemmelserne vexle;
Fornemmelsesevnen, Sensationen, deler sig formedelst Opmærksomheden
imellem den Fornemmelse, vi have havt, og den Fornemmelse, vi nu have.
Den Fornemmelse, vi ved det umiddelbare Indtryk have, er Følelse; den
Fornemmelse, vi have havt, men hvis reelle Indtryk er forsvunden, er
Hukommelse; saaledes er Hukommelsen da heller ikke en fra Fornemmelsen
forskjellig Sjæleevne: Hukommelsen er en forvandlet Fornemmelse. Ved
Modsætning mellem Følelse og Hukommelse forvandler sig igjen
Opmærksomheden, thi her danner sig nu en dobbelt Opmærksomhed, en, der
bestemmes ved umiddelbar Sands\-%
% --- p. 27 ---
\setcounter{footnote}{0}%
ning, og en, der bestemmes ved Hukommelsen. Men, naar her er to Slags
Opmærksomhed, saa er her jo ogsaa Sammenligning; thi at være
opmærksom paa to Ideer og at sammenligne dem, er Eet og det Samme. At
sammenligne er imidlertid umuligt uden at mærke sig Lighed og Ulighed;
men at mærke sig indbyrdes lige Relationer er at dømme o. s. v. Ved
denne Fornemmelsernes Forvandling (\textit{transformation des
sensations}) fremkommer saaledes i væsentlig Affølge: først
Opmærksomheden, deraf Sammenligningen og deraf igjen Dommen; paa denne
Maade skal, ifølge Condillac%
% printed footnote mark: *)
\footnote{\textit{Traité des sensations,}%
% the note ends with a comma where a full stop is expected; so printed --- see the
% errata note in the header.
}, Tænkningen lade sig deducere af Sandsningen.

Men denne hele Deduction beroer aabenbart paa en Misforstaaelse. Det
er ikke Fornemmelsen, der uden nogensomhelst anden Betingelse i Sjælen
forvandler sig til Opmærksomhed; men det er den i Subjectiviteten
slumrende Opmærksomhed, Opmærksomhedsanlæget selv, der vækkes ved
Sandseindtrykket; det er ikke Sandsningen, der ved sig selv
efterhaanden forvandler sig til Tænkning; men det er den i Sandsningen
skjulte Tænkning, der lidt efter lidt udvikler sig af Sandsningen. Er
Sandsningen i Principet sin egen Forudsætning, da er Tænkningen det
ikke mindre. Saalidt som en dyrisk Sandsebevidsthed kan ved gunstige
Betingelser forvandle sig til menneskelig Selvbevidsthed, ligesaa lidt
kan Dyrets sandselige Tænkning forvandles til en fri, fornuftig,
menneskelig Tænkning. Kun den Tænkning, som ifølge Sjælens
Eiendommelighed oprindelig er anlagt i, med og under Sandsningen, kan,
ved Subjectivitetens Selvudvikling, udvikle sig af Sandsningen.

Sæt endog, at man i en vis overfladisk Forstand lod den foregivne
Deduction gjælde: var saa Sensualismens Opgave hermed løst?
Ingenlunde! Hvormeget man end mener at kunne aflede af den umiddelbare
Sandsning, saalænge man
% --- p. 28 ---
\setcounter{footnote}{0}%
ikke paa nogen Maade seer sig istand til at aflede Sandsningen selv af
det Sandselige, er Subjectivitetens Oprindelighed, Selvhedens og
Selvets aprioriske Gyldighed stiltiende forudsat. Men hvorledes skulde
det vel være muligt at aflede Sandsningen selv af det Sandselige? At
aflede Sandsningen af „Sensibiliteten“ og bestemme Sensibiliteten som
en Beskaffenhed ved et „Ganglievæv“, er en psychologisk Udflugt, der i
denne Sammenhæng forraader Tankeløshed. Thi „det første for Sjælen som
Sjæl Eiendommelige, det Mærke, hvorved Sjæleprincipet just adskiller
sig fra det almindelige Livsprincip, er Fornemmelsen. Men hvad er
Fornemmelse? Umiddelbar Selvreflexion, Selvets ubevidste
Selvopdagelse, dets Indtryk af sig selv under Indtryk af Andet, et
Befindende, hvori det finder sig sat.“%
% printed footnote mark: *)
\footnote{Forelæsninger over „Phil. Propæd.“ 1860--61 S. 335.}

Havde Sensualismen virkelig været istand til at løse sin Opgave: at
deducere Sandsningen af det Sandselige, da vilde Materialismen
unegtelig have været dens sande og høiere Conseqvens; nu derimod, da
det har viist sig, at Opgaven ikke kunde løses, efterdi man, for at
forklare det Sandselige, nødvendig maatte forudsætte Sandsningen, er
det ikke Materialismen, men Spiritualismen, der i strengere Forstand
angiver Conseqvensen; thi Spiritualismen begynder jo netop med at
fæste Opmærksomheden paa Sandserne og Sandsningen, for at forstaae de
sandselige Ting. „Sandselige Ting“, siger Berkeley, „ere saadanne, som
umiddelbart (\textit{immediatly}) opfattes ved Sandserne; hvorimod
Tingenes Aarsag, der ikke kan sandses, udtrykkelig maa tænkes og, idet
den tænkes, være i Bevidstheden. Seer jeg, at den ene Deel af Himlen
er rød, den anden blaa, saa slutter jeg, at der maa være en Aarsag til
denne Farveforskjel; men Aarsagen kan ikke siges at være en sandselig
Ting. Sandselige Ting ere Indbegreb af sandselige Egenskaber
(\textit{combination of sensible qualities});
% --- p. 29 ---
\setcounter{footnote}{0}%
disse sandselige Egenskaber ere ikke udenfor Sjælen, men i Sjælen.“ At
Tingene udenfor os forvandles til „Ideer i vor egen Sjæl“, er jo
vistnok ikke at forstaae, som om de sandselige Forestillinger beroede
paa vor Vilkaarlighed; de maae meget mere henføres til en af os
uafhængig Aarsag. Men da Aarsagen og Virkningen maae være ligeartede,
og da en Ting umulig kan meddele os, hvad den ikke selv har, saa er
det ligesaa urimeligt at antage, at en Ting, der ikke selv har
Fornemmelse, skulde kunne vække Fornemmelse, som at et Hjerneindtryk
skulde kunne blive til en Forestilling: Der existerer altsaa ingen
Legemer, men kun Aander; det er ikke en materiel Aarsag, men en
almægtig Aand, der igjennem alle Sandsebilleder, Forestillinger,
Fornemmelser og Indtryk lader de endelige Aander meddele sig til
hverandre.

Er Sensualismen i sit Væsen Spiritualisme, og Spiritualismen i sin
oprindelige Forudsætning Sensualisme, medens de dog i deres
Umiddelbarhed fremtræde som tvende hinanden udelukkende Extremer, da
er det en Selvfølge, at den Subjectivitet, der i den skarpt udtalte
Modsætning ikke kan opdage Eenheden, maa ved at svinge mellem
Extremerne betages af Svimmelhed og ende i Skepticisme. Den Modsigelse
i Aarsagen, at den som Realbegreb synes at udspringe af Sandsningen,
naar man vil aflede den af Tænkningen, og dog ligesaalidt kan
nedstamme fra den rene Sandsning som fra den rene Tænkning, den
Modsigelse, der i det Hele afgiver Grundlaget for Hume's Kritik af
Aarsagsbegrebet, er væsentlig Eet med den Modsigelse, at
Spiritualismen i Grunden er Sensualisme og Sensualismen
Spiritualisme.

\subsubsection*{b) Idealistisk Subjectivitet.}

Beroer Skepticismen nu væsentlig derpaa, at en Subjectivitet, der er
bleven splidagtig med sig selv, maa opgive al Forvisning og dermed al
sand Viden, fordi den har opgivet Forvisningen om sin egen aprioriske
Gyldighed, da maa
% --- p. 30 ---
\setcounter{footnote}{0}%
den første Betingelse for at tilbagevinde den tabte Viden aabenbart
være den, at Selvvisheden atter indfinder sig, at den forsagte
Subjectivitet igjen mander sig op. Subjectiviteten maa ved at tabe sig
selv vinde sig selv; den maa, for at vinde sig selv, forvisse sig om
Nødvendigheden af at holde fast paa sine egne Former og sine egne
Functioner, idet den med det Samme forvisser sig om, at alle mulige
Modsigelser dog i sidste Instans have deres Udspring fra en skjult
Selvmodsigelse.

„Vilde man“, siger Kant, „spørge den koldblodige, til Dommens
Ligevægt egentlig skabte, David Hume, hvad der da bevægede ham til ved
Betænkeligheder, som kun en møisommelig Grublen kunde udfinde, at
undergrave den for Menneskene saa trøstelige og nyttige Indbildning,
at deres Fornuftindsigt slaaer til, naar Begrebet af det høieste Væsen
skal bestemmes, da vilde han vel svare: Intet uden den Hensigt at
bringe Fornuften videre i Selverkjendelse o. s. v.“ Det var som
bekjendt Skeptikeren, „der Zuchtmeister des dogmatischen
Vernünftlers“, der førte Kant til apriorisk Selvbesindelse og derved
til en sund Kritik af Forstand og Fornuft; og det var igjen, som
bekjendt, den rene Fornufts Kritik, der førte Fichte til Jeget som Jeg
og derved til en idealistisk Udvikling af Subjectivitetens uendelige
Selvgyldighed%
% printed footnote mark: *)
\footnote{Saa stor Betydning den kantiske Fornuftkritik end har for
Er\-%
% The Sperrsatz begins only after the printed line break: ``Er-'' at the end of
% the first line of the note is set solid, ``kjendelseslæren'' on the next line
% is letterspaced. Transcribed as printed; not regularised.
\emph{kjendelseslæren}, saa kan den dog i nærværende Sammenhæng kun
nævnes som Overgangsstadium fra Hume's Skepticisme til Fichtes
subjective Idealisme. Naar man har besindet sig paa, hvad den rene
Fornufts Kritik beviser, at der gives: aprioriske Sandseanskuelser,
aprioriske Forstandskategorier og aprioriske Fornuftideer, og hvad det
betyder, at en „Ding an sich“ som den yderste Reduction af en i sig
gjældende Objectivitet blev stillet i skarp Modsætning saavel til
Phænomenerne som til Subjectivitetens aprioriske Erkjendelsesformer,
saa er her i Grunden Intet videre at opholde sig ved; thi med noget
Forsøg paa at deducere den psychologiske Subjectivitet enten af et
Ikke-Subjectivt eller af en guddommelig Forudsætning har
Fornuftkritiken jo ikke indladt sig.}.
% --- p. 31 ---
\setcounter{footnote}{0}%
„I enhver Erfaring“, saaledes begynder den fichteske Tankegang, „haves
to Factorer satte mod hinanden, Jeget og Tingen, Intelligensen og dens
Gjenstand. Abstraherer man fra Jeget, saa faaer man \emph{en Ting i og
for sig}, og maa da opfatte Forestillingen som et Produkt af
Gjenstanden; abstraherer man fra Gjenstanden, faaer man \emph{et Jeg i
og for} sig, og maa da holde sig til den selvvirksomme, sit Indhold
% The letterspacing stops at the printed line break: ``et Jeg i og for'' is
% spaced, the ``sig,'' opening the next line is set solid. As printed.
frembringende Bevidsthed. Hiint er \emph{Dogmatismen}, dette
\emph{Idealismen}. Disse to ere indbyrdes uforenelige; et Tredie gives
ikke: her maa vælges. For at dømme Dogmatismen og Idealismen imellem,
mærke man sig Følgende: Jeget forekommer i Bevidstheden, Tingen i sig
er en Digtning af Bevidstheden. Idet Dogmatismen skal forklare
Forestillingens Oprindelse fra Tingen i sig, udøver den det mislige
Kunststykke at gaae bagved Bevidstheden og paakalde Væren; men Væren
virker kun Væren, ikke Bevidsthed; altsaa kan kun det Standpunkt være
sandt, der vælges, ikke i Væren, men i Bevidstheden, ikke i
Gjenstanden, men i Jeget selv.“

Havde Sensualismen som et Slags Realdogmatisme begyndt med at sætte
Eenheden af Tænken og Væren i Tænkningens umiddelbare Eenhed med
Sandsningen og det Sandselige (\textit{nihil est in intellectu, quod
non antea fuerit in sensu}), saa lægger nu den subjective Idealisme an
paa at opløse alt Sandseligt i Sandsningen, al Sandsning i Tænkningen
og al tænkelig Væren i Jegets Selvtænkning. „Indrømmer man“, saaledes
begynder den subjective Idealisme, „at Sætningen $A = A$ er
indlysende ved sig selv, da indrømmer man med det Samme Bevidsthedens
Evne til at sætte Noget. Man sætter ikke, at $A$ er, men kun, at,
dersom $A$ er, saa er $A$. Ville vi nu, istedenfor at see paa det
Satte, see paa den Sættende, eller rettere: paa Identiteten af den
Sættende og det Satte, da forvandler Sætningen: $A = A$ sig til: Jeg
$=$ Jeg. Medens vi istedenfor $A = A$ ikke
% --- p. 32 ---
\setcounter{footnote}{0}%
kunde sige, at $A$ er, kunne vi nu, istedenfor: Jeg $=$ Jeg, sige: Jeg
er, og i dette paa sig selv Grundede see Grunden til al Handling i den
menneskelige Aand, det oprindelige Særkjende paa
Bevidsthedsvirksomheden i og for sig. Ligesom Sætningen: $A = A$
betegner en Grundhandling, saaledes betegner Sætningen:
$\text{Non-}A$ ikke $= A$ ligeledes en Grundhandling, en efter sin
Form ubetinget Urhandling. Dog, da Sætningen $\text{Non-}A$ ikke
$= A$ forholder sig til Sætningen: $A = A$ som Antithesis til Thesis,
er den allerede herved efter sit Indhold, sin Materie, at opfatte som
betinget; thi Antithesis er nødvendig betinget af Thesis. Ved den
første Grundsætning: $A = A$, var Urhandlingens Form en
\emph{Sætten}, ved den anden er den en \emph{Modsætten}. Hvad
$\text{Non-}A$ er, vides endnu ikke; kun dette vides, at
$\text{Non-}A$ er Modsætning til $A$; men $A$ er sat ved Jeget; her er
kun forudsat et Jeg: det Jeget Modsatte er altsaa \emph{Ikke-Jeg}. Af
de to Grundsætninger følger med Nødvendighed en tredie. Opgaven for
den Handling, der udtrykkes ved den tredie Grundsætning, er given; men
ikke Opgavens Løsning. Paa den ene Side ophæves Jeget fuldstændig ved
Ikke-Jeget: Jeget kan ikke være sat, forsaavidt Ikke-Jeget er sat. Paa
den anden Side er Ikke-Jeget, sat af Jeget, selv et Jeg, sat i
Bevidstheden. Jeget ophæves altsaa ikke ved Ikke-Jeget; Jeget er
ophævet, og dog ikke ophævet: dette er Modsigelsen. Til at løse denne
Modsigelse maae vi da see at finde et $X$, hvori de hinanden
udelukkende Grundsætninger kunne forenes. Men hvorledes kunne Væren og
Ikke-Væren, Realitet og Negation forenes med hinanden uden at
tilintetgjøre hinanden? Ved gjensidig Indskrænkning! Det omspurgte
$X$ er altsaa en Skranke, og Selvindskrænkning Jegets Urhandling.“

Af disse første Grundtræk vil den subjective Idealismes logiske
Charakteer i det Hele være kjendelig. Principet er Jeget, det rene
Jeg, og Begyndelsen skeer dermed, at Jeget paa een Gang sætter,
forudsætter og deducerer sig selv. For\-%
% --- p. 33 ---
\setcounter{footnote}{0}%
holdet mellem Princip og Begyndelse er da med et ganske andet
Eftertryk udtalt i Fichtes rene Jeg end i Hegels rene Væren. At gjøre
det rene Jeg til Princip er væsentlig det Samme som at gjøre
Subjectiviteten til Princip; at begynde med det rene Jeg er med andre
Ord at begynde med den rene Viden. Hvoraf kommer det da, at den
fichteske Begyndelse ikke har kunnet tilfredsstille Systematikens
Fordring? hvori bestaaer da Fichtes oprindelige Feilgreb?

Fra tre Sider er Misligheden med „det fichteske Jeg“ iøinefaldende.
For det Første er her et tvetydigt Forhold mellem Jeget selv og dets
Existens eller, for at bruge de gangse Udtryk, mellem det rene og det
empiriske Jeg. Det empiriske Jeg, min og din virkelige Subjectivitet,
har sin Existens i sin uendelige Begrændsning: mit Jeg falder i denne
begrændsede Existens udenfor dit Jeg, min Tænken udenfor din Tænken,
min Viden udenfor din Vider
% PRINTER'S DEFECT: the print has ``Vider'' for ``Viden.'' --- the final sort is
% an r (calibrated against the certain ``Viden'' earlier in the same line, which
% shows the two stems of n), and the sentence-closing full stop is absent. All
% three witnesses read ``Vider''. Transcribed as printed.
I Existensen selv er her altsaa en existerende Mangfoldighed, en
Mangfoldighed af Jeg-Existenser; men hvor er den existerende Eenhed? I
Jegheden selv, det rene Jeg! Nei, Jegheden selv, det rene Jeg,
existerer ikke. Det rene Jeg er en Abstraction, som hvert existerende
Jeg for sig maa foretage. Den Abstraherende er nødvendig forud for sin
Abstraction; det empiriske Jeg maa altsaa i Virkeligheden være forud
for det rene Jeg. Men, ifølge Principet skulde det jo netop være det
rene Jeg, der var forud for det empiriske: her er en paafaldende Strid
imellem Princip og Virkelighed. For det Andet er Jegets oprindelige
Selvsætten, den Urhandling, hvorved Jeghedsprincipet sætter sig som et
virkeligt Jeg, en uopløselig Gaade. For at løse Gaaden maatte man da
allerførst faae en Besvarelse af det Spørgsmaal, om Jegets oprindelige
Selvsætten er fra Evighed af eller en Handling i Tiden. At lade
Urhandlingen være fra Evighed af er med andre Ord at tillægge Jeget en
mythisk-platonisk Præexistens; at lade den foregaae i Tiden er at lade
det Ubevidste afgive
% The foot of this page carries the signature mark ``Grundideernes Logik.'' + 3;
% a printer's mark, not transcribed.
% --- p. 34 ---
\setcounter{footnote}{0}%
en Forudsætning for det Bevidste, er at lade Væren gaae forud for
Viden. Lægges Urhandlingen ind i Evigheden, da kan det virkelige Jeg
kun erkjende sit Princip, forsaavidt det kan erindre sig en egen
Præexistens; men en saadan Præexistenserindring er i Virkeligheden
umulig. Drages Urhandlingen ned i Tiden, da gaaer der jo hos ethvert
Individ, i hvilket en Bevidsthed skal fremkomme, en Tilstand af Væren
forud for Bevidstheden; men „Væren virker kun Væren, ikke
Bevidsthed“; ved at lade en Bevidsthedstilstand udvikle sig af en
Værenstilstand har man i Grunden ophævet Principet
% PRINTER'S DEFECT: no full stop after ``Principet'' at the end of the printed
% line, though a new sentence begins with ``For det Tredie''. All three
% witnesses agree. Transcribed as printed.
For det Tredie er Bevidsthedens Urhandling, hvad enten man saa
forøvrigt vil henlægge den i Tiden eller i Evigheden eller i Grændsen
mellem Tiden og Evigheden, behæftet med den Grundmodsigelse at være en
ubevidst Bevidsthedsact, en absolut Videnshandling uden Viden%
% printed footnote mark: *). This note runs over to the foot of p. 35; the break
% falls after ``vom Denken des Ich abhängig''.
\footnote{Fichte har selv en Følelse heraf, men hjælper sig ved at lade
„Videnskabslæren“ gjøre Realismen en Indrømmelse. „Die
Wissenschaftslehre ist demnach \emph{realistisch}. Sie zeigt, daß das
Bewußtseyn endlicher Naturen sich schlechterdings nicht erklären
lasse, wenn man nicht eine unabhängig von denselben vorhandene, ihnen
völlig entgegengesetzte Kraft annimmt, von der dieselben ihrem
empirischen Daseyn nach selbst abhängig sind\dots{} Ohnerachtet ihres
Realismus aber ist diese Wissenschaft nicht transcendent, sondern
bleibt in ihren innersten Tiefen \emph{transcendental}. Sie erklärt
allerdings alles Bewußtseyn aus einem unabhängig von allem Bewußtseyn
vorhandenen, aber sie vergißt nicht, daß sie auch in dieser Erklärung
sich nach ihren eignen Gesetzen richte, und so weit sie hierauf
reflectirt, wird jenes Unabhängige abermals ein Produkt ihrer eignen
Denkkraft, mithin etwas vom Ich abhängiges, insofern es für das Ich
(im Begriff davon) da seyn soll. Aber für die Möglichkeit dieser neuen
Erklärung jener ersten Erklärung wird ja abermals schon das wirkliche
Bewußtseyn, und für dessen Möglichkeit abermals jenes Etwas, von
welchem das Ich abhängt, vorausgesetzt: und wenn jetzt gleich
dasjenige, was fürs erste als ein Unabhängiges gesetzt wurde, vom
Denken des Ich abhängig%
% (the note is continued at the foot of printed p. 35)
{} geworden, so ist doch dadurch das Unabhängige nicht gehoben,
sondern nur weiter hinausgesetzt. --- Alles ist seiner Idealität nach
abhängig vom Ich, in Ansehung der Realität aber ist das Ich selbst
abhängig; aber es ist nichts real für das Ich ohne auch ideal zu seyn;
mithin ist in ihm Ideal- und Realgrund eines und eben dasselbe, und
jene Wechselwirkung zwischen dem Ich und Nicht-Ich ist zugleich eine
Wechselwirkung des Ich mit sich selbst.“ Grundlage der
Wissenschaftslehre, Leipzig 1794, S. 261, 272, 273}. Den i
Urhandlingen skjulte
% --- p. 35 ---
\setcounter{footnote}{0}%
Modsigelse giver sig da ogsaa tydelig nok tilkjende derved, at der med
Jegets Selvindskrænkning danner sig en selvmodsigende Selvhed, en
Subjectivitet, som hverken kan undvære eller fordrage det Objective.
Forsaavidt Bevidsthedshandlingen skeer ubevidst, frembringes her i
dunkel Nødvendighed en Verden af Ting: og Tingene tage sig for den
umiddelbare Bevidsthed ud som udvortes tilværende Gjenstande, som
virkelige Objecter. Men hver Gang Jeget bliver sig sin egen ubevidste
Handling bevidst, forsvinde Objecterne; det viser sig da, at Tingene
ikke ere Ting, men Tilsyneladelser af et i det Uendelige virkende
Ikke-Jeg, Skranker i det sig selv indskrænkende Jeg. Den uophørlige
Strid mellem Jeget og Skrankerne er fra dette Synspunkt en uophørlig
Strid mellem det Subjective og det Objective, en oprindelig Strid i
Subjectiviteten selv.

\subsubsection*{c) Subjectivitetsproblemet.}

Modsætningen mellem en sensualistisk og idealistisk Subjectivitet er
qvalitativ, og forsaavidt umiddelbar; ikke desto mindre have de to saa
modsatte Synsmaader dog det tilfælles, at de begge søge
Subjectiviteten i den menneskelige Selvbevidsthed, og til et vist
Punkt ere enige i at bestemme den psychologisk. Men den psychologiske
Subjectivitet har igjen en qvalitativ og forsaavidt umiddelbar
Modsætning i den abstract logiske. Fra det psychologiske Synspunkt er
Subjectiviteten væsentlig sandsende, tænkende, vidende; fra det ab\-%
% The foot of this page carries the signature mark 3*; not transcribed.
% --- p. 36 ---
\setcounter{footnote}{0}%
stract logiske Synspunkt er Subjectiviteten derimod kun en almindelig
Begrebseenhed, den Begrebseenhed nemlig, hvori Subjectet har alle sine
Prædicater saaledes, at det selv er disse Prædicaters Indbegreb. Den
abstract logiske Subjectivitet er da med sin abstracte Form en
Subjectivitet i Sandsebestemmelser uden Sandsning, i
Tankebestemmelser uden Tanker, en selvløs Selvhed, et Subject, der kun
existerer som Object.

Hvor umuligt det er at løse Subjectivitetsproblemet, naar man eensidig
fastholder det psychologiske Udgangspunkt, er allerede viist ved
Bedømmelsen af Sensualismen og den subjective Idealisme; men hvor
umuligt det paa den anden Side er at løse Problemet, naar man eensidig
gaaer ud fra det almeenlogiske Synspunkt, sees, om muligt, endnu
tydeligere af de Misviisninger, hvortil den objective Idealisme, og da
i Særdeleshed den hegelske Panlogisme, allerede har ført. At begynde
med det Almeenlogiske er at begynde med en Bortseen fra alt
Psychologisk; det er at begynde med den almene Tænken i den objective,
almene Væren. Den objective, almene Væren udvikler sig da til det
objective, almene Væsen, som igjen udvikler sig til den objective,
almene Substans, som saa igjen udvikler sig til det objective, almene
Begreb, som derpaa strax begynder med at udvikle sig til den
objective, almene Subjectivitet. At nu en saadan objectiv
Subjectivitet, en selvløs Selvhed, der --- efter at have gjennemgaaet
„Slutningernes Mediation“ --- maa søge sin høiere Concretion i
Objecter som Guld og Jern, skulde ved sig selv bevæge sig videre frem,
først gjennem Planten, og saa igjennem Dyret, indtil den i Mennesket
omsider naaede sig selv i Selvbevidsthedens Form, er en urimelig
Paastand. Det er kun det ovenfor omhandlede frugtesløse Forsøg:
fuldstændig at ville aflede det Subjective af det Objective,
Bevidstheden af det Ubevidste, Viden af Ikke-Viden, vi her møde igjen
under en anden Vending.

% --- p. 37 ---
\setcounter{footnote}{0}%
„Subjectets Bestemmelse som Object“, siger J. L. Heiberg%
% printed footnote mark: *)
\footnote{Forelæsninger over Philosophiens Phil. eller den speculative
Logik 1831--32 S. 102--103. (Prosaiske Skrifter, 1ste Bind 1861, S.
322--325.)}, „er det vigtigste Punkt i hele Logiken. Det er det
Centrum, hvorom al Philosophie, ligefra de ældste til de nyeste Tider,
har dreiet sig; det er det Problem, som alle Philosopher, hver paa sin
Maade, have forsøgt at løse. I nærværende Fremstilling er det
Slutningen, hvorved Subjectet bestemmes som Object. Seet fra det blot
psychologiske Standpunkt har denne Bestemmelse ingen anden Betydning,
end at Virkeligheden maa svare til det Omdømme, som fremkommer ved en
Slutning, eftersom Slutningen er et Middel til Sandhedens
Erkjendelse, men Sandheden sættes i Begrebets Overeensstemmelse med
Virkeligheden, saa at, naar denne afviger fra Begrebet, forfeiles
Sandheden, hvilket ene kan reise sig deraf, at Slutningen har været
urigtig. Men denne Betragtningsmaade er høist eensidig og
indskrænket. Uden at tale om, at Sandheden er selve Begrebet, og at
Virkeligheden, kun forsaavidt som den ligeledes selv er Begrebet, er
sand Virkelighed, eller meer end Phænomen, saa at der ikke kan være
Spørgsmaal om Begrebets Overeensstemmelse med hvad der kun er det
selv, saa forudsættes i denne Bestemmelse Objectiviteten som noget
Givet, hvilket Subjectet skal bemægtige sig, og i denne Modsætning
bliver da Subjectet bestemt som den menneskelige Erkjendelse. Men
seet, ikke fra det psychologiske, men fra det logiske Standpunkt,
faaer Subjectets Bestemmelse som Object en langt dybere og meer
omfattende Betydning. Logiken forudsætter nemlig ikke Objectet, men
lader Subjectet selv bestemme sig dertil, ligesom den heller ikke
bestemmer Subjectet som et menneskeligt, erkjendende Subject, men
snarere som det gudommelige skabende\dots. Maaden, hvorpaa
Philosopherne have opfattet Subjectets Forhold til Objectet,
charakteriserer de forskjellige Systemer. Alle idealistiske, saavel
ældre, som
% --- p. 38 ---
\setcounter{footnote}{0}%
nyere Philosopher have betragtet Subjectet som identisk med Objectet,
men have, ligefra Plato til Schelling, forudsat denne Eenhed, medens
den derimod skal udvikle sig, nemlig som Ideen. Men mellem de
Philosopher, som, istedenfor at lægge denne Eenhed til Grund, have
betragtet enten Subjectet eller Objectet som det Absolute, ere Fichte
og Leibnitz de, som skarpest have fremstillet Modsætningen, hiin ved
at bestemme Subjectet som det Absolute, og Objectet kun som Grændse
for Subjectet; denne derimod ved at bestemme Objectet som det
Absolute, nemlig som Monade, d. e. selve Substansen, betragtet som det
i objectiv Henseende Ene og Udeelbare, men indbefattende
Subjectiviteten i Form af \emph{Forestillings-Evnens forskjellige
Grader}. Da Monaderne, ligesom Atomerne, udgjøre en Fleerhed, og ere
selvstændige, saa ophæver Leibnitz denne Bestemmelse ved de nye
Bestemmelser af en Centralmonade eller \emph{Monadernes Monade, og af
den forudbestemte Harmonie}. Saavel Fichte som Leibnitz ere
Idealister, men hiin er subjectiv, denne objectiv Idealist. I
nærværende Fremstilling er Idealismen begge Dele. Her, hvor Subjectet
bestemmes som Object, er den subjectiv, siden, naar Objectet vender
tilbage til Subjectet (i Ideen), er den objectiv. Det er den første af
disse Bestemmelser, som vi her have at betragte. Omendskjøndt nu denne
vel kan siges at være subjectiv, saa er det dog paa en ganske anden
Maade end Fichtes udelukkende subjective Idealisme; thi den modsatte
Bestemmelse ligger allerede deri, at det er det samme Begreb, der
baade er Subject og Object, saa at Begrebet allerede sees som deres
Eenhed, d. e. som Ideen. Saaledes sees da Objectet ikke her som en
blot Indskrænkning af Subjectets Virksomhed, en Indskrænkning, hvorom
det er tvivlsomt, enten man helst skal betragte den som et Produkt af
Subjectet, eller (ligesom paa det psychologiske Standpunkt) som en
Forudsætning, hvilken Subjectet ikke kan overvinde og bemægtige sig,
da i begge Tilfælde Objectet kun er bestemt som det Indiffe\-%
% --- p. 39 ---
\setcounter{footnote}{0}%
rente, som Grændse, og forsaavidt som Subjectets Qvalitet; tvertimod
er i nærværende Fremstilling Objectet den Væren, som Subjectet, efter
iforveien at have optaget den, nu udvikler af sig selv, og altsaa paa
engang baade Subjectets Forudsætning og dets Produkt. Fremdeles er
Objectet ikke Subjectets Grændse eller Qvalitet; i denne Betydning er
det endnu bestemt som et Andet for Subjectet, og er kun Prædicat i det
umiddelbare Omdømme; men idet Omdømmets Umiddelbarhed er medieret ved
Slutningen, og denne Mediation igjen er hævet ved Systemet af
Slutninger, og vendt tilbage til Umiddelbarheden, saa er Prædicatet
blevet til Object, og Objectet er i dette Forhold den hele objective
Verden, som har samme Realitet som Subjectet, da det ikke er et Andet
for dette, men staaer i Umiddelbarhedens Forhold dertil. Subjectets
Bestemmelse som Object har altsaa her den Betydning, at al Væren
udgaaer af Tanken, og ikke er andet end den synlige Tanke.“

Til en Charakteristik af Subjectivitetsproblemet, seet fra et hegelsk
Standpunkt, er ovenstaaende Udtalelse en af de klareste og meest
uforbeholdne, man i denne Sammenhæng kan ønske sig; vi have her
optaget den i hele dens Udførlighed, fordi den frembyder just
saadanne Hovedpunkter, som ere nødvendige, naar det ret skal blive
tydeligt, hvad det er, hvorom her egentlig handles. Forskjellen
imellem det psychologiske og det logiske Standpunkt sættes altsaa
deri, at hiint fordrer, at det Omdømme, der fremkommer ved en
Slutning, skal svare til Virkeligheden, dette derimod lader Subjectet
selv bestemme sig til Object. Hvorledes et Omdømme fremkommer paa det
psychologiske Bevidsthedstrin, er en bekjendt Sag; men hvorledes
fremkommer det i „den speculative Logik“? Omdømmet kan jo dog umulig
enten dømme sig selv til at staae i Logiken eller dømme Logiken til at
skaffe Subjectet det Prædicat, det behøver for at bestemme sig selv
til Object. Kan her ikke være Tale om „Begrebets
\emph{Overeensstemmelse med}
% The letterspacing and the quotation both run on across the page break into
% printed p. 40; the \emph{} is closed here only because the fragment ends, and
% the opening „ is left unmatched for the same reason.
% --- p. 40 ---
\setcounter{footnote}{0}%
hvad der kun er det selv“, hvordan kan her da være Tale om en
Begrebets \emph{Overgang} til, hvad der kun er det selv? At Subjectet,
idealistisk seet, er identisk med Objectet, maa indrømmes. Men
Identiteten har jo Forskjel; Forskjellen mellem Subject og Object er
meer end Forskjel; den er en udtrykkelig Modsætning, ja, en i
Begrebet totaliseret Modsætning. Men hvordan gaaer det saa til, at
Subjectet og Objectet kunne i „den speculative Logik“ ved lidt
Mediation ombyttes saaledes, at hvad der den ene Gang er Subject,
ganske og aldeles Subject, er den anden Gang Object, ganske og
aldeles Object? Det siges udtrykkelig, at Objectet, der har samme
Realitet som Subjectet, „ikke er et Andet for dette, men staaer i
Umiddelbarhedens Forhold dertil“. Hvordan kan Subjectet da fra sin
Side staae i et Umiddelbarheds Forhold til det, der i Realiteten
„ikke er et Andet for det“, ikke er Andet end det selv? Hvad er det
for et Subject, hvorom her tales?

Det hedder, at Logiken „ikke bestemmer Subjectet som et menneskeligt
erkjendende Subject; men snarere som et guddommeligt skabende“.
Hvormange Subjecter have vi da egentlig her for os? Her er et
menneskeligt erkjendende Subject, der udelukkes af Logiken; et
guddommeligt Subject, der anerkjendes af Logiken; og saa et
udelukkende og anerkjendende Subject, det Subject, der bestemmer
„Subjectet“, ikke som menneskeligt, men som guddommeligt; dette
„Subjectet“ bestemmende Subject er Logiken selv; altsaa er her i Alt
tre Subjecter. Men hvordan forholde disse tre Subjecter sig nu
egentlig til hinanden? Hvilket af dem er det uafhængige, og hvilke
ere de afhængige? Ved første Øiekast skulde man antage, at „det
guddommeligt skabende Subject“ maatte være det uafhængige; thi at et
Subject ved sig selv formaaer at bestemme sig til Object, tyder
unegtelig paa Magtfuldkommenhed og stor Uafhængighed. Men nærmere
beseet, kan det guddommeligt skabende Subject dog kun bestemme sig
til Object,
% --- p. 41 ---
\setcounter{footnote}{0}%
forsaavidt Logiken „lader det bestemme sig dertil“; men er Logiken
virkelig det Subject, der bestemmer „det guddommeligt skabende
Subject“, saa maa Logiken jo være det mægtige, Alt formaaende
Subject, altsaa det uafhængige! Dog heller ikke denne Antagelse vil
ret tilfredsstille. Vistnok kan man ved en Forkortelse i
Udtryksmaaden tale om, hvad Logiken gjør, i Modsætning til, hvad
Mathematiken, Physiken og de øvrige Videnskaber, hver fra sit
Synspunkt, maae og bør gjøre: man kan, med andre Ord, meget vel
opfatte de forskjellige Videnskaber som tænkende, villende og
handlende Subjecter, uden at derved fremkaldes nogen Misforstaaelse,
naar Forkortelsen blot ikke anbringes saaledes, at det figurlige
Udtryk tages bogstavelig, og det just i en Sammenhæng, hvor det
gjælder om at faae Begrebet „Subject“ bestemt med logisk
Nøiagtighed. Subjectet i Logiken kan ganske vist ikke røre sig af
Stedet uden Logikens allerhøieste Tilladelse, kan ikke bestemme sig
selv enten til Prædicat eller til Object, uden at Logiken staaer det
bi med Raad og Daad, befordrer det fra en Ansættelse til en anden,
og ved stadig Forfremmelse bestemmer det til at bestemme sig for
enhver conseqvent og væsentlig Bestemmelse.

Men er Logiken, „den speculative Logik“, dette alle logiske
Subjecter, Prædicater, Objecter og Subject-Objecter bestemmende
Subject da ogsaa virkelig et logisk Selv, et selvstændigt Væsen, en
tænkende, dømmende, vidende Subjectivitet? Dersom vi antoge dette,
saa maatte vi vel ogsaa antage, at „den speculative Logik“ kunde
udvikle sig selv; men dersom den virkelig kunde udvikle sig selv,
saa maatte den vel ogsaa kunne fremstille sig selv; men dersom den
virkelig kunde fremstille sig selv, saa maatte den vel ogsaa kunne
foredrage sig selv, og da sagtens ogsaa skrive sig selv, befordre
sig selv til Trykken o. s. v. Men ved denne „immanente Udvikling“
af Logikens logiske Subjectivitet ere
% --- p. 42 ---
\setcounter{footnote}{0}%
vi jo allerede komme ud over Logiken som Subject til et Logiken med
dens Subjecter oversvævende Subject, nemlig til det Subject, der paa
en Maade bliver overseet, idet det bliver „bestemt som den
menneskelige Erkjendelse“.

Istedenfor at afføre sig den menneskelige Subjectivitet med den
menneskelige Erkjendelse, maa den tænkende Bevidsthed just lægge an
paa at finde Gaadens Løsning i sit eget subjective Væsen. Kan denne
Bevidsthed erkjende sin Endelighed og psychologiske Begrændsning, da
er denne Erkjendelse just et Beviis for, at den med det Samme er
hævet over Endeligheden og den blot psychologiske Begrændsning; kan
den i al sin menneskelige Ufuldkommenhed erkjende Ufuldkommenheden,
da er den ved denne Erkjendelse allerede i sig selv hævet over sig
selv, og kan igjennem det Ufuldkomne opdage det Fuldkomne, kan i det
Menneskelige see det Guddommelige.

Imellem den psychologiske og den almeenlogiske Subjectivitet maa her
absolut, dersom Problemet skal lade sig løse, kunne paavises et
principielt Mellemled, en Grundsubjectivitet, hvoraf det Subjective
i alle relative Subjectiviteter er at aflede, en Princip-Viden,
hvortil al psychologisk Bevidsthed og endelig Viden kan føres
tilbage. Den Tanke ligger da nær, at gjøre Grundsubjectiviteten til
Eet med Guds Aand og aflede den menneskelige Viden af den
guddommelige, Menneskets Selvbevidsthed af Guds Selvbevidsthed:
Mennesket er jo skabt i Guds Billede. Ikke desto mindre vilde en
saadan umiddelbar Opstilling af Modsætningen mellem Guds og
Menneskets Subjectivitet kun bidrage til at forøge Forvirringen. Ved
nemlig at indføre Gud selv som videnskabelig Forudsætning indføre vi
en overnaturlig Factor i Videnskaben, et dogmatisk Princip, der gjør
al videnskabelig Begrundelse umulig.

Vel kan man fra et „theocentrisk Standpunkt“ med et vist Skin af
Grundighed raisonnere saaledes:

% --- p. 43 ---
\setcounter{footnote}{0}%
% The page opens the quoted raisonnement on a fresh line, but the print
% sets that line FLUSH, without the paragraph indent it gives every other
% new paragraph on pp. 40--53; \noindent reproduces the printed setting.
\noindent „Vi Mennesker vide jo vistnok umiddelbart, hvad det er at
vide Noget, at det udtrykkelig er at fornemme, at føle, i det Hele
at sandse Noget, at anskue Noget, at forestille sig Noget og fremfor
Alt at tænke Noget; men vi vide tillige, at al vor Viden er
indskrænket og har en Begyndelse i Tiden. Da nu det Bevidste umulig
kan afledes af det Ubevidste, Viden umulig afledes af Ikke-Viden,
saa maa vor indskrænkede Viden nødvendigviis forudsætte en
uindskrænket og fuldkommen Viden; vor i Tiden og under timelige
Betingelser vaagnende Aand kan ikke være et Produkt af Materiens
blinde Elementer; den maa nødvendig have en evig, over Tid og
Timelighed ophøiet Aand, altsaa Gud selv, til sit Ophav.“ Ved et
saadant Raisonnement synes den psychologisk begrændsede
Subjectivitet jo vistnok ved første Øiekast at ville hæve sig over
sin egen Indskrænkning og ligesom gaae bagved sig selv, for at søge
sin Kilde; men den hele Bevægelse er kun tilsyneladende. Den
raisonnerende Subjectivitet vil, naar vi see nærmere til, i
Virkeligheden blive staaende indenfor endelige Skranker og midt i
Endeligheden vedblive at maale det Uendelige med sin egen
psychologiske Alen.

Hvorledes kommer vel den raisonnerende Subjectivitet til Indsigt i
den guddommelige Viden? \textit{Via negationis}: ved at borttage
alle Skranker; \textit{via eminentiæ}: ved at tænke sig vor Videns
ufuldkomne Former hævede til absolut Fuldkommenhed. Ved \textit{via
negationis} vises vi da hen til en abstract qvalitativ Modsætning af
Endeligt og Uendeligt, Begrændset og Ubegrændset; det Ubegrændsede
maa altsaa blive kjendeligt paa Grændsen, ligesom det Uendelige paa
det Endelige. Ved at negere det Endelige i Bevidstheden kommer man
ganske rigtig til en Bevidsthed, man maa kalde uendelig; men en
Bevidsthed, der er uendelig i den Forstand, at den intet endelig
bestemt Indhold har, er tom. Ved at hæve Videns Grændser kommer man
unegtelig til en Viden uden Grændser; men en Viden uden alle
Grændser, altsaa uden Bestemthed,
% --- p. 44 ---
\setcounter{footnote}{0}%
er en ubestemmelig Viden af et i sig Ubestemmeligt: τὸ ἄπειρον er i
og for sig et ἄλογον.

For at bøde paa den grændseløse Tomhed, hvortil \textit{via
negationis} conseqvent maa føre, slaaer den raisonnerende
Subjectivitet dernæst ind paa \textit{via eminentiæ}; Analogien
mellem den menneskelige og den guddommelige Viden fastholdes: den
guddommelige Viden svarer til den menneskelige, som det Fuldkomne
til det Ufuldkomne. Men nu beroer den menneskelige Viden jo ikke
blot paa en Tænkning, men ogsaa paa en Sandsning. Uden Syn og
Hørelse vilde vor Erkjendelse af Tingene i Verden være deels umulig,
deels ufuldstændig. Til en fuldkommen Viden vil der altsaa kræves et
fuldkomment Syn og en fuldkommen Hørelse. Fra et theocentrisk
Standpunkt har man da heller ikke taget i Betænkning at tillægge Gud
baade Syn og Hørelse. Men hvordan forholder det sig nu med det
guddommelige Syn og den guddommelige Hørelse?

At Religionen i sit eiendommelige Billedsprog taler om Guds Øie,
Guds Øre o. s. v., volder os i denne Sammenhæng ingen Bryderier,
eftersom Religionen ikke er og ikke vil være Videnskab. Men
Theologien vil jo være Videnskab; Theologien maa vel altsaa see sig
istand til at give os et videnskabeligt Begreb om Guds Alvidenhed:
hvilke Bidrag har Theologien da ydet til en dybere Forstaaelse af
Læren om Guds Øie og Guds Øre og de Vidensfuldkommenheder, som
hertil svare? Ingen, slet ingen! At erklære de billedlige Udtryk for
Anthropomorphismer og lade Eftertrykket falde paa den rene
Alvidenhed er en let Sag; men Spørgsmaalet er: hvad bliver det til
med den rene Alvidenhed, naar man gjør Alvor af at fjerne alle
saakaldte Anthropomorphismer?

Lader os betragte Alvidenheden nærmere. Den guddommelige Alvidenhed
maa naturligviis være, ikke blot en Viden om alle Ting, men ogsaa en
Viden af Tingenes sande Væsen: saaledes som Tingene vides af den
Alvidende, saa\-%
% --- p. 45 ---
\setcounter{footnote}{0}%
ledes og kun saaledes ere de i sig selv. Heraf følger, at, dersom
der virkelig existerer en Verden af sandselige Ting, saa maa der i
Alvidenheden selv være en sandselig Viden; thi kun ved en sandselig
Viden kan det Sandselige vides. Ere de sandselige Ting synlige,
følelige, hørlige, saa maa den tilsvarende sandselige Viden være
betinget af en Seen, Følen, Høren o. s. v.

Men hvilken Forskjel maa her nu ikke være mellem den menneskelige og
den guddommelige Seen, mellem den menneskelige og den guddommelige
Høren, kort sagt: mellem den menneskelige og den guddommelige
Sandsning! For det menneskelige Syn er der en uafviselig Modsætning
mellem Lys og Mørke, mellem Hvidt og Sort, mellem Rødt og Gult,
mellem gjennemsigtige og uigjennemsigtige Legemer; denne Modsætning
er afgjørende saavel for det Subjective i Synet, som for det
Objective i de synlige Ting. Men Modsætningen viser tillige en
væsentlig Indskrænkning, en Ufuldkommenhed i vor Seen. Thi det er jo
en Ufuldkommenhed ved vort Syn, at vi ikke see, hvad der er i det
Fjerne, ligesaa tydeligt, som hvad der er i vor
% printed footnote mark: *)
Nærhed\footnote{For at der kan raades Bod paa denne Ufuldkommenhed,
maa Dantes Syn i Paradiset undergaae mangfoldige Forvandlinger.
Saaledes i 31te Sang, hvor han seer Beatrice paa Thronen
\textit{nel terzo giro dal sommo grado}:
% The Italian verse is set in antiqua, indented, one line to a line.
% In line 1 the grave of `sù' is set displaced to the right of the u.
\begin{quote}
\textit{Da quella region, che più sù tuona}\\
\textit{Occhio mortale alcun tanto non dista}\\
\textit{Qualunque in mare più giù si abbandona,}\\
\textit{Quanto li da Beatrice la mia vista;}\\
\textit{Ma nulla mi facea; chè sua effige}\\
\textit{Non discendeva a me per mezzo mista.}
\end{quote}}; det er en Ufuldkommenhed, at vi ikke see Tingene
ligesaa klart om Natten som om Dagen; en Ufuldkommenhed, at vi ikke
ligesaa let kunne see igjennem tykke Mure som igjennem Vindues\-%
% --- p. 46 ---
\setcounter{footnote}{0}%
ruder. Slige Ufuldkommenheder maae altsaa \textit{via eminentiæ}
forsvinde for den Alvidendes Øie.

Men idet Ufuldkommenhederne forsvinde, forsvinde jo med det Samme
alle tilsvarende saavel objective som subjective Modsætninger; den
Alvidendes Øie maa altsaa see alle synlige Ting i lige Nærhed og
lige Fjernhed, lige klare, lige gjennemsigtige, altsaa: uden Skygge,
uden Farve, uden Omrids, uden Skikkelse; for det altseende Øie løbe
de synlige Ting da, reentud sagt, ud i Eet: det altseende Øie lukker
sig; den synlige Verden forsvinder. Til lignende Resultater komme vi
ved at undersøge hvilkensomhelst af de øvrige Sandsninger og
Sandsebestemmelser.

Vel bør vi i denne Sammenhæng heller ikke forglemme, at der er en
speculativ Dogmatik, som ved den Hellig Aands Bistand --- „idet vi
tillægge \textit{testimonium spiritus sancti} ikke blot praktisk,
men ogsaa theoretisk Betydning“ --- randsager Guds Dybheder,
idetmindste „til en vis Grad“ (thi dersom den Hellig Aand virkelig
vilde tage sig af Videnskaben, saa var der vist ingen Tvivl om, at
det jo blev i Dogmatiken, vi kom til at søge
Subjectivitetsproblemets egentlige Løsning); men de halvsande
Bemærkninger, Dogmatiken i denne Anledning kan anføre som Tegn paa
de mange Vidnesbyrd, Aanden her skal have givet sig saavel i „den
christelige Videnskab“ som i „den christelige Verdensanskuelse“ og i
„den christelige Kunst“, beroe tilhobe paa en Sphæreforvexling, idet
Christendommen selv forvexles med det Ideerige, hvori den har fundet
en Afspeiling, og for hvis Udvikling den som verdenshistorisk Magt
har afgivet væsentlige Betingelser.

Havde den „Guds Dybheder randsagende“ Dogmatik nu virkelig saamegen
religiøs Inderlighed, at den var istand til at opstille og behandle
Spørgsmaalet om den evige Salighed som et afgjørende
Livsspørgsmaal, da vilde det Subjectivitetsprincip, den i saa Fald
maatte anvende, jo være saa existentielt subjectivt, at det alene af
den Grund ikke kunde tilegnes af
% --- p. 47 ---
\setcounter{footnote}{0}%
Videnskaben. Denne Vanskelighed bortfalder imidlertid af sig selv.
Thi den objectivt bevægede Dogmatik har ikke engang saamegen ethisk
Inderlighed, at den, selv om den vilde nedlade sig til at benytte
Philosophiens Efterladenskaber, og udstoffere sig med et imponerende
Lærdomsapparat, fra sit objective Standpunkt var istand til at
behandle Spørgsmaalet om Sjælens Udødelighed med noget synderligt
% printed footnote mark: *) -- the note runs on to p. 48, where it closes
Eftertryk\footnote{„Spørgsmaalet om Udødelighed er altsaa et lærd
Spørgsmaal? Ære være Lærdom, Ære være den, som lærd kan behandle det
lærde Spørgsmaal om Udødeligheden; men Spørgsmaalet om
Udødeligheden er væsentlig ikke noget lærd Spørgsmaal, det er et
Inderlighedens Spørgsmaal, som Subjectet ved at blive subjectivt maa
gjøre sig selv. Objectivt lader Spørgsmaalet sig slet ikke besvare,
fordi objectivt lader der sig ikke spørge om Udødeligheden, da
Udødeligheden netop er den udviklede Subjectivitets Potensation og
høieste Udvikling. Først ved ret at ville blive subjectiv kan
Spørgsmaalet ret komme frem, hvorledes skulde det da kunne besvares
objectivt? Socialt lader Spørgsmaalet sig slet ikke besvare, fordi
det socialt ikke lader sig fremsætte, da kun det Subject, der vil
blive subjectivt, kan fatte Spørgsmaalet, og spørge rigtigt: bliver
\emph{jeg} eller er \emph{jeg} udødelig? See flere Ting kan man godt
slaae sig sammen om; saaledes kan flere Familier slaae sig sammen om
en Loge i Theatret, og tre eenlige Herrer slaae sig sammen om en
Ridehest, saa hver rider hver tredie Dag. Men saaledes er det ikke
med Udødeligheden, Bevidstheden om min Udødelighed tilhører mig
ganske alene; netop i det Øieblik, jeg er mig min Udødelighed
bevidst, er jeg absolut subjectiv, og jeg kan ikke blive udødelig i
Compagnie med to andre eenlige Herrer paa Omgang. Subskribentsamlere,
der tilveiebringe en talrig Subskription af Mænder og Koner, der føle
en Trang i Almindelighed til at blive udødelige, erholde heller ingen
Fordele for deres Uleilighed, thi Udødelighed er et Gode, som ikke
lader sig tiltrodse ved en talrig Subskription. Systematisk lader
Udødeligheden sig heller ikke bevise. Feilen ligger ikke i
Beviserne, men i at man ikke vil forstaae, at systematisk seet er
det hele Spørgsmaal Nonsens, saa man istedenfor at søge yderligere
Beviser, snarere maa see at blive lidt subjectiv.
% the note continues here from the foot of p. 47 on to the foot of p. 48
Udødeligheden er Subjectivitetens meest lidenskabelige Interesse, i
Interessen ligger netop Beviset; naar man, systematisk ganske
conseqvent, systematisk abstraherer objectivt fra den: Gud veed,
hvad saa Udødelighed er, eller blot hvad det vil sige, at ville
bevise den, eller blot hvad det er for en fix Idee at bryde sig
videre derom. Om man systematisk kunde faae hængt en Udødelighed op,
ligesom Geßlers Hat, for hvilken vi Alle i Forbigaaende toge Hatten
af, det er ikke at være udødelig, eller at være sig Udødeligheden
bevidst. Systemets utrolige Møie med at bevise Udødeligheden er
spildt Uleilighed og en latterlig Modsigelse: at ville systematisk
besvare det Spørgsmaal, der har den Mærkelighed, at det ikke lader
sig systematisk fremsætte. Det er ligesom at ville male Mars i den
Rustning, der gjør ham usynlig. Pointen ligger i Usynligheden, og
ved Udødeligheden ligger Pointen i Subjectiviteten og i
Subjectivitetens subjective Udvikling.“

Afsluttende uvidenskabelig Efterskrift til de philosophiske Smuler
af Joh. Climacus. Udgiven af S. Kierkegaard, Kjøbenhavn 1846. S.
126--127.}.
% --- p. 48 ---
\setcounter{footnote}{0}%
Derimod har den med den Hellig Aands \textit{testimonium} sig
anbefalende Dogmatik det Betænkelige ved sig, at den, trods sit
videnskabelige Tilsnit, overalt paa Videnskabens Enemærker gererer
sig, som om den stod i et født Overlegenhedsforhold til Videnskaben
selv; den synes i det Hele at betragte sine videnskabelige
Undersøgelser som en Art diplomatiske Forhandlinger, om hvilke den
veed, at de aldrig kunne sætte den som aandelig Stormagt i nogen
alvorlig Forlegenhed; skulde den virkelig her og der geraade i nogen
Forlegenhed, kan den jo altid i Egenskab af Stormagt forevise
\textit{testimonium} og bryde af ved et Magtsprog.

Hvorledes det i videnskabelig Henseende forholder sig med
Dogmatikens foregivne \textit{testimonium}, kan her blandt Andet
sees af dens Behandling af Guds Alvidenhed. Skulde Testimoniet
virkelig sætte den istand til at randsage den guddommelige Videns
Dybder, da maatte den dog vel være istand til at give et nyt og
ganske eiendommeligt Bidrag til at
% --- p. 49 ---
\setcounter{footnote}{0}%
bestemme Begrebet Viden, og belyse os Videns Forhold til den
Vidende. Dette dens eiendommelige Bidrag maatte da være af den
Beskaffenhed, at enhver redelig Tænker kunde indrømme, at det nye
Vidensbegreb virkelig stod i den nøieste Sammenhæng med
Forudsætningerne: Aabenbaring og Tro. Fremdeles maatte det kunne
gjøres indlysende, at den nye paa Aabenbaringstroen byggede Viden var
i sin Art høiere og herligere end den, hvormed Videnskaben, uden Tro
og Aabenbaring, ellers maa lade sig nøie. Først naar dette ret var
blevet indlysende, kunde her være Tale om Modsætningen mellem et
troende og et ikketroende Standpunkt indenfor Videnskaben selv. De,
som forkastede det høiere, af Aabenbaring og Tro betingede,
Vidensbegreb, maatte da, saasandt de vilde gaae videnskabeligt
tilværks, tilstaae, at Vidensbegrebet, for sig betragtet, virkelig
var det høiere; men at de forkastede det, alene fordi de i deres
Vantro ikke kunde fastholde de Forudsætninger, hvorpaa dets Realitet
kom til at hvile.

Nu er der jo for enhver Dogmatik, og da især for den speculative, en
ganske særegen Opfordring til at udfolde et saadant høiere
Vidensbegreb, naar Talen er om Guds Alvidenhed. Hvad har da „den
Guds Dybheder randsagende“ Dogmatik at lære os om den alvidende Gud?
Det er ikke uden en vis Spænding, at vi lytte til en Meddelelse som
denne: „Den alvidende Gud er den sig selv aabenbare Gud, den Gud,
for hvem Alt er aabenbart“. Man venter her at see Begrebet
Aabenbaring udbrede et høiere Lys over Alvidenhedsbegrebet. Men naar
dette høiere Lys nu alene bestaaer i den dogmatiske Oplysning, at
man „fra gammel Tid har forestillet Gud under Billedet af et Øie“,
og det med Grund, efterdi „Gud ikke har et Øie“, men er Øie, efterdi
„hans Væsen er Viden“, saa bliver man tvivlraadig, og har
Vanskelighed med at udfinde, hvori Oplysningen egentlig skal
bestaae. Er det nu Dogmet: Gud er Øie, der skal
% --- p. 50 ---
\setcounter{footnote}{0}%
forklare os Dogmet: Gud er Viden, saa kan man jo heri gjerne see et
Beviis for, at der virkelig maa tillægges \textit{testimonium
spiritus sancti}, ikke blot praktisk, men theoretisk Betydning; thi
hvor ofte der end i den hellige Skrift tales om Guds Øie, ja, om
Guds Øine, saa skal man dog neppe finde noget Sted, hvor enten det
gamle eller det nye Testamente har hævet sig til den speculative
Viden, at Gud ikke har et Øie, men er Øie.

Men hvori finder nu en Sætning som denne: „Gud \emph{har} ikke et
Øie, men han \emph{er} Øie“ sin videnskabelige Begrundelse? Er den
maaskee ved sig selv indlysende for Alle, eller dog for de Troende?
Er den indlysende for alle Troende eller kun for dem, der ikke lægge
saamegen Vægt paa, hvad Bibelen lærer om Guds Øie, altsaa om det
Øie, Gud \emph{har}, som paa den Omstændighed, at man fra gammel Tid
har forestillet Gud under Billedet af et Øie, altsaa ---
speculativt! --- af det Øie, Gud \emph{ikke har}? Er det maaskee
denne sidste praktiske Omstændighed, hvori den theoretiske Betydning
af \textit{testimonium spiritus sancti} særlig maa søges?

„Men“, vil en billigttænkende Læser kunne indvende, „hvorfor nu just
hænge sig i et saa prægnant og aandrigt Udtryk som: „Gud er Øie“?
Betydningen er jo aabenbart den, at hans Væsen er Viden!“ Herpaa
have vi følgende Svar: saasnart det blot indrømmes, at den hele
„dogmatiske Oplysning“ løber ud paa en aandrig Omskrivning af, hvad
der allerede er sat med Alvidenheden selv som Forudsætning, saa
opgive vi strax det prægnante Udtryk. Den alvidende Gud er altsaa,
ifølge Dogmatikens høiere Vidensbegreb, den Gud, „hvis Væsen er
Viden“, den Gud, „hvis Væren er forklaret i evig Tænken“, „den sig
selv aabenbare Gud, for hvem Alt er aabenbart“; den Hellig Aands
theoretiske \textit{testimonium} løber altsaa gjennem adskillige
prægnante Omskrivninger ud paa en dogmatisk Retfærdiggjørelse af den
formelle
% --- p. 51 ---
\setcounter{footnote}{0}%
Logiks bekjendte Grundsætning: enhver Ting er, hvad den er,
\emph{den Alvidende er den Alvidende!}

Dog, thi heller ikke dette tør vi oversee, behøver „den dogmatiske
Oplysning“ just ikke at bestaae i nominelle Definitioner af
Alvidenheden; den „Guds Dybheder randsagende“ Dogmatik skal have Lov
til at anvende saamange prægnante, billedlige, halvbilledlige,
forbilledlige, ubilledlige Omskrivninger af Alvidenhedsbegrebet, det
skal være, naar den saa i Kraft af den „theoretiske Betydning“, der
bliver at tillægge \textit{testimonium spiritus sancti}, blot vil
forhjælpe os til en Løsning af de ovenfor angivne Problemer.
Hvorledes er det muligt, at endelige Forskjelle og sandselige
Modsætninger, saasom Modsætning af Dag og Nat, af Lys og Mørke,
kunne bestaae for den uendelige Viden, „den altgjennemtrængende
Seen“? --- Saaledes som Tingene vides af den Alvidende, saaledes ---
vi gjentage det --- saaledes ere de i sig selv. At den Alvidende maa
kjende baade Lys og Mørke, baade Sandseligt og Usandseligt, dersom
disse Modsætninger have sand Gyldighed, og dersom der virkelig er en
Alvidende, er ikke Problemet; thi det er Noget, der følger af sig
selv, og kan indsees uden „theoretisk“ \textit{testimonium} af den
Hellig Aand. Nei, Problemet er dette: hvorledes er det muligt, at
der kan være en alvidende Gud, dersom der i Sandhed er en endelig,
qvalitativ og væsentlig Modsætning mellem Lys og Mørke, mellem
gjennemsigtige og uigjennemsigtige Legemer? eller: dersom der i
Sandhed er en endelig, qvalitativ og væsentlig Modsætning mellem
gjennemsigtige og uigjennemsigtige Legemer, hvorledes er det da
muligt, at en Alvidenhed kan tænkes at være? At Problemer af denne
Natur ikke fremkunstles med Vilkaarlighed, men paatrænge sig med
Nødvendighed, kan sees af et hvilketsomhelst Exempel. Hvad er
Lyset? Object og Medium for Synet, altsaa for den seende Viden. Og
Mørket? En Skranke for Synet, Object for en ikke seende Viden,
altsaa for en Viden om, at „en altgjennem\-%
% --- p. 52 ---
\setcounter{footnote}{0}%
trængende Seen“ ikke er mulig! Hvordan kan en „altgjennemtrængende
Seen“ da paa een Gang baade see i Mørket og dog see Mørket? Hvordan
kan Alvidenheden, for at anføre et haandgribeligt Exempel, bestaae
med en Kampesteen? Seet i Alvidenhedens Lys, maa den for os
uigjennemsigtige Steen jo blive fuldkommen gjennemsigtig, og alle
dens materielle Egenskaber forvandle sig til Vidensbestemmelser,
Reflexer i den rene Viden; men paa denne Maade forsvinder Stenen som
Steen. Skal Stenen forblive og beholde sine materielle Egenskaber:
Uigjennemsigtighed, Haardhed, Vægt o. s. v., saa vil den som
udvortes endelig Existens fremstille sig som en Gjenstand for Viden;
men en saadan Videns Gjenstand er tillige en Skranke for Viden:
Gjenstanden indskrænker da den Viden, hvis Gjenstand den er; men
Alvidenheden taaler jo ingen Indskrænkning. Vil man nu beraabe sig
paa, at Alvidenheden vel skal kunne magte Kampestenen, efterdi den
Alvidende tillige er den Almægtige, da har en saadan umiddelbar
dogmatisk Paaberaabelse paa Almagtens Eenhed med Alvidenheden, som
det synes, ikke saamegen Lighed med et theoretisk \textit{testimonium
spiritus sancti}, som snarere med et praktisk Humbug, saalænge der
ikke gjøres Anstalter til at oplyse Noget angaaende Forholdet
imellem Alvidenheden og Almagten.

Hvilke Bidrag har da den objective, Guds subjective Dybheder
objectivt randsagende, Dogmatik ydet Videnskaben til nærmere
Forstaaelse af Alvidenheden „i Forhold til Skabningen“? Ingen,
aldeles ingen! Thi Erklæringer som disse: „Idet Gud erkjender Alt i
den evige Eenhed, erkjender han det ligesaavel i dets indre
Modsætninger og Forskjelle. Gud sætter Skilsmisse mellem Lys og
Mørke; han kjender Væsenet som Væsen, Skinnet som Skin“ o. s. v. ere
jo ikke videnskabeligt begrundede Erkjendelser; tvertimod! det er
reent grundløse Paastande, Almeensætninger, der kunne opstilles,
antages og gjentages af Enhver, uden at den, der opstiller dem,
derfor i nogen Maade
% --- p. 53 ---
\setcounter{footnote}{0}%
behøver at have \textit{testimonium spiritus} „theoretisk“; med
mindre det da theoretisk skulde være et \textit{testimonium
paupertatis}.

Afskrækket ved Dogmatikens mislykkede Forsøg med Testimoniet, kunde
man næsten føle sig fristet til med eet Slag at gjøre Ende paa alle
Qvaklerier med Subjectiviteten og dens Viden ved at opstille og
fastholde den Forudsætning, at den absolute Aand er i sin Alvidenhed
ligesaa ophøiet over al menneskelig Tænkning som over al menneskelig
Sandsning. Men, ikke at tale om det Uvidenskabelige i at antage en
guddommelig Viden, der skal falde saa ganske udenfor vor
menneskelige, at vi om en saadan, strengt taget, ingen Viden kunne
have, saa lægger jo den opstillede Forudsætning et saa svælgende
Dyb, ikke blot imellem den menneskelige og den guddommelige Viden,
men imellem den menneskelige Viden og Sandheden, at det paa slige
Vilkaar maa blive Videnskaben umuligt at danne sig et sandt Begreb
om, hvad Sandhed er. Kan man nu paa dette Standpunkt ikke vide, hvad
Sandhed er, saa kan man jo heller ikke vide, hvad der hører til en
sand Viden. Men en Subjectivitet, der saaledes faaer det Sande
udenfor sig og i den Grad farer vild fra al Viden af Sandheden, at
den tilsidst ikke veed, hvad sand Viden er, farer vild i sig selv,
og kan som en i sig vildfarende Subjectivitet umulig unddrage sig
den Selvmodsigelse, hvoraf al anden Modsigelse udspringer. Skal
Selvmodsigelsen hæves, og Subjectivitetsproblemet løses, da maa her
i Principet kunne paavises en Modsætningseenhed af oprindelig og
afledet Subjectivitet, en Forskjelseenhed af guddommelig og
menneskelig Viden. Men Vidensprincipet, den Grundsubjectivitet,
hvori den guddommelige Viden kan blive menneskelig og den
menneskelige guddommelig, kan ikke udvikle sig paa de umiddelbare
Modsætningers Trin; her udkræves høiere Former.
\subsection*{B.\\ Den reflecterende Subjectivitet.}
\addcontentsline{toc}{subsection}{B. Den reflecterende Subjectivitet}

% --- p. 54 ---
\setcounter{footnote}{0}%
Som Viden er, saaledes er den Vidende; som Subjectivitetsudviklingen er,
saaledes er Subjectiviteten. Dersom man til en Charakteristik af den
reflecterende Subjectivitet kun fordrer en Bevidsthed, der reflecterer, da
have vi allerede i Sensualismen og endnu mere i den subjective Idealisme
talende Vidnesbyrd om en Reflexion, der er uadskillelig fra
Subjectivitetens Væsen; men naar der ved Udtrykket „reflecterende
Subjectivitet“ betegnes et eget Trin af Subjectivitetens logiske
Selvudvikling, er det ikke nok, at Subjectiviteten reflecterer, den maa,
for at blive sig selv gjennemsigtig, udtrykkelig reflectere sin
Reflecteren. Og ikke engang det er nok; thi den subjective Idealisme har
tilstrækkelig godtgjort, at en Subjectivitet kan reflectere sin egen
Reflecteren paa en saa abstract, formel og eensidig Maade, at den, ved
uophørlig at fortære sit Indhold, tærer paa sig selv og fortærer sig selv.
Den i Selvreflexion klare Subjectivitet maa vide at udvikle Reflexionens
Former i uopløselig Eenhed med et i Formerne reflecteret Indhold; den maa
med det saaledes reflecterede Indhold ikke blot speile sig selv i sit
Andet, men gjennemskue Modsætningen af sig selv og Andet saa fuldstændig,
at Andetheden fra sin Side omvendt kan sees at reflectere sin
Selvstændighed, idet den reflecteres i Selvheden, og det Objective
stadfæstes af det Subjective som Videns Gjenstand af den Vidende. At just
dette er Opgaven, indlyser, naar man ret begynder at reflectere over
Reflexionen. Hvad er vel Reflexionen?

Ifølge almindelig Sprogbrug falde Tænken og Reflecteren saa temmelig
sammen: at reflectere over Noget er at tænke over Noget; Forskjellen er i
alt Fald kun den, at man ved Reflexionen gjerne betegner en særlig bevidst
Tænken, en Overveielse, en udtrykkelig Anvendelse af Eftertanken. Man gjør
endvidere, ifølge almindelig Sprogbrug, Forskjel imellem
% --- p. 55 ---
\setcounter{footnote}{0}%
at tænke paa Noget og at reflectere paa Noget: herpaa Reflecterende,
o. s. v. Fra dette Synspunkt er Reflexionen med sine forskjellige
Betydninger altsaa en subjectiv Act, ligesaavel som Tænkningen. I en ganske
anden Mening forekommer Ordet Reflexion i Optiken, hvor Talen er om Lysets
Reflexion, i Betydning af Lysstraalernes „Tilbagekastning“; denne Reflexion
er aldeles objectiv. Reflexion i objectiv og Reflexion i subjectiv
Betydning ere, ifølge den sædvanlige Forestillingsmaade, saa aldeles
forskjelligartede, at de neppe engang synes at vedkomme hinanden.
Anderledes i den absolut objective, immanente Begrebsudvikling; her er man
gaaet til den modsatte Yderlighed, at slaae den subjective Reflexion sammen
med den objective. Det Sidste er ligesaa eensidigt som det Første. Fra
Logikens Synspunkt bliver det derfor en væsentlig Opgave at udvikle
Forskjellen mellem Reflexion og Tænkning og derhos paavise Eenheden, nemlig
Forskjelseenheden af den logiske og den optiske, den subjective og den
objective Reflexion.

Betragte vi Forholdet mellem Tænkningen og Reflexionen nærmere, da bliver
det strax et Spørgsmaal, om Reflexionen skal bestemmes som en Ændring af
den almene Tænkning, eller Tænkningen som et særligt Moment, en Side af
Reflexionen selv. Hos Hegel er Reflexionen udtrykkelig bestemt som en
særlig Art af Tænkning, en objectiv Mellemform imellem Umiddelbarheden og
Speculationen. Saaledes svare alle Værenskategorierne til den umiddelbare
Tænkning, alle Væsenskategorier til Reflexionen, alle under Begrebet
hørende Totalitetskategorier til Speculationen. Hvad der gjælder om
Trichotomien i det Store, gjælder tillige i det Mindre og gjentager sig ved
hver enkelt Trehed. Reflexionens Plads er i det hegelske System saa typisk,
at enhver Afvigelse fra Typen forraader sig som en Inconseqvens. „Hos
% printed footnote mark: *)
Hegel“ --- siger Heiberg\footnote{Perseus Nr. 2, S. 44. (Prosaiske
Skrifter, 2det Bind, 1861, S. 165)} --- „er den absolute Væren fremstilt i
Kategorierne:
% --- p. 56 ---
\setcounter{footnote}{0}%
a) Væren, b) Intet, c) Vorden. --- Men denne Orden maa ansees som en lille
% the enumerators a) b) c) are antiqua against the Fraktur, left as plain
% letters per the house convention recorded for the synopsis head „A.“
Uagtsomhed, thi den er i Strid med Systemets hele øvrige Architektur. Vel
er Vorden Eenheden af Væren og Intet, men kun den dialektiske Eenhed af
begge, deres --- som Hegel selv kalder den --- \emph{urolige} Eenhed,
altsaa ikke den resulterende eller speculative. Vorden er i det hegelske
System aabenbart Grundkategorien for alt Endeligt; men netop derfor maa
den, ifølge selve Systemets hele Methode, holdes paa Endelighedens Plads.“
Men Endelighedens Plads er i denne Sammenhæng netop Reflexionens Plads. At
Reflexionen paa denne Maade er bleven en Tænkningens Mellemform, en ved
„det Dialektiske“ betegnet Overgangsform fra den kategoriske Umiddelbarhed
til den udviklede Forskjels resulterende „speculative“ Eenhed, skulde her
slet ikke have været berørt, dersom denne Reflexionens Underordnen under
Tænkningen som det Endelige under det Uendelige, som Forstanden under
Fornuften, ikke havde sin Grund i en umiskjendelig Overvurdering af „den
rene Tænknings“ abstracte Væsen.

I skarp Modsætning til Anskuen, Sandsen og Forestillen er den rene Tænken
for sig eensidig, ligesom den rene Negation er eensidig ligeover for den
umiddelbare Position. Tanken maa, ifølge Hegel, være det Negative, hvori al
Sandsning og Forestilling forsvinder: „den, der ret begynder at tænke,
taber Syn, Hørelse og alle Sandser“; men skal Negationen ikke være Andet
end en Abstraction, vil man virkelig gjøre Alvor af at bortkaste alle
Sandsningens, Forestillingens og Anskuelsens forskjelligartede Elementer,
for at Tanken selv som det Almene kan bevæge sig fra det Almene gjennem det
Almene til det Almene, da kan man spendere saamange Negationer og
Negationers Negationer paa dette Abstractionens \textit{caput mortuum}, som
man vil; det er dødt og bliver dødt. Skal Tænkningens Frigjørelse fra
Sandsning og Forestilling derimod, hvad Hegel jo ogsaa indskærper, forstaaes
saaledes, at de Forestillingselementer, som Tanken ved
% --- p. 57 ---
\setcounter{footnote}{0}%
sin Negation ophæver, bevares i den concrete Tænkning, da er Tænkningen med
det Samme forvandlet til Reflexion, og her er ikke længer Grund til at give
den Forestillingen abstract modsatte Tænkning noget ubetinget Fortrin
fremfor Forestillingen. Ved Reflexionen bliver det Positive en Reflex af
det Negative, men saaledes, at det Negative med det Samme bliver en Reflex
af det Positive; ved Reflexionen bliver Anskuelsen en Reflex af Tænkningen,
men saaledes, at Tænkningen med det Samme bliver en Reflex af Anskuelsen;
ved Reflexionen bliver Væren en Reflex af Intet, men saaledes, at Intet med
det Samme bliver en Reflex af Væren. Istedenfor at bestemme Reflexionen som
en underordnet Form af den rene Tænken og tage Parti for Abstractionen,
maae vi tvertimod erkjende, at det netop er Reflexionen, der er den
udviklede, ved sine Modsætninger integrerede Tænkning, den Tænkning, der
alene er skikket til at overgaae i Concretionen.

Er det allerede fra den subjective Side nødvendigt at minde om, at al
Tænken \textit{in concreto} er en Reflecteren, da er det fra den objective
Side, med udtrykkeligt Hensyn paa den hegelske Misviisning, endnu mere
nødvendigt at minde om, at al Reflecteren, logisk forstaaet, er en Tænken.
Thi hvor strengt Hegel end overalt holder paa „das Denken und das Denken
und das allgemeine Denken“, saa hændes det ham dog, at han i sin objective
Logik træffer paa Tanker, om hvilke han ikke ret veed, hvorvidt de egenlig
ere tænkte eller ikke tænkte. „Wenn man sagt,“ --- hedder det i
% printed footnote mark: *)
„Vorbegriff“ til 1ste Deel af Encyclopædien\footnote{Hegels Werke, 6ter
Band. Berlin 1840, S. 45.} --- „der Gedanke, als objectiver Gedanke, sey
das Innere der Welt, so kann es scheinen, als solle damit den natürlichen
Dingen Bewußtseyn zugeschrieben werden. Wir fühlen ein Widerstreben
dagegen, die innere Thätigkeit der Dinge als Denken aufzufassen, da wir
sagen:
% --- p. 58 ---
\setcounter{footnote}{0}%
der Mensch unterscheide sich durch das Denken vom Natürlichen. Wir müßten
demnach von der Natur als dem Systeme des bewußtlosen Gedankens reden, als
von einer Intelligenz, die, wie Schelling sagt, eine versteinerte sey.
Statt den Ausdruck \emph{Gedanken} zu gebrauchen, ist er daher, um
Mißverständniß zu vermeiden, besser Denkbestimmung zu sagen.“ Det er et
stort Problem, Hegel berører i disse faa Ord, et Grundproblem, der allerede
foresvævede Plato: „Dersom Ideerne,“ hedder det i Platos Parmenides, „kun
ere subjective Begreber (νοήματα), og Tanken nødvendigviis maa have sin
Gjenstand, saa maa man enten antage, at hver Gjenstand bestaaer af Tanker,
og at Alt tænker, eller ogsaa, at der vel er Tanker, men ingen Tænken
% Greek read at 600 dpi. The fount prints smooth-breathing-plus-acute as a
% single hooked sort (calibrated on the certain ἄρα and εἴπερ, p. 63); the
% marks on εἰ, τἆλλα and the two ἢ below could not be separated from
% circumflex/grave at the scan's native resolution, so the standard
% orthography of the Parmenides passage is followed and no variant asserted.
(ἀνάγκη, εἰ τἆλλα φῄς τῶν εἰδῶν μετέχειν, ἢ δοκεῖν σοι ἐκ νοημάτων ἕκαστον
εἶναι καὶ πάντα νοεῖν, ἢ νοήματα ὄντα ἀνόητα εἶναι).“ At Tingene selv
skulde kunne tænke, eller at der i Tingene skulde kunne gives Tanker, der
ikke vare tænkte, fandt Plato i den Grad meningsløst, at han, for at blive
en saa urimelig Antagelse quit, ikke tog i Betænkning at sætte de evigt
værende Ideer istedenfor de utænkte Tanker. Hvad har nu Hegel med sin
uophørlige Indskærpen af „das Denken und das Denken und das allgemeine
Denken“ egentlig gjort for at løse hint ældgamle platoniske Problem? Intet!
Istedenfor at skærpe Problemet, har han tvertimod udvisket det med en
Raisonnementsbemærkning som denne: „Statt den Ausdruck \emph{Gedanken} zu
% both halves of the line-broken „Ge=danken“ are letterspaced, so the
% emphasis is not split; the print has „ist es daher“ here against „ist er
% daher“ in the same quotation above, and both are left as printed
gebrauchen, ist es daher, um Mißverständniß zu vermeiden, besser
Denkbestimmung zu sagen“. Hvad der i denne Henseende har kunnet blende en
ellers saa energisk, aarvaagen og udholdende Tænker som Hegel, er ikke
vanskeligt at paavise; det er aabenbart det uklare Forhold, hvori han fra
Først af sætter Reflexionen til den egentlige Tænkning, og den, paa Grund
af Forholdets Uklarhed, tvetydige Brug, han saa gjør af den objective
Reflexion. Et slaaende Beviis herpaa er Følgende. Om Kategorien Existens
hedder det
% --- p. 59 ---
\setcounter{footnote}{0}%
% „Existenz“ and „Existerende“ below are set in fat-face (bold) Fraktur,
% not letterspaced: verified at 600 dpi against the plain „Existenz“ six
% lines further down and the plain „Existirende“ beside it. This is a third
% emphasis device in the book, not yet in the header's convention list.
i Logiken: „Die \textbf{Existenz} ist die unmittelbare Einheit der
Reflexion- in- sich und der Reflexion- in- Anderes“; Tingen, det
\textbf{Existerende}, bestemmes nu i Conseqvens heraf saaledes: „Die
Reflexion- in- Anderes des Existirenden ist aber ungetrennt von der
Reflexion- in- sich; der Grund ist ihre Einheit, aus der die Existenz
hervorgegangen ist. Das Existirende enthält daher die Relativität und
seinen mannigfachen Zusammenhang mit andern Existirenden an ihm selbst, und
ist in sich als Grund \emph{reflectirt}. So ist das Existirende Ding.“

Det er af Logiken uimodsigelig klart, at Hegel ikke blot lader Tingenes
Existens beroe paa objectiv Reflexion, men endog ligefrem gjør Tingene selv
reflecterende, forsaavidt nemlig enhver Ting reflecterer sig i sig selv ved
at reflectere sig i sine Egenskaber, o. s. v. Men naar Hegel ikke tager i
Betænkning at gjøre Tingene til reflecterende Væsener, medens han dog tager
i Betænkning at gjøre dem til tænkende Væsener, saa sees jo tydelig heraf,
at han stiltiende gjør en væsentlig Skilsmisse imellem Reflexionen og
Tænkningen, og det en Skilsmisse, hvorefter det, trods al Paaberaabelse af
„das Denken und das Denken und das allgemeine Denken“, dog ikke bliver
% printed footnote mark: *)
Tænkningen\footnote{At der ifølge sædvanlig Sprogbrug kan være Grund til at
bruge \emph{Ordet Tænkning} som det mere almindelige Udtryk for den logiske
Aandsvirksomhed, kan gjerne indrømmes; kun maa man erindre, at den
egentlige, bestemte Tænken er en Reflecteren.}, men Reflexionen, nemlig den
optisk objective Reflexion, der fra Først til Sidst opfattes som det
Almene.

Til nærmere Bestemmelse af Forholdet mellem Tænken og Reflecteren og derved
tillige af Forholdet mellem subjectiv og objectiv Reflexion, vil det altsaa
være nødvendigt at lade Reflexionen udklare sig fra den subjective Side som
\emph{Forstandsreflexion}, fra den objective som \emph{Væsensreflexion} og
ved energisk Vexelbestemmelse af det Subjective og det
% --- p. 60 ---
\setcounter{footnote}{0}%
Objective, som \emph{Virkelighedsreflexion}. Modsætningseenheden af
Subjectiviteten og det Objective udvikler sig da med det Samme saaledes, at
Forstandsreflexionen væsentlig sees at være en af Andetheden betinget
Selvreflexion, Væsensreflexionen en i Selvheden begrundet Reflexion af
Andet, og Virkelighedsreflexionen en for Andetheden ved den overgribende
Selvhed bestemt Selvstændighedsreflexion.

% Centred display head, letterspaced Fraktur with the antiqua-free „d)“ in
% Fraktur; the lettering continues the a)/b)/c) of division A. Set at the
% level below \subsection* (the skeleton's A/B/C) so that it does not
% outrank its own division head.
\subsubsection*{d) Forstandsreflexion.}
Al Forstaaelse af Andet er væsentlig Selvforstaaelse; thi Forstaaelsen
beroer just derpaa, at den forstandige Subjectivitet, saa at sige, staaer
foran sin Gjenstand som foran et Andet, et Fremmed, og dog i dette Andet
erkjender sit Eget, i dette Fremmede opdager sig selv: under Forstaaelsen
bliver Andetheden netop reflecteret i Selvheden saaledes, at
Subjectiviteten udtrykkelig veed sig som reflecterende. I Forstaaelsen
reflecterer Forstanden tillige sig selv. Vistnok kan den reflecterende
Forstand kun udvikles ved Udvikling af et Gjenstandsindhold, der med det
Samme bliver forstaaet; men Gjenstandens Udvikling er dog i Reflexionen
forskjellig fra den sig derved udviklende Forstand. Her maa i Forstaaelsen
selv gjøres en forstandig Adskillelse imellem en Forstaaelse af Gjenstanden
og en Forstaaelse af den Forstand, hvorved Gjenstanden bliver forstaaet. Jo
skarpere Adskillelsen er, desto tydeligere er igjen Reflexionen af det
Adskilte. Dersom det var Forstanden en Umulighed at reflectere sine evige
Formbestemmelser i reflecteret Adskillelse fra sit reale Indhold, da vilde
der være en saa uendelig Kløft imellem den menneskelige og den
guddommelige Forstand, at Forstanden derved absolut maatte blive sig selv
en Gaade. Men hvo indseer ikke, at de Formbestemmelser, der vise sig som
absolut nødvendige for den menneskelige Forstand, ogsaa maae have Gyldighed
for den guddommelige, efterdi de høre til Forstandens Væsen? Det er
tankeløst, om Nogen vilde sige: „vor Forstand er i sin menneskelige
Indskrænket\-%
% --- p. 61 ---
\setcounter{footnote}{0}%
hed jo vistnok nødt til at gjøre Forskjel imellem Væren og Vorden, Hvile og
Bevægelse, Væsen og Skin, Grund og Følge, Hele og Dele, Indre og Ydre,
Aarsag og Virkning, Midler og Formaal o. s. fr., men derfor er det jo ikke
sagt, at en saadan Forskjel skulde være nødvendig ogsaa for den
guddommelige Forstand. Vi med vor ufuldkomne Forstand maae naturligviis
indrømme, at det Hele, mathematisk seet, er større end en Deel, og
ligestort med sine Deles Sum; men hvad Beviis have vi for, at den
guddommelige, for os ufattelige, Forstand ikke maaskee vender Forholdet om,
og gjør Delen større end det Hele og Delenes Sum mindre end hver enkelt
Deel? Vi med vor menneskelige Forstand formaae jo vistnok ikke at tænke os
et Væsen uden Phænomen, en Ting i sig (ein Ting an sich) uden Egenskaber,
% „ein Ting an sich“ as printed: German phrase in Fraktur with Danish
% „Ting“ for „Ding“; not normalised.
et absolut Indre uden et tilsvarende Ydre; men en guddommmelig Forstand, en
% PRINTER'S DEFECT: „guddommmelig“ with three m's, verified at 600 dpi.
% Transcribed as printed; the sense-reading is „guddommelig“.
Forstand, der er Eet med en guddommelig Magt, kunde jo derfor gjerne see
sig istand til at kaste Phænomenerne bort og beholde Væsenet,
tilintetgjøre Egenskaberne og conservere Tingen, afskaffe det Ydre og slaae
sig til Ro i det Indre?“ --- At en saadan Adskillelse af guddommelig og
menneskelig Forstand er tankeløs og uforstandig, trænger ikke til særligt
Beviis; Beviset ligger i Sagen selv. Bestræbelsen for at gjøre Forstanden
overmenneskelig og guddommelig er ved den projecterede Adskillelse
forvandlet til en Bestræbelse for at gjøre den guddommelige Forstand til en
Forstandens Uting, en høieste Uforstand. „Ja, for vor indskrænkede
Fatteevne,“ siger man, „maa den guddommelige Forstand nødvendigviis tage
sig ud som Uforstand; thi i modsat Fald vilde den jo ikke absolut overgaae
vor Forstand.“ Men Tvetydigheden i slige Vendinger ophører at være
blendende, naar vi besinde os paa den afgjørende Forskjel imellem: at fatte
Forstandens uendelige Reflexion i sit uendelige Indhold og at fatte
Forstandsreflexionen selv i sin formelle Forstandighed. Ved den første
Opgave betegnes Modsæt\-%
% --- p. 62 ---
\setcounter{footnote}{0}%
ningen, ved den sidste Eenheden af den guddommelige og den menneskelige
Forstand.

Forsaavidt da det hele System af reale Forskjelsbestemmelser, der særligt
angaae Indholdet, bortfalder eller, nøiagtigere, taber sig og forsvinder i
Reflexer af Forstandens ontologiske Selvreflexion, have vi i den saaledes
bestemte ontologiske Selvhed Forstandsprincipet som et
Subjectivitetsprincip. Principets Selvudvikling kan altsaa kun i samme Grad
blive alsidig opfattet, som Forstandsreflexionen selv bliver alsidig
udviklet.

% Centred display head: italic antiqua Greek α) with a letterspaced Fraktur
% rubric.
\subsubsection*{α) Sættende og forudsættende Reflexion.}
Vilde Nogen spørge: „hvor gammel er vel Tilværelsens Forstand? naar blev
den til? hvorlænge kan den vare?“ da vilde naturligviis enhver Forstandig
strax mærke det Urimelige i slige Spørgsmaal. Forstanden selv, Tilværelsens
Forstand, kan jo ikke have nogen endelig Tilbliven: den er og maa være
evig. Men hvad skulle vi nu sige om en Forstand, der er til i en heel
Evighed, uden at forstaae enten sig selv eller Andet, altsaa: uden mindste
Forstand enten paa sig selv eller Andet, altsaa: en evig Forstand uden al
Forstand? Det Syn seer noget uforstandigt ud. Paa timelig, endelig Maade
kan der jo vistnok i Tingene selv være meget Forstandigt, som din og min
Forstand ikke fatter; men her er nu ikke Tale om en Forstand som din og
min, men om Forstanden selv, den evige Forstand. Hvad der gjælder om
Tanken, Anskuelsen og Forestillingen, gjælder i denne Henseende ogsaa om
Forstanden. Uden Tænkning ingen Tanke, uden Anskuen ingen Anskuelse, uden
Forestillen ingen Forestilling, uden Forstaaen og Forstaaelse ingen
Forstand.

Fra Psychologiens Standpunkt gjælder det, at det forstandige Subject er Et
og Subjectets Forstand et Andet; men det Endelige i denne Adskillelse
tilhører kun Endeligheden. Seet i Principet, d. v. s. fra Ontologiens
Standpunkt, er
% --- p. 63 ---
\setcounter{footnote}{0}%
den sig forstaaende Forstand Subjectiviteten selv. Det er den guddommelige
Forstand --- de Gamles νοῦς ---, der evig tænker sig selv, og denne
Tænkning er, som Aristoteles udtrykker det, Tænkningens Tænkning: αὐτὸν
ἄρα νοεῖ, εἴπερ ἐστὶ τὸ κράτιστον, καὶ ἔστιν ἡ νόησις νοήσεως νόησις.

Men naar Forstanden selv er oprindelig Subjectivitet, hvoraf kommer det da,
at man ikke ligesaa umiddelbart indseer Subjectivitetens som Forstandens
Evighed? Det kommer aabenbart deraf, at man her har at stride med en
dobbelt Misforstaaelse, en psychologisk og en ontologisk. Den psychologiske
Misforstaaelse er, at al tænkende Subjectivitet, al Viden og Bevidsthed
uden videre maa være analog med den menneskelige, af Nerve- og
Hjernevirksomhed betingede, Subjectivitet og dens indskrænkede Viden og
Bevidsthed. Den ontologiske Misforstaaelse er, at Forstandens evige Væren
skulde kunne være umiddelbar, som en Væren af et evigt Rum, medens man
% „Rum“: the impression is light at the line start and both OCR witnesses
% read „Ram“; the stroke count at 600 dpi gives u, and the sense (Hegel's
% Raum) agrees. Rejected reading: „Ram“.
derimod ikke kan dølge for sig selv, at Subjectivitetens Væren er og maa
være uendelig reflecteret. At tillægge Forstanden en oprindelig umiddelbar
Væren, være sig i en given Tingenes Orden eller i Materien og dens Kræfter,
er i allerhøieste Grad uforstandigt; thi den uendelige Forstand er kun i en
uendelig Reflexion. Psychologien siger, at Forstanden er en Evne til
Reflexion; Ontologien indseer, at Forstanden selv er Reflexion. Problemet
om Subjectivitetens Oprindelighed er netop Problemet om Forstandens
Oprindelighed, Problemet om den evige νοῦς som ἡ νοήσεως νόησις.

Skulle saadanne kategoriske Udtryk som: „den sig tænkende Tænkning“, „den
sig forstaaende Forstand“, „det sig begribende Begreb“ være mere end
Talemaader, hvorved man søger at undgaae Vanskelighederne med det
Subjective, for uforstyrret at give sig hen i Betragtningen af det
Objective, da er det i Sandhed ikke nok, at man engang imellem kommer i
Tanke om, at der er en tænkende Subjectivitet, og
% --- p. 64 ---
\setcounter{footnote}{0}%
saa i al psychologisk-ontologisk Tvetydighed siger med Hegel: „Das Denken
als Subject vorgestellt ist Denkendes, und der einfache Ausdruck des
% printed footnote mark: *)
existirenden Subjects als Denkenden ist Ich\footnote{Hegels Werke, 6ter
Band. Berlin 1840, S. 34.}.“ Allerede Fichtes Begyndelse til en
Construction af det rene Jeg viser tydelig nok, hvad det vil sige, at den
sig forstaaende Forstand oprindelig er en Subjectivitet. De to første
% The letters A, the equals signs and „Non-A“ are antiqua against the
% surrounding Fraktur; set as in-line mathematics, with Non-A italic.
Sætninger: $A = A$, og \textit{Non-A} ikke $= A$, ere, for sig betragtede,
den formelle Logiks noksom bekjendte Principer, nemlig Identitets- og
Contradictionsprincipet, altsaa rene Forstandssætninger, Grundsætninger,
hvortil den reflecterende Forstand vil føre alle Reflexioner tilbage, for
ret at forstaae sig selv. Men medens den formelle Logik kun opstiller
Grundsætningerne saaledes, at den i en egen, udenfor den logiske Sætten
faldende Forudsætning fastholder det psychologisk tænkende Subject, et
Subject, der bruger sin Forstand uden at indsee sin Væsenseenhed med
Forstanden selv, er det derimod Fichtes Fortjeneste at have fra Grunden af
paaviist, at den Tænkende ikke staaer udenfor Tanken, at det netop er den
sig forstaaende Subjectivitet, der i de allerførste forstandige
Vexelbestemmelser af Sætten, Forudsætten og Modsætten væsentlig har med sig
selv at gjøre.

Ifølge den formelle Logik maa det at sætte og det at forudsætte Noget
betragtes som to Ting, der falde udenfor hinanden. At sætte Noget er at
give det Tilværen; hvad jeg selv sætter, har jeg i min Magt: min Tanke kan
holde det fast eller slippe det, eftersom den finder for godt. At
forudsætte Noget er derimod at lade sig Noget give, at antage det for et
Givet. Det Givnes Realitet viser sig just deri, at det i Tanken Forudsatte
behersker den forudsættende Tanke, og gjør den afhængig af sig. Under sin
Sætten viser Forstanden sig, med andre Ord, som activ, under sin
Forudsætten som passiv. Saalænge Sætten og Forudsætten paa denne
% --- p. 65 ---
\setcounter{footnote}{0}%
Maade stilles imod hinanden, er det ganske vist en Umulighed at fatte
Subjectivitetens oprindelige Selvsætten, den Urhandling, hvorved den
ontologiske Jeghed fra Evighed af er Jeg $=$ Jeg.

For i Sandhed at kunne være et Selv, maa Subjectiviteten sætte sin egen
Selvhed; men for at en Subjectivitet virkelig skal kunne foretage en saadan
Selvsætten, maa den jo nødvendigviis allerede forud være; en ikke værende
Subjectivitet kan umulig sætte enten sig selv eller noget Andet. Den
virkelige Subjectivitet maa, for at kunne sætte sig selv, begynde med at
forudsætte sig selv, begynde med at finde sig som et allerede Sat. Altsaa
Thesis: Subjectiviteten er kun ifølge en oprindelig Selvsætten subjectiv;
dens Sætten gaaer i det Uendelige forud for dens Forudsætten. Antithesis:
den oprindelige Selvsætten forudsætter et sættende Selv; Subjectivitetens
Forudsætten gaaer i det Uendelige forud for dens Sætten. Denne Antinomie
vilde ganske vist være uopløselig, og Problemet om Jegets „Urhandling“ et
hyperspeculativt Problem, hvis den reflecterende Forstand i hiin
tilsyneladende mystiske Urhandling ikke naturligt kunde gjenkjende sine
egne velbekjendte Operationer, sin egen Sætten og Forudsætten.

For den rene Reflexion af Sætten og Forudsætten er det ligegyldigt, hvad
Subjectiviteten iøvrigt formaaer at sætte, om et heelt Solsystem, eller kun
Bogstavet A. Sagen er, at enhver Sætten beroer paa en Forudsætten. Man kan,
for igjen at minde om Fichte, ikke sætte $A = A$, uden at forudsætte Jeg
$=$ Jeg; men man kan heller ikke omvendt sætte Jeg $=$ Jeg, uden at
forudsætte, at, dersom A er, saa er $A = A$; idet $A = A$ beroer paa
Sætten, beroer Jeg $=$ Jeg paa Forudsætten, og omvendt; Sætten og
Forudsætten ere i Reflexionen altsaa ikke to Handlinger, men een:
Forudsætten en begyndende Sætten, og Sætten en sig fuldendende
Forudsætten. Ved at reflectere sig i sin Forudsætten reflecterer enhver
Sætten sig jo vistnok i sit Modsatte; men da dens
% the signature mark „Grundideernes Logik.  5“ at the foot of this page is
% a printer's mark and is not transcribed
% --- p. 66 ---
\setcounter{footnote}{0}%
Modsatte netop er en Forudsætten, og Forudsætten som en begyndende Sætten
ligeledes reflecterer sig i sig ved at reflectere sig i sit Modsatte, maa
enhver Sætten, som saaledes vedblivende reflecterer sig i en sit Modsatte
reflecterende Modsætning, hævde sig en uendelig Selvfordobling, og i denne
Selvfordobling vise sig som oprindelig Selvsætten. Paa den oprindelige
Selvsætten, den i Reflexionen gjennemsigtige Eenhed af Sætten og
Forudsætten, beroer da Selvhedens Fastsætten af sig selv og sit Andet, og
paa denne Fastsætten beroer saa igjen al væsentlig Bestemmen saavel of det
% PRINTER'S DEFECT: „of“ for „af“, verified at 600 dpi against the certain
% a of „saavel“ and „saa“ on the same line and the o of „roer“. Transcribed
% as printed.
Objective som af det Subjective.

% Centred display head: italic antiqua Greek β) with a letterspaced Fraktur
% rubric.
\subsubsection*{β) Modsættende og eenhedsættende Reflexion.}
Fastholdes den oprindelige Selvsætten i Form af Forudsætning, maa den
tilsvarende Sætten af Andet fastholdes som Følgesætning. Al Sætten af Andet
beroer paa Selvsætten; men medens Selvsætten, for sig betragtet, holder sig
i det Eensformige, betegner Sætten af Andet derimod et Mangeformigt og
Forskjelligt; thi Andet reflecteres ikke blot som Selvets Modsatte, men i
Andetheden selv reflecteres tillige en sig gjentagende Modsætning af Andet
og Andet (τὸ ἕτερον) i det Uendelige: den sættende Reflexion er herved
bestemt som modsættende. Hvad er i Tilværelsen selv det oprindeligt Første,
Eenheden eller Toheden, det Endelige eller det Uendelige, Indifferensen
eller det Differente, det Ufuldkomne eller det Fuldkomne? Dette mangesidige
og fleertydige Spørgsmaal, der nu i Philosophiens Historie saa ofte har
gjentaget og fremdeles vil gjentage sig, kan umulig faae nogen
tilfredsstillende Besvarelse, saalænge man endnu ikke ret har gjort sig
Rede for Modsætningseenheden af Princip og Begyndelse; men for at klare sig
denne Modsætningseenhed, maa man fremfor Alt klare sig Modsætningen af den
modsættende og eenhedsættende Reflexion.

% --- p. 67 ---
\setcounter{footnote}{0}%
Gaaer man uden videre ud fra den Forudsætning, at Eenheden, den rene
absolute Eenhed, absolut maa sættes som det egentligt Første, da indsniger
sig let den Biforestilling, at den rene, absolute Eenhed skulde være, om
ikke just en heel Evighed, saa dog i Evigheden selv nogle Øieblikke forud
for den deraf afledede Tohed og den af Toheden igjen afledede
Mangfoldighed. Denne Synsmaade er uholdbar. Det er, naar man besinder sig,
indlysende, at den Eenhed, der skal kunne udvikle Tohed, Mangfoldighed og
Modsætning af sig, allerede maa have Modsætningen, Mangfoldigheden og
Toheden i sig. Saalidtsom Forstanden kan begynde med at sætte en oprindelig
Tohed uden at forudsætte Tohedens Modsatte, en oprindelig Eenhed,
ligesaalidt kan den begynde med at sætte en oprindelig Eenhed, uden at
forudsætte Eenhedens Modsatte, en oprindelig Tohed.

Den Grundmodsigelse, at Toheden altid viser sig som det Første, naar man
vil begynde med Eenheden, og Eenheden som det Første, naar man vil begynde
med Toheden, finder sin Opløsning deri, at Toheden kun kan begynde,
forsaavidt Eenheden er i Principet, efterdi Eenheden kun kan begynde,
forsaavidt Toheden er i Principet. Urhandlingen, den oprindelige
Selvsætten, hvori den sig reflecterende Subjectivitet har evig Bestaaen,
maa for Alting ikke forvexles med en til $A = A$ svarende Sætten af Jeg $=$
Jeg; en saadan første almindelig Sætten er kun en Begyndelsessætning, ikke
en Urhandling, men en Abstractionshandling, et eensidigt Moment i den
alsidige, uendelig concrete Urhandling selv. Heller ikke kan man sige, at
Urhandlingen fra Begyndelsen a skulde være fuldendt i den Forstand, at den
% PRINTER'S DEFECTS, both on this page and both verified at 600 dpi:
% „a“ for „af“ (the f is simply absent, with the word's space left standing)
% and „om“ for „som“ on the next line. Transcribed as printed. A speck of
% ink or dirt stands before „om“ and is not transcribed.
kunde opstilles om et første Led i en uendelig Række af Handlinger; man
kunde med lige saa megen Ret paastaae, at Urhandlingen først blev fuldendt
med Rækkens sidste Led; ja, man kunde maaskee med endnu større Ret
paastaae, at det Led, der betegnede Urhandlingens Fuldendelse, egentlig
maatte udgjøre Midten i den til
% --- p. 68 ---
\setcounter{footnote}{0}%
begge Sider uendelige Række. Sandheden er, at Urhandlingen netop er den i alle Forstandshandlinger sig selv gjentagende Subjectivitetshandling, den i alle Reflexionens Modsætninger reflecterede Eenhedsætten.

Modsætningen af den modsættende og eenhedsættende Reflexion kan da heller ikke afsluttes dermed, at den ene af disse Reflexioner umiddelbart sættes som Negation af den anden; thi i Reflexionen er, som ovenfor viist, det abstract Negative ligesaavel en Reflex som det abstract Positive: den umiddelbare Negation er kun en Grændse for det Umiddelbare, ikke for det i sig Reflecterede. Ved at lade Jeget selv sætte Ikke-Jeget som sin oprindelige Skranke mener Fichte jo vistnok at have een Gang for alle sat den qvalitative Modsætning fast; men det er en Misforstaaelse. Saalænge Jegets Skranke ikke kan faae anden Selvstændighed end den abstract negative: at være Ikke-Jeg, og det et Ikke-Jeg, som, sat af Jeget, falder indenfor Jeget selv, saalænge vil den Conseqvens ogsaa paatrænge sig, at Jeget ligesaavel er en Skranke for Ikke-Jeget, som omvendt; at Jeget følgelig ligesaa vel kan siges at falde indenfor Ikke-Jeget, og at være sat af Ikke-Jeget, som omvendt. Men naar Ikke-Jeget paa denne Maade bliver baade Ikke-Jeg og Jeg, idet Jeget fra sin Side skal være baade Jeg og Ikke-Jeg, saa maa Modsætningen af Jeg og Ikke-Jeg jo med det Samme gaae ud%
% printed mark: *)  — the note is set in Fraktur (German), with antiqua equals-signs
\footnote{„Ist Ich $=$ Ich, so ist alles gesetzt, was im Ich gesetzt ist. Nun soll der zweite Grundsatz im Ich gesetzt seyn, und auch nicht im Ich gesetzt seyn. Mithin ist Ich nicht $=$ Ich, sondern Ich $=$ Nicht-Ich, und Nicht-Ich $=$ Ich.“ J. G. Fichtes Wissenschaftslehre. Leipzig 1794, S. 25.}.

For at fastholde den nominelle Modsætning mellem Jeg og Ikke-Jeg, beraaber Fichte sig umiddelbart paa den Kjendsgjerning, at begge Modsætningens Led ere i Bevidstheden, og
% PRINTER'S DEFECT: "og" is set twice, at the end of one line and the start of the
% next ("i Bevidstheden, og / og at deres Eenhed"). Dittography; transcribed as printed.
og at deres Eenhed som Følge heraf ogsaa maa sættes i Be\-%
% --- p. 69 ---
\setcounter{footnote}{0}%
vidstheden.%
% This page carries a footnote but NO printed call mark anywhere in the body text
% (checked over the whole page at 600 dpi). The note is placed here, at the sentence
% the quotation glosses. The printed note opens with a T sort where German requires D
% ("Tie Gegensätze"): calibrated against "Demnach" on the same line of the same note,
% whose D shows the closed Fraktur bowl. Transcribed as printed.
\footnote{„Tie Gegensätze, welche vereinigt werden sollen, sind im Ich als Bewußtseyn. Demnach muß auch $X$ im Bewußtseyn seyn.“ Anfr. Skr. S. 26.}
Men Bevidstheden som psychologisk Kjendsgjerning er en altfor endelig concret Størrelse, til at Reflexionen uden videre kan sætte den i Eenhed med hiin første abstracte Modsætning. Skal den naturlige Forstand dømme Jeget og Ikke-Jeget imellem, da vil den sikkerlig ikke undlade at tilkjende Ikke-Jeget en Verden af Gjenstande. Ved at paaberaabe sig den umiddelbare Bevidsthed om Modsætningen har Fichte fra Begyndelsen af netop afledet Opmærksomheden fra det Utilstrækkelige i Modsætningen selv; heri ligger aabenbart den egentlige Grund til, at han med al sin Anstrengelse for at udvikle den hele Virkelighed af Jeget, ikke har kunnet drive det videre end til en Række kategoriske Omskrivninger af hiin første, idelig forsvindende og idelig tilbagevendende Strid imellem Jeg og Ikke-Jeg. Det er en Reflexionsfordobling, her fordres; thi det er ikke nok, at her paavises en abstract Modsætning af Jeg og Ikke-Jeg, her maa i Ikke-Jeget tillige paavises en udtrykkelig Modsætning imellem det ene
% "Ikke Jeg" here and twice below is printed WITHOUT the hyphen, unlike the many
% hyphenated "Ikke-Jeg" on the same page; verified at 600 dpi against the plain
% double-oblique hyphen in "Ikke=Jeg" on the same lines. Not regularised.
Ikke Jeg og det andet. Først idet Ikke-Jeget paa denne Maade fordobles ligesaavel som Jeget, vil det af den dobbeltsidige Fordobling sees, at den Modsætning, hvori det ene Ikke Jeg formedelst Jeget staaer til det andet, er af ganske anden Art end den Modsætning, hvori Jeget formedelst Ikke-Jeget oprindelig staaer til sig selv. I Jegets Selvmodsætning formedelst sit Ikke-Jeg er Reflexionen \emph{eenhedsættende}, thi Modsætningen selv er Reflex i den satte Eenhed: den satte Eenhed er den i sit Andet reflecterede Selvhed, Subjectiviteten selv. I Ikke-Jegets Modsætning til Ikke-Jeg formedelst Jeget er Reflexionen derimod \emph{modsættende}, thi Eenheden af Ikke Jeg og Ikke-Jeg er negativ, en Reflex i den sig fortsættende Andethed. Den Selvsætten, der udelukkende gjælder
% --- p. 70 ---
\setcounter{footnote}{0}%
for Jeget, staaer altsaa nu, men ogsaa først nu, i afgjørende Modsætning til den Sætten af Andet mod Andet, o. s. fr., altsaa til den Andethedsætten, hvorved Ikke-Jeget væsentlig charakteriseres. Idet Modsætningen gjennemføres og Andetheden betones, falder Eftertrykket ikke længere paa en eensformig Gjentagelse af Ikke-Jeg og Ikke-Jeg, men paa en udtrykkelig Modsætning af det ene Ikke-Jeg mod det andet. Ved den udtrykkelige Modsætning falder da ogsaa det negative Udtryk: Ikke-Jeg omsider bort; istedenfor at sætte negativt det ene Ikke-Jeg imod det andet Ikke-Jeg, sætter Reflexionen nu positivt det Ene imod det Andet, det ene Positive imod det andet Positive; altsaa --- i Modsætning til den sig selv sættende Subjectivitet --- det ene Objective imod det andet Objective, det ene Object imod det andet.

\subsubsection*{γ) Iboende og oversvævende Reflexion.}

Forsaavidt Subjectiviteten selv væsentlig er Reflexion, er Reflexionen Subjectiviteten \emph{iboende}; for Objectet derimod, der ikke saaledes kan reflectere sig selv, men maa bestemmes ved den paa sin Modsætning reflecterende Subjectivitet, er den bestemmende Reflexion \emph{oversvævende}. Kunde Modsætningen af iboende og oversvævende Reflexion fastholdes i al Umiddelbarhed, da maatte vi ubetinget anerkjende den kantiske Skilsmisse imellem Forstandsreflexionen og Tingen i sig; vi maatte conseqvent benegte Gyldigheden af en objectiv Reflexion. Men Modsætningen af oversvævende og iboende Reflexion er igjen saa reflecteret, at den iboende selv viser sig som oversvævende og den oversvævende som iboende.

Dersom den iboende Reflexion ikke tillige var oversvævende, kunde Subjectiviteten ingen Viden have, hverken af sig selv eller af noget Andet. Der kunde slet ingen Viden gives; thi da Viden, som ovenfor er viist, forudsætter sig selv i det Uendelige, maa den overalt beroe derpaa, at Opfattelsen af et allerede Opfattet igjen bliver opfattet, at Forestillingen om
% --- p. 71 ---
\setcounter{footnote}{0}%
en Forestilling igjen bliver forestillet o. s. v. Forstand paa det Endelige begynder altid med en Forestilling om, hvorledes en Forestilling bestemmes ved en anden Forestilling, og saaledes gaaer det i det Uendelige. Ligesom den forestillende Subjectivitet med hvert Øieblik, ogsaa da, naar den udtrykkelig vil forestille sig selv, gaaer ud over det Forestillede, saaledes gaaer den tænkende Subjectivitet med hvert Øieblik, ogsaa da, naar den udtrykkelig tænker sig selv, ud over den tænkte Tanke. Da nu den reflecterende Subjectivitet har Forestillingen i Reflex af Tanken og Tanken i Reflex af Forestillingen, saa maa Subjectiviteten nødvendigviis have sine iboende Reflexioner saaledes, at den uophørlig oversvæver dem, og derved tillige uophørlig oversvæver sig selv.

Hvad Reflexionen af det Objective, Objecternes Reflexion, angaaer, har Forstandskritiken hidtil stillet sig i en høist charakteristisk Modsætning til „Speculationen“. Medens „Speculationen“ ikke tager i Betænkning at tillægge Tingene selv, Substanserne, Objecterne o. s. v. en saa immanent eller iboende Reflexion, at de derved sættes istand til at reflectere sig i hinanden og med hinanden og hver for sig i sig selv, uden at hertil behøves nogensomhelst subjectiv eller oversvævende Reflexion, saa har Forstandskritiken fra sin Side drevet det til den modsatte Yderlighed, og erklæret al objectiv, Tingene selv iboende Reflexion for en plat Uting. Hegels og Herbarts diametralt modsatte Behandling af Identitetets- og Contradictionsprincipet viser tydeligt nok, hvad Yderlighederne her ville sige.

Sætningen $A = A$ betegner, ifølge Herbart, ikke, at $A$ reflecterer sig i sig selv, men kun, at vor Tanke fordobler sig i en tom Gjentagelse, en Tautologie. \textit{Posito A, nulla ponitur relatio. Posito $A = A$, res non duplicatur, sed cogitatio. Itaque res non refertur ipsa ad sese (quasi Fichtianum Ego), sed duplex vel multiplex cogitandi actus, quem iterare licet vel hodie vel cras vel minimo
% --- p. 72 ---
\setcounter{footnote}{0}%
quoque temporis intervallo interjecto, semper in idem punctum reverti dicitur sine ullo discrimine, cui relationis aliquid affingi debeat vel possit}%
% printed mark: *)  — Herbart's Latin title is antiqua, the rest of the note Fraktur.
% The third initial of Heegaard's name is printed with a B sort (two closed bowls,
% like the B of "Bogstavet" on this page), not the open-armed Fraktur V of "Værende"
% on this page; transcribed as printed, though the man is P. S. V. Heegaard.
\footnote{Herbarts \textit{De principio logico exclusi medii inter contradictoria non negligendi}. Kl. phil. Schr. II. P. 724. Jvfr. P. S. B. Heegaard, den Herbartske Philosophie, Kbhvn. 1861, S. 4 flgd.}.

Denne af Herbart opstillede Synsmaade vil den sunde Menneskeforstand ved første Øiekast naturligviis finde ligesaa uigjendrivelig, som den ved første Øiekast vil finde Hegels objective Reflexion i allerhøieste Grad urimelig. Det rene $A$ kan jo umulig reflectere sig selv, saa lidt som Bogstavet $A$ kan tænke sig selv, eller Størrelsen $A$ dividere sig selv med $B$. Men er det først indrømmet, hvad den sunde Menneskeforstand saaledes ved første Øiekast maa have indrømmet, saa sees Problemet ved næste Øiekast; og strax naar Problemet sees, vil det ogsaa sees, at Forstandskritiken er ligesaa eensidig som Speculationen. Hvad er det nemlig for et Standpunkt, hvorfra Herbart vil afvise den Objectet $A$ iboende Reflexion? Det empirisk-psychologiske! Den reflecterende Subjectivitet stiller sig her i al Endelighed ligeoverfor de tilværende Gjenstande; paa dette Standpunkt kan, som vi ovenfor have seet, den sættende Reflexion ikke gjennemskue sin egen Forudsætten: idet den psychologiske Subjectivitet forudsætter $A$, som et virkeligt Object, forglemmer Reflexionen sig selv i den Grad, at den uden videre betragter dette sit Forudsatte som et Værende, hvis Tilværen dog ikke forud er sat. Realet, det sande Værende, er jo netop et saadant Forudsat, der ikke forud maa være sat. Den Grundmodsigelse, at Realerne forudsættes ikke forud at være satte, gjentager sig i Reflexionen af deres indbyrdes Modsætning. Dersom Realerne ikke vare i indbyrdes absolut Modsætning, kunde der ikke være en Fleerhed af absolute Realer. Men dersom Realerne virkelig ere i Modsætning, maae de jo være hverandre modsatte; men modsatte kunne de
% --- p. 73 ---
\setcounter{footnote}{0}%
dog umulig være uden at være satte. Modsætning er desuden umulig uden Relation, og Relation umulig i en absolut Modsætning, efterdi en absolut Modsætning, kritisk forstaaet, selv er umulig. Modsigelsen er fra alle Sider iøinefaldende: forsaavidt Realernes egentlige Væsen antages at falde udenfor de Relationer, hvori de, ifølge vore „zufällige Ansichten“, tilfældigviis kunne mødes, er hvert Reale forvandlet til en
% The German phrase is printed here with a T sort, "Ting an sich": the glyph is the
% open, top-barred Fraktur T of "Til" two lines below, not the closed-bowl D of
% "Dele" and "Differentiation" on this same page. Compare p. 74, which sets the
% same phrase „Ding an sich“ with the D sort. Transcribed as printed.
„Ting an sich“, og maa forsvinde som Reflex i et Ikke-Jeg; forsaavidt Realernes bestemte Væren derimod skal finde sit Udtryk i bestemte Relationer, maae de nødvendigviis gjennemtrænges af Reflexionen, thi al gjennemført Relation er, væsentlig seet, iboende Reflexion.

Til at opklare Forholdet mellem Relation og Reflexion ere de mathematiske Analyser særdeles skikkede. Skulde Forestillingen om ligegyldige Relationer med tilsvarende „zufällige Ansichten“ finde Anvendelse paa nogensomhelst Art Forholdsled, da maatte det vel være paa de mathematiske Størrelser, naar, vel at mærke, hver Størrelse med sin umiddelbare Begrændsning betragtes i og for sig. Imellem Materiens Dele er der dog en Vexelvirkning og forsaavidt en væsentlig Relation; men imellem de rene mathematiske Enere, de rene mathematiske Punkter og de rene mathematiske Linier er der, forsaavidt ethvert af disse Elementer betragtes i og for sig, ingensomhelst væsentlig Relation: de kunne i saa Henseende gjælde for ægte, tilforladelige Realer. Men hvo indseer nu ikke, at en ubetinget Fastholden for sig af de enkelte Led som rene Elementer i deres rene Enkelthed er en reen Fixidee? Hvad allerede Ligningen $\frac{y}{x} = c$ tilstrækkelig viser, at Størrelserne $y$ og $x$ kun gjælde for sig, hvad de gjælde i Forholdet, bliver, om muligt, endnu tydeligere, naar vi besinde os paa, at en Differentiation af Ligningen $y = f(x) = cx$ jo ligeledes giver os Forholdet, nemlig: $\frac{dy}{dx} = \frac{0}{0} = c$. Kunde man end,
% --- p. 74 ---
\setcounter{footnote}{0}%
bedaaret af den umiddelbare Forestilling, et Øieblik fristes til at fastholde de endelige Størrelser $x$ og $y$ som Led for sig, der skulde have Tilværelse og Betydning ogsaa udenfor Forholdet, saa maa ethvert Forsøg paa at holde Forholdsleddene fast udenfor Forholdet dog vise sig aldeles urimeligt, naar vi see Forholdet, Relationen selv, blive staaende, uagtet de rene Elementer, Relationens umiddelbare Led, forsvinde i det Uendelige og tabe sig i Nul. Det gjælder om Tingene \textit{in concreto}, hvad der gjælder om Størrelserne \textit{in abstracto}: de ere for sig, hvad de ere i og ved deres Relationer; et absolut Reale, en „Ding an sich“ uden Egenskaber, er en ligesaa puur Fiction, som en absolut Ener, en Størrelse „an sich“ uden Relation til andre Størrelser.

Vexelforholdet mellem Størrelsen og Relationen viser sig nu tillige som et Vexelforhold mellem Relationen og Reflexionen. At den empiriske Verden er en Verden af Relationer, maa vel Enhver indrømme; derimod have selv skarpsindige Tænkere ikke altid været opmærksomme paa den heri skjulte Sandhed: at det netop er Relationerne selv, de objective Relationer, hvorpaa en objectiv Verden skal kjendes. Havde Fichte, istedenfor at indskrænke Relationen $A = A$ til Jeg $=$ Jeg, gjort Alvor af at undersøge Relationerne imellem $A = A$, $B = B$, $C = C$, $D = D$ o. s. v., da vilde en Ligning som $\frac{A}{B} = \frac{C}{D}$ allerede have været tilstrækkelig til at bevise den store Tænker, at der i Ikke-Jeg selv er et System af i og for sig gyldige Relationer, at der, med andre Ord, gives en objectiv Verden. At sætte $a = b$ og $b = c$ kan vel siges at beroe paa Jeget; $a$, $b$ og $c$ ere forsaavidt uden objectiv Selvstændighed. Men naar den reflecterende Subjectivitet paa denne Maade har gjort sig til Herre over visse Forudsætninger, saa indstille sig med det Samme Conseqvenser, over hvilke den ikke er Herre. Af Forudsætningerne: $a = b$, $b = c$ fremgaaer jo Følgesætningen: $a = c$ med en saa
% --- p. 75 ---
\setcounter{footnote}{0}%
uafviselig Nødvendighed, at Subjectiviteten selv i slige Følgesætninger maa bøie sig for de objective Relationer som for en objectiv Magt.

Den objective Magt er desuagtet ingen fremmed Magt, ingen blind Skranke, hvormed Jeget igjen skal brydes; tvertimod! det er just i den saaledes bestemte objective Magt, at Subjectiviteten erkjender sin egen subjective Magt. Relationerne ere Conseqvensernes System iboende og blive kun ved deres udtrykkelige Modsætning til den reflecterende Subjectivitet objective; Reflexionerne foregaae i Subjectiviteten og blive kun ved deres udtrykkelige Modsætning til de Conseqvenserne iboende Relationer subjective og oversvævende.

Forstaaelsen begynder med, at de objective Relationer og de subjective Reflexioner Punkt for Punkt svare til hinanden. Men det er kun en Begyndelse. Vilde den reflecterende Subjectivitet definitivt blive staaende ved en endelig Paralleliseren af vore Reflexioner med Tingenes Relationer og indskrænke sig til i from Beundring at henføre Overeensstemmelsen enten umiddelbart til en guddommelig Sanddruhed, der ikke „vil bedrage“, eller til en \textit{harmonia præstabilita}, der ikke kan bedrage, da maatte Forstanden paa disse Vilkaar blive staaende ved en paafaldende Misforstaaelse af sit eget Væsen. Idet Relationens Led forvandles til Reflexer i Relationen selv, forvandles Relationen med det Samme til objectiv Reflexion; Leddenes Relation er kun gjældende for Leddene selv, forsaavidt den er gjældende for en over Leddene reflecterende Subjectivitet, og omvendt. Her er ikke en Række Relationer at stille op ved Siden af en Række Reflexioner; thi ligesom ingen Relation kan være uden Reflexion, saaledes kan heller ingen Reflexion være uden Relation. At antage et System af objective Relationer, uden deri at ville see et System af objective, Ledreflexerne iboende Reflexioner, er ligesaa urimeligt, som at antage et System af objective Reflexioner uden deri at see et System af subjective, Reflexbestemmelserne over\-%
% --- p. 76 ---
\setcounter{footnote}{0}%
svævende Relationer. Den objective, Leddene selv iboende Reflexion er Relationernes og de oversvævende Reflexioners reflecterede Eenhed; den subjective, Relationerne og deres Led oversvævende Reflexion er Relationernes og de iboende Reflexioners reflecterede Eenhed.

% The printed label is "b)", read at 600 dpi (clear ascender and closed bowl; not "e").
% The lettered series had reached "d)" at p. 60, so this label appears to restart;
% transcribed as printed. The book's synopsis puts Væsensreflexion at pp. 77--123,
% but the printed head stands here, on p. 76.
\subsubsection*{b) Væsensreflexion.}

Forstandsreflexionen er overalt ledsaget af den Biforestilling, at Selvhed og Andethed, det Subjective og det Objective, befinde sig i en endelig Adskillelse fra hinanden. Det er det Endelige i Adskillelsen, paa hvis Ophævelse den sig selv forstaaende Forstand uophørlig reflecterer. Allerede Adskillelsen af sættende og forudsættende Reflexion viser, at dersom Andetheden, det Ikke-Subjective, kun har en negativ, forsvindende Væren, da maa det Samme gjælde om Selvheden, det Ikke-Objective: er Andetheden Skin, saa er Selvheden ogsaa Skin, og omvendt; er Selvheden Væsen, saa er Andetheden ogsaa Væsen, og omvendt. Væsenet er, som Hegel træffende har viist, den sig ved Negation af sin egen Negation bekræftende Væren. Væsensreflexionen er i det Hele hos Hegel mesterlig gjennemført, hvad det Formelle angaaer. Det er kun at undres over, at Analyser, der i Fiinhed og Skarphed endnu staae uopnaaede, ikke have kunnet overbevise Tænkeren om, at Væsenet ikke kan reflectere sig i objectiv-abstract Almindelighed; at det i Væsenet selv Oprindelige netop er Selvhed, Subjectivitet. Hvad det vil sige, at Andetheden, det Objective, bliver sat som det afledede Væsen, idet Selvheden, den sig selv forstaaende Subjectivitet, udtrykkelig sætter sig som det oprindelige, kan Forskjellen mellem den modsættende og eenhedsættende Reflexion allerede vise. Af den gjennemreflecterede, i Reflexionen fordoblede Modsætning fremlyser det nemlig, at Andethedsreflexionen, den optisk-objective Reflexion, hvorved det ene Led i Andethedens uendelige Række speiler sine Bestemmelser i Bestemmelser af det andet, er væsentlig forskjellig
% --- p. 77 ---
\setcounter{footnote}{0}%
fra Selvreflexionen, den Væsensreflexion, hvorved den Alt forudsættende Subjectivitet i Modsætningernes Mangfoldighed hævder sin Eenhedsætten, sin uendelige og concrete Selvsætten. Idet enhver endelig Adskillelse af modsættende og eenhedsættende Reflexion igjen ophæves i reflecteret Eenhed, sees det udtrykkelig, at Modsætningen af Selvhedens og Andethedens Reflexion er en Modsætningseenhed. Saaledes fremkommer da tilsidst en endelig Forstandsadskillelse imellem Reflexioner, som, ifølge Biforestillingen, blot ere subjective, og Relationer, som gjælde for det Objective. Det Endelige i denne Forstandsadskillelse er ikke fjernet derved, at Forskjellen mellem subjectiv og objectiv Reflexion fortoner sig i tvetydig Eenhed; det Endelige bliver just opløst og dog bevaret derved, at Modsætningen mellem Reflexionen som et blot Subjectivt og Relationen som et blot Objectivt forvandles til en i Reflexionen selv faldende Modsætning, Modsætningen nemlig mellem den oversvævende og den iboende Reflexion. Istedenfor at falde udenfor Reflexionen som en mod Viden ligegyldig Væren, er Systemet af de objective Relationer tvertimod forvandlet til et System af Mellembestemmelser, hvori den iboende og oversvævende Reflexion udvikle deres reflecterede Eenhed, altsaa til et System af Bestemmelser, der tilhøre Reflexionen selv. Ved alsidig Ophævelse af de endelige Adskillelser efter deres Umiddelbarhed bliver Forstandsreflexionen efterhaanden udklaret til Væsensreflexion.

Med særligt Hensyn paa Forholdet mellem Væsen og Phænomen, Phænomen og Grund, Grund og Existens opfatte vi Væsensreflexionen i denne Sammenhæng som \emph{phænomenal, som grundende og som begrundende} Reflexion.

\subsubsection*{α) Den phænomenale Reflexion.}

At Tænkningen i sin abstract negative Form falder udenfor Anskuelsen, har kun Betydning, forsaavidt Reflexionens ueensartede Momenter adskille sig for Reflexionen selv;
% --- p. 78 ---
\setcounter{footnote}{0}%
men en saadan Adskillelse er da ligesaa uendelig forsvindende som uendelig tilbagevendende: Momenterne ere den reflecterede Eenheds egne Reflexer.

Den phænomenale Reflexion forudsætter naturligviis en umiddelbar Adskillelse af sandselig og intellectuel Anskuelse. Dersom da en reflecterende Forstand, der endeliggjør sig selv ved at holde det, hvorover der reflecteres, udenfor Reflexionen, faaer Lov til at slaae Adskillelsen fast, saa har den kritiske Philosophie jo vistnok Ret, naar den forkaster den intellectuelle Anskuelse som et eget \textit{cognoscendi genus}; men den endelige Adskillelse af sandselig og intellectuel Anskuelse maa ved Tilsyneladelsens Reflexion uophørlig omsættes i Eenheden. Hvad er vel intellectuel Anskuelse Andet end en i det Tilsyneladende reflecteret Eenhed af Tænken og Anskuen? Dersom Anskuelsen ikke i sig væsentlig var intellectuel, kunde den umulig bestemmes ved nogetsomhelst Intellectuelt; dersom det Intellectuelle ikke i sig væsentlig var et Anskueligt, kunde det umulig bestemmes ved Tilsætning af nogetsomhelst Anskueligt. Det Ene af to Ueensartede kan ikke umiddelbart bestemme det Andet; Farvebestemmelser kunne f. Ex. ikke anvendes paa Tonerne, efterdi Farver og Toner ere ueensartede. Forholdt det sig med Adskillelsen af de rene Anskuelser og de rene Forstandsformer, som Kant antager, da vilde det være umuligt, hvad dog Kant fordrer, at bestemme Anskuelser ved Begreber og Begreber ved Anskuelser.

Er Forskjelseenheden af Anskuelse og Tænkning fra subjectiv Side først bleven gjennemsigtig i den phænomenale Reflexion, saa følger fra objectiv Side Reflexionseenheden af Phænomenet selv og Phænomenets Substrat med uafviselig Conseqvens. At der i ethvert Phænomen, ethvert Tilsyneladende, maa gjøres Forskjel imellem Tilsyneladelsen for sig og det, som deri lader sig tilsyne, er en Selvfølge%
% printed mark: *)  — the note is German set in Fraktur, and runs over onto p. 79;
% the printed page break inside it falls after "einer Erscheinung".
% "Erscheinigung" is so printed (Kant has "Erscheinung"); transcribed as printed.
\footnote{„Dies war das Resultat der ganzen transcendentalen Aesthetik, und es folgt auch natürlicher Weise aus dem Begriffe einer Erscheinung überhaupt, daß ihr etwas entsprechen müsse, was an sich nicht Erscheinung ist, weil Erscheinung nichts für sich selbst und außer unserer Vorstellungsart sein kann, mithin, wo nicht ein beständiger Zirkel herauskommen soll, das Wort Erscheinigung schon eine Beziehung auf Etwas anzeigt, dessen unmittelbare Vorstellung zwar sinnlich ist, was aber an sich selbst, auch ohne diese Beschaffenheit unserer Sinnlichkeit (worauf sich die Form unserer Anschauung gründet) Etwas d. i. ein von der Sinnlichkeit unabhängiger Gegenstand sein muß“. Kritik der reinen Vernunft, Leipzig 1838, S. 247--48 Anm.}; men
% --- p. 79 ---
\setcounter{footnote}{0}%
Spørgsmaalet er, hvorledes det i selve Væsensreflexionen forholder sig med denne Fordobling. Det er en Forskjelseenhed af Væsen og Skin, hvorpaa Eftertrykket falder. Kunde der gives Væsensbestemmelser saa absolut væsentlige, at de slet ikke havde med Skinnet, end ikke med et Skin af Skinnet, at gjøre, da maatte man unegtelig indrømme Kant, at der kunde tænkes et Væsen i sig, der holdt sig udenfor ethvert Phænomen. Men hvad bliver der af det rene Væsen, naar ethvert Skin forsvinder? Det bliver selv et Skin. Ligesaa vil Skinnet, naar Væsenet er blevet Skin, fordoble sig til et Skin af Skin, et i det flygtige Skin bestaaende Skin, og saaledes paa eengang vise sig baade som Væsen og Skin. Som Væsenets reflecterede Negation er Skinnet Væsenets Form, som Skinnets reflecterede Position er Væsenet Skinnets Indhold: det af Væsenet opfyldte Skin er i reflecteret Umiddelbarhed et Phænomen.

Her kan da umulig være to Rækker af Bestemmelser: rene Væsensbestemmelser, der ikke skulde vedkomme Phænomenet, og rene Phænomenalbestemmelser, der ikke skulde vedkomme Væsenet selv. Gjælder det nu i objectiv Forstand, at alle Væsensbestemmelser væsentlig ere Phænomenalbestemmelser, saa gjælder det tillige i subjectiv Forstand, at alle væsentlige Phænomenalbestemmelser maae være Væsensbestemmelser. Et Væsen i sig, hvis objective Indhold var aldeles ueensartet med vore subjective Anskuelses- og Forstandsformer, vilde
% --- p. 80 ---
\setcounter{footnote}{0}%
ikke paa nogen Maade kunne udfylde de tomme subjective Former.

Det Urimelige i den Betragtning, at der bagved Phænomenerne skulde skjule sig et utilgængeligt Substrat, et Noumenon, der var saa absolut positivt, at det kun lod sig bestemme negativt, har Fichte fortræffelig paaviist; men idet han forvandlede den positive Ting i sig til det negative Ikke-Jeg, og trak den objective Factor i Form af en dunkel Nødvendighed ind under Jeget selv, forfeilede ogsaa han den phænomenale Reflexion, kun fra et modsat Synspunkt.

For at gjennemskue den phænomenale Reflexion, er det nemlig ikke nok at gjøre Forskjel imellem det Reflecterede og det Umiddelbare, imellem Reflexionen selv og det, hvorover der reflecteres; det maa udtrykkelig bringes til Bevidsthed, at det Umiddelbare, hvori Reflexionen har sin Modsætning, er sat og forudsat af Reflexionen selv, at det concret Umiddelbare, som et i sig Reflecteret, netop maa for Reflexionens Skyld erholde relativ Selvstændighed.

For den kantiske Kritik blev det i Objectet Umiddelbare staaende ligeoverfor den derpaa reflecterende Reflexion; ligeledes blev, hvad det Subjective angaaer, den ene af Reflexionens Hovedsider stillet som et umiddelbart Givet ligeoverfor den anden, Rumsanskuelsen ligeoverfor Tidsanskuelsen, de rene Sandseanskuelser ligeoverfor Forstandskategorierne: det er en rationel Udstykning og en rationel Sammenstykning, der charakteriserer „den rene Fornufts Kritik“; den kommer ikke ud over Stykværket. Det ved Fichte betegnede Fremskridt%
% printed mark: *)  — German note set in Fraktur, with antiqua "a priori".
% "Vst." so printed, for "Vernunft".
\footnote{„Kant geht in der Kritik d. r. Vst. von dem Reflexionspunkte aus, auf welchem Zeit, Raum und ein Mannigfaltiges der Anschauung gegeben, in dem Ich, und für das Ich schon vorhanden sind. Wir haben dieselben jetzt \textit{a priori} deducirt, und nun sind sie in dem Ich vorhanden.“ Wissenschaftslehre 2te Lieferung S. 108. Schluß-Anmerkung.} er
% --- p. 81 ---
\setcounter{footnote}{0}%
en Ophævelse af Stykværket, en Deduction af det umiddelbart Givne; men idet Reflexionen fra alle Sider arbeider paa at assimilere sig det Umiddelbare, fortærer den Umiddelbarheden selv og hilder sig saa i den Modsigelse, at den hverken kan døie det saaledes Fortærede eller være det foruden. Det er Væsensreflexionen i subjectiv Form, den uendelige Selvreflecteren, der i Fichtes Jeg er Eet og Alt; men denne Reflecteren befinder sig endnu i abstract Modsætning til Phænomenernes Mangfoldighed. For at reflectere over sig selv maa Jeget naturligviis tillige reflectere over Phænomenerne; men til en affirmativ Bestemmelse af, hvad Phænomenet er, kan den eensidige Selvreflecteren, trods alle Anstrengelser med „aprioriske Syntheser“ af $X$ og $Y$, af $A + B$ ved $a$ og $b$ o. s. v. dog ikke bringe det. Det gaaer med Phænomenet i Fichtes Ikke-Jeg, som med Ellepigen i Eventyret: det er huult i Ryggen. Var det Fornuftkritiken umuligt at gjennemskue Phænomenet, fordi den vilde Tingen selv i absolut Umiddelbarhed, saa er det nu den subjective Idealisme umuligt at tilkjende Phænomenet relativ Selvstændighed, fordi den ikke indseer, at Selvreflexionen forudsætter en Reflexion af Andet, en phænomenal Reflexion, hvori det Reflecterede selv er umiddelbart og det Umiddelbare reflecteret.

At det Umiddelbare er et i sig Reflecteret, og det saaledes Reflecterede et Umiddelbart, er vistnok en af de haardeste Modsigelser, der kan bydes den reflecterende Forstand. Saalænge Modsigelsen ikke er gjennemskuet, vil Forstanden da ogsaa idelig fristes til at antage Eet af To: enten at Modsigelsen ligger i Reflexionen, eller at den ligger i Umiddelbarheden. Fornuftkritiken antager det Første; den subjective Idealisme det Sidste. Fra begge Synspunkter bliver Phænomenet miskjendt; og dog er det saa langt fra, at Modsigelsen skulde ligge i Phænomenet, at Phænomenet just er Modsigelsens Løsning.
% Signature marks on p. 81, not transcribed: the gathering title
% "Grundideernes Logik." and the sheet numeral "6" at the foot.

% --- p. 82 ---
\setcounter{footnote}{0}%
Som Reflexion af Andet i Andet er Phænomenets Væsen
en reflecteret Umiddelbarhed. Fornuftkritikens Modsigelse er
nu den, at Phænomenet skal være en Tilsyneladelse af en
Ting, der dog ikke lader sig tilsyne. Ved at fastsætte Tingen
i sig har den reflecterende Forstand drevet Reflexionens egen
Modsætning til den yderste Spidse af forstenet Umiddelbarhed:
den har forsaavidt sikkret sig en afgjørende Realitet, som den
har sikkret sig en uopløselig Andethed. Men idet Reflexionen
paa denne Maade vil sikkre sig Realiteten i Tingen, gaaer
Tingen i sig selv tilbage i det Ubekjendte, og lader kun
Phænomenet som et i Reflexionen forvansket Efterbillede blive
staaende. Forstanden har her meget rigtig indseet, at Væsenet
i og for sig er Et, og det Skin, hvori det reflecterer sig, et
Andet; men den har endnu ikke indseet, at den i Reflexionen
umiddelbare Forskjel er reflecteret i en ligesaa umiddelbar
Eenhed, at Forskjellen er tilsyneladende, fordi Eenheden er
tilsyneladende, at her fra alle Sider er Reflexion af Andet i
Andet, og at denne Reflexion just er det Phænomen, hvori
Væsenet selv umiddelbart lader sig tilsyne.

Seet fra den subjective Idealismes Standpunkt er Phænomenet
af dunkel Nødvendighed en Tilsyneladelse, hvori Intet lader
sig tilsyne. Det Skuffende i Phænomenet beroer da ikke
derpaa, at det skulde være som en Maske, hvorunder Tingen
selv bevarede sit Incognito; Skuffelsen er meget mere den, at
den af Umiddelbarheden blendede Reflexion i Phænomenet
bestandig vil søge en Ting, og see, saa er det dog i
Phænomenet ingen Ting. Det sig i Alt reflecterende Jeg har
rigtig indseet, at der ikke kan gives en Umiddelbarhed saa
fast, at den jo lader sig gjennemtrænge af Reflexionen; men,
hildet i eensidig Selvreflecteren, har det endnu ikke indseet, at
der gives en Reflexion af Andet i Andet, en Reflexion, der
affirmerer sin egen Modsætning, idet den indrømmer det
Reflecterede umiddelbar Tilværen og i det Tilværendes
reflecterede Umiddelbarhed erkjender det bestemte Phænomen.
Det
% --- p. 83 ---
\setcounter{footnote}{0}%
er ikke Andetheden og heller ikke Umiddelbarheden, her
ligefrem bliver afviist; tvertimod! idet her idelig tales om
Skranke og mange Skranker, om Realitet og Negation, om
Stød og Anstød, Modstød o. s. v., saa mærkes det tydeligt
nok, at den umiddelbare Andethed ikke paa nogen Maade
lader sig afvise. Men hvorfor ere da Phænomerne saa hule
% printed "Phænomerne" (so both witnesses and the 600 dpi image:
% P-h-æ-n-o-m-e-r-n-e, no "en" between m and r) for "Phænomenerne";
% transcribed as printed.
og væsenløse? Det kommer aabenbart deraf, at Umiddelbarheden
vel bliver opløst af Reflexionen, men ikke tilegnet af
Reflexionen, ikke støbt igjen i Reflexionens Former, ikke præget
til eiendommelig bestemte Phænomener. Det Umiddelbare, hvis
umiddelbare Form idelig forstyrres, uden at det derfor vinder
Selvstændighed og Skikkelse i en ny og høiere Form, har kun
en elementær Tilværelse, kan kun gjælde for et Stof. Det er
just i denne stofagtige Umiddelbarhed, at Kategorierne:
Skranker, Anstød, Modstød o. dsl. bruges i den subjective
Idealisme; just fordi Reflexionen af Andet er tom, er
Selvreflexionen ogsaa tom.

Vexelbestemmelsen af Selvreflexionen og Reflexionen af
Andet faaer først sin Udviklings positive Indhold i den
phænomenale Reflexion. Kun af den phænomenale Reflexion
sees det, at der gives udvortes tilværende Væsener, Væsener,
der reflectere sig i hinanden objectivt, uden dog at kunne
reflectere sig for sig selv subjectivt. Ethvert saadan Væsen
tager sig ved første Øiekast ud som et umiddelbart Tilværende;
dets umiddelbare Tilværen er imidlertid kun en Tilsyneladelse:
det er i sig en Negation af det Umiddelbare, et Reflecteret, et
Væsen. Idet Forstanden nu fatter, at det første Umiddelbare
ikke kan være det sande Objective, kan den jo vistnok meget
let komme til at erklære det umiddelbart Tilværende for lutter
Skin. Men, naar den saa, ifølge Herbarts: „wieviel Schein
soviel Hindeutung auf Seyn“, tillige bilder sig ind at kunne
fange det sande Værende bagved Skinnet, da er det en
Skuffelse: Væsenet er i Skinnet selv, i Reflexionen; det sande
Værende lader sig tilsyne. Mener endelig „den speculative
% signature mark "6*" at the foot of p. 83 (third leaf of gathering 6);
% not transcribed.
% --- p. 84 ---
\setcounter{footnote}{0}%
Fornuft“, at, efterdi Umiddelbarheden fra alle Sider
reflecteres i Tilsyneladelsen, saa maa den phænomenale
Reflexion nødvendigviis være saa alsidig, at den, uden
oversvævende Reflecteren, kan paa immanent Maade reflectere
sig selv, altsaa, med andre Ord, forvandle sig til en
Phænomenerne iboende Selvreflexion, da er ogsaa dette en
Skuffelse. Er Phænomenets Umiddelbarhed i Modsætning til
Reflexionen en Tilsyneladelse; er det objective Væsens
Reflexion i sit Phænomen en Tilsyneladelse; er det
eiendommelig bestemte Phænomens Reflexion i Forskjelligheden
fra andre Phænomener en Tilsyneladelse: da er det sig
reflecterende Phænomens Selvreflexion i Sandhed ogsaa en
Tilsyneladelse. Hvad er det vel, der uafbrudt forhindrer det
som Phænomen sig reflecterende Væsen fra at opnaae
Selvreflexion? Svar: det er Andetheden, det er
Umiddelbarheden. Istedenfor en sig bestemmende Selvhed har
Phænomenet en tilsyneladende Selvhed; denne dets
tilsyneladende Selvhed er just dets umiddelbare Selvstændighed,
dets Existens. Det phænomenale Væsen er en Existens, og det
Existerende en Ting. At Tingene reflectere sig i deres
Egenskaber og formedelst Egenskaberne i hinanden indbyrdes,
vil da med andre Ord sige, at Tingene lade sig tilsyne i deres
Egenskaber, og at de igjennem Egenskaberne lade sig tilsyne
mod hinanden. Men Tilsyneladelsen er fra alle Sider Reflexion
i Andet; Tingene kunne umulig lade sig tilsyne mod hinanden,
uden derved at de lade sig tilsyne for en betragtende
Reflexion: det i Tilsyneladelsen Objective forudsætter det
Subjective, og omvendt.

Kan Tingenes Tilsyneladelse mod hinanden indbyrdes ikke
være abstract objectiv, da kan den enkelte Tings Tilsyneladelse
i sine Egenskaber heller ikke være det: kun ved at lade sig
tilsyne i en Reflecteren, der ikke selv er tilsyneladende, kan
Tingen lade sig tilsyne i sine Egenskaber. En Tilsyneladelse,
der blot er subjectiv, er falsk; det er en Hallucination, et
Skin, en Skuffelse. Men en Tilsyneladelse, der skulde
% --- p. 85 ---
\setcounter{footnote}{0}%
være blot objectiv, er ogsaa falsk; den phænomenale Reflexion
er kun objectiv, forsaavidt den er subjectiv; den er kun
iboende, forsaavidt den er oversvævende: Tingene reflectere sig
i deres Relationer derved, at de reflectere sig for den
reflecterende Subjectivitet.

\subsubsection*{β) Den grundende Reflexion.}

Imellem Selvreflexionen og Reflexionen af Andet eller, mere
anskueligt, imellem Subjectiviteten selv og Phænomenernes
objective Verden har der altsaa udviklet sig en væsentlig og
dybtgaaende Modsætning; skal denne Modsætning ikke
forvandle sig til en Dualisme, der gjør al Forstaaelse umulig,
da maa her nu tillige vise sig en gjennemgaaende
Modsætningseenhed. Det er igjen den reflecterede
Umiddelbarhed, hvorom Spørgsmaalet dreier sig. Vistnok har
den phænomenale Reflexion allerede belyst os denne
Umiddelbarhed; men Lyset falder endnu kun ind fra den ene
Side. Den phænomenale Reflexion tog i en vis Maade Parti
for Phænomenerne og deres Realitet som existerende Ting; den
gav os overalt det Umiddelbare reflecteret i umiddelbar og
positiv Form. Men Reflexionen er kun positiv, forsaavidt den
væsentlig er negativ, positiv ved sit Anskuelsesmoment, negativ
ved sit Tankemoment; den reflecterede Umiddelbarhed maa som
Følge heraf nødvendigviis tillige have sin negative Form. Den
reflecterede Umiddelbarheds negative Form udvikler sig i
\emph{den grundende} Reflexion; thi den grundende Reflexion
% the letterspacing stops at "grundende": "Reflexion" on the same line
% is set plain, and so is the second "den grundende Reflexion".
løser Umiddelbarhedens Indhold op i Selvreflexionen, og danner
saaledes en betegnende Modsætning til den phænomenale. Al
Grunden beroer væsentlig paa reflecteret Eenhed af
Selvfordybelse og Fordybelse i Sagens, d. e. i
Umiddelbarhedens Indhold: den grundende Reflexion søger paa
een Gang Subjectivitetens Væsen i Tingene og Tingenes i
Subjectiviteten.

% --- p. 86 ---
\setcounter{footnote}{0}%
Den grundende Reflexion søger Grunden; men hvad er
Grunden? Fra hvilken Side vi end opfatte Begrebet Grund,
det viser os bestandig hen til en Umiddelbarhed, som har
været, en Umiddelbarhed, der reflecterer sig i en negativ Form.
Saaledes er f. Ex. „Tingenes Grund“ jo vistnok i een
Henseende det Samme som „Tingenes Væsen“, og dog er her
imellem Væsen og Grund en væsentlig Forskjel. I Væsenet er
det Umiddelbare negeret derved, at det er forvandlet til Skin;
i Skinnet som Skin er Umiddelbarheden forsvunden:
Negationen af det Umiddelbare er i sit Skin kun umiddelbar
som Negation, men just derfor ikke selv en negativ, værende
Umiddelbarhed. Hvad her gjælder om Skinnet, gjælder ogsaa
om Væsenet: vil man umiddelbart gribe det i Skinnet, saa
griber man et Negativt, og vises hen til det, der er bagved
Skinnet; vil man umiddelbart gribe, hvad der er bagved
Skinnet, da griber man igjen et Negativt, og vises hen til,
hvad der sees i Skinnet. Væsenet er kun i sig, hvad det er i
sit Skin: det skinner i Skinnet, et Skin af Skin; det er fra
alle Sider negativt bestemt, og følgelig uden Umiddelbarhed.
Anderledes med Grunden. Forsaavidt Væsenet opfattes som
Grund, har det, i Negationen af det Umiddelbare, selv en
Umiddelbarhed. Dette sees tydeligt nok af den Forandring,
der foregaaer med Reflexionen, idet Forskjelseenheden af Væsen
og Skin forvandles til en Modsætningseenhed af Grund og
Følge. Hvor afgjørende Umiddelbarheden er ved Bestemmelsen
af Forholdet mellem Grund og Følge, fremlyser strax deraf, at
Grunden kan opfattes for sig som Forudsætning, og Følgen
som Forudsætningens betingede Conseqvens. Fra dette
Synspunkt sætter Subjectiviteten sit Videns Grundlag i et
System af faste, almeengjældende Forudsætninger, de saakaldte
Grundsætninger. I Grundsætningernes System reflecteres alle
Væsensbestemmelser som umiddelbare Værensbestemmelser;
Grundlaget, den værende Grund, er netop en Reflexion af
Grundens Umiddelbarhed. En tilsvarende Umid\-%
% --- p. 87 ---
\setcounter{footnote}{0}%
delbarhed sees i Conseqvensernes System. Vistnok kan der i
Conseqvensen ikke fremkomme Andet end, hvad der alt har
ligget i Forudsætningen; Conseqvensen er forsaavidt kun en i
Reflexionen omsat, i formel Omskrivning gjentagen
Forudsætning. Ikke destomindre er her en umiddelbar
Adskillelse af Forudsætning og Conseqvens, hvis Betydning
ikke tør oversees. Den samme Forudsætning kan nemlig have
forskjellige Conseqvenser, og den samme Conseqvens fremgaae
af forskjellige Forudsætninger. Imellem Forudsætningen og
dens Conseqvenser falde nemlig, umiddelbart seet, en
Mangfoldighed af vexlende Betingelser; disse Betingelser kunne
igjen umiddelbart deels adskilles fra, deels falde sammen med
den bestemte Forudsætning, o. s. v.

Den her paapegede Umiddelbarhed hentes ikke udenfra, som
fra Indtryk af en blot empirisk Mangfoldighed; den ligger i
Grunden som Grund og udgjør en Hovedside af dens logiske
Charakteristik. De exacte og physiske Videnskaber blive
væsentlig staaende ved denne Umiddelbarhed, og mange
Philosopher, f. Ex. Herbart, mene i Grunden ikke Andet end
en almeen Forudsætning i en Sum af Betingelser, naar de tale
om Grunden. Men saalidtsom Væsenet, for at være Grund,
kan undvære det Umiddelbare, ligesaalidt kan
Umiddelbarheden, for at være Grund, undvære Negationen.
Den blotte Forudsætning er endnu ikke Grund, og den
betingede Conseqvens kun en abstract
Følge%
% printed footnote mark: *)
\footnote{„Da Functionsbegreb netop viser os Forudsætningens og
Conseqvensens, det Uafhængiges og det Afhængiges reflecterede
Eenhed, kan Forskjellen mellem den blotte Forudsætning og den
egentlige Grund i denne Sammenhæng oplyses. Forudsætningen er
et Givet, et Antaget, der opfattes i Forestillingen, og kun
forsaavidt falder
% the printed page break of the note falls here: it runs on at the
% foot of p. 88.
under Reflexionen, som det indsees, at den, for at være
Forudsætning, maa have sine Conseqvenser. Man kan saaledes
antage en Forudsætning, fastholde den i Forestillingen som
noget Bekjendt, og dog i høi Grad overraskes af dens
Conseqvenser. Overraskelsen hidrører derfra, at
Forudsætningens Væsen ikke fra Begyndelsen af blev
gjennemskuet; Forudsætningen blev vel tydeliggjort i en
Forestilling, og forsaavidt opfattet fra en vis Side, men ikke
alsidig opklaret, ikke opfattet efter sit Begreb. Først naar vi
gjennemskue Forudsætningen og opdage de deri skjulte
Conseqvenser, erkjende vi Forudsætningens Væsen, sætte den
som Grund og Conseqvensen som den begrundede Følge.“
Mathematik og Dialektik S. 7--8.}. Det er just ved at klare
Umiddelbarheden i negativ Reflexion, at Philosophien gaaer
videre end Fagvidenskaben, naar det gjælder en Undersøgelse
af „Tingenes Grunde“.

% --- p. 88 ---
\setcounter{footnote}{0}%
For at opfattes som Grund, maa Forudsætningen og dens
Betingelser omsættes som Momenter i negativ Reflexion; uden
en saadan Omsætning er Forholdet mellem Grund og Følge
utænkeligt. „Følgen“ --- siger Herbart --- „skal ligge i
Grunden. Men den skal ogsaa følge af den, det vil sige: den
skal afsondre sig derfra. Ligger den nu virkeligt i den, saa
tilhører den Grunden, og den, som vilkaarligt skiller den
derfra, har ikke udvidet sin Viden; han har snarere blot
gjentaget, hvad han allerede kjendte, eftersom han kjendte
Grunden. Lærer Følgen derimod noget Nyt, saa er dette Nye
ikke det Gamle, og laae ikke i Grunden; det kaldes da med
Urette en Følge af samme.“ Men den Udvei, Herbart saa
vælger: at opfatte Følgen som en Deel af Grundens Totalitet,
d. v. s. identisk med en Deel af samme, saa at den derved
bliver „en Gjentagelse i en afsondret Tanke“, er kun en tom
Udflugt; thi dersom den bestemte Følge ligefrem er en Deel af
Grunden, saa gjælder uimodsigelig Eet af To: enten have vi i
det til en vis Grund svarende Indbegreb af Følger kun et
Indbegreb af deelte Grunde, altsaa i Grunden slet ingen
Følger, men en Sum af Grunde; eller ogsaa, dersom Grundens
Dele ikke selv ere Grunde, men Følger, have vi i Grunden
selv ingen Grund, men lutter Følger, Følger uden al Grund.
Skulde her i Indbegrebet af slige følgeløse Grunde og
grund\-%
% --- p. 89 ---
\setcounter{footnote}{0}%
løse Følger blive endog blot et Skin af Grund tilbage, da
maatte Grunden nødvendigviis søges i Følgernes indbyrdes
Sammenhæng og den ved Sammenhængen bestemte Eenhed.
Men hvo seer nu ikke, at en saadan i Følgernes Totalitet
reflecteret Sammenhængseenhed netop er den Totalitet, hvori
Forudsætningen med sine Betingelser er sat med negativ
Umiddelbarhed?

Idet den ene og samme Totalitet udtrykkes i den reflecterende
Umiddelbarheds to forskjellige Former: Grundens negative og
Følgens positive, maa den naturligviis tage sig ud som to
Totaliteter, thi den væsentlige Forskjel i Form er tillige en
Forskjel i Indhold. At see et Qvadrat være ligestort med to
andre, er at see et Umiddelbart reflecteret i positiv Form; thi
vel ere Qvadraterne selv umiddelbare, men deres Ligestorhed
er kun for Reflexionen. At see Grunden til de to Qvadraters
Ligestorhed med et tredie, er derimod at see det Umiddelbare
reflecteret i negativ Form. Opmærksomheden kan fra dette
Synspunkt ikke længere dvæle ved Qvadraternes
Umiddelbarhed eller ved det Umiddelbare i den Kjendsgjerning,
at den omspurgte Ligestorhed virkelig finder Sted; enhver
umiddelbar Bestemthed forvandler sig til en
Relationsbestemmelse. Saaledes sees da hvert Qvadrats
Størrelse i Relation til Siden; Siderne i de tre Qvadrater sees
i deres indbyrdes Relation, forsaavidt de sees som Sider i en
Trekant; Trekantens Sider sees endvidere i Relation til
Vinklerne, det bliver jo en retvinklet Trekant o. s. v. Da nu
Relationen er en kategorisk Mellembestemmelse imellem
Reflexion og Umiddelbarhed, er det tillige heraf indlysende,
hvor let Grunden selv kan blive forvexlet med Betingelserne og
de umiddelbare Forudsætninger.

At Grunden maa søges i Beviset, er uimodsigeligt; men det
kommer an paa, hvorledes man seer Beviset, og hvad man
seer deri. Der kan være mange Beviser for den samme
Sætning, men der kan i dem alle kun være een totaliseret
Grund. Fordi man kan bevise en mathematisk Læresætning,
% --- p. 90 ---
\setcounter{footnote}{0}%
deraf følger endnu ikke, at man strax gjennemskuer Forholdet
mellem Beviis og Grund. Ligesom Følgen er Conseqvensernes
Eenhed, udtrykkelig sat i den reflecterede Umiddelbarheds
positive Form, saaledes er Grunden Forudsætningernes og
Betingelsernes Eenhed, udtrykkelig sat i den reflecterede
Umiddelbarheds negative Form; er Følgen enkelt, saa er
Grunden det ogsaa, er Følgen som Følge Totaliteten i det
Ydre, saa er Grunden som Grund Totaliteten i det Indre: det
Ydre sees, det Indre indsees; Følgen sees, Grunden indsees.

Naar den totaliserede Følge sættes som Existens, er Grunden
med det Samme forudsat som den reale Grund. At den reale
Grund ikke holder sig bagved Existensen, men gaaer op i den,
kommer til Existens i og med det Existerende, sees
udtrykkelig deraf, at det Existerende, under sin Vorden og
Forandring, uophørlig er i Begreb med at opstaae og forgaae.
Men at opstaae er jo at fremgaae af Grunden, at forgaae er
at gaae tilgrunde, altsaa at gaae tilbage i Grunden. Existensen
er en Mellemform, og det Existerende det reale Mellemled,
igjennem hvilket Grunden i uendeligt Kredsløb udgaaer fra sig
selv og vender tilbage til sig selv. I Modsætning til det
Existerendes umiddelbare og positive Realitet er Grunden
væsentlig Idealitet; men denne Idealitet er Phænomenets
negative Realitet.

Som negativ Realitet er Grunden beslægtet med den
schellingske Indifferens; desuagtet er Grunden ikke uden videre
en Indifferens, og Indifferensen ikke Grund: i Indifferensens
Dæmring dæmrer allerede Differensen; begge sees oprindelig i
samme Tusmørke. En Indifferens i sig uden Differens er en
mystisk „Ugrund“, en hyperspeculativ „Urgrund“, der
ligesaalidt fordrager Reflexionens Lys, som „de Underjordiske“
fordrage Daglyset. Opfatter man den rene Indifferens som en
første ubestemt Grund, en grundløs Grund, hvoraf Differensen
med de bestemte Grunde er fremgaaet, hvad gjør man saa?
Man søger en Grund til Grunden, idet man paa
% --- p. 91 ---
\setcounter{footnote}{0}%
ulogisk Maade slaaer Kategorierne: „Grund“, „Princip“ og
„Begyndelse“ sammen i Eet. Den reale Grund er, for sig
betragtet, hverken Princip eller Begyndelse, men et i negativ
Umiddelbarhed reflecteret Mellemled mellem begge. Vistnok
forekommer Ordet \textit{ἀρχή} hos Aristoteles ogsaa i
Betydning af Grund; men deraf kan man ligesaa lidt udlede,
at Grunden uden videre er Princip, som at Princip er
Begyndelse.

At Grunden, uagtet den har Oprindeligheden og det
oprindelige Indhold tilfælleds med Principet, dog ikke er
Principet selv; at, med andre Ord, Principet, opfattet som
Grund, kun er eensidigt og ufuldstændigt opfattet, skjønnes
allerede deraf, at Grundens Indhold ligesaa væsentlig beroer
paa en Resulteren som paa en Anticiperen. Intet kan
fremgaae af Grunden, uden det, som indeholdes i Grunden;
men Intet kan indeholdes i Grunden, uden det, som forud er
gaaet tilgrunde: Grundens rige Indhold hidrører fra de
utallige Existenser, som uophørlig gaae tilgrunde; Tingenes
Udstrømning (Emanation) af Grunden er betinget af
tilsvarende Tilbagestrømning. Medens Principet, det
Ubetingede, altid virker i sit System af Betingelser saaledes, at
det magter disse Betingelser, forvandler dem til Momenter, og
midt i al Vorden og Forandring reflecterer sig som det
Absolute og Uforanderlige, er Grunden, den reale Grund,
derimod saaledes nedsænket i Vorden, Forandring, Endelighed
og Relativitet, at den netop maa siges at reflectere
Foranderligheden selv i Principets Uforanderlighed.

Reflecteret i Vorden og Forandring er Grunden beslægtet med
Begyndelsen; hvad der i Sandhed skal begynde, maa begynde
fra Grunden af; Grunden til Tingene er Tingenes reale
Begyndelse: Grunden er Begyndelsens Realitet. Men hvad her
siges om Begyndelsen, gjælder ligesaa eftertrykkelig om Enden.
Tingen, der gaaer tilgrunde, ender i Grunden; Tingen
forgaaer; men denne Forgaaen er ikke den absolute
Tilintetgjørelse: Tilintetgjørelsen har sin Realitet i Grunden.

% --- p. 92 ---
\setcounter{footnote}{0}%
Negativt bestemt kan Grunden nu vistnok opfattes som en
Indifferens, men denne Indifferens er ingen „Ugrund“. Det er
den reflecterede Umiddelbarheds negative Form, der giver os
Indifferensen: i denne Form er Grunden indifferent mod
Forskjellen mellem Princip og Begyndelse, efterdi den som
begges væsentlige Mellembestemmelse har lige Meget og lige
Lidt af begge; i denne Form er Grunden indifferent mod
Væren og Vorden, mod Uforanderligt og Foranderligt, mod
Begyndelse og Ende, mod Selvhed og Andethed, mod iboende
og oversvævende Reflexion o. s. v.

Forsaavidt den grundende Reflexion med sin negative Form
behersker den phænomenale, er Subjectiviteten eensidig
fordybet i sig selv, og Reflexionen i denne Selvfordybelsens
Eensidighed en Grublen. Forsaavidt Umiddelbarheden i den
negative Form derimod paa alle Punkter trænger sig
umiddelbart frem med existentielle Elementer, der fordre
Skikkelse uden dog at kunne antage bestemt Skikkelse, faaer
Andetheden Overhaand, og den grundende, i Tingenes Væsen
og Vorden fordybede, Subjectivitet befinder sig i en Gjæring.

\subsubsection*{γ) Den begrundende Reflexion.}

Under Subjectivitetens uklare Grublen og Gjæring vil
Eenheden som Ugrundens Ginnungagap uophørlig sluge
Mangfoldigheden, og Mangfoldigheden med sine Grundpunkter
og Spirer igjen opfylde, overfylde og ligesom overvælde
Eenheden; i dette Punkt er Grunden mystisk: Grunden som
Grund er Mystikens Eet og Alt. Den klare Dunkelhed, som
den mystiske Speculation saa let udbreder over de herhen
hørende Spørgsmaal, har i den Grad virket afskrækkende paa
den kritiske Forstand, at mangen skarpsindig Tænker endnu
anseer Undersøgelser af Tingenes Grundforhold for ørkesløse
Speculationer. Men istedenfor at lade sig afskrække af
„Ugrundens“ Tomhed og Indifferensens Overflødighed, maa
den grundende Reflexion just lægge an paa at klare sig selv;
thi idet den
% --- p. 93 ---
\setcounter{footnote}{0}%
klarer sig selv, arbeider den sig med det Samme ud af
Grundens Uklarhed. Opgaven er at bestemme Grunden ved
det Begrundede saaledes, at Eftertrykket falder paa
Begrundelsen.

Den begrundende Reflexion kan i denne Sammenhæng støtte
sig til følgende tre Sætninger: \emph{1) Alt er i Grunden Eet;
2) Alt er i Grunden forskjelligt; 3) Alt maa have sin
tilstrækkelige Grund.} Men det ligger i Sagens Natur, at
slige Sætninger, enkeltviis opstillede, ikke tør tages for Mere
end betegnende Udgangspunkter; det Væsentlige er, ikke
hvorledes Sætningerne formuleres, men hvorledes de udvikles.

\emph{1) Alt er i Grunden Eet.} Sætningen, saaledes opstillet,
maa ikke forvexles med den eleatiske Sætning, at kun det Ene
(\textit{τὸ ἕν}) er; thi hos Eleaterne, der lade det Ene falde
umiddelbart sammen med det Værende (\textit{τὸ ὄν}), kan der
egentlig ikke være Tale enten om nogen Væren i Grunden eller
Væren af Grunden. Vistnok meente Eleaterne med det
Værende det i Sandhed Værende (\textit{τὸ ὄντως ὄν}), altsaa
det Grundværende, thi ellers kunde de umulig have stillet det
imod det Ikke-Værende i Sætninger saa skarpe som disse:
\emph{kun det Værende er, det Ikke-Værende er ikke.}
„Men hvad skal man vel“, spørger allerede Plato, „dømme om
det Ikke-Værende? Er det Ikke-Værende slet ikke, saa kan det
jo hverken tænkes eller udsiges; er det Ikke-Værende et Skin,
saa maa dette Skin jo dog være. Dette er just
Vanskeligheden ved at tænke det Ikke-Værende, at baade den,
der benegter det, og den, der forsvarer det, er nødt til at
modsige sig selv. Indrømmer man, at der gives en falsk
Mening (\textit{ὄντος δέ γε ψεύδους ἔστιν ἀπάτη}), saa
forudsætter man med det Samme Forestillingen af det
Ikke-Værende, thi kun den Mening kan kaldes falsk, som
enten erklærer det Ikke-Værende for værende, eller det
Værende for ikke-værende
% --- p. 94 ---
\setcounter{footnote}{0}%
(\textit{κινδυνεύει τοιαύτην τινὰ πεπλέχθαι συμπλοκὴν τὸ μὴ
ὄν τῷ ὄντι}).“

Det i denne Sammenhæng Charakteristiske er, at Eleaterne
anvende Begrundelsen til at tilintetgjøre al Grund. Skulde
Phænomenernes Verden være begrundet, saa maatte det Ene jo
være Grunden til det Mangfoldige; men Phænomenernes
Verden er, efter Eleaternes „Mening“ (\textit{δόξα}),
grundløs, det Mangfoldiges Væren er et Ikke-Værende, og det
Ikke-Værende er ikke. Disse for den sunde Sands anstødelige
Sætninger skulle imidlertid bevises, altsaa begrundes; ved
deres Begrundelse komme Eleaterne da paa den meest
bagvendte Maade til at begrunde det Grundløse. Hvorfor kan
det Ene ikke være Grund? hvorfor kan det Mangfoldige ikke
være begrundet? Eleaterne svare correct: „al Grund og
Begrundelse forudsætter Vorden.“ Dette Svar ligger vel ikke
ligefrem i deres Ord, men det ligger i deres Tankegang: det
ligger i deres Angreb paa „Vorden“. „Vorden (Opstaaen og
Forgaaen) er utænkelig; Intet vorder eller forgaaer. Intet kan
opstaae: det maatte da enten opstaae af Noget, der var det
ligt, eller af Noget, der var det uligt. Det kan ikke opstaae
af det Lige, thi det Lige forholder sig til det Lige paa en lige
Maade%
% printed footnote mark: *)
\footnote{Poul Møller anmærker hertil: „Det forstaaer jeg
saaledes: det Lige kan ikke opstaae af noget andet Ligt; thi
da vare de deri ulige, at det Ene var det, hvoraf Noget
opstod, og det andet det, som opstod selv: følgelig havde de
to foregivne Lige ulige Prædicater og vare saaledes ikke
lige.“ Udkast til Forelæsning over d. gl. Phil. Xenophanes.
(Efterladte Skrifter, 4de Bind, S. 59).}. Det kan heller ikke
opstaae af det Ulige, thi da skulde Noget komme af Intet. Det
Svagere kan ikke komme af det Stærkere, det Stærkere ikke af
det Svagere: thi da skulde der komme Noget af en Ting, som
den ikke selv var, altsaa Noget af Intet.“ Paa denne Maade
søge
% --- p. 95 ---
\setcounter{footnote}{0}%
Eleaterne altsaa at begrunde, at der hverken gives Grund
eller Begrundelse.

Men Sætningen: Alt er i Grunden Eet, kan heller ikke
forstaaes, som om der gaves et gnostisk Pleroma, i hvis Dyb
Grundeenheden hvilede. Vel maae vi indrømme Gnostikerne,
at den oprindelige Grund er en indholdsrig Eenhed, og
Eenheden det Alt Begrundende; men det i Grundeenheden
Umiddelbare er reflecteret, og Reflexionen negativ. Det
Dybsindige i Gnostikernes, navnlig i Valentinianernes
anskuelige og naive Fremstilling af Urgrunden og de
ideal-reale Modsætninger, som deraf fremgaae og deri røre sig,
er forøvrigt ikke at miskjende. „I Urgrundens Dybde
(\textit{βυθός, προπάτωρ})“ --- hedder det --- „var
Selvbevidstheden (\textit{ἔννοια}) og Tausheden
(\textit{σιγή}). Den i Dybet skjulte Guddom, Urgrunden selv,
aabenbarede sig i 3 Rækker af Æoner: \emph{Fornuften og
Sandheden} (\textit{νοῦς} og \textit{ἀλήθεια}), \emph{Ordet og
Livet} (\textit{λόγος} og \textit{ζωή}), Mennesket og
Menigheden (\textit{ἄνθρωπος} og \textit{ἐκκλησία});
% the compositor letterspaces the first two pairs of Æons and not the
% third: "Mennesket og Menigheden" is set plain. As printed.
de emanerede parviis (\textit{σύζυγοι}) som Mand og Qvinde,
og fremstille i Forening Ideal-Realitetens „Fylde“
(\textit{πλήρωμα}) ligeover for det væsenløse Chaos, den
uendelige Tomhed (\textit{κένωμα}), o. s. v.“ --- En saadan
dramatisk-anskuelig Fremstilling i storstilet Naivetet gjør jo
vistnok et bedre Indtryk end de Talemaader om
„Kjærlighedens Mangel“ i „Kjærlighedens Overflødighed“, om
Kjærligheden som „Skabelsens Grund“ og Grundens
„Endemaal“ i \textit{creat sibi, creat nobis, creat omnibus}
o. s. v., hvorved en speculativ Dogmatik, der maa sikkre sig
sin Orthodoxie i smaalig Forstandighed, før den tør slaae sig
løs i gnostisk Aandrighed, gaaer fyldig tilbunds i Grundenes
„Fylde“ og aabenbarer sin hule Videnskabelighed.

Den afgjørende Modsætning imellem den gnostisk-dogmatiske
Speculation paa den ene og den begrundende Reflexion paa
den anden Side er imidlertid iøinefaldende nok. Istedenfor en
Begrundelse giver den gnostiske Speculation os
Phantasiebilleder, der forestille en Begrundelse; men en Be\-%
% --- p. 96 ---
\setcounter{footnote}{0}%
grundelse lader sig ikke forestille; den forestillede Begrundelse
lider af den Modsigelse, at det Tilværende, der skulde begrundes,
allerede forefindes som tilværende i Grunden. Urgrunden,
Grundeenheden, er ved denne udvortes Opfattelse forvandlet til en
blot og bar Indfatning for det Mangfoldige, et dybt Bassin, hvoraf
det Forskjellige udstrømmer; men Forskjellen i dette Forskjellige
bliver ikke begrundet.

2) \emph{Alt er i Grunden forskjelligt.} For at forstaae
Grundeenheden er det nødvendigt at gjøre sig Rede for, hvad der
ligger i Forskjelligheden. To Extremer ere her at mærke: enten
opgiver man ganske Eenheden og lader en Mangfoldighed af tilværende
Realpunkter afgive Forudsætning for Begrundelsen af Phænomenernes
Forskjellighed; eller man lader, for at Intet skal være
ubegrundet, det Mangfoldige og Forskjellige fremkomme ved rene
Negationer og Negationers Negationer i Væsenet selv, det almene
Grundvæsen. Det første Extrem charakteriserer Realismen, det
sidste Idealismen.

Begynde vi med Idealismen, da er det unegtelig en meget rigtig
Tanke, at det Forskjellige kun kan begrundes derved, at det
fremgaaer af en Eenhed, hvori det vel har sin Grund, men ikke sin
Tilværelse: denne Grundeenhed er Identiteten. Men hvad er nu
Forskjellen for sig som Forskjel? „En Negation af Identiteten!“ og
Identiteten for sig som Identitet? „En Negation af Forskjellen!“
Skulde nu Identiteten for sig være Grund til Forskjellen, saa
maatte Forskjellen for sig jo ogsaa være Grund til Identiteten.
Forskjellen kunde da paa denne Maade ikke blive begrundet for
Identiteten, og Identiteten ikke for Forskjellen. Forat en
Begrundelse kan blive mulig, maa Forskjellen mellem Identitet og
Forskjel selv være begrundet i begges Eenhed; den begrundende
Eenhed bliver da som Forskjelseenhed selv en Identitet; og saaledes
forstaae vi, hvad Hegel mener, naar han kalder Grunden en Identitet
af Identitet og Forskjel.
% --- p. 97 ---
\setcounter{footnote}{0}%
% signature mark at the foot of this page: „Grundideernes Logik.“ + 7
Men Manglen ved denne Begrundelse er, at det Begrundede mangler
Realitet.

At mangle Realitet er at mangle det umiddelbart Positive. De
anaxagoriske Homoeomerier, de demokritiske Atomer, de herbartske
Realer tildrage sig da i denne Sammenhæng paany Opmærksomheden.
„Hellenerne“, siger Anaxagoras, „antage med Uret Vorden og
Forgaaen, thi ingen Ting vorder eller forgaaer, men Tingene
sammensættes af bestaaende Ting, og saaledes vilde de rigtigere
kalde Vorden Sammensætning og Forgaaen Adskillen“%
% printed note mark: *)  — the range dashes in this note are the short (en)
% sort, plainly shorter than the long dash used in the body text
\footnote{Forelæsninger over „Phil. Propæd.“ fra Universitetsaaret
1860--61 S. 122--167.}. Istedenfor rene Grunde faae vi altsaa rene
Realiteter, d. v. s. tilværende Grunddele: til Grunddelenes
Realitet svarer altsaa Forskjellighedens Realitet. At nu
Grunddelene ogsaa have deres Realitet, skal ikke negtes; kun maa
det erindres, at denne Realitet ikke tør være grundløs. Der er i
Tingenes Verden en real Forskjellighed, der ikke lader sig forklare
af logiske Negationer og fortættede Reflexer. Eller hvorledes
skulde man vel tænke sig Muligheden af, at en Forskjel saa
udtrykkelig bestemt, som Forskjellen mellem Iltens, Brintens,
Svovlets Atomer o. s. v., kunde fremkomme, dersom der ikke i
Altilværelsens ideale Grundeenhed var noget Realt, noget i sig
Bestemt, der oprindelig laae til Grund for slige Forskjelligheder.
Atomikernes Feil er ingenlunde den, at de forudsætte Fleerhed og
Forskjellighed i Grunden; men Grundfeilen er, at de tilkjende de
enkelte reale Led, de enkelte Atomer, en absolut, i sig ubegrundet,
Selvstændighed; Grundfeilen er, at Physikerne ligesaa vel som
Gnostikerne lade det, der skulde begrundes i og ved den Alt
begrundende Sammenhængseenhed, have en umiddelbar, altsaa grundløs,
Tilværelse i Grunden.

Paa det oprindelige Forhold imellem Realitet og Idealitet,
Forandring og Forvandling beroer væsentlig al Begrun\-%
% --- p. 98 ---
\setcounter{footnote}{0}%
delse. Forstandskritiken lader i Almindelighed Negationen falde
udenfor Realiteten; det er en Misforstaaelse: dersom Realitet og
Negation holdt sig udenfor hinanden, kunde den ene Realitet jo ikke
selv være Negation af den anden. Men nu vil dog vel Ingen negte,
at Ilten og Brinten netop holde sig i real Forskjellighed derved,
at Ilten er en Negation, nemlig en Realnegation af Brinten, og
Brinten ligesaa af Ilten. Grundforskjellen, den grundige Forskjel,
beroer ligesaalidt paa umiddelbare Realiteter, som paa rene
Negationer: den grundige Forskjel beroer paa reale Negationer. Som
Reflexer i den begrundende Reflexion ere de reale Negationer
tillige Momenter i den begrundende Eenhed: saaledes, og kun
saaledes, bliver den grundige Forskjel selv begrundet.

At Forskjellen bliver begrundet, sees udtrykkelig deraf, at den ene
Forskjellighed gaaer tilgrunde, for at den anden kan fremgaae af
Grunden. Saaledes gaaer den reale Forskjel mellem Ilten og Brinten
tilgrunde i den chemiske Forening, der begrunder Dannelsen af
Vandet; men de Egenskaber, de reale Forskjelligheder, der
udelukkende beroe paa Vanddannelsen, forsvinde igjen og gaae
tilgrunde, naar Vandet decomponeres, og den reale Forskjel mellem
Ilten og Brinten med alle deraf flydende Forskjelligheder derved
igjen begrundes, o. s. fr. Forskjelseenheden af Realitet og
Negation er kun et bestemtere Udtryk for den reflecterede
Umiddelbarhed; Modsætningen af Negation og Realitet er i
Begrundelsen forvandlet til en Modsætningseenhed af Idealitet og
Realitet.

Modsætningseenheden af Idealitet og Realitet begrunder igjen en
Modsætningseenhed af Forvandling og Forandring og derved tillige en
Modsætningseenhed af den grundende Subjectivitet og Begrundelsens
objective Indhold. At den eensidige Idealisme, hvis Negationer og
Forskjelsbestemmelser altid mangle Realitet, miskjender Andetheden
og den til Andetheden svarende Forandring, er fra mange Sider
indlysende. Den forflygtiger overalt Forandringens Realitet og
holder sig
% --- p. 99 ---
\setcounter{footnote}{0}%
i det Uendelige til Forvandlingen; derfor kan den trods sin
Grunden og Begrunden aldrig faae Noget i Virkeligheden begrundet;
den kommer, med al sin Opmærksomhed for det Objective i det Almene
og det Almene i det Objective, aldrig ud over den grundende
Subjectivitets almene Reflexionsbestemmelser.

Med den eensidige Realisme forholder det sig omvendt; fra dette
Synspunkt er „Forvandlingen“ kun en mystisk Talemaade: al Grund og
Begrundelse beroer her paa lutter „Forandringer“. Ved udelukkende
at fastholde Forandringen forvandler Realismen imidlertid, uden
selv at vide deraf, alle reale Følger til formelle Conseqvenser,
alle reale Grunde til umiddelbare Forudsætninger og endelige
Betingelser. Saaledes kommer da Alt igjen til at beroe paa
Udstykning og Sammenstykning. For at sikkre sig mod den mystiske
Forvandling har den over sine egne Reflexioner uklare Forstand i
denne Henseende tilladt sig den utilladeligste af alle
Forvandlinger: den Forvandling, hvorved det Objective i død
mechanisk Udvorteshed falder udenfor det Subjective; den
Forvandling, hvorved Idealitetens d. v. s. Fornuftens Lys
udslukkes; den Forvandling, hvorved al Grund og Begrundelse
væsentlig afskaffes%
% printed note mark: *)
\footnote{Jvfr. Forelæsninger over „Phil. Propæd.“ 1860--61, S.
160--167, hvor Exempler \textit{in concreto} findes.}.

3) \emph{Alt maa have sin tilstrækkelige Grund.} Sætningen om den
tilstrækkelige Grund (\textit{principium rationis sufficientis})
er, som bekjendt, bleven opstillet og forsvaret af Leibnitz; men dens
logiske Udvikling er i flere Henseender hæmmet og forqvaklet. For
det Første er det en Forqvakling af Grundforholdet, naar Sætningen
om den tilstrækkelige Grund stilles paa udvortes Maade overfor
Identitetssætningen som et Princip for Erfaringssandheder overfor
et Princip for Fornuftsandheder. Herved dannes en Kløft imellem
det
% --- p. 100 ---
\setcounter{footnote}{0}%
% signature mark at the foot of p. 99: 7*
Aprioriske og det Empiriske, det Abstracte og det Concrete, der
gjør al sand Begrundelse umulig. Saalænge den formelle Identitet:
$A$ er $A$, skal gjælde for mere end en reen Begyndelsessætning, en
Abstraction, der er bestemt til at integreres ved et: $A$ er ikke
$A$, men $B$, kan her slet ikke være Tale om at indsee, hvad Grund
og Begrundelse er. Den abstracte Sætning: $A$ er, hvad $A$ er, er
et ligesaa grundløst Udtryk for en formel Overeensstemmelse, som
Sætningen: $A$ er, hvad $A$ ikke er, er et grundløst Udtryk for en
formel Modsigelse. Først ved Hjælp af Contradictionssætningen: $A$
er ikke, hvad $A$ ikke er, viser sig en Skygge af Begrundelse. Det
er let at see, hvorfra Skyggen kommer; den kommer fra den i
Contradictionssætningen modsagte Modsigelse. Det er i høi Grad
charakteristisk, at det saakaldte Contradictionsprincip, --- der
naturligviis ikke er et Princip, men en Abstractionssætning,
--- maa tage Modsigelsen to Gange, altsaa bestaae i en fordoblet
Modsigelse, for at Identiteten, den rene, skjære Identitet, kan
falde udenfor al Modsigelse. At Sætninger som: $A$ er ikke, hvad
$A$ er; $A$ er, hvad $A$ ikke er, ere, for sig betragtede, baade
modsigende og grundløse, forstaaes vel af sig selv. Men hvorledes
ophæve vi nu Modsigelsen? Paa overfladisk Maade ved
Contradictionssætningen; det Overfladiske sees deri, at Tanken
ligefrem vil afvise Modsigelsen, uden at betænke, at det netop er
Modsigelsen, ene og alene Modsigelsen, den har at takke for, at den
dog --- ved nemlig at modsige Modsigelsen --- er kommen i et Slags
formel Bevægelse. Paa en grundig Maade kan Modsigelsen
naturligviis kun hæves ved Grunden selv og den tilsvarende
Begrundelse. $A$ er ikke --- i Grunden, hvad $A$ er --- i
Existensen. Dersom en Ting var i sin Grund, hvad den er i sin
Existens, saa var der hverken Grund eller Existens; og dog er
Tingens Grund identisk med dens Existens, thi Grunden gaaer med i
Begrundelsen, og Existensen er begrundet.

% --- p. 101 ---
\setcounter{footnote}{0}%
At den formelle Logik ved sine saakaldte Principer ikke tænker sig
Andet end Grundsætninger, er tydeligt nok; men charakteristisk er
det i høi Grad, at netop disse Grundsætninger overalt stille sig i
et tankeløst Forhold til Grund og Begrundelse. De passe paa det
Abstracte i dets rene, formelle Abstraction, thi det Abstracte har
endnu ikke erholdt sin Begrundelse; de passe paa det Concrete i
dets umiddelbare Concretion, thi det umiddelbart Concrete er et
Existerende, og det Existerende har allerede erholdt sin
Begrundelse: Begrundelsen falder altsaa i begge Tilfælde udenfor
Reflexionen. Skal en Begrundelse finde Sted, maae de formelle
Grundsætninger haves. Saaledes gjælder \textit{principium exclusi
medii} unegtelig for det Existerende; men for det Existerendes
Grund gjælder \textit{principium inclusi medii}. Imod Existensens
disjunctive Dom: $A$ er enten $B$ eller $C$, stiller Grunden sin
conjunctive Dom: $A$ er --- i Grunden --- baade $B$ og $C$.
Mennesket er i Grunden baade Qvinde og Mand, thi Qvinden er ligesaa
væsentligt Menneske som Manden; men i Existensen er Mennesket enten
Mand eller Qvinde, thi Androgynens Existens er mystisk.

Det Spring fra de rene Fornuftsandheder til Erfaringssandhederne,
hvorved Grund og Begrundelse springes over, falder umiddelbart
sammen med et Spring fra Grund til Aarsag, hvorved de begrundende
Mellembestemmelser ligeledes springes over. Vistnok ere, hvad vi
siden skulle see, alle Causalitetsforhold i sig tillige
Grundforhold; men derfor ere alle Grundforhold dog ikke uden videre
at opfatte som Causalitetsforhold. Det er imod Malebranche's
eensidige Fornegtelse af de naturlige Mellemaarsager, at Leibnitz
har gjort Sætningen om den tilstrækkelige Grund gjældende. Dersom
nu den i Sætningen skjulte Grundtanke var bleven fuldstændig
udviklet fra Grundens og Begrundelsens Synspunkt, da maatte
Udviklingen unegtelig have været frugtbar i to Henseender. For det
Første vilde Læren om den tilstrækkelige
% --- p. 102 ---
\setcounter{footnote}{0}%
Grund have viist os, hvad det vil sige, at Tingene selv ere i
væsentlige, objective Relationer til hinanden; thi de væsentlige
Relationer bestaae jo netop deri, at Tingene begrundes ved hinanden
indbyrdes, at en Ting kun opstaaer derved, at andre Ting gaae
tilgrunde. Begrundelsen maatte da fra denne Side have viist sig som
en objectiv, Phænomenerne iboende, Bevægelse. Men til en saadan
Indsigt kan en Monadelære, der vil ophæve Tingenes Vexelvirkning og
sætte sin kunstig udtænkte \textit{harmonia præstabilita} istedet,
umulig bringe det. For det Andet vilde Læren om den tilstrækkelige
Grund, dersom den var bleven tilstrækkelig udviklet, have aabnet os
et dybt Indblik i Vexelbestemmelsen af Grundenes Mangfoldighed med
Grundeenheden, af Tingenes særskilt betingede, afledede Grunde med
Altilværelsens oprindelige Grund. Thi det forholder sig ikke
saaledes, som Hegel i flygtige Reflexioner har yttret, at enhver
Grund ligefrem maa gjælde for en tilstrækkelig Grund%
% printed note mark: *)
\footnote{„Daß der Grund \emph{zureichend} sey, ist eigentlich sehr
überflüssig hinzusetzen“. Werke. 3ter Bd., Berlin 1833. S. 76.
Jvfr. S. 103.}. Man kan nemlig forqvakle en Begrundelse paa høist
forskjellig Maade: enten ved at holde sig til det abstract
Almindelige eller ved at indskrænke sig til det altfor Particulære.
Hvad er f. Ex. Grunden til, at $A$ er bleven drikfældig? Svarer
man: Grunden er, at vi alle ere Syndere, saa holder man sig for nær
til det abstract Almindelige; svarer man: Grunden er det
Misforhold, $A$ hver Dag sætter imellem det Qvantum, han nyder, og
det Qvantum, han taaler, saa vilde man holde sig for eensidigt til
det abstract Particulære. Fra det particulære Synspunkt maa hver
eiendommelig Ting nødvendigviis have sin eiendommelige Grund; thi
det existerende Phænomen kan umuligt indeholde Andet, end hvad der
ligger i dets Grund.
% --- p. 103 ---
\setcounter{footnote}{0}%
At den kulsure Kalk krystalliserer den ene Gang som Kalkspath, den
anden Gang som Arogonit, har sin særegne Grund i de med
% „Arogonit“ so printed, for Aragonit; the second letter is a clearly
% bowled o, not the a of „Kalkspath“ on the same line
Temperaturen vexlende Betingelser. Denne særegne Grund er
naturligviis en specificeret Udtalelse af den almene Grund, hvortil
Stoffernes Krystallisering i det Hele maa henføres. Men
Krystalliseringens almindelige Grund med det System af Relationer
og Betingelser, hvorunder dens Phænomener fremtræde, er igjen at
bestemme som en særskilt Grund i Forhold til andre, ligeledes hver
i sin Kreds almindelige og ved Kredsenes Modsætning særskilte
Grunde, f. Ex. Grunden til Tilstandsformernes Forandring ved
Varmen, Grunden til Dannelsen og Opløsningen af chemiske
Forbindelser ved Affiniteten%
% printed note mark: *)
\footnote{Saadanne Factorer som Affiniteten og Varmen ere nemlig,
for sig betragtede, ikke umiddelbart Grunde men Aarsager; heller
ikke kan Grunden siges at ligge i Affiniteten for sig eller i
Varmen for sig; thi Grunden bestemmes egentlig ved Forholdet, det
existentielle Væsensforhold mellem Varmen og Tilstandsformerne,
mellem Affiniteten og de for Forbindelserne gjældende Betingelser,
o. s. v.} o. s. fr. Først efterhaanden som den begrundende
Reflexion seer sig istand til at articulere de særlige Grundes
System i hver særlig Kreds, kan den tilstrækkelige Grund komme
tilsyne i Kredsens begrundede Sammenhæng.

I Begrundelsen af det enkelte Phænomen maa Begrundelsen af dets
Stilling i Phænomenernes hele System reflecteres, og omvendt. Dette
har Leibnitz forsaavidt indseet, som han har indseet, at hvert
enkelt Naturphænomen er betinget, at det ene Betingede viser hen
til det andet i det Uendelige, at det Betingedes endeløse Række
altsaa kun kan bestemmes ved en Alt bestemmende Grundenes Grund,
hvori Betingelserne gaae op som Momenter; saaledes kom han da til
at sætte et Begrundelsernes Endemaal, en Grundenes Grund, en
\textit{ratio finalis}. Men var der først gjort et Spring fra
Grunde til Aarsager, saa er det ikke forunderligt, at der blev
gjort lignende Spring fra Causalitet
% --- p. 104 ---
\setcounter{footnote}{0}%
til Teleologie og ved Hjælp af Teleologien fra de naturlige til de
overnaturlige Aarsager, fra Physik til Theologie.

See vi imidlertid bort fra disse Spring og holde os til, hvad der
ligger i Begrundelsens Problem, da overbevises vi om, at Læren om
„den tilstrækkelige Grund“ falder sammen med Læren om den fuldendte
Begrundelse. I den fuldendte Begrundelse faaer det Objective først
tilfulde sin Realitet formedelst det Subjective. Kunde Tingenes
Begrundelse i og ved deres indbyrdes Sammenhæng afsluttes for sig,
da vilde den begrundende Reflexion føre Subjectiviteten, ikke blot
til en Fordybelse, men til en fuldstændig Selvfortabelse i Tingenes
Grunde; der maatte da gives en i Tingene faktisk Fornuft, en
Fornuft saa objectiv, at den vel var istand til at begrunde alle
Ting, men ikke istand til at udgrunde sit eget Væsen. Det er denne
Grundmodsigelse i Fornuften selv, der kommer tilsyne i Tingenes
Verden derved, at det Betingede skal begrundes af det Betingede i
endeløse Rækker. Ifølge den tilstrækkelige Grund gaaer den
objective Begrunden derimod i Eet med en subjectiv Bestemmen; thi
den subjective Grundbestemmen er en Motivering. Systemet af
Tingenes objective Grunde er fra Subjectivitetens, den ontologiske
Subjectivitets, oprindelige Synspunkt et System af Bevæggrunde.
Existensgrunden og Bevæggrunden ere i Begrundelsen Eet og dog
grundforskjellige. Existensgrunden, den objective, betingede Grund
svarer til Tingenes endelige Realitet; Motivet, den subjective
Grund, svarer til Tingenes Betydning for Relationernes System. Kun
i Eenhed med den reale Grund kan Motivet blive tilstrækkeligt;
adskilt fra den objective Begrunden vilde Bevæggrunden for sig
derimod være abstract-subjectiv og vilkaarlig.
% „abstract-“ at the printed line end carries a SINGLE hyphen, not the
% double-stroke break hyphen used elsewhere on this page: a real compound
Kun i Eenhed
med Motivet kan den reale Grund blive tilstrækkelig; thi adskilt
fra Motivet vilde Realgrunden for sig netop i sine umiddelbare,
faktisk givne Forudsætninger og Betingelser have noget Grundløst.
At Motivet er Grund, beviser, at Subjectivitetens Selvfordybelse
% --- p. 105 ---
\setcounter{footnote}{0}%
er en Fordybelse i Tingenes objective Væsen; at Grunden er Motiv,
beviser, at Begrundelsen af Tingenes Væsen maa grunde sig paa
Subjectivitetens uendelige Selvfordybelse.

\subsubsection*{c) Virkelighedsreflexion.}
% printed label is a Fraktur „c)“, read off the image. It follows the
% „d) Forstandsreflexion“ of p. 60 and so looks out of sequence; the head
% is transcribed exactly as printed, not renumbered.

Forsaavidt Tilværelsens Phænomener ere tilstrækkelig begrundede, er
Tilværelsen virkelig; den begrundende Reflexion, der bestemmer
Phænomenerne, bestemmer altsaa Virkeligheden. Da nu en Virkelighed
ikke kan være uden Virksomhed, og en Virksomhed ikke uden Virken,
saa opstaaer her det Spørgsmaal: hvorledes bestemme vi Forholdet
mellem Reflecteren og Virken? „Al Virken er i Grunden en energisk
Reflecteren og følgelig Eet med en objectiv Tænken. Den
gjennemgaaende Reflexion vil altid fordre, at de i hinanden
reflecterede Led ikke alene forvandles til uselvstændige Momenter,
til rene Reflexer, men under Forvandlingen tillige vise sig i deres
Forskjellighed og bevare en vis relativ Selvstændighed. Jo større
nu den Selvstændighed er, hvormed Leddene optræde, desto stærkere
maa ogsaa den Reflexion være, der skal forvandle dem til Momenter
og magte dem som Reflexer; vi have her en Maalestok for
Reflexionens Energie. I den blot intellectuelle Reflexion ere de
reflecterede Led som Billeder, Forestillinger og Begreber kun
flygtige Skin, uselvstændige Reflexer; heraf forklares saa den
rene, d. v. s. abstracte Tænknings Afmagt til at frembringe
nogetsomhelst Existerende. I den naturkraftige Reflexion have de
reflecterede Led en fremtrædende Selvstændighed og blive just ved
denne deres Selvstændighed slaaende Udtryk for den Magt, der
betegner deres Afmagt og holder dem i Uselvstændighed. Reflexionen
er i sin energiske og substantielle Væsenhed Magt. Ved Magt sættes
de i Grunden reflecterede Momenter som virkelige Existenser, som
særskilte Ting: herfra stammer den materielle Discretion, herfra
Atomernes Vrimmel, herfra Aggregaternes og Massernes Tilværen i
Rummet. Ved Magt sættes de virkelige Existenser,
% --- p. 106 ---
\setcounter{footnote}{0}%
de særskilte Ting som Momenter og Reflexer i den reflecterede
Vexelvirkning: herfra al Forening og Opløsning, herfra al Opstaaen
og Forgaaen, herfra Forandringernes Mangfoldighed med alle
Overgange ligefra den flygtige Ændring til den gjennemgaaende
Forvandling! Den Magt, der sætter, og den Magt, der hæver Tingenes
Selvstændighed, er vel i Existensens Form forskjellig, men dog i
Grunden en med sig selv identisk Magt“%
% printed note mark: *)  — the note's opening quote is set tight against
% „over“ with a space after it: over„ Phil. Propæd.“
\footnote{Forelæsninger over „Phil. Propæd.“ fra Universitetsaaret
1860--61. S. 167--177.}, d. e. en uendelig Selvmagt.

I den gjennemgaaende Eenhed af Reflecteren og Virken have vi da en
Forskjelseenhed af Viden og Magt. Denne Eenhed af Viden og Magt
lader sig, hvad Reflexionen selv angaaer, udvikle fra
Synspunkterne af en \emph{substantiel}, en \emph{causal} og en
\emph{totaliserende} Reflexion.
% the letterspacing covers only the three adjectives: „en“, „og“ and
% „Reflexion“ are set normally. Not regularised.

\subsubsection*{α) Den substantielle Reflexion.}
% printed label is a Greek alpha, not a Fraktur „a“; in the printed head
% only „Den substantielle“ is letterspaced, „Reflexion“ is not

Den substantielle Reflexion er i Almindelighed en
Selvstændighedsreflexion, en Reflexion af det i sig Bestaaende.
Ligesom den grundende Reflexion søgte at bestemme Grunden, saaledes
søger nu den substantielle at bestemme Substansen. For at bestemme
Substansen, betragte vi igjen det Existerende, det tilstrækkeligt
begrundede Phænomen, dog fra en ny Side. Det i Existensen
Selvstændige, det umiddelbart Bestaaende, er da Tingen selv, og
Beskaffenhederne ved dette for sig Bestaaende Tingens Egenskaber.
Er nu Tingen som Ting i logisk Forstand at kalde Substans? Nei,
dertil er dens Selvstændighed endnu for umiddelbar og i sin
Umiddelbarhed igjen for flygtig. Skulde Tingen i Sandhed gjælde
for Substans, da maatte den have en væsentlig Bestaaen i sig selv;
men Tingen bestaaer kun i sig ved at bestaae i sine Egen\-%
% --- p. 107 ---
\setcounter{footnote}{0}%
skaber. Tingen i sin Eenhed for sig er Et og Egenskaberne i deres
Mangfoldighed et Andet; ved at bestaae i sine Egenskaber bestaaer
Tingen altsaa i sit Andet, og er ligesom Egenskaberne underkastet
Forandringer.

At Egenskaberne forandre sig, er begrundet deri, at den ene
Egenskab altid maa beroe paa den anden, Egenskaben ved den ene Ting
paa Egenskaben ved den anden. Saaledes kunde Tingen f. Ex. ikke
have Egenskaben: rød, dersom den ikke tillige havde Egenskaben:
Udstrækning, heller ikke kunde den have Egenskaben: rød, dersom
Lyset ikke havde en tilsvarende Egenskab. For en overfladisk
Betragtning tager det sig jo vistnok ud, som om Egenskaberne meget
vel kunde forandre sig, og Tingen endda blive den samme: Bladet,
som før var grønt, er nu guult; men det gule Blad er dog ogsaa et
Blad%
% printed note mark: *)
\footnote{„Hvor længe kan man vedblive at kalde en Gjenstand „En“
eller „Et“? Det vil sige: Hvor længe kan man betragte den som den
samme Gjenstand, tiltrods for de idelige Forandringer, som den er
underkastet? Blive Forandringerne ikke tilsidst saamange og saa
store, at Gjenstanden derved bliver en anden „En“, eller et andet
„Et“? --- I „en Skov“ kan man fælde mange Træer, og den bliver dog
endnu den samme Skov; „et Huus“ kan deels udvides, deels forfalde,
og det bliver dog det samme Huus. Men at der for alle disse
Forandringer gives en Grændse, udenfor hvilken Gjenstanden bliver
til en anden, er indlysende, og er ogsaa udtrykt i den bekjendte
populære Ironie om Kniven, der efter nogle Gange at have faaet et
nyt Blad og andre Gange et nyt Skaft, endnu bestandig gjaldt for
den samme Kniv.“ Heiberg, Det logiske System. Første Afhandl.
Perseus. Nr. 2. S. 38. (Prosaiske Skrifter, 2det Bind, 1861, S.
158.)}. Denne Adskillelse mellem Egenskabernes og Tingenes
Foranderlighed gjælder imidlertid kun deelviis i det Endelige.
Vistnok er her med Hensyn paa Tingen selv en Forskjel mellem
væsentlige, mindre væsentlige og uvæsentlige Egenskaber; men i
logisk Forstand gjælder det, at en Forandring i en hvilkensomhelst
af Egenskaberne \textit{eo ipso} maa være en Forandring i Tingen
selv, efterdi Tingen i sin existentielle Umiddelbarhed
% --- p. 108 ---
\setcounter{footnote}{0}%
ligesaa umiddelbart er Eet med som forskjellig fra sine
Egenskaber.

Som det Foranderlige bestemmes, saaledes maa dets Modsætning, det
Uforanderlige, ogsaa bestemmes. Opfattes Noget i abstract
qvalitativ Modsætning til Andet, da er den uophørlige Overgang fra
Noget til Andet et kategorisk Udtryk for Forandringen selv; men det
Uforanderlige falder da ogsaa, ved denne Overgang, qvalitativt
udenfor det Foranderlige. I Væsensreflexionen er Overgangen fra
Noget i Andet en Overgang af Skin i Skin, en Overgang, som, for at
bruge et Udtryk af Hegel, ingen Overgang er; det Foranderlige, der
paa denne Maade er forvandlet til lutter Skin, har med det Samme
erholdt en saa fuldkommen Gjennemsigtighed, at den i alle
Skinforandringer sig reflecterende Uforanderlighed, den rene
Selvhed, paa ethvert Punkt kan skinne igjennem.

Dersom nu Skinnet i sin Gjennemsigtighed ikke tillige reflecterede
sit umiddelbart Modsatte, det Umiddelbare selv, vilde Selvheden
være ligesaa væsenløs som Andetheden, og Selvreflexionens rene
Klarhed tabe sig i Tomhed. Men med hvilket Eftertryk
Umiddelbarheden paatrænger sig Reflexionen, have
Vexelbestemmelserne i den phænomenale, den grundende og den
begrundende Reflexion tilstrækkelig viist. De flygtige Skin, de
tomme Reflexer fortrængtes saaledes efterhaanden af reale
Phænomener, af existerende Ting; men alt som hvert Led i
Endelighedens Række erholdt forhøiet Selvstændighed, blev ogsaa
Modsætningen mellem det Selvstændige og det Uselvstændige desto
mere paatrængende.

I Existensen, den begrundede Tilværelse, er Forskjellen mellem
Selvstændigt og Uselvstændigt endnu kun relativ; Et er
selvstændigt, et Andet uselvstændigt; det Ene er mere, det Andet
mindre selvstændigt: Overgangen sees i Forandringerne. Men naar
denne Relativitet nu skal gjennemreflecteres, stode
% „stode“ so printed, without the slash, for støde: the o is a plain bowl,
% while „forhøiet“ six lines above carries a clear slash on the same page
vi paa en Modsigelse, der er Virkelighedens allerhaardeste,
% --- p. 109 ---
\setcounter{footnote}{0}%
den nemlig, at det Selvstændige selv er et Uselvstændigt, det
Uselvstændige et Selvstændigt.

Forandren er Virken og al Virken en Reflecteren. I den virksomme
Eenhed af Reflecteren og Virken er Omsætningen af Selvstændigt i
Uselvstændigt som en uendelig Vibreren; hvert Led, hver Deel, hvert
Atom vibrerer i Virksomheden selv mellem Selvstændighed og
Uselvstændighed. Det er denne Vibreren, denne Zittren og Dirren i
punktuel Reflecteren, der i nærværende Sammenhæng særlig
charakteriserer Virkeligheden og det Virksomme. Det er ikke
længere den rene Skinnen, som i Reflexionen af Skin og Væsen, vi
her have for os; heller ikke kan Forskjellen mellem det i Grunden
Tunkle og Dæmrende paa den ene Side og det i Phænomenerne
% „Tunkle“ so printed: the initial is the Fraktur T sort, plainly distinct
% from the D of „Dæmrende“ on the same line. For Dunkle.
Halvklare, Klare og Anskuelige paa den anden længere afgive
Charakteermærker; thi Grundbestemmelser og Phænomenalbestemmelser
ere Sider af det samme Virkelige: Virksomheden selv er paa een
Gang grundbestemmende og existensbestemmende. Det i de virkelige
Forandringer virksomme Væsen fremstiller sig da heller ikke for
Anskuelsen som en i Dybet gjærende Grund; den i Virksomheden selv
reflecterede Eenhed af Vorden og Virken er snarere at sammenligne
med Ilden, den „evig levende Ild“ (πῦρ ἀεὶ ζῶον), om hvilken
Heraklit siger, at den udslukker Maal (ἀποσβεννύμενον μέτρα) og
antænder Maal (ἁπτόμενον μέτρα).

I alle virkelige Forandringer er Virksomheden væsentlig Eet med det
Virksomme selv, og det Virksomme selv det Uforanderlige. Det
Virksomme selv er det i sig Ene; thi ellers kunde det umulig virke
i alle virkelige Forandringer og i dem alle hævde sin
Uforanderlighed. Det Virksomme selv har en absolut Bestaaen; thi
det virker i Alt og paa Alt, og dets Selvvirken er dets Bestaaen.
Men det Virksomme selv har ikke en absolut Bestaaen i endelig
Modsætning til Mangfoldigheden af relativt bestaaende Ting; det
relativt Bestaaende er kun et Virkelighedsmoment, et
Gjennemgangspunkt for den
% --- p. 110 ---
\setcounter{footnote}{0}%
uendelige Virken. Det Virksomme selv, der i Alt udvirker sin egen
Selvstændighed, er en bestaaende Selvhed; thi Bestaaen er Selvhed, og
virkelig Selvhed, Selvvirksomhed, Eet med reflecteret Selvstændighed.

Er nu det i Virkeligheden Uforanderlige, det i sig Bestaaende, Substans,
saa er den absolute Selvbestaaen absolut Substantialitet, og den absolute
Substantialitet paa dette Udviklingstrin identisk med den absolut
virkelige, i Alt virksomme Subjectivitet. Det er ikke, hvad Hegel
antager, Substantialiteten, der er det Oprindelige, og Subjectiviteten,
der er det Afledede; men det er det Omvendte. Det er ikke
Substantialiteten, der er bleven subjectiv; det er tvertimod
Subjectiviteten, der under Begrebsudviklingen er bleven, hvad den evig
er, substantiel. At \emph{aflede} Subjectiviteten \emph{af}
Substantialiteten er ligesaa umuligt som at aflede Viden af Ikke-Viden,
% Letterspacing on p. 110, checked at 600 dpi: only „aflede“ and „af“ are
% spaced in the first clause, and only „udvikle“ and „til“ in the answering
% clause below — „Subjectiviteten“, „Substantialiteten“ and the SECOND
% „aflede ... af“ (this line) are set normally. The emphasis falls on the
% verb and its preposition, which is the point of the antithesis; it is not
% regularised over the whole phrase.
% „Ikke=Viden“ is printed with the Fraktur double-stroke hyphen AT A LINE
% END, so the print cannot distinguish a lexical hyphen from a break
% hyphen. Kept as the compound „Ikke-Viden“ (cf. „Non-$A$“, p. 16); not
% joined to „IkkeViden“, which is not a word.
Selvheden af det Selvløse, det Subjective af det Objective. Derimod er
det i god Overeensstemmelse med Sagens Natur at \emph{udvikle}
Subjectiviteten \emph{til} Substantialitet, thi det er just at udvikle af
Subjectiviteten selv, hvad der evigt ligger i dens Væsen.

Det Selvstændiges reflecterede Eenhed med sit Uselvstændige er Magt;
Momenternes Gjennemsigtighed for det reflecterende Selv er Viden: i
Magtens og Videns reflecterede Eenhed har den substantielle Subjectivitet
sin absolute Bestaaen. Al Bestaaen og Forgaaen beroer paa Reflecteren af
Selvstændigt i Uselvstændigt. Hvad der i Virkeligheden bestaaer, maa
bestaae ved en Magt; hvorfor man jo ogsaa siger om et saadant Bestaaende,
at det staaer ved Magt, og mener dermed, at det gjælder for et Virkeligt,
at det virkelig gjælder. Hvad der i Virkeligheden forgaaer, forgaaer kun
derved, at det viger for en Magt, det synker i Afmagt, det bukker under i
Striden og forsvinder. Paa Virkelighedens Trin træder Magten i Grundens
Sted; istedenfor Grunde og Begrundelser møde vi her overalt
Magthandlinger, d. v. s. Magtsprog,
% --- p. 111 ---
\setcounter{footnote}{0}%
hvori en bydende Nødvendighed udtaler sig. Dette maa dog ikke
misforstaaes, som om Grunden var afskaffet, og Magten grundløs; thi
Magten selv er just den i alle Momenter af Forudsætninger og Betingelser
fuldendte Grund, en Grund saa virkelig, at Begrundelsen og det
Begrundede maae falde sammen i Magthandlingen selv. Magten er i uendelig
Reflecteren; og dog seer det ud, som om den hvilede i evig Væren; Magten
er i sin reflecterede Væren ophøiet over al Vorden: dens Handlinger ere
dens Aabenbarelser. Hvad dette vil sige, bliver klart, naar vi undersøge,
hvorledes Magten selv stiller sig til det Mægtige og til det Afmægtige.
% A large round ink blot stands in the measure after „Afmægtige.“ at the
% end of this line (600 dpi: a solid blob, not a typographic sort). Not
% transcribed.

Som Identiteten er, maa ogsaa Forskjellen være; er Identiteten af Magten
og det Mægtige substantiel, saa maa her være en tilsvarende Forskjel.
Dette har Spinoza overseet. Hos ham er den uendelige Substans en absolut
Eenhed af Magten og det Mægtige: det Mægtige selv, Substansen
\textit{(id quod in se est et per se concipitur)} har ikke Magten; thi
Substansen \textit{(cujus essentia involvit existentiam)} er intet Andet
end Magten selv, og just derfor uendelig mægtig. Er det Mægtige, Magten
selv, det Eneværende, saa følger naturligviis heraf, at de tilværende
Ting i deres Endelighed, Forskjellighed og Mangfoldighed maae i sig være
absolut afmægtige. Den umiddelbare Modsætning af Magten og det Afmægtige
beroer netop paa den umiddelbare Eenhed af Magten og det Mægtige. Men
den umiddelbare Modsætning af Magten og det Afmægtige er en substantiel
Modsigelse. Hvorledes kunne Tingene i deres Mangfoldighed have endog blot
en Skygge af Tilværelse, naar de i deres Afmagt staae i ligefrem
Modsætning til Magten? Fra Spinozas Standpunkt maa hertil svares: at
Tingenes Afmagt ikke falder udenfor, men tvertimod sammen med
Substansens Enemagt; at Tingene netop af den Grund maa bestemmes som
Maader og Modificationer \textit{(modi, modificationes)} af Substansen
selv. Tingen, den enkelte Modus, er altsaa paa een Gang baade
% --- p. 112 ---
\setcounter{footnote}{0}%
Eet med og forskjellig fra Substansen: Eet med Substansen, forsaavidt
dens Magt, d. v. s. Magten i den, er Substansens egen Magt, Formen
Substansens Form, dens Væren Substansens Væren; forskjellig fra
Substansen, fordi den er endelig, Substansen uendelig, fordi den er
afmægtig, Substansen enemægtig.

Den substantielle Modsigelse kan nu nærmere bestemmes saaledes: Det Ene
og det Mangfoldige, det Uendelige og det Endelige ere to Sider af samme
Substans; det Uendelige er det Enemægtige, det Endelige det Afmægtige:
det Afmægtige er altsaa en Sidebestemmelse ved det Mægtige, og det i
Magten selv Mægtige en Sidebestemmelse ved det Afmægtige. Og dog ligger
det i den substantielle Reflexion, at de to Sider tillige maae vise sig
som to Selvstændige. „Det substantielle Forhold bestaaer i Magten;
Substansens absolute Selvstændighed beroer paa den uendelige Selvmagt,
hvormed den sætter og ophæver den accidentelle Totalitet. Men den rene
Selvmagt er en tom Abstraction. Magten har ikke sig selv, og kan
desaarsag ikke hypostaseres; den er Selvstændighedens Attribut, der
forudsætter det Selvstændige. Substansen, der er \textit{sui causa}, er
tillige \textit{sui compos}, og idet den selv er et bestemt Noget, maa
den tillige magte noget Andet. Den substantielle Enemagt beroer altsaa
paa den Magt, der kan indrømmes den accidentelle Mangfoldighed. Det er
netop Virkelighedens og Magtens haarde Modsigelse, at den meddeler
ethvert af Forholdets Led en adamantisk Selvstændighed, idet den dog med
det Samme gjør det ene til en Reflex af det andet. Lad $A$ t. Ex. have en
fuldkommen Magt over $B$, og Modsigelsen vil være iøinefaldende. Thi hvis
$A$ har den hele absolute Magt, saa er $B$ jo kun en Skygge; men nu
bestaaer $A$'s uendelige Selvmagt netop i den Magt, hvormed den magter
sit $B$; er $B$ kun en tom Skygge, saa har $A$ kun Magt over en Skygge,
og dens hele Selvmagt er følgelig en Skyggemagt. Er $B$ derimod mere end
en Skygge, saa har
% --- p. 113 ---
\setcounter{footnote}{0}%
% p. 113 is the first page of gathering 8: the signature mark
% „Grundideernes Logik.“ + „8“ stands at the foot. Not transcribed.
det tillige en virkelig Selvstændighed, altsaa en Magt i sig selv; $A$,
der nu magter $B$, kan da vistnok siges at have en virkelig Magt,
eftersom det magter noget Virkeligt. Men da dette Virkelige selv har en
Magt, og forsaavidt ikke magtes af $A$, saa magter $A$ altsaa sit $B$,
forsaavidt det i Virkeligheden ikke magter det“.\footnote{Den propæd.
Logik af R.\ Nielsen, Kbhvn. 1845 S.\ 166--67.}
% printed mark: *). The quotation opened on p. 112 („Det substantielle
% Forhold …“) closes here; the closing „ is followed by a full stop and
% then the note mark, as printed.

Denne Modsigelse opløser sig i en substantiel Modsætning mellem Magten i
Reflexionens og Magten i Umiddelbarhedens Form. I Reflexionens Form er
Magten ideel, i substantiel Umiddelbarhed er den reel. Det umiddelbart
Reelle har altid sit Andet, sit endeligt Modsatte udenfor sig; det
Ideelle har sit Modsatte i sig. Magten kan i reel Umiddelbarhed aldrig
være absolut; de reelle Magter ere alle relative. Den ideelle Magt kan
derimod være absolut, thi det i Magten Absolute er netop dens Idealitet:
den absolute Magt er oprindelig Selvmagt. De reelle Magter ere
Substanser; den ideelle Magt er ikke Substans, men Substantialitet:
Substantialiteten selv er Subjectivitet.

Eenheden af Selvmagten og de relative Magter, af Substantialiteten selv
og de tilværende Substanser er en reflecteret Eenhed af Viden og Magt.
Viden er Magtens Idealitet; Magten er Videns Realitet: Viden er Magt.
Dersom den absolute Substantialitet ikke var, saa var der ingen
Substanser; dersom Viden ikke var substantiel, saa var der ingen
Substantialitet. Kun Selvmagten, den subjective Magteenhed, den vidende
Magt, den mægtige Viden, formaaer at bære Substanserne, de relative
Magter, den objective Magternes Mangfoldighed, hvori det virkelige Moment
er et selvstændigt Led og det selvstændige Led et Moment.

% --- p. 114 ---
\setcounter{footnote}{0}%
\subsubsection*{β) Den causale Reflexion.}
% Display head as printed: Greek β in italic antiqua, then letterspaced
% Fraktur (structural letterspacing carried by position, so not marked
% \emph{}). The printed letter is β, read off the image at 600 dpi; the
% Greek series α) β) that ran at pp. 62 and 66 under „d) Forstands\-%
% reflexion“ is therefore running again here — β) at p. 114 and γ) at
% p. 123 — so an α) of this second series must stand somewhere in
% pp. 96--113. Transcribed as printed, not renumbered.

Aristoteles har i en af sine Aporier\footnote{Metaphys. III.\ 5.} paa
eiendommelig Maade behandlet det Spørgsmaal, om Tallene, Legemerne,
Fladerne og Punkterne ere Væsenheder (οὐσίαι τινές) eller ikke. „I
Sammenligning med de Ting, man vel meest maatte tillægge Væsenhed,
saasom Vand, Jord og Luft, hvoraf de sammensatte Legemer bestaae, ere
Varme og Kulde og lignende Beskaffenheder (πάθη, Affectioner) jo ikke
Væsenheder; kun Legemet, der har disse Affectioner, forbliver som et
Værende, en Væsenhed (ὡς ὄν τι καὶ οὐσία τις οὖσα). Og dog synes Legemet
at have mindre Væsenhed end Overfladen, denne mindre end Linien, og
Linien mindre end Eenheden og Punktet. Thi med disse er Legemet
begrændset, og det synes, som kunde disse vel bestaae uden Legemet, men
Legemet derimod umuligt uden dem. Indrømmer man imidlertid, at Linierne
og Punkterne have mere Væsenhed end Legemerne, kommer man i Forlegenhed;
thi disse Grændsebestemmelser beroe øiensynlig paa Delinger af Legemet,
(διαιρέσεις ὄντα τοῦ σώματος), deels i Breden, deels i Dybden, deels i
Længden. Desuden maa enten enhver Form indeholdes i Legemet eller ingen:
dersom altsaa Hermes ikke er i Stenen, og Tærningens Halvdeel ikke findes
afgrændset i Tærningen selv, saa vil Overfladen heller ikke findes her;
thi dersom nogen Overflade fandtes i Tærningen, maatte ogsaa den findes,
hvorved Tærningen halveres, o. s. v.“ De Vanskeligheder, som i denne
Henseende reise sig, naar der spørges om, hvad der i Sandhed er det
Oprindelige: Punktet eller Linien, Linien eller Fladen, Fladen eller
Legemet, Legemet selv som existerende Ting eller de Beskaffenheder, der
udgjøre Tingens væsentlige Egenskaber, gjentage sig i nye Vendinger og
med fornyet Eftertryk, naar der udtrykkelig spørges: hvad er i absolut
Forstand det Oprindelige, Selvstændigheden eller det
% --- p. 115 ---
\setcounter{footnote}{0}%
% p. 115 is the third page of gathering 8 and carries the signature „8*“
% at the foot. Not transcribed.
Selvstændige, Virkeligheden eller det Virkelige, Substantialiteten eller
Substansen, Causaliteten selv eller Aarsagen? For at behandle slige
Opgaver maa man holde skarpt Øie med de vexlende Synspunkter og være
meget nøieregnende med de kategoriske Bestemmelser.

Seet fra Endelighedens Synspunkt er det umiddelbart Concrete at opfatte
som det Oprindelige, som det i sig Selvstændige i Modsætning til de deraf
abstraherede Formbestemmelser. I denne Henseende gjælder det uden
Modsigelse, at Selvstændigheden maa beroe paa det Selvstændige, ikke
omvendt. Den aristoteliske Kritik af Platos Ideer viser os fra mange
Sider det Mislige i at selvstændiggjøre den rene Væsenhed, og f. Ex.
lade Bordheden (ἡ τραπεζότης ἐν τῇ τραπέζῃ) være det Oprindelige, altsaa
det i sig Substantielle, Bordet, det virkelige Bord, det Afledede. Men
hvormeget Aristoteles end for Resten har indskærpet, at de existerende
Enkeltvæsener ere de egentlige Substanser, at Substansen (ἡ πρώτη οὐσία)
altid maa være et bestemt Dette (τόδε τι), et Individ, saa har han dog
heller ikke selv været klar over Forholdet mellem det Concrete og det
Abstracte. Saaledes gjør han f. Ex. i sin bekjendte Construction af de
fire Elementer det Abstracte til det Oprindelige, det Concrete til det
Afledede. „De elementære Egenskaber: hed, kold, tør og fugtig, fremhæves
nemlig saaledes, at hver af dem sættes som Væsenhed for sig, en
substantiel Form, der bestemmer Materien; saaledes giver da en samtidig
Forening af Tørhed og Varme Ildelementet, af Varme og Fugtighed
Luftelementet, af Fugtighed og Kulde Vandelementet, af Kulde og Torhed
% printer, p. 115: „Torhed“ here with a PLAIN o, against „Tørhed“ with a
% plainly slashed ø two lines above and again two words below, and „tør“ /
% „tørre“ elsewhere on the same page. Settled at 600 dpi on the positive
% presence of the slash in the one and its positive absence (clean, closed
% counter) in the other. The book's ø-dropping is inconsistent within the
% page; stands as printed.
Jordelementet.“\footnote{Jvfr Forelæsn. over „Phil. Propæd.“ 1860--61,
S.\ 84 o.\ flgd.}
% printed mark: *). „Jvfr“ WITHOUT the abbreviation point, checked at
% 600 dpi (clean white between the r and the F) — against „Jvfr.“ with the
% point in the note on p. 121. As printed.
Men naar den rene Tørhed paa denne Maade kan tørre Tingen, saa kan den
rene Bordhed jo ogsaa fabrikere Bordet.

% --- p. 116 ---
\setcounter{footnote}{0}%
Fra Uendelighedens, d. v. s. Væsenhedens Standpunkt seer Sagen anderledes
ud. Her er det nu Modsætningen, ikke mellem det Abstracte og det
Concrete, men imellem to Former for det Concrete, Reflexionens og
Umiddelbarhedens, hvorom det gjælder. Det synes jo vel, naar man tager
Udtrykket Reflexion i den vulgære Betydning, som om det Concrete absolut
maatte være til paa umiddelbar Maade: „en Steen“, siger man, „er noget
Concret; et Begreb, en Idee er derimod noget Abstract“. Men, hvad denne
Misforstaaelse angaaer, har Hegel tilstrækkelig viist, at det Ideelle
fuldt saavel kan være concret som det umiddelbart Reelle. See vi nærmere
til, da er det saa langt fra, at Reflexionsformen skulde være Formen for
det Abstracte, at den gjennemførte Reflexion, i Modsætning til
Umiddelbarheden, netop er Formen, den eneste Form, for det absolut
Concrete. At Umiddelbarheden ikke kan være Concretionens fuldendte Form,
følger ligefrem deraf, at det ene umiddelbart Concrete altid maa befinde
sig i en endelig Modsætning til det andet. Den endelige Modsætning
betyder Udstykning, og Udstykningen er det Modsatte af Concretion. Skal
Stenen forestille det Concrete, saa have vi i den enkelte Steen kun et
Brudstykke. Og hvor eensidigt er saa ikke Concretionen i et saadant
Brudstykke? Qvartsen f. Ex. er et andet Concret end Feldtspathen, som
igjen er et andet Concret end Hornblenden, o. s. v. Kun paa eet Vilkaar
kan en Concretion af Alt, en uendelig og absolut Concretion, finde Sted;
hertil udfordres nemlig, at de Tilværelses\-%
bestemmelser, der sprede sig i ydre Existens, opfattes i indre
Sammenhæng. Jo mangfoldigere det Udvortes er i sin reelle Udstykning,
desto rigere er da det Indvortes i sin ideelle Eenhed. Den ideelle, i
Reflexion af Realiteternes Mangfoldighed berigede, Eenhed er Inderlighed;
Inderlighed er væsentlig Concretion: jo dybere Inderlighed, desto rigere
Concretion. Seet fra dette Synspunkt er den absolute Substantialitet ikke
blot concret, men alsidig concret, det eneste i
% --- p. 117 ---
\setcounter{footnote}{0}%
sig absolut Concrete; enhver i umiddelbar Form existerende Substans er i
Modsætning hertil kun et eensidigt Concret, et enkelt af de utallige Led,
hvori det ydre Concrete er udstykket.

Endnu en Misforstaaelse er her at tage Hensyn til, naar der spørges om
det Oprindelige og det Afledede. Det gjælder som Lov paa Virkelighedens
Trin, at al Væren er Virken. Den tilværende Substans kan ikke være til
uden at virke paa en anden tilværende Substans; her maa altsaa i
Virkeligheden være idetmindste to Substanser, af hvilke den ene virker,
den anden paavirkes: det i den Virkende Virksomme bestemmes som Aarsag,
Forandringen i den Paavirkede som Virkning. Naturlig seet skulde
Aarsagen altsaa være i Tiden først, og Virkningen følge bag efter.
Tidsbestemmelsen melder sig i dette Forhold med en saadan
Paatrængenhed, at en virkelig Succession let kan synes at være det for
Forholdet selv Charakteristiske. Fra dette Synspunkt har Kant i sin
„Kritik af den rene Fornuft“ da ogsaa opstillet sin „Grundsætning for
Tidsfølge efter Causalitetsloven“\footnote{„Der Begriff, der eine
Nothwendigkeit der synthetischen Einheit bei sich führt, kann nur ein
reiner Verstandesbegriff sein, der nicht in Wahrnehmung liegt, und das
ist hier der Begriff des Verhältnisses der Ursache und Wirkung, wovon die
erstere die letztere in der Zeit, als die Folge, und nicht als etwas, was
bloß in der Einbildung vorhergehen (oder gar überall nicht wahrgenommen
sein) könnte, bestimmt.“ Kritik der reinen Vernft. Leipzig 1838, S.\
196--97.}.
% printed mark: *) — the closing parenthesis of the mark AT THE FOOT of
% the page is under-inked and reads as a comma in both tesseract witnesses;
% it is the same *) sort, defectively printed. The German note is set in
% FRAKTUR throughout (not antiqua), so it takes no \textit{}.
% „Vernft.“ (for Vernunft) with a Fraktur V, calibrated at 600 dpi against
% the certain V of „Verhältnisses“ and the certain B of „Begriff“ on the
% note's own third line; not „Bernft.“.
Skulde det nu virkelig forholde sig saaledes, at Aarsagen som Aarsag
maatte være i Tiden nogle Øieblikke forud for sin Virkning, saa maatte
Kategorierne Aarsag og Virkning ophøre at være correlate Begreber. Men en
Aarsag kan ligesaalidt tænkes uden Virkning, som en Grund uden Følge: i
samme Nu som Aarsagen er, er ogsaa Virkningen; i samme Nu som Virkningen
ophører, ophører ogsaa Aarsagen. Hvad Grunden er
% --- p. 118 ---
\setcounter{footnote}{0}%
i Væsenet, det er Aarsagen i Virkeligheden\footnote{„At Tingen maa have
sin Grund, vil sige, at dens egen indre Forudsætning ikke selv kan være
en Ting, men maa være en i sig reflecteret Eenhed, et System af
virksomme Relationer, et Væsenssystem, hvori Tingsaarsagerne ere
nedsatte til Aarsagsmomenter. Saaledes kan t. Ex. Ilten ikke være
Grunden til Vandet, saalidtsom Brinten. Ilten og Brinten ere, idet
Forbindelsen indgaaes, medvirkende Aarsager, og deres betingede
Vexelvirkning altsaa Vandets nærmeste Aarsag; i denne Henseende gjælder
det, at enhver Ting maa have sin Aarsag. Men Et er Aarsagen, et Andet er
Grunden. Grunden til Vanddannelsen, nemlig Realgrunden, er den
Forvandling, der i Foreningens Øieblik foregaaer med Iltens og Brintens
Atomer. Uden en saadan Forvandling er Vandets Tilblivelse --- naar det,
vel at mærke, skal være et virkeligt, ikke blot for „vore ufuldkomne
Sandser“ tilsyneladende Vand --- umulig.“ Forelæsninger over „Phil.
Propæd.“ Aar 1860--61, S.\ 165--66.}. Aarsagsforholdet kan derfor ikke
opstilles for sig ved Siden af Grundforholdet; thi Aarsagsforholdet er i
sig selv ikke Andet end Grundforholdets egen
Virkeliggjørelse\footnote{Om det kategoriske Forhold mellem
\emph{Grund} og \emph{Aarsag} bemærker Heiberg i sin Logik (Anfr. Skr.
S.\ 47) med sædvanlig Skarphed og Klarhed: „Naar man giver \emph{Grunde}
for Noget, saa har man endnu ikke angivet \emph{Grunden}. Herved
adskiller Grunden sig fra \emph{Aarsagen}. Grunden kan have Vorden, og
saaledes er det ubestemt, om den skal faae Tilværen eller ikke. (Man kan
have Grund til at foretage en Handling og dog lade det være, thi det er
vilkaarligt, om og hvor man vil afbryde den uendelige Række --- af Grunde
\textit{pro et contra}). Men Aarsagen er selv Tilværen, nemlig af
Nødvendigheden, og maa derfor nødvendigen have Virkning. En Grund til at
handle behøver saaledes ikke at være virksom: dette bliver den først ved
at optages i en Villie, som giver den Tilværen, eller maaskee endog gjør
den til Aarsag.“ --- Senere hen (Anfr. Skr. S.\ 75) bemærker han
angaaende „Vexelvirkningen“: „Naar man betragter Aarsagen som blot activ,
og Virkningen som blot passiv med Hensyn paa Aarsagen, skjøndt vel som
activ med Hensyn paa en følgende Virkning, saa bliver man
% the note runs over the printed page break here: p. 118 | p. 119
i den uendelige Række af Aarsager og Virkninger, altsaa i den endelige
Causalitet, i hvis Ophævelse Causalitetens Sandhed bestaaer. I denne
Ophævelse er nemlig Aarsagen ogsaa passiv, idet den ophæves i
Virkningen, og herved bliver denne selv activ, eller reagerer som Aarsag
imod sin egen Aarsag.“ Endvidere: „Ved Reaction maa man ikke her tænke
sig den, som finder Sted mellem to Substanser, af hvilke enhver virker
som Aarsag (f. Ex. naar en Kugle, skudt imod en haard Gjenstand, springer
tilbage, eller naar Solens Varme, formedelst de jordiske Substansers
Reaction, bliver Aarsag til Planternes Fremkomst), thi her er endnu ikke
Tale om een Substans, og endnu mindre om to eller flere. Den Reaction,
som her omtales, er den, som følger af en eneste oprindelig Aarsag, idet
dens Virkning blot som dens egen Virkning, ikke som ny Aarsag,
reagerer.“ --- Imod det Allersidste kunde her imidlertid reise sig
adskillige Betænkeligheder. Til den endelige Causalitet maa jo
nødvendigviis fordres, at Aarsagen ikke blot adskiller sig formelt fra
sin Virkning, men at Aarsagen har Tilværelse i en, og dens Virkning i en
anden Existens. Om man nu, med Hegel, vil have denne Existens iforveien
udtrykkelig bestemt som Substans eller indtil videre lade sig nøie med at
opfatte den som Ting, kan vistnok være ligegyldigt; men, naar Virkningen
i „den endelige Causalitet“ har sin Existens for sig, ligesaa vel som
Aarsagen, er det da ikke heraf en uundgaaelig Consequens, at Virkningens
Existens tillige maa fungere som Aarsag og Aarsagens som Virkning?
Reactionen maa jo i det Endelige netop være betinget af to
Exi\-%
stenser, „af hvilke enhver virker som Aarsag“.}.
% p. 118 CARRIES TWO FOOTNOTES, printed *) and **) — the first page in the
% batch to do so. Both are set here with the plain \footnote{}, which the
% preamble renders *) in each case; the second mark is NOT invented.
% Letterspacing inside the **) note confirmed at 600 dpi: „Grund“ and
% „Aarsag“ in the opening line, and „Grunde“, „Grunden“, „Aarsagen“ in the
% Heiberg quotation. The following „Grunden kan have Vorden“ is NOT spaced.
% „pro et contra“ is antiqua.
At der nu, hvad
% --- p. 119 ---
\setcounter{footnote}{0}%
Grundforholdet angaaer, ikke kan være nogen endelig Tidsadskillelse
mellem Grunden og dens Følge, er iøinefaldende. Men hvoraf kommer det da,
at Forestillingen om en Tidsfølge trænger sig saa stærkt frem i
Aarsagsforholdet? Det kommer deraf, at det Skin af Succession, der
danner sig i enhver Tanke- og Reflexionsbevægelse, erholder en desto
væsentligere Betydning, jo større Selvstændighed Momenterne opnaae, jo
mere den gjennemførte Reflecteren gaaer op i Virken. For en overfladisk
Betragter af Virkeligheden seer det ud, som om Aarsagen selv var en
Substans, eller med et
% „Tanke= og Reflexionsbevægelse“ is a SUSPENDED compound (checked at
% 600 dpi) and keeps its printed hyphen and its space; it is not joined.
% --- p. 120 ---
\setcounter{footnote}{0}%
mere gængse Udtryk, en Ting. Men at gjøre Aarsagen til en Ting er
ligesaa urimeligt, som at gjøre Kraften til en Materie. Tingsaarsagen
virkeliggjør kun Aarsagsbegrebet paa en aldeles udvortes, endelig og
eensidig Maade. Vel er Tingen Aarsag, forsaavidt den er Krafternes
materielle Eenhed, det existentielle Midtpunkt, hvorfra Virkningen
udgaaer; men Aarsagen selv er desuagtet ikke en Ting. Tingen forgaaer,
men Kraften, det Foraarsagende, forgaaer ikke; den virker kun i anden
Sammenhæng og paa anden Maade, identisk med sig selv som en anden Kraft,
naar Tingen er gaaet tilgrunde.

At Aarsagen umulig kan have Tilværelse for sig i den ene Ting og
Virkningen i den anden, kan da blandt Andet sees af Stødet. Det Virkende
paavirkes, og det Paavirkede virker, Stødet er Vexelvirkning: i
Vexelvirkningen er Aarsagen Virkning og Virkningen Aarsag. Men ihvorvel
Aarsagen er Eet med sin Virkning i Stødet, derfor blive to sammenstødende
Kugler dog ikke til een Kugle. Derimod bliver Aarsagen i den ene ikke
alene Eet med Virkningen i den anden; men i enhver af dem bliver Aarsagen
formedelst Reactionen sin egen Virkning og Virkningen omvendt sin egen
Aarsag\footnote{Under Sammenstødet af $A$ og $B$ svarer Aarsagen i $A$ til
Virkningen i $B$, og omvendt, Aarsagen i $B$ til Virkningen i $A$. Dette
maa, naar den causale Reflexion skal gjennemføres, igjen udtrykkes
saaledes: Aarsagen i $A$ forholder sig --- formedelst Reactionen fra $B$
--- til Virkningen i $A$, ligesom Aarsagen i $B$ --- formedelst
Reactionen fra $A$ --- forholder sig til Virkningen i $B$.}.
% printed mark: *). The antiqua letter-symbols are set in math mode. The
% second sentence of the note reads „… Aarsagen i A forholder sig …
% til Virkningen i A …“ — the second A checked twice at 600 dpi against
% the certain B of the same line; it is an A, not a B, and is not
% „corrected“ here.

Forskjellen mellem Tingen og Tingsaarsagen er endvidere kjendelig
derpaa, at den ene og samme Ting, f. Ex. Svovl, kan i flere forskjellige
Henseender blive Aarsag, Aarsag til Gnidning, til Lyd, til Varme, o. s.
v.; der kan paa denne Maade være mange og forskjellige Aarsager i een
Ting. Omvendt, kan den ene og samme Tingsaarsag gjentage sig i
% --- p. 121 ---
\setcounter{footnote}{0}%
utallige Ting; skal et Tryk f. Ex. være Aarsag til en Bevægelse, da kan
et Pund Vand jo udøve samme Tryk, altsaa afgive samme Aarsag, som et Pund
Bly, et Pund Jern o. s. v.

Som tilværende, i Tilværelsen virksomme, Substanser ere de virkelige Ting
altsaa paa een Gang Aarsagsmomenter og existerende Aarsager.
Tingsaarsagen er i denne Belysning baade identisk med Tingen og
forskjellig fra Tingen. Fremhæves nu Forskjellen, saa forvandle de
existerende Aarsager sig til Aarsagsmomenter, substantielle
Gjennemgangspunkter for den uendelig concrete Selvaarsag. Hertil svarer
saa den kantiske Thesis, at vi foruden den naturlige Række af betingede,
hinanden i det Uendelige forudsættende Aarsager nødvendig maae
forudsætte en i sig ubetinget, sig selv bestemmende Aarsag, „en absolut
Aarsagernes Spontaneitet“, en „transcendental Frihed“. Lægger man derimod
det hele Eftertryk paa Eenheden af Ting og Tingsaarsag, saa forsvinder
Selvaarsagen; vi ledes da til at antage den mod Thesis opstillede
Antithesis, at der nemlig „ingen Frihed er, men at Alt i Verden skeer
alene efter Naturens Love“\footnote{Jvfr. Forelæsninger over „Phil
Propæd.“ 1861--62, S.\ 139--43.}.
% printed mark: *). „Phil Propæd.“ WITHOUT the abbreviation point after
% „Phil“, checked at 600 dpi (clean white gap, no point at the baseline) —
% against „Phil. Propæd.“ with the point in the note on p. 115. As printed.
% „Thesis“ and „Antithesis“ are FRAKTUR here, not antiqua (600 dpi).
% A small round ink speck stands in the left margin beside the second line
% of this paragraph; not transcribed.

Det er en aldeles rigtig Tanke, Kant har udtalt i Beviset for sin Thesis,
at der nemlig ingen sand Realitet vilde være i de relative Aarsager,
dersom der ikke gaves en første oprindelig Aarsag; men denne i og for sig
sande og sunde Tanke er i Beviset for Antithesis bleven forqvaklet og
gjort ukjendelig ved det logiske Monstrum af en transcendental Frihed,
der ikke er Frihed, men abstract Vilkaarlighed. Skal $A$ være en Virkning
af $B$, $B$ af $C$, $C$ af $D$ o. s. fr., saa faae vi ganske rigtig, hvad
Kant paa sin Maade gjør gjældende, en Virkning af en Virkning af en
Virkning --- uden Aarsag. Men den i det Uendelige bortflygtende Aarsag
vender jo strax tilbage, naar vi blot lade Virkningerne danne et
Kredslob.
% printer, p. 121: „Kredslob.“ with a PLAIN o, settled at 600 dpi on the
% clean closed counter — against „Kredsløbet“ with a plainly slashed ø
% four times on p. 122 immediately following. Stands as printed.
% --- p. 122 ---
\setcounter{footnote}{0}%
For at den udvortes Mangfoldighed ikke skal virke forstyrrende, vil det i
denne Sammenhæng være tilstrækkeligt at lade hele Kredsløbet afsluttes
med tre Virkninger. $A$ være saaledes en Virkning af $B$, som igjen er en
Virkning af $C$, som igjen er en Virkning --- hvoraf? Af $A$! $A$ er
altsaa i Kredsløbet sin egen Aarsag og sin egen Virkning, og forsaavidt
et selvstændigt Led; men $A$'s hele Virken er i det uendelige Kredsløb
tillige et i det Uendelige Bevirket. At $A$ har sin Aarsag i sig selv,
vil da egentlig sige, at det har sin Aarsag udenfor sig selv, nemlig i
$B$ og $C$; som et saaledes Bevirket er $A$ kun et Moment i Kredsløbet.
Hvad der gjælder om det enkelte Led, gjælder om dem alle: de ere
selvstændige og dog Momenter.

Med Reflexionseenheden af Momentet og det selvstændige Led maa
Vexelvirkningen nødvendig fordoble sig. Her er i Kredsløbet en dobbelt
Vexelvirkning, en endelig eller reel og en uendelig eller ideel. Den
endelige Vexelvirkning bestemmer de selvstændige og dog betingede Leds
indbyrdes Afhængighed: den viser os Tingsaarsagernes endeløse Rækker, den
fremstiller den causale Reflexion endeliggjort i Leddenes substantielle
Umiddelbarhed. Den ideelle Vexelvirkning sees derimod i den uendelige
Reflecteren selv, der bærer Kredsløbet og giver det sit
allestedsnærværende Midtpunkt. For denne Reflecteren ere de substantielle
Led kun Gjennemgangspunkter; her have vi det i de betingede
Aarsagsrækker Oprindelige, ikke som et enkelt tilværende Led, men som
Selvaarsag, som absolut Causalitet.

Den absolute Causalitet er altsaa ikke, hvad Kant antager, et Ubetinget
for sig i endelig Modsætning til de endelige Betingelser; tvertimod! det
Ubetingede virker i et uendeligt System af Betingelser og er forsaavidt
selv betinget. Men betinget og betinget betyder for det Endelige og det
Uendelige noget aldeles Forskjelligt. Det endelige Led er betinget
saaledes, at dets Betingethed netop udtrykker dets
% --- p. 123 ---
\setcounter{footnote}{0}%
Afhængighed, efterdi det altid er betinget af sit Andet. Det Uendelige er
vel ogsaa betinget; men dog saaledes, at det igjennem Betingelsernes
Kredsløb forholder sig til sig selv, og bestemmer de Betingelser, hvoraf
det betinges. Det Ubetingedes Betingethed er Virkelighedens Udtryk for
dets ideelle Concretion, dets Oprindelighed og absolute Selvstændighed.
Selvaarsagen er med andre Ord de endelige Aarsagers Idealitet, og disse
ere igjen Selvaarsagens Realitet. Den causale Idealitet er en mægtig
Viden; derfor ere alle virkelige Aarsagsforhold intellectuelle, alle
Causalbestemmelser oprindelige Vidensbestemmelser. Den causale Realitet
reflecterer en vidende Magt; derfor ere de intellectuelle
Aarsagsmomenter tilværende Substanser, existerende Ting i en virkelig
Tingenes Orden.

\subsubsection*{γ) Den totaliserende Reflexion.}
% Display head as printed: Greek γ in italic antiqua, then letterspaced
% Fraktur (structural, not marked \emph{}). Read off the image at 600 dpi.

„Der er i vor umiddelbare Viden Meget, der jo forsaavidt maa falde
udenfor Tænkningen, som det hører ind under Sandsningen, og, hvad
Sandsningen angaaer, Meget, der falder udenfor Synet, forsaavidt det f.
Ex. falder ind under Hørelsen o. s. v.; altsaa“ --- kunde Nogen sige ---
„hvormeget maa der ikke være til i den virkelige Verden, som falder
aldeles udenfor al Reflexion og derved udenfor al Viden!“ Vi svare: der
er naturligviis mangfoldige Ting, der falde udenfor din og min Reflexion
og derved udenfor din og min Viden; men hvad den ontologiske Reflexion
angaaer, er der Intet til, Intet i Verden, Intet hverken i Himlene eller
paa Jorden eller i Afgrundene, der enten falder eller kan falde udenfor
Reflexionen i sig selv, den substantielle Reflexion, der evig virker som
Causalitetens uendelige Selvreflecteren. Det gaaer med den ontologiske
Reflexion i dens Forhold til Virkeligheden, som med vor psychologiske
Reflexion i Forhold til de Gjenstande, vi blive os bevidste: ethvert
Udenfor gaaer i Eet med sit Indenfor, og dog maa et Uden\-%
% --- p. 124 ---
\setcounter{footnote}{0}%
for og et Indenfor bestandig adskilles. Synets Gjenstand falder virkelig
udenfor Synet, forsaavidt den falder udenfor den Seende; og dog er
Gjenstanden i sig selv kun en Synsgjenstand, forsaavidt den falder i Synet.
Billedet i Speilet er forskjelligt fra den afspeilede Gjenstand, og dog falder
Gjenstanden i sin Afspeiling sammen med Billedet: vi kunne ikke blive os
Billedet bevidste som Billede, uden at blive os en til Billedet svarende
Gjenstand bevidste som Gjenstand. Reflexionen har overalt den uendelige
Smidighed, at den under sine Fordoblinger optager i sig, hvad den sætter
udenfor sig, og tilegner sig, hvad der stilles imod den. Sættes
Umiddelbarheden mod Reflexionen, da er det ikke det Umiddelbare, der
bestemmer sig uden Reflexion, men det er omvendt Reflexionen, der bestemmer
det Umiddelbare ved at bestemme sig selv. Skal en høiere, d. v. s. rigere
Umiddelbarhed kjendes derpaa, at Reflexionen, hvoraf den fremgik, igjen er
ophævet, da vedbliver den ophævede Reflexion dog endnu at bevare Reflexen af
et Ophævet; og saaledes i alle andre Henseender\footnote{Træffende bemærker
Heiberg i sin Logik (Anfr. Skr. S.\ 36), at „i følgende Kategorier, af hvilke
de To, som udgjøre hvert Par, staae i Reflexionsforhold til hinanden, saasom
Tingen og Egenskaberne, det Hele og Delene, Materie og Form, Aarsag og
Virkning, Substans og Accidens, \emph{har} det første Led det andet, og
omvendt. F. Ex. a) Aarsagen har Virkning. Saaledes er den resulterende Grund,
hvori Virkningens Væren er ophævet, d. e. Aarsagen har selv været Virkning,
eller en tidligere Aarsag har havt den. Den Virkning, den selv skal
frembringe, er endnu immanent i den, hvorved den ganske falder sammen med den
allerede ophævede Virkning. Men b) Virkningen har Aarsag; thi naar Virkningen
har faaet Væren, er Aarsagen ophævet i den. Og endelig er c) Vexelvirkningen
den Eenhed, hvori Aarsagen har Virkning, og Virkningen har Aarsag. --- At
Grunden har Væren, og Væren har Grund, kan ogsaa udtrykkes ved, at Grunden
\emph{har været}. Den har nemlig været a) som umiddelbar Væren, inden denne
blev ophævet i Grunden, og b) som Grund, inden denne
% the note runs over the printed page break here: p. 124 | p. 125
blev ophævet i den umiddelbare Væren. Begge Dele høre til det Forbigangne, og
udtrykkes derfor ved at \emph{have været}. Ordet „have“ betegner, som vi have
seet, den ophævede Væren, og denne er det Forbigangne. Derfor er det en rigtig
Conseqvens af visse Sprog, f. Ex. det danske, at udtrykke den forbigangne Tid
med Verbum „have“, medens andre f. Ex. det tydske, gjøre det med Verbum
„være“. („Det \emph{har} været“; „es ist gewesen“).“

Men naar man nu ret besinder sig paa, hvad hermed er sagt, saa maa det dog
omsider indlyse, at al objectiv logisk Reflexion nødvendig maa foregaae i en
reflecterende Subjectivitet. Objecterne i deres existentielle Umiddelbarhed
kunne umulig enten have eller bevare enten „den ophævede Væren“ eller „den
ophævede Reflexion“. Vel har Tingen et Væsen, altsaa „en ophævet Væren“; men
den har kun dette Væsen i og formedelst den substantielle Reflexion, der evig
bæres af Subjectiviteten.}.
% Printed footnote mark *), set with the plain \footnote{}. The note is the
% only one on p. 124; its text runs on into the footnote area of p. 125, and
% the second paragraph („Men naar man nu ret besinder sig …“) is Nielsen's own
% continuation of the same note — there is no second printed mark.
% The enumerators a) b) c) inside the Heiberg quotation are ANTIQUA in the
% print, like the Greek γ) of the display head on p. 123; transcribed plainly,
% not as \textit{}. Letterspacing confirmed at 600 dpi: „har“ (before „det
% første Led“), „har været“ (p. 124), „have været“ and the „har“ of „Det har
% været“ (p. 125). The closing „været“ of „Det har været“ is NOT spaced.
% „es ist gewesen“ is Fraktur, not antiqua — no \textit{}.

% --- p. 125 ---
\setcounter{footnote}{0}%
Er det nu indlysende, at alle Virkelighedsbestemmelser maae, som Reflexer og
Reflexionsbestemmelser, gaae op i Reflexionen, saa er det med det Samme
indlysende, at hele Virkelighedens uendelig rige Indhold maa falde ind under
den reflecterende Subjectivitet. Det subjective Indhold er i Reflexionen selv
objectivt: Reflexionen er \emph{totaliserende}.

Det er den substantielle Eenhed af de utallige relative Substanser,
Aarsagseenheden af alle relative Aarsager i deres alsidige Virken og
Vexelvirken, hvormed Totaliseringen fuldendes. Her er da i Virkeligheden ---
efter hvad vi have seet --- to virkelige Totaliteter, en virkelig indvortes og
en virkelig udvortes; og dog er den udvortes kun virkelig formedelst den
indvortes, og omvendt: de to Totaliteter ere i Virkeligheden een.

I Totaliteternes Eenhed sees Principet. Principet selv, det absolute Princip,
er ikke et Subjectivitetsprincip for sig, saa lidt som det er det abstract
Objectives eensidige Princip; thi Subjectiviteten har i sin Reflecteren baade
det Subjective og det Objective, og Subjectiviteten er Principet. Men saa
uendelig er den totaliserende Reflexion, at Subjectiviteten,
% --- p. 126 ---
\setcounter{footnote}{0}%
uagtet den er Principet selv, tillige kan oversvæve Principet, idet den kan
oversvæve sin egen dobbeltsidige Udvikling. Principet, hvis ubetingede Virken
gaaer op i Udviklingen og det Udviklede, adskiller sig da med særligt Hensyn
paa de to Totaliteters forskjellige Udviklingsmaader i to relativt
forskjellige Principer, Principet for den subjective eller ideelle og
Principet for den objective eller reelle Totalitet.

I den indre, subjective og ideelle Totalitet blive alle
Existensbestemmelser reflecterede som rene Vidensbestemmelser: denne Totalitet
er i Eenhedens Form identisk med den oprindelig concrete Selvhed, og maa
derfor kaldes en Selvhedstotalitet. I den ydre, objective og reelle Totalitet
blive Vidensbestemmelserne reflecterede som reale Magtbestemmelser; denne
Totalitet er i Mangfoldighedens Form et System af Andet mod Andet, og maa
kaldes en Andethedstotalitet. Den principielle Vexelbestemmelse af de to
Totalitetsudviklinger er mere end en tom Speilen; den er en i alle Momenter og
Led articuleret Virken, en Vexelvirken, hvis principielle Eiendommelighed
noiagtig lader sig bestemme.

Principet for Selvhedsudviklingen bestemmes ved en Bevægelse af det Samme
gjennem Andet tilbage til det Samme, en Bevægelse, hvori Reflexionen idelig
hæver Successionen. Saaledes er vel Synet af Dette forskjelligt fra Synet af
Hiint; men Synet af Dette er kun dette Syn, Synet af Hiint kun hiint Syn;
dette og hiint Syn ere Forskjelsbestemmelser i den samme Seen, i Synernes
reflecterede Eenhed, i den seende Selvhed. Paa tilsvarende Maade er Tanken af
Dette og Tanken af Hiint, det vil sige: denne og hiin Tanke, to Tanker i een
Tænken; Forestillingen om Dette og Forestillingen om Hiint, d. v. s. denne og
hiin Forestilling, to Forestillinger i een Forestillen. Det udvortes Andet
reflecterer sig i alle Punkter som et indvortes, og det indvortes som en
Forskjelsbestemmelse i den sig articulerende Eenhed. Vistnok maa Dette
bestemmes ved en Seen, naar Hiint bestemmes
% --- p. 127 ---
\setcounter{footnote}{0}%
ved en Tænken, Dette beroe paa en Forestillen, naar Hiint beroer paa en
Fornemmen; men den Seende tænker, den Tænkende forestiller, den Forestillende
fornemmer: Syn og Tanke, Forestilling og Fornemmelse ere Sidebestemmelser ved
den samme Viden, Reflexer i den samme Reflecteren. --- Principet for
Andethedens Udvikling gaaer i omvendt Retning. Her sees den oprindelige
Eenhed kun ved en Speilen i den umiddelbart fremtrædende Forskjel; det
indvortes Andet er som et Skin mod det udvortes, thi Andethedens Eftertryk
falder paa det Udvortes. Idet Andethedens existerende Led sættes som Udtryk
for den reale Magt, bliver Magtmangfoldigheden reflecteret i den ydre
Totalitet, og den ydre Totalitet igjen adskilt i en Mangfoldighed af relative,
ved deres indbyrdes Modsætning forskjelligartede Totaliteter.

I den meest udvortes af alle udvortes Former, den mechaniske nemlig, har det
existerende Led en saa umiddelbar Selvstændighed, at Selvstændigheden, seet
fra Endelighedens Side, virkelig seer ud, som om den var absolut. Atomer og
Masser antage i denne Belysning Skin af en evig, paa sig selv beroende
Existens. Det i denne umiddelbare Selvstændighed Selvløse har sit kategoriske
Udtryk i Materiens Inertie. „Materien“, siger Physiken, „er villielos“, d. v.
s. den kan ikke bestemme sig selv til nogen Forandring. Ved sin døde
Selvstændighed, sin Ligegyldighed mod den Tilstand af Hvile eller Bevægelse,
hvori det er sat, er det materielle Led fra alle Sider underlagt
Andethedsforholdet: det kan ikke bevæge sig selv, men det kan bevæge et andet
og bevæges af et andet; det tiltrækker ikke sig selv, men det tiltrækker et
andet, og tiltrækkes af et andet, o. s. v. Magtaabenbarelsen er i denne
udvortes Form saa endelig, Udstykningen saa fremherskende, at ethvert Hele
synes at være samlet og sammensat af forud existerende Dele.
% „villielos“ is printed with an UNSLASHED o, while „Selvløse“ four lines
% above it, on the same page, has a plainly slashed ø; both verified at 600
% dpi and transcribed as printed (cf. „selvlos“ elsewhere in the book).

Der gjøres som bekjendt den Adskillelse imellem et mechanisk og et organisk
Hele, at Delene i hiint altid ere det
% --- p. 128 ---
\setcounter{footnote}{0}%
Første, medens Delene i dette derimod udvikles af det Heles Anlæg, saa at
Heelheden i Anlæget selv er det Første. Denne Adskillelse er, hvad det
Logiske angaaer, vildledende. Der gives i logisk Forstand ingen Totalitet,
altsaa heller ingen mechanisk Totalitet, hvorom det med Grund kan siges, at
Delene i Virkeligheden ere det Første og at det Hele kun fremkommer som et
Resultat af Delenes Sammensætning. De Exempler, man i denne Anledning pleier
at anføre, forraade tilhobe en øiensynlig Miskjendelse deels af den Maade,
hvorpaa et Hele bliver dannet, deels af Heelhedens eget Begreb. „Et Uhr“,
siger man, „er et mechanisk Hele“; og det maa indrømmes. Men derimod maa det
ikke indrømmes, at nogensomhelst af Uhrets Dele skulde være dannet i
fuldkommen Uafhængighed af det Hele; det er jo tvertimod vitterligt, at den
Form og Beskaffenhed, hver enkelt Deel erholder, ene og alene bestemmes ved
den Maade, hvorpaa den skal tjene sit Hele. Vil man hente Exempler fra Stene,
Krystaller og Jordarter, da egne slige Exempler sig kun til at vise, hvor
overfladisk og tankeløst man fatter Totalitetens Begreb. Vistnok kan en Steen,
en Sandsteen f. Ex., siges at være et Resultat af forud existerende Deles
Sammensætning; men et saadant Aggregat er jo kun et Brudstykke, ikke et Hele.
Begrebet Totalitet\footnote{Her tænkes udtrykkelig paa den reale Totalitet,
ikke paa Kategoriens abstracte Formudvikling. Fra den formelle Side kan
Totalitetsbegrebet fremkomme f. Ex. ved en Udvikling af Ting og Egenskaber,
Grund og Betingelser o. s v. Enhver logisk Concretion kan i Forhold til de
Momenter, den optager i sig, fra Formsiden bestemmes som Totalitet. --- „Den
Eenhed, at Tingen kun er (som ophævet) i Egenskaberne, og at Egenskaberne kun
ere, (som ophævede) i Tingen, er \emph{Realitet} eller \emph{Eenhed af Væsen}
(svarende til \emph{Eenhed af Væren}“, --- nemlig Eet). --- „Naar
Totaliteten, som er Eenheden af den phænomenale Væren, ligesom Realiteten er
Eenheden af den væsentlige, udvikler sine Momenter, Indhold og
% the note runs over the printed page break here: p. 128 | p. 129
Omfang, til et nyt og høiere Kredsløb, saa bliver Omfanget bestemt som
\emph{det Hele}. Indholdet derimod bestemmes som \emph{Delene}; og i denne
Reflexion har det Hele Dele, og Delene have (udgjøre) det Hele.“ Heiberg.
Anfr. Skr. S.\ 49; S.\ 63--64.

„Eenhed af Væsen“ kan i Modsætning til „Eenhed af Væren“ kaldes Totalitet med
samme Ret som „Eenheden af den phænomenale Væren“ i Modsætning til „Eenheden
af Væsen“.

Hvor abstract og formelt Heelhedsbegrebet her er opfattet, sees blandt Andet
af følgende Forklaring. „Delenes Kategorie bestemmer det Tilfældige i
Virkeligheden saaledes, at det Tilfældige er de \emph{Dele, som falde udenfor
det Hele}. I følgende Exempel
\begin{center}
\begin{tabular}{@{}c@{~}c@{~}c@{~}c@{~}c@{~}c@{~}c@{}}
\multicolumn{7}{c}{$C$}\\
\cline{1-7}
\multicolumn{3}{c}{$A$} & & \multicolumn{3}{c}{$B$}\\
\cline{1-3}\cline{5-7}
1. & 2. & 3. & 4. & 5. & 6. & 7.
\end{tabular}
\end{center}
er Tingen $A$ deelt i Delene 1, 2, 3, og Tingen $B$ i Delene 5, 6, 7; men
Delen 4 hører hverken til $A$ eller $B$, men falder udenfor dem begge. Delen
4 sees altsaa slet ikke som Deel, undtagen forsaavidt som den er Deel af den
hele Række fra 1 til 7. Skal den altsaa sees som Deel, saa maa det være af et
høiere Heelt $C$, der indbefatter baade $A$ og $B$ og den af begge udelukkede
Deel. Delen 4 er saaledes tilfældig for de to Heelheder $A$ og $B$, thi den
falder udenfor begge, men den er ikke tilfældig i den høiere Heelhed $C$,
hvoraf den er en Deel.“ Anfr. Skr. S.\ 64. (Prosaiske Skrifter, 1ste Bind, S.
241.)} fordrer udtrykkelig, at der maa
% Printed footnote mark *), the only note on p. 128; its text fills the whole
% footnote area of p. 129 as well, so p. 129 carries no call of its own.
% Printer's defect logged at the site: the print gives „o. s v.“ with NO
% period after the s (verified at 600 dpi; the same abbreviation on p. 127
% above is correctly „o. s. v.“). Transcribed as printed.
% Printed placement of the closing quotation mark kept as it stands: the „ of
% „Den Eenhed …“ is closed by the “ immediately after „Væren“, so that the
% parenthesis opened by „(svarende til“ closes OUTSIDE the quotation, after
% Nielsen's interpolated „--- nemlig Eet“.
% Letterspacing at 600 dpi: „Realitet“, „Eenhed af Væsen“, „Eenhed af Væren“;
% „det Hele“ and „Delene“ (p. 129); and the run „Dele, som falde udenfor det
% Hele“. The intervening „eller“ is NOT spaced.
% The figure: C stands over a rule spanning the whole numeral row 1--7 (the
% rule is printed as two abutting segments with a hair gap at the centre), and
% A and B stand below it over shorter rules spanning 1--3 and 5--7
% respectively; 4 carries no rule. A, B, C and the numerals are antiqua, so
% the letters are set in math mode as elsewhere in this edition.
% --- p. 129 ---
\setcounter{footnote}{0}%
være, ikke blot en qvantitativ, men ogsaa en qvalitativ Forskjel mellem det
Hele og dets Dele. I denne Henseende kan ikke engang Krystallen siges at være
et sandt Hele; thi enhver Deel, hvori Krystallen deles, er selv en Krystal,
altsaa i Qvalitet det Hele; hvoraf igjen følger, at det Hele, i Qvalitet, er
en Deel. Men, naar Delen er som det Hele og det Hele som Delen, saa er her i
væsentlig Forstand hverken Deel eller Hele. Overalt derimod, hvor
Heelhedsbegrebet fyldestgjøres, viser det sig, at enhver Totalitetsudvikling,
ogsaa den mechaniske, er en Principudvikling.
% „fyldestgjøres“: the y is blotted in the scan and both tesseract witnesses
% read „fnldest-“; the sense-reading is adopted (playbook §2, rule 3).

Men Principet er jo væsentlig Selvhed; hvorledes er det da muligt, at en
Selvhed kan constituere sig i det virkelige Selvløse? Det er Realiteten og det
umiddelbart Reelle,
% --- p. 130 ---
\setcounter{footnote}{0}%
hvorpaa Eftertrykket falder. For sig betragtede, danne Realitetens enkelte Led
i deres udvortes Mangfoldighed jo vistnok en, saa at sige, blind Modsætning
til den i dem alle virkende Selveenhed; men i Totaliteten selv, det mechaniske
System af Bevægelser og bevægende Aarsager, er Modsætningens Eenhed motiveret.
Den mechaniske Selvhed, reflecteret i Magtens Form, beviser just, hvad den
totaliserende Reflexion formaaer. Thi saa almægtig er den uendelige
Reflecteren, at Selvaarsagen, det Ubetingede selv, kommer til Bestaaen i en
Magtmangfoldighed af Betingelser og relative Aarsager; og saa reflecteret er
igjen den uendelige Magt, at de relative Aarsager med deres Betingelser
forvandle sig til Gjennemgangspunkter for den Totalitet af mechaniske
Bestemmelser, hvori Selvmagten objectiveres.

I den mechaniske Totalitet er Modsætningen af Selvhed og Andethed abstract; de
modsatte Bestemmelsers Eenhed er derfor ligeledes abstract. Vi tillægge Jorden
f. Ex. en Selvbevægelse, forsaavidt den under Bevægelsen i sin Bane tillige
bevæger sig om sin Axe; men Selvheden i denne Bevægelse er umiddelbart Eet med
Andetheden; Jordens Bestemmelse ved sig selv er fra alle Sider en Bestemmelse
ved Andet; dens Selvbevægelse er kun en Egenbevægelse. Det Abstracte og
forsaavidt Haarde i den mechaniske Form maa da ophæves derved, at de endnu
tilbagetrængte Mellembestemmelser komme til Udvikling. Fra dette Synspunkt
danner den organiske Totalitet i Totaliteternes objective Række en
eiendommelig Modsætning til den mechaniske.

Hvad Forholdsleddene i deres Endelighed angaaer, er Forskjellen mellem
mechaniske Sammensætninger og chemiske Forbindelser allerede charakteristisk.
I den mechaniske Sammensætning beholde Leddene deres udvortes endelige
Selvstændighed; hvor noie de end forenes, er Foreningen dog kun en
Sammenhængseenhed: det ene Led falder i Sammenhængen udenfor det andet. Først
i den chemiske Forbindelse bliver
% --- p. 131 ---
\setcounter{footnote}{0}%
det paa reel Maade udtalt, at de existerende Forholdsled virkelig ere Momenter
for hinanden gjensidig. Den Mangfoldighed af Overgange fra existerende Led
(Stof, Produkt, Edukt) til Momenter og fra Momenter til existerende Led, som
bevirkes ved de chemiske Processer, afgiver derfor ogsaa et Mellemtrin af
elementære Dannelser, i hvilke Totalitetsudviklingen gjennemføres med større
Inderlighed og rigere Articulation end i de abstract mechaniske
Kredsbevægelser. Tage vi det jordiske Hele for sig som geologisk Totalitet, da
har den reale Selvhed, Jordudviklingens Princip, jo allerede her begyndt at
bearbeide de elementære Former; her er i denne Bearbeidelse en saa
gjennemgaaende Vexelvirkning af Heelhed og Dele, af almene og særskilte
Totaliteter, at Jordlegemet vel maa kaldes en elementær Organisme.

Den elementære Organisme er imidlertid behæftet med den Modsigelse, at være et
organisk Uorganisk; Organiseringens Ufuldkommenhed sees deri, at Delene endnu
have Charakteer af Brudstykker og udvortes adskilte Masser. De
forskjelligartede Dannelser foregaae overalt paa en saadan Maade, at
Frembringelserne og det Frembragte falde udenfor hinanden; Selvheden i dette
Hele er kun elementær. Skal Modsigelsen løses og det Elementære articuleres,
maa Principet ved nye Gjennembrud anlægge Udviklingen saaledes, at det Hele
kan reflectere sig ideelt i hver enkelt Deel, og hver enkelt Deel være Moment
for det Hele. Dette skeer med den gjennemformede Organisme.

Men Organismen i sin Umiddelbarhed er dog behæftet med den Modsigelse: at være
paa een Gang Andetheden og Selvheden i lige Grad underlagt. Leddenes Forening
er fra den udvortes Side en Sammenhængseenhed, og det Heles Tilværen betinget
af gjennemgaaende, uophørlige Forandringer; seet fra den indvortes Side er
Leddenes Forening derimod en Selveenhed, en articulerende Virken, der magter
Forandringerne og bestemmer det Ydre til en Afspeilen af det Indre. Efter
% --- p. 132 ---
\setcounter{footnote}{0}%
det Ydre ere de enkelte Led og Dele med deres Eiendommelighed i endelig
Modsætning saavel til hinanden indbyrdes som til det Hele; thi det ydre Hele,
der bestaaer i Delenes Sammenhæng, er kun deelviis, d. v. s. tildeels i hver
særegen Deel: Totaliteten er en Andethedstotalitet. Men efter det Indre
reflecteres de enkelte Led i den virksomme Eenhed saaledes, at Heelheden er
udeelt i hver Deel: Totaliteten er en Selvhedstotalitet.

Som totaliseret Eenhed af to ueensartede Totaliteter sætter Organismen en
afgjørende Spænding mellem det Indvortes og det Udvortes, Selvheden og
Andetheden. Ifølge Selveenheden og efter sit Indvortes er Organismen et
Levende; ifølge Sammenhængseenheden og efter sit Udvortes er det Levende en
Organisme. I den abstract levende eller blot vegetative Organisme er
Modsætningen af det Indvortes og Udvortes imidlertid kun antydet ved to Sider
af det samme Levende; men med den besjælede Organisme blive de to Sider til to
Totaliteter, og Spændingen indtræder. Den Bestemthed, hvori de to indbyrdes
modsatte Eenheder, Sammenhængseenheden og Selveenheden, her træffe sammen, er
en Fornemmelse. I Fornemmelsen er Selvheden et umiddelbart Selv, og dette Selv
reflecteret i en umiddelbar Eenhed med og en ligesaa umiddelbar Forskjel fra
sit Andet. Fornemmelsen er den sjælelige Selvheds første eiendommelige
Yttring; her er Principet for den psychologiske Subjectivitet.

Nedsænket i Strommen af de umiddelbare Fornemmelser er den psychologiske
Subjectivitet en dæmrende Bevidsthed; denne dæmrende Bevidsthed er en dyrisk
Sandsebevidsthed. Men den dyriske Sandsebevidsthed er behæftet med den
Modsigelse, at være en ubevidst Bevidsthed, en Bevidsthed af et Selv, der dog
ikke er sig selv bevidst. Denne Modsigelse løses ved et nyt Gjennembrud af
Principet; thi Bevidsthedens Princip er paa Menneskelivets Trin et
Selvbevidsthedsprincip. Først paa dette Trin formaaer den reflecterende
Subjectivitet
% „Strommen“ printed with an unslashed o (verified at 600 dpi against the
% plainly slashed ø of „løses“ and „Først“ two lines below); as printed.
% The sentence begun by „Først paa dette Trin formaaer den reflecterende
% Subjectivitet“ is NOT finished on p. 132: it completes at the TOP of p. 133
% („igjen at sætte Andetheden som gjennemsigtigt Moment, og udvikle Systemer
% af Vidensbestemmelser, der svare til Systemer af Magtbestemmelser.“), and
% one further paragraph of division B follows before the „C.“ display head
% stands on p. 133. Division B therefore ends part-way down p. 133, NOT at
% the foot of p. 132.
% --- p. 133 ---
\setcounter{footnote}{0}%
% Tail of division B standing at the TOP of printed p. 133, above the display head
% "C. / Den begribende Subjectivitet." (which stands about three-fifths down that
% page, not at its top). This text belongs immediately BEFORE the skeleton's
% \subsection*{C.\\ Den begribende Subjectivitet.} line, i.e. at the end of the
% pp. 124--132 marker's slot. It deliberately carries NO page marker of its own:
% the marker for printed p. 133 lives in pp133-146.texfrag, just below the C head,
% so that check.py sees exactly one marker per printed page.
igjen at sætte Andetheden som gjennemsigtigt Moment, og udvikle Systemer
af Vidensbestemmelser, der svare til Systemer af Magtbestemmelser.

Men naar den psychologiske Subjectivitet udvikler sig til en alsidig
totaliserende Reflecteren, falde de objective Magtbestemmelser udenfor de
subjective Vidensbestemmelser; det indlyser, at Reflexionen med de
abstracte Idealitetsformer maa blive formel og afmægtig ligeoverfor det
System af Existenser, hvori Magtmangfoldigheden har sin Realitet. Den
principielle Eenhed af Viden og Magt, hvorpaa den ontologiske
Subjectivitets Oprindelighed beroer, er da paa det psychologiske
Standpunkt forvandlet til en saa endelig Adskillelse, at et System af
Magtbestemmelser uden Viden maa reflectere sig for det reflecterende Selv
i et System af Vidensbestemmelser uden Magt. Reflexionen er gjennemført;
men den herved indtraadte Dualisme medfører nu Problemer, der først
efterhaanden kunne belyses fra et høiere Synspunkt.
\subsection*{C.\\ Den begribende Subjectivitet.}
\addcontentsline{toc}{subsection}{C. Den begribende Subjectivitet}

% NOTE: the display head "C. / Den begribende Subjectivitet." does NOT stand at the
% top of p. 133. It stands about three-fifths of the way down the page, and above it
% p. 133 carries three lines plus one whole paragraph that CLOSE division B. That
% text is delivered separately as .parts/p133-tail-of-B.texfrag and belongs BEFORE
% the skeleton's \subsection*{C.} head. The p. 133 page marker is kept here (so this
% fragment carries one marker per page in its range); if the edition would rather
% have the marker at the true top of the printed page, move these two lines above
% the B-tail.
Allerede paa Umiddelbarhedens Trin stræber den tænkende, psychologisk
reflecterende Subjectivitet i to Retninger: deels at forstaae sin Omverden
med dens Kjendsgjerninger, deels at forstaae sig selv. Men begge
Bestræbelser maae mislykkes, saalænge Standpunktet er eensidig
psychologisk, d. v. s., saalænge det i ontologisk Forstand endnu er
umiddelbart. Den sensualistiske Bevidsthed, der med al sin Tænkning kun
retter sig efter Sandseindtryk, er en i sit Andet fortabt Selvhed, og kan
desaarsag hverken begribe Sandsningen eller det Sandselige. Den i sin
Tanken opgaaende idealistiske Bevidsthed gjør jo vistnok den Opdagelse,
at alt sandseligt Givet kan opløses af Reflexionen; men som en sit Andet
for\-%
% --- p. 134 ---
\setcounter{footnote}{0}%
tærende Selvhed er den psychologiske Subjectivitet uophørlig i Begreb med
at udvikle det Ontologiske, uden dog at have nogen Anelse om, at den med
dette Ontologiske igjen vil møde sig selv i en høiere Subjectivitet. Den
objective Idealisme kan jo vistnok paa en Maade siges at have udviklet sig
med formel Conseqvens af den subjective; men i Realiteten har den gjort et
Spring, idet den i Logiken selv er sprunget af fra sit subjective Grundlag.
Stille vi Hegels „Logik som Videnskab“ imod Fichtes „Videnskabslære“, saa
kan man jo vistnok sige, at begge Værker i Grunden gaae ud paa det Samme,
forsaavidt som de begge gaae ud paa en Analyse og immanent Deduction af
aprioriske Kategorier; og dog hvilken Forskjel! I „Videnskabslæren“ er det
Jeget først og Jeget sidst, hvorom Alting dreier sig; Kategoriernes
Udvikling gaaer her i en saa subjectiv Retning, at Udviklingen selv skal
være Beviset for, at et Objectivt ikke kan gives. I „Logiken som
Videnskab“ er Jeget, den tænkende Subjectivitet, i den Grad blevet
overflødigt, at de objective Kategorier her udvikle hinanden paa objectiv
Maade, og opføre deres logiske Skuespil for hinanden indbyrdes,
ligegyldige ved, om der er nogen Tilskuer eller ikke.

Den objective Logiks Tvetydighed sees blandt Andet deraf, at den kan
gjælde for en almindelig Væsensphilosophie, en Metaphysik i Ordets
strængere Mening, og for en blot Ordphilosophie, en Art dialektisk
Synonymik, eftersom man tager det til. Til en Belysning af denne
Tvetydighed ville vi i nærværende Sammenhæng henvise til dens Behandling
af tre i forrige Afsnit udviklede Kategorier: \emph{Grunden, Substansen og
Causaliteten}.

\emph{Grunden.} Anlægges den logiske Udvikling paa følgende Maade:
„Grunden kan tages 1) abstract eller \emph{adskilt} fra sit Indhold;
saaledes er den \emph{Væsen}; 2) concret, eller som værende i sit Indhold;
saaledes er den \emph{Phænomen}; 3) som den Eenhed, hvori Væsenet har
Phænomen, og Phænomenet
% --- p. 135 ---
\setcounter{footnote}{0}%
har Væsen, men ingen af dem er, undtagen i og ved den anden; saaledes er
den \emph{Nødvendighed}“\footnote{J. L. Heiberg, Ledetraad ved
Forelæsningerne over Philosophiens Philosophie eller den speculative Logik
1831--32, S. 28 flgd. (Prosaiske Skrifter, 1ste Bind S. 183.)}% printed mark: *)
, --- da er man i lige Grad berettiget til at vente saavel en
Væsensforklaring som en Ordforklaring. Begge Dele kunne da ogsaa til en
vis Grad gives i Forening, og det med en Klarhed, Skarphed og Ligelighed,
der gjør Fremstillingen logisk plastisk, uden at det derfor kommer til en
Udvikling af det, der egentlig er Problemets metaphysiske Knude. See vi
paa Ordbetydningen og det Formelle, maae vi jo vistnok indrømme, at flere
Grunde ere Kjendetegnet paa, at man endnu ikke har fundet den eneste
tilstrækkelige, og at, „naar man giver Grunde for Noget, saa har man endnu
ikke angivet Grunden“; men lægge vi særlig Vægt paa Væsensforklaringen,
altsaa paa det i Sagen selv Objective og Ontologiske, hvilket forunderligt
Spring er der da ikke fra et Raisonnement af Grunde, der foregaaer i den
psychologisk subjective Reflexion, til en Begrundelse af det Objective
selv, der kategorisk betegnes ved Ordet \emph{Existens} og ved det i
Existensen Existerende: \emph{Ting}! Vel maa det indrømmes, at „den
tilværende Grund betinger eller er Betingelse for Existensen“; men naar
det for fuld Alvor skal undersøges, hvorledes Grunden da egentlig faaer
Tilværelse, eller, hvorledes det Tilværende bliver begrundet, da er det i
Sandhed ikke nok at lade „den tilværende Grund“ fremkomme som Resultat af
et mellem $+$ og $-$ svævende Raisonnement; ikke nok, at man ved en formel
Omskrivning forvandler den saaledes „tilblevne Grund“ til en
„Betingelse“, og derpaa gaaer videre med en Forklaring \textbf{af} Ordene
\emph{Ting og Egenskab}.
% the "af" above is set in halbfette (fat-face) Fraktur: calibrated at 600 dpi
% against the plain "af" of "Resultat af et mellem" nine lines earlier on the
% page, whose strokes are plainly lighter.
Vil man virkelig gaae ind paa Spørgsmaalet om Tilværelsens ontologiske
Grundforhold, d. e. om Forholdet imellem de existerende Phænomeners
Mangfoldighed og Grundens op\-%
% --- p. 136 ---
\setcounter{footnote}{0}%
rindelige Eenhed, da bliver Grundproblemet, som ovenfor er viist, ikke
blot et objectivt Existensproblem, men et væsentligt
Subjectivitetsproblem.

\emph{Substansen.} Den dialektiske Synonymik kan om Substantialitet og
Substans med Grund anføre f. Ex. Følgende: „Substantialiteten er
Nødvendighedens \emph{Eenhed af Væren}, altsaa den Eenhed, som først
betegnes ved Eet, men dernæst ved Flere, fordi Eet ved sin Repulsion eller
udelukkende Selvstændighed betinger de Flere. Paa samme Maade betinger den
ene Substans \emph{Fleerheden af Substanser}; enhver Aarsag er
\emph{Substans}, og der gives ligesaa mange Substanser, som Aarsager \ldots{}
Da de mange Substanser, hver for sig, ikke ere Andet, med Hensyn paa
Substantialiteten, end hvad den ene Substans er, eller, da den Virkning,
som den ene Substans har paa den anden, bliver, formedelst
Vexelvirkningen, dens Virkning paa sig selv, saa flyder Fleerheden af
Substanser sammen i Identitetsbestemmelsen, ikke af den ene, men af den
eneste Substans. Den eneste Substans kan, fordi den er den eneste, ikke
virke paa nogen anden end sig selv. Dens Virkninger falde altsaa ikke
udenfor den, men ere immanente. I denne Bestemmelse ere Virkningerne dens
\emph{Accidenser} \ldots{} Fremdeles ere Accidenserne ikke Andet end de
mange endelige Substanser, hvilke udgjøre Forskjelsbestemmelsen imod
Identitetsbestemmelsen af den eneste Substans, paa samme Maade som i
Realiteten de mange Materier bleve Former i Modsætning til den eneste
Materie, og som i Totaliteten Mængden af Dele bleve Heelheder udenfor den
eneste Heelhed.“\footnote{J. L. Heiberg. Anfr. Skr. S. 77--78. (Prosaiske
Skrifter, 1ste Bind S. 271--273.)}
% printed mark of the note above: *)
Ved denne Forklaring er her unegtelig ikke blot seet paa den logiske
Ordbetydning, men paa Kategoriernes ontologiske Grundbetydning. Men, naar
vi nu virkelig anlægge Tilværelsens egen substantielle Maalestok ved
Bedømmelsen af denne Forklaring, hvorledes er det da muligt, at
% --- p. 137 ---
\setcounter{footnote}{0}%
„Accidenserne“ kunne være „endelige Substanser“, naar her i Virkeligheden
kun er „en eneste Substans“? Er det kun nogle Accidenser, der kunne blive
endelige Substanser, hvilke mon det da være? Eller er det maaskee alle
Accidenser? Men saa forvandles jo i Endeligheden alle Substanser til
lutter Accidenser, saa at der i Virkeligheden ikke bliver nogen Forskjel
imellem Substanser og Accidenser. Virkeligheden skulde altsaa endnu
virkelig kunne bestaae, efter at Substantialitetsforholdet er opløst?
Eller skulde Substantialitetsforholdet maaskee endnu ikke være opløst,
uagtet den substantielle Modsætning imellem den eneste Substans og de
mange Substanser, imellem Substanserne og deres Accidenser er forsvunden?
Nei, skal man virkelig gjøre sig Rede for Vexelbestemmelsen af den
substantielle Eenhed og den substantielle Mangfoldighed, istedenfor at
løse Substantialitetsforholdet op, ombytte Ordet Substans med et nyt Ord,
og gaae videre i Texten, da vil man, ifølge vort Foregaaende, vel blive
nødt til at indrømme, at Substantialitetsproblemet væsentligt er et
Subjectivitetsproblem.

\emph{Causaliteten.} Som særligt Bidrag til en Charakteristik af Logikens
ontologiske Ordforklaringer eller den dialektiske Synonymik skal i denne
Sammenhæng endnu anføres Følgende med udtrykkeligt Hensyn paa Hegel selv.
I sin „objective Logik“ har Hegel, som bekjendt, indledet Læren om
Virkeligheden med, hvad han kalder „en Udlægning af det
Absolute“.\footnote{Wissenschaft der Logik. 2te Abtheilung. Berlin 183. S.
185 flgd.}
% printed mark of the note above: *). The year stands as "183." in the print, a
% digit short (Hegel's 2te Abtheilung is Berlin 1834); transcribed as printed.
Imod det Absolute staaer nu det Relative; denne Modsætning gjentager sig
i ethvert af „de absolute Forhold“, saavel i Substantialitets- som i
Causalitets-Forholdet. Man maatte altsaa, hvad navnlig Causaliteten
angaaer, være berettiget til at vente en gjennemgaaende Udvikling af
Forholdet imellem Mangfoldigheden af de relative Aarsager og den absolute
% --- p. 138 ---
\setcounter{footnote}{0}%
Aarsagseenhed. Til en saadan Udvikling var her en saameget større
Opfordring, som Kant i den ovenfor omhandlede Antinomie --- og Hegel
behandler jo ellers gjerne, hvor Leilighed gives, de kantiske Antinomier
--- har paa en skarp og bestemt Maade stillet Problemet. Det kantiske
Problem lader sig, hvad Modsætningseenheden af det Absolute og Relative
særlig angaaer, i al Korthed udtrykke saaledes: „Causaliteten kan umulig
være relativ, uden at den tillige maa være absolut; da nu en absolut
Causalitet modsiger sig selv, er Causalitetsbegrebet uden objectiv
Gyldighed.“ Men uden at indlade sig videre med dette Problem, slutter
Hegel af med en skizzeret Fremstilling af den almindelige Vexelvirkning,
og mener saa i Vexelvirkningens Kategorie at naae til en grundig Løsning
af Causalitetens Knuder. Hele Løsningen er summarisk indeholdt i følgende
Forklaring: „Die Kausalität ist bedingt und bedingend; das Bedingende ist
das Passive, aber ebenso sehr ist das Bedingte passiv. Dieß Bedingen oder
die Passivität ist die Negation der Ursache durch sich selbst, indem sie
sich wesentlich zur Wirkung macht, und ebenso dadurch Ursache ist. Die
Wechselwirkung ist daher nur die Kausalität selbst; die Ursache
\emph{hat} nicht nur eine Wirkung, sondern in der Wirkung steht sie als
\emph{Ursache} mit sich selbst in Beziehung. Hierdurch ist die Kausalität
zu ihrem \emph{absoluten Begriffe} zurückgekehrt, und zugleich zum
Begriffe selbst gekommen.“\footnote{Til yderligere Beviis for, at det ikke
er af Uagtsomhed eller Forglemmelse, at det kantiske Problem her er
forbigaaet, men at dets Løsning virkelig skal søges i den almindelige
Vexelvirkning, henvise vi til Hegel's philos. Propædeutik, Berlin 1840 S.
144.}% printed mark: *)

Dersom man nu virkelig tør stille den Fordring til Logiken som Videnskab,
at den udtrykkelig skal orientere os med Hensyn paa Forskjellen imellem de
samme Kategoriers Gyldighed for det Absolute og for det Relative, da er
hiin Løsning af Causalitetens Problem unegtelig meget overfladisk. At de
% --- p. 139 ---
\setcounter{footnote}{0}%
endelige Aarsager maae være Momenter i den uendelige Vexelvirkning, er
forstaaeligt; men hvordan gaaer det saa til, at disse Aarsager, uagtet de
ere Momenter, dog i Virkeligheden vedblive at bestaae som existerende Led,
og derved bevise, at det Relative ikke vil lade sig logisk ophæve? Naar de
relative Aarsager i al deres Relativitet vedblive at existere, hvad bliver
der da af den absolute Causalitet? Mon ikke ogsaa den maa have en
Virkelighedsform, hvori den i al Evighed vedbliver at magte de relative
Aarsager? Disse Spørgsmaal kunne umulig besvares, uden at
Causalitetsproblemet, som vi paa sit Sted have seet, maa blive et
Subjectivitetsproblem.

Skal Logiken som Videnskab derimod kun være Væsensphilosophie, forsaavidt
den i al Almindelighed kan være det i Egenskab af Ordphilosophie,
dialektisk Synonymik, da maa den hegelske Løsning af Causalitetsproblemet
i det Hele ansees for tilfredsstillende. Er det tydeligt indseet, hvilken
logisk Betydning der maa forbindes med Ordene \emph{Aarsag, Virkning} og
\emph{Vexelvirkning}, da er det jo ganske i sin Orden, at man gaaer over
til nye Ord, for i Sammenhæng med det allerede Udviklede at udvikle deres
logiske Betydning.% the letterspacing breaks at the line end: "Aarsag, Virkning"
% is spaced, the following "og" is not, and "Vexelvirkning," is spaced again.
% Transcribed as printed.

„Men“ --- kunde Nogen sige --- „det Problem, hvis Løsning her fra saamange
Sider fordres, nemlig Subjectivitetsproblemet, har da ogsaa, ifølge Hegels
Methode, paa sit Sted fundet sin Løsning, nemlig under første Afsnit i
Læren om \emph{Begrebet}. Overgangen fra \emph{Virkelighed til Begreb},
der udtrykkelig er betegnet ved en Overgang fra \emph{Causalitet til
Subjectivitet}, viser jo tydelig nok, at Hegel har henført, ikke blot
Substantialitets- og Causalitetsproblemet\footnote{Som Beviis for, at Hegel
virkelig har erkjendt Eenheden af Causalitets- og Subjectivitetsproblemet,
og i denne Eenhed Betydningen af den kantiske Antinomie om Frihed og
Nødvendighed, kunde man henvise til hans Wiss. d. Logik 2ter Theil S. 214.
Men et Blik paa S. 215 nederst vil være tilstrækkeligt til at vise, hvor
lidt herved er vundet.}% printed mark: *)
, men paa
% --- p. 140 ---
\setcounter{footnote}{0}%
Udviklingens fremadskridende Strøm alle Værens- og Væsensproblemer under
Subjectivitetsproblemet. „Begrebet“, siger han, „er Sandheden af det
substantielle Forhold, hvori Væren og Væsen have opnaaet deres fyldige
Selvstændighed og Bestemmelse ved hinanden gjensidig“, men Begrebet
begynder jo netop som Subjectivitet. --- At Hegel maa forstaaes og
virkelig ogsaa er af Sagkyndige bleven forstaaet saaledes, kan let bevises.
„Som Substantialitetens Overgang“ --- siger Heiberg\footnote{Anfr. Skr. S.
80. (Prosaiske Skrifter, 1ste Bind S. 276--277.)}
% printed mark of the note above: *)
--- „er Begrebet \emph{Subjectivitet}, som Nødvendighedens er det
\emph{Frihed}, og som hele Grundens eller Reflexionens (da disse ere
Vorden af Begrebet) er det \emph{Tilværen} af den abstracte Tanke. Med
Begrebet er Subjectiviteten umiddelbart given. Begrebets første
umiddelbare Standpunkt er derfor \emph{Subjectet}. Og heri ligger
ligeledes \emph{Friheden}; thi da Substansen er Aarsag til sig selv, saa
mangler den ikke Andet til Friheden end Subjectivitet. Ved denne hæves
Substansen fra Nødvendighed til Frihed; thi Subjectet er det, som har
udviklet sig af Nødvendigheden, som sin Forudsætning, altsaa det Frie.
Subjectiviteten udtrykkes ved Jeg, men dette er endnu ikke Bevidsthed; det
er kun den for Naturens og Aandens Skikkelser fælleds (altsaa logiske)
Subjectivitet, og saaledes den Rod, hvoraf Bevidstheden udvikler sig,
skjøndt endnu ikke denne.“% the single closing mark here has to serve both the
% Heiberg quotation and the objector's „det Problem (opened on p. 139), which the
% print never closes separately. Transcribed as printed; hence a „ surplus of one.

Disse og lignende Oplysninger bevise just det Modsatte af, hvad de
tilsigte. Er der noget Afsnit i den hele hegelske Philosophie, der viser,
i hvilken Grad Subjectivitetsbegrebet her er bleven mishandlet og
forqvaklet, da er det netop det af Logiken anførte Afsnit om
Subjectiviteten. Hvor tvungen og struet den hele Opfattelse \textbf{er},
fremlyser umiddelbart af Yttringer som denne: „Die Gestalt des
\emph{unmittelbaren} Begriffes macht den Standpunkt aus, aus welchem der
Begriff ein subjectives Denken, eine der Sache äußerliche Reflexion ist.
% "er" set in halbfette (fat-face) Fraktur, calibrated at 600 dpi against the
% plain "er det Tilværen" eighteen lines above.
% --- p. 141 ---
\setcounter{footnote}{0}%
Diese Stufe macht daher die \emph{Subjectivität} oder den \emph{for\-%
mellen Begriff} aus.“ Det subjective Begreb, den logiske Subjectivitet
fremkommer altsaa, efter denne Forklaring, derved, at den objective
Tankebevægelse pludselig lægges ind i det psychologisk subjective Medium.
Den immanente Methode, som paa Værens og Væsenets Trin har kunnet undvære
alt Subjectivt, endogsaa det ontologisk Subjective, faaer paa Begrebets
Høidepunkt pludselig en saadan Trang til en Subjectivitets Bistand, at
den, som i en snever Vending, endog tager til Takke med den psychologiske
Subjectivitet.

Undseelig over at skulle behøve en slig Nødhjælp søger den immanente
Methode naturligviis paa bedste Maade at give Begrebets Subjectivitet et
mere objectivt Anstrøg. Subjectiviteten skal da egentlig ikke hidrøre
derfra, at dets „Skikkelse“ just „udgjør det Standpunkt, hvorfra Begrebet
er en subjectiv Tænken, en for Sagen udvortes Reflexion“; det skal med
andre Ord ikke være Standpunktet i det Subjective, men Skikkelsen i det
Objective, hvorpaa Eftertrykket falder. Det subjective Begreb er altsaa i
strengeste Forstand ikke subjectivt, thi det er objectivt; men hvoraf
kommer det? Naturligviis deraf, at der, reentud sagt, ingen væsentlig
Forskjel er imellem det Subjective og det Objective. „Saavel i Naturens
ubevidste, som i Aandens selvbevidste Skikkelser“ --- hedder
det\footnote{Heiberg. Anfr. Skr. S. 80. (Prosaiske Skrifter, 1ste Bind S.
277.)}
% printed mark of the note above: *)
--- „har Causaliteten bestemt sig som Substantialitet og Spontaneitet;
heri ligger Subjectiviteten, og enhver Substans er et Jeg, hvad enten
dette Jeg er kommet til Bevidsthed om sig selv eller ikke. Bestemmelsen af
Jeg ligger allerede i Reactionen; thi at enhver Substans reagerer, skeer
netop i Kraft af dens Subjectivitet. (Et Bord, en Stol, et Skab o. s. v.
ere alle af Træ; men i hine tilfældige Former, som en Snedker har givet
Træet, ligger ikke nogen Subjectivitet eller noget Jeg; thi det er ikke
som Stol, Bord o. s. v., at de
% --- p. 142 ---
\setcounter{footnote}{0}%
reagere, men som Træ, altsaa ved deres Substans; og Jeget af denne
Substans viser sig i Træets naturlige, organiske Form, ligesom allerede
den uorganiske Verdens Jeg viser sig i Krystallisationen).“

Ved at tillægge den endelige Substans (Træet, Krystallen) i dens
umiddelbare og materielle Existens en Subjectivitet, et Jeg, har den
objective Logik fra Roden af forqvaklet Modsætningen mellem Selvhed og
Andethed, mellem det Subjective og det Objective; thi er der Noget, der i
Virkeligheden maa siges at existere i Andethedens Form, at udtrykke det
Selvløse og derved i sig Objective, da maa det dog vel være Stofsubstanser
som Krystallerne og Træet i Stol og Bord. At paavise „den for Naturens og
Aandens Skikkelser fælleds (altsaa logiske) Subjectivitet“ kan dog umulig
være det Samme som at udviske den logiske Forskjel mellem Subjectiviteten
og det Objective saavel i „Naturens som i Aandens Skikkelser“. Men at det
Sidste netop skeer i den hegelske Logik, kan sees paa følgende Maade. Hvad
er en Krystal: Subject eller Object? Fra Hegels Standpunkt maa man svare:
„i Begrebets første Afsnit er Krystallen Subject, i Begrebets andet Afsnit
er den derimod Object“. Men nu er der dog vel ingen Hegelianer, der vil
paastaae, at Krystallen som Subject først maa gjennemgaae det Cursus af
subjectiv Logik om Domme og Slutninger, der afhandles i første Afsnit, før
den i andet Afsnit kan indlemmes i Rækken af de virkelige Naturgjenstande
og gjælde for et ordentligt Object. Her er, hvad Krystallen selv angaaer,
ikke Skygge af Begrebsbevægelse fra Subject til Object; hvorimod der i vor
Opfattelse af, hvad en Substans som Krystallen er, ganske rigtig er en
Begrebsbevægelse, og det en sand Begrebsbevægelse: fra at opfatte
Krystallen under Synspunktet af et logisk grammatisk Subject til at
opfatte den under Synspunktet af et virkeligt, i Naturen forekommende
Object. Men naar Begrebsbevægelsen, hvordan man end vender og dreier det,
altid falder og maa falde --- ikke i
% --- p. 143 ---
\setcounter{footnote}{0}%
Objectet som Object, --- men i den tænkende Subjectivitet, saa maa det jo
nødvendig blive Logikens Opgave at paavise, hvorledes det i den tænkende
Subjectivitet Ontologiske og Psychologiske bestemme hinanden gjensidig;
med andre Ord: at bringe Subjectiviteten til logisk alsidig
Selvforstaaelse ved at fremstille den som \emph{begribende Subjectivitet}.

Saa afgjørende det logisk objective Indhold nu vistnok er for
Begrebsformernes eiendommelige Udvikling, saa tør vi dog, med
Subjectivitetsproblemet for Øie, ikke paa noget Punkt skille Begrebet fra
den begribende Virksomhed; de for Udviklingen charakteristiske Synspunkter
betegne vi derfor ved: \emph{Begriben og Begreb, Begribningens Phaser,
Begriben og Viden}.

\subsubsection*{a) Begriben og Begreb: logisk Ideebevægelse.}
% The book's Indhold (p. 4) gives C's first limb as "a) det Logiske og det
% Antilogiske"; the printed head here reads as set above. Printed head kept in the
% body, per the convention. The label "a)" is Fraktur a; the α) of p. 146 below is
% Greek.
Der gives ingen Tanke, som ikke er fremkommen ved en Tænken, ingen
Reflexion, som ikke er fremkommen ved en Reflecteren, intet Begreb, som
ikke er fremkommet ved en Begriben; hører Tænkningen op, saa maa Tanken
forsvinde, hører Begribningen op, saa maa Begrebet forsvinde: Begrebet er
ligesaa uadskilleligt fra Begribningen, som Tanken fra Tænkningen. Hvad
der ikke kan tænke sig selv, har heller ikke Tanken i sig selv; hvad der
ikke kan begribe sig selv, har heller ikke Begrebet i sig selv. Vistnok
kunne baade Tanker, Reflexioner og Begreber opbevares f. Ex. i Skrifter,
uden at Skrifterne læses; men Tankerne i et Skrift ere kun til efter
Mulighed, ikke efter Virkelighed; denne Tilværen efter Mulighed er
udtrykkelig betegnet derved, at Skriftet, som Tankernes Udtryk, er et
logisk Mellemværende mellem Forfatteren og Læseren.

Hvorledes Tænkningen i Almindelighed forholder sig til Reflexionen, er
allerede viist; men hvorledes forholder nu Reflexionen sig til den
egentlige Begriben? En vis Reflecteren er vel mulig uden Begriben, men
ingen Begriben er mulig
% --- p. 144 ---
\setcounter{footnote}{0}%
uden en Reflecteren; thi Reflexionen er Begrebets Vorden, og som Begrebets
Vorden selv en vordende Begriben. Da nu Begribningen paa ethvert Trin maa
rette sig efter det, der skal begribes, saa indsees heraf, at der til
ethvert Trin, som den reflecterende Subjectivitet gjennemløber, maa svare
en egen, om end eensidig, Begrebsudvikling. Efterhaanden som Reflexionen
bliver mere totaliserende, bliver ogsaa Begribningen mere alsidig; først
med den gjennemførte Reflecteren begynder den egentlige Begriben. Og dog
er den egentlige Begriben væsentlig forskjellig fra Reflexionen, om denne
end gjennemføres nok saa alsidig. Reflexionen culminerer i Momenternes
Gjennemsigtighed; selv den meest gjennemførte Reflexion naaer aldrig ud
over Systemet af de rene Reflexer. For den rene Reflexion bliver Alt
logisk, Indholdet ligesaa vel som Formen. Overalt, hvor det blot kommer an
paa en formel Begriben, er det nok at at gjennemføre Reflexionen;
% "at at" is the print's own reading, verified at 600 dpi: a dittography of the
% compositor's. Transcribed as printed.
men naar det gjælder om Realitet, saa viser Modsætningen sig imellem den
blotte Reflecteren og den egentlige Begriben. For at adskille Reflexionens
Sphære fra Begrebets, gjør den hegelske Logik, og det med Rette,
opmærksom paa Forskjellen imellem Kategorierne, de eensidige Begreber, og
Begrebet selv, Kategoriernes Totalitet. Enhver af Reflexionens Kategorier
for sig er kun en Side af Begrebet, men i Begrebets egen Sphære er enhver
af Begrebets Sider selv det hele Begreb. Men med rene logiske Totaliteter
kan Logiken aldrig naae ud over Reflexionen; og naar en Logik ikke kan
naae ud over Reflexionen og derved ud over det blot Logiske, saa kan den
hele tilsvarende Philosophie heller ikke naae ud over det blot Logiske,
men bliver en Panlogisme. Den hegelske Philosophie er Panlogisme, ikke
fordi den holder paa „Begrebet“, men fordi den endnu ikke er kommen til
den egentlige Begriben.

I den egentlige Begriben er der nemlig et Noget, der danner en afgjørende
Modsætning til Reflexionen og derved til det blot Logiske; hvad
\textbf{er er} vel dette for et Noget? Det
% "er er": the word is doubled in the print, and both sorts are halbfette
% (fat-face), calibrated at 600 dpi against the plain "er der nemlig" two lines
% above. Transcribed as printed.
% --- p. 145 ---
\setcounter{footnote}{0}%
Ulogiske! Nei, det Ulogiske kan umulig danne Modsætning til Reflexionen og
det Logiske, da det jo ligefrem maa forkastes af den logiske Reflexion.
Det Umiddelbare! Vistnok danner det Umiddelbare, som vi i det Foregaaende
have seet, en Modsætning til Reflexionen; men denne Modsætning er fra et
logisk Synspunkt saa bevægelig, at den just falder for Reflexionen,
eftersom det jo netop er Reflexionen selv, der baade sætter og hæver den.
Det Reale! Men naar det skal indsees, hvad det Reale er, saa kommer
Reflexionen igjen og forvandler os dette Reale til en logisk Reflex i
Realitetens Kategorie; og saa ere vi lige nær. For da at gjøre
Modsætningen til Reflexionen og det Logiske saa stærk, at ingen
Reflecteren skal kunne opløse os den og forvandle den til en blot logisk
Reflex, uagtet Modsætningen selv naturligviis maa vedblive at falde under
Reflexionen og det Logiske, betegne vi den negativt ved Kategorien
\emph{det Antilogiske}. Først idet den logiske Reflexion reflecterer det
Antilogiske, uden at kunne forvandle det til en blot logisk Reflex,
indtræder den egentlige Begriben, og først med den egentlige Begriben
kommer Begrebet.

For den gjennemførte Reflexion er der altid en Modsætning mellem
Totalitetens Form og dens Indhold; men da Modsætningen er logisk, saa er
Identiteten det ogsaa; er Formen logisk, saa er Indholdet det ogsaa, og
omvendt. At den mechaniske Form bestemmes som logisk, vil da paa dette
Standpunkt sige, at Indholdet, det mechaniske Hele, selv bestemmes som
logisk; det Samme gjælder om de halvorganiske og organiske Totaliteter: de
kunne tilhobe gjennemreflecteres saaledes, at Indholdets Speiling af
Formen og Formens af Indholdet gjør dem begge til logiske Reflexer. Med
den egentlige Begriben indtræder derimod den Vending, at Totalitetens
Indhold er antilogisk, naar Formen er logisk, og Indholdet logisk, naar
Formen er antilogisk. At begribe det Logiske ved det Logiske er formelt;
at begribe det Antilogiske
% signature mark at the foot of this page: „Grundideernes Logik.“ + 10 --- not
% transcribed.
% --- p. 146 ---
\setcounter{footnote}{0}%
ved det Antilogiske er umuligt; at begribe det Logiske ved det
Antilogiske, og omvendt, er reelt. Naar det Logiske begribes ved det
Antilogiske, faaes Begrebet i logisk Form, og alle Bestemmelserne ere
Negationsbestemmelser; naar det Antilogiske begribes ved det Logiske,
faaes Begrebet i antilogisk Form, og alle Bestemmelserne ere
Positionsbestemmelser. Kun i Begrebets Idee, de logiske og antilogiske
Bestemmelsers concrete Eenhed, kan Begribningen fuldendes.

\subsubsection*{α) Negationsbestemmelser: Begrebet i logisk Form.}
Den rene Tænken er negativ, og enhver Negation en subjectiv Act. At der
uafhængigt af Subjectiviteten skulde gives objective Negationer, at en
existerende Ting, et virkeligt Object, virkelig skulde kunne negere enten
sig selv eller noget Andet, er, naar man ret besinder sig, en ligesaa
urimelig Phantasie, som, at en Dør skulde være istand til at negte enten
sig selv eller Husets Beboere hjemme. „Negationerne“, siger Herbart, „ere
kun i dens Forestilling, der i Gjenstanden søger Noget, han ikke finder;
Manglen selv er Intet.“ Ja, saa subjectiv er Negationen, at den ikke blot
forudsætter en Bevidsthed, men en tænkende Bevidsthed, der kan bruge
Sproget: Negationen har kun Tilværelse i Ordet. Har end den subjective
Logik havt Uret imod Hegel, naar den ikke vilde indsee, at dens
Contradictionsprincip, ligesom enhver anden formel Grundsætning, trængte
til en Begrundelse, saa har den dog, i en vis Forstand, havt Ret i at
holde paa det i al Tænken, Reflecteren og Begriben Subjective. Hvorledes
den subjective Logik imidlertid fra sin Side dog ogsaa misforstaaer det i
Subjectiviteten Logiske, forsaavidt den i al Bornerthed fastholder sit
psychologiske Standpunkt, er allerede tilstrækkelig omhandlet; her skal
kun tilføies, at dens Forblindelse i saa Henseende snarere maatte
forstørres end formindskes ved Polemiken mod den immanente Methodes
mystiske Objectivitet.

% --- p. 147 ---
\setcounter{footnote}{0}%
Den hegelske Lære om Negationen og det Negative er i
Philosophiens Historie et mærkeligt, maaskee et af de allermærkeligste
Exempler paa, hvor forunderlig en stor videnskabelig Opdagelse kan
falde sammen med en stor videnskabelig Misforstaaelse. Det er Hegels
udødelige Fortjeneste at have opdaget „Negationens Negation“, „den
immanente Negation“, „den uendelige Negativitet“%
\footnote{% printed mark: *)
„Man hører ikke sjældent den Meningsyttring, at den negative Viden
ikke tilfredsstiller, at Philosophien bør være positiv. Det kommer nu
an paa, hvorledes man forstaaer det Positive. Opfattes det Positive
saa umiddelbart, at det har det Negative udenfor sig, opfattes det i
endelig Form, og er da i denne sin Endelighed hjemfalden til
Modsigelsen, og derved til det Negative. En Philosophie, der er
dogmatisk positiv ved at holde det Negative ude, er ideeløs, og den
ideeløse Positivisme just af Ideen bestemt til Opløsning i det
Negative. Derfor bliver Dogmatismen altid angreben af Skepticismen.
Dogmatismen er ideeløs, fordi den mangler det Negative. Skepticismen
er ideeløs, fordi den mangler det Positive. Det Positive er i Ideen
uendelig negativt, og dette Negative just derfor det sande Positive.
Den Indvending, at en Philosophie ikke er positiv nok, er da i
Begrebet at forstaae, at den ikke er negativ nok. Skepticismen er i
samme Grad, som den har det Negative, nærmere ved at naae det sande
Positive, end Dogmatismen.“
% the bibliographical reference is set in halbfette (fat-face) Fraktur,
% calibrated against the note's own text on the same page
\textbf{Phil. Propæd. i Grundtræk. Kjøbenhavn 1857, S. 152--153.}}, og
dog har den immanente Negation afstedkommet megen logisk Forvirring.

Pointen er, at „Negationen er immanent“. Dersom Negationen ikke
opfattes som „immanent“, saa er „Negationens Negation“ et tomt
Ordspil. Saaledes maa Contradictionssætningen i den formelle Logik:
$A$ er ikke Ikke-$A$, nødvendigviis blive komisk, naar den skal
anvendes paa bestemte Tilfælde eller, med andre Ord, tjene til at
udvikle noget Bestemt. Hvad er et Speil? „En Negation af et Bord, thi
et Speil er ikke et Bord, altsaa et Ikke-Bord.“ Og hvad er saa en
Elephant? „En Negation af et Speil; altsaa en Elephantnegation af en
Speilnegation af et Ikke-Speil.“ Ligesom
% --- p. 148 ---
\setcounter{footnote}{0}%
man i et Kaleidoskop kan af de samme Glasstumper udbringe de
allerforskjelligste Combinationer, saaledes kan man ved Hjælp af de
opstillede Negationer drive det meest vilkaarlige Spil med Negationens
Negation. For at undgaae Vilkaarligheder, klamrer den formelle Logik
sig da saa fast til Contradictionssætningens dobbelte Negation, at det
ved Negationens Negation bebudede Fremskridt kommer i en latterlig
Strid med Tautologien: $A$ er ikke, hvad $A$ ikke er; den hele
Tankebevægelse er fingeret.

% the compositor letterspaces only "være" at the end of the line and only
% "objectiv" in the next; "immanent" and the second "være" stand unspaced.
% Transcribed as printed.
„Men“, saa sluttede Hegel, „naar Negationen skal \emph{være}
immanent, saa maa den nødvendig være \emph{objectiv}; det Immanente er
altsaa det Objective, og dette Objective uafhængigt af
Subjectiviteten.“ Her er Misforstaaelsen, en uhyre Misforstaaelse!
Misforstaaelsen hidrører fra Tvetydigheden i det Objective. Det gaaer
med Objectet og det Objective, som med Begreb og Tanke: Objectet er i
Virkeligheden kun til derved, at det objectiveres: tænker man al
Objectiveren bort, saa er der ingen Objecter og intet Objectivt. Men
gjælder nu ogsaa det Omvendte, at, dersom man tænker Objecterne bort,
saa er der ingen Objectiveren? Ja, saalænge der udtrykkelig
reflecteres paa Virkeligheden, den virkelige Existeren; men saa
overgribende er Subjectivitetens Magt, at den kan objectivere sig ikke
blot det Existerende, men tillige det Ikke-Existerende, og derhos
endog objectivere sig denne Objectiveren. Dersom der ikke var en
Objecterne objectiverende Subjectivitet, og det en Subjectivitet, der
tillige kunde objectivere sin egen Objectiveren, vilde Negationen og
det Negative aldrig være blevet til. Negationen er ikke blot af
subjectiv Oprindelse, men det reneste, abstracteste Udtryk for
Subjectivitetens Frigjørelse fra alt Andet, ja fra sin egen
Umiddelbarhed. Thi der er intet Værende saa positivt, ingen Sandhed
saa slaaende, at jo den oversvævende Subjectivitet kan, ved at
anbringe sin Negation, berede den en formel Tilintetgjørelse. Men just
idet den subjective Frihed naaer sin yderste Spidse i den
% --- p. 149 ---
\setcounter{footnote}{0}%
vilkaarlige Negation, maa Subjectiviteten, saaret af Selvmodsigelsen,
reagere mod Negationens Vilkaarlighed ved en ny Negation; saaledes
bliver da Negationen immanent. Den immanente Negation er objectiv; men
her er just det Vendepunkt, hvor det Objectives Tvetydighed maa falde
i Øinene.

Det virkelige Object objectiveres derved, at det Objective bestemmer
det Subjective og bestemmes ved det Subjective. Reflecteres nu paa, at
det er det Objective, der bestemmer det Subjective, da er Objectets
Væren antilogisk; reflecterer man derimod paa, at det er det
Subjective, der bestemmer det Objective, da er Objectets Væren logisk.
Det Objective er altsaa baade logisk og antilogisk; men kun det
Antilogiske kan stilles selvstændigt mod Subjectiviteten: det Logiske,
hvor objectivt det end er, tilhører Subjectiviteten, ganske og aldeles
den tænkende, begribende Subjectivitet. Naar Hegel nu lægger Vægt paa,
at den immanente Negation er objectiv, da give vi ham gjerne Ret; men
naar han saa lader, som om det faldt af sig selv, at det logisk
Objective maatte være ligesaa ligegyldigt mod den tænkende
Subjectivitet, som de virkelige Gjenstande mod en tilfældig Tilskuer,
da maae vi heri see en for Logik og for al sund Philosophie yderst
farlig Misforstaaelse.

Hvorledes den immanente Negation forholder sig til
Vexelbestemmelsen af det Logiske og det Antilogiske, kan sees ved et
Blik paa den hegelske Logiks Begyndelsesbegreb: Væren. Vi have i
Sætningen „Væren er“ en umiddelbar Eenhed af det Logiske og det
Antilogiske. Det Antilogiske er underforstaaet, forsaavidt her tænkes
paa en Væren af et i Virkeligheden Existerende, et Detteværende. Det
Existerende i dets Enkelthed og Punktualitet byder den tænkende
Subjectivitet Spidsen, og modstaaer det Logiske; men dette Antilogiske
bliver igjen logisk bestemt derved, at Subjectiviteten opfatter og
fastholder det i en bestemt Kategorie. I Kategorien Detteværende
findes igjen Kategorien et Værende, og i Kategorien
% --- p. 150 ---
\setcounter{footnote}{0}%
et Værende atter Kategorien Væren. Naar nu Væren, saaledes afsondret
ved Abstractionen og frigjort fra alle antilogiske Biforestillinger,
betragtes for sig, da er denne rene Væren, seet i antilogisk
Belysning, ganske rigtig det Samme som et Ikke-Værende, et Intet. I
den saameget omtvistede Sætning: „Væren er Intet“ giver den hegelske
Logik os altsaa den første Prøve paa immanent Negation. I hvilken
Betydning denne Negation er objectiv, fremgaaer af følgende
Overveielse. Reflecterede vi paa det Logiske alene, maatte Kategorien
Væren nødvendigviis vedblive at være, hvad den er, en Bestemmelse ved
Kategorien det Værende, altsaa kun i meget uegentlig Forstand et
Intet. Reflecterede vi paa det Antilogiske alene, maatte Bestemmelsen
Væren gaae umiddelbart op i sit Detteværende, og her vilde end ikke
fremkomme Skygge af Negation. Kun idet her fra den ene Side sees paa
den yderste Abstraction af en logisk Væren og fra den anden paa den
antilogiske Værens bestemte Fordring, er Modsætningen mellem Væren og
Intet dreven til sit Yderste. Med Hensyn paa det Antilogiske er Intet
et absolut Ikkeværende; Intet har slet ikke med det Værende at gjøre;
et værende Intet er ligesaa meningsløst, som en
% printed "gron", not "grøn": the o carries no slash and matches the plain o
% of "som" on the same line, while the ø of "meningsløst" on the same line
% shows its bar clearly. Transcribed as printed.
gron Tone; fra dette Synspunkt gjælder det: af Intet kommer Intet. Med
Hensyn paa det Logiske er Intet i sin objective Ikke-Væren et negativt
Udtryk for Væren selv; den samme subjective Tænken, der bærer Væren,
den yderste Abstraction i positiv Form, formedelst den rene Anskuen,
bærer ogsaa Intet, den yderste Abstraction i negativ Form, formedelst
den ophævede Anskuen; Abstractionens negative Side har altsaa en
ligesaa objectiv Væren som den positive. I denne Forstand gjælder det,
at Intet er Væren, efterdi Væren er Intet.

Det er da uimodsigeligt, at den Vorden, hvorom den hegelske Logik
taler, naar den lader den fortsatte Afvexling af Væren og Intet
ophæve sig selv i sit tredie Led: 1) Væren $= A$, 2) Intet $=$
% the first hyphen of "Ikke-Ikke-A" is only partly inked; read as a hyphen
% on the evidence of the ink present at the hyphen site (rejected: "Ikke Ikke-A")
Ikke-$A$, 3) Væren $=$ Ikke-Ikke-$A$,
% --- p. 151 ---
\setcounter{footnote}{0}%
er en Bevægelse, der vel har objectiv Gyldighed, men dog foregaaer i
Subjectiviteten; den gjælder som en logisk Bestemmelse for det
Antilogiske, men er derfor ikke selv antilogisk.

I nøie Sammenhæng med Hegels Misforstaaelse af den immanente
Negations Objectivitet staaer hans gjennemgribende Misforstaaelse af
det objective Forhold mellem det \emph{Almene} og det \emph{Enkelte}.
Hvad er det Almene i sit Væsen, subjectivt eller objectivt? Af Hegels
Logik er det umuligt at faae noget bestemt Svar paa dette
Spørgsmaal. „Det rene Begreb“, hedder det%
\footnote{% printed mark: *)
Hegel, Wissenschaft der Logik, 2ter Theil, Berlin 1834, S. 36 flgd.},
% "Væsenet", "Væren," and "Begrebet" are letterspaced; the second
% "Væsenet" (broken "Væ-senet" across the line) is not. As printed.
„er det absolut Uendelige, Ubetingede og Frie. \emph{Væsenet} har sin
Vorden fra \emph{Væren}, \emph{Begrebet} fra Væsenet, altsaa fra
Væren. Men denne Begrebets Vorden har Betydning af et Modstød i det
selv, saa at det ved Vorden Fremkomne netop er det Ubetingede og
Oprindelige. Væren er i sin Overgang til Væsen blevet til et Skin,
dens Vorden eller Overgang i Andet til en Sætten; omvendt har denne
Sætten eller Væsenets Reflexion igjen ophævet sig selv, og er vendt
tilbage til en oprindelig Væren. Begrebet er disse Momenters
Gjennemtrængen af hinanden saaledes, at det Qvalitative og oprindelig
Værende kun er en Reflexion i sig selv, og denne Reflexion i sig selv
den Vorden af Andet, der udtrykker Bestemtheden: denne sig selv
bestemmende Bestemthed er uendelig. Begrebet er altsaa for det Første
den \emph{absolute Identitet med sig}, men saaledes, at denne
Identitet udtrykkelig er Negationens Negation, Negativitetens
uendelige Eenhed med sig selv. Denne Begrebets
% printed "Omboining" without ø: the o shows no slash, unlike the ø's
% elsewhere on this page. Transcribed as printed.
Omboining i sig selv og Beroen paa sig selv formedelst Negativiteten
er Begrebets \emph{Almeenhed}. Det er netop det \emph{Almenes} Natur,
at være et saadant Simplex, hvis eenfolde Væren, i Kraft af den
absolute Negativitet, indeholder den høieste Forskjel og Bestemthed.
Værens Eenfoldhed er dens Umiddelbarhed; man kan antyde den ved en
Mening, men ikke bestemt sige, hvad
% --- p. 152 ---
\setcounter{footnote}{0}%
den er; den er derfor umiddelbart Eet med sit Andet, med Ikke-Væren.
Just dette er dens Begreb: at være et Simplex, der umiddelbart
forsvinder i sit Modsatte: den er \emph{Vorden}. Det \emph{Almene}
derimod er det \emph{Eenfolde}, der dog ligesaa meget er det i sig
selv Rigeste, fordi det er Begrebet.“

For at forstaae den hegelske Bevægelse fra Væren gjennem Væsen til
Begreb, er det tilstrækkeligt at forstaae Bevægelsen i den allerførste
Trilogie: \emph{Væren, Intet, bestemt Væren} eller \emph{Tilværen}. I
denne Trilogie forekommer Umiddelbarheden \textbf{to} Gange, den
immanente Negation \textbf{to} Gange, den i Negativiteten synlige
Almeenhed eller det udviklede Begreb \textbf{een} Gang. Saaledes som
det Objective er i denne Trilogie, saaledes er det i den hele Logik;
og det Samme gjælder om det Subjective. Som det Almene er i denne
Trilogie, saaledes er det i den hele Logik; og det Samme gjælder om
det Enkelte. Reflectere vi nu paa Umiddelbarheden, da indlyser det,
at det Almenes Objectivitet i tredie Led \textbf{er} ligesaa langt fra
det Antilogiske, som i første Led. Heraf sees, at den immanente
Bevægelse fra abstractere til concretere Begreber i logisk Form
\textbf{ikke} udviser nogen Tilnærmelse til Begrebsbestemmelser i
antilogisk Form. Udtrykker nu det reelt Existerende sit
Begrebsindhold i antilogisk Form, saa er det jo herved beviist, at
Begrebet i den hegelske Logik \textbf{er} paa sit sidste og høieste
Udviklingstrin ligesaa langt borte fra det reelt Existerende, som i
sin første abstracte Begyndelse. Reflectere vi dernæst paa det
Negative, da indlyser det, at, naar den første immanente Negation:
Ikke-Væren $=$ Intet, er, som vi have seet, en reen og immanent
Subjectivitetsbestemmelse, saa maa det Samme gjælde om den
Negationens Negation, hvorved det tredie Led i høiere og rigere
Umiddelbarhed gjengiver det første. Er nu Negationen selv med samt
Negationens Negation i den første Trilogie en immanent
Subjectivitetsbestemmelse, saa gjælder det Samme om alle andre
Negationer og Negatio\-%
% --- p. 153 ---
\setcounter{footnote}{0}%
ners Negationer paa de følgende Trin lige til Logikens Slutning.

Det er en reen Indbildning, naar man ad denne Vei mener at komme det
i Objectiviteten Objective nærmere ved Hjælp af Begreber som Grund og
Phænomen, Ting og Egenskaber, Aarsag og Virkning. Den Negation,
hvorved Aarsagen omsættes i sin Virken, er i logisk objectiv Væsenhed
ligesaa subjectiv, som den Negation, hvorved Væren ombyttes med Intet;
den Negationens Negation, hvorved Aarsagen gjennem sin Virkning
restitueres til Eenhed med sig selv som Vexelvirkning, befinder sig i
en ligesaa uendelig Afstand fra det objectivt Virkelige, som den ved
Negationens Negation bestemte Tilværen fra det objectivt Tilværende.
% PRINTER'S DEFECT: the print has "og!" where the sense requires "og".
% The mark is an unmistakable exclamation sort (stem + dot below).
% Transcribed as printed.
Det er Altsammen i Begrebet subjectivt og! dog Altsammen meent
objectivt. Her er i den objective Logik en naiv Collision imellem
Begreb og Mening.

„\emph{Det Enkelte}“ --- siger Hegel --- „er det Almene.“ Denne
Sætning er, hvad vi senere skulle oplyse fra et andet Synspunkt,
aldeles falsk, dersom det Enkelte skal opfattes i umiddelbar Existens;
thi i umiddelbar Existens er det Enkelte antilogisk, medens det Almene
i sine Negationer altid er og forbliver logisk. Ville vi derimod ikke
fordre det umiddelbare Enkelte, men holde os til dets logiske Reflex,
da kan Sætningen: det Enkelte er det Almene, med fuld Ret gjælde for
et Begrebets Udtryk. For imidlertid ikke at forqvakle et saadant
Begrebsudtryk, tør vi dog heller ikke i Copula „\emph{er}“ see en
altfor umiddelbar „Tilbagevenden til den oprindelige Væren“. Her
gjælder jo ikke blot endnu, hvad Dommen i det Hele angaaer, den
Reflexionsmodsætning, at Identiteten af det Enkelte i Subjectet og det
Almene i Prædicatet maa have sin væsentlige Forskjel, men her er
tillige, hvad det \emph{Enkelte} særlig angaaer, en for Begrebet selv
charakteristisk Modsætning mellem det Antilogiske og det Logiske, en
Modsætning, hvis Eenhed just maa gribes i den egentlige Begriben.
Totali\-%
% --- p. 154 ---
\setcounter{footnote}{0}%
teten af Enkeltbestemmelserne i det existerende Subject vil fra den
logiske Side reflectere sig i Totaliteten af det fuldstændige
Prædicats Almeenbestemmelser, medens Enkeltbestemmelserne selv fra den
antilogiske Side ville modsætte sig Reflexionen: \emph{saaledes gribes
og begribes et antilogisk Indhold i logisk Form.}

Skal det Antilogiske derimod ganske forsvinde, og det Enkelte gaae
fuldstændig op i det Almene, da maae vi ikke med Hegel tænke paa de
grammatiske Subjecter, hvori Begrebet har „umiddelbar Tilværen som et
Noget eller en Ting“; her maa da vælges en høiere Form. Thi der er
virkelig en Form, hvori det Enkelte er identisk med det Almene, og det
er den uendelige Negativitet, Subjectivitetens egen Form. Ifølge den
uendelige Negativitet kommer det Almene allerede frem som Begrebet af
den første Trilogie. Thi hvad er vel her det Almene? Det Almene er
hverken Væren for sig eller Intet for sig eller Tilværen for sig. Ved
at være et enkelt af disse Led for sig vilde det Almene jo udelukke de
andre Led, og ved denne Udelukkelse være, ikke det Almene, men det
Særskilte. At det Almene heller ikke kan indskrænke sig til at negere
et enkelt Led og affirmere de andre, er ligeledes indlysende: den
uendelige Negativitet viser sig deri, at det Almene ligelig negerer og
ved Negationens Negation affirmerer alle enkelte Led, idet det i dem
alle ligelig negerer og affirmerer sig selv. I den uendelige
Negativitet er ganske rigtig det Almene et Enkelt, thi den Negation,
der forvandler alle Led til Bestemmelser ved det Almene selv,
\textbf{er som} Subjectivitetens meest subjective Virken det
kategoriske Udtryk for den rene Selvhed, den logiske Enkelthed.
Ligesaa er det logisk Enkelte det Almene; thi den negativt virksomme
Selvhed holder sig ikke tilværende som et umiddelbart Individ, der
% a small round mark stands between "at" and "fortære"; taken for plate or
% scanner dirt (this page and p. 160 both carry loose specks), not a sort
mættes ved at fortære de enkelte, ved deres indbyrdes Forskjellighed
særskilte Led; tvertimod! den negative Selvhed negerer just sin egen
Umiddelbarhed, idet den negerer de
% --- p. 155 ---
\setcounter{footnote}{0}%
umiddelbare Led, og affirmerer igjen de umiddelbare Led som
Bestemmelsesudtryk, hvori den affirmerer sig selv. Under disse
% printed "Omboininger" without ø, as on p. 151. Transcribed as printed.
Omboininger blive da alle særskilte Bestemtheder almene Udtryk for det
ene og samme Enkelte, der i dem alle bestemmer sig som det Almene.

Fra dette Synspunkt skulde man altsaa vente, at den hegelske Logik ved
sin væsentlige Dom: \emph{det Enkelte er det Almene}, maatte aabne os
et rigt Indblik i den begribende Subjectivitets uendelige
Selvfordybelse, og derved udlede den logiske Bestemmelse af
Modsætningen mellem det Subjective og det Objective; men, hvor blive
vi ikke skuffede! Kun paa et og andet enkelt Sted tages der ligesom
Tilløb til en Opfattelse af det Almene i dets Subjectivitet. „Das
Allgemeine“ --- hedder det%
\footnote{% printed mark: *)
% set in halbfette (fat-face) Fraktur, unlike the plain footnote on p. 151
\textbf{Hegel. Anfr. Skr. S. 39.}} --- „ist die \emph{freie} Macht; es
ist es selbst, und greift über sein Anderes über; aber nicht als ein
\emph{Gewaltsames}, sondern das vielmehr in demselben ruhig und
\emph{bei sich selbst ist.} Wie es die freie Macht genannt worden, so
könnte es auch die \emph{freie} Liebe und schrankenlose Seligkeit
genannt werden, denn es ist ein Verhalten seiner zu dem
\emph{Unterschiedenen} nur als zu \emph{sich selbst}, in demselben ist
es zu sich selbst zurückgekehrt.“ At dette, der saa udtrykkelig lyder
paa det Subjective, dog ikke desto mindre er meent om det Objective,
at her i Udviklingen af det Enkeltes Forhold til det Almene er en
Confusion af det Subjective med det Objective, det Logiske med det
Antilogiske, vil endnu yderligere opklares, naar vi i det næst
Følgende betragte Sagen fra det objective Standpunkt.

\subsubsection*{β) Positionsbestemmelser: Begrebet i antilogisk Form.}

Den formelle Logiks Grundsætning: $A$ er $A$, peger paa een Gang i
logisk og antilogisk Retning. Forstaaes Sætningen
% --- p. 156 ---
\setcounter{footnote}{0}%
om det Almindelige, bliver den abstract logisk: \emph{det Almindelige
er det Almindelige}; forstaaes det om det Enkelte, bliver den abstract
antilogisk: \emph{det Enkelte er det Enkelte.} I begge Tilfælde er den
begrebløs; thi hvor formelt man end vil tage Begrebet, saa vil den
formelle Begriben dog altid kræve, at det Almindelige udvikler sig til
en Totalitet af Forskjelsbestemmelser. Det Almindelige maa da under en
saadan Udvikling vise sig baade som Eet med og forskjelligt fra sig
selv. Men meest iøinefaldende er dog Begrebløsheden fra Synspunktet af
det Enkelte; thi en Sætning som: \emph{Dette er Dette}, peger saa
direct paa det i Objectiviteten Objective, det Antilogiske, at det
Enkelte i sin Enkelthed endog ligefrem bliver unævneligt.

Det er heraf forklarligt, at den abstracte Tænkning kan ved strengt
at holde paa Tautologien: $A$ er $A$, forivre sig med det Logiske i
den Grad, at den, uden selv at begribe hvorledes, seer sig fangen af
det Antilogiske. Ogsaa i denne Henseende frembyder Eleaternes og
Megarikernes Logik mærkelige Exempler. Om Stilpon berettes der, at han
kun ansaae det Almindelige, ikke det Individuelle, for det Sande. Idet
han da af Iver for den rene Identitet: \emph{det Almene er det
Almene}, ret vilde holde paa det Logiske, kom han, som bekjendt, til
at opstille den Paastand, at man ikke kunde sige Andet end identiske
Sætninger, f. Ex. „Menneske er Menneske, Hest er Hest“; men ikke:
„Hesten løber“, efterdi „Hest“ og „løbe“ ikke ere identiske; --- han
blev da paa denne Maade kastet fra det Logiske over i det Antilogiske,
idet han kom til at negte Gyldigheden af den logiske Dom, der sætter
Prædicatet som Eet med og dog forskjelligt fra Subjectet. Det Samme
gjælder om Eleaterne med Hensyn paa det „\emph{Ene}“. Realiteten af
det Ene bygges, som vi ovenfor have seet, paa de to tautologiske
Sætninger: „det Værende \textbf{er}“, „det Ikke-Værende \textbf{er
ikke}“; men tautologiske Sætninger bevirke let et pludseligt Omslag af
det Logiske i det Antilogiske. Det
% --- p. 157 ---
\setcounter{footnote}{0}%
Mangfoldige skal her være det Utænkelige, altsaa det Ulogiske, det
Ene derimod det Tænkelige, det eneste Tænkelige, altsaa det eneste
Logiske; og see, saa viser det sig, at dette eneste Logiske just er
antilogisk. Hvad det Antilogiske, det i Objectiviteten Objective
egentlig betyder, har Plato vel ikke indseet; men han har dog med
dialektisk Kunst og beundringsværdig Skarpsindighed, navnlig i sin
„Parmenides“, tilstrækkelig viist, at det eleatiske Ene fra alle Sider
staaer i Strid med enhver mulig Form af det Logiske, at det ifølge
Antagelsen absolut Positive fra alle Sider opløser sig i det reent
Negative. „Antages,“ --- hedder det --- „at det Ene er (εἰ ἕν ἐστιν)
i skarp Adskillelse fra al mulig Fleerhed, da kan dette Ene for det
Første ikke være et Hele; thi dersom det var et Hele, maatte det have
Dele, og kom derved til at bestaae i Delenes Fleerhed. Har det Ene
ingen Dele, da har det heller ikke enten Begyndelse (ἀρχή) eller Ende
(τελευτή) eller Midte (μέσον), thi slige Momenter vilde jo da være
dets Dele. Men Begyndelse og Ende ere jo enhver Tings Grændser
(πέρας); har det Ene hverken Begyndelse eller Ende, maa det være
ubegrændset (ἄπειρον), og uden Skikkelse (ἄνευ σχήματος). Det kan da
hverken være deelagtigt i det Runde eller i det Lige, kan hverken
omgives af et Andet eller af sig selv, kan hverken være i Bevægelse
eller i Hvile, kan i sin uendelige Selvmodsigelse overhovedet ikke
taale nogetsomhelst Prædicat, kan, med eet Ord, slet ikke være.“%
\footnote{% printed mark: *)
\textbf{Phil. Propæd. S. 131 flgd.}}

Hvad der gjælder om det absolut Ene, forsaavidt dette Ene skal ---
ved udtrykkelig Udelukkelse af alle Negationer --- være det absolut
Positive, gjælder ogsaa for en Fleerhed af slige Positive, gjælder,
med andre Ord, for de absolute Atomer og de herbartske Realer; kun at
Sætningen: \emph{det Almene er det Almene}, helst ombyttes med
Sætningen: \emph{det Enkelte er det Enkelte.} Hvert Værende
\textbf{er} her et Dette\-%
% --- p. 158 ---
\setcounter{footnote}{0}%
værende, dette Enkelte, dette Positive. Det ene absolut Positive kan
følgelig slet ikke have med det Andet at gjøre; om her i
Virkeligheden er utallige eller kun eet eneste, er for det enkelte
Positive ligegyldigt: Forskjellen mellem Eet og Flere angaaer kun en
udvortes Betragtning, en oversvævende Reflexion. Den absolute
Modsigelse, der er i det enkelte absolute Atom, det enkelte absolute
Reale, forsvinder ikke, fordi her forudsættes en uendelig
Mangfoldighed af slige absolut Positive; det beviser tvertimod, at her
er en uendelig Mangfoldighed af absolute Modsigelser. Ved at sætte
Negativiteten i Bevægelse, fordre Sammenhæng og dermed
Forskjelseenhed og Relativitet, reagerer det Logiske overalt mod det
abstract Antilogiske. Og saa stærk er Reactionen, at selv en Tænker,
saa skarpsindig og saa egensindig som Herbart, maa gjøre Virkeligheden
Indrømmelser, maa slaae lidt af paa Conseqvensen og gaae ind paa den
Forudsætning, at Realerne med al deres Positivitet dog ere i et vist
Forhold til det Negative, at de med al deres Egenhed dog optage
Reflexer af Andethed, at de i deres absolute Absoluthed dog ere
relative, hvilket Altsammen vil sige: at Realerne, trods det i deres
Væren Antilogiske, dog \textbf{ere logiske}.

Men ligesaa afgjørende som det Logiske reagerer, naar man eensidig
vil fastholde det abstract Antilogiske, ligesaa afgjørende reagerer
igjen det Antilogiske, naar man vil opløse dets eiendommelige
Bestemthed i det blot Logiske. Hegels Behandling af den sandsende
Bevidsthed i hans „Phänomenologie des Geistes“ er i denne Henseende
charakteristisk%
\footnote{% printed mark: *)
\textbf{Jvfr. Philos. Propæd. S. 37.}}. Tankebevægelsen er følgende:
„den sandsende Bevidsthed er umiddelbar Vished om en \emph{udvortes
Gjenstand.} Udtrykket for en saadan Gjenstands Almindelighed er, at
den \emph{er, og er denne}, nu efter Tiden, her efter Rummet.
Gjenstanden er altsaa derved aldeles forskjellig fra alle andre
% --- p. 159 ---
\setcounter{footnote}{0}%
Gjenstande, og for sig tilstrækkelig bestemt. Men dette Nu er ikke
mere, et andet er traadt i dets Sted, men som dog ogsaa er et Nu: Nuet
er i sin Forsvinden et Blivende, nemlig det Almene. Dette
\emph{Her}, jeg mener og paapeger, har et Høire, et Venstre, et
Foroven, et Forneden, et Bagved, et Foran i det Uendelige. Dette
\emph{Her} er altsaa ikke det enkelte, umiddelbart bestemte, men: nu
her, nu her, nu her! et Indbegreb af Mange, og i disse Mange dog det
Ene, altsaa det Almene.“ Men saa fortræffeligt den hegelske Tænken
\textbf{end} forstaaer at opdage det Almene i det umiddelbart
Individuelle og finde det Logiske \textbf{der}, hvor det synes at
blive borte i sandselige Forestillinger og tilfældige Meninger, saa
eensidig er igjen Hegels Bestræbelse for at opløse Alt, selv det
Antilogiske, i logiske Reflexer. Den Panlogisme, Hegel allerede i
første Afsnit af sin Aandsphænomenologie fik grundlagt, og som han i
Logiken og i Philosophiens øvrige Dele gjennemførte med en saa
forbausende Alsidighed, Conseqvens og Udholdenhed, maatte omsider
fremkalde en Reaction. Denne Reaction, der i L. Feuerbach havde sin
talentfuldeste Repræsentant, gik nu til den modsatte Yderlighed, lod
det Almene forsvinde i det umiddelbare Enkelte og endte med at beraabe
sig paa „det Uudsigelige“. „Das Sein“ --- siger Feuerbach%
\footnote{% printed mark: *)
Grundsätze der Philosophie der Zukunft \S~26 flgd.} --- „ist an und
für sich selbst Nichts. Aber eben deswegen ist auch das Sein, wie es
die speculative Philosophie in ihr Gebiet hereinzieht und dem Begriffe
vindicirt, ein pures Gespenst, das absolut im Wiederspruch steht mit
dem wirklichen Sein und dem, was der Mensch unter Sein versteht.“
Fremdeles (\S~28): „Welch ein gewaltiger Unterschied ist zwischen dem
Diesen, wie es Object des abstracten Denkens, und eben demselben, wie
es Object der Wirklichkeit ist! Dieses Weib z. B. ist mein Weib,
\emph{dieses Haus mein Haus}, obgleich jeder von seinem Weibe
\textbf{und} seinem Hause wie ich sagt: dieses Haus,
% --- p. 160 ---
\setcounter{footnote}{0}%
dieses Weib\dots{} Das Sein, gegründet auf lauter solche
Unsagbarkeiten, ist darum selbst etwas Unsagbares.“ --- Her staae vi
da igjen ved det Antilogiske.

At Tænkningen, efter saamange systematiske Forsøg, kan vedblive at
svinge imellem to saa afgjorte Extremer som det hegelske og det
feuerbachske, er et Beviis paa, at Begrebet af en virkelig Begriben
endnu ikke er bragt til klar Bevidsthed. Ere de to Sætninger:
\emph{det Almene er det Almene}, og: \emph{det Enkelte er det
Enkelte}, hver for sig begrebløse; er hiin den abstract logiske, denne
den abstract antilogiske, saa indsees heraf, at Begribningen maa beroe
paa en Vexelbestemmelse af det Almene og det Enkelte, altsaa paa en
Vexelbestemmelse af det Logiske og det Antilogiske.

Sætningen: \emph{det Enkelte er det Almene}, \textbf{kan ikke}
udtrykke enten Begrebets eller Begribningens Realitet, saalænge det
antilogiske Element, der skjuler sig i „det Enkelte“, bliver overseet.
Men dette Element faaer kun sin logiske Gyldighed derved, at det
modsatte Extrem betegnes, ikke ved Sætningen: \emph{det Enkelte er det
Enkelte}, men ved Sætningen: \emph{det Almene er det Enkelte.} Først
naar det bliver klart, hvad der egentlig ligger i denne Sætning, vil
det kunne blive klart, hvad Vexelbestemmelsen af det Logiske og
Antilogiske har at betyde.

Om det Enkelte i dets umiddelbare Existeren gjælder det, at det baade
er og dog ikke er det Almene. Hvad er f. Ex. dette Enkelte? Et Bord!
Som Existens er Bordet det Enkelte; men som Kategorie er det det
Almene. Skal det nu ret indlyse, hvad Forholdet er mellem det Enkelte
og det Almene, Existensen og Kategorien, da er det i Sandhed ikke nok,
at vi gaae ud fra det Enkelte og spørge efter det Almene; her fordres
udtrykkelig, at vi tillige vende Forholdet om, gaae ud fra det Almene
og spørge efter det Enkelte. Hvad er altsaa et Bord? Just et Saadant
som dette Existerende! Det Almene er dette Enkelte. Vistnok bryder nu
% --- p. 161 ---
\setcounter{footnote}{0}%
Modsigelsen frem fra begge Sider; thi det Enkelte vil i sin Enkelthed
ikke være det Almene: Existensen protesterer mod sin Kategorie, det
Antilogiske mod sit Logiske. Paa den anden Side vil det Almene i sin
Almindelighed heller ikke være dette Enkelte; thi dette Enkelte falder
jo udenfor hiint Enkelte (dette Bord har Intet at gjøre med hiint
Bord), og just derfor falder Forskjellen imellem det Enkelte og det
Enkelte udenfor det Almene: Kategorien protesterer mod sin Existens,
det Logiske mod det Antilogiske. Men idet Modsigelsen saaledes paa een
Gang bryder løs fra to modsatte Sider, opløser den sig selv, og bukker
under for Begrebet.

Den logisk positive Dom: \emph{det Enkelte er det Almene}, hviler paa
et negativt Grundlag; den forudsætter nemlig, at det Enkelte og det
Almene ere to forskjellige Begrebsudtryk for samme uendelige
Negativitet, den forudsætter, at det Enkelte dog ikke er det Enkelte,
det Almene ikke det Almene. Den logisk negative Dom: \emph{det Enkelte
er ikke det Almene}, hviler paa et positivt Grundlag; den forudsætter,
at Enkeltbestemmelserne kun negere Almeenbestemmelserne, fordi de selv
ere, ikke blot antilogiske, men logiske, altsaa almene. Saavel fra den
positive som fra den negative Side begribes det Enkelte ved det
Almene; men Begrebets Realitet beroer udtrykkelig derpaa, at det
Antilogiske begribes ved det Logiske.

Men naar Begrebets Realitet fordrer, at det Antilogiske begribes ved
det Logiske, da maa Begrebets Realitet jo ogsaa fordre, at det Logiske
begribes ved det Antilogiske, saaledes komme vi da til den antilogiske
Dom: \emph{det Almene er det Enkelte}. At den antilogisk positive Dom
ligeledes hviler paa et negativt Grundlag, er let at see. Det Almene
er ikke det Enkelte, forsaavidt Positionerne i det Almene ere logiske,
i det Enkelte antilogiske. Men dersom det Antilogiske nu forsvinder i
det Logiske, maa dette selv, formedelst Tautologien, blive antilogisk.
Hertil svarer, hvad Stilpon siger om Kaalen: „Den Kaal, som her
faldbydes, er ikke; thi Kaal var allerede
% ved sidefoden signaturmærket „Grundideernes Logik.“ med arknummeret 11; ikke transskriberet
% --- p. 162 ---
\setcounter{footnote}{0}%
for 1000 Aar siden, og denne Kaal er ei den, som var for 1000 Aar
siden.“ Skal det Logiske derimod forsvinde i det Antilogiske, maa
dette, trods dets „Uudsigelighed“, dog formedelst Tautologien blive
abstract logisk. Hertil svarer det Feuerbachske: „\emph{Dieses Weib
ist mein Weib, dieses Haus mein Haus}“. Begrebets Realitet maa da, som
Følge heraf, altid fordre, at dets Form er logisk, naar dets Indhold
er antilogisk, og omvendt. Kun i antilogisk Form faaer Begrebets
Indhold objectiv Existens; dette har Kant fra sit Standpunkt udtalt i
den bekjendte Sætning, at hundrede virkelige Daler ikke indeholde det
Mindste mere end hundrede mulige. Men, det forstaaer sig, Begrebet
kunde jo umulig blive udtrykt i antilogisk Form, dersom dets logiske
Indholdsbestemmelser ikke selv vare bestemmende for det existentielle
Udtryk. Det er denne Side af Sagen, Hegel saa eensidig fremhæver, naar
han idelig og idelig indskærper, at Begrebet realiserer sig ved sig
selv.

Kant og Hegel misforstaae, hver paa sin Maade, det oprindelige Forhold
mellem Realbegrebet og dets Existens. Vistnok har Kant Ret, naar han
antager, at Begrebets logiske Form, for sig betragtet, maa være
subjectiv, og at alle Bestemmelser ved denne subjective Form kunne
være givne, uden at Noget endnu dermed er afgjort angaaende Begrebets
Existens. „Et Begreb mister ikke en eneste af sine Egenskaber, fordi
dets Gjenstand ikke er til, og vinder ikke en eneste Egenskab derved,
at dets Gjenstand findes existerende.“ Dette kan med andre Ord
udtrykkes saaledes: det Antilogiske, hvorved enhver af Begrebets
logiske Bestemmelser skal integreres, for at Begrebet selv kan erholde
sin fuldkomne Realitet, kan ikke naaes ved Tilvæxt af nye
Bestemmelser, der igjen hentes fra det Logiske. Men herved overseer
saa Kant, at et Realbegreb umulig kan slaae sig til Ro i et System af
blot subjective Bestemmelser, at det paa ingen Maade kan nøies med en
intellectuel Tilværen i Subjectivitetens rene Forstand; thi
% --- p. 163 ---
\setcounter{footnote}{0}%
hvad er et Realbegreb? Svar: et Begreb, hvis Indhold altid er
antilogisk, naar dets Form er logisk, og omvendt.

Forsaavidt nu Indholdet er forudsat i virkelig Existens, er Begrebet i
sin logiske Form et Abstractionsbegreb. I denne Forstand er Zoologiens
Begreb \emph{Hest} et Abstractionsbegreb. I den endelige,
psychologiske Subjectivitet kunne vistnok slige Abstractionsbegreber
siges at have et Slags abstract Tilværelse som rene logiske Former.
Men de logiske Positioner have dog kun Begrebsgyldighed i og ved de
antilogiske; den abstracte Formbestemmelse er en halveret
Begrebsbestemmelse; en saadan abstract Formbestemmelse faaer først
Realitet som Indholdsbestemmelse i Eenhed med sit virkelige Indhold.
Abstractionsbegrebet \emph{Hest} er altsaa ikke, som Kant antager, et
fuldstændigt Begreb, saa at hundrede virkelige Heste ikke skulde
indeholde mere end hundrede mulige. Abstractionsbegrebet er tvertimod
et Skyggebegreb: det er saa langt fra, at et saadant Skyggebegreb
skulde være ligegyldigt mod sin Existens, at dets Realitet netop
beroer paa dets Existens; det er saa langt fra, at hundrede mulige
Heste skulde i Begrebet indeholde Ligesaameget som hundrede virkelige,
at Begrebet i een eneste virkelig Hest tvertimod indeholder uendelig
Mere, end det rene Begreb, Abstractionsbegrebet af alle mulige Heste,
efterdi et Realbegreb i dets Realitet altid indeholder uendelig Mere
end dets logiske Afskygning\footnote{Derom veed saavel Zoologien som de
øvrige reale Videnskaber ypperlig Besked.}. % trykt notemærke: *)

Forsaavidt Realbegrebets Indhold er forudsat i en mulig Existens, er
Begrebet i sin logiske Form et Constructionsbegreb. I denne Forstand
er f. Ex. Begrebet \emph{Huus} for Bygmesterens Subjectivitet et
Constructionsbegreb. At nu Constructionsbegrebet ligesaa lidt kan
blive staaende ved den blot logiske Form, som Abstractionsbegrebet, er
indlysende. Men medens Abstractionen altid bestemmes ved den logiske
Dom: \emph{det Enkelte er det Almene}, maa Con\-%
% ved sidefoden signaturmærket 11* ; ikke transskriberet
% --- p. 164 ---
\setcounter{footnote}{0}%
structionen bestemmes ved den antilogiske Dom: \emph{det Almene er
dette Enkelte}. Naar En bliver spurgt, hvad en Trekant er, og besvarer
Spørgsmaalet ved at tegne en Trekant, saa svarer han ved en
Construction, d. v. s. han udtrykker sin Dom om Trekanten, ikke paa
logisk, men paa antilogisk Maade. Formedelst Abstractionen blive de
antilogiske Positioner omsatte i logiske; formedelst Constructionen
blive de logiske Positioner omsatte i antilogiske. Forsaavidt
Abstractionen maa lade Indholdet blive, hvad det er i sin Existens, og
kun kan forandre Formen, vil den af Abstractionen betingede Begriben
gribe det antilogiske Indhold i logisk Form. Hvad Constructionen
derimod angaaer, da kan den jo ikke skee med fuld Indsigt, dersom der
under Planens Udførelse indblander sig noget Uvedkommende og for
Begrebet Forstyrrende; de logiske Positionsbestemmelser maae da fra
dette Synspunkt opfattes som Begrebets egne Indholdsbestemmelser; de
antilogiske derimod som Formbestemmelser. Den af Constructionen
betingede Begriben vil derfor altid gribe et logisk Indhold i
antilogisk Form.

Saavel fra Abstractionens som fra Constructionens Synspunkt vise
Begribningen og Begrebet os en udtrykkelig Modsætningseenhed af det
Subjective og det Objective. Her komme vi nu igjen til Hegels
Misforstaaelse. Thi det forholder sig aldeles ikke, som Hegel mener,
at det subjective Begrebs Realitet skulde i Logiken falde sammen med
Subjectivitetens egen Forsvinden i det Objective. Vistnok forsvinder
Subjectiviteten hos Hegel i det Objective, for at vende tilbage i
Teleologien; men da denne dens Tilbagevenden i den hegelske Logik er
ligesaa mystisk og gaadefuld, som dens Forsvinden, kan det ikke hjælpe
at beraabe sig paa den.%
% TRYKFEJL: s. 164 bærer noten *) ved sidefoden, men INTET notemærke er sat
% i brødteksten; hele siden er gennemgået ved 600 dpi. Noten er her hængt
% paa afsnittets slutning; dens indhold peger paa „Teleologien“ ovenfor.
\footnote{„Es ist“ --- bemærker Hegel i Anledning af Kants Adskillelse
mellem „objective Sætninger“ og „subjective Maximer“ med Hensyn paa
det Teleologiske (Logik. 2ter Thl. S. 215) --- „auf diesem
% her falder det trykte sideskift midt i noten: s. 164 slutter, s. 165 fortsætter
ganzen Standpunkte dasjenige nicht untersucht, was allein das
philosophische Interesse fordert, nämlich welches von beiden
Principien an und für sich Wahrheit habe; für diesen Gesichtspunkt
macht es keinen Unterschied, ob die Principien als \emph{objektive},
das heißt: hier äußerlich existirende Bestimmungen der Natur, oder als
bloße Maximen eines \emph{subjektiven} Erkennens betrachtet werden
sollen“ \dots{} Men hvad er „an und für sich Wahrheit“, naar det
Subjective og det Objective løbe ud i Eet?}

% --- p. 165 ---
\setcounter{footnote}{0}%
Hvor umuligt det er, ifølge den hegelske Methode, at naae til
nogensomhelst Charakteristik af det i Objectiviteten Objective, er
allerede viist i det Foregaaende og kan end yderligere sees deraf, at
den med alle sine Variationer over Umiddelbarhedens rige Thema, ikke
veed at gjøre Forskjellen imellem logisk og existentiel Umiddelbarhed
logisk kjendelig. „Es sind, wie bereits erinnert worden, schon mehrere
Formen der Unmittelbarkeit vorgekommen; aber in verschiedenen
Bestimmungen. In der Sphäre des Seyns ist sie das Seyn selbst und das
Daseyn; in der Sphäre des Wesens die Existenz und dann die
Wirklichkeit und Substantialität, in der Sphäre des Begriffs außer der
Unmittelbarkeit, als abstrakter Allgemeinheit, nunmehr die
Objektivität \dots{} Seyn ist überhaupt die \emph{erste}
Unmittelbarkeit, und \emph{Daseyn} dieselbe mit der ersten
Bestimmtheit. \emph{Die Existenz} mit dem Dinge ist die
Unmittelbarkeit, welche aus dem \emph{Grunde} hervorgeht, --- aus der
sich aufhebenden Vermittelung der einfachen Reflexion des Wesens.
\emph{Die Wirklichkeit} aber und die \emph{Substantialität} ist die
aus dem aufgehobenen Unterschiede der noch unwesentlichen Existenz als
Erscheinung und ihrer Wesentlichkeit hervorgegangene Unmittelbarkeit.
\emph{Die Objektivität} endlich ist die Unmittelbarkeit, zu der sich
der Begriff durch Aufhebung seiner Abstraktion und Vermittelung
bestimmt.“

Ved at betragte denne Umiddelbarhedernes Skala, vil man let overbevise
sig om, at det Umiddelbare paa ethvert
% --- p. 166 ---
\setcounter{footnote}{0}%
Trin, fra det nederste til det øverste, er af reen logisk Oprindelse.
Dannelsen af enhver saadan objectiv logisk Umiddelbarhed beroer
aldeles ikke paa noget i sig Existentielt, men udelukkende paa den
tænkende Subjectivitets almene Begrebsudvikling. De nævnte
Umiddelbarheder ere, hver paa sit Stadium, fremkomne ved en Negations
Negation, altsaa ved en indenfor Reflexionen ophævet Reflexion; det er
lutter Affirmationer uden en eneste existentiel Position. Den
Umiddelbarhed, som her skal forestille Objectiviteten, er desaarsag i
Virkeligheden længere fra at være noget Objectivt, end den
Skuespiller, der i „en Sommernatsdrøm“ skal forestille Væggen, er fra
at være Væg.

Den existentielle Umiddelbarhed er antilogisk; det Antilogiske er det
i Objectiviteten Objective. Istedenfor at abstrahere fra dette
Objective og udvikle et „objectivt Begreb“ uden Objectivitet, maa
Logiken just lægge an paa at klare Modsætningseenheden af det
Subjective og det Objective ved fra alle Sider at klare
Modsætningseenheden af det Logiske og det Antilogiske.

\subsubsection*{γ) Begrebets Realitet: den fuldkomne Begriben.}

Modsætningseenheden af det Subjecttve og det Objective % TRYKFEJL: „Subjecttve“ (t-sort for i) som trykt; jf. 600 dpi
er i Begrebet selv bestemt ved en Modsætningseenhed af Negations- og
Positionsbestemmelser. Alle Negationsbestemmelser ere
Vidensmomenter, alle Positionsbestemmelser Magtmomenter. Realbegrebet
i sin Fuldendelse, d. v. s. Begrebet selv efter sin absolute Realitet,
er en totaliseret Eenhed af negative og positive Bestemmelser, af
Vidensmomenter og Magtmomenter. Den fuldkomne Begriben er altsaa ikke
en blot menneskelig, men en guddommelig Sag; her sees Modsætningen
imellem den psychologiske og den ontologiske Subjectivitet.

Paa det psychologiske Standpunkt er Magten skilt fra Viden. Systemet
af Magtbestemmelserne afgiver den virkelige
% --- p. 167 ---
\setcounter{footnote}{0}%
Omverden, Systemet af Vidensbestemmelserne findes i Subjectiviteten;
hiint fremstiller paa eensidig Maade Realiteten, dette Idealiteten;
hiint udtrykker Begrebet i antilogisk, dette i logisk Form.
Skilsmissen imellem det Subjective og det Objective er her saa
gjennemgaaende, at Existensbegrebet paa een Gang maa synes baade at
fordre Existens og at modsætte sig Existensen. „Sættes det hele System
af sandselige Intellectualbestemmelser i den intellectuelle Form, saa
forsvinder det Sandselige som et i det Tænkelige ophævet
Reflexionsmoment, det vil sige: det Existerende forvandler sig til
Existensbegreb, men Existensen gaaer ud. Sættes omvendt det hele
Begreb af intellectuelle Sandsebestemmelser i den sandselige Form,
svarende til umiddelbare Sandseindtryk, saa træder Gjenstanden i
Existens, men Begrebet gaaer ud.“\footnote{Jvfr. Forelæsn. ov. „Phil.
Propæd.“ 1860--61, S. 24--30.} % trykt notemærke: *)

Begrebet kan som Begreb umulig staae stille: en Steen kan
\emph{ligge} stille, en Muur kan \emph{staae} stille; men et Begreb
kan ligesaa lidt indgaae til sandselig Hvile, som det kan indlade sig
med den sandselige Forskjel imellem at ligge, at gaae og at staae. Og
dog er Begrebet paa een Gang baade i Hvile og i Bevægelse; dets
ideelle Bevægelse er Eet med dets ideelle Hvile, thi Bevægelsen er en
Begriben. Det er en haardnakket Fordom, en Fordom, som „den objective
Logik“ fra sin Side har givet ligesaa rigelig Næring, som „den
subjective“ fra sin, at Begrebet, „det i og for sig Værende“, skulde
befinde sig objectivt i en objectiv Væren, istedenfor at beroe paa en
virkelig Begriben.

Vistnok seer det fra det psychologiske Synspunkt ud, som om her maatte
gjøres en endelig Adskillelse imellem Begrebernes Vorden og deres
Tilværen; Begreberne vorde jo, idet vi logisk udvikle dem, og de have
som Existensbegreber en virkelig Tilværen i de Gjenstande, hvoraf vi
abstrahere dem, og hvori vi gjenkjende dem. Saaledes har f. Ex.
Begrebet Feldtspath
% --- p. 168 ---
\setcounter{footnote}{0}%
sin Existens i den tilsvarende Steen, og i Stenens sandselige Existens
en sandselig Hvile. Fremdeles har Begrebet Feldtspath en Vorden i
Mineralogens Bevidsthed, naar han enten selv udvikler det eller faaer
det udviklet af en Anden. Hvad der paa denne Maade gjælder om
Abstractionsbegreberne, maa da ogsaa gjælde om
Constructionsbegreberne. I Bygmesterens Bevidsthed kan Husets Begreb
med Grund siges at have en ideel Tilværen, medens det under Byningens
% TRYKFEJL: „Byningens“ for Bygningens (g-sorten mangler); kalibreret mod
% det korrekt satte „Bygningen“ tre linier længere ned paa samme side
Opførelse befinder sig i en reel Vorden; er Bygningen omsider færdig,
finder det realiserede Begreb en sandselig Hvile i den hvilende
Bygning. Her er altsaa, hvad Begrebet selv angaaer, en endelig
Adskillelse mellem ideel Tilværen, reel Vorden og reel Tilværelse.

Men hvor vaklende og slibrige slige Adskillelser ere, vil let sees ved
en nærmere Betragtning. Den existerende Gjenstand er ikke det Samme
som det existerende Begreb: thi dersom Begrebet uden videre var
Gjenstanden, maatte man ved at bryde et Stykke af Gjenstanden bryde et
Stykke af Begrebet; den røde Gjenstand vilde da være et rødt Begreb,
den tunge et tungt, den vaade et vaadt Begreb o. s. v.; hvilket
Altsammen er urimeligt. Det gaaer, fra dette Synspunkt, med det
realiserede Begreb, som med den realiserede Mulighed: det forsvinder i
sin Realitet, dets umiddelbare Existens er dets Undergang. At
abstrahere Begrebet af tilværende Gjenstande er altsaa med andre Ord:
at kalde det i Gjenstandens Existens uddøde Begreb tillive, at tænde
paany den Idealitetens Gnist, der udslukkes i Realiteten; det vil
sige: at begynde Begribningen forfra, thi Begrebet er kun i en
Begriben. At realisere Begrebet i en til dets Fordringer svarende
Construction er at ophæve Begrebet ved Begrebet selv, og sætte
Realiteten i Idealitetens Sted.

Realbegrebet kan paa denne Maade aldrig komme til Ro: naar dets Form
er logisk, saa har det, som Abstractionsbegreb, sit Indhold udenfor
sig; thi saa er Indholdet anti\-%
% --- p. 169 ---
\setcounter{footnote}{0}%
logisk. Naar dets Indhold er logisk, saa har det, som
Constructionsbegreb, sin Form udenfor sig; thi saa er Formen
antilogisk.

Dersom denne Uro skal ophøre, da kan det ikke skee ved at standse
Bevægelsen, thi hvorved skulde den vel standse? Ved Existensen! Men,
naar Existensen indtræder, gaaer Begrebet jo ud, thi Begrebets Form er
logisk. Ved Begrebets Form! Men, naar Begrebet ifører sig sin logiske
Form, kommer det i Splid med Existensformen, thi Existensformen er
antilogisk. Istedenfor at standse Bevægelsen, maae vi tvertimod
uendeliggjøre den, og lade den fuldende sig i et Kredsløb.

Det logiske Kredsløb: fra Begreb gjennem Existens til Begreb, fra
Existens gjennem Begreb til Existens, o. s. fr. kan da nærmere
betegnes som en Bevægelse: fra Construction gjennem Abstraction til
Construction, fra Abstraction gjennem Construction til Abstraction.
Det gjælder her blot om at faae Kategorierne Abstraction og
Construction til at gaae op i fyldigere, mere betegnende Udtryk.
Forsaavidt Abstractionen særlig betinger den Begribningens Act,
hvorved Indholdets antilogiske Form forvandler sig til logisk, maa
denne Kategorie integreres ved Alt, hvad der hører under Dannelsen,
Tilegnelsen og Bearbeidelsen af Realbegrebets ideelle Momenter og
fuldendes i, hvad vi her ville kalde \emph{den logiske Erindren}. I
Kraft af Tilegnelsen og den logiske Erindren besidder Subjectiviteten
altsaa sine Realbegreber paa ideel Maade, og kun under Forudsætning af
en saadan intellectuel Besiddelse kunne Constructionsbegreber danne
sig. Seet i sit Forhold til den fordrede Construction er Begrebet et
subjectivt eller anticiperet Formaal; Construeringen selv med sine
Betingelser er Formaalets Udførelse ved tilsvarende Midler; det ved
Constructionen Frembragte, det realiserede Formaal, er endelig den
Existens, hvori Begrebet døer, for igjen at afgive Stof for en ny
Abstraction og opstaae forklaret i den logiske Erin\-%
% --- p. 170 ---
\setcounter{footnote}{0}%
dren. Constructionens eensidige Kategorie maa her integreres ved Alt,
hvad der logisk hører med til intellectuel-reel Frembringelse.

Hvor væsentlig den intellectuelle Frembringelse er for den formelle
Begriben, derom mindes vi ved den bekjendte Sætning: \textit{scimus,
quia facimus}; hvor nødvendig en intellectuel-reel Frembringen hører
med i en reel Begriben, er charakteristisk betegnet ved Schellings
geniale, men dumdristige Ord: „Ueber die Natur philosophiren heißt die
Natur schaffen“.

Paa det psychologiske Standpunkt kan den til Begribningen svarende
Ideebevægelse vel erkjendes \textit{in abstracto}, men ikke
gjennemføres \textit{in concreto}. For den psychologiske Subjectivitet
gaae Abstractionsbegreberne i en Retning, Constructionsbegreberne i en
anden; her falder Realiteten, Begrebets Gjenstand, udenfor Begrebet
selv; her falder Magten udenfor Viden. I den ontologiske Subjectivitet
derimod er Magten identisk med Viden; og paa denne Identitet beroer
den fuldkomne Begriben. Da al Begriben væsentlig beroer paa Erindren
og Frembringen, maa den reale Begriben nødvendig grunde sig paa en
Formaaen, en Magten. Begrebet i sin Realitet er Totalitet, i
Totaliteten ere de enkelte Dele paa een Gang Led og Momenter. Denne
Dobbeltsidighed kan, hvad Begrebet angaaer, ikke længere forklares ved
en reen Reflexionsadskillelse mellem det Umiddelbare og dets Reflex:
Totalitetens tilværende Led have som Udtryk for den umiddelbare og
reelle Magt en antilogisk Selvstændiged. % TRYKFEJL: „Selvstændiged“ for Selvstændighed (h-sorten mangler); som trykt
Med denne umiddelbare Selvstændighed ere disse Led nu Gjenstande for
den fuldkomne Begriben; men hvorledes blive de Momenter? De blive
gjennemsigtige for den uendelige Viden formedelst Magten; i den
fuldkomne Gjennemsigtighed ere de Momenter. Men hvad bliver det saa
til med den antilogiske Selvstændighed? Den staaer ved Magt, men sees
i Videns Lys; thi den uendelige Magt er selv gjennemsigtig for Viden.
Som
% --- p. 171 ---
\setcounter{footnote}{0}%
Videns endelige Gjenstande ere de existerende Led Elementer i det
reale Begreb, idet de ere Momenter for den ideale Magt, den uendelige
Selvmagt: Magtens Idealitet er Videns Realitet. Ideebevægelsen er en
uophørlig Omsætten af det Logiske i det Antilogiske, det Ideelle i det
Reelle, og omvendt: \emph{en uendelig Vexelbestemmelse af Erindren og
Frembringen}! Men denne Bevægelse er igjen Hvile; Vidensmomentet skal
ikke paa endelig Maade foruroliges ved at lede om sit tilsvarende
Magtelement, eftersom det jo allerede har sit integrerende Element i
sin Magt.

Paa Ideebevægelsens Identitet med Hvilen beroer den endelige
Adskillelse af en Hvile, hvori Bevægelsen er hævet, og en Bevægelse,
hvori Hvilen er hævet. Den Hvile, hvori Ideebevægelsen er hævet, er
Kjendemærket for de udvortes tilværende Ting; her er Begrebet
umiddelbart fortabt i sin Existens; Begrebet er med andre Ord ikke
Begreb, men Gjenstand. Dersom nu denne udvortes, antilogiske Hvile
virkelig var absolut, en absolut Negation af al Ideebevægelse, saa
vilde den sandselige Forstand have Ret, naar den paastaaer, at
Tingene ikke ere realiserede Begreber, at der i Tingenes Verden er et
System af Kræfter og Betingelser, der falder udenfor al Viden. Men
hvor overfladisk denne Synsmaade er, vil kunne indsees, naar vi
besinde os paa, at Tingene umulig kunde være virkelige Objecter,
dersom de ikke vare realiserede Begreber, og at de umulig kunde have
Bestaaen som realiserede Begreber, dersom de ikke bares af en uendelig
Begriben.

Den existerende Ting er et virkeligt Object; det virkelige Object
beroer paa en virkelig Objectivering; men den i Virkeligheden
gjennemførte Objectivering beroer igjen paa en virkelig Begriben. Saa
sikker den psychologiske Subjectivitet end er paa sin Adskillelse af
det Subjective og det Objective, saalænge den kan holde sig til
umiddelbare Indtryk, ligesaa usikker og tvivlraadig bliver den, naar
Reflexionen sættes i
% --- p. 172 ---
\setcounter{footnote}{0}%
Bevægelse. --- „Den sunde Forstand fordrer, at jeg umiddelbart skal
blive mig bevidst, at der er Noget udenfor min Bevidsthed; men, idet
jeg bliver mig Noget bevidst, er det jo i min Bevidsthed; forsaavidt
det er udenfor, kan jeg ikke blive mig det bevidst: altsaa skal jeg
blive mig bevidst, at det, der er i min Bevidsthed, ikke er der;
hvilken Modsigelse! Naar Solen skinner og Træerne speile sig i Aaen,
mener den sunde Forstand, at Lyset er Klarhed, enten der er Nogen, der
seer det eller ikke, og at Træerne virkelig speile sig i Aaen, enten
saa denne Speiling afspeiler sig i et seende Øie eller ikke. Men
Spørgsmaalet er, hvorledes man skal overbevise sig herom. Ved Synet?
Det er en Modsigelse, da her jo handles om Phænomenet, ikke forsaavidt
det sees, men forsaavidt det ikke sees. Ved Tanken? Det er en
Modsigelse, Phænomenet lader sig som Phænomen ikke tænke; thi i saa
Fald kunde den Blindfødte jo ogsaa tænke det. Ved Andres Vidnesbyrd?
En tom Udflugt, da det jo her maa gaae den Ene, som den
Anden.“\footnote{Phil. Propæd. S. 41.} % trykt notemærke: *)

For at undersøge, hvad heri ligger, ville vi i denne Sammenhæng
henføre det almindelige Objectiveringsproblem til de tre særskilte:
Bevægelsesproblemet, Forandringsproblemet og Mulighedsproblemet.
Altsaa:

\emph{Bevægelsesproblemet.} Blandt de flere Beviser, Eleaterne have
udtænkt mod Bevægelsens Realitet, nævne vi her „den flyvende Piil“.
„Skulde et Legeme bevæge sig, saa vilde deraf følge, at det paa een
Gang maatte være i Bevægelse og i Hvile, thi ethvert sig bevægende
Legeme maa i hvert Øieblik være i et Rum, der netop er ligestort med
det selv. Men saaledes er Legemet i Hvile; thi hvad der, om end kun
for et Øieblik, befinder sig i et tilsvarende lige Rum, er med Hensyn
paa dette Rum i Hvile. Tænke vi os en Piil, der med største Hurtighed
gjennemflyver et Rum, saa
% --- p. 173 ---
\setcounter{footnote}{0}%
maa den hvert Øieblik være i et andet Rum og, forsaavidt den indtager
dette, hvert Øieblik være i Hvile. En saadan flyvende Piil, der altid
er: her, her, nu, nu (\textit{ἐν τῷ νῦν κατὰ τὸ ἴσον}), maa da tænkes
at være paa een Gang baade i Hvile og i Bevægelse, hvilket er en
Mosigelse.“\footnote{Jvfr. Mathm. og Dialekt, S. 137 flgd.} % trykt notemærke: *)
% TRYKFEJL: „Mosigelse“ for Modsigelse (d-sorten mangler); som trykt
Man paaviser i Almindelighed det Sophistiske i Beviset ved, efter
Aristoteless Exempel, at lægge Vægt paa Vexelbestemmelsen af
Continuitet og Discretion; men ved at indskrænke sig hertil overseer
man dog i Grunden den egentlige Pointe. Idet Opmærksomheden fæster sig
paa den flyvende Piil, fæster den sig jo ikke paa et formelt Begreb,
men paa et virkeligt Object. Det er ikke Bevægelsesbegrebet i dets
logiske Form, men det i Bevægelsen Objective, Bevægelsesbegrebet i
dets antilogiske Form, den virkelige Bevægelse, der ifølge Eleaternes
Synsmaade er det Utænkelige. Fra den logiske Side kan en Bevægelses
Continuitet meget vel forenes med sin tilsvarende Discretion,
forsaavidt det Forbigangne formedelst Erindringen altid kan forenes
med det Nærværende; men det, hvorpaa det her absolut kommer an, er
netop den Omstændighed, at „den flyvende Piil“ er uden Erindring; den
har i sit Nærværende intet Forbigangent. Det i Objectet Objective kan
umulig bevæge sig uden Overgang fra en foregaaende til en
efterfølgende Tid, fra et foregaaende til et efterfølgende Sted; men
Relationen mellem et Foregaaende og et Efterfølgende er en Reflexion,
der falder udenfor Objectet selv; derfor kan Objectet hverken komme
frem i Rummet eller i Tiden; dets Sted er altid her, her! dets Tid er
nu, nu! Naar det nu desuagtet gjælder, at et virkeligt Object under
sin Bevægelse virkelig tilbagelægger et bestemt Rum i en bestemt Tid,
da er Forklaringen den, at det Tilbagelagte, det Forbigangne, bevares
i den subjective Reflexion, i Betragterens Erindren. Den virkelige
Bevægelse kan med andre Ord
% --- p. 174 ---
\setcounter{footnote}{0}%
kun være objectiv derved, at den objectiveres; men, naar den i
Virkeligheden objectiveres, har den objectiv Gyldighed.

\emph{Forandringsproblemet.} Kan Bevægelsesproblemet kun løses paa
den Betingelse, at det Objectiveringen bestemmende Objective bestemmes
ved Objectiveringen selv, da maa det Samme naturligviis gjælde om
Forandringsproblemet, eftersom enhver Forandring jo maa forudsætte en
Bevægelse. Den meest objective Opfattelse af det i Forandringen
Objective støtter sig til den Forudsætning, at Forandringen beroer paa
Delenes, altsaa, naar vi gaae til Grændsen, paa Atomernes Omflytning.
Men hvad enten Omflytningen nu iøvrigt bestaaer deri, at forbundne
Dele skilles ad, eller deri, at adskilte Dele forenes, saa maa
Forandringen i begge Tilfælde forudsætte en Relation, en virkelig og
virksom Relation mellem Delene ,indbyrdes. % TRYKFEJL: løst kommategn foran „indbyrdes“; som trykt
Her er i Forandringen en Tiltrækken og Frastøden, der maae henføres
til Delene selv. Men nu har det jo fra mange Sider viist sig i vort
Foregaaende, at den objective Relation væsentlig er en Reflexion, og
Reflexionen oprindelig subjectiv. Den Modsigelse, at al objectiv
Relation beroer paa subjectiv Reflexion, finder sin fuldstændige
Løsning i Objectiveringens Ontologie.

Fra det psychologiske Standpunkt er det ganske vist en Urimelighed, om
Nogen vilde paastaae, at der ikke kunde foregaae Forandringer med de
virkelige Ting, uden naar der var en Betragter tilstede, der kunde
objectivere sig de stedfindende Forandringer; men, seet fra et
ontologisk Synspunkt, har Sagen sin Rigtighed. I den virkelige
Forandring er Vexelbestemmelsen af det Logiske og det Antilogiske
iøinefaldende. $A$ forandrer sig ved at gjennemløbe Tilstandene:
$A_0$, $A_1$, $A_2$\dots{} Sees her nu paa den abstract logiske
Identitet: $A = A$, er Forskjellen mellem $A_0$ og $A_1$ abstract
factisk, altsaa antilogisk; bestemmes Forskjellen af $A_0$ og $A_1$
abstract logisk, saa er Identiteten: $A = A$ abstract faktisk, altsaa
antilogisk.

% --- p. 175 ---
\setcounter{footnote}{0}%
Forskjellen imellem den psychologiske og den ontologiske
Subjectivitet er paa eiendommelig Maade at charakterisere ved
den tilsvarende Forskjel imellem en psychologisk og en ontologisk
Objectiveren. Den blot psychologiske Objectiveren forudsætter en
endelig Adskillelse af Viden og Magt, og staaer desaarsag i et
endeligt og tilfældigt Forhold til Objecterne; den ontologiske
Objectiveren, der forudsætter en oprindelig Eenhed af Viden og
Magt, staaer derimod i et saa væsentligt, inderligt og nødvendigt
Forhold til de virkelige Existenser, at her ligesaa lidt kan være
Tale om Objecter uden Objectiveren, som om Følger uden
Begrundelse, Virkninger uden Causalitet, Begreber uden Begriben.

\emph{Mulighedsproblemet.} Den ontologiske Objectiveren er
immanent, d. v. s. en Objecterne selv iboende Bestemmen; men som
en fra Subjectiviteten udgaaende Bestemmen er den immanente
Objectiveren tillige en transcendent Handlen, en de virkelige
Objecter oversvævende Bestemmen. Kun ved denne sin
Dobbeltsidighed kan Objectiveringen magte det Dobbeltsidige i
Objectiviteten selv, nemlig det Nødvendige og det Tilfældige. Fra
dette Synspunkt er Mulighedsproblemet væsentlig et
Objectiveringsproblem. Det er ret charakteristisk for Kategorien
Mulighed, at naar den fastsættes objectivt, saa peger den mod det
Subjective, og naar den fastsættes subjectivt, saa peger den mod
det Objective. Efter Megarikernes Logik kunde der i Virkeligheden
ikke gives nogen Mulighed. „Intet“, siger Diodoros Kronos, „er
sandt, uden det, som er virkeligt eller engang bliver det. Følgelig
er heller Intet muligt uden det, som er virkeligt eller engang
bliver det. Det, som ikke bliver virkeligt, er umuligt, men det,
som bliver virkeligt, var nødvendigt.“ Er Noget i Naturen muligt,
da er det, fra dette Synspunkt, tillige virkeligt og nødvendigt, er
Noget ikke virkeligt, da er det, fordi dets Virkelighed ikke er
mulig (ὅταν ἐνεργῇ, μόνον δύνασθαι, ὅταν δὲ μὴ ἐνεργῇ, μὴ
δύνασθαι). Men hvad beviser saa dette? At
% --- p. 176 ---
\setcounter{footnote}{0}%
Mulighedsbegrebet er af reen subjectiv Oprindelse. Det begynder,
psychologisk seet, som et Uvishedsbegreb, et Begreb, hvorved
Subjectiviteten udtrykkelig giver tilkjende, at den maa lade det
Objective selv raade. Det i Objectiviteten Objective er altsaa ikke
et Muligt; det er tvertimod et Virkeligt. Det Virkelige er, fordi
det er; men hvad der saaledes er, fordi det er, er fra den ene Side
et Tilfældigt, fra den anden et Nødvendigt. Men dersom det
Tilfældige skal være objectivt, saa maa det Mulige ogsaa være det,
thi Tilfældene vexle, og det Mulige ligger i Tilfældenes Vexlen.
Siger man nu: det Mulige er dog subjectivt; det ligger ikke i
Tilfældenes Vexlen, thi Tilfælde følger paa Tilfælde, som et
Virkeligt paa et Virkeligt, og Følgen skeer med Nødvendighed, ---
saa er det altsaa det i Virkeligheden Nødvendige, der er det
Objective. Men hvis Nødvendigheden, nemlig den hypothetiske:
dersom $a$ er, saa maa $b$ være, dersom Tilfældet $a$ er
indtraadt, saa maa Tilfældet $b$ indtræde, virkelig skal være det
Objective, saa maa Muligheden selv være objectiv; thi Muligheden
er da den forventede, men endnu ikke indtraadte Nødvendighed. Den
betingede Nødvendighed er i en bestandig Vorden, og Muligheden er
just denne Nødvendighedens Vorden. „Det er overalt et mærkeligt
Forhold, der er imellem den reale Mulighed og dens Betingelser.
Kun med Hensyn paa den fordrede Virkelighed er Muligheden real:
saalænge her altsaa endnu mangler en Betingelse, kan Muligheden
ikke blive til Virkelighed og er forsaavidt umulig; men, naar alle
Betingelser ere tilstede, er Muligheden jo heller ikke Mulighed,
thi saa er den Virkelighed“%
\footnote{Jvfr. Forelæsn. ov. „Phil. Propæd.“ 1860---61 S. 98---100.}%
, det vil sige: Tilværelse af en betinget Nødvendighed.

Som Nødvendighedens Vorden er Muligheden objectiv og immanent, og
dog er Muligheden tillige subjectiv og oversvævende, thi hvor Noget
blot er muligt, er dets Modsatte
% --- p. 177 ---
\setcounter{footnote}{0}%
jo ogsaa muligt. At det Virkelige hviler i Muligheden og fremgaaer
af den, vil sige: det Objective bestemmes subjectivt; at det Mulige
hviler i Virkeligheden og betinges af den, vil sige: det Subjective
bestemmes objectivt. Som den Kategorie, der betegner Overgangen fra
det Subjective til det Objective, er Muligheden den concrete
Objectiveringskategorie, der i sig har optaget Bevægelsen og
Forandringen. \textit{Le présent est gros de l'avenir}, siger
Leibnitz; dette Svangerskab beviser just den antilogiske Formæling
af det Objective med det Subjective i den uendelige Objectivering;
men denne uendelige Objectiveren er Eet med den fuldkomne Begriben.

% Display sub-head; the printed label is an antiqua "b)".  The book's
% Indhold calls this limb "Idealsystemet og Realsystemet"; the printed
% head reads as below and is what stands in the body.
\subsubsection*{b) Begreb og Idee: Begribningens Phaser.}

„Ideen“, hedder det i den hegelske Logik, „er Eenheden af det
Subjective og det Objective, den Eenhed, der udtrykker
Subjectivitetens Tilbagevenden til sig selv, efter at den objective
Verdens Modsætning er overvunden o. s. v.“ Denne Bestemmelse af,
hvad „Ideen“ er, kan meget vel anerkjendes for rigtig i formel
Henseende; men hvad hjælper det at henvise til Eenheden af det
Subjective og det Objective, naar disse Begreber selv ere tomme
Former? „Die Idee“, siger Hegel, „ist die \emph{Wahrheit}; denn die
Wahrheit ist dieß, daß die Objectivität dem Begriffe entspricht“;
% letterspacing inconsistent: the FIRST "Wahrheit" is spaced, the
% second is not.  Transcribed as printed.
men naar Objectiviteten er Begreb og Begrebet Objectivitet, saa er
jo begges Eenhed blot formel, og Idee-Udviklingen en reen
Formalisme. Som Begrebet bestemmes, saaledes ogsaa Ideen; som
Modsætningseenheden af det Subjective og det Objective udvikles,
saaledes udvikles ogsaa Ideen. Siger man: Begreb og Idee forholde
sig, som den absolute Form til sit absolute Indhold, kan denne
Opfattelse ligeledes gjælde; men her forudsættes da en dybtgaaende
Udvikling af Forholdet mellem det Absolute og det Relative, og en
saadan Udvikling maa da falde sammen med Idee-Udviklingen selv. Den
Indsigelse, der gjælder mod det hegelske Begreb, at det med alle
sine ideelle Concre\-%
% signature mark at foot: "Grundideernes Logik." + sheet number 12
% --- p. 178 ---
\setcounter{footnote}{0}%
tioner dog kun er formelt, fordi det Logiske mangler sit
Antilogiske, den samme Indsigelse gjælder ogsaa mod den hegelske
Idee. For at afvise Formalismen ville vi i denne Sammenhæng
bestemme Ideen som den absolute Eenhed af det Logiske og det
Antilogiske; vi have da med det Samme bestemt den som den absolute
Eenhed af det Absolute og det Relative, af Subjectiviteten og det
Objective, af Realbegrebet og den fuldkomne Begriben.

Den absolute Eenhed af det Logiske og det Antilogiske, af
Subjectiviteten og dens objective Indhold, kan ikke tænkes som en
umiddelbart værende; dens Væren er en Ideebevægelse; kun forsaavidt
vi ere istand til at fremstille Ideebevægelsen, ere vi istand til
at tydeliggjøre Begrebet ved Ideen, Ideen ved Begrebet. Skal
Ideebevægelsen fremstilles, maa Begribningen bevæge sig i en
tredobbelt Kreds, og derunder frembyde følgende tre forskjellige
Phaser. Det Antilogiske bliver Moment for det Logiske, derved at
Magten opgaaer i Viden: denne Kredsbevægelse viser os Videns
Oprindelighed. Det Logiske bliver Moment for det Antilogiske, idet
Viden opgaaer i Magten: denne Kredsbevægelse viser os Magtens
Oprindelighed. Det Logiske og det Antilogiske articuleres ved
alsidig Vexelbestemmelse, idet Begribningens sidste og hoieste
Kredsbevægelse bestemmes ved en Fordoblingseenhed af de tvende
foregaaende: Subjectiviteten sees i Ideen selv.

% Display sub-head; label is an antiqua Greek alpha, the Latin tag is
% antiqua against the surrounding Fraktur.
\subsubsection*{α) Videns Oprindelighed: \textit{universalia ante rem.}}

Da ingen Viden er mulig uden en Vidende, følger det af sig selv, at
Videns Oprindelighed maa forudsætte en oprindelig Vidende, at her
altsaa i logisk Forstand vises hen til Videns Princip, det
Oprindelige selv, den vidende Subjectivitet. Med den vidende
Subjectivitet er det Almene, det i sig Logiske, tillige sat som det
Overgribende; men hvorfra kommer saa det Antilogiske?
Naturligviis, ifølge vort Foregaaende, fra Andetheden, fra det
Enkeltes Umiddelbarhed, fra det Objective!
% --- p. 179 ---
\setcounter{footnote}{0}%
Bevægelsen i denne Kreds er desaarsag at bestemme som en Overgang
fra Selvhed gjennem Andethed tilbage til Selvhed, fra det Almene
gjennem det Enkelte tilbage til det Almene. At det Almene, der i
denne Bevægelse er det Første og det Sidste, maa opfattes som det
Evige, det i sig Selvstændige og absolut Gyldige, er i Overskriften
antydet ved det scholastiske, for Phasen betegnende Udtryk:
\textit{universalia ante rem.}

Betegnelsen \textit{ante rem} maa dog her ikke misforstaaes, som om
\textit{universalia} virkelig kunde være uden \textit{res}; her
vækkes tvertimod en logisk Forventning om \textit{res}, naar
\textit{universalia} ere forudsatte; saaledes vækkes der jo ogsaa,
logisk seet, en Forventning om Følgen, naar Grunden er forudsat, om
Virkningen, naar Aarsagen er forudsat, o. s. v.

Den næste Misforstaaelse, her maa afværges, er den, at et Udtryk
som: \textit{universalia sunt ante rem} skulde tyde paa en
umiddelbar Væren af det Almene, en evig Hvilen, uden Bevægelse og
uden Virken. Denne Misforstaaelse, hvori Middelalderens eensidige
Realister især gjorde sig skyldige, hæves derved, at \textit{res},
det existerende Enkelte, opfattes som det Mellemled, hvorved
Universalet, det sig i Alt bestemmende Almene, sætter og bestemmer
sin Forskjelseenhed med sig selv.

Endnu en forstyrrende Misforstaaelse maae vi see at fjerne, den
nemlig, at det Almene i sin selvvirksomme Idealitet skulde være
tomt og eensformigt. For at fjerne denne Misforstaaelse, er det
imidlertid ikke nok at holde logiske Lovtaler over det Almene,
erklære det for det evigt Værende, det Sande, det Guddommelige
o. s. v. Opgaven er at paavise, hvorledes det ved Begribningen
gaaer til, at det Almene, endskjøndt det overalt maa være det
Samme, kan være det i sig uendelig Indholdsrige, og, skjøndt i sig
selv det uendelig Indholdsrige, dog vedblive at berige sig i det
Uendelige. At sætte det Almene som det Absolute er i det Hele et
Særkjende for Idealismen; --- thi Middelalderens scholastiske
Realisme svarer just til den philosophiske Retning, som efter
% signature mark at foot: 12*
% --- p. 180 ---
\setcounter{footnote}{0}%
Nutidens Sprogbrug kaldes den idealistiske; --- her vil altsaa til
Opgavens Løsning fordres en Charakteristik af Idealismens logiske
Eiendommelighed.

Med Hensyn paa de ovenfor anførte Misforstaaelser er det især to
Former af eensidig Idealisme, der i nærværende Sammenhæng maae
tiltrække sig Opmærksomheden, nemlig den platoniske og den
hegelske. Med al den Forskjel, her iøvrigt er imellem dem, have
disse to Former af Idealisme dog det tilfælleds, at de begge
opfatte det Almene som det i sig Objective. Hvad den hegelske
Misforstaaaelse i dette Punkt
% printer's error: "Misforstaaaelse" with three a's, verified at 600 dpi
% against "Misforstaaelser" and "Misforstaaelse." elsewhere on this page.
angaaer, er Sagen allerede oplyst i vort Foregaaende. Her skulle vi
endnu blot kaste et Blik paa den platoniske Misforstaaelse.
„Forvisningen om Ideernes Realitet“ --- siger Poul Møller%
\footnote{Phil. Hist. S. 161---62.}
i Afsnittet om Plato --- „hviler i Forvisningen om Videnskabens
Realitet, der gaaer ud paa Tænkning af Tingenes Væsen. Der kunde
ikke gives Videnskab, dersom man med Protagoras og andre Sophister
antog, at de forsvindende sandselige Phænomener var det eneste
Sande. Ideerne ere det Samme som Tingenes Væsen. Men der gives
flere Ideer, fordi Eleaternes Lære om Eenheden, der udelukker al
Mangfoldighed, er urigtig. Dog gives der Eenhed, Sammenhæng i
Fleerheden af Ideer; thi den fornuftige Tænkning maa forbinde
Begreber i høiere almindeligere Eenheder, og til denne Sammenhæng
mellem Begreber maa ogsaa svare en Sammenhæng i det Reale. Ideerne
ere ikke isolerede Eenheder, der ere for sig uden Forhold til
hinanden. Deres indbyrdes Forhold bestaaer deri, at underordnede
Ideer sammenfattes i høiere som deres fælleds Eenhed. Den efter
Eenhed søgende Tænkning kræver da til sin fuldkomne
Tilfredsstillelse, at alle Ideer maae være forenede i een hoieste
Idee.“ --- Ideernes System er altsaa, ifølge denne Synsmaade, Eet
med de videnskabelige Begrebers System; den
% --- p. 181 ---
\setcounter{footnote}{0}%
Objectivitet, der maa tillægges det videnskabelige Begreb, som
Tænkningens Udtryk for Tingenes Væsen, bliver derfor hos Plato som
hos Hegel at tillægge det Almene. Da nu det, der objectiveres i
Viden, unegtelig med Hensyn paa det Almeengjældende er i ganske
anden Forstand objectivt end det, der umiddelbart opfattes ved
Sandsning og bestemmes ved den blotte Menen, saa er det alene heraf
forklarligt, at Modsætningen mellem Menen og Viden har kunnet
forlede Tænkere som Plato og Hegel til at forvexle det Almene med
det Objective, uagtet det Almene selv ikke blot er det Subjective,
men det eneste i sig absolut Subjective.

Hvor absolut subjectivt det Almene maa være, sees allerede af, hvad
der i en tidligere Sammenhæng er oplyst om dets Forhold til det
Negative og dets Eenhed med Subjectiviteten; thi er der Noget, der
er absolut subjectivt, maa det da vel være Subjectiviteten selv.
Paa Reflexionens Trin var det Reflexionen, den rene Reflecteren,
der var Subjectivitetens Form; paa Begrebets Trin er det nu det
Almene. Lader os derfor endnu engang besinde os paa Kategorien det
Almene; hvad betyder da egentlig det Almene? Man kunde sige: det
Almene er Eenheden, Tilværelsens sande, Alt omfattende Eenhed. Det
Almene maa altsaa være den Eenhed, der har det Mangfoldige i sig,
men saaledes, at den tillige har Mangfoldigheden og
Forskjelligheden imod sig. Det Almene maa, i kategorisk Betydning,
nødvendigviis tilfredsstille følgende Fordringer:

1) Som al Mangfoldigheds Eenhed maa det Almene gjennemtrænge det
Mangfoldige saaledes, at det i Alt og over Alt er Eet med sig selv;
det Almene maa omfatte og gribe det Mangfoldige saaledes, at det i
den hele Mangfoldighed fatter og griber sig selv. Men en saadan Alt
i sig selv og sig selv i Alt gribende og begribende Eenhed er en
\emph{Selveenhed}; det Almene er som Eenhed altsaa Selveenhed,
altsaa Selvhed, altsaa Subjectivitet.
% the two following occurrences of Selveenhed/Selvhed are NOT spaced

% --- p. 182 ---
\setcounter{footnote}{0}%
2) Som Tilværelsens uendelige Selveenhed maa det Almene i alle
Forskjelligheder og endelige Modsætninger nødvendig være den
\emph{absolute Identitet}. Det er kun i overfladisk Forstand, at
det Almene kan siges at være det Fælleds, det Lige i det ellers
Ulige, det Generelle. Opfattet som et blot Fælleds er det Almene
kun et tomt Schema, et Abstractum; men et Abstractum kan ikke være
det Almene, det er ved sin Eensidighed tvertimod kun noget
Særskilt.

3) Identiteten af Alt, den uendelig concrete Identitet, er altsaa
Selvhed, Subjectivitet. Men som Identiteten af Alt maa det Almene i
alle Ting jo overalt og paa alle Maader være det Samme, altid det
Samme; men denne Fordring kan alene opfyldes paa den Betingelse, at
det Almene udgaaer fra sig selv og gjennem alt Andet vender tilbage
til sig selv, at det, med andre Ord, er i \emph{uendelig
Bevægelse}, og at denne Bevægelse er Subjectivitetens egen
Ideebevægelse.

Men uagtet Plato og Hegel, af Interesse for Viden og det i Viden
Objective, begge, saa at sige, ere enige i at oversee Videns
Princip, den vidende Subjectivitet, saa er dog deres indbyrdes
Afvigelse, hvad Mystification af det Almene og dermed af Selvheden
angaaer, en nødvendig Conseqvens af deres indbyrdes Afvigelse i
Bestemmelsen af Andetheden. At bestemme, hvorledes \textit{res}
afhænge af \textit{universalia}, er i logisk Forstand at bestemme,
hvorledes Andetheden oprindelig forholder sig til Selvheden, thi
med Andetheden fremkomme \textit{res}. Hvad Rolle spiller
Andetheden da hos Plato?

„Andetheden udtrykker, at det Værende i uendelig Selvadskillelse
forholder sig til sig selv ved uophørlig at falde ud fra sig selv.
Det Andet er da ikke et Værende, men et sig i det Uendelige
Forandrende, en i sine Deles Mængde ubegrændset Masse (ὁ ὄγκος
αὐτῶν ἄπειρός ἐστι πλήθει). I denne Masse kan hver mindste Deel
for sig vel synes at være et Ene; men det er som i et Drømmesyn
(ὥσπερ ὄναρ ἐν ὕπνῳ): pludselig viser sig istedenfor det
tilsyneladende
% --- p. 183 ---
\setcounter{footnote}{0}%
Ene igjen det Mangfoldige, og istedenfor et Mindste et i Forhold
til den fortsatte Udstykning uendelig Stort. Med Andethedens
Udvikling i „Parmenides“ er allerede Sandselighedsmomentet antydet.
Dette Moment bliver særlig fremhævet i „Philebus“, der skjelner
mellem tre Arter af Væren, nemlig: Grændsen (πέρας), det
Ubegrændsede (ἄπειρον) og det af begge Sammensatte (τὸ ἐξ ἀμφοῖν
τούτοιν τι συμμισγόμενοι), hvortil endnu føies et fjerde Princip,
% printed "συμμισγόμενοι" (final iota, verified at 2400 dpi against the
% v-shaped nu of the preceding "-μενο-"); Plato has -ον.  As printed.
nemlig Blandingens Aarsag (αἰτία); ligeledes i „Timæus“, hvor der
skjelnes mellem det i Sandhed Værende, der ikke vorder, og det
altid Vordende, der ikke har nogen sand Væren (τὸ ὂν ἀεί, γένεσιν
δὲ οὐκ ἔχον, καὶ τὸ γιγνόμενον μὲν ἀεί, ὂν δὲ οὐδέποτε). Det
altid Vordende er hos Plato, hvad vi kalde Materien, en usynlig
Tingenes Amme (πάσης γενίσεως τιθήνη), som paa ubegribelig Maade
% printed "γενίσεως" (iota, positively distinct from the epsilon two
% letters earlier); Plato has γενέσεως.  As printed.
er bleven deelagtig i det Fornuftige (ἀνόρατον εἶδός τι καὶ
ἄμορφον, πανδεχές, μεταλαμβάνον δὲ ἀπορώτατά πη τοῦ νοητοῦ), den
Masse (ἐκμαγεῖον), hvoraf alle Sandsephænomenerne ere dannede. Fra
Materien, det Ikke-Værende, der hverken kan sandses eller tænkes,
men kun opfattes ved en Art uægte Slutning (μετ’ ἀναισθησίας
ἁπτὸν λογισμῷ τινι νόθῳ), have Tingene altsaa deres sandselige
Beskaffenhed; fra Ideerne have de deres Grundform og Realitet; som
Ideernes ufuldkomne Afbilleder i det Sandselige ere Tingene
Mellemled imellem Væren og Ikke-Væren.“%
\footnote{Phil. Propæd. S. 133---134.}

Det er af Ovenstaaende umiskjendeligt, at Platos Ideelære har i
Andetheden sit antilogiske Element (ἕτερον δέ γέ ποί φαμεν τὸ
ἕτερον εἶναι ἑτέρου, καὶ τὸ ἄλλο δὴ ἄλλο εἶναι ἄλλου); men det
% printed "ποί" (iota with acute, verified at 2400 dpi); Plato has πού.
% As printed.
er tillige øiensynligt, at dette antilogiske Element ikke betegner
en medbestemmende Magt, men tvertimod en reen Afmagt. Vistnok skal
denne Afmagt, ifølge den platoniske Synsmaade, ikke vedkomme
Ideerne i
% --- p. 184 ---
\setcounter{footnote}{0}%
deres evige Væren, men kun være charakteristisk for det i Tingene
Sandselige, Timelige, Ubestandige og Forgængelige. En nøiere
Betragtning lærer os imidlertid snart, at den Afmagt, der
udelukkende skulde falde paa Tingene, væsentlig falder paa Ideerne
selv. Just fordi Ideerne i sig selv mangle Andethedselementet, den
reale Negation, der skulde give dem Drift til Virksomhed og
Selvudvikling, ere de, som Aristoteles træffende har bemærket%
\footnote{Jvfr. Phil. Propæd. S. 137---138.}%
, uden Princip for Vorden og Bevægelse. De evige Ideer ere i denne
Henseende ikke blot ligegyldige mod de timelige Ting, men i deres
ideale Umiddelbarhed ere de endogsaa ligegyldige mod hverandre
indbyrdes. I denne blanke Umiddelbarhed er enhver Idee et Ubetinget
for sig: Lighed for sig, Ulighed for sig, Bevægelse for sig, Hvile
for sig o. s. v.; saaledes fremkommer en Mangfoldighed af
substantielle Enere (ἑνάδες, μονάδες), som, uagtet enhver af dem
er, med Hensyn paa de vexlende Phænomener, et Almindeligt (τό ἐπὶ
πᾶσι κοινόν) og har en af sandselige Ting uafhængig Væren (χωρὶς
ἔστι), dog udelukke hinanden gjensidig, umulig kunne gaae op som
Momenter i en eneste, absolut Eenhed. Vistnok fremstiller Plato
Ideernes Totalitet i deres hoieste Eenhed som Guddommen, ei som
Gud; her skulde vi da vente en høiere Belysning af Videns Indhold i
den Vidende selv. Men den naive Maade, hvorpaa Forholdet mellem Gud
og Ideerne overalt beskrives, viser tydeligt nok, hvorledes den
geniale Anskuelse træder i Begribningens Sted, og hvad der egentlig
er Grunden til, at den altid, hvor Talen er om Ideernes Idee, faaer
Overvægt over det Philosophiske, det Æsthetiske over det Logiske.
Det hedder i Republiken: „Ligesom Solen baade er Aarsag til, at
Tingene sees ved Lyset, og at de voxe og trives, saaledes er det
Gode (Ideernes Idee) Aarsag for Sjælen til Videnskab og Sandhed i
Alt, hvad der er Gjenstand for Videnskab. Ligesom Solen ikke er
% --- p. 185 ---
\setcounter{footnote}{0}%
Synet eller det Sete, men staaer over begge, saaledes er det Gode
ikke Sandhed eller Videnskab, men staaer over dem begge, idet de
begge ere beslægtede med det Gode.“ Det er unegtelig et træffende
Billede; men det er kun et Billede.

Da Ideernes logiske Væren oprindelig mangler det Magtens
antilogiske Element, der betinger Energien, kunne Kategorierne
Grund og Aarsag kun i billedlig Betydning anvendes paa dem.
Ligesom Andethedselementet, det Ubegrændsede (ἄπειρον), er for
afmægtigt til selv at gribe Ideerne, saaledes ere Ideerne, der
skulle afgive Grændsen, fra deres Side ogsaa for afmægtige til at
gribe det Ubegrændsede. Foreningen af Grændsen og det Ubegrændsede
kommer derfor istand paa en temmelig udvortes Maade. Foreningen
betragtes nemlig, hvad der udtrykkelig tyder paa et Udvortes og
Tilfældigt, under Synspunktet af en Blanding, og henføres til et
fjerde Princip, en Blandingens Aarsag, en Aarsag, der altsaa
vilkaarlig oversvæver sine egne Virkninger. Skjøndt Ideerne skulde
udgjøre Tingenes Væsen og følgelig i Væsen være Eet med Tingene,
have de dog i deres evige Væren et fra Tingene selv forskjelligt
Væsen; det er denne løse Forbindelse af Ideerne med Tingene,
Aristoteles dadler, idet han gjør gjældende, at Platonikernes
Ideer, uagtet de skulde være selvstændige for sig bestaaende
Eenheder, dog mangle særeget Indhold. „Den ideale Tohed og den
empiriske Tohed have samme Indhold. Man taler om Selvmennesket, om
Selvhesten, om Selvsundheden (αὐτὸ ἄνθρωπόν φασιν εἶναι καὶ
ἵππον καὶ ὑγίειαν); men ved at sætte Selv foran disse Hovedord
vindes Intet. Den virkelige Hest har væsentlig samme
Begrebsbestemmelser som Selvhesten; at definere det virkelige
Menneske er det Samme som at definere Selvmennesket o. s. v.“

Vende vi os nu mod den hegelske Idealisme, da har Ideen, den
absolute Negativitet i positiv Form, unegtelig nok af det Negative;
men dette Negative er, som vi alt have seet,
% --- p. 186 ---
\setcounter{footnote}{0}%
det abstract Logiske: Realnegationen, Begrebets antilogiske
Element, er ganske forsvunden. At den logiske Undermineren af
Andetheden, Enkeltheden, det i Objectiviteten Objective ikke saa
let mærkes i Logiken, er forklarligt, naar man betænker, at det
Surrogat, der stilles i den antilogiske Andetheds Sted, nemlig
Umiddelbarheden, den logiske Umiddelbarhed i sin
Vexelbestemmelse ved Negationen jo paa mange Maader kan frembyde en
skuffende Efterligning af det Objective selv. Hvorfra kommer vel al
Materie? Fra Umiddelbarheden! Og det i Begrebet Mechaniske? Fra
Umiddelbarheden! Og Modsætningen mellem Naturen og Aanden? Fra
Umiddelbarheden! Umiddelbarheden har just den Beqvemhed, at den kan
forestille det Sandselige og dog være det Tænkelige, angive det
Positive og dog forvandle sig til det Negative, sætte det Endelige,
det Udvortes, og dog ophæves i det Uendelige og reflectere sig som
et Indvortes.

Med disse Forudsætninger er den hegelske Ideelære i logisk
Henseende naturligviis ganske anderledes gjennemformet end den
platoniske. Istedenfor med Plato at lade ethvert nok saa abstract
Begreb optræde med ubetinget Selvstændighed som en Idee for sig,
gjør Hegel det hele System af eensidige Begreber til et System af
Momenter for Begrebet selv; han tildeler ikke de enkelte logiske
Idealformer en evig Tilværen i Ro og Hvile; tvertimod! han ordner
Kategorierne i store Kredse fra de allerabstracteste til de
allerconcreteste, han vækker i hver enkelt Kategorie den slumrende
logiske Drift, han viser ved sin immanente Methode, hvorledes de
følgende Kategorier udvikle sig af de foregaaende, som de høiere af
de lavere, de rigere af de fattigere, og forvandler paa denne Maade
hele Logiken til en Afspeiling af Ideens logiske Vorden. „I den
platoniske Fremstilling hersker det Subjective --- nemlig det
psychologisk Subjective --- forsaavidt der vilkaarlig antages snart
dette (εἰ ἓν ἐστι), snart hiint (εἰ ἓν μὴ ἔστι); en saa subjectiv
Form kan umulig være Ideens adæqvate Form. Den
% --- p. 187 ---
\setcounter{footnote}{0}%
hegelske Methode er derimod objectiv --- indenfor det abstract
Logiske --- og kan ved sin systematiske Udfoldning af Indholdets
Momenter afgive en objectiv klar Ideeform. Med sin anatomerende
Distinctionsdialektik kunde Plato ikke faae de enkelte Ideer til at
opgaae som Momenter i Ideens Eenhed; med sin immanente
Identitetsdialektik kan Hegel paavise Ideens absolute Eenhed i alle
særskilte Ideer, idet disse paa een Gang ere selvstændige og dog
Momenter.“%
\footnote{Phil. Propæd. S. 157.}

Men med disse formelle, for Philosophiens videre Udvikling store og
umiskjendelige Fortrin har den hegelske Ideelære i sit Væsen en
Huulhed og indvortes Tomhed, som kun til en given Tid kan tilsløres
af en i saamange Henseender aandrig og storstilet Formudvikling.
Grundskaden ligger i det Logiske, skjøndt den i Logiken selv er
mindst paafaldende, da den, som sagt, beroer paa en Udelukkelse af
det Antilogiske.

At Platos Ideer som \textit{universalia ante rem} ikke vare istand
til at magte de virkelige \textit{res}, er forstaaeligt af deres
Logik og tydeligt nok af Platos Physik. „Naturtingene ere et
Produkt af to Principer: Ideernes Totalitet, Gud, der er alle Tings
Fader, og det passive Princip, Tingenes Moder eller Amme, der
modtager Skikkelse af det dannende Princip. Naturen er ligesom et
Barn (ἔκγονος) af disse Tvende.“ Paa hvilken sindrig, poetisk
skjøn og barnlig Maade den hele Verdensbygning med alle
Naturgjenstandene, fra dette Synspunkt, lader sig opfatte som
Afbillede af Ideeriget, hvorledes, med andre Ord, den idealistiske
Physik kan forvandles til Natursymbolik, derpaa afgiver Platos
Timæus et mærkeligt Exempel. Men formaaer den platoniske Idealisme
ikke at meddele sine \textit{universalia} Evne til at magte
\textit{res}, og det af den gode Grund, at Realnegationen, det
Antilogiske selv, kun er erkjendt som et Billede paa Afmagten, da
kunne \textit{universalia} i den hegelske Idealisme, hvor stærkt
Dialektiken end,
% --- p. 188 ---
\setcounter{footnote}{0}%
om jeg saa maa sige, pidsker paa det Almene, endnu mindre forskaffe
de virkelige \textit{res} nogen Realitet, eftersom denne Idealisme
fra Roden af har bortskaaret det Element, hvorpaa det Reale selv i
al Realitet beroer, nemlig det Antilogiske. „En Philosoph, der
ringeagter Naturlærens udtrykkelige Paaviisning af Kildernes
Dannelse og Vulkanernes Virksomhed
% the "e" of "Vulkanernes" is all but unprinted (only ink fragments
% survive); reading certain from sense and from the fragments.
og med en Bemærkning som denne: „wie die Quellen die Lungen und
Absonderungs-Gefäße für die Ausdünstung der Erde sind, so sind die
Vulkane ihre \emph{Leber}, indem sie dieß Sich an ihnen selbst
Erhitzen darstellen“ mener at træffe den rene Sandhed i den rene
Idee, har intet Blik for Tingsaarsager og ingen Sands for
Virkelighed; den, der seer Bitterjorden i et saa speculativt Lys,
at han deri opdager
% "opdager" is badly battered in the forme; reading certain
„Saltets Subject, en Mellemsmag, der er blevet til Ildprincip“,
behøver ligesaa lidt at ændse, hvad der staaer i Chemien, som den,
der i Granitens metaphysiske Væsen seer sig istand til at paavise
„die irdische Dreieinigkeit, welche sich nun nach ihren
verschiedenen Seiten entwickelt“ ---, fordi det ligger i
„Granitprincipet“, at denne „Kern der Erde“ skal, med logisk,
dialektisk Nødvendighed, bestaae af Qvarts, Glimmer og Feldtspath.
--- behøver at raadføre sig med Geognosien.“%
\footnote{Forelæsn. over „Phil. Propæd.“ 1861---62 S. 71.}

Den hegelske Idealisme kan, forsaavidt den i Naturphilosophien
bliver nødt til at indlade sig med det sandseligt Concrete,
naturligviis ikke unddrage sig Fornemmelsen af, at her er et
antilogisk Element; men da Hegel i Logiken ingen Plads har for dets
Udvikling, da det rene Begreb og den rene Idee kunne blive
fuldkommen færdige uden at opløse dette Element, er det forsaavidt
nok, at dette Element maa i Naturriget selv blive betragtet som
noget Ureent, der let kan smitte den rene, speculative Tænkning.
Misforholdet imellem Idee og Natur bliver desaarsag i Hegels
Naturphilosophie beskrevet paa en Maade, der i logisk Forstand maa
vække
% --- p. 189 ---
\setcounter{footnote}{0}%
Betænkeligheder. „Naturen er i sig selv, i Ideen, guddommelig; men
saaledes som den er, svarer dens Væren ikke til Begrebet; den er meget
\textbf{mere} den \emph{\textbf{uopløste Modsigelse}}. Dens
% fat-face (halbfette) Fraktur: `mere' plain-width, `uopløste Modsigelse'
% fat-face AND letterspaced; calibrated against `meget' and `den' on the
% same line, which are markedly lighter.
Eiendommelighed er den satte Væren, det Negative; som de Gamle jo i det
Hele have anseet Materien for et \textit{non-ens}. Saaledes er Naturen
ogsaa en Udtalelse af Ideens Affald fra sig selv, idet Ideen som denne
Skikkelse af Udvorteshed er i en med sit Væsen uadæqvat og modstridende
Tilstand, o. s. v.“

Den sande, alsidig opfattede Idealisme maa nødvendig lægge an paa at
vise, hvorledes det antilogiske Element ikke som et Afmagtens, men
tvertimod som et Magtens Moment betinger Viden og muliggjør den reale
Begriben. Idet vi da, som Phasen fordrer, stille Begribningen selv og
den begribende Subjectivitets uendelige Selvudvikling som høieste
Formaal, betragte vi dette Formaal først fra det phychologiske, dernæst
% PRINTER'S ERROR, transcribed as printed: `phychologiske' for
% `psychologiske'. At 600 dpi the second sort is unambiguously h (bowl
% with a descending leg); calibrated against `psychologiske' four lines
% below, where the long s is straight and has no leg.
fra det ontologiske Standpunkt.

Hvad Magtmomentet gjælder, naar Fordringen lyder paa en bestandig
rigere Udvikling af et i sig uendeligt Begrebsindhold, kan fra den
psychologiske Subjectivitets Standpunkt sees af Sandsningen og det
Sandselige. Der er i det Sandselige ikke blot en Umiddelbarhed, men en
antilogisk Umiddelbarhed; det er fra denne antilogiske Umiddelbarhed,
at den uendelige Rigdom af Bestemmelser udgaaer, der alene kan
articulere Tænkningen og hæve Reflexionen til en virkelig Begriben.
Hvad er det nemlig, der gjør, at de mange Fornemmelser ikke lade sig i
Begrebets Navn reducere til en fad Almeenbestemmelse, og denne
Almeenbestemmelse igjen, efterdi den er sandselig, reducere til en
logisk Reflex i det logiske Almene, ligesom Nu og Her i den hegelske
Aandsphænomenologie? Aabenbart det antilogiske Element! Med
Fornemmelsen gives os jo i Virkeligheden en uendelig Mangfoldighed af
særskilte Fornemmelser; hver saadan Fornemmelse har i sig det
Dobbeltsidige, at den i samme Nu, som den
% --- p. 190 ---
\setcounter{footnote}{0}%
viser hen til en antilogisk Umiddelbarhed, der modstaaer Reflexionen,
just afgiver en Reflex, der paa eiendommelig Maade bestemmer og beriger
Reflexionen. Man tænke paa Lyden og Tonerne, man tænke paa Lyset og
Farverne, man tænke paa den punktuelle Vexelbestemmelse af Farvernes og
Tonernes æsthetiske Virkninger med de tilsvarende
mathematisk-physiske Functioner; man tænke paa den Vrimmel af
rationelle Vidensbestemmelser, der indeholdes i slige mathematisk
% The print sets `mathematisk-physiske' with the hyphen on the first
% occurrence and `mathematisk physiske' without it on the second, two
% lines below; both verified at 600 dpi. Transcribed as printed.
physiske Functioner; man tænke paa den uendelige Tilfredsstillelse, som
den Vidende har, der ved Udviklingen af alle Vidensbestemmelser
udvikler sig selv, bevarer Eenheden med sig selv, og forvandler de
utallige, ved Jagttagelser bestemte Realitetsindtrnk til
% PRINTER'S ERROR, transcribed as printed: `Realitetsindtrnk' for
% `Realitetsindtryk'. At 600 dpi the sort is n --- two arches, straight
% right stem, no descender --- while every certain y on this page
% (`Lys', `bryde', last line) carries a plain descending tail. All three
% OCR witnesses read n as well.
Begrebsbestemmelser i et stedse voxende Videns Indhold, og man vil
sikkerlig indrømme, at Tilværelsens antilogiske Element er afgjørende
for en fremadskridende Begrebsudvikling.

Indvender man paa Idealismens Vegne, at den rene Begriben paa denne
Maade jo maa blive belæmret med saamange Tilfældighedselementer,
saamange Indblandinger af Sandsebilleder og endelige Forestillinger, at
Ideen som Idee umulig kan speile sig i saa grove, forplumrede Media, da
bør man vel betænke, at Indblandingen af det Tilfældige og sandselig
Dunkle ingenlunde hidrører derfra, at her i Begribningen er et
antilogisk Element; men meget mere derfra, at dette antilogiske Element
beroer paa en Magt, der falder udenfor den endelige, psychologiske
Viden.

Seet fra et ontologisk Synspunkt gaaer Magtmomentet paa uendelig Maade
op i Vidensmomentet; den sande, Alt begribende Viden er en vidende
Magt, en almægtig Viden. Den oprindelige Eenhed af Viden og Magt i
Videns Form viser os her den begribende Subjectivitet i et System af
virksomme Almeenformer, overgribende \textit{universalia}, der sætte og
forudsætte \textit{res}, idet de overalt sætte og forudsætte reale
Media, hvori Ideens Lys kan bryde sine Straaler.

% --- p. 191 ---
\setcounter{footnote}{0}%
At \textit{res}, de umiddelbare Existenser i deres Enkeltheder og
uendelige Mangfoldighed, staae ved Magt, vil altsaa sige: 1) at
Tingene selv blive Gjenstande, ikke flygtige Skin, men virkelige
Gjenstande for den ideale Viden; 2) at den ideale Viden væsentlig
behøver slige virkelige Gjenstande, fordi Begribningens logiske Former
Punkt for Punkt maae integreres ved antilogiske Elementer; 3) at den
begribende Subjectivitet absolut magter Tingene, eftersom den
antilogiske Magtbestemmelse, der fra Magtens Side er Tingenes Realitet,
punktlig svarer til den logiske Magtbestemmelse, der, fra Begribningens
Side, er identisk med Idealiteten selv.

Paa disse Vilkaar er Phasen bestemt ved Videns og den vidende
Subjectivitets uendelige Energie; Ideebevægelsen gaaer fra Selvhed
gjennem Andethed tilbage til Selvhed, fra det Subjective gjennem det
Objective tilbage til det Subjective, fra det Almene gjennem det
Enkelte tilbage til det Almene. Men giver denne Kredsbevægelse nu ikke
en tom Cirkel, en ufrugtbar Gjentagelse af det Samme og det Samme i det
Uendelige?

I abstract Almindelighed er Ideebevægelsen jo vistnok en Gjentagelse;
men Gjentagelsen er i Begrebet en uendelig Fornyelse, fordi
Begribningen selv er en uendelig Vexelbestemmen af Erindren og
Frembringen, og den ideelle Vexelbestemmen altsaa Eet med en uendelig
Produktivitet. Alt Nyt er gammelt, „der gives intet Nyt under Solen!“
dette er den ideeløse, triviale Opfattelse af Gjentagelsens Kategorie,
en Opfattelse, der altid er slagen med logisk Ufrugtbarhed, efterdi den
er begreblos. I Forhold til en saadan begreblos Gjentagelse har Lessing
Ret, naar han etsteds udtrykker sig omtrent saaledes: „dersom Gud holdt
al Sandhed i sin høire Haand og Sandhedens Undersøgelse i sin venstre,
og bød mig vælge, jeg valgte den venstre.“ Var det muligt at besidde
Sandheden og det Sande i en død Væren, en dogmatisk afsluttet Form, da
maatte Ideen forsvinde, og Spotteren faae
% --- p. 192 ---
\setcounter{footnote}{0}%
Ret, naar han finder det kjedsommeligt at være en Alvidende, efterdi en
% The printed line breaks `Al | vidende' with NO hyphen at the line end
% (verified at 600 dpi to the right margin); the word is joined here.
Alvidende jo aldrig vilde kunne faae noget Nyt at vide. I Kraft af den
ideale Produktivitet er det Gamle derimod bestandig nyt; saaledes er
Gjentagelsen et væsentligt Moment i den oprindelige Begriben, hvis
Slutningspunkt overalt falder sammen med sit nye
Begyndelsespunkt\footnote{% printed mark: *)
I sin Afhandling: „Det astronomiske Aar“ har J. L. Heiberg (Prosaiske
Skr. 9de Bd. S. 79 flgd.) med særligt Hensyn paa Udtalelser af
S. Kierkegaard og Göthe belyst „Gjentagelsens Dialektik“ paa følgende
Maade. „Snart lyder den Klage, at Gjentagelsen er kjedsommelig, fordi
den bestandig bringer os det Samme, snart den modsatte Klage, at hvad
den bringer, er flygtigt og forgængeligt, og at den derfor idelig maa
sætte det Nye istedenfor det Gamle. Men denne Modsigelse ligger i selve
Gjentagelsen og udgjør dens Natur. Hvad den bringer, maa være
forgængeligt, thi ellers kunde det ikke gjentages; men det maa tillige
være uforgængeligt, thi ellers kunde det Nye ikke være det Samme som
det Gamle. Hvo som har Tilbøielighed til at see Alt i et hypochondert
Lys, kan da ved Naturens Gjentagelser baade sørge over, at Alt er
foranderligt, og ærgre sig over, at Alt er uforanderligt. Men paa denne
Maade lader man det Foranderlige og det Uforanderlige falde fra
hinanden i Betragtningen, holder dem adskilte som to selvstændige, af
hinanden uafficerede Egenskaber, og isaafald kan man vel have Grund til
at sørge over det Ene og ærgre sig over det Andet. Men med en sund
moralsk Constitution sætter man dem i indbyrdes Forbindelse, som
modsatte Sider af det samme Begreb, og har nu ligesaa gyldig Grund til
at glædes og oplives ved begge \dots{} Men hvad det fornemmelig kommer
an paa, det er, at Naturens objective Gjentagelse møder en saadan
Subjectivitet i Opfattelsen, som gjør den til sit Eget og bearbeider
den til noget Nyt. Thi hvad Naturen skjænker, er ligesom de gamle
Eventyrs Talismane, der have vidunderlige Kræfter, men ere moralsk
indifferente, saa at de blive til Velsignelse i den Enes Haand, i den
Andens til Fordærvelse.“}. For den psychologiske Subjectivitet i dens
menneskelige Begrændsning ere Videns Betingelser indskrænkede, derfor
% --- p. 193 ---
\setcounter{footnote}{0}%
kan den ideelle Produktivitet paa dette Trin saa let udtømmes, og det,
der skulde gjøre Nyhedens, d. v. s. Oprindelighedens Indtryk, tage sig
ud som en trættende Gjentagelse af det Gamle. For den ontologiske
Subjectivitet i dens principielle Eenhed med den guddommelige Idee ere
den produktive Videns Muligheder uudtømmelige; her bliver det
Forbigangne væsentlig forynget i det Nærværende; her er det Gamle evigt
et Nyt --- under Solen.

\subsubsection*{β) Magtens Oprindelighed: \textit{universalia post rem}.}

Kun igjennem Enkeltbestemmelsen kan det Almene beriges i det Uendelige;
men Enkeltbestemmelsen er i sin Umiddelbarhed en Magtbestemmelse:
saamange Vidensbestemmelser i Begrebet, saamange Magtbestemmelser i den
til Begrebet svarende Existens. Har nu det Enkelte i den reale
Begrebsudvikling samme Oprindelighed som det Almene, saa ligger heri en
Opfordring til at betragte Ideebevægelsen fra et Idealismen modsat
Synspunkt, søge dens Begyndelse og Slutning, ikke i det Almene, men i
det Enkelte, ikke i Viden, men i Magten, og derved see, hvorledes den
ved sig selv forvandler sig til en Realitetsbevægelse. Det Enkelte
danner altsaa i denne Sammenhæng ikke en Modsætning til det
Sammensatte, men til det Almene.

Det Enkelte er i sin existentielle Umiddelbarhed kun dette Enkelte,
ikke det Almene; og dog skal det Enkelte, just dette Enkelte, bestemmes
ved det Almene, opfattes, som om det virkelig var det Almene: her er
Problemet, det Problem, hvis eiendommelige Løsning Middelalderens
Nominalister betegnede ved Formlen \textit{universalia post rem}. Det
nominalistiske Standpunkt er væsentlig psychologisk; her begyndes med
en endelig Adskillelse af Subject og Object. Den erkjendende
Subjectivitet staaer som existerende Individ midt i en Verden af
existerende Gjenstande; hver Gjenstand er et Detteværende, et
Individuelt og Begrændset, ligesom Subjectiviteten selv.
% signature mark at the foot of p. 193: `Grundideernes Logik.' + 13
% (first page of the tenth gathering); not transcribed.
% --- p. 194 ---
\setcounter{footnote}{0}%
Erkjendelsen er betinget af, at en Fleerhed af enkelte Gjenstande
betegnes ved et fælleds Navn; dette er et Huus, dette og dette ligesaa
o. s. v. Men uagtet de mange enkelte Gjenstande kunne sammenfattes
under en Fælledsbenævnelse, saa har dog den ene Gjenstand i sin
existentielle Enkelthed Intet tilfælleds med den anden. Vel faaer det
Udseende af, at her maatte være noget Fælleds, noget Almindeligt, der
gik igjennem de Enkelte; men det er en Indbildning: det Fælleds, det
Almene vedkommer kun Benævnelsen, ikke Gjenstanden, har kun en nominel,
ikke en reel, kun en subjectiv, ikke en objectiv Gyldighed. De
existerende Gjenstande --- \textit{res ipsæ} --- ere da, ifølge denne
Synsmaade, ikke blot uafhængige af og ligegyldige imod, men ogsaa
ueensartede med de \textit{universalia}, hvorved de betegnes. At
\textit{universalia} ere \textit{post rem}, betyder, at vore Begreber
om Tingene vel ere abstraherede af Tingene, men at Tingen selv er
ligesaa ligegyldig mod sit Begreb, som Legemet mod sin Skygge. Langt
fra at høre blandt \textit{realia}, ere \textit{universalia} selv kun
\textit{nomina, flatus vocis}.

Det Modsigende i den nominalistiske Anskuelse kommer dog strax tilsyne,
naar Forholdet mellem Begrebet og dets Existens betragtes lidt nærmere.
Begrebet kunde jo umulig være en Fælledsbenævnelse for visse bestemte
Gjenstande, dersom der i Gjenstandene selv intet Fælleds var, der lod
sig fastholde ved Benævnelsen; Begrebet kunde umulig abstraheres af
Gjenstanden, dersom der ikke i Gjenstanden selv var Noget, der svarede
til Begrebet. Med den Erkjendelse, at Abstractionsbegrebet væsentlig
maa svare til de Gjenstande, hvoraf det er abstraheret, at det Almene,
der i Begrebet fremhæves som det for Gjenstandene Fælleds, ogsaa
virkelig findes hos Gjenstandene, at Enkeltexistenserne trods al deres
Enkelthed nødvendig maae gaae ind under det Almene, bestemmes og
charakteriseres ved det Almene, er Nominalismen forvandlet til
Empirisme.

% --- p. 195 ---
\setcounter{footnote}{0}%
Uden at indlade sig paa nogen egentlig Løsning af Nominalismens logiske
eller metaphysiske Problem om Forholdet mellem \textit{universalia} og
\textit{res}, gaaer Erfaringslæren ligefrem ud fra den Forudsætning, at
Tingene maae lade sig bestemme ved tilsvarende Abstractionsbegreber,
naar der, vel at mærke, ved nøiagtige Jagttagelser sørges for, at
Begreberne selv bestemmes ved Tingene. Den erkjendende Subjectivitets
hele Opmærksomhed er altsaa rettet mod Tingenes Eiendommelighed. Det
Eiendommelige i Tingene og det Almene i Begreberne afgive her
Betingelser for en saa rig Vexelbestemmelse af det Subjective og det
Objective, at en heel Verden af Kjendsgjerninger kan afspeiles i et
System af tilsvarende Abstractionsbegreber.

Den empiriske Begrebsudvikling begynder jo vistnok forsaavidt
theoretisk, som den begynder med at tage Tingene, som de nu engang ere
givne, og rette sine Tankebestemmelser efter de stedfindende
Existensbestemmelser; men Subjectivitetens eiendommelige
Selvvirksomhed kan ikke udeblive. Det viser sig snart, at de Begreber,
der paa eiendommelig Maade ere blevne bestemte ved Tingene, reagere
bestemmende for Tingene, og hævde det Almene en ideel Overlegenhed over
det Enkelte. Naar Empirismen paa denne Maade er udviklet til Realisme,
ere Begreberne i Realiteten blevne til noget ganske Andet end subjective
Navneværdier; de ere blevne til Betegnelser for substantielle Former,
for objective Magter, for de uforanderlige Love, hvilke Tingene ere
underlagte; Resultatet er, at det Almene behersker det Enkelte. Paa
Realismens Standpunkt kan dette imidlertid ikke forstaaes saaledes, at
det Enkelte paany skulde forvandles til en Vidensreflex, et Realmoment
i den almene Viden; tvertimod! Realismen lægger Eftertrykket paa
\textit{res} og derved paa Realiteten af det Enkelte som Enkelt.
Begrebsbevægelsen er altsaa denne: det Enkelte bliver formedelst det
Almene at bestemme i dets Enkelthed; ved denne Bestemmelse af det
Enkelte i dets Enkelthed faaer den begribende
% signature mark at the foot of p. 195: 13* (third page of the tenth
% gathering); not transcribed.
% --- p. 196 ---
\setcounter{footnote}{0}%
Subjectivitet det Objective paa objectiv Maade i sin Magt:
Realerkjendelsen er en praktisk Erkjendelse.

„Vi lære“, siger Naturforskeren, „at beherske Naturen ved vor stedse
voxende Indsigt i Kræfternes Sammenhæng, ved vore Opdagelser af
Naturlovene“; og det maa indrømmes, at Naturlæren har i det sidste
Aarhundrede ført et storartet Beviis for Sandheden i den gamle Sætning,
at Viden er Magt. Men den i Naturlæren stedfindende Begrebsudvikling
kan i denne Sammenhæng tillige tjene til Beviis for Umuligheden af at
opnaae fuldstændig Indsigt i Forholdet mellem Begreb og Existens,
saalænge den vidende Subjectivitet ikke kan opgive sit psychologiske
Standpunkt.

Lægge vi Mærke til den Maade, hvorpaa Erfaringsvidenskaben i det Hele
bearbeider sit Stof og udvikler sine Realbegreber, da frembyder sig her
% On this page several line-break hyphens are dropped or barely inked in
% the print (`Real | begreber', `indivi | duelle', `Gjen | stande');
% verified at 600 dpi. The words are joined here.
en naturlig Adskillelse mellem Billedet, Beskrivelsen, Charakteristiken
og det egentlige Begreb. Billedet, den umiddelbare Betegnelse af det
Enkelte i dets Enkelthed, forestiller jo Gjenstanden, Realbegrebet i
dets individuelle, antilogiske Form. Istedenfor at frembringe
Gjenstanden selv maa den vidende Subjectivitet, naar vi tage Opgaven
fra den theoretiske Side, altsaa indskrænke sig til at efterligne den
ved et Billede. Idet Billederne af flere ligeartede Gjenstande optages
i Forestillingskredsen og sættes i Bevægelse, bearbeides de af
Reflexionen, og afgive Elementer for Beskrivelsen. I Naturbeskrivelsen,
der allerede fordrer en Adskillelse af Væsentligt og Uvæsentligt, sees
Overgangen til Charakteristiken; Charakteristiken forener Fremstillingen
af det Individuelle med Fremstillingen af det Almene paa en saa
eiendommelig Maade, at en egentlig Begriben derved kan siges at være
indledet. Hvorpaa beroer saa den egentlige Begriben? Aabenbart paa
Indsigt i de Analogier og causale Forhold, hvorved en
Sammenhængseenhed med dens forskjellige Forgreninger lader sig
bestemme.

% --- p. 197 ---
\setcounter{footnote}{0}%
At Begrebet som Begreb paa denne Maade maa blive fattigere end den
tilsvarende Existens, er en Selvfølge. Hvorvidt Undersøgelserne end
drives, vil det intellectuelle Fremskridt altid blive en Tilnærmelse.
Og naar Fagvidenskaben med sine Realundersøgelser ikke kan komme ud
over den blotte Tilnærmelse, da kan Philosophien, hvad
Begrebsconcretionerne angaaer, det heller ikke. „I sin meest endelige
Realitet er Sagen et umiddelbart Tilværende og dette Tilværende er i
Forhold til sit Begreb et Exemplar. I denne tilværende Steen have vi
saaledes et Exemplar af en Steen, i denne Plante et Exemplar af en
Plante, i dette Dyr et Exemplar af et Dyr, o. s. v. I Sagen som
Exemplar have vi Begrebets Indhold, men ikke dets Form; i Sagen som
Begreb have vi Formen, men ikke det virkelige Indhold. Imellem Sagen
som Exemplar og Sagen som Begreb ligger nu Exemplet. I Modsætning til
Exemplaret er Exemplet det Almene og forsaavidt beslægtet med
Begrebet; i Modsætning til Begrebet er Exemplet det Individuelle og
forsaavidt beslægtet med Exemplaret. Skal her altsaa kunne gives en
reel Eenhed af Sagudvikling og Begrebsudvikling, da maa en saadan
Forskjelseenhed nødvendig være betinget af en
Exempeludvikling.“\footnote{% printed mark: *)
Forelæsn. over „Phil. Propæd.“ 1861--62 S. 41.} Men
Exempeludviklingen underligger en uendelig Tilnærmelse.

Tage vi Opgaven fra den praktiske Side, vil en lignende Relativitet
vise sig. Hvad Constructionsbegrebet betyder og hvorledes det forholder
sig til den antilogiske Dom: \emph{det Almene er det Enkelte}, er i det
Foregaaende allerede omhandlet. Her skal endnu tilføies en Angivelse af
Relativitetens Mellembestemmelser. Er Abstractionsbegrebet
ufuldkomment, maa Constructionsbegrebet naturligviis ogsaa være det.
Det Materiale, der skal benyttes, maa jo forefindes som et Givet; dette
Givne har da i sin existentielle Concretion en
% Between `sin' and `existentielle' the print carries a small single
% stroke at x-height. It is NOT the Fraktur hyphen sort, which is a
% double stroke (cf. `mathematisk-physiske', p. 190), so it is read as
% dirt or a damaged sort and no hyphen is transcribed.
% --- p. 198 ---
\setcounter{footnote}{0}%
Mangfoldighed af væsentlige Egenskaber, hvoraf kun nogle enkelte hæves
frem og fastholdes i tilsvarende Begrebsbestemmelser. Betingelserne og
de nærmere Omstændigheder, under hvilke Frembringelsen foregaaer, ere
her underkastede saamange Afvexlinger, Ændringer og Forandringer, at
der vilde behøves et uendeligt Antal af Forsøg, inden et enkelt Stofs
Egenskaber, endog i en ganske særegen Kreds af Vexelvirkninger, f. Ex.
den chemiske, fuldstændig kunde bringes for Lyset. Den ved
Naturbegrebet bestemte Frembringelse, f. Ex. Frembringelsen af Vand,
svarer jo vistnok til Begrebet selv, forsaavidt her blot tages Hensyn
til den Existenslov, der paa givne Betingelser opfyldes ved
Frembringelsen; men hvilken Afstand er der dog ikke imellem den
Indsigt, der er tilstrækkelig for at bringe Ilt og Brint til at forene
sig til Vand, og den Indsigt, der vilde udkræves for at gjennemskue
enten Iltens eller Brintens Anlæg til chemiske Forbindelser med alle
mulige Stoffer under alle mulige Betingelser! Alene Betragtningen af et
enkelt, indskrænket Exempel kan næsten bringe den efter strenge
Begreber higende Forstand til at svimle. „Med Antallet af Elementerne,
med Antallet af Atomerne, der forenes til en Gruppe“, siger Liebig,
„mangfoldiggjøres Tiltrækningens Retninger; hvisaarsag Tiltrækningens
Styrke da ogsaa aftager i samme Grad, som Retningernes Mangfoldighed
tiltager. En Kogsaltdeel, en af de mindste Zinnoberdele fremstiller en
Gruppe af ikke mere end to Atomer; et Sukkeratom derimod indeholder
sexogtredive, den mindste Olivenoliedeel flere hundrede Enkeltatomer
\dots{} Uden at Noget kommer til eller tages fra, kunne vi tænke os
Sukkeratomets sexogtredive Enkeltatomer ordnede paa tusinde
forskjellige Maader; med enhver Forandring ved et eneste af disse
Atomers Leie ophører det sammensatte Atom at være et Sukkeratom
o. s. v.“

Naturlærens Begreber ere i den Grad Facticitetsbegreber, at
Videnskaben i Grunden helst undgaaer ethvert Forsøg med
% --- p. 199 ---
\setcounter{footnote}{0}%
at forklare sig Virksomhedernes indre Hvorledes, tilfreds, naar den
blot kan paavise de factisk gjældende Love; et Beviis for, at
Almeenbestemmelserne selv dog væsentlig opfattes som givne
Kjendsgjerninger.

Er Idealismen i det Hele tilbøielig til at forvexle sin abstracte Viden
af det Guddommelige med den guddommelige Viden selv, saa er Realismen
med sine vidtforgrenede Kundskabsmasser og sin mangfoldige Viden af det
Endelige tilbøielig til at negte Muligheden af en sand Idee-Erkjendelse.
Men begge Eensidigheder ville forsvinde, eftersom Forholdet mellem den
psychologiske og den ontologiske Subjectivitet nærmere bliver belyst.
Det er Videns Opgaaen i Magten, hvorpaa det i nærværende Sammenhæng
kommer an. Ligesom det fuldkomne Vidensbegreb, i Begribningens første
Phase, forudsatte det fuldkomne Magtbegreb, saaledes forudsætter nu
Magtbegrebet her, i Begribningens anden Phase, sit fuldkomne
Vidensbegreb.

Dersom Magtbegrebet imidlertid var saaledes Eet med sit Vidensbegreb,
at her i denne Eenhed ingen Forskjel var, saa vilde det være en
temmelig ligegyldig Formalitet, om vi foretrak den Betegnelse, at det
er Magten, der gaaer op i Viden, eller den omvendte, at det er Viden,
der gaaer op i Magten. Men ogsaa her gjælder det, at den totaliserede
Eenhed bestaaer i en totaliseret Modsætning. Den uendelige Tilnærmelse,
hvorved Realbegrebet paa den menneskelige Videns Standpunkt kan
fuldkommengjøres til stedse større Overeensstemmelse med sine
Exemplarer, tyder ingenlunde hen paa en fuldendt Begriben, der skulde
forudsætte, hvad den hegelske Logik saa falsk forudsætter, at det
adæqvate Begreb umiddelbart faldt sammen med sin Existens. Tvertimod!
Realbegrebernes trinvise Udvikling beviser just, at f. Ex. Begrebet
Hest kunde i alle Bestemmelser være fuldkommen adæqvat med Existensen
Hest, og Modsætningen af Begreb og Existens
% --- p. 200 ---
\setcounter{footnote}{0}%
dog ligefuldt vedblive, forsaavidt Begrebet som Begreb er logisk,
Existensen som Existens antilogisk.

I den fuldkomne Begriben, i Ideen selv, danner Vidensbegrebet, Begrebet
som Begreb, altsaa en gjennemgaaende Modsætning til Magtbegrebet,
Begrebet som virkelig Existens; og dog er Vidensbegrebet i Ideen Eet
med Magtbegrebet. Modsætningseenheden maa da udvikle sig i en
Bevægelse. I Begribningens første Phase, hvor det Antilogiske blev
Moment for det Logiske, var Bevægelsen en reen Ideebevægelse; her i
Begribningens anden Phase, hvor det Logiske bliver Moment for det
Antilogiske, maa Ideebevægelsen forvandle sig til en Realbevægelse.

Hvilken Forvirring der opkommer, naar Ideebevægelsen og Realbevægelsen
ifølge Hegels immanente Methode uden videre slaaes sammen under Navn af
„Sagbevægelse“, er allerede oplyst fra mere end een Side:
„Mystificationen med „Sagbevægelsen“, „Begrebets objective
% The quotation opened here with „ runs on to p. 201, where it closes
% with “ after `det Værende.“' --- hence the quote-mark imbalance on
% this page (4 low, 3 high).
Selvbevægelse“, hidrører navnlig ingenlunde, som Saamange have antaget,
derfra, at Methoden er dialektisk, men meget mere derfra, at den
methodiske Dialektik kun er eensidig gjennemført. Dialektiken af Begreb
og Begreb kjender Methoden ypperligt, men Dialektiken af Værens Begreb
og det Værende kjender den kun maadeligt: istedenfor at udvikle Tingens
Begreb, det Begreb, der gaaer op i Tingen, udvikler Methoden overalt
Begrebets Begreb, Begrebet af Tingens Begreb, det Begreb, hvori Tingen
selv ikke er. Et Par Exempler ville oplyse Sagen bedre end vidtløftige
Forklaringer. For at Væren kan være, maa et Værende være; for at et
Værende kan være, maa der være et Detteværende. Grundstoffer som Ilt,
Svovl, Kalium o. s. v. ere hver for sig et Værende, derved at de, hvert
for sig, ere et Detteværende. Af Udtrykket \emph{Detteværende} sees,
at det Værende selv kun lader sig paapege, ikke udsige. Men det, der
ikke lader sig udsige, lader sig jo heller ikke tænke. Skal nu Eenheden
af den tænkte Væren og det Værende selv
% --- p. 201 ---
\setcounter{footnote}{0}%
ikke desto mindre retfærdiggjøres derved, at det Ikke-Tænkelige bliver
et negativt Element, en Tankens Grændse, som overalt hører med i den
bestemte Tænkning, da følger jo heraf, at Eenheden af Tænken og Væren
maa være ligesaa dialektisk, som Eenheden af Væren og det Værende.“

„Men er Dialektiken af Værens Begreb og det Værende selv forskjellig
fra Dialektiken af Værensbegreberne indbyrdes, af Væren og Ikke-Væren,
af Vorden og Tilværen, saa følger da ogsaa heraf, at Tankebevægelsen
maa være forskjellig fra Værensbevægelsen. Den \emph{Vorden}, hvorom
Logiken handler, idet Væren forener sig med Intet, saa at Tilværen
vorder, er ueensartet med den Vorden, hvorom Chemien handler, naar
Svovl forener sig med Ilt, saa at Svovlsyren vorder, Kalium med Ilt,
saa at Kali vorder, og Svovlsyren igjen med Kali, saa at det svovlsure
Kali vorder.“\footnote{% printed mark: *)
Forelæsn. over „Phil. Propæd.“ 1861--62 S. 55--56.}

Ved at tydeliggjøre Modsætningen mellem Begreb og Existens som en
eiendommelig Modsætning mellem Vidensbegrebet og Magtbegrebet lære vi
at indsee, hvad der fra et realistisk Standpunkt nødvendig hører til en
fuldkommen Begriben. Det er for det Første 1) Tingen med dens
Egenskaber, dernæst 2) det Udvortes og Tilfældige i Tingenes dog
væsentlige Sammenhæng, og endelig 3) Modsætningen mellem det Værende og
det Vordende, det Hvilende og det Bevægede, hvorpaa Opmærksomheden her
maa fæste sig.

1) Hvad Tingen med dens Egenskaber angaaer, maa den naturligviis,
forsaavidt den sees under Begrebets Synspunkt, opfattes som en
Totalitet af Existensbestemmelser. At denne Totalitet, naar den
udtrykker Magtbegrebet, ikke alene maa fremstille sig i en anden
Belysning, men tillige frembringes paa en ganske anden Maade, end naar
den skal udtrykke Vidensbegrebet, vil indsees ved nøiere Eftertanke.

% --- p. 202 ---
\setcounter{footnote}{0}%
Stille vi Sølvets Begreb mod det existerende Sølv, saa maa Identiteten
af denne Modsætning, dersom Begribningen skal være fuldkommen,
nødvendig haves i to Former: i Begrebets Form som Vidensbegreb, i
Existensens Form som Magtbegreb. See vi nu paa Vidensbegrebet, da yder
Mineralogien os jo ved sin Fremstilling af Sølvets Egenskaber Bidrag
nok, idetmindste til et Tankeexperiment, som kunde antyde, hvorledes
Begrebet Sølv vel maatte udvikles, naar det i alle Henseender skulde
tilfredsstille Videns Fordringer. Enten vi tænke paa Hvidheden og
Glandsen eller paa Klangen eller paa Krystallisationsformen eller paa
den svage Tiltrækning til Ilten, den stærke Tiltrækning til Chloret,
eller paa hvilkensomhelst af Sølvets Egenskaber, vi end ville nævne saa
vil, fra Videns Synspunkt, Indsigten i enhver af disse Egenskaber være
betinget af, at det Enkelte gaaer op som Moment i Bestemmelsen af det
Almene. Sølvets Glands forudsætter, for at blive til en
Begrebsbestemmelse, en Viden af, hvad Metalglands er; men en Viden af,
hvad Metalglands er, faaer kun Betydning for Begribningen under den
Forudsætning, at her allerede er en Viden af, hvad Glands i det Hele
væsentlig er, altsaa en Viden, der kunde gjøre Vexelvirkningen mellem
Lyset og Legemerne fuldkommen gjennemsigtig, en Viden, der først vilde
være fuldendt, naar Optiken var fuldendt. Tage vi dernæst Affiniteten
for os, da maa Sølvets Tiltrækning til Chloret nødvendig reflectere sig
som særlig Begrebsbestemmelse i Bestemmelsen af dets Tiltrækning til
Svovl, til Brom, til Jod, til Fluor, Cyan o. s. v. Bestemmelsen af de
enkelte Tiltrækninger vil altsaa i Begrebet falde sammen med en
Bestemmelse af Sølvets chemiske Tiltrækning i dens concrete
Almindelighed, men denne concrete Almindelighed vil igjen være et
særskilt Moment af den almene chemiske Tiltrækning i dens gjennemførte
Articulation, dens alsidige Concretion. En fuldkommen Indsigt i Sølvets
Tiltrækning til Chloret vilde forudsætte en fuldkommen Indsigt
% --- p. 203 ---
\setcounter{footnote}{0}%
i det hele System af chemiske Tiltrækninger, altsaa en fuldkommen
chemisk Viden.

Den fuldkomne Viden gaaer altsaa --- thi det er Videns Lov --- hvad
Existensbegreberne angaaer, \emph{fra det Enkelte til det Almene,
gjennem det Enkelte til det concret Almene}, ligesaa vel som den
ufuldkomne. Men denne Bevægelse, der er saa afgjørende for
Vidensbegrebet, passer aldeles ikke for Magtbegrebet. Idet Viden
støtter sig til Magten, forudsætter den Existensen, dens logiske
Bestemmelser ere tilsammen Reflexer af de antilogiske
Existensbestemmelser; de individualiserede Almeenbegreber ere i denne
Kreds \textit{univer\-%
salia post rem}. Anderledes med Magtbegrebet. Forsaavidt
Magtbestemmelserne oprindelig ere identiske med Vidensbestemmelserne,
har Magtbegrebet i Systemet af disse Vidensbestemmelser en skjult
Viden, et logisk Indhold, og kan med Hensyn paa sit ideelle Udspring
vel siges at begynde med et System af \textit{universalia ante rem}.
Men Magtbegrebets Charakteer er, som vi allerede have seet, jo netop
kjendeligt derpaa, at det forvandler den logiske Dom: \emph{det Enkelte
er det Almene}, til den antilogiske: \emph{det Almene er det Enkelte,
ja dette Enkelte}. \textbf{Medens Begrebet i Videns Form} begynder med
% halbfet Fraktur „Medens Begrebet i Videns Form“, kalibreret mod plant
% „Videns Lov“ og „Vidensbegrebet“ paa samme Side (600 dpi)
at forudsætte det existerende Object som et Dommens Subject, og saa
udvikler sig ved at udvikle Systemet af alle Prædicaterne, begynder
Magtbegrebet omvendt med at forudsætte et System af Prædicater, og
træder saa i Kraft ved at lade disse Prædicater concentrere sig i deres
Subject som i et existerende Object. I det existerende Object
forvandles Prædicaterne til Existensbestemmelser, til Egenskaber ved
Tingen, til Attributer og til Accidenser for Substansen. Disse
Accidenser, disse Egenskaber og Attributer kunne nu i Magtbegrebernes
Sphære ikke mere forvandles til Prædicater, til logiske
Almeenbestemmelser; deres Charakteer er i den Grad antilogisk, at
ethvert Forsøg paa at opfatte dem som \textit{universalia} maa tabe sig
i det Urimelige. Som
% --- p. 204 ---
\setcounter{footnote}{0}%
Magtbegreb existerer Sølvet i et System, ikke af logiske Prædicater,
men af virkelige Egenskaber, individuelle Attributer, factiske
Accidenser, Existensbestemmelser saa antilogiske, at hver enkelt af
Egenskaberne er ligesaa umiddelbart Eet med, som umiddelbart
forskjellig fra alle de øvrige. Her kan nu ikke tænkes paa Sølvets
Glands i al Almindelighed, ikke paa Metalglands i særskilt
Almindelighed, ikke paa Sølvglands i individuel Almindelighed; nei, her
kan slet ikke tænkes paa Glands i nogensomhelst Art af Almindelighed,
men alene paa denne Glands af dette existerende Sølv. Det existerende
Sølv med denne bestemte Glands har, ikke en almindelig Tiltrækning til
Chlor i al Almindelighed, men under \emph{disse Betingelser denne
Tiltrækning} o. s. v. I det existerende Sølv have den bestemte Glands
og den chemisk bestemte Tiltrækning deres factiske Identitet; men denne
Identitet \textbf{er} saa antilogisk, at ingen Guddom, endnu mindre
noget Menneske skal paatage sig at deducere \emph{denne Glands af denne
Tiltræknig}, eller
% TRYKFEJL: „Tiltræknig“ som trykt (tabt n; for „Tiltrækning“) --- utvetydigt
% ved 600 dpi; det samme Ord staaer rigtigt to Linier ovenfor og een Linie nedenfor
\emph{denne Tiltrækning af denne Glands}\footnote{At her fra
Videnskabens Standpunkt naturligviis kan være Tale om ---
idetmindste til en vis Grad --- at forklare den ene Egenskab ved den
anden, kan ikke komme i Betragtning, saalænge her udelukkende tilsigtes
en Bestemmelse af Magtbegrebets Charakteer. Ved en Forandring af
Synspunkt forvandles Egenskaberne igjen til Prædicater, og Reflexionen
tilegner sig et System af \textit{universalia post rem}.}. At begribe
er da i nærværende Sammenhæng ikke indskrænket til at tænke eller
reflectere; at begribe er: i skjult Eenhed og factisk Forstaaelse med
en uendelig Tænken, en uendelig bestemt Reflecteren, at
% Spatieringen er inconsequent over Linieskiftet: „at“ i Linieenden staaer
% uspatieret, Resten af Rækken spatieret. Gjengivet som trykt.
\emph{objectivere, at gribe, at magte.}

2) \textbf{Af Forholdet imellem Vidensbegrebet og Magtbegrebet forklares
nu Vexelbestemmelsen af Tingenes} ydre og indre, tilfældige og
nødvendige Relationer paa en meget naturlig Maade. Som existerende
Objecter træde Tingene saaledes op
% --- p. 205 ---
\setcounter{footnote}{0}%
mod hinanden, at det ved en overfladisk Betragtning seer ud, som
\textbf{om} den ene Ting aldeles ikke vedkom den anden; enhver Ting
\textbf{er}, hvad den \textbf{er}; her støder det Logiske igjen paa det
Antilogiske. Men det maa bestandig fremhæves, at Tingene som
Magtbegreber dog i Grunden ere Begreber; at Magtbegrebernes Indhold er
logisk, skjøndt Formen er antilogisk; at det Enkelte jo ikke har
løsrevet sig fra det Almene, men tvertimod fremstillet sig som det
Almenes egen individuelle Bestemthed, en Bestemthed, uden hvilken det
Almene selv ikke paa nogen Maade kunde faae Realitet i det Objective:
trods Tingenes ophævede Sammenhæng er her altsaa dog en virkelig
Sammenhæng. Denne Tingenes Sammenhæng frembyder nu en dobbelt Side, en
antilogisk og en logisk. Seet fra den antilogiske Side er Sammenhængen
\emph{udvortes}; i den udvortes Sammenhæng med dens vexlende
Omstændigheder hersker Tilfældigheden: Tilfældigheden er just den
individuelle Bestemtheds og den virkelige Sammenhængs antilogiske
Element. Fra den logiske Side er Sammenhængen \emph{indvortes}; i den
indvortes Sammenhæng med dens væsentlige Betingelser hersker
Nødvendigheden, med Nødvendigheden Lovene og med Lovene det System af
\textit{universalia}, der i Kjendsgjerningernes Verden erkjendes som
\textit{objective} Magter. I Lovenes Realitet har Magten sin Idealitet;
derfor ere Lovene væsentlige Mellembestemmelser mellem Vidensbegrebet
og Magtbegrebet. Kun i Videns Form \textbf{er} Loven Lov, thi Loven
\textbf{er} det Almene; kun i Magtens Form er Loven Lov, thi Lovens
Gyldighed \textbf{er} dens Opfyldelse, Opfyldelsen betinges af
Existensen, Existensen af en begribende Magt.
% Halbfet „er“ er inconsequent sat: bolden i „Videns Form er“, i „Loven er det
% Almene“ og i „er dens Opfyldelse“, men PLAN i den mellemste, ordlydende
% samme Vending „kun i Magtens Form er Loven Lov“. Kontrolparret staaer paa
% to Linier over hinanden og er utvetydigt ved 600 dpi. Gjengivet som trykt.

3) \textbf{Modsætningen imellem Realbevægelsen og den rene
Ideebevægelse} kan da heller ikke længer være tvivlsom. Som det
fuldkomne Vidensbegreb har Begrebet ideel Tilværen i en mægtig
Reflecteren; den mægtige, alle Begreber begribende Reflecteren
\textbf{er en med Hvilen identisk} Bevægelse, en Bevægelse i den rene
Idee. Fra den menneskelige Videns Standpunkt
% --- p. 206 ---
\setcounter{footnote}{0}%
kan en saadan Bevægelse sees i en systematisk Udvikling af de
ontologiske Kategorier, skjøndt en slig Udvikling naturligviis er
uendelig langt fra at congruere med den Begrebsbevægelse, hvori Ideen
selv udvikler og bestemmer sit logiske Indhold. Som eiendommeligt
Magtbegreb er Begrebet en Ting i Tingenes Række, og som existerende
Ting udtrykkelig bestemt for en endelig Modsætning af Hvile og
Bevægelse. For at belyse denne endelige Modsætning maae vi ikke
indskrænke os til at betragte det Mechaniske alene: enhver Naturproces
er et i sin Art betegnende Exempel. Realitetsbevægelsen forudsætter
Rummet og Tiden; hvad der i Videns Form er forenet i reen Reflexion,
udlægges i Magtens Kreds formedelst Udvorteshedens Former ved Systemer
af virkelige Overgange. Det gjælder for Begribningen om at kunne i
disse Overgange erkjende de existentielle Udtryk for en i Grunden
væsentlig og uendelig Reflecteren.

„Existensens Strid med Grunden er en Modsigelse, som ved Hjælp af
Udvorteshedens Former: Rummet og Tiden finder sin Løsning. Spørge vi,
hvad Vandet som Stof væsentlig er, da er Vandet et i Grunden
reflecteret Væsen; thi Vandet er jo som Produkt af Ilt og Brint
reflecteret i sine Bestanddele og disse Bestanddele omvendt
reflecterede i deres Produkt. Det saaledes reflecterede Produkt er
derhos --- i Grunden --- behæftet med den Modsigelse, at det ikke kan
være til, forsaavidt dets Bestanddele have Tilværen, og Bestanddelene
ikke være til, forsaavidt Produktet har Tilværen: Produkt og
Bestanddele forudsætte og ophæve hinanden gjensidig. Endvidere er dette
i Grunden reflecterede Produkt behæftet med den Modsigelse, at dets
Aggregationsform baade er luftformig, draabeflydende og fast, medens
enhver af disse Former dog gjensidig udelukke hinanden \dots\ Det er
umuligt at sige, hvad der meest charakteriserer Vandets Væsen, enten
dets Dannelse af eller dets Adskillelse i Ilt og Brint; thi i Grunden
maa Adskillelsen jo være reflecteret i Dannelsen, og Dannelsen omvendt
i Adskillelsen: Adskillelse og Dannelse ere med deres Forskjel i
% --- p. 207 ---
\setcounter{footnote}{0}%
Grunden Eet. Anderledes med Existensen. I den umiddelbare Existens kan
eet og samme Produkt umuligt være paa een Gang baade dannet og opløst;
men hvad der ikke kan skee paa een Gang, kan jo skee efterhaanden,
derfor have vi Tiden; og hvad der ikke til samme Tid kan skee paa eet
Sted, kan jo samtidigt skee paa forskjellige Steder, derfor have vi
Rummet.“\footnote{Forelæsn. over „Phil. Propæd.“ 1860--61 S. 173--74.}
Det er Magtbegreberne, der udfylde Tiden, idet de opfylde Rummet.

\subsubsection*{γ) Magtens og Videns oprindelige Eenhed:\\
\textit{universalia in re.}}
% Overskriften er spatieret Fraktur med anden Linie i spatieret Antiqua;
% Spatieringen i Overskrifter mærkes ikke, jf. husnormen.

I sin Modsætning til Magten begynder Viden med det Almene; Formlen
\textit{universalia ante rem} er forsaavidt betegnende for Videns
Oprindelighed. I Modsætning til Viden begynder Magten med at sætte og
forudsætte det Enkelte; Formlen: \textit{universalia post rem}, er med
Hensyn hertil betegnende for Magtens Oprindelighed. Et Tilbageblik paa
de to Phaser vil imidlertid snart overbevise os om, at
Vexelbestemmelsen af det Almene og det Enkelte ikke ligefrem lader sig
angive ved scholastiske Formler. Betragte vi den for Viden gjældende
Ideebevægelse, da er det indlysende, at det Almene i Principet
forudsætter det Enkelte, at det under sin uophørlige Selvbestemmen
bestandig maa have det Enkelte til Mellemled ved de Fordoblinger, hvori
det har sin uendelige Rigdom, at det Almene, med andre Ord, ligesaavel
maa resultere af det Enkelte, som gaae forud for det Enkelte, at altsaa
\textit{universalia ante rem} oprindelig maae forudsætte \textit{uni\-%
versalia post rem}. Betragte vi den for Magten gjældende
Realbevægelse, da kan det Almene jo kun abstraheres af det Enkelte
under den Forudsætning, at det Almene selv er sat i og med dette
Enkelte; for at sætte det Enkelte som det Almenes existentielle Udtryk,
maa Magtbegrebet forudsætte Videns\-%
% --- p. 208 ---
\setcounter{footnote}{0}%
begrebet: \textit{universalia post rem} \textbf{er kun} en Stadfæstelse
af \textit{universalia ante rem}. Men, naar Extremerne: \textit{ante
rem} og \textit{post rem} saaledes uophørlig forudsætte og frembringe
hinanden, saa ligger det nær at søge Udtrykket for det Sande i den
forenende Midte, altsaa at vælge Betegnelsen \textit{universalia in
re}.

I formel Henseende har Betegnelsen \textit{universalia in re} ganske
vist ogsaa sin Betydning; men det Charakteristiske \textbf{er}, at den
logiske Betydning ligesaavel ved denne sidste som ved de to
foregaaende Betegnelser just er i sit Begreb det Modsatte af, hvad den
er i den scholastiske Mening. Efter den scholastiske Mening vil
Sætningen: \textit{universalia sunt in re} aabenbart forstaaes
saaledes, at Tanken skal slaae sig til Ro ved den Væren,
\textit{universalia} have i \textit{res}, altsaa opfatte Eenheden af
\textit{universalia} og \textit{res}, ikke blot som en Værenseenhed,
men som en værende Eenhed. Men her viser sig igjen Misforstaaelsen.
Sætningen: \textit{universalia sunt in re} \textbf{hører} ingenlunde
til de rolige, men tvertimod til de urolige, af allehaande skjulte
Modsigelser hjemsøgte Sætninger; Betegnelsen \textit{in re} sætter jo
paa een Gang Identitet og Forskjel af \textit{uni\-%
versalia} og \textit{res}. Hvad er vel \textit{res}? Systemet af de
Existensbestemmelser, der ere forenede i Objectet, det existerende
Enkelte. Og \textit{universalia}? Systemet af de Prædicater, der i
Dommen svare til de for Subjectet gjældende Existensbestemmelser. Men
nu har det jo i vort Foregaaende viist sig, at Existensbestemmelserne
ligesaalidt kunne være Prædicater, som Prædicaterne
Existensbestemmelser, efterdi Prædicaterne ere logiske,
Existensbestemmelserne antilogiske; hvori da ligger: at
\textit{universalia}, forsaavidt de ere \textit{in re}, ikke kunne være
\textit{universalia}, og, forsaavidt de ere \textit{universalia}, ikke
kunne være \textit{in re}. Vexelbestemmelsen af \textit{universalia} og
\textit{res} fordrer altsaa paany en Bevægelse.

\emph{I den evige Fornuft er Viden og Magt oprin\-%
delig Eet.} Det er mærkværdigt, hvor naturlig og let man faaer den
videnskabelige Bevidsthed, især, hvis man omskriver
% --- p. 209 ---
\setcounter{footnote}{0}%
Vidensbestemmelserne ved „Tankebestemmelser“, endog fra de
forskjelligste Synspunkter til at indrømme denne Sætning.
„Naturhandlingerne“, siger H. C. Ørsted\footnote{Naturlærens
mechaniske Deel. Kbhvn. 1844. Første Indledning.}, „rette sig efter
Tankehandlingerne. Naturlovene kunne derfor ogsaa kaldes Naturtanker.
Samtlige Naturlove udgjøre derfor et Hele, ligesom samtlige Tanker i
Fornuften. Naturen maa saaledes opfattes af os som et Fornufthele.
Mangfoldigheden i Naturen er derfor at betragte som en Tankeudfoldning,
Omskiftningerne som en Tankebevægelse. Men ligesom der i al Tænken er
en Virksomhed, og Tanken tænkt uden denne kun er en Tænkningsform,
saaledes er der i enhver tilværende Ting en Virksomhed, ved hvilken den
hævder sin Tilværelse og yttrer den mod andre Ting. Uden denne
Virksomhed har den ikke Virkelighed. I Tænkningen kalde vi dette
Virksomme Villie, i de udvortes Gjenstande Kraft\footnote{Det
kategoriske Forhold mellem Kraft og Magt er i nærværende Afsnit kun
flygtig berørt, fordi Kræfternes Analyse, forsaavidt den gives i
Grundideernes Logik, væsentlig henhører under Magtens Kredse. At Magt
og Kraft ere beslægtede, kan sees af Sprogbrugen. Heiberg har i sin
„speculative Logik“ (Anfr. Skrft. S. 72) angivet Forholdet saaledes:
„Magten er den absolute Nødvendighed. Dens Momenter ere Kraft og
Yttring, af hvilke den Ene er Grændsen for den Anden, og for saa vidt
dens Qvalitet. Kraften er Mulighed af Yttring, og Yttringen er
Umulighed af Kraft, d. e. Kraftens Ophør, thi i Yttringen udtømmes
Kraften, og ophører, som i sin Grændse. Men Kraften, betragtet som
Mulighed, er \emph{Betingelsen} for Yttringen, og Yttringen er den
\emph{Existens}, som fremgaaer, hvor Betingelsen finder Sted; medens
\emph{Nødvendigheden} er den Virksomhed, som, ved at ophæve
Betingelsen, giver det Betingede Existens, og (da Betingelsen eller
Forudsætningen ophæves) fører det Existerende tilbage til umiddelbar
Væren, ikke, saaledes som de øverste logiske Grundsætninger, ved at
tilintetgjøre Reflexionen, men tvertimod ved at fuldføre denne, saa at
den kommer til sin Ende, og ophæves. Det Nødvendige er vel nødvendigt
formedelst et Andet,
% her falder det trykte sideskifte: Noten fortsætter uden Mærke ved Foden
% af trykt S. 210
som er dets Betingelse, og denne Betingelse er selv
existerende eller selvstændig, men da det Nødvendige ikke selv kommer
til Existens, førend Betingelsens Existens er ophævet, saa er det
Nødvendige ligesaavel ubetinget d. e. absolut, eller: det er, fordi det
er.“}. Men ligesom Villen og Tænken
% Signaturmærket „Grundideernes Logik.“ + „14“ staaer ved Foden af denne
% Side; ikke transskriberet.
% --- p. 210 ---
\setcounter{footnote}{0}%
er Eet, kun betragtet fra to Sider, saaledes ogsaa Kraft og Naturlov;
de ere Yttringer af den guddommelige Villie og den guddommelige
Tænken.“

Men naar man i logisk Forstand skal gjøre Rede for Videns og Magtens
oprindelige Eenhed og derved for Muligheden af, at „Naturhandlingerne“
kunne „rette sig efter Tankehandlingerne“, saa viser sig et stort
Problem, det Problem, med hvis Løsning Subjectivitetens Logik maa
culminere. Et Forsøg paa at besvare det Spørgsmaal, hvad det egentlig
vil sige, at „Naturhandlingerne“ rette sig efter „Tankehandlingerne“,
vil allerede gjøre Problemet kjendeligt. Vælge vi som Exempel paa en
Naturhandling Faldet paa Skraaplanet, da er det for det Første
indlysende, at de Tankehandlinger, der hos Betragteren udfordres til at
forstaae en saadan Naturhandling, ingenlunde falde sammen med en
Opfattelse, der umiddelbart og ligefrem retter sig efter
Naturhandlingen selv. Hele Systemer af logiske Bestemmelser ere i det
tilværende Skraaplan, i det faldende Legeme og som Følge heraf i
Faldbevægelsen selv forenede til antilogisk Eenhed; disse i
Naturhandlingen forenede Systemer maae ved særskilte Tankehandlinger
adskilles og forbindes paa eiendommelig Maade, naar Handlingen skal
sees at udtrykke bestemte Naturtanker.

„Et Legeme“, hedder det,\footnote{H. C. Ørsted. Anfr. Skr. S. 202.}
„behøver saa mange Gange mere Tid til at falde ned ad Skraaplanets
Længde end til at falde igjennem dets Høide, som hiin er større end
denne.“ Herved er da en Naturlov udtalt som en Naturtanke; men i
hvilken Grad Naturhandlingerne i de phænomenale Bevægelser, hvorved
Loven opfyldes, afvige fra de Tankehandlinger, hvor\-%
% --- p. 211 ---
\setcounter{footnote}{0}%
ved Lovopfyldelsen indsees, fremlyser, naar man kaster et Blik paa de
mathematiske Udtryk\footnote{„Betegnes Skraaplanets Høide ved a, dets
Længde ved l, vil \dots\ Grundhastigheden for Faldet paa Skraaplanet
blive $g\frac{a}{l}$, og Bevægelsen er bestemt ved Udtrykkene
\[ h = \frac{gat}{l} \qquad r = \frac{ga}{2l}t^2 \]
Varigheden af Faldet paa Skraaplanet findes af det andet Udtryk ved at
sætte $r = l$, og betegnes den tilsvarende Tid ved $t_1$, faaer man
\[ t_1 = l\sqrt{\frac{2}{ga}} \]
Den Tid $t_2$, et Legeme vilde bruge til at falde igjennem
Skraaplanets Høide, er“ --- her henvises til et Foregaaende ---
„bestemt ved
\[ t_2 = \sqrt{\frac{2a}{g}} \]
og af dette Udtryk i Forbindelse med det foregaaende findes
\[ \frac{t_1}{t_2} = \frac{l}{a} \]
eller, med andre Ord, Faldet paa Skraaplanet fordrer i samme Forhold
længere Tid, som Længden af Skraaplanet er større.“ H. C. Ørsted,
Naturlærens mechaniske Deel 3die Udgave ved C. Holten. 1859 S. 187.},
der bestemme Lovopfyldelsen fra dens forskjellige Hovedsider. Vistnok
mødes i alle disse Hovedpunkter „Naturtanken“ med den mathematiske
Tanke; men Naturhandlingen er derfor ligesaalidt en mathematisk
Operation, som den mathematiske Operation er en Naturhandling; der
faldes ikke i en Ligning, ikke i en Qvotient, ikke i en Qvadratrod.

For Naturhandlingen gjælder det, at \textit{universalia ere in re}, om
Tankehandlingen gjælder det derimod, at \textit{universalia} saavel
\textit{post rem} som \textit{ante rem ere in mente}. Ethvert
Experiment forudsætter et System af \textit{universalia ante rem}; thi
Experimentet faaer just sin Betydning derved, at der i
% Signaturmærket „14*“ staaer ved Foden af denne Side; ikke transskriberet.
% --- p. 212 ---
\setcounter{footnote}{0}%
det givne Tilfælde forelægges Naturen et bestemt Spørgsmaal til
Besvarelse, og den Spørgende maa nødvendig have et Indbegreb af
allerede udviklede Begreber \textit{in mente}. Ligeledes forudsætter
ethvert Experiment, at der kan udvikles et System af
\textit{universalia post rem}; thi Forsøget anstilles jo netop i det
Øiemed, at faae Existensbegreberne oversatte i Begrebets Form, og
bevare de af Kjendsgjerningerne abstraherede \textit{univer\-%
salia in mente}. De Tankehandlinger, der udfordres for at anstille
Experimentet, og det Indbegreb af Naturhandlinger, hvori Experimentet
fuldstændig udtales, ere, vel at mærke, som Handlinger, ikke parrallelt
løbende.
% TRYKFEJL: „parrallelt“ som trykt (for „parallelt“); bogstaverne
% p-a-r-r-a-l-l-e-l-t er utvetydige ved 600 dpi
Kunde Tankehandlingen ligefrem og i al Umiddelbarhed følge
Naturhandlingen, saa maatte Reflexionen falde sammen med Sandsningen.
Der vilde altsaa til at forstaae en Naturbegivenhed ikke behøves Andet,
end at Iagttageren fra alle Sider fik seet, hvad her var at see, hørt,
hvad her var at høre, kort: sandset, hvad her var at sandse; men
ligesaa lidt som Sandsningen ligefrem kan følge Tænkningen, saa at
Sandsebevægelsen Punkt for Punkt skulde svare til Tankebevægelsen,
ligesaalidt kan Naturhandlingen ligefrem følge Tankehandlingen.

Af den her paaviste Charakteerforskjellighed følger da med uafviselig
Conseqvens, at Tankehandling og Naturhandling ligesaa umulig kunne være
Eet i „den evige Fornuft“, som i „den menneskelige Erkjendelse“.
Vistnok maa det indrømmes, at den evige Fornuft ikke behøver en saadan
Udstykning og Sammenstykning af Forkundskaber, Forberedelser og
Overveielser for at erkjende Faldet paa Skraaplanet, som den
menneskelige; man kan umulig tænke sig den evige Fornuft i Begreb med
at addere, multiplicere, uddrage Qvadratroden, indprænte sig Formler
for Hastigheden, det gjennemløbne Rum, Tiden o. s. v. Men derfor staaer
det ligefuldt fast, at de Idealitetsbestemmelser, der paa det
psychologiske Standpunkt haves paa ufuldkommen Maade \textit{in mente
humana}, maae
% --- p. 213 ---
\setcounter{footnote}{0}%
paa det ontologiske Standpunkt, saa sandt der gives en evig Viden,
haves paa fuldkommen Maade \textit{in mente divina}.

At Naturlæren som Fagvidenskab ikke kan finde sig opfordret til at
anstille logiske Undersøgelser angaaende Tankehandlingens
Eiendommelighed i Modsætning til Naturhandlingen, er forklarligt nok;
men mærkværdigt er det, at Problemet hidtil endog synes at have
unddraget sig Philosophiens Opmærksomhed. Den „immanente Methode“ er i
denne Henseende saalangt fra at ane nogen Vanskelighed, at den
tvertimod, som vi have seet, i al dogmatisk Umiddelbarhed gaaer ud fra
den Forudsætning, at Tankebevægelsen og Sagbevægelsen overalt maae
falde sammen. „Slutningen\footnote{J. L. Heiberg. Specul. Log. (Anfr.
Skr. S. 93 flgd.) (Prosaiske Skrifter, 1ste Bind, S. 303--304).} maa
følgelig, ligesaalidt som Omdømmet i Almindelighed, betragtes blot som
en menneskelig Tankeoperation; tvertimod maa Slutningen betragtes som
Sandheden af alt Existerende, som Fornuften i alt Tilværende; og Alt er
saaledes en Slutning. Eller: Alt er Begreb, og dets Tilværelse er
Forskjellen mellem Begrebets Momenter, hvilken Forskjel bestaaer deri,
at enhver Tings Enkelthed staaer under en Particularitet af Egenskaber
eller Væren udenfor sig selv, hvorved den hæver sig til sin almindelige
Natur, og optager denne i sig selv. Saaledes er f. Ex. hele det
menneskelige Liv en Slutning. Individet er det Enkelte. Dets
Fornødenheder, dets Virksomhed, dets Eiendom, med eet Ord Alt, hvorved
det staaer i Forbindelse med den udvortes Verden, er Mellemledet mellem
det selv og det Almindelige, eller \emph{Midler} (d. e. Mellemled,
\textit{media}) til at forædle og uddanne det, d. e. til at forene dets
Individualitet med det Almeengyldige. Betragtede som menneskelige
Tankeoperationer, ere Omdømmer og Slutninger Udtrykket for den
subjective Erkjendelse af hiin Grundfornuft i Tingene, eller Tankens
Congruens med Virkeligheden.“ Man skulde altsaa, ifølge denne
Udtalelse, antage, at den fornuftige Tanke\-%
% --- p. 214 ---
\setcounter{footnote}{0}%
bevægelse var congruent med Bevægelserne i de virkelige Begivenheder;
at her i Virkeligheden altsaa var et virkeligt Almeent og et virkeligt
Særskilt, ligesaavel som et virkeligt Enkelt.

Den ideale Sammensmeltning af det Subjective med det Objective, vi ved
at opstille Subjectivitetsproblemet allerede have fremhævet, og siden
fra forskjellige Sider og under forskjellige Vendinger nærmere søgt at
belyse, maa da i nærværende Sammenhæng endnu engang underkastes en
særegen Prøvelse. Det er Forholdet imellem det Enkelte, det Særskilte
og det Almindelige, hvorpaa Opmærksomheden fæster sig. Hvorledes bliver
da det Enkelte, det Særskilte og det Almene, ifølge den speculativt
immanente Methode, bestemt ved Slutningen?

„I Slutningen,“ siger Heiberg,\footnote{Anfr. Skr. S. 93. (Pros. Skr.
1, S. 304).} „bliver det Enkelte, det Særskilte og det Almindelige
bestemt som \emph{Individ, Art og Slægt}\dots\ Betegner man det Enkelte
med $E$, og det Almindelige med $A$, saa kan Omdømmet fremstilles ved
Udtrykket $E - A$, hvilket vil sige, at det Enkelte og det Almindelige
ere forenede, eller at det Enkelte er det Almindelige. Betegnes nu
endvidere det Særskilte med $S$, saa fremstilles det fuldstændige
Omdømme eller Slutningen ved Udtrykket $E - S - A$, hvilket vil sige,
at det Særskilte er Mellemleddet, hvorved det Enkelte er forbundet med
det Almindelige, eller at det Enkelte er det Almindelige formedelst det
Særskilte\dots\ Men\footnote{Anfr. Skr. S. 97--98. (Pros. Skr. 1, S.
313--314).} idet det Enkelte i hiin første Sætning bestemmes som det
Almindelige, følgelig som det, der er det Fælleds for begge de andre
Led, saa bliver det selv til den forbindende Midte, hvilket giver den
anden Slutning, $S - E - A$, hvis Betydning er, at det Særskilte er det
Almindelige formedelst det Enkelte, eftersom dette selv er baade
særskilt og
% E, A og S er Antiqua-Versaler; Bindetegnet imellem dem er en lang
% Tankestreg, her sat som minus i mathematisk Sats. Afstanden om Stregen
% er uregelmæssig i Trykket (snart spatieret, snart tæt); normaliseret.
% --- p. 215 ---
\setcounter{footnote}{0}%
almindeligt, ifølge den første Sætning. --- Da fremdeles, ifølge den
første Sætning, det Enkelte er bestemt som det Almindelige, og ifølge
den anden det Særskilte ligeledes er bestemt som det Almindelige, saa
er det Almindelige Bindeledet, som forbinder det Enkelte med det
Særskilte, hvilket giver den tredie Slutning, $E - A - S$, hvis
Betydning er, at det Enkelte er det Særskilte formedelst det
Almindelige. Og da nu endelig, ifølge denne sidste Slutning, det
Enkelte er det Særskilte, og ifølge den foregaaende det Særskilte er
det Almindelige, saa bliver det Særskilte paany Bindeleddet, og
frembringer den første Slutning, $E - S - A$, der først var fremkommen
uden foregaaende Slutning. Herved er Slutningernes System fuldendt,
eftersom ethvert af Omdømmerne \textbf{nu er}
medieret.“\footnote{Som Exempel paa det Fornuftige i, at Alt „er et
System af tre Slutninger“, vedføies (Anfr. Skr. S. 99) følgende
Anmærkning: „Individet, som ved sin sandselige Natur forbinder sig med
Fornuften, staaer under Formen af den første Slutning, $E - S - A$. Men
Individet er selv Foreningen af den sandselige Natur og Fornuften;
dette giver den anden Slutning, $S - E - A$. Og endelig er det
Fornuften, hvori saavel Individets Enkelthed, som dets sandselige Natur
er grundet; dette giver den tredie Slutning, $E - A - S$. Mange
Uenigheder og Forvirringer i Videnskaben have deres Grund deri, at man
holder sig til den enkelte Slutning, istedenfor til det hele System.
Herved fremkomme da forskjellige, ja modsatte Paastande, af hvilke
enhver er rigtig, forsaavidt som den er en Side af Sandheden, men
urigtig, forsaavidt som dens Sandhed er eensidig.“}

Kom det \textbf{nu an} paa at udføre det Formelle i denne Systematik
videre, da vilde det her være en let Sag at vise, hvorledes Slutningen
$E - A - S$ væsentlig svarer til \textit{universalia ante rem},
Slutningen $S - E - A$ til \textit{univer\-%
salia post rem} og Slutningen $E - S - A$ til \textit{univer\-%
salia in re}; men Formsiden er allerede saa fremherskende, at
% --- p. 216 ---
\setcounter{footnote}{0}%
Systematiken, endog i betænkelig Grad, nærmer sig til Schematisme. Skal
„Systemet af de \textbf{tre} Slutninger“ faae nogen sand Betydning ved
Bestemmelsen af Forholdet mellem Tankehandlingen og Naturhandlingen,
med andre Ord, ved Bestemmelsen af Forholdet mellem Ideen og
Virkeligheden, da maa enhver af Systemets \textbf{tre} Slutninger igjen
afgive Typen for en Sphære af Systemer, hvori den eiendommelige
Bevægelse kan sees at være bestemt ved Beskaffenheden af det Led, hvis
Udvikling Bevægelsen tjener. Sættes saaledes Udviklingen og
Articulationen af det Almene som Opgave, ifølge Schemaet: $E - A - S$,
faae vi et Idealsystem; sættes omvendt Udviklingen og Articulationen af
det Enkelte som Opgave, efter Schemaet: $S - E - A$, faae vi et
Realsystem; sættes endelig Vexelbestemmelsen af det Almene og det
Enkelte eller Udviklingen af det Særskilte som Opgave efter Schemaet:
$E - S - A$, vil det være begge de eensidige Systemers teleologiske
Eenhed, der articulerer sig i et fuldstændigt System.

1) \emph{Idealsystemet.} Det Enkelte er det Almene, $E - A$, denne Dom
er logisk, altsaa udtrykkelig en Tankehandling. Den væsentlige Forskjel
mellem Subject og Prædicat beroer, ifølge vort Foregaaende, ikke paa
nogen Forskjel i Fuldstændigheden af Bestemmelserne, thi den samme
Totalitet, der findes i Subjectet, kan ogsaa findes i Prædicatet; men
Modsætningen, den afgjørende og gjennemgaaende Modsætning mellem
Subjectet og Prædicatet, beroer derpaa, at Prædicatets Bestemmelser ere
logiske, Subjectets antilogiske. I sin logiske Form vil
Prædicatstotaliteten som det Almene udtrykke Eenheden, medens
Subjectstotaliteten i sin existentielle, antilogiske Form vil kunne
gjentage sig i det Uendelige, og saaledes udtrykke Mangfoldigheden.
Saavel paa Mangfoldighedens som paa Eenhedens Side have vi her det
Absolute selv; nemlig paa Eenhedens Side, altsaa i det Almindeliges
Form, det Absolut-Absolute; paa Mangfoldighedens Side,
% --- p. 217 ---
\setcounter{footnote}{0}%
altsaa i det Enkeltes individuelle Form, det Relativ-Absolute: under
begge Former haves Begrebet, forsaavidt det er et Vidensbegreb, i sit
absolute Indhold som Idee.

Det Enkelte i dets existentielle Umiddelbarhed kan den guddommelige
Viden --- medmindre den da skulde kunne see Fiirkanten som en Cirkel ---
ligesaalidt gjennemskue umiddelbart, som den menneskelige Viden kan det:
det Existerende er og bliver Videns Gjenstand og som Videns Gjenstand en
Skranke for Viden. Derimod har den guddommelige Viden det uendelige
Fortrin fremfor den blot menneskelige, at den Vidende seer det
existerende Enkelte i dets adæqvate Idee. Individualideen, det Enkeltes
adæqvate Idee, gjentager sig altsaa i Mangfoldigheden med det Enkelte;
saaledes faaer da Dommen $E - A$ den for Ideelæren afgjørende Betydning,
at den samme Totalitet, det samme Absolute, den samme Idee, kan paa een
Gang sees i Eenhedens og i Mangfoldighedens, i Absoluthedens og
Relativitetens Lys. I den ene og samme Idee kan den vidende
Subjectivitet altsaa have utallige bestemte Ideer, og i de utallige
Ideer een eneste bestemt Idee.
% „Fiirkanten“ with two dotted i's, checked at 600 dpi; „adæqvate“ (qv)
% both times. The judgement-formulae of these pages are set in upright
% antiqua capitals joined by an em-length dash; set in math mode with a
% minus, following the house treatment of $A$, $A=A$ etc. (pp. 16, 174).

Det vil, som Følge heraf, da heller ikke kunne være til Hinder for den
fuldkomne Begriben, at her, istedenfor eensformig Gjentagelse af det
samme Enkelte, indtræder en existentiel Totalitetsforskjel mellem
Enkeltexistenserne indbyrdes, saa at de Enkelte adskille sig i to eller
flere Rækker. Tvertimod! Totalitetsforskjelligheden er jo netop et
eiendommeligt Udtryk for, at Magten tjener Viden, idet Vidensbegrebet
articuleres ved det varierende Magtbegreb. Med Existenstotaliteternes
Forskjellighed fremkommer det Særskilte. Men ligesom det Særskilte i sin
Begrebseenhed altid udtrykker Fordoblingen, saa at, hvor der er eet
Særskilt, er der ogsaa to, saaledes viser det Særskiltes Oprindelse sig
strax dobbeltsidig: det er Særskiltheden i det Enkelte, der gjør det
Almene til et Særskilt, hvilket her kan betegnes ved Dommen: $E - S$.
% --- p. 218 ---
\setcounter{footnote}{0}%
Men da Magtbegrebet, som Begreb, kun kan forvandle Vidensbegrebets Form,
ikke dets Indhold, kan Særskiltheden umulig fremkomme i det Enkelte,
uden under den Forudsætning, at den allerede er i det Almene; hvilket her
kan betegnes ved Dommen: $S - A$. Almeenideens Eenhed med
Individualideerne er da herved articuleret til et System af særskilte
Ideer; i hver særskilt Idee er det hele System gjennemsigtigt.

I Ideernes Sphære er Særskiltheden ikke blot et Mellemled mellem det
Almene og det Enkelte, men tillige begges absolute Eenhed. Her have vi
Slutningen: $E - S - A$, og Heiberg har forsaavidt Ret, naar han lægger
Vægt paa, at det Enkelte, det Særskilte og det Almene blive i Slutningen
bestemte som \emph{Individ, Art og Slægt}. Naar det Særskilte er Art, saa
er det Almene Slægt; men det kommer i denne Sammenhæng fornemmelig an paa
at bestemme Bevægelsen som Tankehandling, og fremstille \emph{Slægt og
Art}, ikke som blot formelle Kategorier, men som Ideer. Det er det Almene
selv, paa hvis Articulation Idealsystemet fra alle Sider lægger an. Hvad
enten her iøvrigt handles om Individet eller om Arten eller om Slægten,
saa er det dog lutter Vidensbestemmelser eller Ændringer, Formationer og
Berigelser af det Almene, hvorom det gjælder: herved, og kun herved, er
Bevægelsessystemet eiendommelig charakteriseret som et System af
Tankehandlinger.

Af disse Tankehandlinger sees udtrykkelig, at den Totalitet af
Bestemmelser, der paa existentiel Maade svarer til Individerne i en vis
Række, netop afgiver de antilogiske Momenter, hvorved Kategorien Art
integreres til Eenhedsudtryk for et relativt Absolut, og bliver en
særskilt Idee. Fremdeles vil det af disse Tankehandlinger være
indlysende, at de særskilte Totalitetsbestemmelser, hvorved de under
samme Slægt hørende Arter adskille sig fra hverandre indbyrdes, ikke
falde udenfor, men sammen med det System af Bestemmelser, der
charakteriserer Slægten som Slægt. Saalænge \emph{Slægten}
% --- p. 219 ---
\setcounter{footnote}{0}%
blot opfattes som Kategorie, er den naturligviis fattigere paa
Bestemmelser end de Arter, der høre under den; men som absolut
Eenhedsudtryk for de forskjellige, paa Arterne fordeelte
Eiendommeligheder, er Slægten Idee. „Her danner sig ikke et
Slægtsbegreb, fordi her forekomme forskjellige Arter; men her forekomme
forskjellige Arter, fordi Slægtens Idee kun kommer til Existens ved at
differentiere sig i Artsideer, ligesom Artsideen kun kommer til Existens
ved at differentiere sig i
Individualideer.“\footnote{Forelæsn. ov. Phil. Propæd. 1860--61 S.
324--25.}
% printed mark: *)  — the letterspacing of „Slægten“ ends with the word
% on p. 218; the sentence runs on, so no blank line at the page break.

Endelig vil det af de for Idealsystemet afgjørende Tankehandlinger kunne
forstaaes, at det Almene, der i en bestemt Sammenhæng er articuleret og
bestemt som \emph{Slægt i Modsætning til Art}, \textit{eo ipso} er
bestemt som \emph{Art} i Modsætning \emph{til Familien}, der i logisk
Forstand kun er en Slægt af høiere Orden, o. s. v., saa at Systemet paa
denne Maade kan articuleres til Systemer af Systemer i det Uendelige.
Bevægelsen fra Individualideerne gjennem Arterne til Slægten, gjennem
Slægt-Arterne til Slægternes Slægt i Almeenideens absolute,
indholdsrige Eenhed, er at opfatte som Subjectivitetens Fordybelse i sin
Videns Indhold. „Den for Ideeudviklingen afgjørende, med Selvhedsloven
identiske Concretionslov er allerede bestemt, naar Slægtsideen er naaet.
Istedenfor at bestemme Gjennemgangen gjennem Individerne til Arten,
gjennem Arterne til Slægten, gjennem Slægten til Familien o. s. v. som en
Overgang fra det Concretere til det Abstractere, maae vi da, ifølge den
paaviste Lov, omvendt bestemme denne Gjennemgang som en Fordybelse i
stedse rigere, stedse inderligere
Concretioner.“\footnote{Forelæsn. ov. Phil. Propæd. 1860--61 S. 325.}
--- Allerede den grundende Reflexion gav os en Selvfordybelse; men dens
Eensidighed viste sig deri, at den førte til Gjæring og Grublen. Denne
Uklarhed er nu fuldstændig overvunden:
% printed mark: **) — this page carries TWO footnotes; the per-page
% counter gives *) and **) as printed. The emphasis is broken up as
% printed: „Slægt i Modsætning til Art“ is letterspaced throughout, but
% in the answering clause only „Art“ and „til Familien“ are spaced and
% the intervening „i Modsætning“ is not. Not regularised.
% --- p. 220 ---
\setcounter{footnote}{0}%
den absolute Ideernes Eenhed er i Idealsystemet ligesaa klar og bestemt
som de relative Ideers uendelige Mangfoldighed; thi Eenheden er med
Mangfoldigheden overalt seet og sat ved den samme Tankehandling af den
samme Tænkende.

2) \emph{Realsystemet.} Det Almene er det Enkelte, $A - E$, denne Dom er
antilogisk, altsaa betegnende for en Naturhandling. Det Almene er i sin
rene Almindelighed altid en Tanke; men i dette System er Tanken identisk
med Magten: det Almene er en mægtig Tanke, en Tankemagt, altsaa en Lov.
Men Loven er kun virkelig i sin Opfyldelse; Opfyldelsen er i det Enkelte,
bestemt ved det Enkelte og Eet med det Enkelte. Alt hvad der i
Virkeligheden er til, er til i Kraft af en Lovopfyldelse, og just derfor
til som et Enkelt: hvad der hviler, hviler ifølge en Lovopfyldelse; hvad
der bevæger sig, bevæger sig ved en Lovopfyldelse. Det Samme gjælder om
al Vorden, Opstaaen, Forgaaen o. s. v. Loven, den almene Magt, udgaaer
ikke fra de evige Tankers Verden saaledes, at den paa umiddelbar Maade
griber ind i det Endelige og Enkelte, nei, Loven, det mægtige Almene, er
tvertimod bestemt ved Vexelbestemmelse af de enkelte Tings enkelte
Egenskaber. Det er overalt den eiendommelige, punktuelle Vexelvirkning af
Enkelte med Enkelte indtil de allermindste Enkeltheder, hvorpaa
Fyldestgjørelsen af det Almenes Fordringer her maa beroe: Lovens
Opfyldelse er nøiagtig bestemt ved den Tingenes Sammenhæng, den selv
bestemmer. Den antilogiske Dom $A - E$ er fra dette Synspunkt at betragte
som et Udtryk for Lovopfyldelse i Almindelighed; den er forsaavidt et
logisk Symbol paa alle virkelige og mulige Naturhandlinger.
% „2) Realsystemet.“ is a run-in numbered head, letterspaced, kept in the
% paragraph; the numeral is antiqua. „nøiagtig“ has a slashed o at 600
% dpi (the counter is filled where the plain o's on the same page are
% open). „logisk Symbol“ is Fraktur throughout, not antiqua.

Af det Almenes Eenhed med det Enkelte i Lovopfyldelsen følger, at
Opfyldelsen af en enkelt Lov kan \textit{in concreto} være forenet med en
Opfyldelse af flere indbyrdes forskjelligartede Love. Det er ikke blot en
eensidig abstract Almeenbestemmelse, der fra dette Standpunkt sees at
gaae i Eet med en abstract
% --- p. 221 ---
\setcounter{footnote}{0}%
Existensbestemmelse; hele Systemer af \textit{universalia} udgjøre her
det logiske Indhold, som Enkeltexistensen, \textit{res}, det umiddelbare
Magtbegreb, udtrykker i antilogisk Form. Concretionen af
Existensbestemmelserne i deres Idealitet giver os igjen det Almene i
Tingenes Væsen under Form af Kategorier som \emph{Grund og Mulighed}; dog
med den forhøiede Eiendommelighed, der er et Udtryk for, at Reflexionens
Former nu ere integrerede ved Virkelighedens antilogiske Momenter. Den
substantielle Reflexion gav os Magtens Kategorie med dens existentielle
Betydning; men Et er Magtens Kategorie i logisk Reflexion, et Andet er et
System af Magtbestemmelser, hvis Forening factisk udtrykker den factisk
givne Existens. Den nødvendige Eenhed af Virkeligheden og den reale
Mulighed er fra Begrebets Standpunkt allerede paaviist som en nødvendig
Eenhed af det Objective med det Subjective i den ontologiske
Objectivering; men i nærværende Sammenhæng er det ikke nok med Eenheden i
dens Almindelighed; hele Eftertrykket falder paa den for det Virkelige
afgjørende Bestemthed, at det hver Gang just er \emph{denne} Eenhed under
\emph{disse, just disse} Omstændigheder, paa \emph{disse, just disse}
Betingelser. Var Loven som Lov det Almene, saa er den bestemmende Grund,
den reale Mulighed, nu det Særskilte: for at opfyldes i det Enkelte, maa
Loven væsentlig bestemmes ved det Særskilte. Den kategoriske Bevægelse
fra Begreb til Existens gaaer altsaa fra det i Magtens concrete Idealitet
(den bevægende Grund, den virkelige Mulighed) hvilende \emph{Al\-%
mene} gjennem Magtbestemmelsernes \emph{særlige Sammen\-%
voxen} (denne paa disse Betingelser tilstrækkelige Grund; denne i disse
Vexelvirkninger til Anlæg sig fuldendende Mulighed), og kan forsaavidt
betegnes ved Dommen: $A - S$.
% „universalia“ and „res“ are antiqua; „Omstændigheder“ and
% „Betingelser“ are NOT letterspaced although the „disse, just disse“ on
% either side of them are — as printed.

Dommen: $A - S$, erholder imidlertid først sin rette existentielle
Gyldighed, naar den sees i Forbindelse med Grundens Opgaaen i den
bestemte Existens, Mulighedens
% --- p. 222 ---
\setcounter{footnote}{0}%
Omsætning i den individuelle Virkelighed, Anlægets Udførelse, altsaa med
det Særskiltes Overgang i det Enkelte; denne Overgang betegne vi ved den
antilogiske Dom: $S - E$. Naturhandlingerne ere kategoriske
Magthandlinger, de forudsætte Ideen, og fyldestgjøre Ideen, men falde
ikke sammen med Tankehandlingerne. Saaledes er i Forhold til
Naturfrembringelsen Slægtens Idee Artens Mulighed, Artens Idee Individets
Mulighed; men Frembringelsen selv viser, at det ikke er Ideen som Idee,
men Individerne, der i Virkeligheden fuldbyrde Naturhandlingen. Idet
Magtbegrebet overalt har det Enkelte, det Individuelle til Realitet, har
det i den Grad lagt Beslag paa alle Vidensbestemmelser, at de alle maae
finde deres logiske Undergang i Existensens antilogiske Bestemthed. At
denne Undergang ikke desto mindre er Eet med en stadig Gjenopvækkelse, at
Idealiteten i kategorisk bestemte Former gjennemvirker og bærer
Realiteten, saa at vi paa alle Punkter have \textit{universalia in re},
er en ligefrem og nødvendig Følge af, at det i Objectiviteten Objective
maatte forsvinde som et Phantom, dersom det ikke blev ontologisk
objectiveret. Den ontologiske Objectiveren er som Naturhandling --- thi
Naturhandlingen er ikke en abstract objectiv Begivenhed, der paa abstract
objectiv Maade udfører sig selv --- ligesaalidt en umiddelbar Anskuen som
en reen Tænken eller almindelig Reflecteren; den har alle logiske
Forudsætninger, alle almene og særskilte Begrebsbestemmelser nærværende i
Virksomhedens og Anvendelsens Øieblik. Men disse logiske Bestemmelser
frigjøres ikke til logisk Idealitet; tvertimod! de bindes, tilsættes og
udtømmes i den objectiverede Gjenstands reale Frembringelse: $S - E - A$.
% „Øieblik“ has the Fraktur capital O with a slash (600 dpi); both OCR
% witnesses read it as D. The r of antiqua „re“ carries a heavy inked
% serif; read as r against the antiqua r's of „universalia“.

At Naturhandlingen svarer til Tankehandlingen, vil da med andre Ord sige,
at den har Tankehandlingen paa een Gang baade som iboende og
oversvævende. Forsaavidt Vidensbegrebet er bestemt og ved Bestemmelsen
forvandlet til Magtbegreb, er Tankehandlingen iboende; det vil sige: den
er ikke
% --- p. 223 ---
\setcounter{footnote}{0}%
Tankehandling, men en reel Frembringelse, der i alle Punkter udtrykker
Tanken ved at ophæve Tanken, og saaledes svarer til Tankens Fordring.
Forsaavidt Magthandlingen kategorisk virkeliggjør et System af
Vidensbestemmelser, maa Tankehandlingen nødvendig være oversvævende. Det
var jo umuligt, at Magtbegrebet kunde være Begreb, dersom det ikke var et
reelt Udtryk for en begribende Magt; men en begribende Magt maa absolut
være en vidende Magt; men en vidende Magt er en mægtig Viden, thi den
reale Begriben forudsætter det Logiske, og det Logiske beroer paa Viden.
Er Idealsystemet med sine Tankehandlinger et tomt Skin, naar det Reale
mangler, saa er Realsystemet med sine Magthandlinger en meningsløs
Kjendsgjerning, naar den Idealitetens Frigjørelse, der er nødvendig for
Begrebernes logiske Udvikling, tænkes borte.

Ogsaa i Realsystemet gjør Forskjellen imellem Ideen som Kategorie og
Kategorien som Idee sig gjældende. Seet fra det empiriske Synspunkt kan
Ideen kun vise sig som Kategorie. Sagen, det Concrete, det Virkelige selv
er Individet. „Naturen“, hedder det, „giver os overalt kun Individer,
ikke Arter, endnu mindre Slægter.“\footnote{Jvfr. Forelæsn. over Phil.
Propæd. 1860--61 S. 250. flgd.} Slægter og Arter ere da, ifølge denne
Synsmaade, kun Abstractionsbegreber, \textit{universalia}
\emph{\textit{post rem}}, der ikke have Noget med de reale
Frembringelser selv at gjøre. Men hvor tankeløs en saadan Opfattelse er,
kan allerede sees deraf, at den gjør Naturhandlingen selv tankeløs. Det
er ved Identiteten af det Logiske og det Antilogiske i Videns Form, at
Kategorien bliver Idee; det er ved den tilsvarende Identitet i Magtens
Form, at der gives en Realitet, og at denne Realitet udtrykker Ideen.
Dersom der Intet var, der bar Realitetens Anlæg i sig, saa kunde der jo
heller Intet realiseres, og saa var der ingen Realitet; men tomme Tanker,
abstracte Begreber, subjective Indfald og vil\-%
% printed mark: *)  — „universalia“ is plain antiqua, „post rem“ is
% letterspaced antiqua, so spaced antiqua as printed; „Jvfr.“ takes the
% J sort per the house convention. The word „vilkaarlige“ is broken at
% the foot of the page and joined.
% --- p. 224 ---
\setcounter{footnote}{0}%
faarlige Forestillinger bære ikke i sig Realitetens Anlæg, og kunne
derfor, strengt taget, heller ikke realiseres. Kun naar Begrebet i sit
logiske Væsen saaledes magter de antilogiske Momenter, at det formedelst
denne ideelle Magt har Realitetens Anlæg i sig, kan det --- i en
Naturhandling --- realiseres; men et saadant Begreb er meer end en
Kategorie, det er en Idee. Det er ligesaalidt de formelle Begreber som de
begrebløse Elementer, der, enten hvert for sig eller naar de lægges
sammen, have Realitetens Anlæg; nei, det er i Begrebernes, Elementernes,
Kræfternes oprindelige, absolute Eenhed, altsaa i Ideen, at Anlæget
væsentlig maa søges: kun Ideerne have Realitetens Anlæg; kun de
guddommelige Ideer lade sig i egentlig Forstand realisere\footnote{At det
Taabelige og Forvirrede dog til en vis Grad lader sig realisere, beroer
derpaa, at Ideemomenterne kunne udstykkes i det Endelige, og Realitetens
Anlæg til en vis Grad lade sig forvirre.}.
% printed mark: *)  — printed „realisere *).“, the point after the mark.

3) \emph{Ideal- og Realsystemets teleologiske Eenhed.} I hvilken Grad
den hegelske Logik, saasnart man seer bort fra Idealitetens mystiske
Dæmring i prægnante Udtryk, synker ned til en blot og bar Formlære, er
især iøinefaldende, naar man betragter dens Behandling af „den
teleologiske Virksomhed“. „Saavel i den mechaniske, som i den dynamiske
Virksomhed“ --- hedder det\footnote{Heiberg. Specul. Logik. (Anfr. Skr.
S. 107. flgd.; Pros. Skrifter, 1ste Bind S. 333 flgd.)} --- „er
Begrebet, som reen Objectivitet, udvortes for sig selv, hist som det
umiddelbart Ydre, her som den relative og gjensidige Yderlighed af det
Positive og det Negative. Men i den teleologiske Virksomhed, som er
Begrebets Virksomhed, (ligesom den mechaniske er Umiddelbarhedens, og den
dynamiske Reflexionens) vender det tilbage til sig selv, og da Begrebet
baade er Subject og Object, fremstiller det paany den samme Modsætning,
og sætter sig, som subjectivt, imod den objective udvortes Verden,
% printed mark: **) — this page carries TWO footnotes. The run-in head
% „3) Ideal- og Realsystemets teleologiske Eenhed.“ is letterspaced
% throughout, including the suspended compound „Ideal- og“, which keeps
% its printed hyphen. „Pros. Skrifter“ (= Prosaiske Skrifter) has a
% long s, not an f, at 600 dpi. The quotation opened at „er Begrebet“
% closes on p. 225.
% --- p. 225 ---
\setcounter{footnote}{0}%
men saaledes, at det betragter denne som den Modsætning, der skal
ophæves i det selv. Paa det subjective Standpunkt er Begrebet, i
Modsætning til den objective Verden, bestemt som den Hensigt eller Plan,
der vel er fattet, men endnu ikke udført, d. e. som \emph{subjectivt
Øiemed}. I denne Slutning, hvori det subjective Øiemed er bestemt som det
Enkelte, og Øiemedets Opfyldelse som det Almindelige, hvormed det Enkelte
skal sammensluttes, er \emph{Midlet} det Særskilte eller det Bindeled,
som forener Hensigten eller Planen med dens Udførelse. Resultatet af
Slutningen er saaledes den udførte Hensigt eller Plan, eller det opfyldte
Øiemed, d. e. \emph{Øie\-%
medet som objectivt}; eller Resultatet er det Omdømme, hvori det
Subjective sammensluttes med det Objective: $E - S - A$. Den anden
Slutning, $S - E - A$, giver til Resultat det Omdømme, at Midlet selv er
det Objective, d. e. at Hensigten allerede er opfyldt i Midlet. Og i den
tredie Slutning, $E - A - S$, udtrykkes dette endnu bestemtere ved det
resulterende Omdømme, at det subjective Øiemed har sin Objectivitet i
Midlet.“
% „Øiemed“, „Øiemedets“, „Øiemedet“ and „Omdømme“ all have the slashed
% Fraktur capital O / o at 600 dpi; both OCR witnesses read them as D.
% Closes the quotation opened on p. 224.

Betragter man Sammenstillingen af Kategorierne mechanisk, dynamisk og
teleologisk Virksomhed udelukkende fra Formsiden, og holder sig til den
vedtagne Betegnelse, ifølge hvilken „det subjective Øiemed“ skal falde
sammen med „det Enkelte“, og „det objective“ med „det Almindelige“, da
kan det ikke negtes, at „den teleologiske Virksomhed“ her er bleven
belyst paa en saadan Maade, at Begrebet i Teleologien tydelig sees at
staae høiere end Begrebet i den mechaniske og dynamiske Form. Men naar
det saa skal undersøges, hvorledes det egentlig gaaer til, at Begrebet,
skjøndt det synes at forblive paa hver Station, hvor det er, i det
Mechaniske som „mechanisk Object“, i det Dynamiske som „dynamisk
Object“, --- dog formaaer at hæve sig opad fra Trin til Trin og trods de
mange „Urdelinger“, for ikke at sige Udstykninger, bevare Eenheden med
sig selv, saa staae vi igjen ved Formalismen.
% „Urdelinger“ read at 600 dpi: U-r-d-…, the third sort being a Fraktur
% r and not a second d (Nielsen's rendering of the Hegelian pun
% Urtheil / Ur-Theilung). Not transcribed: the signature mark
% „Grundideernes Logik.“ + 15 at the foot of this page.

% --- p. 226 ---
\setcounter{footnote}{0}%
Det er virkelig en stærk Personification, en mod et saa ætherisk Væsen
som det rene logiske Begreb nærgaaende Anthropomorphisme, naar der
fortælles --- thi Fremstillingen nærmer sig her meget til en Fortælling
--- om Begrebet, at det „sætter sig, som subjectivt, imod den udvortes
Verden, men saaledes, at det betragter denne som den Modsætning, der skal
ophæves i det selv“. Det „betragtende“ Begreb, der her har stillet sig
„paa det subjective Standpunkt“, kommer \textit{mirabile dictu!}
ovenikjøbet til at betragte sig selv „som den Hensigt eller Plan, der vel
er fattet, men endnu ikke udført“; hvilken vidunderlig Smidighed et
saadant logisk Begreb dog maa besidde!\footnote{For at Læseren ikke skal
troe, at Hegel selv taler tydeligere, hidsættes Følgende til en
Sammenligning. „Der Zweck ist nämlich der an der Objektivität zu sich
selbst gekommene Begriff; die Bestimmtheit, die er sich an ihr gegeben,
ist die der \emph{objektiven Gleich\-%
gültigkeit und Aeußerlichkeit} des Bestimmtseyns; seine sich von sich
abstoßende Negativität ist daher eine solche, deren Momente, indem sie
nur die Bestimmungen des Begriffs selbst sind, auch die Form von
objektiver Gleichgültigkeit gegen einander haben \dots{} Der Zweck ist in
ihm selbst der Trieb seiner Realisirung; die Bestimmtheit der
Begriffs-Momente ist die Aeußerlichkeit, die \emph{Ein\-%
fachheit} derselben in der Einheit des Begriffs ist aber dem, was sie
ist, unangemessen und der Begriff stößt sich daher von sich selbst ab.
Dieß Abstoßen ist der \emph{Entschluß} überhaupt, der Beziehung der
negativen Einheit auf sich, wodurch sie \emph{ausschließende Ein\-%
zelnheit} ist; aber durch dieß \emph{Ausschließen entschließt} sie sich,
oder schließt sich \emph{auf}, weil es \emph{Selbstbestimmen, Setzen
seiner selbst ist}.“ Wissenschaft d. Logik 2ter Theil S. 219--20.}
% printed mark: *)  — printed „besidde ! *)“. The German of the note is
% Fraktur, so it is NOT italicised; its emphasis is letterspacing and is
% given with \emph{}. No fat-face was found: the heavy look of the
% capitals in „Entschluß“ and „Ausschließen“ is the ordinary weight of
% the Fraktur capitals, checked against the plain „Abstoßen“ and
% „Selbstbestimmen“ on the same lines. Only „auf“ is spaced in „schließt
% sich auf“ — as printed, not regularised. „mirabile dictu!“ is antiqua.

Feilen er ikke den, at her tales prægnant om det sig bevægende Begreb;
thi dersom Udtrykket Begreb udtrykkelig var valgt og fastholdt saaledes,
at det virkelig kunde betegne den i Alt virksomme, Alt i sig begribende
Væsenseenhed, hvorfor skulde man saa ikke ligesaavel kunne tale om det
alvirksomme Begreb, som om den alvirksomme Fornuft? Nei, Feilen, den for
en Logik høist betænkelige Grundfeil, er netop
% „høist“ as printed: the o's counter is filled at 600 dpi where the
% plain o's of „for“ and „Logik“ on the same line are open. (pp. 1--67
% have „hoist“ without the slash; the book is not self-consistent and is
% not made so here.) The final e of the first „alvirksomme“ carries an
% inked blemish that both tesseract witnesses read as an acute; it is
% joined to the letter body and is not a separate sort.
% --- p. 227 ---
\setcounter{footnote}{0}%
den, at Begrebet selv ikke begribes saaledes, at Begrebsudviklingen
retfærdiggjør den Selvstændighed, der i de prægnante Talemaader tillægges
det. Efter Talemaaderne at dømme skulde man troe, at Logikens Begreb
virkelig var et guddommeligt Væsen, ja Gud selv --- thi Gudsbegrebet er
jo allerede med i den rene Væren ---; men Begrebsudviklingen beviser
tydeligt nok, at Begrebet fra Først til Sidst har sin Vorden, Tilværen og
Bevægelse i en menneskelig Tænkning, i en Tænkning, der ikke blot er
menneskelig, men i den Grad menneskelig, at den ved at abstrahere fra sig
selv endog mener at kunne gjøre sig guddommelig. Grundfeilen er, at
Tankehandlingen uden videre bliver en Naturhandling, den logiske
Formudvikling en reel Sagudvikling, Idealsystemet Realsystem, og omvendt.

Lægges Eftertrykket paa Modsætningen mellem Ideal- og Realsystemet, vil
hiint eensidig fremstille Friheden, dette Nødvendigheden. Den logiske
Frihed beroer da derpaa, at hver Kategorie, hver Begrebsbestemmelse og
hvert særskilt Begreb udvikles og klares i Belysning af det hele
uendelige Begrebssystem. Saaledes kan f. Ex. Begrebet Væren ligesaalidt
være Begreb i den guddommelige som i den menneskelige Tænkning,
medmindre det ved en særegen Tænkningsact hæves frem for sig i
Modsætning til Ikke-Væren; saaledes kan Begrebet Grund ikke være Begreb,
medmindre Grunden udtrykkelig begribes i sin Forskjelseenhed med Begrebet
Existens; kort: Idealsystemet kan fra denne Side betragtes som en
guddommelig Logik, et System af Tanker, Begreber og Ideer, der bevæge sig
i evig Tænken og hvile i evig Viden.

I Modsætning til den logiske Frihed beroer den reale Nødvendighed paa
Begrebets umiddelbare Eenhed med sin Existens, altsaa paa dets
Forvandling fra et Vidensbegreb til et Magtbegreb: Realsystemet er et
System af Kategorier i antilogisk Form; her er i objectiv Forstand intet
Andet end
% Not transcribed: the signature mark 15* at the foot of this page.
% --- p. 228 ---
\setcounter{footnote}{0}%
Enkeltexistenser og deres Vexelvirkninger. Disse Vexelvirkninger med
tilsvarende Dannelser, Omdannelser og Forandringer ere deels successive,
deels simultane; her er Bevægelser i Tid og Rum; her er endelig
Modsætning af Bevægelse og Hvile.

Da Formaalet for Idealsystemet er en Udvikling af Begreberne som
Begreber, og Realsystemets Formaal en Virkeliggjørelse af Begreberne i
Existens, saa er Modsætningen mellem Ideal- og Realsystemet en
teleologisk Modsætning; hiint er et Videns, dette et Magtens Rige. Den
teleologiske Modsætning hviler i den teleologiske Eenhed; denne Eenhed er
Subjectivitetens altomfattende Idee. Existensen kunde umulig være det
realiserede Begreb, dersom der intet Begreb var, der kunde realiseres.
Men i det Existerende er Begrebet ophævet; Magtbegrebet, det factiske
Begreb, er et Existerende, og dette Existerende et Begreb, som intet
Begreb er. For at Begrebet skal kunne realiseres, maa det være; for at
det kan være, maa det begribes: Idealsystemet med sine Begreber er
Middel, d. e. væsentligt Medium for Realsystemet. Omvendt kan intet
Begreb udvikles, articuleres eller sættes som Existensbegreb, medmindre
det i sin Udvikling oprindelig har optaget det Indbegreb af antilogiske
Reflexer, hvorved dets existentielle Eiendommelighed ene og alene kan
bestemmes: Realsystemet er i denne Henseende Middel, d. v. s. virkeligt
Medium for Idealsystemet.

Eenheden af den reale Nødvendighed og den logiske Frihed er, paa Grund af
den uendelige Vexebestemmelse af det væsentlige og det virkelige Medium,
saa gjennemgaaende, at Nødvendigheden paa ethvert Punkt udtrykker
Friheden, og Friheden forklarer Nødvendigheden. Thi hvad er vel i denne
Sammenhæng Friheden? Den ontologiske Objectiverings subjective
Væsenhed. Derfor maae alle Existenser og existentielle Vexelvirkninger,
skjøndt de i den virkelige Existeren ere antikategoriske, dog opfattes i
kategorisk Form, bestemmes som Kjendsgjerninger, der fyldestgjøre
kategoriske Fordringer. Fri\-%
% „Vexebestemmelse“ as printed, without the l of „Vexelbestemmelse“; all
% three witnesses agree and the image shows no l. Transcribed as printed.
% --- p. 229 ---
\setcounter{footnote}{0}%
heden i Idealsystemet er uden Vilkaarlighed; her tænkes ikke andre
Tanker, her udvikles ikke andre Begreber, her klares ikke andre Former,
her haves ikke andre Ideer, end saadanne, der --- i den uendelige
Mægling af Middelbart med Umiddelbart --- deels middelbart, deels
umiddelbart fordres, for at Existens kan fremgaae af Existens ved at
fremgaae af Grunden, Virkeligheden af det Virkelige ved at fremgaae af
Muligheden, og Naturhandlingen bestemmes som kategorisk Magthandling.
Nødvendigheden i Realsystemet er meer end tilfældig. Den tilfældige eller
blot hypothetiske Nødvendighed (dersom $a$ er, saa er $b$) viser allerede
Vexelbestemmelsen af Kategorie og Kjendsgjerning: er Forudsætningen
factisk, saa er Følgesætningen kategorisk. Den virkelige Sammenhæng af
existerende Ting beroer paa en væsentlig Sammenhæng af
Existenskategorier; Tingenes System er væsentlig et Grundenes System, thi
Tingen er den tilstrækkelige Grunds existentielle Udtryk;
Virkelighedernes System er paa tilsvarende Maade et Mulighedernes System,
Individernes System et Arternes og Slægternes System o. s. v. Men
Grundenes, Mulighedernes, Slægternes og Arternes Systemer med deres
saavel reelle som ideelle Sammenhæng klares i Idealsystemet, altsaa i
Friheden.
% „Friheden“ is broken „Fri= / heden“ across the p. 228/229 break and
% joined; the sentence runs on, so no blank line at the marker.

I Idealsystemet ere alle Naturhandlinger reflecterede som mægtige
Tankehandlinger, hvis Bevægelse er logisk; i Realsystemet ere alle
Tankehandlinger reflecterede som kategoriske Magthandlinger, hvis
Bevægelse er antilogisk. Modsætningseenheden, den articulerede
Vexelbestemmelse af Tankehandling og Naturhandling, en
Vexelbestemmelse, hvori det Antilogiske er Middel, naar det Logiske er
Maal, og omvendt, er \emph{teleo\-%
logisk}. Den klare, teleologisk gjennemsigtige Eenhed af Videns og
Magtens Rige er Ideen selv, et Sandhedens Lys i Subjectiviteten.
% „a“ and „b“ are antiqua letter-symbols, set in math mode after the
% house treatment of $A$, $B$. The letterspacing of „teleologisk“ runs
% across the printed line break and both halves are spaced; the second
% „teleologisk“ two lines later is NOT spaced. The rest of the page
% below this line is blank: the division closes here.

% --- p. 230 ---
\setcounter{footnote}{0}%
\subsubsection*{c) Det logiske Ideal: Begriben og Viden.}
% Display head as printed, letterspaced, with an antiqua „c)“. The
% synopsis (p. 5) describes this third limb of C only by its matter
% („den fuldendte Begriben er teleologisk“) and gives no title; the
% printed head stands in the body. Sub-head level: \subsubsection*,
% matching a) on p. 143.

At det i Logiken som i enhver anden Videnskab væsentlig kommer an paa
tilforladelige Udgangspunkter og paalidelige Kjendemærker, naar ellers
Undersøgelser skulle kunne anstilles, forstaaer sig af sig selv. Det er
den psychologiske Subjectivitet med dens Forudsætninger og Betingelser,
hvorfra vi i Grundideernes Logik bestandig maae gaae ud, og hvortil vi
igjen maae vende tilbage. Det kommer ved de logiske Analyser Altsammen an
paa, hvorledes de som Kjendsgjerninger givne Udgangspunkter benyttes: at
blive i Endeligheden tilfredsstiller ikke; at abstrahere fra alt Endeligt
og lade Tankerne svæve i det tomme Uendelige tilfredsstiller heller ikke.
Her er ved Behandlingen af Forholdet mellem den fuldkomne Viden og Videns
psychologiske Factorer i Logiken Noget, der minder om Behandlingen af de
saakaldte „ubestemte Former“ i Mathematiken; ingen af de endelige
Forudsætninger kan ligefrem beholdes som Kjendsgjerning, men enhver
Kjendsgjerning kan og maa benyttes som Moment: Opgaven er som at regne
med uendelige Størrelser.

Umiddelbart og ligefrem kan Ingen sige, hvormeget $0 . \infty$ gjælder,
eller hvilken Værdi man maa tillægge $\infty - \infty$; ligesaa kan
heller Ingen umiddelbart og ligefrem afgjøre, hvorledes et Realbegreb
gaaer op i sin Existens, og hvilke Bevægelser Tankehandlingerne maae
fremkalde i den rene Idee. Ikke ved et Spring, men ved en til
Functionernes Eiendommelighed svarende Overgang maa den begribende
Subjectivitet i hvert Tilfælde bevæge sig fra det Endelige til det
Uendelige, dersom det ellers skal lykkes den at finde Tilbageveien fra
det Uendelige til det Endelige. Til Belysning af de Fremskridt, som den
menneskelige Aand, trods alle Feilgreb, Vildfarelser og Eensidigheder,
ogsaa i denne Henseende har gjort, vil Aandsphænomenologien med
Philosophiens Historie bestandig vedblive at afgive nye og væsentlige
Bidrag. Triviale Betragtninger som disse: at Philosophien endnu er
% „ubestemte“ with a lower-case u, as printed (the ABBYY witness has it
% right, tesseract capitalises it). The mathematics: the print gives
% „0. ∞“ with a low multiplication point on the baseline, and „∞ — ∞“
% with an em-length dash; both read off the image at 600 dpi and set in
% math mode.
% --- p. 231 ---
\setcounter{footnote}{0}%
i sin Barndom, eller at den af Alderdom gaaer i Barndom, at den endnu
ikke har ført til noget Vist, eftersom den ene Philosoph mener Et og den
anden et Andet, naar den Ene mener, at man skal holde paa
Contradictionsprincipet og den subjective Logik, den Anden, at kun
Dialektiken og den objective Logik formaae at udtrykke det Absolute;
naar den Ene mener, at det er Empirien, den Anden, at det er
Speculationen, der alene fører til Sandhedens Erkjendelse, o. s. v. ---
kan Videnskaben neppe finde det Umagen værd at
gjendrive\footnote{„Kjære Jensen! De forundrer Dem over, at jeg, uagtet
Philosophernes vedvarende Stridigheder, tillægger Philosophien Realitet.
Troer De, vil jeg til Svar spørge Dem, at nogen Philosoph kunde komme i
Strid med Dem? De er vist for beskeden til at vente det. Mener De da ei,
at Philosophernes Stridigheder maae grundes paa en gjensidig
Forstaaelse, en Enighed; at deres Stridigheder ere Kjendsgjerninger, som
forudsætte, at de have en fælleds Grund at staae paa?“ (Poul
Møller.)}.% printed mark: *)

Den nyere Philosophie er, hvad Subjectivitetsproblemet angaaer,
orienteret efter en ganske anden Maalestok, end den gamle nogensinde har
kunnet være det. Med Subjectivitetsproblemets alsidige Belysning følger
naturligviis en tilsvarende Belysning af Objectivitetens Problemer. Naar
en Fichte i alle Kategorier kun seer paa det Subjective, medens en Hegel
i dem alle væsentlig seer det Objective, da er det ikke Tankeløsheden,
men tvertimod den vedholdende, fra et eiendommeligt Udgangspunkt
gjennemførte Tænkning, der bevirker slige Eensidigheder. Det er i
Grunden een og samme Uendelighedsanalyse, her er anstillet fra to
modsatte Sider. Eller hvorledes vilde det være muligt at lade Videns
almene Begreber fremstaae i Subjectivitetens Lys, dersom en subjectiv
Tænken, Reflecteren og Virken ikke paa uendelig Maade gjennemtrængte dem
alle? Hvorledes vilde det være muligt at udfolde Kategorierne i et saa
objectivt System, at det liv\-%
% --- p. 232 ---
\setcounter{footnote}{0}%
agtig seer ud, som om Sagens Begreb var Sagen selv, og
Tankebevægelsen Eet med den virkelige Sagbevægelse, dersom disse vor
Videns aprioriske Objectiveringsformer ikke i uendelig Forstand vare
ligesom gjennemlynede af Virkelighedsglimt fra den objective Verden?

Naar nu fremdeles, i skarp Modsætning saavel til den objective som til
den subjective Gjennemførelse af Uendelighedsanalysen, Endeligheden og
det Endelige paany accentueres, --- hvad enten det iøvrigt skeer paa en
finere Maade ved Hjælp af „Realerne“ og den mathematiske Discretion
eller i grovere Maneer ved Hjælp af Atomer og andre materialistiske
Reqvisiter saa er en saadan Reaction paa de endelige Kjendsgjerningers
% the em-dash that opened this parenthesis after „accentueres,“ is not
% answered by a closing dash before „saa er“: as printed.
Vegne just egnet til paa en forfriskende og velgjørende Maade at
frembrage de Virkelighedspunkter og articulerende Mellembestemmelser,
som hinc for tidlig afsluttede, af et eensidigt Princip ledede
% PRINTER'S DEFECT: „hinc“ as printed --- a wrong sort, c for e; the
% sense requires „hine“. Read at 600 dpi and calibrated on the same
% page against the certain c of „Discretion“ (continuous top arm) and
% the certain e of „Maade“ and „Mellembestemmelser“ (detached top-right
% flag, closed bowl). The glyph is unambiguously the c sort.
Uendelighedsanalyser endnu ikke have formaaet at assimilere. Ogsaa i
denne Anledning maae vi minde om Poul Møllers Ord, „at den sande
philosophiske Forsken er et gjennem flere Sekler sammenhængende Værk, en
uafbrudt Kjæde, hvis følgende Led forudsætte de foregaaende som deres
nødvendige Betingelse“. --- „Maalet, den absolute Viden eller den sig
som Aand vidende Aand“ --- saaledes slutter Hegel sin
„Phä\-%
nomenologie des Geistes“ --- „har til sin Vei Erindringen om Aanderne,
% German in Fraktur, so not italicised.
som de i sig selv ere, og fuldføre deres Riges Organisation. Dens
Opbevaren i Henseende til deres frie, i Tilfældighedens Form
fremtrædende Tilværelse er Historien, men i Henseende til deres begrebne
Organisation den phænomenale Videns Videnskab; begge tilsammen, den
begrebne Historie, danne den absolute Aands Erindring og Domsted, den
\emph{Thrones} Virkelighed, Sandhed og Vished, uden hvilken den vilde
% „Thrones“ letterspaced in Fraktur (clear gaps T-h-r-o-n-e-s against
% the tight setting of „den“ and „Virkelighed“ on the same line);
% spacing.py did not flag it.
være det Livløse, Eensomme; kun af dette Aanderiges Kalk fremsprudler
Uendeligheden for den.“

I nøieste Sammenhæng med den Opgave: af det Psycho\-%
% --- p. 233 ---
\setcounter{footnote}{0}%
logiske at udklare det Ontologiske i den vidende Subjectivitet, staaer
en anden, ligesaa væsentlig Opgave, den nemlig: at bestemme Forholdet
mellem det Menneskelige og det Guddommelige i den ontologiske Viden
selv. Denne Opgave har sine Vanskeligheder deri, at enhver Usikkerhed
angaaende Vexelbestemmelsen af det Endelige og det Uendelige indenfor
den menneskelige Bevidstheds egne Grændser strax medfører en tilsvarende
Usikkerhed i Bestemmelsen af Forskjellen mellem den menneskelige og den
guddommelige Viden.

Det logiske Ideal, Vidensidealet, i dets rene Almindelighed maa da,
saavel fra Endelighedens som fra Uendelighedens eensidige Standpunkt,
nødvendig blive misforstaaet og forqvaklet. Ved at blive siddende midt i
Endeligheden vil Subjectiviteten kunne see feil af Idealet paa to
Maader: deels ved en tankeløs og selvmodsigende Opfattelse af den
guddommelige Videns Fuldkommenheder, deels ved en borneret Fastholden af
den menneskelige Videns psychologiske Betingelser. I første Tilfælde
bliver Idealet meningsløst, i sidste Tilfælde bliver det forkastet. Er
en Viden ikke mulig uden paa menneskelige Vilkaar, maa enhver vidende
Subjectivitet absolut have en Hjerne, ligesom Mennesket, fem Sandser,
som Mennesket, o. s. v., saa følger heraf, at Viden efter sin Natur er
og maa være relativ, at der kun kan gives en Gradforskjel mellem lavere
og høiere Videns Trin, men ingen absolut eller fuldkommen Viden. Bliver
den vidende Subjectivitet derimod ved Hjælp af Abstractioner og
speculative Reductioner for let færdig med det Endelige, og kommer i den
rene Uendelighed til en Forvexlen af sin tomme Skuen med den ideale
Viden, da vil den i mystisk Begeistring kun altfor let sammenblande det
Guddommelige med det Menneskelige og tage Idealet forfængeligt.
„Speculation er Alt, er i Gud at skue, at betragte det, som er.
Videnskaben selv har kun forsaavidt Værd, som den er speculativ, det er
Contemplation af Gud, som han er. Men den Tid vil komme,
% --- p. 234 ---
\setcounter{footnote}{0}%
da Videnskaberne mere og mere ville ophøre, og den umiddelbare
Erkjendelse indtræde. Kun i den høieste Videnskab lukker det dødelige
Øie sig, idet ikke mere Mennesket seer, men den evige Seen selv er
seende i Mennesket.“ (Schelling).

Hvad disse Misforstaaelser i logisk Henseende have at betyde, vil i det
Følgende bedst kunne belyses, naar Vidensidealet charakteriseres med
særligt Hensyn til \emph{den intuitive og den discursive Viden}.
% the letterspacing begins at „den intuitive“ --- „til“ is not spaced
% --- and runs on across the line break through „Viden“.

\subsubsection*{α) Den intuitive Viden.}

Det vil i nærværende Sammenhæng --- til Forklaring af Vidensidealet ---
ikke være uden logisk Interesse at kaste et Blik paa de strengt
orthodoxe Dogmatikeres Opfattelse af Guds Alvidenhed, og i denne
Henseende gjøre sig bekjendt med \textit{testimonium spiritus} af en
ældre Dato. Det Held, hvormed den nyere speculative Dogmatik har vidst
at samle og benytte philosophiske Guldkorn, --- og naar det blot gjælder
om at overtrække de gamle Dogmer med tynd Forgyldning, kan man hjælpe
sig med lidt Guld --- er jo allerede paaskjønnet i vort Foregaaende; her
skulle vi endnu, til en Bekræftelse paa, hvorledes den Hellig Aand den
Gang, da det kirkelige Lærebegreb endnu holdt sig reent fra enhver
philosophisk Besmittelse, maa have taget sig af de „dogmatiske
Oplysninger“, anføre nogle Træk.

Til Bestemmelse af Vidensidealet frembyder Dogmatiken fra det 17de
Aarhundrede interessante Hovedpunkter. Alvidenheden“ --- hedder
% PRINTER'S DEFECT: the opening „ before „Alvidenheden“ is dropped ---
% the word follows the full stop of „Hovedpunkter.“ with nothing but a
% word-space, verified at 600 dpi. Transcribed as printed, so this page
% carries one unpartnered closing “.
det\footnote{\textit{Hutterus Redivivus.} Leipzig 1836. S.
138--139.}% printed mark: *)
\ --- „er med Hensyn paa sit Object: 1) \emph{en nødvendig} Viden
(\textit{scientia necessaria, qua Deus semetipsum atque in se ipso
omnium rerum necessitatem perspicit}), 2) \emph{en fri Viden}
(\textit{scient. libera [repræsentatio visionis s. intuitionis], qua
Deus omnes res præter ipsum vere existentes vere novit}), 3) \emph{en
betinget}
% the compositor's letterspacing is inconsistent here: in 1) only „en
% nødvendig“ is spaced and „Viden“ is not, while in 2) „en fri Viden“ is
% spaced throughout. Transcribed as printed.
% --- p. 235 ---
\setcounter{footnote}{0}%
\emph{Viden} (\textit{scient. media [sc. de futuro conditionato s.
futuribili, hypothetica, simplicis intelligentiæ], qua Deus perspicit
omnia, quæ, positis quibusdam conditionibus, evenire potuissent}). Med
Hensyn paa det Fuldkomne i Erkjendelsesmaaden er Guds Viden: 1)
\emph{intuitiv} (\textit{intuitiva [pura, immediata] i. e. sine sensu,
imaginibus, abstractione, discursu et ratiocinio}), 2) \emph{simultan}
(\textit{simultanea}), 3) \emph{tydelig} (\textit{distinctissima}) og
følgelig 4) \emph{fuldkommen sand} (\textit{verissima}).“

Underkaste vi nu disse under eet Overblik sammenstillede
Vidensfuldkommenheder en nærmere Prøvelse, saa bliver det os ret
indlysende, at „det dogmatiske Talent“ er kommet til at stille disse
Fuldkommenheder op, ikke ifølge nogen egentlig Indsigt i Videns
Functioner, men ved et umiddelbart Spring fra Endeligt til Uendeligt. At
dette Spring, logisk forstaaet, er skeet i al mulig Tankeløshed, derfor
have vi det \textit{testimonium spiritus}, --- idet vi saaledes tillægge
\textit{testimonium spiritus} en „theoretisk Betydning“ --- at den ene
af de opstillede Fuldkommenheder ligefrem tilintetgjør den anden. „Med
Hensyn paa sit Object“ skal den guddommelige Viden vide de samme Ting
paa tre forskjellige Maader. „Gud gjennemskuer sig selv og i sig selv
alle Tings Nødvendighed.“ Her synes det at have foresvævet
Dogmatikerne, hvad Proklus lærer i sin „platoniske Theologie“: „Gud
erkjender det Deelte udeelt, det Timelige tidløst, det Ikke-Nødvendige
nødvendigt, det Foranderlige uforanderligt og overhovedet alle Ting
fortræffeligere end efter deres Orden.“ Men hvad betyder nu \emph{den
frie Viden}? Svar: en Viden af frie, d. e. selvstændige, udenfor Gud
existerende Ting. Men hvorledes kunne Tingene udenfor Gud vides af Gud
at være med Frihed, naar Gud i sig selv dog veed, at de ere med
Nødvendighed? Er Friheden i den guddommelige Viden uden videre det
Samme som Nødvendigheden, hvorfor bruger man da to saa forskjellige
Udtryk som Frihed og Nødvendighed? Er Friheden Et og
% --- p. 236 ---
\setcounter{footnote}{0}%
Nødvendigheden et Andet, hvorledes gaaer det da til, at de samme Ting
kunne være baade frie og nødvendige? Siger man: „Tingene ere i deres
Grund og Begrundelse nødvendige; men i deres udvortes og endelige
Existens tilfældige og i denne Tilfældighed frie“, saa maae vi igjen
spørge, hvorledes Tilfældigheden og Friheden da forholde sig til
Grunden? Thi dersom der i det for Menneskenes Øine Tilfældige virkelig
er Noget, om hvilket Gud veed, at det i Grunden er løsrevet fra
Nødvendigheden og sat til alene ved hans frie Villie
(\textit{libertas}), da ere de udenfor Gud existerende Ting jo med det
Samme løsnede fra Nødvendighedens Grund; den \emph{nødvendige} Viden,
den Viden, hvormed Gud i sig gjennemskuer alle Tings Nødvendighed
(\textit{omnium rerum necessitatem}), maa altsaa være en usand Viden.

Ved Bestemmelsen af „den frie Viden“ er Dogmatiken udtrykkelig gaaet ud
fra den Forudsætning, at der er en Verden af virkelige ligeoverfor Gud
existerende Ting (\textit{res præter ipsum vere existentes}) og at Gud
erkjender disse Ting i deres Endelighed, idet han erkjender dem,
saaledes som de virkelig ere (\textit{vere novit}). Dette den
guddommelige Videns Hensyn til Endeligheden og dens Betingelser træder
endnu tydeligere frem under Synspunktet af den hypothetiske Viden;
herved fremkommer da en endelig Adskillelse ikke blot imellem Mulighed
og Virkelighed, men i Muligheden selv imellem de Muligheder, der under
givne Betingelser blive virkelige, og de Muligheder, der ikke blive det,
fordi Betingelserne ikke indtræde. I alt dette er her en vis
Naturlighed. Dogmatikeren tager Virkeligheden ligefrem, saaledes som den
viser sig for menneskelige Øine, fremhæver dens forskjellige Sider, og
saa forstaaer det sig jo som af sig selv, at Gud, som er alvidende, maa
kjende alle mulige og alle virkelige Ting fra alle mulige og virkelige
Sider. Hertil svarer da ogsaa den under „den frie Viden“ faldende
Underafdeling: Erindring (\textit{reminiscentia}), Syn
(\textit{visio}) og Forudvidenhed (\textit{præscientia}).

% --- p. 237 ---
\setcounter{footnote}{0}%
Men Alt, hvad der saaledes stilles op under de to Rubriker: „den frie“
og „den betingede Viden“, maa rives ned igjen, naar
„Erkjendelsesmaaden“ i og for sig skal nærmere bestemmes. Her er nu
„den intuitive Viden“ den allervigtigste Kategorie: hvorledes
charakteriseres da „den intuitive Viden“? Den er, hedder det, aldeles
umiddelbar. Umiddelbarheden er da her nærmere at forstaae saaledes, at
„den intuitive Viden“ for det Første er uden Sandsning og uden Billeder
(\textit{sine sensu, sine imaginibus}); men allerede disse negative
Bestemmelser maae vække Betænkeligheder. Er den guddommelige Viden uden
Sands og Sandsen (\textit{sine sensu}), hvorledes kan den da bestemmes
ved en Seen, et Syn (\textit{visio})? Skal den guddommelige Skuen være
umiddelbart Eet med det Skuede, saa at her slet ingen Forskjel er, da
maa man nødvendig antage Eet af To: enten maa det Skuede forsvinde i den
rene Skuen, som da bliver en Skuen, der kun skuer sig selv, en Skuen af
en Skuen, en uendelig tom Skuen; eller ogsaa maa Skuen gaae op i det
Skuede saaledes, at Videns Gjenstand træder i Videns Sted, og Viden
% a stray raised mark stands immediately after the comma of „Sted,“ ---
% a small slanted stroke at x-height, apparently a piece of loose type
% or dirt, not identifiable as a sort. Not transcribed.
forsvinder. Lader os engang endnu forsøge det med den intuitive Viden af
en Steen: enten maa Stenen ganske forsvinde for Intuitionen, eller
Intuitionen for Stenen; altsaa: dersom Stenen er, saa er den intuitive
Viden ikke; dersom den intuitive Viden er, saa er Stenen ikke. Hvad der
gjælder om Stenen, gjælder naturligviis om enhver endelig og virkelig
existerende Ting; altsaa: dersom den existerende Ting er, saa er den
intuitive Viden ikke; dersom den intuitive Viden er, saa er den
existerende Ting ikke. Men naar den guddommelige Viden er intuitiv, og
en intuitiv Viden af existerende Ting en Modsigelse, hvorledes skal man
saa forstaae en Sætning saa dogmatisk som denne: \textit{scientia libera
est repræsentatio visionis, qua Deus omnes res præter ipsum vere
existentes vere novit}?

Skal der forbindes nogen Mening med det, Dogma\-%
% --- p. 238 ---
% FOLIO MISPRINT: this page's folio is set „233“ instead of „238“ (a 3
% for an 8, read at 600 dpi). The text is continuous with p. 237 before
% and p. 239 after, so only the numeral is wrong; this IS printed p. 238.
\setcounter{footnote}{0}%
tiken lærer om Guds intuitive Viden, da maa man nødvendig forudsætte, at
der imellem Skuen og det Skuede, imellem Viden og Videns Gjenstand er en
Forskjel, og at denne Forskjel har umiddelbar Gyldighed; Et er altsaa
Skuen, et Andet det, som derved skues: Et er Viden, et Andet er Videns
Gjenstand. Under denne Forudsætning er det da, idetmindste til en vis
Grad, muligt at forbinde en Mening med den Antagelse, at Gud har en
intuitiv Viden af alle existerende Ting. Forbilledet er igjen
menneskeligt. Vi Mennesker kunne jo paa en Maade ogsaa have en intuitiv
Viden af visse, let overskuelige Gjenstande (af denne Steen, denne
% the opening parenthesis before „af denne Steen“ is a broken sort:
% only the upper arc printed, leaving a short slanted stroke at cap
% height. Transcribed as „(“.
Plante, dette Huus), naar de stilles i en heldig Belysning. Forskjellen
imellem den menneskelige og den guddommelige Viden er naturligviis den,
at hiin i hvert enkelt Tilfælde indskrænker sig til at opfatte en enkelt
Gjenstand fra en enkelt Side, medens denne paa een Gang
(\textit{scientia simultanea}) opfatter alle mulige Gjenstande lige
klart og lige tydeligt (\textit{scientia distinctissima}) fra alle
mulige Sider. Men komme vi nu ikke i Forlegenhed med den dogmatiske
Erklæring, at Guds intuitive Viden er uden Sands og Sandsning
(\textit{sine sensu}), uden Billede og Forestilling (\textit{sine
imaginibus})?

Dersom den intuitive Viden absolut forudsætter Forskjelseenhed af Viden
% a stray raised dot stands in the space between „Dersom“ and „den“;
% ink or loose type, not punctuation. Not transcribed.
og Videns Gjenstand, saa er en intuitiv Viden uden Sands og Sandsning en
meningsløs Modsigelse. Hvad er vel Sandsning Andet end en umiddelbar
Opfattelse, altsaa en umiddelbar Viden af Noget, der umiddelbart
bestemmer denne Viden og derved udgjør dens Gjenstand? At den
guddommelige Sandsning maa være væsentlig forskjellig fra den
menneskelige, forstaaer sig af sig selv; men dersom Dogmatiken ikke tør
tillægge Gud Sandsning, fordi den guddommelige Sands er saa uendelig
% „uendelig“: the second e failed to print altogether and the d printed
% only in part, so the impression reads „uen?~lig“ with a lowercase-wide
% gap. Read as „uendelig“ --- the same phrase, „uendelig forskjellig fra
% den menneskelige“, stands complete in the next sentence on p. 239.
forskjellig fra den menneskelige, hvorledes tør Dogmatiken da
overhovedet tillægge Gud
% --- p. 239 ---
\setcounter{footnote}{0}%
nogen Viden, eftersom den guddommelige Viden jo er og maa være uendelig
forskjellig fra den menneskelige?

Dogmatiken lærer, at Gud i intuitiv Viden erkjender Tingene uden
Billeder (\textit{sine imaginibus}). Men de, der uden videre opstille en
saadan Paastand, synes ikke saameget at være tilskyndede ved noget
\textit{testimonium spiritus sancti}, som snarere, ved en dunkel Følelse
af Theologiens videnskabelige Trang, fristede til at aflægge et
theoretisk Vidnesbyrd om, at de i Grunden hverken have tænkt over den
menneskelige eller den guddommelige Videns Betingelser. Hvordan skulde
der i al Verden kunne tænkes en Skuen af en tilværende Gjenstand, hvis
den ikke tør tænkes saaledes, at et Gjenstandens Billede formedelst
denne Skuen danner sig i den Skuende? Opgiver man denne Fordobling, og
lader Guds intuitive Viden anskue Gjenstanden som Gjenstand uden
mellemtrædende Billede, saa falder man med uundgaaelig Nødvendighed
tilbage til det Urimelige, at det, som allerede bemærket, enten maa være
Viden, der i Intuitionen fortærer Gjenstanden, eller Gjenstanden, der
absorberer Viden.

Det Urimelige i den opstillede Erklæring kan desuden, fra Dogmatikens
eget Standpunkt, sees paa flere Maader. Hvad bliver der, for ikke at
tale om den guddommelige Erindring (\textit{reminiscentia}), af Guds
Forudvidenhed (\textit{præscientia}), dersom der i den guddommelige
Viden ikke gives en Forskjel, og det en væsentlig Forskjel, imellem
Tingens Billede og Tingen selv? Falder Tingen absolut sammen med sit
Billede, da kan man vende og vride sig saa dogmatisk, som man vil: man
har aabenbart fornegtet den guddommelige Forudvidenhed; thi
Forudvidenheden beroer just paa, at Synet af Tingen, d. v. s. Tingens
Billede, gaaer forud for Tingen selv.

Er det virkelig en Sandhed og ikke en tom høitidelig Talemaade, at den
Hellig Aand --- „idet vi saaledes tillægge \textit{testimonium spiritus
sancti} ikke blot praktisk, men ogsaa theoretisk Betydning“ --- har
hjulpet saavel de ældre som de
% --- p. 240 ---
\setcounter{footnote}{0}%
nyere Dogmatikere med deres videnskabelige Opfattelse af „Guds intuitive
Viden“, „den altgjennemtrængende Seen“, da kommer det for os skrøbelige
Mennesker, som uden at kunne beraabe os paa overnaturlige Indskydelser
maae see Sagen med blot menneskelige Øine, unegtelig til at tage sig ud,
som om den Hellig Aand var noget skjødesløs med sine theoretiske
Vidnesbyrd; ja det seer undertiden ud, som om den i de Vidnesbyrd, den
aflægger i den nyere Tid, reentud sagt, havde glemt en heel Deel af de
Vidnesbyrd, den med yderlig Strenghed fastholdt i en ældre Periode.
Dette gjælder navnlig om Bestemmelsen af Forholdet mellem Tingene og
Tingenes Billeder. Var Guds „intuitive Viden“ i ældre Dage \textit{sine
imaginibus}, uden Phantasie og følgelig uden Phantasier, saa har den nu
i disse sidste Dage, til Erstatning for hine Mangler, undfanget ved
„Speculationen“ og bekommet sig et „Pleroma, et faderligt Pleroma“ af
sønlig-speculative Phantasier. „Det faderlige Pleroma, der i Sønnen
aabenbares som et nødvendigt oprindende Ideerige, forklares ved Aandens
frie kunstneriske Virken til et Herlighedens Rige (\textit{δόξα}), hvor
% Greek in upright antiqua; set as in the earlier batches.
de evige Muligheder lege for Guds Aasyn som magiske Virkeligheder, som
en himmelsk Hærskare af Syner, af plastiske Forbilleder for en
Aabenbaring \textit{ad extra}, til hvilken de ligesom begjære at
løslades.“

Lignende Vanskeligheder paatrænge sig, naar vi, istedenfor at see paa
det Anskuelige i Syn og Billede, særlig holde os til Forstandssiden, til
det, der angaaer den rene Tænken, Begriben og Indseen. Her vil nu
Dogmatiken have Abstractionen og den discursive Viden ligefrem udelukket
fra det Intuitive; den guddommelige Intuition er \textit{sine
abstractione, sine discursu, sine ratiocinio}. Hvad herved tilsigtes, er
klart; den guddommelige Viden kan ikke skride frem paa samme endelige og
middelbare Maade som den menneskelige. „Vor Viden er Stykværk“;
Sammenstykningen betinges af Abstractionen, af det Discursive i vore
Undersøgelser og Overveielser,
% --- p. 241 ---
\setcounter{footnote}{0}%
det Middelbare i vore Slutninger. Men Et er det, at den guddommelige
Viden ikke bruger det Middelbare paa samme endelige Maade som den
menneskelige, noget ganske Andet derimod, at den i sin Umiddelbarhed
skulde udelukke det Middelbare. Dersom Guds intuitive Viden virkelig er
uden al Abstraction, hvad følger saa heraf? At den maa være uden al
Tænken, altsaa tankeløs! Al Tænken beroer i den Grad paa Abstraheren og
paa Integreren af det Abstraherede, at Tænkningen nødvendig maa standse,
naar disse reent intellectuelle Virksomheder ophøre.

Uden Abstraheren og Tænken er Anskuelsens Umiddelbarhed en dyrisk
Umiddelbarhed. At anskue en Gjenstand saaledes, at Anskuelsen er fortabt
i dyrisk Umiddelbarhed, er, hvad man kalder at gloe. Dyret gloer paa den
fremmede Gjenstand; Koen gloer paa den nye Port o. s. v. Den dyriske
Gloen er netop charakteristisk derved, at dens Umiddelbarhed mangler den
for al Tænken og Forstaaen nødvendige Reflexion. Dyret seer Tingen uden
at kunne hæve sig til en Seen af denne sin Seen; dets Seen er i den
Forstand umiddelbar, at det ikke kan frigjøre sig fra Umiddelbarheden.
Paa hvilken Betingelse kan en saadan Frigjørelse da finde Sted? Paa
Abstractionens, kun paa Abstractionens! thi Abstractionen betinger
Negationen.

Hvormange Abstractioner maae ikke forudsættes, for at den, der betragter
en Gjenstand, skal kunne sige til sig selv: „her seer jeg en Plante“;
„her seer jeg en Steen“. Hvad Aandedrættet er for vort organiske Liv,
det er Abstractionen og dens Ophævelse for al Tænken og Begriben; al
Tænken og Begriben --- deri har Hegel ubetinget Ret --- gaaer igjennem
Mediationer; men Mediationen er en uendelig Mægling, hvorunder det
Middelbare kommer frem, for at gjennemformes, ophæves og forsvinde i en
høiere Umiddelbarhed. Overalt, hvor Talen er om en umiddelbar Viden, maa
man vel gjøre Forskjel imellem den Umiddelbarhed, hvormed der begyndes,
og den Umiddelbarhed, hvormed der
% signature mark at the foot of this page: „Grundideernes Logik.“ + 16.
% Not transcribed.
% --- p. 242 ---
\setcounter{footnote}{0}%
gjennem mange Mæglinger og Mellembestemmelser fuldendes. Jo flere
Mediationer en Viden har gjennemløbet, desto høiere er igjen den
Umiddelbarhed, som deraf vil resultere.

At denne Lov, efterdi det er en Videnslov, ligesaavel maa gjælde for den
guddommelige som for den menneskelige Viden, vil uden Vanskelighed kunne
tydeliggjøres. Antage vi med vore Dogmatikere, at „Guds intuitive Viden“
er en klar og tydelig Viden (\textit{scientia distinctissima}) af alle
enkelte existerende Ting i alle Enkeltheder; antage vi endvidere med
Dogmatikerne, at Guds intuitive Viden af hver enkelt existerende Ting er
umiddelbar (\textit{pura, immediata}), hvad følger saa heraf? Seet fra
Dogmatikens eget Synspunkt, vil heraf nødvendig følge, at der i Guds
intuitive Viden maa være ligesaa mange Intuitioner, ligesaa mange
umiddelbare Blikke, med andre Ord, ligesaa mange guddommelige Øine, som
der er existerende Ting og Enkeltbestemmelser ved de existerende Ting i
Himlen og paa Jorden. Enten Dogmatiken nu helst vil udtrykke sig mere
bibelsk, og tale om Guds Øie, altsaa om det Øie, Gud „\emph{har}“,
eller mere speculativt om det Øie, Gud „\emph{ikke har}“, saa bliver dog
Resultatet dette, at der i det ene guddommelige Øie maa være en uendelig
Mangfoldighed af Øine, i den ene almene Intuition en uendelig
Mangfoldighed af særskilte Intuitioner, i det ene altgjennemtrængende
Guddomsblik en uendelig Mangfoldighed af guddommelige Blikke.

Vil man nu aldeles tankeløst fastholde Mangfoldigheden af særskilte
Intuitioner i umiddelbar Modsætning til den umiddelbare Intuitionens
Eenhed, eller, da Modsigelsen er lige skjærende, enten man vælger den
ene eller den anden „dogmatiske Oplysning“, lade den umiddelbare
Intuitionens Eenhed gaae op i den umiddelbare Mangfoldighed af særskilte
Intuitioner, da forvandler man blasphemisk det guddommelige Øie til et
facetteret Øie, et uhyre Oldenborre-Øie, og forvexler Umiddelbarheden af
Guds intuitive Viden med Umiddelbarheden
% --- p. 243 ---
\setcounter{footnote}{0}%
i en dyrisk Gloen. Skulle Umiddelbarheder af denne Natur den Dag idag
endnu gjentages, indpræntes og beundres som videnskabelige
Indskydelser, da kan et \textit{testimonium spiritus} saare let
forvandle sig til et \textit{testimonium stupiditatis}. Slige
\textit{testimonia stupiditatis}, idet vi saaledes tillægge
\textit{testimonia stupiditatis} „ikke blot praktisk men ogsaa
theoretisk Betydning“, hidrøre aabenbart ikke fra religiøs
Dybsindighed, men, hvad navnlig de nye Opkog af gamle Randsagelser
angaaer, fra logisk
Letfærdighed.\footnote{„Den, der vil gjøre Lykke hos den litterære
Pøbel, maa ikke hæve sig til den rene Speculations Region, ikke heller
krybe paa den prosaiske Jord, men være liig en halvstækket Fugl, der
flyver lidt i Veiret af og til, men dog idelig plumper ned igjen.“
Denne fra Poul Møllers Strøtanker hentede „Klogskabsregel“ burde vel
især anbefales det „dogmatiske Talent“, der ærlig stræber, ikke just at
udgrunde en Sandhed, men --- ved Afbenyttelse af adskillige
\textit{testimonia} --- at opretholde en Opinion.}% printed mark: *)

\subsubsection*{β) Den discursive Viden.}

Stærk formedelst sin stærke Tro og fremfor Alt baaren af sit
Testimonium, kommer den theologisk-dogmatiske Bevidsthed let, som ved en
Luftning af „Inspirationen“, til en Art klimisk Svæven i det høie
stjernerige Himmelrum, hvor Vidensidealet viser sig omtrent som det
hvide Skjær af en uendelig blaa Tone; til Gjengjæld have vi da paa det
sandselig-realistiske Bevidsthedstrin et jordfast Standpunkt.

Den sandselig-realistiske Synsmaade lader sig, hvad Vidensidealet
angaaer, henføre til følgende tre Punkter: 1) det Eiendommelige ved vore
Sandser og deres Forhold til vor Tænkning; 2) den til det eiendommelige
Forhold imellem vor Sandsning og Tænkning svarende Eiendommelighed i vor
endelige og middelbare Viden; 3) Umuligheden af at tænke sig en Viden,
der skulde være af høiere Art end vor egen menneskelige.
% signature mark at the foot of this page: 16*. Not transcribed.

% --- p. 244 ---
\setcounter{footnote}{0}%
1) Hvad nu for det Første Eiendommeligheden ved vore Sandser og deres
Forhold til Tænkningen angaaer, kan Sensualismen beraabe sig paa den
Kjendsgjerning, at det Umiddelbare fra Begyndelsen af maa opfattes paa
umiddelbar Maade. „Hvad en Fornemmelse er, lære vi ved at fornemme;
hvad et Syn er, ved at see; hvad en Høren er, ved at høre. Men heraf
følger da, at vor Viden af, hvad Fornemmelse er, ganske og aldeles maa
rette sig efter Beskaffenhederne ved de menneskelige Fornemmelser.
Hvilke Fornemmelser et Dyr, f. Ex. en Snegl, vil kunne have, derom ere
vi kun istand til at danne os en Forestilling, forsaavidt vi forudsætte,
at der imellem Dyrets og Menneskets Fornemmelser maa være en væsentlig
Lighed. Det Samme gjælder om Syn og Hørelse; vi ere nødte til at
bedømme al anden Seen efter vor egen Seen, al anden Hørelse efter vor
egen Hørelse. Hvor ubetinget vi ere bundne til det Umiddelbare i hver
eiendommelig Sands, sees endvidere deraf, at hver Sands opfatter Sit;
hvad der sees, kan ikke høres, hvad der høres, ikke føles, o. s. v.
Hvormeget man end uddanner, skærper og væbner den enkelte Sands, saa er
det dog ikke muligt at tilveiebringe en jevn Overgang fra den ene Sands
til den anden: en skærpet Hørelse nærmer sig ikke til at blive Syn, et
skærpet Syn ikke til at blive Hørelse. Vistnok kan man paa en Maade
gjøre det Hørlige synligt, Tonerne f. Ex. ved Klangfigurer, men dette
skeer dog kun derved, at den samme Sag frembyder to forskjelligartede
Sider, af hvilke den ene svarer til den anden, den hørlige til den
synlige, eller omvendt. At Sandserne, der ere saa ueensartede med
hverandre indbyrdes, ogsaa ere ueensartede med Tænkningen, er ligesaa
afgjort. Psychologerne maae længe nok gjentage den Forsikkring, at
Forstanden er en Sands; de ville dog ikke være istand til at vise os,
hvad det er for et Sandseligt, Forstanden sandser. At gjøre Forstanden
til en Sands for det Usandselige, altsaa til en Sands for det, som ikke
kan sandses, er omtrent det
% the quotation opened at „Hvad en Fornemmelse er“ is still open at the
% foot of this page; it closes on p. 245 („\ldots\ end den
% psychologiske.“), outside this batch.
% --- p. 245 ---
\setcounter{footnote}{0}%
Samme som at give Afkald paa den: enhver Sands maa have sit
Sandselige; en Sands for noget Usandseligt er en Modsigelse. Der
tænkes over det Sandselige, ja der kan og maa tænkes ved Hjælp af det
Sandselige, men det Sandselige selv kan ikke tænkes.
% After „tænkes“ at the end of the third printed line there stands a faint
% detached mark at comma height, separated from the word by a full
% word-space; read as dirt on the plate, not as punctuation.
Som Tænkningen er, maa Forstanden være, og omvendt. En Forstand, der
enten skulde kunne undvære det Sandselige, eller gjøre sig til Eet med
det Sandselige, eller gjennemtrænge det Sandselige, er vort Væsen i
den Grad fremmed og uforstaaelig, at vi reent ud maae benegte dens
Realitet. Som Følge heraf maae vi da ogsaa benegte, at der skulde
gives Vidensbetingelser af anden Art end de menneskelige, eller nogen
vidende Subjectivitet af høiere Orden end den psychologiske.“
% This „…“ closes a quotation opened before p. 245; the partner „ falls in
% the preceding batch (pp. 231--244).

2) Paa det Eiendommelige i Sandsernes Forhold til hverandre indbyrdes
og til Tænkningen beroer, og dette er det næste Hovedpunkt, vor Videns
Eiendommelighed. „Vil man uden videre kalde den umiddelbare Viden
intuitiv, da kan man gjerne sige, at vor Viden paa alle Punkter
begynder med det Intuitive; men den intuitive Viden kommer da paa denne
Maade blot til at danne Forudsætningerne for den discursive; saaledes
maa der altid noget Umiddelbart til, der kan afgive Grundlaget for det
Middelbare. Det Middelbare forudsætter Mellemled; ethvert Mellemled er
et Middel, men Midlet er, for sig betragtet, umiddelbart. Som nu det
Middelbare forholder sig til det Umiddelbare, saaledes forholder sig,
hvad vor Viden angaaer, det Discursive til det Intuitive. Er Videns
Gjenstand en Totalitet, saa er vor Viden af det Hele middelbar,
forsaavidt den kommer istand ved en umiddelbar, men successiv
Opfattelse af de enkelte Dele og deres indbyrdes Sammenhæng: det
Deelvise og Middelbare falder her sammen med det Discursive. Omvendt
kan der ogsaa gives en vis umiddelbar Opfattelse af det Hele, nemlig et
Totalindtryk; det Middelbare og Discursive kommer da til at beroe paa
den deelvise Undersøgelse af Enkelthederne.“

% --- p. 246 ---
\setcounter{footnote}{0}%
„Forholdet mellem det Intuitive og det Discursive kan ifølge heraf
bestemmes ved Modsætningen imellem Analysen og Synthesen; er Synthesen
umiddelbart given for Anskuelsen og vor Viden forsaavidt intuitiv, saa
er Analysen discursiv; er det enkelte, ved Analysen fremhævede Element
umiddelbart tilværende for Anskuelsen og altsaa at opfatte intuitivt,
saa giver Synthesen, Delenes deelvise Forbindelse, Elementernes
successive Forening, os det Discursive. Det Samme kan sees af Forholdet
mellem Præmisserne og Conclusionen i enhver Slutning. Forsaavidt en
Præmis er indlysende for sig, opfattes den intuitivt; men netop den
Omstændighed, at her er en Præmis, at man altsaa maa gaae enten fra
Præmis til Præmis eller dog fra Præmis til Slutning, viser det
Discursive. Da nu det Discursive overalt er kjendeligt paa Fleerheden
af de forskjellige Udgangspunkter, paa det Deelvise, det Successive og
Middelbare, og det Intuitive med al dets Umiddelbarhed ligeledes er
deelviist, successivt og middelbart, nemlig forsaavidt den deelvise og
successive Umid\-%
% The m of „Umid-“ printed with its first stroke failed: the scan shows
% „U.nid=“, a low blob then a clean n. The blob and the n together span
% 51 px where a sound m on the same page spans 53, so the sort is an m
% with an unprinted left stroke, not a wrong sort.
delbarhed altid deelviis maa være betinget, saa indsees heraf, at det
Intuitive, langtfra at have nogen selvstændig Betydning ligeoverfor det
Discursive, tvertimod er at opfatte som den discursive Videns eget
Element, et Element, der bestandig viser ud over sig selv, og først
faaer sin Betydning derved, at de forskjelligartede Bidrag belyse
hinanden. Vi have altsaa, naar vi see paa Totaliteten, ikke en intuitiv
Viden for sig ligeoverfor eller ved Siden af en discursiv, men al vor
Viden er med sine forskjelligartede Umiddelbarhedselementer væsentlig
discursiv.“\footnote{% printed mark *)
Hos Kant svarer Adskillelsen af Sandselighed og Forstand, af Anskuelse
og Begreb, paa eiendommelig Maade til Adskillelsen mellem intuitiv og
discursiv Erkjendelse. „Reflectiren wir auf unsere Erkenntnisse in
Ansehung der beiden wesentlich verschiedenen
% the printed page break p. 246 | p. 247 falls here, inside the note
Grundvermögen der Sinnlichkeit und des Verstandes, woraus sie
entspringen, so treffen wir hier auf den Unterschied zwischen
Anschauungen und Begriffen \ldots{} Dieses ist der \emph{logische
Unterschied} zwischen Verstand und Sinnlichkeit, nach welchem diese
nichts, als Anschauungen, jener hingegen nichts, als Begriffe liefert
\ldots{} Auf den hier angegebenen Unterschied zwischen
\emph{intuitiven und discur\-%
siven} Erkenntnissen, oder zwischen Anschauungen und Begriffen gründet
sich die Verschiedenheit der \emph{ästhetischen} und der
\emph{logischen Vollkommenheit} des Erkenntnisses \ldots{} Die logische
Vollkommenheit des Erkenntnisses beruht auf seiner Uebereinstimmung mit
dem Objecte; also auf \emph{allgemeingültigen} Gesetzen, und läßt sich
mithin auch nach Normen \textit{a priori} beurtheilen. Die ästhetische
Vollkommennheit besteht in der Uebereinstimmung des Erkenntnisses mit
dem Subjecte, und gründet sich auf die besondere Sinnlichkeit des
Menschen.“ Kants Logik von G. B. Jäsche. 1800. Werke von
G. Hartenstein. Leipzig 1838 1ster Bd. S. 360--61. --- Men vil man paa
denne Maade fastholde en intuitiv Viden ligeoverfor en discursiv, saa
maa man idetmindste tilstaae, at den intuitive Viden er uden al Indsigt
og Begreb. „Es finden daher bei der ästhetischen Vollkommenheit keine
objectiv- und allgemeingültigen Gesetze Statt, in Beziehung auf welche
sie sich \textit{a priori} auf eine für alle denkende Wesen überhaupt
allgemeingeltende Weise beurtheilen ließe.“ Selv naar der, saaledes som
i Kritik der Urtheilskraft, gjøres Rede for, hvad Skjønhed er, bliver
Erkjendelsen discursiv.}
% „Vollkommennheit“ is set with a doubled n on p. 247; so read by both
% tesseract witnesses and confirmed at 600 dpi. Transcribed as printed.

% --- p. 247 ---
\setcounter{footnote}{0}%
3) „Er vor Viden i sin Totaludvikling betinget af lutter partielle
Udviklinger og altsaa væsentlig discursiv, da maa det Samme gjælde om
al mulig Viden. Vi indrømme, at der paa andre Kloder i Himmelrummet
kunne gives Fornuftvæsener, hvis Viden er langt fuldkomnere end den
% „hvis“ is badly under-inked: only the h and the final s print clearly,
% and no witness reads it. Four letters, and the relative clause requires
% it; rejected alternative: „hvis“ as two words.
menneskelige; men Fuldkommenheden beroer da paa en Gradforskjel, ikke
paa en Sphæreforskjel. Er vor Viden væsentlig discursiv, saa hører det
Discursive oprindelig og nødvendig med til Videns Væsen, og saa maa de
høiere, høieste og allerhøieste Fornuftsvæseners høiere, høieste og
allerhøieste Viden nødvendig være discursiv. Hvad dette vil sige, kan
let oplyses
% --- p. 248 ---
\setcounter{footnote}{0}%
ved et Par Exempler. Til vor Videns mange Ufuldkommenheder hører
blandt Andet den, at vi i det Hele taget lære saa langsomt, og let
glemme det Ene over det Andet. Hvorlænge ere vi ikke om at gjennemgaae
et mere omfattende, være sig mathematisk eller logisk, Cursus; hvilke
Anstrengelser koster det os ikke, naar vi skulle arbeide os igjennem
store Systemer, og forene Indsigt i det Enkelte med Oversigt over det
Hele! Da nu Erfaringen ikke desto mindre lærer, at der indenfor
Kredsen af de menneskelige Individer kan være en stor og mærkværdig
Forskjel i Henseende til intellectuel Begavelse, saa er her Intet til
Hinder for den Antagelse, at der paa andre Kloder maa kunne gives
Fornuftvæsener med langt større og fortrinligere Aandsevner end
Menneskenes. Disse Væsener kunde da være begavede med et Syn, der i
Styrke, Skarphed og Fiinhed overgik det menneskelige; men de kunne
umulig være begavede med et Syn, der gjorde det muligt for dem paa een
Gang at see i alle Retninger og see alle Ting i eet og samme Nu lige
tydelig. Lad dem, hvad Aandsevnerne angaaer, være saa rigt udstyrede,
at de paa en Dag, ja i en Time, kunne sætte sig ind i Systemer, hvis
Tilegnelse vilde koste det største menneskelige Genie en halv Levetid,
saa vilde de dog ikke være istand til at tænke to Tanker, endsige da to
Læresætninger eller to Beviser, lige tydeligt paa een Gang. Lad dem
være nok saa store Physikere, lad deres Sandsen, Forestillen og Tænken
i den Grad følges ad, at de ere nærved at tælle Luftsvingningerne
ligesaa let som de høre Tonerne, og Æthersvingningerne ligesaa hurtig
som de see Farverne: en simultan Høren, en simultan Seen, en simultan
Tællen vil dog være dem umulig. At see hele Fladen af en Tavle med eet
Blik er en Umulighed; dog foregaaer Successionen, endog efter vore
Synsvilkaar, saa hurtig og let, at vi, naar Fladen ikke er for stor,
have en Fornemmelse, som om vi saae det Hele paa een Gang. Er Fladen
f. Ex. rød, saa vil det Eensartede endog bidrage til at forhøie Ind\-%
% --- p. 249 ---
\setcounter{footnote}{0}%
trykket af det Simultane. Men forestille vi os nu, at hver af de røde
Straaler tillige skulde opfattes for sig, saa begynder Svimmelheden;
fordre vi endvidere, at Blikket skal følge, og Tanken tælle
Svingningerne i hver Straale (altsaa, for hver Straale 477 Billioner
Svingninger i Sekundet), saa sortner det for vore Øine; føie vi saa
endelig den Fordring til, at den, der seer Farvefladen og tæller
Æthersvingningerne, med det Samme skal høre Ouverturen til Figaros
Bryllup og mærke sig hver Tone og tælle Luftsvingningerne, saa have vi
det Meningsløse.“

„Hvad heraf kan udledes, er noget meget Vigtigt; thi heraf kan udledes,
at ikke blot den menneskelige Viden, men enhver iøvrigt nok saa
fuldkommen Viden er og maa være indskrænket og ifølge sin Indskrænkning
ufuldkommen. En fuldkommen Viden, en Viden uden Eensidighed, uden
Skranker og Indskrænkninger er en Trekant uden Grændser, en Cirkel med
uendelig Radius.“

Den saaledes motiverede Fornegtelse af Vidensidealet egner sig ret til
at gjøre os opmærksomme paa det eiendommelige Vexelforhold, hvori
Erkjendelseslæren og Logiken nødvendig maae staae til hinanden.
Abstrahere vi fra Erkjendelseslæren, vil det være umuligt at udvikle
Subjectivitetens Logik, eftersom vi da maae savne de Udgangspunkter,
der alene bestemmes ved den subjective Videns psychologiske
Betingelser; holde vi os udelukkende til Erkjendelseslæren, altsaa til
den subjective Videns psychologiske Forudsætninger, komme vi consegvent
til en Fornegtelse af Videns og dermed af Sandhedens absolute
Gyldighed. Hvilken Forskjel det gjør, om Erkjendelsesproblemerne
behandles som i en Indledning til Ideernes Logik, eller opstilles paa
Grundlaget af Subjectivitetens ontologiske Selvudvikling, sees let ved
et Tilbageblik.

Den under første Synspunkt paapegede Eiendommelighed ved vore indbyrdes
ueensartede Sandser og deres Forhold til Tænkningen vilde, dersom vi
udelukkende holdt os til Er\-%
% --- p. 250 ---
\setcounter{footnote}{0}%
kjendelseslæren, ganske vist være uforenelig med Begrebet af en
fuldkommen, Alt omfattende og gjennemtrængende Viden. Men erindre vi,
hvad der er udviklet paa de allerede tilbagelagte Stadier af den
\emph{umiddelbare}, den \emph{reflecterende} og den \emph{begribende}
Subjectivitet, faaer Sagen et andet Udseende. Medens \emph{den
umiddelbare Subjectivitet} kun havde Blik for Sandsernes psychologiske
Eiendommelighed, indsaae den \emph{reflecterende} klart, at denne
psychologiske Eiendommelighed i alle Punkter maatte have sin
tilsvarende ontologiske Bestemthed. Det er nu denne Bestemthed, som for
\emph{den begribende Subjectivitet} udtrykkelig articuleres ved
Vexelbestemmelsen af det Logiske og det Antilogiske. Det er ikke
Sandselighedens psychologiske Kjendsgjerninger, men Sandsningens og
Sandsernes Ontologie, hvorpaa Eftertrykket nu maa falde.

Hvad her under det andet Synspunkt er blevet fremhævet for at vise, at
al vor Viden er og nødvendig maa være discursiv, vedbliver ligeledes at
have uimodsigelig Gyldighed, saalænge vi holde os til den endelige
Adskillelse af det Subjective og det Objective. Men allerede den Maade,
hvorpaa \emph{Forstandsreflexionen} efterhaanden gik over i
\emph{Væsens\-%
reflexion}, gjorde det jo indlysende, hvorledes det Endelige i
Forholdet Punkt for Punkt opløste sig i eiendommelig bestemt
Uendelighed; det uendelig bestemte Vexelforhold mellem Subjectiviteten
og det Objective blev derefter paa den begribende Subjectivitets
Standpunkt tydeliggjort ved en Paaviisning af Forskjellen mellem den
psychologiske og den ontologiske Objectivering, en Forskjel, som, naar
der spørges om Forholdet mellem den intuitive og discursive Viden, er
af afgjørende Betydning. Erkjendelseslæren kjender kun den
psychologiske Objectivering; den objective Logik „lader Subjectet selv
bestemme sig til Object“, men synes forøvrigt ikke at have nogen Anelse
om, at der i ontologisk, d. v. s. i subjectiv-objectiv logisk Forstand
gives et Objectiveringsproblem. Saa\-%
% --- p. 251 ---
\setcounter{footnote}{0}%
længe Objectiveringsproblemet endnu ikke er blevet til, vil der ganske
vist udfordres deels en mystisk Contemplation deels et dogmatisk
Testimonium, for at hæve Bevidstheden fra den discursive Erkjendelse op
til den intuitive Videns, den salige Skuens guddommelige Region.

Hvad endelig det tredie Punkt angaaer, da vil det fremdeles være saavel
Erkjendelseslæren som Aandsphænomenologien en Umulighed at bestemme
Forholdet imellem den ufuldkomne og den fuldkomne Viden. For at
behandle Spørgsmaalet om en absolut fuldkommen Videns Mulighed, maa
Subjectiviteten under en fremadskridende logisk Selvudvikling ikke blot
have lært, hvilke Forvandlinger den endelige Viden med dens
forskjelligartede Erkjendelser kan og maa undergaae; den maa ved i alle
Videns Opgaver at see paa Vidensprincipet tillige have vundet en saadan
Indsigt i Forholdet mellem det Oprindelige og det Afledede, at den i
sin afledede Viden overalt kan see en Afglands af den oprindelige.

Enten Bevidstheden, som paa det nylig omhandlede Standpunkt,
udelukkende holder sig til det Endelige, og seer i Idealet en uklar
Phantasie, en Modsigelse; eller den, som paa det theologiske
Standpunkt, gaaer ud fra en umiddelbar og følgelig dualistisk
Modsætning af det Menneskelige og det Guddommelige, og saa kommer deels
til en abstract og intetsigende, deels til en tankeløs og
selvmodsigende Opfattelse af Alvidenheden, saa bliver Vidensidealet
krænket. Men Idealet bliver da ogsaa krænket derved, at det Absolutes
Philosophie giver sig Mine af at have naaet en absolut Viden af det
Absolute.\footnote{% printed mark *)
„In dem Wissen hat der Geist die Bewegung seines Gestaltens
beschlossen, insofern dasselbe mit dem unüberwundenen Unterschiede des
Bewußtseyns behaftet ist. Er hat das reine Element seines Daseyns, den
Begriff, gewonnen \ldots{} Indem also der Geist den Begriff gewonnen,
entfaltet er das Daseyn und Bewegung in diesem
% the printed page break p. 251 | p. 252 falls here, inside the note
Aether seines Lebens und ist \emph{Wissenschaft}\ldots{} Wenn in der
Phänomenologie des Geistes jedes Moment der Unterschied des Wissens und
der Wahrheit und die Bewegung ist, in welcher er sich aufhebt, so
enthält dagegen die Wissenschaft diesen Unterschied und dessen Aufheben
nicht, sondern indem das Moment die Form des Begriffs hat, vereinigt es
die gegenständliche Form der Wahrheit und des wissenden Selbsts in
unmittelbarer Einheit. Das Moment tritt nicht als diese Bewegung auf,
aus dem Bewußtseyn oder der Vorstellung in das Selbstbewußtseyn und
umgekehrt herüber und hinüber zu gehen, sondern seine reine, von seiner
Erscheinung im Bewußtseyn befreite Gestalt, der reine Begriff und
dessen Fortbewegung hängt allein an seiner reinen \emph{Bestimmt\-%
heit}.“ Phänomenologie des Geistes. Hegel's Werke. 2te Bd. Berlin 1832
S. 609.}

% --- p. 252 ---
\setcounter{footnote}{0}%
At Vidensidealet ligesaavel krænkes ved den sidste af disse
Synsmaader, som ved den første, er indlysende. Hvormeget vi end
anstrenge os for en intuitiv Viden af det Absolute, saa gjør det
Discursive, det Betingede og Relative sig dog gjældende overalt ikke
blot i Videnskaberne, men i Videnskaben. Skal den Viden, Philosophien
paa et givet Udviklingstrin har naaet, udgives for en absolut Viden i
den Forstand, at Idealet selv er naaet, da udspringe heraf to for
Videnskaben skadelige Følger. For det Første vil Misforholdet imellem
den virkelige Præstation og Videns ideale Fordringer, selv om man
tænkte sig en Philosophie saa fuldendt, som den, det største
philosophiske Genie henimod Dagenes Ende engang maatte kunne udvikle,
dog være saa iøinefaldende, at den maa friste Skarpsindigheden til at
anklage Philosophien for utilbørlig Mangel af Selvkritik. Er der nogen
Videnskab, der strengt er opfordret til at mærke sig Grændserne for den
menneskelige Viden, da er det nok Philosophien.

Dernæst, og denne Mislighed er ingenlunde ringere end den foregaaende,
vil den Uklarhed, som altid fremkommer, naar Philosophien forvexler sin
Viden af, hvad
% --- p. 253 ---
\setcounter{footnote}{0}%
der hører til en fuldkommen Viden, med den fuldkomne Viden selv, virke
sløvende og bedøvende paa den philosophiske Forskning. Istedenfor at
tilsløre det i al vor Viden Mangelfulde og Ufuldkomne, er Forskeren
tvertimod opfordret til at drage Problemerne frem. Det forholder sig
ikke med vor Videns Fremskridt saaledes, at de resterende Problemer
skulde blive færre, lettere og ubetydeligere, jo mere vi nærme os
Idealet; tvertimod! jo klarere Idealet sees, desto høiere blive
Problemerne, og desto dybere gaaer Undersøgelsernes Strøm.

\subsubsection*{γ) Vidensidealet.}

Naar den intuitive Viden paa umiddelbar Maade opfattes og bestemmes som
Vidensideal, saa maa den fuldkomne Viden udelukke al Begriben; thi
ingen Begriben er tænkelig \textit{sine abstractione, sine discursu,
sine ratiocinio}: Dogmatikernes Alvidende veed Alt, men begriber Intet.
Skal Viden derimod udvikles til en Begriben, saa bliver den discursiv;
men naar den discursive Viden sætter det Intuitive som et Underordnet,
drager det ned i Endeligheden og benytter det som Element, saa bliver
Begribningen selv saa endelig, middelbar og urolig, at Idealet
forsvinder. Hvad er da her det Almindelige, Viden eller Begriben? Det
er aabenbart Viden; thi der kan gives en Viden uden Begriben, men ikke
omvendt en Begriben uden Viden. Men naar Viden er det Almindelige og
Begribningen det Særskilte, saa maa man nærmere undersøge, hvorledes
dette Særskilte da forholder sig til sit Almindelige. Som det for Viden
Særskilte danner Begribningen en Mellemform mellem Videns to Extremer:
Begyndelsen og Fuldendelsen. Med Hensyn til Begyndelsen staaer
Begribningen høiere end Viden; med Hensyn til Fuldendelsen staaer Viden
igjen høiere end Begribningen. At Viden staaer høiere end Begribningen,
betyder da ingenlunde, at Begribningen er udelukket, men derimod, at
det Middelbare er hævet og forklaret i høiere Umiddelbarhed. „Vilde
imidlertid Nogen
% --- p. 254 ---
\setcounter{footnote}{0}%
sige: vor Erkjendelse er altsaa kun umiddelbar, fordi den er saa
uendelig middelbar, skulde vel en saadan Yttring let blive opfattet som
Beviis paa Spidsfindighed og dialektisk Ordspil. Dog Gyldigheden af den
opstillede Sætning vil netop indlyse umiddelbart, naar Forstaaelsen
først er bleven mæglet. I en Kjæde af sammenhængende Led gjælder det
ganske vist, at medens første og andet Led ere i umiddelbar
Sammenhæng, ere første og tredie Led kun i en middelbar, første og
fjerde i en endnu mere middelbar Forbindelse. Men med en saadan
Udstykning af Umiddelbart og Middelbart er Begrebet „indre
Uendelighed“ uforeneligt. I den intellectuelle Sammenhæng taler man
vel ogsaa om Yderleds Forbindelse ved Mellemled; men jo dybere
Indsigten gaaer, desto mere opløses de endelige Mellemled som
Begrebsbestem\-%
melser i Yderleddenes inderlige Eenhed. For ikke at blendes af
Middelbarheden i de indtrædende Mellemled, maa man see paa
Lysvirkningen og klare sig Forskjellen mellem det Gjennemsigtige og det
Uigjennemsigtige.“\footnote{% printed mark *)
Philosophie og Mathematik. Kjøbenhavn 1857 S. 41.}

For at klare Forskjellen mellem det Gjennemsigtige og det
Uigjennemsigtige, maae vi fremfor Alt lægge Mærke til, at
Gjennemsigtigheden ikke blot retter sig efter Gjenstanden, der skal
sees, men ogsaa efter Synet. Gives der forskjelligt Slags Syn, saa maa
der ogsaa kunne gives forskjelligt Slags Gjennemsigtighed. Saaledes kan
Noget meget vel være gjennemsigtigt for Tanken og den rene Indsigt,
uden at det derfor er gjennemsigtigt for Øiet. Naar Botanikeren
betragter et Træ, da er han vistnok ligesaa lidt istand til enten at
see igjennem Stammen eller at see dets Rod igjennem Jorden, som enhver
anden Betragter; ikke desto mindre vil Træets hele Organisme være
ganske anderledes gjennemsigtig for Botanikeren end for den Ukyndige;
det Samme gjælder, naar Mineralogen betragter et Mineral, o. s. v. Men
slige relative
% --- p. 255 ---
\setcounter{footnote}{0}%
Forskjelligheder ere dog, hvormange Gradationer de end gjennemløbe,
utilstrækkelige til at give os et Begreb om den Viden, der bærer alt
Værende. Grundspørgsmaalet er: hvilke Forvandlinger maa den os
bekjendte Viden undergaae, for at Idealet skal kunne tænkes realiseret,
og paa hvilke Betingelser ere vi istand til at indsee Muligheden af
slige Forvandlinger? Besvarelsen af dette Spørgsmaal vil væsentlig
afhænge af, hvorledes vi nærmere opfatte og bestemme, hvad vi ovenfor
have kaldt Sandsningens og Sandsernes Ontologie.

Hvilken ontologisk Betydning kan der da tillægges Sandsningen?
Sandsningen er den nødvendige og derfor almeengyldige
Mellembestemmelse mellem Videns rene Idealitet og Gjenstandens
umiddelbare Realitet. At den umiddelbare Existens er begrundet i
Reflexion, og Reflexionen i den reflecterende Subjectivitet, er i det
Foregaaende blevet belyst fra saamange Sider, at vi i denne
Sammenhæng kunne fastholde det vundne Resultat som en uomtvistelig
Kjendsgjerning. Den subjective Reflecteren og Existensens objective
Umiddelbarhed forholde sig som det Logiske til det Antilogiske, og
afgive forsaavidt ueensartede Extremer. Det Mellemled, der skal forene
disse Extremer ved at udtrykke de ueensartede Factorers naturlige,
d. e. umiddelbare Forskjelseenhed, maa altsaa fordoble sig saaledes, at
det umiddelbart og paa een Gang udtrykker det Antilogiske i logisk og
det Logiske i antilogisk Form; her have vi Sandsningen og det
Sandselige. Hvilken af vore Sandser vi end vælge som Exempel,
overbevises vi om, at Sagen forholder sig saaledes. I existentiel
Umiddelbarhed staaer Synets Gjenstand udenfor Synet; kun Gjenstandens
Billede er i Synet. Her er i Synet som i enhver Sandsning et
eiendommeligt Forhold imellem Væsen og Virkelighed. Væsentligt seer
den Seende kun Billedet, ikke Gjenstanden; men i Virkeligheden seer han
Gjenstanden, ikke Billedet: Gjenstanden er synlig, og dog er dette
Synlige ikke Gjenstanden. For
% --- p. 256 ---
\setcounter{footnote}{0}%
at see, hvorledes denne Modsigelse løses, maa man see, hvorledes den
fremkommer. Det er en \textbf{ureflecteret} Reflexion, en umiddelbar
Reflecteren, hvorpaa al Sandsning beroer. Enhver Sandsning bestemmes
umiddelbart, forsaavidt den bestemmes ved Indtrykket;
Sandseindtrykkets Realitet er i det Antilogiske: i denne Henseende
gjælder det, at det, der sees, er Gjenstanden selv, ikke dens Billede.
Men da Sandsningen foregaaer i Subjectiviteten og beroer paa en
umiddelbar, og derfor ubevidst Reflexion, saa kan Sandsen ved
Indtrykket kun bestemmes til selv at bestemme; Sandsningens Idealitet
er i det Logiske: i denne Forstand gjælder det, at det, der sees, ikke
er Gjenstanden selv, men dens Billede. Paa den umiddelbare Omsætten af
det Antilogiske i det Logiske, og omvendt, beroer da saavel
Modsigelsen som dens Løsning. Af det i denne Omsætten Eiendommelige
forklares det tillige, at Sandserne kunne være saa indbyrdes
ueensartede: det Ueensartede hidrører nemlig fra den umiddelbare
Bestemthed, der er uadskillelig fra det antilogiske Element.

Vi have her et Udgangspunkt for Bedømmelsen af Sandsningens
Forvandlinger. At denne Verden i sin Mang\-%
% the printed line-break hyphen after „Mang“ is not inked
foldighed af sandselige Ting kun kan have virkelig Gyldighed,
forsaavidt den har Gyldighed i og for Ideen, er indlysende. Men er der
nu et System af sandselige Existensbestemmelser, som i Videns Idee har
Gyldighed for den Vidende, saa følger heraf med uafviselig Conseqvens,
at der i Videns Idee ogsaa maa være et System af eiendommelige
Sandsninger og Sandsningsbestemmelser, der har Gyldighed for den
Vidende; thi Sandsningen og det Sandselige ere Kategorier, der staae og
falde med hinanden.

Man har fra Idealismens logiske Synspunkt ikke taget i Betænkning at
tillægge Ideen Tænken, Begriben og Viden: „Die absolute Idee allein ist
Seyn, \textbf{unvergängliches Leben}, sich wissende Wahrheit und ist
alle Wahrheit\footnote{% printed mark *)
Hegel Die subjektive Logik. Werke 5te Band. Berlin 1834 S. 238.}“;
fremdeles „Die
% --- p. 257 ---
\setcounter{footnote}{0}%
Methode ist daraus als der \emph{sich selbst wissende}, sich als das
Absolute, sowohl Subjektive als Objektive, zum \emph{Gegen\-%
stande habende Begriff}, somit als das reine Entsprechen des Begriffs
und seiner Realität, als eine Existenz, die er selbst ist,
hervorgegangen.“ Men gjør man Alvor af at lægge en absolut Viden ind i
Ideen som Idee og lade denne Viden have sin Gjenstand i den
existentielle Virkelighed, da vil det ved nærmere Overveielse vise sig,
at Viden i Form af Tænken og Begriben umulig kan opfatte den virkelige
Existens, medmindre Sandsningen danner en mæglende Overgang fra
Gjenstandens antilogiske Realitet til den rene Videns logiske
Væsenhed. Tør man ikke indrømme Sandsningen ontologisk Gyldighed, fordi
den ifølge vor egen umiddelbare Erfaring er bunden til psychologiske
Betingelser, da er det kun Phraser og opskruede Talemaader, naar man
paa Grund af Tænkningens ontologiske Gyldighed vil erkjende „den
vidende Sandhed“, „det sig i Alt vidende Begreb“; thi dersom man
udelukkende holder sig til den umiddelbare Erfaring, da forudsætter
Tænkningen jo ligesaavel sine psychologiske Betingelser, som
Sandsningen.

Hvad er det da for en Forvandling af de psychologiske Betingelser,
Sandsningens Ontologie væsentlig forudsætter? Det er, efter hvad vi i
en tidligere Sammenhæng allerede have paaviist, den endelige
Adskillelse af Viden og Magt, paa hvis Ophævelse Eftertrykket her maa
falde. Paa det psychologiske Standpunkt ere de forskjellige
Sandsninger væsentlig betingede af forskjellige „Affectioner“.
% a short diagonal stroke is inked in the word space between „ere“ and
% „de“; a scanner or press artefact, not a character
Disse Affectioner vise allesammen hen til en Dualisme af Viden og Magt.
Magten er i den os afficerende Gjenstand; Viden bestemmes ved den
afficerede Sands. I denne Henseende maa det indrømmes, at Sandsningen
og Sandseligheden hører til det Lavere i vor Natur, Tænkningen med sine
„Functioner“ og Begreber til det Høiere.

Men paa det ontologiske Standpunkt er denne Dualisme
% --- p. 258 ---
\setcounter{footnote}{0}%
hævet. For den ontologiske Subjectivitet er Viden Magt og Magten
Viden; men Sandsningen er herved ingenlunde ophævet; tvertimod! idet
Eenheden af Magt og Viden tillige udtrykker en Modsætning, betegner
% a small raised mark is inked in the word space between „herved“ and
% „ingenlunde“; it is not the printed double hyphen and is read as
% spacing material that took ink
Sandsningen det i denne Modsætningseenhed Umiddelbare, seet fra Videns
Synspunkt. Hvilke Mellembestemmelser den ontologiske Sandsning her
afgiver, kan let paavises. For sig udtrykker Magtbestemmelsen
Existensens antilogiske Realitet; saaledes er det enkelte Atom, for sig
betragtet, et Existenspunkt; dette Existenspunkt er endvidere et
materielt, et sandseligt Punkt, altsaa et Sandsningspunkt. Dersom det
sandselige Existenspunkt ikke tillige var et Sandsningspunkt, saa var
her ingen Eenhed af Magt og Viden; men i Conseqvens af Magtens og
Videns Eenhed gjælder det: saamange sandselige Existenspunkter,
saamange ontologiske Sandsningspunkter. Existenspunkterne ere
fremdeles i forskjellige Relationer til hverandre indbyrdes (Tyngde,
Cohæsion o. s. v.); men det har i det Foregaaende fra mange Sider viist
sig, at alle objective Relationer beroe paa subjective Reflexioner.
Relationernes og Reflexionernes umiddelbare Eenhed bestemmes da i denne
Sammenhæng fra Magtens Side ved et System af Sandsningspunkter, fra
Videns Side ved et System af Sandsninger; til hver sandselig Relation
(et Tryk, en Faldbevægelse) svarer en ontologisk Sandsning. At det ikke
er sandselige Forholdsled, der i deres indbyrdes Relationer tør antages
at sandse hinanden eller at have Sands for hinanden, er forstaaeligt
af, hvad der under „den reflecterende Subjectivitet“ er afhandlet
angaaende Vexelbestemmelsen af Relation og Reflexion. Det er „de
oversvævende Reflexioner“, som her, ifølge den paa Begrebets Standpunkt
optagne antilogiske Factor, have faaet en saa eiendommelig Tilvært af
existentiel Umiddelbarhed, at de netop fra denne Umiddelbarheds Side
blive at opfatte som ontologiske Sensationer.

Vi kunne oplyse Sagen endnu fra en anden Side. Blandt de Philosopher,
der have været allernærmest ved at
% --- p. 259 ---
\setcounter{footnote}{0}%
opdage Sandsningens og Sandsernes Ontologie, maae vi særlig nævne
Leibnitz. Som bekjendt, lader Leibnitz Materien bestaae, objectivt
seet, af „nøgne Monader“, subjectivt seet, af dunkle, forvirrede
Forestillinger. Det er aabenbart, at Leibnitz i de dunkle
Perceptioner (\textit{perceptions petites}) er kommen til at tangere
den ontologiske Eenhed af Sandsningen og det Sandselige; denne
Tangeren er her saameget mere at fremhæve, som den er punktuel og
ved Punktualiteten skikket til at afgive Forudsætninger for
speciellere Analyser; men to i Monadelæren begrundede
Misforstaaelser følge med og gjøre Antydningen ubrugbar. Den
første Misforstaaelse er den, at Materien som Produkt af de dunkle,
forvirrede Forestillinger skulde være et Slags Resultat af
Idealitetens og Videns Afmagt, at Materien, som Følge heraf, ikke
skulde have Gyldighed for den guddommelige Viden; det er den gamle
Misforstaaelse, der hidtil har gjentaget sig paa alle Spiritualismens
og Idealismens Trin. Den anden Misforstaaelse er den, at de til
Sensationerne svarende dunkle Forestillinger skulde være som
Kraftyttringer af Existenspunkterne selv. En saadan Synsmaade maa
føre til de største Urimeligheder. Skulde Existenspunkterne være
sig selv forestillende Forestillinger, saa maatte jo heraf følge, at
hver Gjenstand bestod af Forestillinger, og at Alt forestillede; men
denne Sætning er ligesaa urimelig som den, „at hver Gjenstand
bestaaer af Tanker, og at Alt tænker“. Forskjellen mellem
Kraftyttringer og ontologiske Sandsninger er i Umiddelbarheden, hvad
Forskjellen mellem Relationer og Reflexioner er i Reflexionen.

Existenspunktet, det ved Magten bestemte Sandsningspunkt, er som
Videns Gjenstand en Skranke for Viden; al Viden, ogsaa den i Ideen
guddommelige, har en Skranke, forsaavidt den har en Gjenstand. Der
maa i al Viden være
% a small high ink speck stands in the right margin after "være" at the
% end of this line; it is scanner/ink dirt (the page carries no
% footnote), and is not transcribed.
Noget, som vides; er det, som vides, kun Viden selv, saa er Viden
tom, og den tomme Viden i Begreb med at forvandle
% signature mark at the foot of p. 259: "17*" (not transcribed)
% --- p. 260 ---
\setcounter{footnote}{0}%
sig til Ikke-Viden. For at Viden kan have Realitet, maa det, der
vides, have Realitet; men naar det, der vides, skal
% "skal": the l is only partly inked, its ascender alone showing
have Realitet, maa her paa een Gang være et Videns \emph{Ind\-%
hold} og en Videns \emph{Gjenstand}. En nærmere Undersøgelse heraf
vil lede os til en Indsigt i den ontologiske Vexelbestemmelse af
Sandsning og Reflexion og derved til en klarere Forstaaelse af
Forholdet mellem Begriben og Viden.

Paa Sandsningens nødvendige Eenhed med Reflexionen beroer den
ontologiske Objectiveren. Denne Objectiveren har, som i en tidligere
Sammenhæng er antydet, to Hovedsider, den idealistiske, svarende til
\emph{Idealsystemet}, den realistiske, svarende til
\emph{Realsystemet}. Fra den idealistiske Side er Objectiveringen
logisk og fri for alle Skranker; hver Skranke er i denne Sphære
forvandlet til eiendommeligt Bestemmelsesmoment i Videns uendelig
gjennemsigtige Indhold; her er i denne Kreds ingen Gjenstande: de
adæqvate Ideer træde overalt og fra alle Sider i Gjenstandenes Sted.
I denne Sphære er her da heller ingen Sandsninger; thi de særskilte,
forskjelligartede og eiendommelige Sandsninger forsvinde overalt og
fra alle Sider som Lysninger, Ændringer og Afskygninger, hvorved
Indholdets Idealitet articuleres i indre Uendelighed. Anderledes,
naar vi see paa Realsystemet. Fra den realistiske Side er
Objectiveringen antilogisk og ufri, bunden og bestemt ved utallige
Skranker. I denne Sphære har Viden ikke et immanent Indhold, et
System af klare Ideer, men et uendeligt Indbegreb af reale
Modsætninger, en Verden af Gjenstande, altsaa en Verden af Skranker;
istedenfor de rene Reflexioner faae vi her overalt og fra alle Sider
lutter Sandseyttringer, Fornuftfornemmelser, ontologiske Sensationer.

De ontologiske Sensationer forholde sig til de rene Reflexioner som
Gjenstandsmomenterne til de væsentlige Indholdsmomenter. Ved hvert
Indholdsmoment faaer Viden en ideel Bestemthed, ved hvert
Gjenstandsmoment en reel Skranke. Ligesom Indholdet bestaaer i et
System af
% --- p. 261 ---
\setcounter{footnote}{0}%
Indholdsmomenter, saaledes bestaaer Gjenstanden i et System af
Gjenstandsmomenter. Som det paagjældende Moment er, saaledes er den
tilsvarende Totalitet: just fordi Indholdsmomentet er i sin ideelle
Bestemthed gjennemsigtigt, er Indholdet selv gjennemsigtigt; omvendt!
just fordi Gjenstandsmomentet er og, ifølge sin Charakteer, maa være
uigjennemsigtigt, er Gjenstanden selv som Gjenstand uigjennemsigtig.
At anprise et Vidensideal, der i sin Art skulde være saa ophøiet, at
den Vidende umiddelbart gjennemskuede enhver Videns Gjenstand, er at
anprise sin egen Tankeløshed.

„Men“ --- vil man maaskee indvende --- „naar de virkelige Gjenstande
i deres Virkelighed ikke kunne Andet end afgive Skranker endog for
den guddommelige Viden, saa bliver her jo dog en Mangfoldighed af
Gjenstandsbestemmelser, altsaa en Mangfoldighed af
Værensbestemmelser, der maae falde udenfor Viden, saa er Idealet jo
dog krænket, den ontologiske Subjectivitet
% two faint specks stand between "Subjectivitet" and "fornegtet"; they
% are not the two aligned dots of this fount's colon (cf. "Totalitet:"
% five lines above) and are read as dirt, so nothing is set here.
fornegtet, og Vidensprincipet opgivet.“ Denne Indvending beroer paa
en Misforstaaelse, der uden Vanskelighed lader sig hæve. Vistnok maa
enhver Gjenstand, ja ethvert Gjenstandsmoment sætte Viden en Skranke;
saaledes maa ethvert materielt Atom sætte endog Alvidenheden selv en
Skranke; men man maa ikke forglemme, at den Indskrænkning, Viden ved
ethvert Gjenstandsmoment saaledes erholder, er en Negation, der hører
med i den reale Liden
% PRINTER'S ERROR, transcribed as printed: "Liden" for "Viden". The
% capital is a Fraktur L, not V: calibrated at 600 dpi against "Viden"
% eight lines above on this page and against the L of "Løsning"
% (p. 262) and "Lyset" (p. 263). All three OCR witnesses read L.
selv, og just derved tilfredsstiller en af Videns væsentlige
Fordringer. Skranken er en Magtbestemmelse; men den Magt, som
saaledes sætter Viden Skranker, er i Modsætningseenhed med Viden et
kategorisk Udtryk for Videns egen Realitet. Hver
Gjenstandsbestemmelse, hvorved Viden indskrænkes, er identisk med en
tilsvarende Indholdsbestemmelse, hvorved Viden udvides: at Videns
Indskrænkning i alle Punkter betinger dens Udvidelse, er et talende
Vidnesbyrd om den reale Videns uendelig overgribende Idealitet.

Forsaavidt Gjenstandsmomentet i sin Umiddelbarhed er
% --- p. 262 ---
\setcounter{footnote}{0}%
en Skranke for Viden, vil Materien med dens Atomer og Masser vistnok
være uigjennemtrængelig for den Alt i Ideen Vidende; men forsaavidt
ethvert Gjenstandsmoment er et Magtmoment og som Magtmoment en
Realnegation, en Væsensbestemmelse i og ved den guddommelige Viden
selv, kan den mørke, uigjennemtrængelige og uigjennemsigtige Materie
dog i Ideen gjennemskues af den mægtige Viden, erkjendes og begribes
af den vidende Magt.

Den Modsigelse, at det Uigjennemsigtige dog er gjennemsigtigt, finder
paa alle Punkter sin Løsning i den teleologiske Vexelbestemmelse af
Ideal- og Realsystemet. Af denne Vexelbestemmelse følger nemlig, og
det, vel at mærke, ikke blot i abstract Almindelighed, men tillige i
alt det Specielle, det Enkelte og Punktlige, at hvad der i
Existensens antilogiske Form er Gjenstandsmoment og Skranke, er i
Videns ideelle Form Indholdsmoment og afklarende Bestemthed.

Modsætningen imellem det gjennemsigtige Vidensindhold og den for
Viden uigjennemtrængelige Gjenstand erholder endnu en ny og rigere
Belysning, naar vi see paa de Afskygninger, Ændringer og Overgange,
som de eiendommelige Vexelbestemmelser af Reflexioner og Sensationer
afgive til Udligning af enhver mulig Spænding mellem Videns og
Magtens Fordringer. Ere Videnskategorierne bundne og forsaavidt
skjulte i Magtbestemmelserne, saa ere Magtkategorierne derimod
frigjorte og i Idealiteten klarede som gjennemsigtige
Vidensbestemmelser; men imellem Extremerne af fuldkommen ideel
Klarhed og fuldstændig materiel Dunkelhed falder nu det System af
Mellembestemmelser, der netop articuleres ved Reflexionernes og
Sensationernes Overgange i hinanden.

Jo sandseligere Relationen er, desto mere umiddelbart falder
Fornuftfornemmelsen, den ontologiske Sandsning, sammen med
Kraftyttringen; jo mere elementær Kraftyttringen er (f. Ex. de faste
Legemers Cohæsion), desto dunklere er
% --- p. 263 ---
\setcounter{footnote}{0}%
den tilsvarende Sensation: Fornemmelserne kunne, hvad et saadant
System af materielle Relationer og Formbevægelser angaaer, tabe sig i
Differentialvirkninger af forskjellig Orden. Seet fra Realsystemets
Standpunkt kan Lyset i den guddommelige Viden altsaa gjennemløbe
utallige Gradationer ligefra det fuldstændige ved Gjenstandens
Uigjennemtrængelighed bestemte Mørke gjennem Sensationernes Dæmring
og Sandsereflexionernes Dagskjær til Idealitetens høieste Klarhed.

Herved naaes nu Eenheden af Begriben og Viden. Realsystemet godtgjør
i denne Forstand den \emph{discursive Videns} væsentlige Betydning
for Vidensidealet selv. Fra Idealsystemets Standpunkt derimod have
de angivne Modsætninger af Lys og Skygge kun den Virkning, at de i
det Uendelige tilføre Ideeriget nye Vidensbestemmelser, saadanne
nemlig, som overalt og fra alle Sider udvikle sig ved de existentielle
Dannelser. Her er altsaa i Magtens Kreds, i Tingenes Verden, ikke en
Gjenstand, ja ikke et Gjenstandsmoment, som jo i Idealsystemet har
sin tilsvarende Vidensbestemmelse i fuldkommen Klarhed. Da nu
Idealsystemet med sit hele Indhold af Kategorier, Begreber,
Ideemomenter og Ideer er den vidende Subjectivitet fuldkommen
gjennemsigtigt, saa sees heraf, at Vidensidealet ogsaa i sig antager
Charakteer af \emph{intuitiv Viden}.

I den teleologiske Vexelbestemmelse af intuitiv og discursiv Viden
har Vidensidealet uendelig Realitet, og med en logisk Paaviisning af
denne Realitet maa Subjectivitetens Logik slutte.
% p. 263 closes division C and with it "Første Deel: Videns Idee".
% Beneath the last line stands only a short centred rule (not
% transcribed); no part of "II. Psychologie" appears on this page.
%% ============================================================
%% II. PSYCHOLOGIE  (pp. 264--276)
%% Sub-heads printed within: Indledende Sætninger (264), Plantelivet
%% (266), Dyrelivet (268), Menneskelivet: Psychologien i Grundtræk (273).
%% ============================================================

\part*{II.\\ Psychologie.}
\addcontentsline{toc}{part}{II. Psychologie}
\markboth{Psychologie}{Psychologie}

% --- p. 264 ---
\setcounter{footnote}{0}%
% The part head "II." / "Psychologie." stands above this on p. 264 and is set by the
% skeleton; only the division head below is set here.
\chapter*{Indledende Sætninger.}
% Head centred and letterspaced; period as printed. Agrees with the Indhold (p. 6).

Natur og Aand komme overeens deri, at de begge have Ideen til
absolut Indhold: i Nodvendighedens og den umiddelbare Virkeligheds
% "Nodvendighedens" as printed: plain o, verified at 600 dpi against the plainly
% slashed ø of "Løsning" (p. 266) and "nødvendigt" (p. 265). p. 264 has no ø at all.
Form er Ideen Natur; i Frihedens, Inderlighedens og Selvbevidsthedens
Form er Ideen Aand.

1) Naturens og Aandens Rige forholde sig, hvad den væsentlige
Forskjellighed angaaer, som Andetheden (det Objective) til Selvheden
(det Subjective), som den sandselige, materielle Tilværelse til den
ideelle, i usandselig Virken reflecterede Væren. Saalidt som
Andetheden, det udvortes Tilværende, kan tænkes uden Selvheden,
ligesaa lidt kan den fodende, formende, frembringende Natur tænkes
% "fodende" as printed: plain o, verified at 600 dpi.
uden Aanden. Den reelle Magt, der bærer Materien og ved Lovenes
System bestemmer Naturkrafterne, er en Aabenbarelse af den ideelle
% "Naturkrafterne" as printed: plain a, calibrated against the æ of "tænkes" and
% "mægtige" on this page; the page also prints "Kræfter" with æ. Not normalised.
Magt, der er Eet med Aandens evige Viden og selvbestemmende Virken.

2) Eenheden af Natur og Aand er grundet paa Modsætning,
ligesom Modsætningen paa Eenhed; Grundeenheden udtræder i virksom
Fordobling af en sandselig og oversandselig Verden. Som Aandens
evige Indhold har Ideeriget væsentlig Bestaaen i mægtige Syner
(Urbilleder) og skabende Tanker (Grundbegreber, Ideebegreber,
særskilte Ideer); men i Naturriget er Idee-Indholdet udstykket til
materielle Totaliteter, Elementer og Kræfter. Alle relativt
fuldstændige Naturheelheder (Solsystemer, Jordkloder, levende
Organismer) ere Afbilleder af aandige, men naturkraftige Forbilleder,
Virkeliggjorelser af Aandens virksomme, Naturen iboende, Eenheden af
% "Virkeliggjorelser" as printed: plain o, verified at 600 dpi.
den ideelle og den reelle Magt aabenbarende Ideer.

% --- p. 265 ---
\setcounter{footnote}{0}%
Uagtet enhver nok saa relativ Realidee i sit Væsen er Selveenhed,
og ethvert Ideebegreb forsaavidt et Selvhedsbegreb, saa bære dog de
tilsvarende Naturheelheder alle, om end paa de forskjellige Trin i
forskjellig Grad, Andethedens, Udstykningens og Betingethedens Præg.
De naturlige Heelheder vise sig som sammensatte; Delene kunne endog
særlig fremtræde som Brudstykker (Jordlag, Stene) og Masser
(Luft- og Vandmasser); slige Dele, hvis Kategorier ere physiske
Andethedsbegreber, kunne vel ikke umiddelbart udtrykke nogen Idee,
men ere dog formedelst deres Sammenhæng med det paagjældende Hele
middelbart undergivne Ideens Magt. Eenheden i det fuldstændige
Naturhele er paa een Gang \emph{Sammenhængseenhed og Selveenhed.}
% Letterspacing runs across the line break (Sammenhængs-/eenhed) and takes in "og".
Jo mere abstract physisk Heelheden er (Solsystemet), desto mere vil
Sammenhængseenheden skjule over Selveenheden; jo inderligere
gjennemformet Heelheden er, desto mere vil Selveenheden skinne
igjennem (Organismerne).

3) Paa Modsætningen mellem Sammenhængseenhed og Selveenhed
beroer Modsætningen imellem Tingsaarsager (physiske Aarsager) og
Selvaarsager (Principer), og herpaa igjen Modsætningen imellem
physisk og teleologisk Causalitet. Den teleologiske Causalitet gaaer
overalt i Naturriget sammen med den physiske, og denne har til
Gjengjæld sin oprindelige Forudsætning og Fuldendelse i hiin.

4) Modsætningen mellem Aandens og Naturens Rige er ingenlunde
at indskrænke til Modsætningen imellem Guddommens Aand og denne
nærværende, os sandselig bekjendte Natur. Tvertimod! Det bliver til
Forklaring af de virkelige Livsphænomener nødvendigt at forudsætte
en Mellemsphære, hvori de ellers uforligelige Extremer (af det
Ideelle og det Reelle) gjennemtrænge hinanden. Da en saadan
Mellemsphære imidlertid ligger udenfor al umiddelbar Erfaring, kan
% "kan" very lightly inked here; the a is barely printed. Read from the sense.
den i Ideevidenskaben kun benyttes som metaphysisk Grundlag, og faaer
alene sin Betydning fra de Bidrag, den yder til en alsidig Udtydning
af vitterlige, fra Livet selv uadskillelige Functioner.

% --- p. 266 ---
\setcounter{footnote}{0}%
De angivne Forudsætningers videnskabelige Gyldighed bliver nu at
prøve ved deres Anvendelse paa en naturphilosophisk Løsning af
Livsproblemet.

\chapter*{Plantelivet.}
% Head centred, fat-face Fraktur. Agrees with the Indhold (p. 6).

Et Blik paa Planteorganismens charakteristiske Dele (Rod, Stængel,
Blade af forskjellig Orden) siger os, at disse Dannelser i deres
væsentlige Sammenhang og Følgerække maae henføres til en ganske
% "Sammenhang" as printed: plain a, calibrated against the æ of "Sammenhæng" on
% p. 265; both forms stand in this book. Not normalised.
eiendommelig Vexelvirkning af stofformende Aarsager. Livsaarsagen,
det i Formudviklingen første Virksomme, er en Selvaarsag
(Idee-Aarsag, Princip). Men de fra Selvaarsagen udgaaende Virkninger
forudsætte et Materiale, en Flerhed af forskjelligartede Stoffer, som
ved Paavirkningen sættes i Virksomhed; de under Principets
Paavirkninger virksomme Stoffer ere Tingsaarsager: den levende
Organisme betinges af en \emph{dobbelt Vexelvirkning}, Vexelvirkning
mellem Princip og Tingsaarsager og Vexelvirkning mellem
Tingsaarsagerne indbyrdes.

For at Selvaarsagen kan træde i Vexelvirkning med Tingsaarsagerne,
maa der forudsættes et væsentligt Mellemværende. Dette Mellemværende
har en ideel Side, hvorved det bliver Eet med Selvaarsagen, og en
reel, hvorved det bliver Eet med Tingsaarsagerne. Medens Selvaarsagen
udelukkende bestaaer i Formen og Virken, og forsaavidt har en indre,
usandselig, reflecteret Væren, have Tingsaarsagerne umiddelbar
Bestaaen i de tilsvarende Stoffer og vise sig med udvortes,
umiddelbar, materiel Tilværen. Den levende Organisme antager herved
en \emph{dobbelt Skikkelse}, en ydre, beskuelig i \emph{Exemplaret},
og en indre, i Selvheden reflecteret, en principiel Eenhed af Midler
og Formaal, \emph{Ideen}.

Den gjennemgaaende Modsætning af Livet (Ideen) og den levende
Organisme (Exemplaret) er begrundet i en gjennemgaaende Eenhed. Saa
umuligt det er at aflede Livet af en physisk Vexelvirkning mellem de
blotte Tingsaarsager,
% --- p. 267 ---
\setcounter{footnote}{0}%
ligesaa umuligt er det at tillægge Ideen en
umiddelbar, af Vexelvirkning med Tingsaarsagerne uafhængig
Selvstændighed. Kun ved at paavirke sit Andet og bearbeide sit
Materiale, kan den virksomme Selvhed udvikle sit System af
stofformende Functioner og derunder udforme sin egen oprindelige
Væsenhed (indvortes og udvortes Andet; Materiens Formabilitet;
Plantens Deling).

Muligheden af at dele Planten saaledes, at Livet fortsætter og
udvikler sig typisk af den enkelte Deel, kan ikke modbevise Ideens
Selvgyldighed. Tvertimod! Den fra de forskjellige nydannede
Centralpunkter udgaaende Livsvirksomhed gjentager den samme Kreds af
Dannelser i den samme Grundform, og afgiver et Beviis for, at Ideen,
trods de vexlende Betingelser, maa være den samme. Hvorledes
Selveenheden udvikler sig gjennem den punktuelle Forskjellighed, sees
af Cellelivet. Hver Celle er som Selvhedspunkt oprindelig et Levende
for sig, en særlig Organisme; men under den typiske Vexelvirkning af
de enkelte Celler indbyrdes, synker Celleorganismen ned til
Gjennemgangspunkt, organisk Moment. De mangfoldige Selvhedspunkter
reflectere sig som Sidebestemmelser i en videre omfattende,
overgribende Selvhed. Efter den herved angivne Selvhedslov kunne ikke
blot de enkelte Celler nedsættes til Dannelseselementer for Stængel og
Blad, men Stængel og Blad ved Omdannelser igjen afgive Organer af
høiere Orden (Frugtblad, Frugtbund, Frugtknude, o. s. v.).

Ideen, som udadtil arbeider med Stoffet, altsaa under stadig
Indvirkning af et udvortes Andet, udvikler ved denne udadgaaende
Virksomhed tillige en indadgaaende Forskjellighed, et indvortes Andet,
et System af ideelle Selvhedsbestemmelser, og kommer til eiendommelig
Bestaaen i Naturideernes fælles ontologiske Medium. I dette Medium
har nemlig Plantelivets almeenconcrete Idee trinviis articuleret sig
til et systematisk Indbegreb af over- og underordnede (til Rækker,
Classer,
% --- p. 268 ---
\setcounter{footnote}{0}%
Ordener, Familier, Slægter og Arter svarende) Grundformer
(særskilte Ideer), der virkeliggjøres som Individualideer i
tilsvarende Exemplarer: kun som organiske Exemplarer have Ideerne
umiddelbar Existens.

Af Vexelvirkningen mellem den ideelle og materielle Organisme
(Individualideen og Exemplaret), kan Forholdet mellem Evigt og
Timeligt i Plantelivet nærmere forklares.

Nedsænket og fortabt i stofformende Virksomhed har det individuelle
Liv en endelig Begyndelse og maa da ogsaa have en tilsvarende Ende;
med de organiske Betingelser indtræder Virksomheden, og med
Betingelsernes Ophør ophører den igjen. Men da Selvheden, langt fra at
fremkomme enten af Stofferne eller af Intet, begynder som
Væsenseenhed (Princip), kan den ved sit Ophør heller ikke blive enten
til et stofagtigt Andet eller til Intet: Livets saavel deelvise som
totale Ophør er dets Tilbagevenden til sit principielle Anlæg, til
Inderliggjørelse (ontologisk Erindring) i Arten. Imellem Begyndelsen
og Enden ligger Selvudviklingen, hvis ideelle Udbytte optages og
bevares i Artens ontologiske Erindring: Exemplaret har ikke levet
forgjæves. Den metaphysiske Udodelighed, hvortil det individualiserede
% "Udodelighed" as printed: both o plain, calibrated against the slashed ø of
% "Ophør", "tilhøre" and "høiere" on this same page. Not normalised.
Liv indgaaer i Arten, kan nu vistnok igjen tænkes at afgive
Forudsætning for en ny Spiren i en Kreds af stofagtige Betingelser,
der tilhøre en høiere Region; men en Erkjendelse heraf falder udenfor
den menneskelige Videns Grændser: at Livet, der i enhver af dets
Aabenbarelser opfylder Selvhedens Lov, hverken kan blive til Andet
eller til Intet, er den Grundtanke, ved hvis Udvikling og alsidige
Begrundelse Videnskaben bliver staaende.

\chapter*{Dyrelivet.}
% Head centred, fat-face Fraktur, mid-page. Agrees with the Indhold (p. 6).

Medens Plantelivet ikke kan naae ud over den stofformende
Virksomhed, har Dyrelivet i sit Princip derimod Anlæg til en saa
overgribende Selvvirksomhed, at det i Form af Fornemmelse formaaer at
give Selvet en sandselig Til\-%
% --- p. 269 ---
\setcounter{footnote}{0}%
væren og concentrere sig til Sjæl: det
besjælede Individ har et Legeme. Den udtrykkelige Modsætning af Sjæl
og Legeme bestemmes derved, at Dyreorganismen, idet Systemet af de
vegetative Functioner nedsættes til at afgive Grundlag for et i
Sandsning og vilkaarlig Bevægelse sig yttrende Bevidsthedsliv,
fordobler sig i tilsvarende Organsystemer og forstærker
Concentrationen i den Grad, at de forskjelligartede Organer og dermed
den organiske Individualitet faaer en eiendommelig afsluttende
Begrændsning. Modsætningen imellem de vegetative og de animale
Organer i det dyriske Legeme er den organiserende Selveenhed saa
oprindelig underlagt, at den endog kan henføres til Fosterdannelsens
allerforste Stadier. Allerede Kimen i Ægget sondrer sig i to Lag,
% "allerforste" as printed: plain o, verified at 600 dpi; the same page prints
% "øverste", "henføres" and "Tarmrøret" with a plainly slashed ø.
hvoraf det øverste, mere gjennemsigtige, kaldes det \emph{animale}
eller \emph{dyriske} Blad, det nederste, mere melkehvide, kaldes det
\emph{vegetative} Blad eller \emph{Sliimbladet}. Af denne Begyndelse
sees det udtrykkeligt, at det sjælelige og det vegetative Liv i Dyret
er af een Oprindelse, at denne Oprindelse ikke kan forklares af
organiske Stoffers physiske Vexelvirkning, men beroer paa principiel
Eenhed af vegetative og animale Anlæg. Ligesom Organerne for
\emph{Ernæringslivet} (Tarmrøret, Lungerne, Kjertelgangene) udvikle
sig af Sliimbladet, saaledes fremkomme Organerne for det
\emph{sjælelige Liv} (Nervesystemet, Sandse- og
Bevægelsesredskaberne) af det animale Blad. De tvende saaledes
adskilte Organsystemer ere som to Extremer, der trænge til en Midte,
hvori de kunne sammensluttes; der opstaaer derfor imellem dem et
tredie Blad, nemlig Karbladet, hvoraf Hjertet, Blodkarrene og selve
Blodet, altsaa det for de modsatte Organklasser fælles System, bliver
dannet.

Imellem den stofformende Selvvirksomhed, Charakteermærket for det
vegetative Liv, og Fornemmelsen, Dyrelivets væsentlige Særkjende,
synes Afstanden vel ved forste Øiekast at være ligesaastor som
% "forste Øiekast": the adjective has a plain o, the noun a plainly slashed
% capital Ø (calibrated against the plain capital O of "Oprindelse" and
% "Organerne" on this page). Both as printed.
Afstanden imellem den blot stofagtige
% --- p. 270 ---
\setcounter{footnote}{0}%
Virksomhed og Selvvirksomheden; men
nærmere seet betegner Fornemmelsen, det første umiddelbare
Selvbefindende, netop en inderlig bestemt Selvhedsact og afgiver som
saadan en nodvendig Forudsætning for Vilkaarligheden,
% "nodvendig" here as printed: plain o. The same page prints "nødvendig" with a
% plainly slashed ø lower down ("enhver Sandsning nødvendig forudsætter"), and
% "første", "frigjøres", "skjøndt" likewise. Not normalised.
Selvbestemmelsens første udtrykkelige Yttring.

Under Fornemmelser, Sandsninger og Vilkaarlighedsyttringer staaer
Sjælen gjennem alt Andet i Forhold til sig selv og frigjøres til ideel
Selvstændighed ikke blot ligeoverfor Omverdenen, men i inderlig
Forstand endog overfor sit eget Legeme. Da Selvvirksomhed imidlertid
paa ethvert Trin er betinget af Vexelvirken med Andet, vil det heraf
forstaaes, at det dyriske Sjæleliv maa i sine Functioner være ligesaa
afhængigt af Bevidsthedsorganerne (Ganglier og Nerver) som
Ernæringslivet af Ernæringsorganerne. Jo stærkere Anlæg til
selvsættende Virksomhed og derved til forhøiet Sandsebevidsthed det
dyriske Livsprincip indeholder, desto mere vil Nervesystemets
Centralorganer være fremherskende. Medens Bloddyrene kun have et
spredt og uregelmæssigt Nervesystem, og Leddyrene et
Bugknudenervesystem med en Flerhed af relative Midtpunkter (en dobbelt
Nerveknude, i hvert Led, der afgiver Nerver til de omgivende Dele og
staaer i Forbindelse med det foran og bagved liggende Leds
Nerveknuder), have Hvirveldyrene i Rygmarv og Hjerne
(Rygmarvsnervesystemet), og da navnlig i Hjernen, et for
Bevidsthedslivet afgjørende Centralorgan.

Bevidsthedsfunctionerne, skjøndt afhængige af organisk materielle
Betingelser, afgive ved deres Frigjørelse til Inderlighed et særegent
Beviis for Selvhedens principielle Selvstændighed som sjæleligt Selv.
Vel er Selvopfattelsen under den dyriske Sandsning umiddelbart fortabt
i Opfattelsen af Andet; men hvilken af de fem Sandser man end vil
vælge, saa skal en physiologisk Analyse af Functionerne tydelig vise,
at enhver Sandsning nødvendig forudsætter et sandsende Selv. Vistnok
er Vilkaarligheden paa dens laveste Trin endnu for
% --- p. 271 ---
\setcounter{footnote}{0}%
uvilkaarlig til, at
det dyriske Selv kan i en Vælgen gribe sig selv som den Vælgende; men
derfor er den vilkaarlige Bevægelse ligefuldt at ansee som det første
umiskjendelige Tegn paa en activ, om end naturbestemt, Yttring af
Selvet.

Som Erfaringsvidenskab maa Physiologien undersøge og nøiagtig
oplyse os om Forskjellen imellem Planteorganismen og det dyriske
Legeme i alle saavel indvortes som udvortes Enkeltheder; men
Grundforholdet selv, den oprindelige Væsensforskjel, kan kun udvikles
i Psychologien. Efter sit Væsen er Livet under alle organiske
Virksomheder en sig udformende, til eiendommelig Bestaaen
fremarbeidende Selvvirksomhed; et dybere Indblik i denne
Selvvirksomhed vil altsaa beroe paa en methodisk Opklaring af den
Reflecteren og Virken, hvorved Selvheden gjennem sit udvortes Andet
organiserer det Materielle og paa dette Grundlag tillige gjennem sit
indvortes Andet articulerer det Ideelle: Opgaven er saavel fra
Synspunktet af det vegetative som af det animale Liv: at
tydeliggjøre Umiddelbarheden og Reflexionen i \emph{organisk} Eenhed,
i \emph{organisk} Modsætning, samt i \emph{organisk reflecteret}
Eenhed af Eenhed og Modsætning.
% The compositor letterspaces only "organisk" (three times, the second broken
% across the line end) and "organisk reflecteret"; "Eenhed" and "Modsætning" are
% left unspaced each time. Transcribed as printed, not regularised.

\subsubsection*{a) Det dyrisk-vegetative Liv.}
% Centred display head, letterspaced; the label "a)" is not letterspaced.

\emph{Umiddelbarhed og Reflexion i organisk Eenhed.} Sættes
Næringsmiddel mod Ernæring, lyder det Første paa Umiddelbarhed, det
Sidste paa Reflexion, men da Næring kun er i det Nærende, gaae de
umiddelbart sammen i organisk Eenhed. For Reflexionen forudsætter
Ernæringen et mod Organismen udvortes Andet; den organiske
Reflecteren virker som Forvandlen af Organismens Andet til Organismen
selv. Under den reflecterede Virken kommer det Umiddelbare særligt
tilsyne paa ethvert Trin af den trinvise Omdannelse. Næringsstoffet,
efterat det i Mundhulen er bearbeidet mechanisk og i
Fordøielsescanalen chemisk, deler sig som tilberedet Masse (Chymus) i
et modstridende Andet, et Ubrugeligt, der føres bort, og et til
Forvandling
% --- p. 272 ---
\setcounter{footnote}{0}%
tjenligt Andet (Chylus), der opsuges i Karrene
(Melkekar, Lympekar) og gaaer over i Blodet. Bloddannelsen afgiver en
% "Lympekar" as printed, without the h; verified at 600 dpi.
levende Midte imellem to hinanden reflecterende Extremer: Næringen som
Stof og Stoffet som Næring. Blodet forener Andetheden med Selvheden;
fra Blodet udgaaer den sidste Forvandling: hvert Organ ernæres derved,
at det af Blodet ernærer sig selv.

\emph{Umiddelbarhed og Reflexion i organisk Modsætning.} Under
Ernæringsvirksomheden er Reflexionen bunden i det Umiddelbare: den
frie Reflecteren er kun til for en betragtende Subjectivitet. Den
bundne Reflexion udlægger sine Overgange i en Række af organiske
Midler og Mellemled (Fordøielsesorganer, Blodets Kredsløb,
Aandedrættet). I Middelbarhed og Mellemled hersker Andetheden og med
denne igjen det Umiddelbare; ikke desto mindre er her i hver organisk
Andethed et Moment af Selvhed, og i Selvheden sættes Reflexionen: det
Frembragte er selv et Frembringende, det Dannede selv et Dannende,
det Ernærede selv et Errærende. Andetheden i de umiddelbart
% PRINTER'S DEFECT: "Errærende" is set with a second r where the word requires n
% ("Ernærende"); at 600 dpi the third sort shows the free rightward arm of r and
% not the second stem of n, calibrated against "Ernærede" two words earlier and
% "Ernæringsvirksomheden" above. Transcribed as printed.
tilværende Led er Middel, Selvheden, reflecteret i den hele Kreds af
typiske Formvirksomheder, er Maal: sættes Maalet, det sig udformende
Hele, som umiddelbart, saa ere Midlerne reflecterede; sættes Midlerne
umiddelbart, saa er Reflexionen i Maalet; men en organisk reflecteret
Eenhed af Umiddelbarhed og Reflexion kan først komme frem i

\subsubsection*{b) det dyrisk-sjælelige Liv.}
% Centred display head, letterspaced; "det" is lower-case as printed, because the
% sentence above ends without a full stop and runs grammatically into the head.
% The label "b)" is not letterspaced.

Erfaringsphysiologien maa i det Enkelte undersøge Nervesystemets
anatomiske, chemiske og physiologiske Beskaffenhed saa nøiagtig den
vil: uden Indsigt i det Reflecteredes Forhold til det Umiddelbare,
Selvhedens til Andetheden, Grundaarsagens til Tingsaarsagerne vil den
umuligt kunne danne sig noget Begreb om Nerve- og
Hjernevirksomhedens Eenhed med Bevidsthedslivet. Tages Sagen eensidig
fra det organisk-physiske
% --- p. 273 ---
\setcounter{footnote}{0}%
Synspunkt, vil man ligesaa lidt kunne
forklare en Fornemmelse af Tilstandsforandringer i Nerver og Ganglier,
som man har været istand til at forklare Livet selv af den materielle
Vexelvirkning mellem Atomerne. Betragtes Sagen udelukkende fra et
metaphysisk Standpunkt, kan man jo vel \textit{in abstracto}
eftergjøre de til Fornemmelse og uvilkaarlig Selvbestemmelse svarende
Reflexioner; men den stofagtige Andethed er udelukket, den formelle
Reflecteren falder udenfor Virkeligheden: istedenfor at give os et
concret Begreb af det Ideelles Selvreflecteren gjennem det Reelle,
giver den eensidige Speculation kun Skyggerids af abstract sjælelige
Væsensformer.

Grundsporgsmaalet om Dyresjælenes Udodelighed er allerede besvaret
% "Grundsporgsmaalet", "Udodelighed", "Doden", "hoiere" and "godtgjore" on this
% page all print a plain o, while "eftergjøre", "gjør", "første", "først" and
% "afgjørende" on the same page print a plainly slashed ø. All as printed.
ved den Maade, hvorpaa Livets Væsen i vort Foregaaende er blevet
opfattet. Udviklingen, der gjør Resultatet forskjelligt fra
Begyndelsen, er i det Sjælelige anderledes inderligt end i det
umiddelbart Organiske. Selvudviklingens Resultat maa Sjælen altsaa
beholde i den Inderlighed, hvormed den ved Doden gaaer til Grunde; men
da Sjælefunctionerne Punkt for Punkt ere organiske Betingelser
underlagte, vil Muligheden af en psychisk Fortsættelse af Dyrelivet
ligeledes beroe paa, hvorvidt det i sin første Legemlighed allerede
udviklede Sjæleprincip kan gjennem Ideemediet indgaae til ny
Vexelvirkning med en stofagtig Andethed af hoiere Orden og derved til
Organiseren af en ny Legemlighed: Antydninger af slige Muligheder have
propædeutisk Betydning, forsaavidt de tjene til at belyse
Videnskabens Grundlag og Videns Grændse.
% A small raised mark stands after "Videnskabens"; the page carries no footnote,
% so it is taken for scanner dirt or a stray sort and is not transcribed.

\chapter*{Menneskelivet.}
% Head centred, fat-face Fraktur. The Indhold (p. 6) reads
% "Menneskelivet: Psychologien i Grundtræk"; the print divides this into two
% separate display heads of identical face and size, this one and the one on
% p. 274. Transcribed as printed, at the pages where they stand.

Det er et Misgreb at ville godtgjore Menneskets umiskjendelige
Fortrin for Dyret ved en endelig Opregning være sig af legemlige eller
sjælelige Fortrin; thi hver enkelt Eiendommelighed faaer først sin
Betydning ved det paagjældende Væsens Grundcharakteer. Menneskets ene
afgjørende Hovedfortrin viser sig deri, at det som Menneske er anlagt
paa Selv\-%
% --- p. 274 ---
\setcounter{footnote}{0}%
bevidsthed, medens Dyret som Dyr kun er anlagt paa
Sandsebevidsthed.

\chapter*{Psychologien i Grundtræk.}
% Centred display head in the same fat-face Fraktur, and at the same size, as
% "Plantelivet.", "Dyrelivet." and "Menneskelivet."; see the note on p. 273.

Det er Bevidsthedslivet, betragtet fra dets individuelle Naturside,
som i sine Overgange fra Sandsebevidsthed til Selvbevidsthed er
Psychologiens Gjenstand. Vi inddele med Hensyn hertil Psychologien i
Læren om det selvbevidste Livs naturlige Grundyttringer, om dets
absolute Indhold og dets absolute Gyldighed.

a) \emph{Det selvbevidste Livs naturlige Grundyttringer:
Sjæleevnerne.}
% Not a display head: an indented paragraph of its own, letterspaced, the label
% "a)" unspaced. Left run-in, with the emphasis the print gives it.

I skarp Modsætning til den registeragtige Opregnen af Sjæleevner,
hvormed man saa tidt overlæsser de psychologiske Lærebøger, har
Herbart ligefrem negtet Begrebet Sjæleevne Realitet. Men i Princip og
Anlæg er Sjælen lutter Evne. Opgaven er at paavise de forskjellige
Former og Retninger, hvori den ene og samme Grundevne naturligen
individualiserer og forgrener sig.

Selvhedens alleralmindeligste Aabenbaringsform er
\emph{Fornemmelse}: den fornemmer sig selv, idet den fornemmer sit
% Letterspacing runs across the line break (For-/nemmelse); both halves spaced.
Andet, og den bestemmer sit Andet ved at fornemme sig selv.
Fornemmelsen er i passiv Umiddelbarhed en \emph{Sandsning}, i ophævet
Umiddelbarhed en \emph{Reflexion}, i activ Umiddelbarhed en
\emph{Selvbestemmen}.

Som \emph{Sandsning} er den bestemte Selvfornemmelse en
\emph{Følelse} (Læren om Følelserne), Fornemmelsen af Andet en
umiddelbar Objectivering gjennem forskjelligartede \emph{Sandser}
(Syn, Hørelse o. s. v.), som Eenhed af Selvfornemmelse og Fornemmelse
af Andet omsætter den sin Umiddelbarhed i \emph{Hukommelse og
Phantasie}, og hermed begynder Reflexionen.

I Reflexionen frigjøres Sandsebillederne ved Ordet (Sprogets
Betydning) til Forestillinger og Almeenbegreber (Forstandslivet) og
forklares til Eenhed af det Subjective og
% --- p. 275 ---
\setcounter{footnote}{0}%
det Objective i Tænkningen
af det Almeengyldige (Fornuft; Forhold imellem Forstand og Fornuft).
Reflexionen fuldender sig i \emph{Tænkningen}. Men allerede den
Maade, hvorpaa Vexelbestemmelsen mellem Selvfornemmelse og
Fornemmelse af Andet, Selvbestemmelse og Bestemmelse af Andet
forklares i Reflecteren og Tænken, viser, at Fornemmelsens passive
Umiddelbarhed oprindelig er activ, at med andre Ord Selvfornemmelse
nødvendig forudsætter Selvbestemmelse.

Oprindelig er Selvbestemmelsen, den naturlige \emph{Selvyttring},
at bestemme som \emph{Drift}, Vilkaarligheden fremgaaer af
% Comma, not a full stop, before the capitalised new clause: as printed.
uvilkaarlig Tilskyndelse. Idet Driften peger mod det Objective, sætter
den Selvet i Bevægelse; Driften er Villiens Grund Men
% PRINTER'S DEFECT: the full stop after "Grund" is not printed; the compositor
% leaves an inter-sentence space and sets "Men" with a capital. As printed.
Selvbestemmelsen er reflecteret og klarer sig med Reflexionen til
\emph{Villie}: Villien, der grunder i Driften, ophæver igjen Driften
og grunder i sig selv, (Egoisme, Egoitet). Den reflecterede Drift er
\emph{Tilbøielighed}, og den hæftige \emph{Tilbøielighed Lidenskab}.
Forholdet mellem Drift og Villie sees i Vexelbrydning af Lidenskab og
Villie (Heroisme, Asketik).

b) \emph{Det selvbevidste Livs absolute Indhold: Ideernes
Psychologie.}
% Run-in like the a) rubric on p. 274: its own indented, letterspaced paragraph.

Da Bevidstheden vækkes og udvikles i Tiden ved sandselige Indtryk,
endelige Fordringer og vexlende Interesser, er det naturligt, at den
fra Begyndelsen af maa fyldes med endeligt Indhold. Som
Selvbevidsthed har den ikke desto mindre et oprindeligt Selv og just
derfor et oprindeligt Krav paa et selvgyldigt, absolut Indhold: kun
Ideen kan som Idee mætte den hungrige Subjectivitet (Aanden og det
Aandelige).

Hvor naturlig den opvaagnende Selvbevidsthed trænger til Ideens
Indhold, sees af den Oprindelighed, hvormed ethvert begavet Naturfolk
begynder med religiøse Forbilleder (Mythologie), sædelige
Grunderkjendelser, Stræben efter at give det Oversandselige sandselig
Form og bestemme det Almeengyldige. Det absolute Indhold antager
nemlig en grundforskjellig
% --- p. 276 ---
\setcounter{footnote}{0}%
Form, eftersom det oprindelig opfattes
af Phantasien, af Tænkningen eller af Villien. Form og Indhold
gjennemtrænge hinanden saaledes, at den ene i sig absolute Idee
formedelst de tre Grundformer bliver til tre Grundideer:
\emph{Skjønhed, Sandhed, Godhed}.

c) \emph{Det selvbevidste Livs absolute Gyldighed: Sjælens
Udødelighed.}
% Run-in like the a) and b) rubrics: its own indented, letterspaced paragraph.

Som Organ for evige Ideer maa Selvbevidstheden nødvendig have det
Evige i sit eget Væsen; herpaa beroer dens absolute Gyldighed, herpaa
grunder sig Læren om Sjælens Udødelighed d. v. s. om Personlighedens
evige Liv. At Sjælen er udødelig i metaphysisk Forstand, kan allerede
sees af Plante- og Dyrelivet; men Forvisningen om en existentiel, for
Menneskets Vedkommende personlig, Udødelighed kan først opnaaes ved
klar Indsigt i Forholdet mellem Ideen, navnlig det Godes Idee, og den
menneskelige Selvbevidsthed. Forsaavidt dette Forhold i sin
Almindelighed tydeliggjøres i Videnskaben, kan man jo vistnok sige, at
Sjælens Udødelighed lader sig videnskabelig bevise; men da
Virkelighedsbetingelserne for et tilkommende Liv falde udenfor al
Erfaring, og det rene Tankebeviis følgelig ikke rækker ud over det
Metaphysiske, indskrænker Beviset sig, strengt taget, til at bevise
Fordringens ethiske Gyldighed, ikke dens Opfyldelse. Den fuldkomne
Forvisning om Sjælens Udødelighed er af religios ethisk Natur og
% "religios" as printed: plain o, while "Skjønhed", "nødvendig", "udødelig",
% "Udødelighed", "tydeliggjøres", "følgelig" and "først" on this same page all
% print a plainly slashed ø. As printed.
tjener just til at vise Forskjellen imellem Tro og Viden, imellem
praktisk og theoretisk Erkjendelse (psychologisk nødvendig Adskillelse
af Religion og Videnskab).
% The short centred rule closing part II is not transcribed.
%% ============================================================
%% III. BIOLOGIE OG PSYCHOLOGIE  (pp. 277--476)
%% The Indhold (p. 6) calls this part "Psychologie og Biologie";
%% the head over p. 277 reads "Biologie og Psychologie". Both stand.
%% ============================================================

\part*{III.\\ Biologie og Psychologie.}
\addcontentsline{toc}{part}{III. Biologie og Psychologie}
\markboth{Biologie og Psychologie}{Biologie og Psychologie}

\chapter*{a) Det organiske Formaal.}
\addcontentsline{toc}{chapter}{a) Det organiske Formaal}

% --- p. 277 ---
\setcounter{footnote}{0}%
% The part head "III." / "Biologie og Psychologie." and the division head
% "a) Det organiske Formaal." stand above this on p. 277 and are set by the
% skeleton. The division head is printed in fat-face (halbfette) Fraktur;
% that is a display-head convention here, not running-text emphasis.
Den physiske Aarsag har ikke nogen særlig af Grundaarsagen uafhængig
Selvstændighed. Det er saaledes ikke Fornuften, som for at realisere sine
Fordringer, maa forefinde de physiske Aarsager givne og betjene sig af dem,
ligesom Værkmesteren af sit Værktøi; nei, det er omvendt de physiske
Aarsager, der have deres Tilværelse og Betydning i og ved Fornuftaarsagen.
Staaer det nu fast, at Naturlovens Opfyldelse vilde være umulig, dersom
Fornuften kun var en ideel Grund (\textit{ratio}) uden at være en reel Aarsag
% "ratio" and "causarum causa" are set in upright antiqua against the
% surrounding Fraktur.
(\textit{causarum causa}) dersom den i Tingenes Orden kun havde en
legislativ, ikke tillige en executiv Myndighed: da er det nu ogsaa let at
vise, hvad det vil sige, at den almindelige Fornuftaarsag kan forvandle sig
til en i sin Form eiendommelig Idee-Aarsag.

Under Andethedens Tryk er det Teleologiske tilbagetrængt; Fornuften sees i
uforanderlige Love; men Lovenes abstracte Nødvendighed staaer i Modsætning
til Tingenes abstracte Tilfældighed: den paa tilfældig Nødvendighed (dersom
$a$ er, saa er $b$) grundede Lovopfyldelse er endnu ufuldkommen. Med
% The letter-symbols a and b are lower-case antiqua run into the Fraktur text.
Selvheden sættes det Teleologiske; i den ideelle Concretion af Midler og
Formaal forsvinder den abstracte Tilfældighed, thi det Tilfældige bliver
Middel; med det Samme ophæves da ogsaa den abstracte Nødvendighed, thi det
Nødvendige bliver Maal\footnote{% printed mark: *)
Det er en Misforstaaelse, vi i denne Sammenhæng særlig maae søge at
forebygge. Det er nemlig en Fordom, som endnu kan spøge i klare Hoveder, at
Naturteleologiens Betydning kun skulde bestaae deri, at vi Dødelige
underlægge den uendelige Natur
% the printed page break inside the note falls here: p. 277 | p. 278
Planer og Øiemed, som ere afpassede efter vor egen subjective indskrænkede
Forestillingsmaade. --- „Ist es“, siger Schleiden, „vom Wahnsinn groß
% The German quotation is set in Fraktur like the surrounding Danish, so it
% is not italicised.
verschieden, wenn ein Mensch sich einbildet, in dem Innern seiner zeitlichen
Erscheinungsform irgend einen Maßstab zu finden, mit dem er das Ewige, das
Zeitlose ausmessen könnte? --- Das aber ist die Teleologie.“ Die Beseelung
der Pflanzen. Studien. Leipzig, 1855. S. 149--160.
\par
Men slige opbyggelige Expectorationer, der vel kunde klæde en tænkeløs
% "tænkeløs": the vowel after t is plainly neither the wide closed Fraktur a
% of "saa" nor the o of "nok"/"godt" on the same line; read as æ.
Katechismuspræst nok saa godt, som en genial Naturforsker, have ikke det
Ringeste med Naturteleologiens objective Begreb og de derunder hørende
Problemer at gjøre. I sin fromme Iver for at forvandle det Fornuftige til et
Uendeligt, et ἄπειρον, som ved at fornegte det Endelige og overskride Maal
og Grændse, bliver \emph{begrebløst}, og som et Begrebløst da rigtignok
% "begrebløst" letterspaced; the following "Begrebløst" is not.
hverken kan begribes af Guder eller Mennesker, glemmer Forfatteren ganske,
at den uendelige Fornuft just aabenbarer sin Viisdom og beviser sin
Fornuftighed ved at indgaae i Endelighed og derved i indre
Selvbegrændsning.}.

% --- p. 278 ---
\setcounter{footnote}{0}%
At det Tilfældige bliver Middel, vil med andre Ord sige, at det Particulære
individualiseres ved at gjennemvirkes af det Almene; at det Nødvendige
bliver Maal, sees deraf, at det ved at udvikle sit Anlægs og sine
Præformationers særskilte Bestemmelser ligeledes individualiseres;
\emph{den organiske Individuation, Lovens fuldkomne Opfyldelse, er altsaa
det sande naturlige Udtryk for Fornuftens Indhold i Fornuftens Form,
Fornuften som Idee.} Det Teleologiske, hvori de physiske Aarsager ere
% Letterspacing runs from "den organiske" to "Idee."; the sentence that
% follows is plain, though the justification of its lines is very open.
Idee-Aarsagens Momenter, er altsaa ikke blot naturligt og fornuftigt, men
som den oprindelige Eenhed af Natur og Fornuft det eneste i Sandhed
Fornuftige: Livet, den existerende Fornuft, er i sig Formaal, organisk
Formaal, thi Livet er Idee.

Physiken opfatter Naturlovene paa en saa abstract og eensidig Maade, at det
seer ud, som om disse Love ikke blot
% --- p. 279 ---
\setcounter{footnote}{0}%
forholdt sig ligegyldige til, men endog vare i Strid med det Teleologiske.
Saaledes afhandler man i Physiologien Lovene for Livets Phænomener, men ikke
for Livet selv. Og dog har Livet sin egen eiendommelige Lov ligesaavist som
det Livløse. Er det Livløse væsentlig Eet med det Selvløse og derved
Andetheden og de ligegyldige Forandringer priisgivet, saa er Livet væsentlig
Eet med Selvheden og begrundet i Selvudvikling. Har nu det Selvløse, i Kraft
af Andetheden, sin Lov, da maa ogsaa Livet, i Kraft af Selvudviklingen, have
sin Lov: kun i Ideen, hvis Teleologie underordner den lavere Lov under den
høiere, kan Naturlovenes Opfyldelse tilfulde erkjendes.

\emph{Andethedsloven er Materiens Lov, Elementernes Grundlov.} De
elementære Forandringer have i sig noget Blændende, Noget, der let vildleder
Betragteren. Selv den, der ellers vel kan forstaae, at ingen Ting kan enten
fremkomme af Intet eller blive til Intet, vil dog ved at holde sig til det
Sandselige, let hildes i den falske Forestilling, at Bevægelser og
Kraftyttringer kunne til en given Tid fremkaldes og saa igjen forsvinde saa
sporløst, som om de aldrig havde fundet Sted. Denne Indbildning hidrører
derfra, at „de forbigaaende Virksomheder“ ombyttes med andre, tabe sig i
Andet og saaledes skjule deres uendelige Fortsættelse under Forandringens
Slør. Physiken, der ved alle Love væsentlig tænker paa Andethedsloven, har
imidlertid selv omsider opdaget Vildfarelsen, idet den har opdaget, at ingen
Kraft gaaer tabt, men at Kræfterne ombyttes, forvandles og erstatte
hinanden. Det staaer altsaa i Kraft af Andethedsloven fast, at Følgerne af
en engang indtraadt Virksomhed aldrig kunne tilintetgjøres, at det Skete med
andre Ord aldrig kan gjøres uskeet. Denne Opdagelse er af gjennemgribende
Betydning i Læren om Livet. I Kraft af Forandringsloven er enhver
Tingsaarsag, som i de chemiske Kategoriers Analyse allerede er viist, en
Betingelse, og enhver virksom
% --- p. 280 ---
\setcounter{footnote}{0}%
Betingelse igjen en medvirkende Aarsag. Vel er Betingelsen et Moment, men et
Moment, der udenfra bringes sammen med en anden Betingelse, et andet Moment.
Selv naar det ene Moment med physisk Nødvendighed medfører det andet, vil
det beroe paa Omstændighederne, altsaa paa andre udenfra tilstødende
Betingelser, hvorvidt alle i en given Forudsætning indeholdte Betingelser
skulle træde i Virksomhed eller ikke. I Betingelsernes Verden hersker
Andethedsloven paa alle Punkter: at forgude Andethedsloven, at bevæge sig i
Betingelsernes uendelige Kredsløb, uden deri at opdage det Ubetingede, er at
tænke sig ind i det Tankeløse, at leve sig ind i det Livløse.

\emph{Selvhedsloven er Livets Lov, Organismernes Grundlov.} Ligesom det med
indre Nødvendighed følger af det Elementæres materielle Natur, at det maa
være Andethedsloven underlagt, saaledes følger det med indre Nødvendighed af
Livets ideelle Natur, at det maa være Selvhedsloven underlagt. Thi hvad er
det Ideelle? Det negativt Reelle, det, der ikke paa endelig Maade lader sig
enten sætte sammen med Andet eller adskille fra Andet, men gjennemtrænger
Andet, er i væsentlig Forskjel fra Andet Eet med Andet, har i sig sit Andet,
og just af den Grund forholder sig gjennem Andet til sig selv. Men at
forholde sig gjennem Andet til sig selv og derved udvikle Andetheden i sig
selv er at opfylde Selvhedsloven.

For at Andetheden og Selvheden skulle kunne bestaae i afgjørende
Modsætning, maae Modsætningens Factorer forudsætte og bekræfte hinanden
gjensidig: \emph{det er Andethedens Bestemmelse at tjene Selvheden; det er
Selvheden væsentligt at trænge til Andetheden.} At Planten visner, at Dyret
døer, viser tydeligt nok, hvad det vil sige, at ogsaa det Levende er
Andethedsloven og dermed Forkrænkeligheden underlagt. Livsideen behøver det
udvortes Andet for derigjennem at udvikle sit Anlæg, sit ind\-%
% --- p. 281 ---
\setcounter{footnote}{0}%
vortes Andet. Men det udvortes Andet er kun en materiel Betingelse;
Selvheden, der kun betinges derved, at den med indre Nødvendighed betinger
sig selv, er det i Betingelsernes System Ubetingede: Livsudviklingen er
væsentlig Selvudvikling, og Selvudviklingen, teleologisk seet,
Idee-Udvikling. Paa Selvudviklingens Lov beroer Livsideens
Uforkrænkelighed.

Livet selv, det Liv, der bortviger fra den udgaaende Plante, Sjælen selv,
den Sjæl, der forlader det døende Dyr, kan ikke blive til Intet. Var
Livsvirksomheden end i Eet og Alt Andethedsloven underlagt; den kunde dog,
som vi alt have seet, umuligt blive til Intet. Staaer det nu allerede i
Kraft af Andethedsloven fast, at Følgerne af en engang indtraadt Virksomhed
aldrig kunne tilintetgjøres, saa staaer det i Kraft af Selvhedsloven
ligesaa fast, at Resultatet af den under Organisationen opnaaede
Selvudvikling hverken kan udlignes ved Andet eller omsættes i Andet. Men
naar Selvudviklingens Resultat hverken kan blive til Intet eller omsættes i
Andet, maa den evig forblive i Selvhedens Erindring.
% "Erindring." carries a blot of scanner dirt over the i; the reading is not
% in doubt.

Det rene, hvide Lys kan brydes i Prismet og skille sig i Farvestraaler; men
naar Prismet borttages, og Brydningen ophører, forsvinde Farverne, og vi
have igjen, som før, det hvide Lys. Anderledes med Livet! Den ideelle
Straalebrydning, den indre Selvudvikling, Principet under sin organiserende
Virksomhed opnaaer, bliver ikke forstyrret, fordi den materielle Organisme
forstyrres. Det Skete kan ikke gjøres uskeet; den Rigdom af indre
Bestemmelser, der fremkom under Articulationen af den indvortes, Selvheden
selv iboende Andethed, er en ufortabelig, uforkrænkelig Eiendom, som
Principet evig maa beholde i Ideen. For Lyset, der brydes i Prismet, er
baade Prisme og Brydning et Tilfældigt og Ligegyldigt; derfor kan det Hele
ligesom gaae tilbage igjen, Brydningen ophøre, Farverne forsvinde, og Lyset
efter Brydningen blive,
% --- p. 282 ---
\setcounter{footnote}{0}%
hvad det var før Brydningen, --- det rene, hvide Lys. Men saaledes kan
Livsprincipet, for hvis indre Selvudvikling den materielle Organisme
ingenlunde er noget Tilfældigt og Ligegyldigt, men tvertimod noget
Væsentligt og Nødvendigt, umuligt blive igjen, hvad det var før
Organisationen, --- det rene almindelige Princip.

At gjøre Andethedsloven til Livets Grundlov er Nonsens; at lade
Selvhedsloven gjælde, og dog antage, at Udviklingens ideelle Resultat skulde
kunne forsvinde, og Livet tabe sig i Intet, fordi den materielle,
Andethedsloven underlagte Organisme hjemfalder til Forkrænkelighed, er en
meningsløs Modsigelse.

Som Teleologiens og Livets Lov er Selvhedsloven da ogsaa \emph{Erindringens
Lov.} Saalænge de organiske Individers høieste Bestemmelse sættes deri, at
% Letterspacing begins only at "Erindringens Lov." on the second line; the
% clause before it is plain.
de ved Forplantning skulle sikkre Artens Vedligeholdelse, er Formaalet kun
eensidig opfattet. De til Modenhed udviklede Individer leve her for
Tilblivelsen af andre Individer, som igjen skulle modnes til at leve for
Tilblivelsen af andre, o. s. fr. Istedetfor at have sit Formaal i sig selv,
har Selvheden her sit Formaal udenfor sig selv --- i en anden Selvhed: det
seer fra dette Synspunkt ud, som skulde Selvhedens høieste organiske Act
just indføre Andethedens Herredømme. Denne Modsigelse løses, naar
Individernes høieste Bestemmelse sættes, ikke i Forplantningen, men i
Erindringen. Ved Forplantningen gaaer Selvheden ind paa Andethedens
materielle Vilkaar og kommer saa paa disse Vilkaar til at existere. Da hver
Existens er forkrænkelig, maa den ene efterhaanden træde istedetfor den
anden; men forsaavidt Selvhedens Opholdelse ene og alene skal bestaae i
Forplantningen, er den sig saaledes fortsættende Selvopholdelse væsentlig en
sig fortsættende Selvopløsning. Den organiske Selvhedsact beseirer
Forkrænkeligheden ved at forplante det Forkrænkelige; Selvopholdelsens
fortløbende Række viser en Gjen\-%
% --- p. 283 ---
\setcounter{footnote}{0}%
tagelse af Individer, hvis indskrænkede Natur er incommensurabel med
Selvopholdelsens Fordring: Individet i det titusinde Led er ligesaa
forkrænkeligt som Individet i det tiende Led; her er i den uendelig
fortløbende Række intet væsentligt Fremskridt, intet Endemaal.

Den organiske Udvikling er væsentlig Selvudvikling; men en Existens, hvis
Væsen er Selvudvikling, kan kun fuldendes i den fuldendte Inderliggjørelse
af sit udviklede Indhold d. e. i en tilsvarende Existenserindring. Først ved
Erindringen bliver det i Existensen Forkrænkelige hævet; men kun ved Døden
bliver den gjennemførte Erindring mulig. Den saakaldte „moleculære Død“, der
er en Betingelse for den Reproduction, den Foryngelse ved Assimilation, uden
hvilken Livsvirksomheden ikke kan fortsættes, viser os allerede en organisk
Sammenhæng mellem Døden og Erindringen. I Erindringen gaaer den umiddelbare
Existens med dens ydre og materielle Skikkelse tilgrunde; men den saaledes
ophævede Existens er ikke tilintetgjort; den er tvertimod fuldendt --- i
organisk Erindring; Existenserindring er Livsideens Endemaal: \emph{kun i
substantiel Erindring gaaer det Typiske sammen med det Teleologiske.} ---
Som Forbillede kan Paradigmet jo først virkeliggjøres i et fuldstændigt
Kredsløb af successive Afbilleder. Dersom nu det ene af disse sandselige
Afbilleder under Livets Løb uden blivende Resultat fortrængte det andet, saa
vilde Erindringen af de forsvundne Billeder kun foregaae i vor Bevidsthed,
og den sig igjennem de mange vexlende Skikkelser aabenbarende
Grundskikkelse være et af vor subjective Tænkning abstraheret Schema, uden
objectiv Realitet. For at Forbilledet i Sandhed kan realiseres, maa ethvert
af de forsvindende Afbilleder vedblivende reflectere sin Forsvinden i den
fremadskridende Udvikling, hver Skikkelsesreflex i den organiske
Selvhedserindring være bevaret som objectivt Moment, naar det sidste
Stadium, det Stadium, der betegner Forbilledets
% --- p. 284 ---
\setcounter{footnote}{0}%
synlige Fuldendelse, er naaet\footnote{% printed mark: *)
% The whole note is German, set in a narrower German Fraktur than the
% surrounding Danish; being Fraktur, it is not italicised.
Zunächst ist hier darauf ein entschiedenes Gewicht zu legen, daß man unter
dem Begriffe der thierischen Individuen nur materiell abgeschlossene
morphologische Facta subsumiren darf, daß man also wohl von verschiedenen
Formen der Individuen sprechen, aber nicht einzelne Formzustände eines
Körpers als eben so viele Einzelindividuen unter einem die ganze
Formenreihe begreifenden Gesammtindividuum begreifen darf, wie es
neuerdings \emph{Reichert} that, wenn nicht der wesentlichste Gesichtspunkt
bei Beurtheilung thierischer Individualität, welchen der morphologisch
bestimmte \emph{Körper des Thieres} bietet, verrückt werden soll. Victor
% "Reichert" and "Körper des Thieres" are letterspaced; "Victor Carus" is not.
Carus. System der thierischen Morphologie. Leipzig 1853. S. 253--54.}. Men
Fuldendelsens sidste synlige Billede er jo selv et Afbillede, et
forsvindende Sandsebillede; \emph{altsaa først naar den sidste materielle
Skikkelse opløser sig og forsvinder i den paradigmatiske Inderlighed, den
sin Udviklings Indhold erindrende Selvhed, er Kredsløbet sluttet, og Maalet
naaet.}

Som det absolute Endemaal, Ideens fuldkomne Form, skulde Existenserindringen
altsaa være den Kategorie, hvori Tilværelsens Modsigelser, Livets og Dødens,
omsider fandt deres Opløsning. Men forholder dette sig nu virkelig saaledes?
Istedetfor at naae et Vendepunkt, hvor alle Modsigelser løses, synes vi jo
tvertimod --- dette er idetmindste det foreløbige Indtryk --- at støde paa
et Knudepunkt, hvori alle Modsigelser ere samlede, og hvorfra de igjen, naar
Reflexionen sættes i Bevægelse, myldre frem i uoverskuelig Mængde. „Skal en
Existenserindring have Realitet, saa maa en Erindringsexistens jo ogsaa have
Realitet. Men hvad er en Erindringsexistens? Et Tilværende, som ikke er til!
Et saadant Tilværende kan t. Ex. ikke være i Hvile, thi det Hvilende har kun
Tilværen ved en umiddelbar Modsætning til det Bevægede, en endelig
Modsætning, der kun finder Sted i den udvortes Verden. Ligesaa lidt kan her
i et saadant Til\-%
% --- p. 285 ---
\setcounter{footnote}{0}%
værende være nogen Bevægelse; thi Bevægelse udtrykker Overgang, men en
virkelig Overgang forudsætter en umiddelbar Andethed, som Erindringen ikke
kan have. Existenserindringen kan ikke bestaae i en reen Uvirksomhed; i den
rene Uvirksomhed er der ingen Reflecteren, og hvor der ingen Reflecteren er,
er der heller ingen Erindren. Men Existenserindringen kan heller ikke tænkes
at bestaae i Virksomhed; Virksomhed forudsætter, at virkelige Modsætninger
sættes og overvindes; men virkelige Modsætninger med tilsvarende Spændinger
findes kun i en virkelig, ikke i en ophævet Existens. Den i Erindringen
omsatte og derved ophævede Existens kan ikke være til i Tiden, thi Tiden
fordrer en endelig Succession af vexlende Tilstande; men en Erindring, hvori
Alt er sluttet og fuldendt, stiller jo Alt i Beroe. Ligesaa lidt kan en
Existenserindring foregaae udenfor Tiden, i det Tidløse; Erindringen maa jo
bestaae i en idelig Tilbagekalden af det Forsvundne, en uophørlig Gjentagen
af det, som har været; men en saadan Tilbagekalden, en saadan Gjentagen maa
nødvendigviis have Successionsmomenter og kan umulig foregaae i et tidløst
Medium.“ --- Saaledes kunde man da vedblive med Modsigelse paa Modsigelse i
det Uendelige.

Det er imidlertid ved nogen Eftertanke ingenlunde vanskeligt at opdage,
hvorfra alle disse Modsigelser egentlig have deres Udspring; de hidrore
% "hidrore" as printed: plain o, verified at 600 dpi against the plainly
% slashed ø of "uophørlig", "nødvendigviis", "Tidløse", "opløse", "søge" and
% "løste" on this same page --- and against "hidrører" with a clear slash on
% p. 279. Not normalised.
tilhobe fra overfladisk Kritik, fra Misforstaaelse af Ideens dialektiske
Væsen. Ved at stille de forskjellige Kategorier eensidig og kritisk op imod
hinanden: Hvile imod Bevægelse, Uvirksomhed imod Virksomhed o. s. v., kan
man jo sagtens opløse den ideelle Concretion i lutter Modsigelser. I Ideen
selv er der hverken Hvile for sig eller Bevægelse for sig; men Bevægelse i
Hvile, Hvile i Bevægelse. Det er den kritiske Philosophies Grundvildfarelse,
at den, istedetfor at søge Forstandsbegrebernes Modsigelser løste i den
concrete Idee, tvertimod holder Begreberne fast i endelig Modsætning til
Ideen og lader de Antinomier, der kun have
% --- p. 286 ---
\setcounter{footnote}{0}%
Gyldighed udenfor Ideen, optræde, som om de udtrykkelig angik
Ideens Væsen. Ved at paatvinge Ideen, hvad der ligefrem strider
imod dens Væsen, er det en let Sag at vække Tvivl om Ideens
objective Realitet.
\chapter*{b) Livets Trin.}
\addcontentsline{toc}{chapter}{b) Livets Trin}

% The head `b) Livets Trin.' stands MID-PAGE on p. 286: four lines closing
% `a) Det organiske Formaal' precede it. Those lines, together with the
% page marker for 286, are in .parts/p286-tail-of-a.texfrag and belong
% at the END of the pp. 277--285 block, above the \chapter* head.
% This fragment therefore opens BELOW the head and carries no p. 286 marker.
%
% The whole Propædeutik quotation that opens the division is letterspaced.
\emph{„Forsaavidt Naturen er Idee, er Udviklingen af Naturbegrebet en
Opgave for Philosophien; men da Ideen som Natur kun er, hvad den er i
sine Frembringelser, Phænomener og Love, maa Udviklingen af Begrebets
Indhold nødvendigviis henhøre under Empirien. Det er just en
Eiendommelighed ved Naturbegrebet, at det med lige Kraft adskiller og
forbinder Philosophie og Empirie“.} (Propæd. S. 191).

Som Livets Grund er Ideen Natur, som organisk Natur er Ideen Liv.
Erkjendelsen af, hvad Livet er, forudsætter væsentlig, fra hvilken Side
Opgaven end betragtes, en tilsvarende Erkjendelse af, hvad Naturen er;
de, der i Livet see en uopløselig Gaade, maae da ogsaa tilstaae, at
Naturen, det Naturlige selv, er en uopløselig Gaade.

\emph{Er det muligt for den menneskelige Bevidsthed at trænge ind i
Naturens Indre, eller er det umuligt?}

Et humant videnskabeligt Svar vil, hvad dette Spørgsmaal angaaer, neppe
lyde enten paa afgjørende Benegtelse eller paa afgjørende Bekræftelse.
„Hvad er det, som gjør, at mit Legeme voxer og trives, saalænge jeg kun
tilfredsstiller min Hunger, min Trang til at trække Veiret og til at
holde mit Legeme varmt? som gjør, at mit Hjerte og min Puls slaaer sin
jævne Takt, uden at jeg endog formaaer at hindre det? Hvad gjør, at jeg
derimod er Herre over mine Lemmer? Allerede det spæde Barn stirrer paa
sine afklædte Been, og
% --- p. 287 ---
\setcounter{footnote}{0}%
glæder sig synligen over, at de bevæge sig efter dets Behag, og flux som
det behager det . . . . Mere maaskee skjønne vi alle paa at kunne see,
høre og tale; Enhver kjender noget til Øiets og Ørets underfulde
Bygning, og hører den ikke ugjerne forklaret noget noiere. Besvarelsen
af alle disse Spørgsmaal, Oplysningen om alle disse Gjenstande henhøre
under Physiologien. Vel maae De ikke troe, at Besvarelsen eller
Oplysningen altid er afgjørende; ind i Naturens egentlige Indre trænger
intet menneskeligt Øie. Men ogsaa at see, hvad der er givet, og ane,
hvad der endnu ligger skjult, yder en Tilfredsstillelse, som Ingen, der
har Sands for Naturen, skulde negte sig selv at erholde.“ (Eschricht).

Fra Dannelsens almindelige Standpunkt er en slig moderat Synsmaade
vistnok i det Hele at anbefale; men naar Erkjendelsesspørgsmaalet
forvandler sig til et videnskabeligt Principspørgsmaal, er
Middelveistheorien med alle sine Moderationer ubestemt, svævende,
utilstrækkelig. Hvad vil det sige, at det menneskelige Øie vel paa en
Maade kan trænge ind i Naturens Indre, men dog ikke ind i „Naturens
egentlige Indre“? Hvilken Forskjel er her da imellem Naturens egentlige
og uegentlige Indre? Det lyder fra den ene Side:
% The two German verse blocks are Fraktur, not antiqua: NOT italicised.
\begin{quote}
„Ins Inn're der Natur dringt kein erschaffner Geist,\\
Zu glücklich, wem sie nur die äuss're Schale weis't!“
\end{quote}
fra den anden Side:
\begin{quote}
„Ins Inn're der Natur ---“\\
O Du Philister!\\
„Dringt kein erschaffner Geist.“\\
Mich und Geschwister\\
Mögt Ihr an solches Wort\\
Nur nicht erinnern;\\
Wir denken: Ort für Ort\\
Sind wir im Innern.“
% The closing „ of this last mark has no opening partner anywhere on
% pp. 286--299: the print sets a bare high quote after `Innern.'
% Transcribed as printed; it leaves this batch's „/“ count off by one.
\end{quote}

% --- p. 288 ---
\setcounter{footnote}{0}%
Her er i disse Stemningers Udtalelse fra begge Sider Noget, der
tiltaler, og Noget, der frastøder den sunde Sands.

At negte den menneskelige Aand al Evne til at erkjende Tingenes Væsen
gaaer ikke an (Propæd. S. 63); er end Naturens Rige et
Hemmelighedernes Rige, saa er det dog et Rige af lovbundne
Hemmeligheder\footnote{% printed mark: *)
Disse --- Psychologiens --- Hemmeligheder ville formodentlig aldrig
blive opklarede; imidlertid have Physiologernes Bestræbelser gjort det
muligt idetmindste at skabe en Videnskab om de Livs%
% The compositor breaks `Livsphænomener' across the line here and sets NO
% hyphen: the line ends bare `Livs'. Joined; the omission is logged.
phænomener, der staae i Forbindelse med Nervesystemet; og „saaledes“,
for at bruge Professor Huxley's smukke Udtryk, „er der fra den regelløse
Hemmelighedsfuldhed, der er Uvidenhedens Enemærker, lagt endnu en stor
Provinds til Videnskaben, den \emph{lovbundne Hemmelighedsfuldheds
Rige}“. (Lewes).
% The letterspacing begins at `lov-' at the line end, not at `den':
% transcribed as printed, i.e. the whole word `lovbundne' is spaced.
}. At indrømme den menneskelige Aand Evne til en absolut Viden, en
fuldkommen Indtrængen i Naturens Allerinderste, gaaer heller ikke an
(Propæd. S. 66---67); den speculative Philosophie, der har talt det
dristigste Ord om Naturens Forklarelse i Aanden --- „So ist die Natur
die Braut, mit der der Geist sich vermählt“ (Jvfr. Propæd. S. 193) ---
har ogsaa, hvad de for Natur%
% As with `Livsphænomener' above: `Naturerkjendelsen' is broken at the
% line end with NO hyphen in the print. Joined.
erkjendelsen afgjørende Kjendsgjerninger og Love angaaer, gjort sig
skyldig i de største Misgreb. Selv Goethe, der med sin geniale
Naturanskuelse: „wir denken: Ort für Ort sind wir im Innern“, synes at
stille sig saa ubetinget paa den „absolute Videns“ Standpunkt, skifter
dog Tonart, naar han i sin Elegie: „Die Metamorphose der Pflanzen“,
dybsindig minder os om „ein geheimes Gesetz, ein heiliges Räthsel.“

Indrømmes det, at det menneskelige Øie kun til en vis Grad, en vis
Grændse kan trænge ind i Naturens Indre, saa spørges: \emph{til hvilken
Grad, til hvilken Grændse?} Den Videnskab, der principielt besvarer
dette Spørgsmaal, er
% --- p. 289 ---
\setcounter{footnote}{0}%
ikke Naturlæren, men Philosophien\footnote{% printed mark: *)
Philosophien er nemlig en videnskabelig Ideelære (Propæd. S.
118---198) under den Forudsætning, at den tillige er en videnskabelig
Erkjendelseslære. (Propæd. S. 5---67).
}. „Vi møde her det Forunderlige, at den Videnskab, som ene blandt alle
ret forstaaer at undersøge Naturens Gjenstande, Naturens Kræfter,
Naturens Love, kort: Alt, hvad Naturens er, ikke desto mindre bliver
% `hvad Naturens er' as printed (600 dpi: the `s' is unmistakable);
% the sense wants `Naturen'. Not corrected.
forlegen og begynder at stamme og læspe, naar den skal sige bestemt,
hvad Naturen er . . . At Naturvidenskaben ikke kan sige, hvad Naturen
er, maa dog ikke misforstaaes, som om den enten ikke anerkjendte
Naturbegrebets Realitet, eller ikke var istand til at forbinde en bestemt
Tanke med Ordet: Natur; tvertimod! hvad Naturen er, kan Naturlæren paa
sin Viis tydelig nok fremstille, men, vel at mærke, derved, at den
fremstiller, ikke Naturen selv, men Naturens Kjendsgjerninger og Love,
altsaa ikke ligefrem, men middelbart, og det af den gode Grund, at
Naturen er Idee, og ingen Erfaringsvidenskab istand til at fremstille
Ideen som Idee.“ (Phil. Propæd. S. 191). „Enhver Videnskab fremstiller
et Ideens Indhold, en Side af Ideen, en særegen Idee, men middelbart;
kun Philosophien lægger an paa at fremstille Ideen som Idee.“ (Propæd.
S. 118).

Ligesom Erkjendelsen af Ideens logiske Væsen beroer paa Indsigt i de
logiske Kategoriers System, saaledes beroer Naturideens Erkjendelse paa
klar og tydelig Indsigt i Naturkategoriernes System. Hvor umuligt det er
for Philosophien at sikkre sig nogen sand Opfattelse af Naturbegreberne,
naar den ikke Skridt for Skridt bliver veiledet af Empirien, er i det
Foregaaende fra mange Sider og paa mange Maader indskjærpet. I denne
Sammenhæng skal det nu ligesaa eftertrykkeligt indskjærpes, at alle
Naturvidenskabens Opdagelser, om end disse glimrende Opdagelser
fortsættes gjennem Aartusinder, dog i al deres Mangfoldighed ere
utilstrækkelige til
% --- p. 290 ---
\setcounter{footnote}{0}%
at aabne Menneskets Øie for Naturens „egentlige Indre“, dersom
Philosophien ikke følger med og bearbeider Systemet af de til
Kjendsgjerningerne svarende Kategorier saaledes, at Ideen speiler sig
deri\footnote{% printed mark: *)
Saa sikker Naturvidenskaben er, naar der handles om Kjendsgjerninger og
Love, ligesaa usikker er den, naar det gjælder en Bedømmelse af
Kategoriernes objective Gyldighed (Jvfr. Propæd. S. 182). Ikke nok, at
vi, ifølge Physiologiens Tilstaaelse, „ikke vide, hvad Livet er;“ vi
„vide jo“, ifølge Physikens Forklaring, „egentlig heller ikke, hvad
Elektricitet er, hvad Kraft er“, o. s. v. Skulde man i denne Henseende
tage Naturlærens Talemaader for ramme Alvor og drage Conseqvenserne, da
vilde en egentlig Naturvidenskab, trods alle Opdagelserne, være en
Umulighed. Det seer saa beskedent og dog saa skarpsindigt, saa resigneret
og dog saa høividenskabeligt ud, naar man saadan uden videre erklærer:
„vi vide ikke, hvad Kraft er.“ Hvilke strenge Fordringer til den
menneskelige Viden bliver her ikke antydet ved en saadan Erklæring! „Die
Kraft ist nichts als eine verstecktere Ausgeburt des unwiderstehlichen
Hanges zur Personification, der uns eingeprägt ist; gleichsam als ein
rhetorischer Kunstgriff unsres Gehirns, das zur tropischen Wendung
greift, weil ihm zum reinen Ausdruck die Klarheit der Vorstellung fehlt.
In den Begriffen von Kraft und Materie sehen wir wiederkehren denselben
Dualismus, der sich in den Vorstellungen von Gott und Welt, von Seele und
Leib hervordrängt.“ (Du Bois-Reymond. Untersuchungen über thierische
Electricität. Vorrede). --- Den dybere Grund til en saa fortvivlet
Opgiven af al menneskelig Viden er imidlertid ikke vanskelig at opdage.
Sagen er simpelthen denne, at Naturforskeren, der er i Vane med at
henføre alle Erkjendelser til umiddelbare Sandseindtryk, løber vild i
Bestemmelsen af det, der kan sandses og forestilles, og det, der efter
sin Natur hverken skal sandses eller forestilles, men skal tænkes. At
Kraften ikke umiddelbart lader sig forestille, er ligesaa vist som, at
Kraften ikke er blaa, og at man ikke under Mikroskopet er istand til at
paavise særegne enten for Tiltrækning eller Frastøden synlige
Kraftstriber; Kraften er som Kraft, i Modsætning til Materien som det i
Phænomenets Umiddelbarhed Sandselige, et Reflexionsbegreb. Kommer nu
Tænkningen, blændet og blindet af den uophørlige
% the note runs on over the page break here, and is continued at the
% foot of printed p. 291
Sandsen, Mikroskoperen og Forestillen, i den Grad fremmed til sig selv,
at den end ikke i et Elementærbegreb som Kraften kan kjende sit eget
væsentlige Object, hvormeget mindre vil en saadan forblindet Tænkning da
være istand til at erkjende Kategorier som: \emph{Natur, Liv, Sjæl}, der
jo ere langt mere reflecterede. Ved at forkaste de i sig reflecterede
Naturkategorier, har man imidlertid ikke alene brudt Staven over
Philosophiens „Ideer“, men selv over Naturvidenskabens „Love“; thi ogsaa
Lovene ere jo Reflexionskategorier: man slipper aldrig forbi Hume!
}.

% --- p. 291 ---
\setcounter{footnote}{0}%
Hvad her er sagt om Naturideen i Almindelighed, gjælder i særegen
Forstand om Livsideen. Hvorvidt Philosophien her er istand til at
forvandle de virkelige Kjendsgjerninger til ideelle Momenter, vil da
bedst kunne sees, naar Indholdet articuleres, naar Opgaven tages concret,
naar vi, med andre Ord, forsøge en Fremstilling af Livsideen paa dens
forskjellige Trin som \emph{Planteliv, Dyreliv, Menneskeliv}.

% Display sub-head, letterspaced in the print, set flush with the
% paragraph indent; Greek label, so it is a \subsubsection* like the
% lettered heads of A/B/C.
\subsubsection*{α) Plantelivet.}

Som organisk Natur har Ideen sin første umiddelbare Existens i den blot
vegeterende Organisme, i Planten. Forsaavidt den vegetative Organisme er
Ideens Afpræg, maa den ved Ideen selv bestemte Formvirksomhed i alle
væsentlige Træk nødvendigviis være typisk-teleologisk. Men da Telos
umiddelbart er en vegeterende Leven, og allerede Cellen vegeterer som
levende Celle, saa indsees heraf, at det i Planteverdenen ikke saameget
er det Teleologiske, som derimod det Typiske, hvorpaa Mangfoldigheden af
charakteristiske Former væsentlig beroer. Af Selvhedens Paradigme
fremgaaer, efter hvad vi i det Foregaaende have seet, Plantens
eiendommelige Skikkelse; med den eiendommelige Skikkelse er Telos sat
eiendommeligt; med det eiendommelige Telos ere de vegetative Midler, der
jo, væsentlig seet, ikke betyde Andet end den ved Endeligheden,
Relativiteten, Middelbarheden, motiverede Grunddeling af det iboende
Telos selv, da ligeledes blevne eien\-%
% --- p. 292 ---
\setcounter{footnote}{0}%
dommelige. Formedelst Paradigmernes charakteristiske Mangfoldighed
realiseres altsaa Plantelivets almindelige Formaal paa saa mangfoldige
Maader, at den ene Vegetationens Grundidee derunder sees at udfolde sig
i et System af særskilte Ideer. For at tydeliggjøre det i denne
Udfoldning Principielle, ville vi i nærværende Sammenhæng henvende
Opmærksomheden paa: \emph{Planteverdenens Grundformer, Metamorphoser og
Individuationer}.

αα) \emph{Grundformer.} Man kan i Overeensstemmelse saavel med „det
kunstige“ som med „det naturlige“ System inddele Planterne i to store,
omfattende Grupper: \emph{Kryptogamer} (de Blomsterløse) og
\emph{Phanerogamer} (Blomsterplanter); men Inddelingens videnskabelige
Værd vil altid afhænge af det eiendommelige Lys, hvori de udhævede
Kjendemærker betragtes.

At her imellem de kryptogame og phanerogame Dannelser, saavel i det Hele
som i det Enkelte, viser sig en Forskjellighed, kan ikke negtes; men
hvad betyder nu egentlig en saadan Forskjellighed? Forskjelligheden kan
indskrænke sig til en puur og bar Ulighed, hvorved den saa rigtignok
taber al videre Interesse, med mindre den da skulde tjene en blot
heuretisk Classificationsinteresse\footnote{% printed mark: *)
Til Brug ved den blot heuretiske Classification griber man ethvert
Mærke, selv det meest charakteristiske, i dets rene Tilfældighed; det
linneiske System (Sexualsystemet), der er bygget paa
Befrugtningsorganernes Antal og øvrige Forhold, afgiver herpaa et
storartet Exempel. Ifølge dette System kunne Planter, der iøvrigt ere
hinanden aldeles ulige, komme til at staae hinanden ganske nær, naar de
med Hensyn paa det valgte Mærke tilfældigviis ere lige; saaledes høre
t. Ex. en Berberisse (\textit{Berberis}) og en Vintergjæk
(\textit{Galanthus}) begge til 6te Klasse, en Ask (\textit{Fraxinus}) og
en Salturt (\textit{Salicornia}) begge til 2den Klasse, skjøndt Naturen
ved en hoist forskjellig Organisationsmaade har skilt Berberisien vidt
fra Vintergjækken, Asken vidt fra Salturten. Paa den anden
% the note runs on over the page break here and is finished at the foot
% of printed p. 293
Side kunne Planter, som i alt Øvrigt vise sig indbyrdes beslægtede,
blive fjernede langt fra hinanden, naar det træffer sig, at de med
Hensyn paa det valgte Kjendemærke ere forskjellige; saaledes hører,
ifølge det linneiske System, Agermaane (\textit{Agrimonia}) til 11te,
Mjødurt (\textit{Spiræa}) til 12te, Løvefod (\textit{Alchemilla}) til
4de, og Bibernelle (\textit{Poterium}) til 21de Klasse, uagtet de efter
deres Natureiendommelighed alle ere Medlemmer af een og samme Familie,
nemlig Rosenfamilien. --- At Støvdragere og Støvveie forøvrigt ere ikke
alene sikkre Kjendemærker, men væsentlige Individualformer, faaer her
ingen Betydning. Det i et „kunstigt System“ Kunstige er, at et naturligt
Mærke ved Abstractionen udrives af sin naturlige Forbindelse og
fastholdes for sig til Brug ved Classificationen.
}. Forskjelligheden kan gaae dybere, indtil
% --- p. 293 ---
\setcounter{footnote}{0}%
Ueensartethed, uden at deraf endnu følger, at de Ueensartede forholde
sig til hinanden som det Ufuldkomnere til det Fuldkomnere, det Lavere
til det Hoiere. Men sæt endog, at den kryptogame Dannelse forholdt sig
til den phanerogame, Sporehuset til Frøgjemmet, Sporen til Kimen,
o. s. v., som det Ufuldkomnere til det Fuldkomnere: fra hvilket
Synspunkt og efter hvilken Maalestok skal da den til eiendommelig
Modsætning forvandlede Forskjel nærmere bestemmes?

Da Plantelivet væsentlig gaaer op i organisk Formen, umiddelbar
Organiseren, maa den paalidelige Maalestok for Lavere og Hoiere
upaatvivlelig søges i det Organiseredes indre organiske Beskaffenhed.
For at vurdere denne Beskaffenhed maa Forholdet imellem Anlæg og
Udførelse staae klart for Betragtningen. Et i sig lavere Anlæg kan jo i
et givet Tilfælde have fundet en rig og alsidig Udførelse, medens
Udførelsen af et i sig hoiere Anlæg er bleven staaende paa et lavere
Trin; til at bedømme Modsætningen mellem det Lavere og det Hoiere er det
altsaa ikke nok at see paa det Udvortes, den umiddelbare Skikkelse; her
maa væsentlig sees paa det Indvortes, paa \emph{Grundformen}.

Vi sætte det kryptogame Sporegjemme (\textit{sporangium}) imod det
phanerogame Frøgjemme (\textit{pericarpium}). Skulde
% --- p. 294 ---
\setcounter{footnote}{0}%
vi nu dømme efter det Udvortes alene, da vilde det i sin Grundform
Lavere vistnok meget ofte fremstille sig for os, som om det var det
Hoiere. Saaledes have Mosserne et af et stort Cellecomplex dannet
Sporehuus; Sporehuset aabner sig paatværs, og i ikke faa Tilfælde
saaledes, at en mellem Laaget og Huset værende Springring
(\textit{elater}) kaster Laaget af; det saakaldte Peristom har snart en
enkelt, snart en dobbelt Krands af Tænder, o. s. v.; her er et rigt
Udstyr. I Sammenligning med saa rigt udstyrede Sporehuse maa Frøgjemmet
hos mangen phanerogam Plante, især naar dette Frøgjemme, som hos
\textit{Lathyrus latifolius}, bliver staaende ved Formen af et
almindeligt fladt Blad, unegtelig tage sig fattigt ud og forekomme os
baade simpelt og ubetydeligt. Anderledes, naar vi see paa Grundformen,
det eiendommelige Anlæg, og det til samme svarende morphologiske
Grundlag.

Hvorpaa skal nu det morphologiske Grundlag kjendes? Naturligviis paa det
Element, hvortil det i Dannelsen Eiendommelige væsentlig støtter sig. Et
saadant Element er, hvad Sporehusets Dannelse angaaer, \emph{Cellen};
hvad Frøgjemmet angaaer, er det \emph{Bladet}. Kalde vi for Organismens
Vedkommende Cellen det primære Element, et Dannelseselement af første
Orden, saa maae vi kalde Bladet, den Dannelse, hvori Celleelementet er
nedsat til Moment, et hoiere Element, et Element af anden Orden. Hvor
hoit Dannelsen af Sporehuset end stiger\footnote{% printed mark: *)
Hos Lavarterne (\textit{Lichenes}) ere Sporerne i et bestemt Antal
(4 à 8), omgivne af en enkelt forlænget Celle (\textit{theca}), som i
dette Tilfælde er Sporegjemme; et som oftest stort Antal af slige
sammenstillede Sporegjemmer danner et Sporeleie
(\textit{apothecium}). Hvad Dannelsen af et saadant Sporeleie angaaer,
er Organisationen altsaa endnu saa svag, at den næsten taber sig i
Aggregation. At en saadan organisk Aggregation imidlertid er noget Mere
end en blot mechanisk, kommer saa efterhaanden tilsyne i hoiere
Udførelse
% the note runs on over the page break here and is finished at the head
% of the footnote block of printed p. 295
af det kryptogame Anlæg; men hverken Mossernes eller Brægnernes
Sporehuse naae Bladformen.
}, dets indre Bygning beroer dog kun paa
% --- p. 295 ---
\setcounter{footnote}{0}%
Celler og Aggregationer af Celler, hvorimod det laveste Frøgjemme altid
som Frugtblad oprindelig er et Blad. Det laveste phanerogame Frugtblad
staaer altsaa dog i Kraft af Grundformen hoiere end det hoieste
kryptogame Sporehuus\footnote{% printed mark: *)
For at see, hvad der paa et hoiere Grundlag kan naaes under
fremadskridende Dannelser og Omdannelser, behøve vi blot at kaste et
Blik paa de mere udviklede Frøgjemmer. I Tamarindens Frugt, der dog
endnu kun er dannet af et eneste Blad, er Frøgjemmet differentieret i
indre og ydre Overhud; det mellemliggende Bladkjød (Diachym) er udviklet
til en Frugtmarv (\textit{pulpa}), hvis Celler bugne af rigt formede
Stofsubstanser. Imod en saadan Articulation er den kryptogame Herlighed
et Forsvindende. Betragte vi et af flere Bladelementer sammensat
Frøgjemme, Appelsinens Frugt t. Ex., saa finde vi her en endnu finere,
endnu rigere Articulation, en Overhud, hvori et stort Antal ætherisk
Olie indeholdende Kjertler fremtræde, en læderagtig Masse (Diploe,
Diachym), der har afgivet
% `afgivet': the final `t' is a broken/half-inked sort and prints as two
% detached dots, so both OCR witnesses read `afgive:'. Set as the word.
et nærende Grundlag for den Besætning af tæt sammenstablede, kjødede,
Appelsinsaften indeholdende Sække, der udgaae fra Perikarpiets
Inderflade o. s. v.
}.

Til en nærmere Bestemmelse af Modsætningen imellem det Lavere og det
Hoiere er det fremdeles nødvendigt at opfatte Grundformen
\emph{teleologisk}. Vi sætte Kryptogamens Sporedannelse mod
Phanerogamens Kiimdannelse. Skulde vi nu see bort fra Grundformernes
Forskjellighed og tage Formaalet abstract, vilde vi snart blive færdige
med Modsætningen. Sporedannelsen og Kiimdannelsen ere jo begge anlagte
paa eet og samme Formaal:
Forplantningen\footnote{% printed mark: **) --- second note on p. 295
Hvad Modsætningen mellem Antheridiernes Sædfiim (Spermatozoer) og
Frugtlegemets (Archegoniets) Centralcelle er hos Kryptogamerne, er
Modsætningen mellem Støvsækkens mandlige Sædstof (Pollen, navnlig det i
Pollinarrøret, det forlængede Pollenkorn, Organiserede) og den i Æggets
Kiimsæk indesluttede Kiimcelle hos Phanerogamerne.
}. At
% --- p. 296 ---
\setcounter{footnote}{0}%
nu Forplantningen, Formaalet selv, er noget Væsentligt, forstaaer sig;
men --- siger man maaskee --- „Forplantningens Midler, Maaden, hvorpaa
Forplantningen skeer, synes jo at være noget aldeles Uvæsentligt; eller
hvorfor skulde en Forplantning ved Spore ikke være ligesaa god, som en
Forplantning ved Kiim?“ Denne Synsmaade er ideeløs. Ja, dersom det
organiske Formaal virkelig lod sig saaledes adskille fra de tilsvarende
Midler, maatte vi jo vistnok blive staaende ved den Tautologie, at
Forplantning er Forplantning; men da „det organiske Naturprodukt“, for
ved disse Kants Ord at minde om det tidligere Udviklede, just „er et
saadant, i hvilket enhver af dets Dele er Øiemed og tillige Middel eller
Redskab“, saa bliver det almene Formaal for hver Grundform
individualiseret i Midlerne, saa bliver det Teleologiske bestemt ved det
Typiske, saa bliver Forplantning forskjellig fra
Forplantning\footnote{% printed mark: *)
Den Forskjel imellem Forplantning og Forplantning, hvorom her tales,
forudsætter naturligviis, at Forplantningens Eiendommelighed fastholdes,
forsaavidt „den overalt i den levende Natur udtalte Forplantningslov“
forudsætter „Kjønnets Fordeling paa to og Modsætning i to forskjellige
Væsener.“ --- „Vel angives det almindeligt i enhver Plantelære, at
Hermaphroditismen er den normale, ɔ: hyppigste og naturligste, Form for
Væxternes Forplantningsforhold; men nærmere beseet er dette dog vel
snarest en Talemaade, nedarvet fra tidligere Tider, da Læren om
Planternes saakaldte „Metamorphose“ endnu ikke havde yttret sin rensende
og velgjørende Indflydelse paa Opfattelsen af Planten. Det Hele er her
afhængigt af, hvad der skal betragtes som det egentlige Individ i
Planten, som det egentlige Væxt-Ene.“ J. J. Sm. Steenstrup.
Undersøgelser over Hermaphroditismens Tilværelse i Naturen. Kjøbhvn
1845.
% First occurrence of the id-est mark `ɔ:' in the book (600 dpi: an
% open-o followed by a colon). The book has written `d. e.' hitherto.
% NOTE FOR THE PREAMBLE: U+0254 has no T1 slot, so pdflatex stops with
% `Unicode character ɔ (U+0254) not set up for use with LaTeX' until a
% mapping is declared, e.g.
%   \DeclareUnicodeCharacter{0254}{\reflectbox{c}}   % needs graphicx
% With that line added, this fragment compiles clean.
}.

Det er, vel at mærke, ikke en Betragtning af Phænomenerne i deres
Umiddelbarhed, men en til de betragtede Phænomener svarende
Ideeerkjendelse, hvorom Talen er. At opfatte Modsætningen mellem et
Hoiere og et Lavere som en
% --- p. 297 ---
\setcounter{footnote}{0}%
Modsætning mellem det, der tilfredsstiller Ideen, og det, der ikke
tilfredsstiller den, forraader vistnok en betænkelig Misforstaaelse; thi
Ideen tilfredsstilles ligesaavel ved det Lavere som ved det Hoiere, just
fordi den i umiddelbar Forstand hverken tilfredsstilles ved dette eller
ved hiint, men først i Dialektiken af Modsætningens Led finder sin
væsentlige Tilfredsstillelse. Men, naar Tilfredsstillelsesdialektiken nu
viser sig deri, at det Hoiere i sin Form just tilfredsstiller ved dets
charakteristiske Modsætning til det Lavere, og dette Lavere omvendt ved
dets charakteristiske Modsætning til det Hoiere; naar begge ved at være
lige for Ideen --- i samme Betydning som naar vi sige, at den fattige og
den rige Borger begge ere lige for Loven --- have denne væsentlige
Lighed reflecteret i en for Ideen ligesaa afgjørende Ulighed --- som
naar vi sige, at Uligheden mellem Rige og Fattige er retfærdiggjort ved
den Lov, der hævder Eiendomsretten ---; naar det altsaa tydelig indsees,
at det Laveres Lighed med det Hoiere i Ideen selv maa være begrundet paa
Uligheden: saa bliver det en, idetmindste foreløbig Opgave, at klare sig
de i Forplantningsideen indeholdte Fordringer som ved et
Tankeexperiment, for dernæst at prøve, hvorvidt de givne Phænomener
stemme overeens med det i et saadant Tankeexperiment Udviklede.

Tør vi nu forudsætte, at den phanerogame Grundform indenfor
Planteverdenen hører til en hoiere Region end den kryptogame, da tør vi
ogsaa forudsætte, at den phanerogame Forplantning maa forholde sig til
den kryptogame, som det Hoiere til det Lavere. Sætter Forplantningens
Idee en væsentlig Adskillelse og dog inderlig Forening af Moderplanten
og dens Afkom, saa vil --- det er et Tankeexperiment ---
Grundfordringens hoiere Tilfredsstillelse jo typisk beroe derpaa, at
Befrugtningen skeer direct, at Embryet faaer en egen af Moderlegemet
særlig tilberedt Næring, og at Embryonaldannelsen, fra Begyndelsen af
betegnet ved et be\-%
% --- p. 298 ---
\setcounter{footnote}{0}%
stemt Midtpunkt, giver sig tilkjende i primitive Articulationer.

Har nu denne Tankegang Realitet, saa kan det neppe være tvivlsomt, at
Sporen maa forholde sig til Kimen, som det Lavere til det Hoiere. Medens
Befrugtningen af Phanerogamens Æg skeer direct, saa at den befrugtede
Kiimcelle selv afgiver det nye Centrum, Embryonaldannelsens udtrykkelige
Midtpunkt, skeer den kryptogame Befrugtning derimod indirect, forsaavidt
det nemlig ikke er Sporen\footnote{% printed mark: *)
% The whole note is German, set in Fraktur like the surrounding text and
% therefore NOT italicised; the print gives it no quotation marks.
Die Sporen sind einer vorgängigen Befruchtung nicht bedürftig; die Samen
dagegen sind ohne eine solche nicht entwicklungsfähig; und wenn sie auch
äußerlich die gewöhnliche Gestalt und Ausbildung erreichen, so fehlt
ihnen doch der Embryo, auf dessen Gegenwart alle spätere Entwicklung
beruht. Man hat deshalb denselben Gegensatz dieser beiden
Pflanzengruppen durch die Unterscheidung von geschlechtslosen,
Kryptogamen, und Geschlechtspflanzen, Phanerogamen, bezeichnet, indem
man den in der Thierwelt bekannten Begriff der Sexualität auf die
Pflanzen übertrug . . . Diese Anschauungsweise hat in neuerer Zeit viele
Ungunst erfahren . . . . Es ist nahmentlich die eigenthümliche
Befruchtungstheorie Schleidens, welche der Parallelisirung thierischer
und pflanzlicher Verhältniße entgegentrat. Nach Schleiden ist es das
Pollenkorn, der bisher für das männliche Organ gehaltene Theil, aus
welchem der Embryo des Samens entsteht. Gleich einer Spore oder einer
andern einzelnen Fortpflanzungszelle ist der Pollenschlauch die
entwicklungsfähige, und einer Befruchtung im Sinne der thierischen
Zeugung unbedürftige Grundlage der spätern neuen Pflanze. Aber seine
Lebensfähigkeit ist nicht von der Art, daß er unmittelbar der äußern
Einflüßen ausgesetzt, sich zu einem neuen Organismus entwickeln könnte;
er bedarf es vielmehr, zuerst in eine bestimmt angeordnete organische
Umgebung aufgenommen zu werden, die ihm die zu den ersten Schritten
seiner Bildung nothwendigen Bedingungen gesetzlich und abgemeßen
gewähre. So begiebt sich denn diese keimfähige Grundlage der Pflanze in
jene für weiblich gehaltene Organe des Fruchtknotens, um hier, wie an
eine Brutstelle, so lange gepflegt zu werden, bis sie die nöthige Masse
und Festigkeit innerer Verhältniße gewonnen
% the note runs on over the page break here and is finished at the head
% of the footnote block of printed p. 299
hat, um in unmittelbarer Wechselwirkung mit der Außenwelt sich weiter zu
entfalten. Ich kann aber das Bestreben, die sexualen Verhältniße der
beiden Reiche zu parallelisiren, nicht für eine grundlose abergläubische
Richtung der Phantasie ansehen . . . . Es versteht sich übrigens von
selbst, daß diese methodologisch richtige Vermuthung analoger
Einrichtungen bescheiden genug auftreten muß. Lotze. Allgem. Physiol.
S. 558---61 o. flgd.

Den sexuelle Modsætning er --- ogsaa for Planteorganismerne --- af
gjennemgaaende Betydning. At Antagelsen af en „Coelebogyne“ vel maa
beroe paa en Illusion, synes allerede nu at være saa temmelig afgjort.
„Da wir nun die Coelebogyne noch viel zu wenig kennen, auch bisher für
sie nur ein sehr beschränktes Material zur Beobachtung vorhanden war, so
ist die Frage der Parthenogenesis im Pflanzenreiche noch keineswegs
entschieden und sprechen die von Regel mitgetheilten Beobachtungen nicht
sehr zu ihrem Vortheil.“ H. Schacht. Lehrbuch der Anatomie und
Physiologie der Gewächse. Berlin 1856. 2ter Bd. S. 596. --- Hvad derimod
Parthenogenesien i Dyreverdenen angaaer, jfr. Th. E. Siebold.
\emph{Wahre Parthenogenesis} bei Schmetterlingen und Bienen. Leipzig
1856.
}, men Archegoniets
% --- p. 299 ---
\setcounter{footnote}{0}%
Centralcelle, der befrugtes umiddelbart\footnote{% printed mark: *)
Den befrugtede Centralcelle udvikler sig saa, navnlig hos Brægnen, til
et Sporangium, i hvilket de spiredygtige Sporer fremkomme. „Forkimen,
der udvikler sig af Sporen, er et hjerteformigt, bladagtigt Legeme, der
ved Hæftetraade er fæstet til Jorden; paa Forkimens Underside findes
mange Antheridier og enkelte Frugtlegemer; de sidste befrugtes af
Sædtraade, som udgaae fra Antheridierne, og som Følge deraf udvoxer der
af Frugtlegemet en Brægne, hvis nederste Deel bliver til Mellemstokken,
og udsender Birødder.“ Chr. Vaupell. Planter. Naturhist. S. 133.

Sporen stilles udtrykkelig imod Frøet, fordi det typiske Synspunkt er,
at det nye System af Individuationer tager sin Begyndelse med Sporen.
}. I den directe Befrugtning er altsaa Kimens større Primitivitet
fremfor Sporen allerede begrundet. Denne Forskjel træder nu tydeligere
frem, naar Spørgsmaalet er om Ernæringen. For et Uselvstændigt, en Væxt,
der snarere er en jævn Fortsættelse af et allerede Givet end noget
eiendommeligt Nyt, indføres der naturligviis ingen særegen
Huusholdning; skal en Huusholdning
% --- p. 300 ---
\setcounter{footnote}{0}%
indføres, en egen Næring tilberedes, maa der lægges særlig Vægt paa
Selvstændigheden af det, som derved skal ernæres. Her have vi igjen
Modsætningen imellem det Kryptogame og det Phanerogame. Det kryptogame
Frugtlegemes Centralcelle er --- hos de hoiere Former ---
% "hoiere" with plain o as printed (calibrated at 600 dpi against the slashed ø
% of "øvrige"/"foreløbige" two lines below); see the note on p. 302
henviist til at suge sin Næring af hele den øvrige Plante (hos Brægnerne af
Forkimen), hvorimod Embryonalcellen hos Phanerogamerne begynder med at
ernære sig af Modercellens Indhold (\textit{protoplasma}) og af foreløbige
Dannelser indenfor Modercellen; er Modercellen med samt dens Indhold
fortæret, kan ogsaa, hvad der navnlig hos Frø uden Frøhvide er det
Sædvanlige, Kjernens øvrige Masse (\textit{nucleus ovuli}) tages i Beslag.
Den Indvending, at den kryptogame Moderorganisme jo ligesaavel træder i
Modsætning til sit Afkom som den phanerogame, overseer, at her ikke spørges
om Modsætningen \textit{in abstracto}, men om Modsætningens typiske Udtryk.
Den Indvending, at en Næring, der „suges af hele Planten“, jo til sit Brug
kan være ligesaa kraftig og sund, som en Næring, der særlig tilberedes enten
i „Modercellen“ eller i „Kjernens øvrige Masse“ til sit Brug, en Indvending,
der altsaa begrundes paa den Tautologie, at en Næring altid er en Næring,
ligerviis som et Afkom altid er et Afkom, maa i en Sammenhæng, hvor det
netop gjælder en typisk Udtalelse af de for en selvstændigere
Embryonaldannelse charakteristiske Bestemmelser, afvises som ideeløs og
ubeføiet.

Hvad Dannelsens Udgangspunkt, det primordiale Midtpunkt, og
Articulationens dertil svarende Primitivitet angaaer, er Phanerogamens
Fortrin fremfor Kryptogamen øiensynlig. Sporen, der forøvrigt kan være
enten en enkelt Celle eller et Celleaggregat, forraader sin svagt udviklede,
indifferente Tilstand derved, at den, saa langt Iagttagelserne gaae, aldrig
udviser nogen Polaritet; hvorimod Phanerogamens Kiim, selv om den, som
f.~Ex. hos \textit{Melampyrum}, kun bestaaer af to Celler, udtrykker
Articulationen i
% --- p. 301 ---
\setcounter{footnote}{0}%
en polær Modsætning mellem et Opad og et Nedad\footnote{Den Celle, som har
været vendt imod Kiimhullet (\textit{micropyle}), afgiver nemlig Rodenden,
medens den fra Kiimhullet bortvendte Celle afgiver den opadgaaende Deel af
Planten.}. Kryptogamen begynder i det Eensformige; Phanerogamen indgaaer fra
Begyndelsen af i en Fleerhed af Formbestemmelser. Idet den unge kryptogame
Plante som Forkiim (\textit{prothallium}) kommer frem af Sporen, er den
hverken Andet eller Mere end den Celle, den forhen var, kun større og ændret
i Omrids; hos de fleste Phanerogamer er Kimens Articulation derimod saa
bestemt udtalt i Axe og Blad, ja i Blade af forskjellig Orden: Kiimblade,
Primordialblade, at man endog af den Grund har kaldt Kimen selv den lille
Plante (\textit{plantula}).

Er den kryptogame Forplantning i sit Væsen saa forskjellig fra den
phanerogame, da maae jo ogsaa de Organismer, der ere af et saa forskjelligt
Udspring, naar de betragtes i deres fuldstændige Udvikling, i alle
Hovedtræk bære denne Grundforskjel til Skue. Man sætte t.~Ex. Lavarter mod
Græsarter, Brægner mod Acacier, og den typisk-teleologiske Modsætning vil,
fra Ideens Synspunkt, være iøinefaldende. Hvilken Modsætning, for at blive
ved de høiere Repræsentonter,
% PRINTER'S DEFECT: the print reads "Repræsentonter" (an o sort for a); the
% round o is unmistakable at 600 dpi beside the two a's of "danner" on the
% same line. Transcribed as printed.
danner ikke Brægnen mod Acacien! Roden, Stammen, Bladet, Habitus, hele det
Ydre, den indre Structur af hver enkelt Deel, Alt er i denne Sammenstilling
paa hver Side mærket af Grundformens typiske Eiendommelighed. Da Sporen
mangler den polære Modsætning mellem Opad og Nedad, har Brægnen naturligviis
ingen Hovedrod og kan ingen have; hos Acacien derimod er her, lige fra
Spiringens forste Øieblik af,
% "forste" with plain o, "Øieblik" with slashed Ø, both as printed
en skarpt udtalt Adskillelse imellem den oprindelige Rod og det, der skal
blive til Stængel. Brægnens Stamme, ofte mangeaarig, høi og slank, viser en
ikke ringe Differentiering i det Anatomiske (mægtige Karbundter omgive
Marven
% --- p. 302 ---
\setcounter{footnote}{0}%
o.~s.~v.); men da Brægnestammen --- naar vi see paa det Almindelige --- ingen
over Jorden udskydende Grene frembringer, da Endeknoppen som Følge heraf
alene afgiver de levende Løv, maa Kryptogamernes hoitstillede Repræsentant
% "hoitstillede" with plain o, against the plainly slashed ø of "høieste" and
% "høiere" further down this same page. The compositor mixes hoi- and høi-;
% every instance in this batch was read off the image at 600 dpi.
dog i Habitus mangle den for Planteskikkelsen saa charakteristiske
Forgrening. En saadan Forgrening, typisk og rig, udgaaer derimod fra
Acaciens træagtige, men dog rigt articulerede Stamme, en Stamme med de
høieste anatomiske Reqvisiter: Overhud, Bark, Bast, Splint, Kjerne og Marv.
Af en Kryptogam at være, er Brægnen i Henseende til Bladdannelsen rigelig
udstyret; det forste fra Forkimen udskydende Løv er rigtignok ubetydeligt,
men hvert følgende tiltager i Styrke, og som oftest i Indskjæring, indtil
den høieste vegetative Fylde er naaet; at disse Løv have Spalteaabninger,
altsaa særlig udviklede Respirationsbetingelser, er et umiskjendeligt Tegn
paa den høiere Rang, Brægnen indtager blandt „de Blomsterløse“. Men
\emph{blomsterløs} er Brægnen, det ligger i dens kryptogame Natur: dens
høieste Frembringelser ere de sporebærende Løv (\textit{frondes fertiles}).
Paa det phanerogame Trin naaer Bladdannelsen derimod en ganske anden, baade
fleersidigere og høiere Udvikling. Hos Acacien er Bladet udtalt i alle dets
Dele: Skededelen er ofte sondret endog i flere Tommer lange, mangeaarige
Stipler, Pladen virkelig sammensat og paa saa mangfoldige Maader, som ellers
sjældent i Planteriget; ja Formbevægeligheden er her saa stor, at f.~Ex. de
nyhollandske Acacier ofte ikke engang have udviklede Bladplader, men
Phyllodier, dannede, hvert især, af en Bladstilk (\textit{petiolus}), medens
de rigeste og finest articulerede Bladformer forekomme t.~Ex. hos
\textit{Mimosa pudica}. At Acacien, ifølge den phanerogame
Forplantningsmaade, maa foruden vegetative Blade ogsaa have Blomsterblade,
er en Selvfølge; den Alsidighed, hvormed Blomstdannelsen her foregaaer,
charakteriserer imidlertid særligt det Høidepunkt, hvorpaa Acacien staaer.
Efter mange Svingninger gjennem det Uregelmæssige
% --- p. 303 ---
\setcounter{footnote}{0}%
(\textit{Papilionaceæ}) naaes i Blomstdannelsen her det Regelmæssige; den
regelmæssige \emph{undersædige} Blomst --- just den Form, hvori
Blomstindividerne træde tydeligst frem --- har et meget stort, men ubestemt
Antal Støvdragere; ja, den ellers eenfolde Frugt kan i nogle Former
(\textit{Affonsia}) endog være fleerfold: Alsidigheden er altsaa her i alle
Dele paa sit Hoieste\footnote{Indenfor denne Familie ere da ogsaa de tvende
Vitalphænomener: Søvnen og Sensitiviteten (\textit{Mimosa sensitiva,
Smithia sensitiva}) meget fremtrædende.}.
% "Hoieste" with plain o, verified at 600 dpi; "(Affonsia)" so spelt in the
% antiqua, with an ff ligature.

„Men“ --- indvender man fra den empiriske Forsknings kritiske Standpunkt ---
„det er dog en mislig Sag med disse naturphilosophiske Ideebelysninger;
inden man ret veed deraf, har den videnskabelige Maalestok forvandlet sig
til en blot æsthetisk, og Phantasiens poetiske Idealer sat Farve paa de
prosaiske Kjendsgjerninger“\footnote{Das Unternehmen, die Gattungen der
Thiere und Pflanzen in eine aufsteigende Stufenreihe zu ordnen, setzt nicht
nur die Kenntniß aller für beide Reiche charakteristischen Merkmale, sondern
auch eine mögliche Beurtheilung des Werthes voraus, den jedes einzelne
besitzt. Die erste dieser Aufgaben könnte vielleicht durch die aufgezählten
mannigfachen Unterschiede zwischen Thieren und Pflanzen erledigt scheinen,
deren jeder zugleich einen positiven Beitrag zu dem Bilde dessen ausmacht,
das er von dem andern trennt; die zweite ist unläugbar schwieriger, und die
Entscheidung, welche Züge des thierischen Daseins z.~B. in einer
Stufenleiter der Vollkommenheit immer deutlicher hervortreten müßten, welche
andere im Gegentheil durch ihr Zurückbleiben oder Verschwinden zur
Vervollkommnung beitragen, bleibt meist einer ästhetisch erregten Phantasie
überlassen, die in der Wahl ihrer Gesichtspunkte häufig fehlgreift. Diese
Uebelstände werden stets vorhanden sein, so lange die Begriffe von Pflanze
und Thier nur aus der Beobachtung gewonnen sind, die immer nur einen
gewissen Thatbestand der Verbindung und Gruppirung von anschaulichen
Elementen nachweist, ohne uns die Form oder die Gleichung, durch welche
Werth und Beziehung derselben unter einander bestimmt wird,
% the printed page break of this note falls here, between p. 303 and p. 304
neben der anschaulichen Erscheinung noch einmal ausdrücklich zu
wiederholen.“ Lotze. Allgem. Physiol. S.~507--8.}.
% PRINTER'S PRACTICE: the German quotation in this note carries a closing
% high mark after "wiederholen." but no opening low mark; the note simply
% begins "Das Unternehmen". Hence one unmatched closing mark in this fragment.
Før vi imidlertid indlade os
% --- p. 304 ---
\setcounter{footnote}{0}%
paa en Besvarelse af herunder hørende Spørgsmaal, ville vi, for at vinde
behørige Forudsætninger, benytte de for en Bedømmelse af Grundformerne
væsentlige Oplysninger, der endnu tilbyde sig fra en anden Side.

ββ) \emph{Metamorphoser.} Var det muligt, at der paa et Acacietræ kunde
findes Brægneløv, og paa en Brægnestængel Acacieblade, var det muligt, at
Acacien kunde bære afvexlende baade Blomster og Sporer, og Brægnen igjen
afvexlende Sporer og Blomster; ja, dersom Sligt var muligt, saa var jo
vistnok den teleologiske Grundform en Phantasie, men saa var ogsaa Botaniken
som Videnskab en Phantasie. At de forskjellige til en bestemt Organisme
hørende Dannelser fremgaae af eet Anlæg og udgjøre et teleologisk System,
sees af Metamorphosen.

I sit berømte Skrift: „Die Metamorphose der Pflanzen“, har Goethe udtalt den
i Botaniken epochegjørende Grundtanke, at nemlig alle en phanerogam Plantes
forskjellige Dele, ligefra Cotyledonen til Støvveien, fremkomme derved, at
eet og samme Element paa de forskjellige Hoider af Axen iforer sig
forskjellig Skikkelse og saaledes fremstiller sig som forvandlet.
% "Hoider" and "iforer" with plain o, verified at 600 dpi
Ved mangfoldige Sammenligninger, ved at sammenligne Blade paa samme Plante i
forskjellig Hoide, ved at sammenligne de paa begge Sider liggende Blade med
det mellemliggende Blad af monstrøs Beskaffenhed (Anamorphose) og endelig
ved at sammenligne Blade af samme Orden paa forskjellige Trin i Planteriget,
gjorde det intuitive Genie den Opdagelse, at, hvad der paa et Stadium var
Frøblad (Cotyledon), var paa et andet Løvblad, paa et tredie Dækblad
(Braktee), paa et fjerde Bægerblad, paa et femte Kronblad, paa et sjette
Støvblad, paa et syvende Frugtblad. --- Fordres endnu et yderligere Beviis
for, at Plantens Dele ikke fremkomme ved endelig
% --- p. 305 ---
\setcounter{footnote}{0}%
Sammenstykning, men udvikles indenfra af eet sig differentierende Anlæg, saa
at Organernes Forskjellighed og indbyrdes bestemte Forhold væsentlig grunder
sig paa Paradigmets teleologiske Fordringer: da behøve vi kun at minde om
den Maade, hvorpaa navnlig Schleiden har prøvet og gjennemført det af Goethe
Anskuede. Istedetfor blot at sammenligne Phænomenerne i deres umiddelbare
Tilværen, spørger Schleiden tillige om deres Vorden, gaaer tilbage til det
Oprindelige, sammenligner t.~Ex. en Cotyledon i dens første Tilsyneladelse
med de første Tilsyneladelser af andre Blade paa samme Plante, og finder
saa, at alle Bladene fra Begyndelsen af ligne hverandre, at de alle
oprindelig ere eensdannede, hver især „en vorteformig Fremragen nedenfor
Axespidsen\footnote{Det maa udtrykkelig fremhæves, at der i Botaniken ikke
er Tale om metamorphoserede Blade i samme Betydning, som naar der i
Geognosien tales om metamorphoserede Stene. At Marmor og Alabast ere dannede
ved Metamorphose, forudsætter, at Marmoren har været simpel Kalksteen, før
den blev Marmor, og Alabasten Gips, før den blev Alabast; men en Støvtraad
er ikke et metamorphoseret Kronblad i den Betydning, at den virkelig skulde
have været Kronblad, før den blev Støvtraad. „So lange man jedoch Stengel
und Blatt nur nach ihrem weniger morphologischen als vielmehr descriptiv
geometrischen Begriffe unterscheidet, ist es schwer zu sagen, ob in der
Ansicht, welche verschiedene Theile der Blüthe als Umwandlungen des Blattes
betrachtet, mehr Tautologie oder mehr unklare Mystik liegt.“ Lotze. Allgem.
Physiol. S.~535.}.“

Forsaavidt det oprindelig Ligeartede hvert paa sit Sted og efter sin
Bestemmelse forvandles til et Uligeartet, hersker Idee-Aarsagen; forsaavidt
Idee-Aarsagen, det teleologiske Princip, hersker, er en Parallelisme imellem
Organernes Sammenhæng i den enkelte Plante og Typernes Sammenhæng i
Planterigets „naturlige System“ umiskjendelig. Her gives da ikke blot en
Organernes, men ogsaa en Grundformernes
% --- p. 306 ---
\setcounter{footnote}{0}%
Metamorphose; kalde vi hiin den \emph{reelle}, maae vi kalde denne den
\emph{ideelle}: begge have i Grunden samme Oprindelighed, samme væsentlige
Gyldighed. Betegner den reelle Metamorphose, at de forskjellige Organer paa
samme Organisme ere Omdannelser af eet og samme Grundelement, saa betegner
den ideelle Metamorphose, at de særskilte Typer alle ere Omdannelser og just
derved Gjentagelser af een og samme Grundtype. Den Forskjel, at der imellem
Organerne paa samme Organisme finder en materiel Forbindelse Sted, medens
der imellem de forskjellige Typer ingen saadan tilsvarende physisk
Vexelvirkning lader sig paavise, tør vistnok ikke oversees; men overalt hvor
Idee-Aarsagen gjælder, har den ideelle Sammenhæng større Realitet end den
blot materielle: paa den ideelle Sammenhængs oprindelige Realitet er
Parallelismen imellem den reelle og den ideelle Metamorphose begrundet. En
sand Indsigt i denne Parallelisme er betinget af sand Indsigt i
Tingsaarsagernes, de physiske Mellemaarsagers, dialektiske Eenhed med
Grundaarsagen.

Den speculative Naturphilosophie, der altfor let abstraherer fra de physiske
Mellemaarsager, tager Ideens Navn forfængelig og misbruger Parallelismen,
idet den, for at følge den opadskridende Bevægelse, overlader sig til en
vilkaarlig æsthetiserende Anskuelses Veiledning; den kritiske Forskning, der
fornegter Idee-Aarsagen og betragter Parallelismen blot som et Værk af en
sammenlignende Abstraction, tager Kjendsgjerningens Navn forfængelig og kan,
trods sine mange Oplysninger om „Dannelsesmaterialet“, om „Kulsyre“ og
„Ammoniak“, om „Cellulose“ og „Æggehvide“ og „Materiens plastiske Kraft“ og
„de physiske Betingelser“, o.~s.~v., neppe engang finde Rede i sine egne
Hypotheser, endsige da i de teleologiske Former.

„Men“ --- indvender man fra den kritiske Forsknings empiriske Standpunkt ---
„hvad nytter det vel at beraabe sig paa Idee-Aarsagen og den teleologiske
Grundform, naar man
% --- p. 307 ---
\setcounter{footnote}{0}%
i Virkeligheden dog hverken kan eftervise, hvorledes de physiske Aarsager i
det Enkelte hænge sammen med Grundaarsagen, ei heller forklare
Charakteermærkernes indbyrdes nødvendige Sammenhæng af det i Grundformen
Teleologiske? At der imellem Brægnens Rod, Stamme og Blade er en væsentlig
Sammenhæng, vil Ingen negte; den organiske Vexelvirkning medfører jo netop,
at en saadan Stamme maa have saadanne Blade, og omvendt. At
Brægneorganismens Eiendommelighed lader sig henføre til de ved dens
Forplantning samvirkende Aarsager, maa ligeledes indrømmes; men hvad skal nu
det Teleologiske forklare, hvad Brug har man for Idee-Aarsagen? At Brægnen
ikke kan have Blomster, naar Forplantningen skeer ved Sporer, og ikke have
Hovedrod, naar Sporen mangler Polaritet, er en Selvfølge; men af hvilken
Teleologie skulde det vel kunne udledes, at Brægnerne i vort Klima ere
urteagtige Planter, medens de træagtige Stammer, hvormed de ofte optræde i
tropiske Egne, have en saa fremherskende Lighed med Palmestammer, og det
uagtet Palmerne høre til Blomsterplanterne? Eller gives der maaskee nogen
særegen Idee-Aarsag, hvoraf man kan forklare, hvorfor Brægnens Blad absolut
skal have sit yngste Parti i Spidsen, medens derimod, hvad jo Schleiden har
paaviist, Planternes Blade ellers ere morphologisk begrændsede derved, at de
i Modsætning til Axen, det morphologisk Ubegrændsede, typisk have det ældste
Parti i Spidsen? Kan man ved nogen Ideebelysning opdage Grunden til, at
Befrugtningen for Brægnernes Vedkommende, just fordi her skal dannes en
saadan kryptogam Organisme, maa foregaae paa Forkimen, medens den hos
Mosserne foregaaer paa Planten selv, uagtet Mosserne jo ogsaa ere
kryptogame?“

„At Karbundterne i den phanerogame Stængel ere snart spredte og mangle
Kambialring, snart derimod samlede, krandsstillede og i Vexelforhold til
Kambialringen, er en Kjendsgjerning; men hvor finde vi saa den Idee, som
siger os, hvoraf det
% --- p. 308 ---
\setcounter{footnote}{0}%
egentlig kommer, at det just er de tofrøbladede Planter, der have
Kambialringe, ikke de eenfrøbladede? Saalænge man tager Sagen i al
Almindelighed og holder den i det Svævende, kan man jo vistnok meget let
hævde Ideens Gyldighed; men ved at gaae i det Enkelte og see paa
Kjendsgjerningerne støder man overalt paa eet af to Tilfælde: enten lader
Sammenhængen imellem en Grundforms typiske Mærker sig forklare af
\emph{physiske} --- physiologiske, anatomiske, morphologiske ---
Aarsager\footnote{Paa hvormange Maader alene Tyngdekraften er istand til at
bevirke charakteristiske Ændringer i Planteorganismen, bevise de nyeste
herhen hørende Undersøgelser. „Bei Untersuchung der Mechanik der
geocentrischen Krümmungen pflanzlicher Organe ist vor Allem der Umstand
scharf im Auge zu behalten, daß von gewißen verschiedenen Organen völlig
gleichen anatomischen Baues die einen bei Einwirkung der Schwerkraft
aufwärts, die anderen abwärts sich krümmen. Sichtbare Unterschiede der
Structur eines Stengels und einer Wurzel von \textit{Nitella},
eines auf- und eines abwärts wachsenden Sproßes von
\textit{Equisetum palustre} sind
nicht vorhanden; die Differenzen der Anordnung und Vertheilung der Gewebe in
Stengel und Wurzel derselben Pflanze sind häufig verschwindend klein. Wohl
aber zeigen Stengel und Wurzel in ihren geocentrischer Krümmung fähigen
Theilen ein sehr verschiedenes Maaß von Differenzen der \emph{Spannung} der
einzelnen Gewebsmassen.“ W.~Hofmeister. Ueber die durch die Schwerkraft
bestimmten Richtungen von Pflanzentheilen. Jahrbücher für wissenschaftliche
Botanik von Pringsheim. 3ter Bd. 1stes Heft. Berlin 1861. S.~77 o.~flgd.},
og da er Idee-Aarsagen ganske overflødig, eller ogsaa lade slige physiske
Aarsager sig ikke paavise, og da er det kun tom Tale, naar man, for at
skjule sin Uvidenhed, paakalder det Teleologiske og beraaber sig paa en
intetsigende Idee-Aarsag.“

Det er den gamle indgroede Misforstaaelse, at en Idee-Aarsag Intet er og
Intet betyder, naar den ikke er til som et physisk Element (og altsaa heller
ikke kan skimtes under Mikroskopet), som her har udtalt sig i nye Vendinger.
Saa\-%
% --- p. 309 ---
\setcounter{footnote}{0}%
længe denne Misforstaaelse hersker, vil endog den skarpsindigste Forskning
være hildet i den Modsigelse, at skulle i eet Aandedræt anerkjende Naturen
og dog fornegte Naturen: Naturen er Idee, dens „Dannelsesdrift“, dens
„plastiske Kraft“, dens „typiske Former“ ere Ideens Yttringer,
Idee-Aarsagens Virkninger.

Betragte vi, for just at have noget Bestemt for Oie, t.~Ex. den Diagnose,
% "for Oie" with plain O as printed
hvorved Botanikerne adskille den røde og den blaa Anagallis\footnote{Jvf.
J.~Lange, Haandb. i den danske Flora. Kbhvn. 1856--59. S.~165.}, saa viser
sig her den meest paafaldende Lighed: „Stænglen udstrakt“, „Bladene
modsatte“, „Bægerfligene sylformede“ o.~s.~v., forenet med en ligesaa
paafaldende Ulighed; thi hos den første (\textit{A. arvensis}) er „Kronen
rød“, hos den sidste (\textit{A. coerulea}) er „Kronen blaa“. Det er i det
Hele en Kjendsgjerning, som under mange Ændringer gjentager sig i
Planteriget, at en iøinefaldende Ulighed i nogle Dele findes forenet med en
mærkværdig Lighed i andre\footnote{En paafaldende indbyrdes Lighed have,
hvad Græsarterne angaaer, \textit{Poa nemoralis} med \textit{Poa fertilis};
men med denne Lighed: „Straaet trindt, jævnt“, --- „Toppen med rue, efter
Blomstringen oprette Grene“, --- „\emph{Smaaaxene} 2--5 blomstrede“, ---
forener sig en ligesaa bestemt Ulighed, idet „Bladskederne“ hos hiin ere
„kortere end de tilsvarende Ledstykker, den øverste kortere end sin
Bladplade; Skedehinden kort, afstumpet“, --- hos denne: „Bladskederne lige
lange med Ledstykkerne, den øverste af lige Længde med sin Plade;
Skedehinden forlænget“; --- hos hiin er „Toppen smal“, hos denne „udbredt“;
hos hiin ere „Inderavnerne spidse“, hos denne „budte“. See J.~Lange. Anfr.
Skr. S.~79--80.

Sammenligne vi det \emph{Almindelige Jordbær} (\textit{Fragaria vesca}) med
det saakaldte Bakke-Jordbær (\textit{F. collina}), er Ligheden mellem begge
i alle Dele saa gjennemgaaende, at det sidste kun adskiller sig fra det
første (Jvfr. J.~Lange. S.~353--54) derved, at dets „Bæger ved Modenhed“ er
„tiltrykt til den mørk-rosenrøde, søde Frugt“.}; saa\-%
% --- p. 310 ---
\setcounter{footnote}{0}%
ledes kunne mangfoldige Korsblomstrede udvise stor Lighed i Henseende til
Blomst og Stængelblade, medens Skulperne ere typisk ulige; saaledes finder
man ikke blot indenfor hver Klasse og hver Orden, men i hver Familie og i
hver Slægt det Lige i de forskjelligste Retninger combineret med det Ulige:
hvad betyder nu dette?

Det betyder, at to Planter, som i eet eneste Punkt ere hinanden typisk
ulige, maae i alle organiske Punkter være indbyrdes ulige; det betyder, at
det i Ligheden Ulige beroer paa en ideel Metamorphose. At Kronen paa det ene
Exemplar er blaa, paa det andet rød, viser, at den organiske Virksomhed i
den ene Krone er forskjellig fra den tilsvarende Virksomhed i den anden; men
de organiske Functioner i Kronen ere jo, ifølge Organismens Begreb, i den
inderligste Vexelvirkning med Functionerne i Bægeret, i Dækbladene, i
Stængelbladene, med eet Ord: i alle Plantens øvrige Dele. Det tilsyneladende
Lige maa altsaa dog være et Ulige, siden det i et enkelt Organ frembringer
et Uligt, og de i Phænomenet saa Ulige maae have væsentlig Lighed, siden de
fremkomme i organisk Sammenhæng med alt det øvrige Lige. Den røde og den
blaa Anagallis, der navnlig med Hensyn paa Kronens Farve ere hinanden saa
ulige, at man som Følge heraf har kunnet antage en Artsforskjel, ere derhos
i alt Øvrigt saa indbyrdes lige, at den sidste endog kan betragtes som en
Afart af den første. Her er da, hvad Kronerne angaaer, ikke simpelthen
Ulighed i en enkelt Henseende og for Resten Lighed i alt det Øvrige; men
Uligheden imellem Kronerne forudsætter en, hvor vanskeligt det end er at
paavise den i Phænomenet, alle øvrige Dele gjennemgaaende Ulighed, og
Ligheden i alt det Øvrige forudsætter saa igjen, at Uligheden i Kronernes
Farve maa beroe paa en Ændring af et forøvrigt Fælleds og altsaa dog være
begrundet i Lighed.

Brægnen og Palmen \textbf{ere ikke} hinanden i nogle Henseender lige, i
andre derimod ulige; men \textbf{de ere} i alle væsent\-%
% --- p. 311 ---
\setcounter{footnote}{0}%
lige Punkter hinanden lige, uagtet de i alle Punkter dog ere hinanden ulige.
Dersom Brægnen ikke i alle væsentlige Punkter var identisk med Palmen, og
dersom denne Identitet ikke maatte gjælde bogstavelig, saa kunde Planten
ikke være Plante. Thi hvad \textbf{er} vel Brægnen Andet end Plante, og hvad
\textbf{er} Palmen Andet end Plante? Men dersom Plantens Identitet med sig
selv som Brægne og Palme ikke i alle Punkter var bestemt ved Forskjel og
Modsætning,
% PRINTER'S DEFECT: the o of "Modsætning" did not print; the print shows "M"
% followed by the empty width of the sort and then "dsætning". The word is
% given here in full and the failure recorded.
maatte Brægnen være Palme, og Palmen Brægne; at det samme Væsen, den samme
Plante, er blevet til to Væsener, to indbyrdes forskjelligartede Planter,
deri sees Metamorphosen. Her er da, ifølge Metamorphosen, ikke en deelviis
Identitet, vexlende med en deelviis Forskjel, men en ved Forskjellen bestemt
Identitet, en i Identiteten begrundet Forskjel: Brægnen \textbf{er ikke} i
disse og hine Dele, t.~Ex. i Sporerne, som en Kryptogam, i andre Dele,
t.~Ex. i Bladene eller i Stammens Karbundter, som en Phanerogam; nei,
Brægnen \textbf{er} i alle Dele og Deles Dele kryptogam, ligesom Palmen i
alle Punkter og Dele \textbf{er} phanerogam.

Kun ved at fæste Blikket paa Ideen \textbf{er} det muligt at faae
Sammenligningens og Abstractionens Misviisninger corrigerede, muligt at
fatte Metamorphosen. Saalænge Tingsaarsagerne afgive Erkjendelsens hoieste
Forudsætninger, vil Identiteten af organisk Identitet og organisk
Modsætning altid fremstille sig som en Begrebsmystification.
% "hoieste" with plain o, verified at 600 dpi against the slashed ø of
% "ifølge" on this same page
Langt fra at see Brægnens og Palmens punktuelle, gjennem alle Enkeltheder
udtalte, Forskjellighed reflecteret i Palmens concrete Væsenseenhed med
Brægnen, vil den kritiske Forstandighed, saalænge den har
Endelighedssplinten i Oiet, d.~v.~s. saalænge de physiske Mellemaarsager
endnu ikke ere forvandlede til Ideemomenter, bestandig hjælpe sig med „en
deelviis Lighed“ og „en deelviis Forskjellighed“, med „Lighed i Noget“ og
„Ulighed i Andet“, med „Fælledsskab i det Almene“ og
% --- p. 312 ---
\setcounter{footnote}{0}%
„Eiendommelighed i det Enkelte“; saalænge Tænkningen ikke har arbeidet sig
igjennem til Metamorphosens Kategorie, og derved til Indsigt i det Almenes
ideeltmaterielle Concretion med det Særegne og Enkelte, vil det være en
Umulighed at anerkjende Plantens Indre, Grundformen selv, for Mere end en
Abstraction.

At der med Erkjendelsen foregaaer en Forvandling, en væsentlig Forandring,
naar de physiske Aarsager opfattes som Momenter i Idee-Aarsagen,
\textbf{er umiskjendeligt}. I Modsætningen af Aarsag og Virkning have vi da,
fra Ideens Synspunkt, en virkelig, virksom Modsætning af
Causalbestemmelsernes System i dets indre, ideelle Totalitet paa den ene
Side og de tilsvarende, ved Stoffets Articulation bestemte, Virkninger i
deres reelle Totalitet paa den anden Side. Falder nu Eftertrykket paa den
saaledes gjennemførte Modsætning, saa gjør ethvert af Leddene, Planten selv
og Plantens Idee, Brægnen og Brægnens Idee, Palmen og Palmens Idee, i lige
Grad Fordring paa Virkelighed. Det kan nemlig ikke være Ideen for sig, der
umiddelbart er det Virkelige, thi Ideens Virkeliggjørelse skeer jo netop
formedelst den materielle Existens; endnu mindre kan den materielle Existens
for sig tilegne sig Virkelighedens Prioritet, da den jo ikke frembringer sig
selv, men realiseres af den sig deri realiserende Idee: Ideens og Plantens
Virkelighed \textbf{er} een og den samme; her \textbf{er ikke to, men eet
Virkeligt}.

Hvad \textbf{er} saa Ideen? Den virkelige Plante! Og Planten? Den virkelige
Idee! Den Identitet, hvori Modsætningen nu er opløst, afgiver et Correctiv,
som den videnskabelige Erkjendelse umulig kan savne. For den
sandselig-forstandige Betragtning fremstiller den umiddelbare Existens, den
existerende Plante, sig som det Positive, det i Virkeligheden Tilværende;
derfor siger den sunde Sands: „jeg seer nok Planten, men Plantens Idee seer
jeg ikke.“ Deri har den sunde Sands nu forsaavidt fuldkommen Ret, som det jo
ikke er
% --- p. 313 ---
\setcounter{footnote}{0}%
det umiddelbart Positive, der \textbf{er} Ideen. Dersom det umiddelbart
Positive var Ideen, saa maatte Dyret jo æde Ideen, naar det æder Planten,
saa maatte man da ogsaa kunne sønderrive Ideen, naar man sønderriver
Planten, koge Ideen, naar man koger Planten o.~s.~v.; dersom det umiddelbart
Positive var det absolut Positive, saa var Ideen uden Realitet, og den
kritiske Forstandighed beholdt Ret. Men dette \textbf{er} just det for
Erkjendelsen uundværlige Correctiv, at det, der fremstiller sig for den
sandselige Beskuelse, det umiddelbart Positive, ikke \textbf{er} det
positive Sande selv, at det i sin phænomenale Umiddelbarhed Beskuelige netop
\textbf{er} saa forgjængeligt, saa flygtigt, fordi det væsentlig
\textbf{er} det Negative, det Uselvstændige, der tjener et Andet, det Stof,
hvori Idealiteten brænder, det Skin, hvori Ideen skinner.

Ikke derved, at den umiddelbare Existens, den existerende Plante,
destrueres og giver Plads for andre ligesaa umiddelbare Existenser, andre
Planter, men derved, at det i Existensen Sandselige, det umiddelbart
Positive, fremstiller sig som det i Reflexionen Gjennemsigtige --- det
Negative, der væsentlig bekræfter det af den negerede
Negation\footnote{Jvfr. Phil. Propædeutik. S.~152--58.} resulterende
Positive, Ideen selv, det teleologisk Virksomme ---, kan Sandheden
erkjendes.

γγ) \emph{Individuationer.} Som existerende Organisme, som Brægne, som
Palme, som Piil o.~s.~v., fremstiller Planten sig for den umiddelbare
Betragtning i Skikkelse af naturlig Individualitet; men hvor \textbf{er} nu
i denne Individualitet det egentlige Individ? Tænker man sig ved Individet
den organiske Dannelse, der har Evne til enten selv at bestaae som
selvstændig Organisme eller dog idetmindste til at frembringe en saadan,
opdager man let, at Planten, det vegeterende Exemplar, ikke er et enkelt
Individ, men en Yngel
% --- p. 314 ---
\setcounter{footnote}{0}%
af Individer. Hvert Skud (\textit{turio}), hver Knop (\textit{gemma}),
hvert Blad\footnote{Alexander Braun. Das Individuum der Pflanze. Berlin 1853.}, ja selv en Deel af et Blad, kan jo under gunstige
Betingelser afgive en reel Forudsætning for en selvstændig voxende
% printed mark: *)  --- first of two notes on this page
Plante. Selv om man med Schleiden gaaer tilbage til Cellen, finder
man dog ikke det absolute Individ, eftersom jo en eencellet Alge, en
\textit{Caulerpa t.} Ex., kan rives itu, og Delene endda voxe.
% antiqua covers "Caulerpa t." -- the abbreviation "t." is set in antiqua
% with the genus, while "Ex." is Fraktur; transcribed as printed.

Spørgsmaalet om Planteindividet og dets Realitet er her af en saa
meget mere afgjørende Betydning, som der under dette Spørgsmaal
synes at skjule sig en betænkelig Spænding imellem den paradigmatiske
Selvhedslov og Forplantningsloven. --- „Efter den Tendens, der i de
nyere Forskninger gjør sig meer og meer gjældende, kan dette den hele
Plante sammensættende Blad, der idelig omdannet træder frem i
fuldkomnere og fuldkomnere Former paa Planten, ikke længer være at
betragte som et enkelt Redskab eller „Organ“, i den Forstand et af
vore Lemmer eller Organer er det, men er et afsluttet Hele og Ene,
med den hele Væxts Natur i sig. I Sammenligning med saadanne
individuelle Planter eller Væxtener er den hele Plante en Samvæxt, et
organisk Samfund af Blade, der efter indre Harmonie stræbe hver til
sit Maal, og dog alle til et fælleds og hoiere, hvoraf hver enkelts Vel
% "hoiere" as printed: plain o, unslashed at 600 dpi, beside plainly
% slashed ø in "Spørgsmaalet"/"afgjørende"/"gjør" on this page.
er afhængigt. Denne Betragtning af Bladet som et Ene og af Planten
som en Stat anseer jeg for den eneste naturlige, og for den eneste, der
tillader at see Plantelivet i dets rette Lighed og Ulighed med
Dyrelivet . . . . \emph{Men ere Bladene virkelige Ener, saa ere ogsaa
Støvbladene og Frugtbladene i deres gode Ret som saadanne og de to
Kjøn ligesaa lidt her som andetsteds forenede i eet Væsen.}“
(Steenstrup\footnote{Unders. over Hermaphroditismen. S. 85--86.}).
% printed mark: **)

% --- p. 315 ---
\setcounter{footnote}{0}%
Kritisk stiller nu Sagen sig saaledes. Afviser man --- hvad her
vistnok er al mulig Grund til --- Læren om Hermaphroditismen, da
maa man ogsaa conseqvent indrømme, at Støvblade og Frugtblade ere
Individer af modsat Kjøn; men, naar Planten i sin Heelhed,
istedetfor at være et eneste Individ, er et organisk Samfund, en Stat
af Individer, hvad bliver der saa af den paradigmatiske Eenhed, af
den i det hele System sig articulerende Selvhed? Opfatter man fra
den anden Side --- hvad her dog vistnok ogsaa er al mulig Grund til
--- den hele Plante som en levende Eenhed og reducerer de foregivne
Bladindivider til Momenter i dens Individualitet, saa falder her jo
vistnok et afgjørende Eftertryk paa den paradigmatiske, de mangfoldige
Led gjennemtrængende Selvhed; men, naar de enkelte Blade, ifølge
denne Synsmaade, ophøre at være „Væxtener“, hvorledes undgaae vi
da at gjøre Brud paa Forplantningsloven og at falde tilbage til
Hermaphroditismen? Den Udvei, at lade det ene til hele Planten
svarende Paradigme opløse sig i et System af særskilte Paradigmer, er
altfor mislig. Skal hvert Blad virkelig repræsentere en apriorisk
selvstændig Selvhed, maa hver Celle, ja vel endog hver Cellehinde og
hver Cellesaftdraabe jo ligeledes forudsætte sit særskilte Paradigme,
sin apriorisk selvstændige Selvhed; men naar en saadan Conseqvens
gjennemføres, hvad er saa den hele Forklaring med Paradigme og
Selvhed og Anlæg og Præformationer vel Andet end et tomt
Skyggespil, et metaphysisk Tankespil?

Dersom Eet af To gjaldt ubetinget, at Planten nemlig i sin Heelhed
maatte være \emph{enten} en levende Eenhed \emph{eller} en levende
% only the FIRST "enten"/"eller" pair is letterspaced; the second pair,
% on the following line, is set plain. Transcribed as printed.
Fleerhed, enten et absolut Totalindivid eller en Sammensætning af
absolut selvstændige Partialindivider: saa vilde her unegtelig være et
Knudepunkt. Men det saaledes opstillede Enten-Eller har ingen
Gyldighed. Partialindividet er i vor „organiske Stat“ ligesaa lidt
absolut som Totalindividet. Eftertrykket falder her, paa hvilken Side
man saa
% --- p. 316 ---
\setcounter{footnote}{0}%
end vil søge det Individuelle, ikke saameget paa „Væxtener\-%
nes“ afsluttede Tilværen som paa deres eiendommelige Virken og den
Virksomhed, hvorunder de dannes, altsaa: ikke saameget paa
Individerne som paa Individuationen. Men Individuationen gaaer,
naar vi see den i Ideens Lys, i tvende hinanden indbyrdes modsatte
Retninger, en centrifugal og en centripetal. I Kraft af den
centrifugale Stræben bliver Eenheden Middel for Fleerheden;
Totalindividet faaer fra denne Side en Betydning som „en Samvæxt“,
som „den organiske Stat“, hvori Partialindividerne ere selvstændige
Borgere. I Kraft af den centripetale Stræben bliver Fleerheden
omvendt Middel for Eenheden; de vegeterende Borgere ere fra denne
Side kun Statens vegetative Organer, Realmomenter, hvori den
leddeelte Heelhed reflecterer sin oprindelige Eenhed.

Hvilken af disse Individuationens Retninger man skal betragte som
den overordnede, vil for ethvert enkelt Tilfælde beroe paa, hvorvidt
den paagjældende Opgave fordrer, at der skal lægges udtrykkelig Vægt
enten paa de enkelte Functioner i deres Særskilthed eller paa
Grundfunctionen, den det hele System bestemmende Functionseenhed.
Handles her saaledes om „den overalt i den levende Natur udtalte“,
Hermaphroditismen udelukkende, „Forplantningslov“, da handles her
jo udtrykkelig om en af Livets særskilte Functioner, en Function, der
just betinges af en ganske særskilt Modsætning; her er da tilvisse al
Grund til at lægge Eftertrykket paa det Særskilte og see de modsatte
Kjøn repræsenterede af modsatte Individer. Er Spørgsmaalet derimod
om Selvhedens Realitet, om Paradigmets aprioriske Gyldighed, da maa
den Functionerne gjennemgaaende Almeenfunction nødvendigviis blive
det Overordnede, Enkeltfunctionerne i deres Fleerhed og Særskilthed
det Underordnede; det følger da af Sagens Natur, at hele
Eftertrykket maa falde paa Eenheden selv, den sig i Heelheden
localiserende Eenhed, Principeenheden.

% --- p. 317 ---
\setcounter{footnote}{0}%
Istedetfor et eneste absolut Individ eller en Samvæxt af absolute
Individer have vi altsaa i Planten \emph{et System af organiske
Individuationer}. Dette System giver os da, forsaavidt det omfatter de
for den særlige Plantes hele Levnetsløb charakteristiske
Metamorphoser, Plantens naturlige Individualitet, den Skikkelse, hvori
Paradigmet er realiseret\footnote{Fra et æsthetisk Synspunkt vil den
høiere staaende Plante naturligviis altid gjøre Indtryk af at være et
forgrenet Naturindividuum. --- „So stellen sich denn von der Wurzel
bis zur Frucht eine Reihe bedeutsamer an die Innigkeit anklingender
Gestalten hin, welche jedoch erst dodurch zu ihrer wahren Bedeutung
% "dodurch" as printed (for dadurch); all three witnesses agree and the
% o is unambiguous at 600 dpi. Printer's error, transcribed as printed.
gelangen, daß sie als das Ganze des Pflanzenindividuums
constituirende Momente erfaßt werden. Denn so wie die Stimmungen
erst dann in ihrer Wahrheit erkannt werden, wenn man sie auf die
Einheit der Innigkeit zurückführt, so erhalten auch die Zustände der
Pflanzengestaltung erst dann ihren wahren Charakter, wenn sie zum
Einen Ganzen, zur Physiognomie des Pflanzenindividuums
zusammenstimmen. Für sich genommen, mag jede Stimmung und
jeder Zustand des Reizenden Manches enthalten, die Wahrheit des
Lebens offenbart sich aber doch erst in und aus der Geschlossenheit
des Individuums“. F. Th. Bratranek. Beitr. z. e. Æsthet. der
Pflanzw. S. 224. --- Jvfr. A. S. Ørsted, Et Bidrag til Træernes
Architektonik. Tidsskr. for popul. Fremstill. af Naturvidskb. 2den
Række. 4de Bd. Kbh. 1862. S. 78. flg.}. I denne Skikkelse er
Forskjelseenheden af det Hele og dets Dele, det Indre og det Ydre,
Totalindividet og dets Partialindivider, da ogsaa, efter hvad vi i det
Foregaaende have seet, alsidig bestemt ved den articulerede Modsætning
af Selvhed og Andethed, forsaavidt disse Modsætningens Factorer i
Functionernes System vedblivende hævde hver sin Gyldighed, optræde
hver som eiendommelig Magt.

Spørge vi, hvilken af de to Magter der er den stærkere, Selvheden
eller Andetheden, vil en Dialektik af Magt og Afmagt vise sig paa
begge Sider. Væsentlig forstaaet er Selvheden, der jo principielt
bærer Andetheden i sig, det Oprindelige, det Høiere, og forsaavidt det
Mægtigere; Andetheden, der har
% --- p. 318 ---
\setcounter{footnote}{0}%
sit Integrerende udenfor sig, er derimod det Afledede, det Lavere, og
forsaavidt det Afmægtige, det Svagere. Men falder nu, hvad
Virkeligheden jo ligeledes fordrer, Eftertrykket paa Phænomenet, det
Factiske, den i Kjendsgjerningen slaaende Umiddelbarhed, saa vender
Forholdet sig om, saa maa Andetheden, der jo netop har sin Styrke i
det udvortes Umiddelbare, nodvendigviis vise sig som det Oprindelige,
% p. 318 sets the ø-less forms throughout: nodvendigviis, gronne,
% sonderrive, Oine, skjonne, fuldestgjor, skjore -- all with plainly
% unslashed o at 600 dpi. Transcribed as printed.
det Stærkere, og det ikke blot \emph{uagtet} den er, men \emph{just
fordi} den er det efter sit Væsen Afledede, det Svagere; ligeledes maa
Selvheden, hvis Gyldighed beroer paa Reflexionsenergie, nodvendigviis
fremstille sig som det Afledede, det Abstraherede, en i subjectiv
Reflexion mat Afskygning, og det ikke blot \emph{uagtet} den er, men
\emph{just fordi} den er det i Grunden selv Oprindelige, det i
objectiv Reflexion Evige og Uforkrænkelige.

I den vegeterende Plante, en Zvibel t. Ex., viser Livet sig som en
forunderlig Tvetydighed, en uendelig Vexelbestemmelse af Virkeligt og
Uvirkeligt. Denne Blomsterrigdom og disse saftige gronne Blade med
de Millioner af Celler, hvortil Bestanddelene hentes fra Vand og
Luft, give os ret et anskueligt og farverigt Billede af Principets
allestedsnærværende Virken, af den Overlegenhed, hvormed Selvheden
beseirer Atomernes elementære Andethed, af den Idealitetens executive
Magt, der i alle Punkter fuldestgjor Paradigmets teleologiske
Fordringer. Men just som vi uvilkaarlig glæde os ved at have dette
Selvhedens udtrykte Billede lige for vore Oine, og træde nærmere til
for ret at iagttage den i det skjonne Værk sig aabenbarende
Værkmester, bliver Idealiteten pludselig borte. Ideen svinder som et
Phantom, det ideelt Virkelige forvandler sig til et Uvirkeligt. Man
kan jo forstyrre Ideens Værk, man kan afrive Blomsterne,
sonderrive Bladene; ja det gaaer endogsaa lettere med at sonder\-%
rive disse saftige gronne Blade og bryde denne skjore Stængel, end
med at sonderrive en Snor eller bryde et Glas: „dersom den levende
Plante virkelig er Eet med den sig deri
% --- p. 319 ---
\setcounter{footnote}{0}%
realiserende Idee, saa“ --- siger man --- „maa Ideen rigtig\-%
nok være et afmægtigt Skyggevæsen, siden man med saa ringe
Uleilighed kan tilintetgjøre dens Værk!“\footnote{„Wollten wir irgend
eine Idee als die wirkende Macht ansehen, die in einer Erscheinung
thätig ist, so wäre wenigstens dies gewiß, daß wir uns zu einer
unthätigen Stellung ihr gegenüber verurtheilten. Denn Ideen der Natur
zu verändern, ihre Wirkungen aufzuheben oder zu begünstigen, besitzen
wir doch wohl kein Mittel, wir müßten denn den Sinn dieser Ansicht
entschlossen so vervollständigen, daß wir uns solchen Mächten
geistiger Art mit ebenso geistreichen Beschwörungen und
Zauberformeln gegenüberstellten.“ Lotze. Allgem. Physiol. S. 36. ---
Denne Synsmaade kan i Korthed sammenfattes i disse Ord: Dersom
Naturen ikke var ideeløs, kunde Menneskene ikke ved naturlige Midler
indvirke paa dens Frembringelser. Heraf følger saa med det Samme,
at der heller ingen virksom Fornuft kan være i Naturen; thi dersom
Naturen virkelig er ideeløs, maa den ogsaa være fornuftløs.} Men Eet
er Reflexionens, et Andet Umiddelbarhedens Magt.

Ifølge den vegetative Organismes naturlige Eenhed med Ideen, er
Planteriget at opfatte som et paa Jorden realiseret, i materielle Former
udpræget, Ideerige. For at gjøre et Indblik i dette Riges
Organisation, vil det være nodvendigt at undersøge, hvorvidt „det
naturlige Systems“ Arter, Slægter, Familier, Ordener, Klasser og
Rækker lade sig erkjende som Ideer af Ideer, Individuationer af een
Grundidee.

Individualideen, det i Exemplaret virkeliggjorte Paradigme, er ikke
et ubetinget og abstract Enkeltvæsen, men en levende concret Eenhed,
der reflecterer sin Bestaaen i en levende Fleerhed\footnote{Overalt
hvor Talen er om Exemplarets og Individualideens dialektiske
Forskjelseenhed, maa det udtrykkelig fastholdes, at Exemplaret
forsaavidt maa congruere med Ideen, som dets Existens kun er et
materielt Udtryk for det System af articulerende Bestemmelser, der
udgjøre Selvhedens paradigmatiske Indhold. Det
% the printed page break p. 319/320 falls here, inside this note
Urimelige, som altid viser sig, naar man paa en aldeles umiddelbar
Maade vil overføre Sandselighedens Kategorier paa det Ideelle, kan da
i denne Sammenhæng ikke overraske os. Imidlertid --- thi den
literære Bestialitet er ofte, hvor man mindst venter det, tilrede med
„Bidrag til en Kritik“ --- kunde et par Exempler paa plumpe
Misforstaaelser dog maaskee her være paa rette Sted. „Falder“ ---
saaledes kan man jo vistnok fra „Kritikens“ Standpunkt raisonnere ---
„et Piletræes Individualidee sammen med det virkelige Piletræ, saa
river man jo et Stykke af Individualideen, naar man afriver en
Pileqvist. Kan det ene Pileexemplar af samme Art være snart mindre,
snart større end det andet, saa maa den ene tilsvarende Individualidee
jo ogsaa kunne være snart mindre, snart større end den anden;
Ideerne maae altsaa kunne groe og voxe ligesom Exemplarerne. Men
dersom en Individualidee voxer med sit voxende, saa maa den ogsaa
visne med sit visnende, forraadne med sit forraadnende Exemplar,
o. s. v.“ Har man nu ved den Slags „Bidrag til en Kritik“ faaet det
beviist med høi Røst, og altsaa tilgavns beviist, at Exemplaret umulig
kan fremstille nogen Individualidee, eftersom man da under en saa
urimelig Forudsætning ved at gjennemskjære Exemplaret absolut maatte
gjennemskjære Ideen, saa behøver man jo blot at drage en Conseqvens,
der minder „Kritiken“ om, at en Kugles Tyngdepunkt heller ikke kan
have Realitet, (Mathem. og Dialekt. S. 129--30) eftersom man da,
% the initial R of "Realitet" is very faintly inked in the scan; the
% reading is not in doubt.
under en saa urimelig Forudsætning, bogstavelig maatte \emph{kløve
Tyngdepunktet}, naar man halverede Kuglen.}. Gaaer Individualideens
Existens paa
% --- p. 320 ---
\setcounter{footnote}{0}%
denne Maade op i den tilsvarende Exemplarexistens, da er det en
Selvfølge, at Paradigmet \textit{in concreto} maa, paa saa Viis, være
ligesaa individuelt som det individuelle Exemplar. Her er da ikke en
Mangfoldighed af Exemplarer, hvori det ene og samme Paradigme har
gjentaget sig; et saadant almindeligt Paradigme er kun et Schema, en
abstract Form, ikke en Idee: Individualideen kan ikke gjentage sig.
Mener man imidlertid dog at see en Gjentagelse af Individualideen
deri, at det ene Paradigme saa absolut ligner det andet, da maa man
jo ogsaa ansee de hverandre i alt Væsentligt lignende Exemplarer for
Gjentagelser af eet og samme Exem\-%
% --- p. 321 ---
\setcounter{footnote}{0}%
plar. Er det nu en Modsigelse, at eet og samme Exemplar skulde
kunne existere som en Fleerhed af Exemplarer, da er det jo ogsaa en
Modsigelse, at een og samme Individualidee skulde existere som en
Fleerhed af Ideer. De mange hverandre indbyrdes lignende
Exemplarer fremkomme ved ligesaa mange hverandre indbyrdes lige
Forplantningsacter; de mange hverandre indbyrdes lige
Forplantningsacter forudsætte ligesaa mange hverandre indbyrdes lige
Selvhedsprinciper; de mange hverandre indbyrdes lige Selvhedsprinciper
ere betingede Concentrationer af eet og samme, i sig ubetingede
Grundprincip: kun saaledes er en existentiel Gjentagelse tænkelig.

„Men“ --- kunde Nogen spørge --- „er det da ikke ligesaa urimeligt,
at eet og samme Princip skulde ved existentiel Gjentagelse kunne blive
til en Fleerhed af Principer, som at een og samme Idee skulde ved en
flere Gange gjentagen Existens blive til en Fleerhed af Ideer?“
Besvarelsen af et saadant Spørgsmaal er ikke vanskelig. At det samme
Existensprincip skulde, enten paa een Gang eller efterhaanden, opgaae
og virke i en Fleerhed af Existerende, er --- saalænge det afgjorende
% "afgjorende" here with plain, unslashed o, beside plainly slashed ø in
% "spørge"/"Spørgsmaal"/"uopløselig" on the same page. As printed.
Eftertryk falder paa „Existens“ og „Existe\-%
rende“ --- ganske vist en uopløselig Modsigelse; Existensen beroer paa
\textit{principium exclusi medii}: enten $a$ eller $b$; det ene
% "principium exclusi medii" antiqua; "enten"/"eller" are Fraktur, the
% letter-symbols a and b antiqua. Settled at 600 dpi.
Existerende udelukker det andet. Naar de Existerende nu saaledes
forholde sig exclusivt til hverandre indbyrdes, maae jo vistnok
Existensprinciperne, forsaavidt de, hvert især, virkeliggjøres i
tilsvarende Existenser, ligeledes forholde sig exclusivt til hverandre
indbyrdes. Men ligesom enhver Existens med indre Nødvendighed
forudsætter sin særegne Grund, saaledes maa \emph{ethvert
Existensprincip} paa tilsvarende \emph{Maade forudsætte sit
% the letterspacing lapses across the line break: "paa tilsvarende" is
% set plain between two letterspaced stretches. As printed.
Grundprincip}. De særligt Existerende forholde sig exclusivt i
Existensen; men de særlige Grunde forholde sig ingenlunde exclusivt i
Grunden. Tvertimod! de særlige Grunde fremkomme just paa givne
Betingelser
% --- p. 322 ---
\setcounter{footnote}{0}%
af een for dem alle fælleds Grund, Grundeenheden, den almene
Grund. Kun idet Grundeenheden, den almene Grund, gjennem
Betingelserne reflecterer sig i de særlige Grunde, have vi den for
Existensen nødvendige og tilstrækkelige Grund (\textit{ratio
sufficiens}).

I den uendelige Virksomhedsreflexion, der lader det Ubetingede
gjennemtrænge Betingelserne, gjælder nu det Samme om Principerne.
Idet de organiske Existensprinciper ved deres Identitet med de
existerende Organismer, de sandselig givne Exemplarer, stille sig i
Tilværelsen exclusivt, vise de som særlige Grundaarsager tillige deres
Identitet med Grunden, og ere saaledes i al deres Ubetingethed
betingede Virkninger af den i dem alle reflecterede Principeenhed,
Grundprincipet. Da nu Individualideen ikke er Andet end det i
Organismen articulerede Existensprincip, \emph{den indre, i sit
materielle Ydre energisk reflecterede Organisme}, nærmere bestemt som
absolut Forskjelseenhed af Midler og Formaal: saa forstaaes heraf, at
Individualideerne i deres existentielle Sondring absolut maae være
Ideer af een og samme Grundidee. Denne i Individualideerne sig
realiserende Grundidee er nu \emph{Arten}.

Hvilken Forskjel det giver, om vi opfatte Arten som en oprindelig
virksom Realidee, eller blot som en Form, et ved Exemplarernes
Sammenligning abstraheret Begreb, er ioinefaldende: kun en
% "ioinefaldende" -- the o is very lightly inked but its bowl prints
% complete with no ink anywhere in the slash region, against plainly
% slashed ø in "nødvendige"/"dødt" on this page. As printed.
overfladisk Tænkning kan forvexle den levende Art med det tomme
Schema. Sige vi, at alle Klodeskyer ere af samme Art, saa opfatte vi
Artsbegrebet som Schema; thi den typiske Form er, hvad
Skydannelsen angaaer, ikke Princip, men Resultat. Det er den i de
mange Dannelsesresultater fremherskende Lighed, vi abstrahere som
ligesaa mange indbyrdes lige Former; disse mange indbyrdes
uafhængige, mod hverandre ligegyldige Former tage sig saa i
Abstractionen ud, som om det var een eneste, de enkelte Dannelser
begrundende Form, en Grundform. En saadan Form er et dødt
Abstractum; hvor overfladisk at betragte den levende Art
% --- p. 323 ---
\setcounter{footnote}{0}%
paa samme Maade! I den levende Art er Formvirksomheden
væsentlig Principvirksomhed; den eiendommelige Form, hvori Arten
som Art udtaler sin Existens, er oprindelig bestemt ved Anlæg og
Præformation. Den i Anlæget præformerede Form maa, for at være
Artens Form, principielt differentieres til særskilte Individualformer:
\emph{kun formedelst Individualideen kan Arten komme til Existens
med det existerende Exemplar.}

Forsaavidt Individualideerne ere af samme Art, maae de ikke alene
ligne hverandre i væsentlige Punkter, men, hvad her i Særdeleshed
bliver at fremhæve, være hverandre i alle Henseender absolut lige.
Dersom nu alle Planteverdenens Individualideer vare hverandre absolut
lige: saa maatte den ideelle Concretion standse ved Arten som dybeste
Grundidee. Exemplarernes Mangfoldighed paa Jorden kunde da være
saa stor, den vilde; hele Rigdommen bestod alligevel i en blot
udvortes Overflødighed: Concretionen kunde ikke gaae videre. Vilde
Nogen herimod gjøre den Indvending, at to Ideer umulig kunne være
hinanden absolut lige, uden at de nødvendigviis maae smelte sammen
og ifølge \textit{principium indiscernibilium} blive til Eet: da svares
hertil, at en Existens\-%
forskjel meget vel kan bestaae med en fuldkommen
Lighed\footnote{Der Satz des Nichtzuunterscheidenden gründete sich
eigentlich auf die Voraussetzung, daß, wenn in dem Begriffe von einem
Dinge überhaupt eine gewisse Unterscheidung nicht angetroffen wird,
sie auch nicht in den Dingen selbst anzutreffen sei; folglich seien alle
Dinge völlig einerlei (\textit{numero eadem}), die sich nicht schon in
ihrem Begriffe (der Qvalität oder Qvantität nach) von einander
% "Qvalität"/"Qvantität" as printed -- the Fraktur v after Q is
% unambiguous at 600 dpi (the book's own qv convention, cf. qvalitativ).
unterscheiden. Weil aber bei dem blossen Begriffe von irgend einem
Dinge von manchen nothwendigen Bedingungen einer Anschauung
abstrahirt worden, so wird durch eine sonderbare Uebereilung das,
wovon abstrahirt wird, dafür genommen, daß es überall nicht
anzutreffen sei, und dem Dinge nichts eingeräumt, als was in seinem
Begriffe enthalten ist. Kant. Von der Amphibolie der
Reflexionsbegriffe. Krit. d. r. Vern. S. 268 o. fl.},
% --- p. 324 ---
\setcounter{footnote}{0}%
og at det just er Existensforskjellen, der væsentlig angaaer Ideen som
Individualidee.

At Individualideer af samme Art ikke alene kunne, men
nødvendigviis maae, trods deres indbyrdes absolute Lighed, adskille sig
i existentiel Mangfoldighed, at denne existentielle Mangfoldighed
ingenlunde kan skrives paa Exemplarernes Regning, men maa afledes
af Ideen selv, er en ligefrem Conseqvens af Selvhedsloven: saamange
principielle Selvhedspunkter, saamange Existensprinciper, saamange sig
actualiserende Paradigmer, saamange Individualideer. Skal
Concretionen altsaa gaae dybere, maa det skee derved, at der danner
sig Rækker af indbyrdes ulige, væsentlig forskjellige Individualideer,
hvilket saa igjen forudsætter indbyrdes forskjellige Artsideer.

Den existentielle Forskjel imellem Artsideerne er da med det Samme
reflecteret i en indre Væsensforskjel, en gjennemgaaende
Begrebsforskjel: en saadan voxende Rigdom af indre
Forskjelsbestemmelser er en væsentlig Forudsætning for en stigende
Concretion. Imellem Exemplarer af \textit{Fragaria vesca} t. Ex. kan
der, sandselig seet, vistnok vise sig ikke faa Forskjelligheder; men
disse Forskjelligheder angaae tilhobe det Udvortes, det Tilfældige og
Forsvindende, og vedkomme aldeles ikke Individualideerne; hvorimod
Forskjellen imellem en \textbf{\textit{Fragaria vesca}} og en
\textbf{\textit{Fragaria collina}}, hvormegen Lighed Exemplarerne end
% the second pair of species names is set in a fat-face antiqua: same
% cap-height as the plain antiqua "Fragaria vesca" ten lines above, but
% 44 % more ink over the same measure. The first instance is plain.
have med hinanden, er en Individualideerne gjennemtrængende Forskjel,
en Forskjel, der stiller den ene Artsidee imod den anden. Den
Grundidee, hvoraf Artsideerne udvikle sig, og hvori de igjen reflectere
deres udviklede Indhold, er Slægtsideen. Som de concrete Artsideers
concrete Eenhed repræsenterer Slægtens Idee altsaa nu den rigere
Concretion.

Her danner sig ikke et Slægtsbegreb, fordi her forekomme
forskjellige Arter; men her forekomme forskjellige Arter, fordi
Slægtens Idee \textbf{kun kommer til Existens ved at differen}\-%
% --- p. 325 ---
\setcounter{footnote}{0}%
\textbf{tiere sig i Artsideer,} ligesom Artens Idee kun kommer til
% the fat-face run begun on p. 324 ends with "Artsideer,"; the parallel
% clause that follows is set plain. As printed.
Existens ved at differentiere sig i Individualideer. Langt fra at være
abstractere end Individualideerne, er Artsideen som den causale Grund,
hvori de existentielle Differenser præformeres og efter Udviklingen
erindres, tvertimod concretere; langt fra at være abstractere end de
tilsvarende Artsideer, er Slægtsideen, i hvis causale Grundeenhed de
gjennem Arterne udviklede Forskjelsbestemmelser reflectere sig som
Indholdsmomenter, den concretere Idee.

Saalidt som Artsideen kan være en Abstraction af Individualideer,
saalidt kan Slægtsideen være en Abstraction af Artsideer. Ikke ved
formel Abstraction, men ved dialektisk Reduction naaer Tanken
gjennem de særskilte Ideer til de almindelige. Jo almindeligere et
abstraheret Begreb er, desto indholdslosere maa det naturligviis være;
% "indholdslosere" with plain o, beside plainly slashed ø in
% "Selvfølge"/"opløste"/"afgjørende"/"ifølge" on this page. As printed.
det er da en Selvfølge, at den formelle Logik, der kun holder sig til
Abstractionen, maa lade et Begrebs Indhold staae i omvendt Forhold
til dets Omfang. Jo almindeligere den ved Reductionen fundne Idee
er, desto indholdsrigere er den; den har jo alle ved Reductionen
opløste Momenter i sit Indhold. Dialektiken, der gaaer Reductionens
Vei, sætter desaarsag Idee-Indholdet, ikke i et omvendt, men i et
ligefremt Forhold til Omfanget.

Den for Idee-Udviklingen afgjørende, med Selvhedsloven identiske
Concretionslov er allerede bestemt, naar Slægtsideen er naaet.
Istedetfor at betragte Gjennemgangen gjennem Individerne til Arten,
gjennem Arterne til Slægten, gjennem Slægterne til Familien o. s. v.,
som en Overgang fra det Concretere til det Abstractere, maae vi da,
ifølge den paaviste Lov, omvendt bestemme denne Gjennemgang som en
Fordybelse i stedse rigere, stedse inderligere Concretioner.

Hvorledes de organiske Exemplarer gjennem Individual\-%
ideerne i hver Art forholde sig til Arterne, hvorledes Ar\-%
terne i hver Slægt væsentlig forholde sig til Slægten, Slægten
% --- p. 326 ---
\setcounter{footnote}{0}%
gjennem alle sine Arter til Familien o. s. v., kan dialektisk bestemmes
ved Concretionsloven; derimod gives der ikke nogen Udviklingslov,
hvorved det \textit{a priori} kan bestemmes, hvormange og hvilke Arter
der nødvendig maae høre til hver Slægt, eller hvormange og hvilke
Slægter der oprindelig høre til hver Familie: den philosophiske
Erkjendelse gaaer ikke gjennem abstracte Ideer til virkelige
Kjendsgjerninger, men omvendt gjennem de virkelige Kjendsgjerninger
til de concrete Ideer. (Jvfr. Phil. Prop. S. 158--178).

Dersom de Forsøg, man i den nyere Tid har gjort paa at forklare
Arternes Tilblivelse af Varieteter, som i Tidernes Løb ere blevne
constante, skulle forstaaes saaledes, at det egentlig er
Omstændighederne og de ydre Betingelser, der bestemme de organiske
Typer: da maa en saa principløs og ideefornegtende Opfattelse af
Livsvirksomheden naturligviis møde Modsigelse i Philosophien. Men
det bør i denne Sammenhæng dog heller ikke glemmes, at Ideernes
Realitet er uendelig dialektisk, og at en dialektisk Realitet ikke lader
sig hypostasere i dogmatisk positiv Umiddelbarhed. Forholder det sig
virkelig saaledes, at den organiske Selvhed oprindelig forudsætter den
elementære Andethed, saa følger jo ligefrem heraf, at Livsudviklingen
ogsaa maa medføre en mangfoldig Vexelvirkning af det Typiske og det
Elementære. Varieteter, Aborter, Forskydninger og andre
Abnormiteter vise tydelig nok, hvilken Indflydelse de elementære
Betingelser maae have paa Paradigmets Udførelse\footnote{Hvad de
herhen hørende Undersøgelser angaaer, henvise vi til Charles Darwin
über die \emph{Entstehung der Arten im Thier- und Pflanzenreich.}
% the letterspacing of the German title begins at "Entstehung":
% "Charles Darwin über die" is set plain. As printed.
Aus dem Englischen übersetzt und mit Anmerkungen versehen von Dr.
H. G. Bronn. Stuttgart 1860. --- „Es sind“ --- hedder det i
„Schlußwort des Übersetzers“, S. 508--9 --- „die Existenz-Bedingungen,
welchen sich die Organismen anpassen müssen, und welche eine so
grosse Rolle in diesem Buche spielen. Sie sind theils organische und
theils unorganische, und die
% the printed page break p. 326/327 falls here, inside this note
ersten nach Hrn. Darwin weitaus die zahlreichsten und wichtigsten und
daher auch an und für sich geeignet, die manchfaltigsten Folgen zu
veranlassen. Doch eben diese organischen Prinzipien haben für Darwin
wieder den grossen Vortheil, daß, indem er sich auf ihre
Manchfaltigkeit und auf den Kampf ums Daseyn überhaupt beruft, er
der Nothwendigkeit überhoben ist, Rechenschaft von ihrer
Wirkungs-Weise im Einzelnen zu geben und nachzuweisen, welche
spezielle Folgen diese oder jene spezielle organische Bedingung auf die
Struktur und Entwickelung der ihrem Einfluß unterliegenden
Organismen überhaupt, oder auf einzelne insbesondere ausübt. Hr.
Darwin beruft sich auf jeder Seite darauf, daß nur solche
Abänderungen Aussicht auf Erhaltung haben, welche dem Individuum
und somit der künftigen Spezies nützlich sind; und theoretisch muß
man zugestehen, daß, woferne es eine natürliche Züchtung gebe, die
Sache sich nicht anders verhalten könne. Aber wir müssen gestehen,
doch in fast allen unseren aus angeblich innern Ursachen
hervorgegangenen Varietäten gar nicht finden zu können, worin denn
der Nutzen ihrer Abänderung bestehe . . . . Warum bekommt z. B. in
diesem Kampfe ums Daseyn eine Pflanzen-Art ovale statt lanzettlicher
und die andre lanzettliche statt ovaler Blätter? warum die eine einen
Dolden-artigen und die andre einen Rispen-förmigen Blüthenstand?
warum die eine fünf und die andre vier Staubgefässe, die eine eine
geschlossene und die andre eine weit geöffnete Blüthe? Wozu nützt
der einen Dieß und der andern das Gegentheil? Warum bewirken die
organischen Bedingungen Dieß? Mit welchen Mitteln fangen sie es an?
und wie müssen sie beschaffen seyn, um es zu können? Und wie kann
die eine Art der andern dadurch überlegen werden? Wir gestehen,
keinen Zusammenhang zwischen diesen Erscheinungen zu erkennen, und
Hr. Darwin würde uns antworten, daß es möglicher Weise so oder so
zugehen \emph{könne}. Wir gestehen ferner, uns vergeblich um
positive Beweise oder auch nur Belege in dieser Beziehung umgesehen
zu haben, manche spezielle Fälle eigenthümlicher Art etwa
ausgenommen; denn wir sind weit entfernt davon, allen solchen Einfluß
überhaupt läugnen zu wollen.“}.

% --- p. 327 ---
\setcounter{footnote}{0}%
„Den rene Fornufts Kritik“ forvandler Ideen til et subjectivt
Ubetinget; den dogmatiske Positivisme hypostaserer Ideen til et for sig
bestaaende Urvæsen. Et saadant positivt

% [text to be added: pp. 328--341]

% [text to be added: pp. 342--355]

% [text to be added: pp. 356--369]

% [text to be added: pp. 370--383]

% [text to be added: pp. 384--397]

% [text to be added: pp. 398--411]

% [text to be added: pp. 412--422]
\chapter*{c) Sjæl og Legeme.}
\addcontentsline{toc}{chapter}{c) Sjæl og Legeme}

% [text to be added: pp. 423--436]

% [text to be added: pp. 437--450]

% [text to be added: pp. 451--464]

% [text to be added: pp. 465--476]
%% ============================================================
%% III. PHILOSOPHIEN OG DE AKADEMISKE STUDIER  (pp. 477--482)
%% Printed under a SECOND "III." and absent from the Indhold (p. 6),
%% whose coverage stops at 477. See §2(c) of the header.
%% ============================================================

\part*{III.\\ Philosophien og de akademiske Studier.}
\addcontentsline{toc}{part}{III. Philosophien og de akademiske Studier}
\markboth{Philosophien og de akademiske Studier}{Philosophien og de akademiske Studier}

% [text to be added: pp. 477--482]

\end{document}
